\documentclass[11pt,oneside]{article}
\usepackage[margin=1in]{geometry}
\usepackage{amsmath,amssymb,amsthm,mathtools}
\usepackage{booktabs,longtable,array}
\usepackage{enumitem}
\usepackage[colorlinks=true,linkcolor=blue,urlcolor=blue,citecolor=blue]{hyperref}

\newtheorem{theorem}{Theorem}[section]
\newtheorem{lemma}[theorem]{Lemma}
\newtheorem{proposition}[theorem]{Proposition}
\newtheorem{corollary}[theorem]{Corollary}
\newtheorem{assessment}[theorem]{Assessment}
\newtheorem{firewall}[theorem]{Claim firewall}
\newtheorem{openproblem}[theorem]{Open problem}
\newtheorem{record}[theorem]{Record}
\newtheorem{programme}[theorem]{Programme}
\newtheorem{audit}[theorem]{Audit}
\newtheorem{decision}[theorem]{Decision}
\newtheorem{countertheorem}[theorem]{Countertheorem}
\theoremstyle{definition}
\newtheorem{definition}[theorem]{Definition}
\theoremstyle{remark}
\newtheorem{remark}[theorem]{Remark}

\newcommand{\Tr}{\operatorname{Tr}}
\newcommand{\Ran}{\operatorname{Ran}}
\newcommand{\MGclosed}{\textsc{closed}}
\newcommand{\MGopen}{\textsc{open}}
\newcommand{\MGpartial}{\textsc{partial}}

\title{Renewal Yang--Mills Mass-Gap Research Notes\\
Post-Tenth-Archive Compact Lineage, Version V100 (Archive f11)}
\author{Proof-indexed continuation after archive \texttt{V100\_f10}}
\date{11 September 2026}

\begin{document}
\maketitle
\tableofcontents

\section{Context, archive, and present objective}
\label{sec:v1-context}

The objective of this programme is a rigorous construction of
four-dimensional quantum Yang--Mills theory for a compact nonabelian
gauge group and a proof that the physical Hamiltonian has a strictly
positive spectral gap above its vacuum.  The proposed route begins
with exact gauge-compatible finite regulators, constructs a
volume-uniform renormalization flow and a cofinal positive observable
state, removes all auxiliary regulators, performs
Osterwalder--Schrader reconstruction, and proves quantitative
Euclidean clustering strong enough to imply a Hamiltonian mass gap.

That objective has not been reached.  The archives contain rigorous
finite-stage identities, estimates, conditional closure theorems,
and no-go results.  This compact restart records only the proved
interfaces needed by the live argument.  Its new work is to place the
weak-coupling bounds on the physical cofinal scale clock and to
separate geometrically stable transverse errors from the genuinely
marginal scalar Ward obstruction.  No step below promotes a
finite-cutoff estimate to a continuum theory.

\begin{record}[Frozen predecessor]
\label{rec:v1-frozen}
The complete preceding active lineage is frozen as
\[
\texttt{Renewal\_Yang\_Mills\_Mass\_Gap\_Research\_Notes\_V100\_f10.tex}.
\]
Its validated record is
\[
\begin{array}{c|c}
\text{lines}&21498\\
\text{bytes}&787493\\
\text{SHA--256}&
\texttt{277621520204F943E357AD1E3B971580F17251E750C5F85C7758D49BD10A47A3}.
\end{array}                                                   \tag{1.1}
\]
The files with suffixes \(\texttt{\_f1},\ldots,\texttt{\_f9}\)
remain earlier proof archives.  Successors of the present V1 append
to the complete active source and replace the immediate active
predecessor only after structural validation and compilation.  Every
fifth active version audits the current Standard--Model--Einstein and
Navier--Stokes notes.  At V100 the active source is frozen with the
next unused archive suffix and another contextual V1 begins.
\end{record}

\begin{firewall}[No mass-gap claim]
\label{fire:v1-no-gap}
There is not yet a constructed four-dimensional continuum
Yang--Mills measure, an Osterwalder--Schrader Hilbert space, or a
proved positive Hamiltonian gap.  The exact fibre and torus
calculations recorded below do not replace regulator removal,
reflection positivity, clustering, or the spectral argument.
\end{firewall}

\section{Corrected inherited proof ledger}
\label{sec:v1-ledger}

This section is a status map, not a reproving of archived results.
Every item retains the hypotheses of its source.

\begin{record}[Degree-faithful analytic forests]
\label{rec:v1-forest}
The ninth archive already corrected the degree-free tree estimate
which was mistakenly listed as open in active V99.  For labelled
trees on \([n]\), the unavoidable analytic vertex degrees are summed
by the exact Pruefer identity
\[
 \sum_{T\in\mathcal T_n}\prod_{i=1}^{n}d_T(i)!
   =(n-2)!\binom{3n-3}{n-2}.                         \tag{2.1}
\]
After a single Taylor--Wiener radius loss, the resulting connected
series converges under the declared one-scale condition
\[
 8zK_{\mathrm{HK}}<1,
 \qquad
 K_{\mathrm{HK}}=\frac{8h_*G_\mu}{(R-r)^2}.          \tag{2.2}
\]
Thus a degree-free fixed-tree inequality is not a legitimate live
target.  V100 formally superseded the V99 dependency rows which
treated it as one.

V100 also allowed scale-dependent regulator constants.  If
\[
 z_\beta\le C_{\mathrm{an}}
      \bigl(e^{c_*\beta^{-1/5}}-1\bigr),
 \qquad
 K_{\mathrm{HK}}\le C_K\beta^p,
 \qquad c_*=7680\pi^3,                               \tag{2.3}
\]
then (2.2) is eventually uniform when \(p<1/5\).  At the critical
exponent \(p=1/5\), the sufficient smallness condition recorded in
V100 is \(8eC_{\mathrm{an}}c_*C_K<1\).  Fixed regulator constants
are the special case \(p=0\).
\end{record}

\begin{record}[Exact reduced-fibre and all-torus estimates]
\label{rec:v1-torus}
On the declared finite Wilson regulator, the archive constructs the
reduced parity-tree fibre, stationary conditional centres, and the
complete reduced transverse momentum torus.  V97 proves, for both
signs and every active momentum,
\[
 \frac43 H_W\pm\nabla_1H_W\succeq\theta_{97}I,
 \qquad
 \theta_{97}=
 \frac{23784808996629854813771747874389}
 {241558005344776833024000000000000000000}.           \tag{2.4}
\]
Together with \(H_W\succeq(49/1250)I\), V100 obtains, with
\(\overline M_W\) the explicit finite Laurent-coefficient ceiling,
\[
 \left\|H_W^{-1/2}(\nabla_1H_W)H_W^{-1/2}\right\|
 \le \kappa_W
 :=\frac43-\frac{\theta_{97}}{\overline M_W}<\frac43, \tag{2.5}
\]
and the identical normalized estimate for the differentiated inverse
covariance.  These are exact finite-dimensional statements; no
volume or continuum conclusion is implicit in them.
\end{record}

\begin{record}[The actual marginal obstruction]
\label{rec:v1-edge}
The ninth archive reduces the marginal response problem successively
to the anomaly-decoupling mean condition, the scale-cohomology
interface condition, and finally one constant Fourier mode.  In the
notation of that archive, the edge response has the form
\[
 f_{j,t}=m_{j,t}+(I-P)a_{j,t},                         \tag{2.6}
\]
where \(P\) is the dyadic alias transfer.  The analytic nonconstant
part is a coboundary.  The remaining scale-cohomology condition is
equivalent to vanishing of the appropriate Cesaro average of the
scalar means \(m_{j,t}\).  Under the archived \(L^1\)-in-\(t\)
convergence hypothesis, it is equivalent to vanishing of the
integral of the limiting mean.  An exponentially accurate finite-grid
mean certificate is also proved there.

Consequently the live marginal row is not tree enumeration.  It is
the evaluation, or a cancellation certificate, for the
\emph{complete Ward-recombined scalar mean}, together with the
propagated reset and sector boundary needed by the anomaly-decoupling
mean condition.
\end{record}

\begin{audit}[V100 cross-programme audit]
\label{aud:v1-cross}
The V100 audit used the SM--Einstein active source \(\texttt{V96}\)
(777314 bytes) and the Navier--Stokes active source \(\texttt{V26}\)
(278283 bytes).  Their exact fingerprints were, respectively,
\begin{center}
\scriptsize
\texttt{91A272E2A240AE6B5D339A08017418FCFE5A705F47AC74D3444270C4FBDE4669},
\par\medskip
\texttt{82C887F8C2C69E4F28FAD7C3DC8DE5B9099CB5A4BF431EBA1FFF3E91935978AD}.
\end{center}
The useful SM lesson is structural: neither a chart-local floor nor a
Dirichlet floor gives a global inverse.  Thus the complete periodic
Ward scalar must remain in one global trace/dual norm.  No SM constants
transfer to Yang--Mills.  The NS rank-one
Oseen bounds do not control the compact Haar and determinant anomaly
terms and are not imported.
\end{audit}

\section{The cofinal continuum scale clock}
\label{sec:v1-clock}

The forest estimates are expressed in inverse coupling \(\beta\),
whereas regulator removal is expressed in lattice spacing.  The next
elementary result fixes the conversion and prevents a logarithmic
weak-coupling gain from being mistaken for a physical power law.

\begin{definition}[Geometric cofinal clock]
\label{def:v1-clock}
Fix \(a_0>0\) and a block factor \(B>1\).  Put
\[
 a_j=a_0B^{-j},\qquad j\ge1.                           \tag{3.1}
\]
Let \(\beta_j>0\) obey
\[
 \frac{\beta_j}{j}\longrightarrow b_{\mathrm{YM}}>0, \tag{3.2}
\]
and define the bare scale coupling by \(g_j^2=\beta_j^{-1}\).
\end{definition}

\begin{theorem}[Exact logarithmic clock conversion]
\label{thm:v1-clock}
Under Definition~\ref{def:v1-clock},
\[
 \frac{\beta_j}{\log(a_0/a_j)}
   \longrightarrow\frac{b_{\mathrm{YM}}}{\log B},
 \qquad
 g_j^2\log(a_0/a_j)
   \longrightarrow\frac{\log B}{b_{\mathrm{YM}}}.    \tag{3.3}
\]
If \(p>0\) and a nonnegative function \(w\) satisfies
\(w(\beta)\le C\beta^{-p}\) along this sequence, then
\[
 \limsup_{j\to\infty}
 [\log(a_0/a_j)]^p w(\beta_j)
 \le C\left(\frac{\log B}{b_{\mathrm{YM}}}\right)^p.\tag{3.4}
\]
\end{theorem}

\begin{proof}
Equation (3.1) gives the exact identity
\[
 \log(a_0/a_j)=j\log B.                               \tag{3.5}
\]
Dividing (3.2) by \(\log B\) proves the first limit in
(3.3).  Positivity of \(\beta_j\) and of its limiting ratio permits
inversion for all sufficiently large \(j\), proving the second.
Finally,
\[
 [\log(a_0/a_j)]^p w(\beta_j)
 \le C\left(\frac{j\log B}{\beta_j}\right)^p,
\]
and the right-hand side converges to the constant in (3.4).
\end{proof}

\begin{countertheorem}[A logarithmic estimate is not a mass scale]
\label{counter:v1-log-not-power}
The hypothesis of Theorem~\ref{thm:v1-clock} alone does not imply
\(w(\beta_j)=O(a_j^q)\) for any \(q>0\).
\end{countertheorem}

\begin{proof}
Take \(w(\beta)=\beta^{-p}\).  Then (3.2) gives
\[
 \frac{w(\beta_j)}{a_j^q}
 =a_0^{-q}\beta_j^{-p}B^{qj}
 \sim a_0^{-q}b_{\mathrm{YM}}^{-p}j^{-p}B^{qj}
 \longrightarrow\infty.                              \tag{3.6}
\]
Thus even the model sequence saturating the assumed weak-coupling
bound is not bounded by a positive power of the lattice spacing.
\end{proof}

\begin{remark}[Forest exponent]
\label{rem:v1-forest-clock}
In particular, a proved good-polymer estimate
\(u_{\mathrm{good}}(\beta)\le C\beta^{-1/5}\) becomes
only
\[
 u_{\mathrm{good}}(\beta_j)
   =O\bigl([\log(a_0/a_j)]^{-1/5}\bigr).               \tag{3.7}
\]
This is useful smallness for a stable recursion, but by
Countertheorem~\ref{counter:v1-log-not-power} it is not itself
exponential clustering or a mass gap.
\end{remark}

\section{Stable transverse tails under nonsummable forcing}
\label{sec:v1-stable-tail}

Weak-coupling errors of order \(j^{-p}\) need not be summable.
Uniform contraction nevertheless prevents their accumulation.  The
following result records the exact asymptotic ceiling needed to keep
that fact separate from the marginal scalar mode.

\begin{theorem}[Geometrically forced stable-tail theorem]
\label{thm:v1-stable-tail}
Let \(0\le\kappa<1\), \(C\ge0\), and \(p>0\).  Suppose
\(\beta_j>0\), \(\beta_j/j\to b>0\), and a nonnegative sequence
\((x_j)_{j\ge j_0}\) satisfies
\[
 x_{j+1}\le\kappa x_j+C\beta_j^{-p}
 \qquad(j\ge j_0).                                    \tag{4.1}
\]
Then
\[
 \limsup_{j\to\infty}j^p x_j
 \le\frac{Cb^{-p}}{1-\kappa}.                         \tag{4.2}
\]
For the geometric clock (3.1), this implies
\[
 \limsup_{j\to\infty}
 [\log(a_0/a_j)]^p x_j
 \le\frac{C}{1-\kappa}
       \left(\frac{\log B}{b}\right)^p.               \tag{4.3}
\]
\end{theorem}

\begin{proof}
Iteration of (4.1) gives
\[
 x_j\le \kappa^{j-j_0}x_{j_0}
   +C\sum_{m=j_0}^{j-1}\kappa^{j-1-m}\beta_m^{-p}.    \tag{4.4}
\]
The first term, multiplied by \(j^p\), tends to zero.  In the sum
put \(\ell=j-1-m\).  For each fixed \(\ell\),
\[
 j^p\beta_{j-1-\ell}^{-p}\longrightarrow b^{-p}.      \tag{4.5}
\]
It remains to justify passage through the growing geometric sum.
Choose \(M\) so that \(\beta_m\ge(b/2)m\) for \(m\ge M\).
When \(0\le\ell\le j/2\) and \(j\) is large enough, one has
\(m=j-1-\ell\ge M\) and
\[
 j^p\beta_m^{-p}
 \le (2/b)^p(j/m)^p\le(4/b)^p.                        \tag{4.6}
\]
Hence dominated convergence against \(\sum_{\ell\ge0}\kappa^\ell\)
applies to this half of the sum.  Because \(\beta_m/m\to b>0\),
the positive sequence \((\beta_m^{-p})_{m\ge j_0}\) is bounded;
write its supremum as \(D\).  The other half satisfies
\[
 j^p\sum_{\ell>j/2}\kappa^\ell
       \beta_{j-1-\ell}^{-p}
 \le \frac{D j^p\kappa^{j/2}}{1-\kappa}\longrightarrow0, \tag{4.7}
\]
with the evident direct interpretation when \(\kappa=0\).
Therefore the sum in (4.4), after multiplication by \(j^p\), has
limit \(b^{-p}/(1-\kappa)\).  Taking a limsup proves (4.2), and
(3.5) gives (4.3).
\end{proof}

\begin{corollary}[Stable pieces disappear from the Cesaro edge test]
\label{cor:v1-cesaro}
Let \((\mathcal I,\nu)\) be a probability space.  Suppose a
Banach-valued response splits as
\[
 F_{j,t}=m_{j,t}e+r_{j,t},\qquad
 \sup_{t\in\mathcal I}\|r_{j,t}\|\le x_j,             \tag{4.8}
\]
where \(e\) spans the marginal scalar direction and \(x_j\) obeys
the hypotheses of Theorem~\ref{thm:v1-stable-tail}.  Then
\[
 \left\|\frac1N\sum_{j=j_0}^{N-1}
       \int_{\mathcal I}r_{j,t}\,d\nu(t)\right\|
 \longrightarrow0.                                   \tag{4.9}
\]
More precisely, the left side is
\[
 \begin{cases}
 O(N^{-p}),&0<p<1,\\
 O((\log N)/N),&p=1,\\
 O(N^{-1}),&p>1.
 \end{cases}                                           \tag{4.10}
\]
Thus, once an actual uniform contraction \(\kappa<1\) is proved for
the transverse block, its polynomial weak-coupling forcing cannot
create a new Cesaro obstruction; only the scalar means
\(m_{j,t}\) remain.
\end{corollary}

\begin{proof}
Theorem~\ref{thm:v1-stable-tail} gives \(x_j=O(j^{-p})\).
The triangle inequality and (4.8) bound the left side of (4.9) by
\(N^{-1}\sum_{j=j_0}^{N-1}x_j\).  Integral comparison for
\(\sum j^{-p}\) gives (4.10), hence (4.9).
\end{proof}

\begin{firewall}[The V100 response bound is not yet a contraction]
\label{fire:v1-not-contraction}
The strict estimate \(\kappa_W<4/3\) in (2.5) does not imply
\(\kappa_W<1\).  It therefore does not populate the hypothesis
\(\kappa<1\) of Theorem~\ref{thm:v1-stable-tail}.  The theorem is a
proved interface: it shows exactly what follows once the full reset
map supplies a stable transverse block.  Establishing that block,
with the correct global norm and sector boundary, remains open.
\end{firewall}

\section{Post-archive dependency decision}
\label{sec:v1-decision}

\begin{proposition}[What the new estimates remove]
\label{prop:v1-remove}
Assume that the complete scale response admits the splitting (4.8),
that its transverse component satisfies (4.1), and that the scalar
component is the constant alias mode in Record~\ref{rec:v1-edge}.
Then all transverse forcing of size \(O(\beta_j^{-p})\), including
the exponent \(p=1/5\) supplied by the good-polymer estimate, vanishes
in the Cesaro scale average.  The scale-cohomology condition is then
equivalent to the archived scalar mean condition alone.
\end{proposition}

\begin{proof}
Corollary~\ref{cor:v1-cesaro} removes the transverse term from the
Cesaro average.  The remaining term is \(m_{j,t}e\), and
Record~\ref{rec:v1-edge} identifies its vanishing Cesaro mean as the
necessary and sufficient scalar condition in the archived analytic
alias class.
\end{proof}

\begin{decision}[V2 target]
\label{dec:v1-v2}
The next note will not repeat the alias-cohomology proof.  It will
assemble the actual finite-regulator scalar \(m_{j,t}\) before taking
absolute values.  The assembly must retain, in one global periodic
trace/dual pairing,
\begin{enumerate}[label=\textup{(\roman*)},leftmargin=3.3em]
\item the covariance derivative controlled by the V97--V100 torus
pencil;
\item the score and determinant derivatives, Haar/Jacobian terms,
toron and topological-sector pieces, and the differentiated connected
cumulants required by the Ward ledger;
\item the owner/reset and sector-boundary contributions propagated
between adjacent scales; and
\item the exact constant-mode extraction used by the archived
finite-grid certificate.
\end{enumerate}
The first admissible closure is an exact cancellation of this complete
mean.  The alternatives remain the archived vanishing Cesaro mean or
a decomposition
\[
 e_j^{\mathrm{edge}}=d_{j+1}-d_j+\varepsilon_j,
 \qquad d_N=o(N),
 \qquad \frac1N\sum_{j<N}\varepsilon_j\longrightarrow0. \tag{5.1}
\]
No chart-local inverse and no separate absolute-value estimate of
terms expected to cancel will be substituted for this calculation.
\end{decision}

\begin{openproblem}[Nearest unpaid row]
\label{open:v1-row}
Derive an explicit, regulator-correct formula for the complete scalar
mean \(m_{j,t}\), and prove one of the three closure alternatives in
Decision~\ref{dec:v1-v2}.  Simultaneously prove that the complementary
reset map has a uniform contraction factor below one so that
Theorem~\ref{thm:v1-stable-tail} applies.  Only after these marginal
and stable interfaces are joined can the programme return to
cofinal-state construction, reflection positivity, quantitative
clustering, and the Hamiltonian spectral gap.
\end{openproblem}

% =============================================================================
% BEGIN APPENDED POST-TENTH-ARCHIVE V2 MATERIAL
% =============================================================================

\section{V2: endpoint reduction and a certified negative quarter-turn card}
\label{sec:v2-endpoint-card}

\subsection{Upstream reconciliation of the scalar ledger}

Decision~\ref{dec:v1-v2} deliberately required every measured and
stationary row to be accounted for before the scalar mean was used.
The older archive contains the needed accounting, but its conclusion
must be stated at the correct scope.

\begin{proposition}[Commuting-toron endpoint completeness]
\label{prop:v2-complete-endpoint}
On the one-colour commuting-toron background used to extract the
quadratic gauge response, the following statements hold for both the
Wilson action and the positive comparison endpoint.
\begin{enumerate}[label=\textup{(\roman*)},leftmargin=3.3em]
\item The translated stationary centre is identically zero, including
its first two background derivatives.
\item Product Haar measure is invariant under the translating chart,
so its logarithmic Jacobian has zero first and second background
derivatives.
\item The toron lift and harmonic router agree at the two endpoints.
\item The direct cubic and quartic banks therefore exhaust the
measured stationary endpoint vertices.
\end{enumerate}
Consequently the order-two response
\[
 F_e(k):={\cal R}^{[2]}_e(k),\qquad e\in\{{\rm W},*\},         \tag{6.1}
\]
formed by the V60 Laurent oracle is the complete Gaussian endpoint
integrand in its declared normalization.  No extra Haar, stationary,
or router term may be added to (6.1).
\end{proposition}

\begin{proof}
The exact-flatness theorem in archived compact V26 shows that the
lifted one-colour toron remains in a commuting flat valley for both
actions.  In translated fibre coordinates the vertical first
variation vanishes at zero for every background parameter.  The
positive vertical Hessian gives local uniqueness of the stationary
branch, hence the branch is identically zero and so are its first two
derivatives.  Translation is factorwise left multiplication on the
compact group, and product Haar invariance makes its Jacobian one.
The archived common-lift theorem identifies the two toron routers.
These facts prove that the complete vertices are precisely the direct
sums displayed in archived equation (26.11).

Compact V58 supplies both quartic Laurent banks, compact V59 supplies
both orientations of both cubic Laurent banks, and compact V60 joins
them to the two Hessians, their inverse jets, and every tadpole and
bubble contraction.  Its exact sixteen-corner regression and inverse
identities verify the joined interface.  Thus (6.1) is complete at the
stated Gaussian endpoint scope.
\end{proof}

\begin{firewall}[Endpoint completeness is not SCIC completeness]
\label{fire:v2-endpoint-scope}
Proposition~\ref{prop:v2-complete-endpoint} concerns the Gaussian
commuting-toron endpoint integrand.  The adjacent-scale rechart,
owner/reset boundary, topological-sector mixture, and nonlinear
interacting remainder belong to the later scale comparison.  They are
not missing summands of (6.1), but they remain unpaid rows of SCIC,
ADMC, and the full continuum construction.
\end{firewall}

\subsection{The homotopy scalar is an endpoint difference}

For a periodic integrable scalar \(f\), write
\[
 \langle f\rangle:=\frac1{(2\pi)^4}
   \int_{[0,2\pi]^4}f(k)\,dk.                                \tag{6.2}
\]

\begin{theorem}[Endpoint reduction of the constant edge class]
\label{thm:v2-endpoint-reduction}
Let the map $t\mapsto W_t$ be a $C^1$ finite-regulator response homotopy
whose derivative is jointly integrable in \((t,k)\).  Then
\[
 \int_0^1\langle\partial_tW_t\rangle\,dt
   =\langle W_1-W_0\rangle.                                  \tag{6.3}
\]
The identity is unchanged if one subtracts periodic divergences,
moment-null contacts, or alias coboundaries, because each has zero
constant class under the archived hypotheses.

If, after the physical edge rescaling, the two endpoint banks are
independent of the scale label \(j\), their Gaussian edge scalar is a
single number
\[
 \mathfrak m_{\rm G}:=\langle F_*-F_{\rm W}\rangle.           \tag{6.4}
\]
Its Cesaro average equals \(\mathfrak m_{\rm G}\) at every length.
In particular a scale-stationary nonzero endpoint mismatch cannot be
converted into a sublinear scale coboundary.
\end{theorem}

\begin{proof}
The fundamental theorem of calculus gives
\(W_1(k)-W_0(k)=\int_0^1\partial_tW_t(k)\,dt\) for almost every
\(k\).  Fubini's theorem and (6.2) give (6.3).  Periodic divergences
integrate to zero; the archived contact and alias-coboundary theorems
give zero constant class for the other two listed terms.

Scale independence makes every scalar in the Cesaro sum equal to
\(\mathfrak m_{\rm G}\), hence
\[
 \frac1N\sum_{j=0}^{N-1}\mathfrak m_{\rm G}
 =\mathfrak m_{\rm G}.                                      \tag{6.5}
\]
If instead
\(\mathfrak m_{\rm G}=d_{j+1}-d_j+\varepsilon_j\), with
\(d_N=o(N)\) and
\(N^{-1}\sum_{j<N}\varepsilon_j\to0\), sum this identity from
zero to \(N-1\) and divide by \(N\).  The right side tends to zero,
whereas the left side is \(\mathfrak m_{\rm G}\).  Therefore such a
sublinear representation requires \(\mathfrak m_{\rm G}=0\).
\end{proof}

\begin{remark}[Normalization]
\label{rem:v2-normalization}
Equation (6.4) uses the action-normalized V60 convention.  Restoring
the positive overall factor \(3/2\) in the V16 Ward density changes
neither vanishing nor sign.  It is nevertheless to be restored when a
numerical coefficient is compared with a continuum beta-function
normalization.
\end{remark}

\subsection{A floor-free residual theorem for complete trace words}

The archived V61 value of the \(M=4\) mean used floating inverses.  A
direct replacement by the global Hessian floor is rigorous but loses
many orders of magnitude.  The next theorem retains the trial inverse
inside every complete trace factor.

\begin{theorem}[Multiplicative residual enclosure of an arbitrary
trace word]
\label{thm:v2-word-residual}
Let $H=H^*>0$, let $G=H^{-1}$, and let $R=R^*$.  Put
\[
 E:=I-RH,
 \qquad \|E\|_{\rm F}\le\eta<1,
 \qquad S:=(I-E)^{-1}.                                      \tag{6.6}
\]
For arbitrary square matrices \(X_1,\ldots,X_m\), set
\(A_i=RX_i\).  Then \(G=SR\) and, for every \(m\ge1\),
\[
\boxed{
 \left|
 \Tr(GX_1\cdots GX_m)-\Tr(A_1\cdots A_m)
 \right|
 \le
 \left[(1-\eta)^{-m}-1\right]
 \prod_{i=1}^m\|A_i\|_{\rm F}.}                              \tag{6.7}
\]
No lower spectral bound for \(H\) occurs in (6.7).
\end{theorem}

\begin{proof}
From \(RH=I-E\), right multiplication by \(G\) gives
\(R=(I-E)G\), hence \(G=SR\).  The Neumann series gives
\[
 \|S\|_{\rm op}\le s:=(1-\eta)^{-1},
 \qquad
 \|S-I\|_{\rm op}\le s-1.                                 \tag{6.8}
\]
For a one-factor word, the sharper Frobenius estimate
\[
 \|S-I\|_{\rm F}=\|SE\|_{\rm F}
 \le s\eta=s-1                                             \tag{6.9}
\]
and Hilbert--Schmidt duality prove (6.7).

For \(m\ge2\), write the actual word as
\(\Tr(SA_1SA_2\cdots SA_m)\) and replace the \(m\) copies of \(S\)
by \(I\), one at a time.  In the term generated at the \(r\)-th
replacement, use operator norm on all \(S\) and \(S-I\) factors,
Hilbert--Schmidt norm on two \(A_i\)'s, and operator norm bounded by
Frobenius norm on the remaining \(A_i\)'s.  Schatten--H\"older then
bounds that term by
\[
 (s-1)s^r\prod_{i=1}^m\|A_i\|_{\rm F}
\]
for one of \(r=0,\ldots,m-1\).  Summing the geometric series gives
\[
 (s-1)\sum_{r=0}^{m-1}s^r=s^m-1,
\]
which is (6.7).
\end{proof}

For a fixed action and momentum, abbreviate
\[
 A_X:=RX,\qquad a_X:=\|A_X\|_{\rm F},\qquad
 \rho_m(\eta):=(1-\eta)^{-m}-1.                             \tag{6.10}
\]
The response identity (60.11), after inserting
\[
 G_1=-GH_1G,\qquad
 G_2=2GH_1GH_1G-GH_2G,                                     \tag{6.11}
\]
has the following trial centre:
\[
\begin{split}
 \Phi_R={}&\frac12\Tr(A_{Q_2})+c_e\bigl\{
 \Tr(A_{P_2}A_{M_0})
 +2\Tr(A_{P_0}A_{H_1}A_{H_1}A_{M_0})\\
 &-\Tr(A_{P_0}A_{H_2}A_{M_0})
 +\Tr(A_{P_0}A_{M_2})
 -2\Tr(A_{P_1}A_{H_1}A_{M_0})\\
 &+2\Tr(A_{P_1}A_{M_1})
 -2\Tr(A_{P_0}A_{H_1}A_{M_1})\bigr\},
\end{split}                                                   \tag{6.12}
\]
where $c_{\rm W}=1/3$ and $c_*=3$.

\begin{corollary}[Explicit response radius]
\label{cor:v2-response-radius}
With (6.10)--(6.12), the exact response obeys
\[
 |F_e(k)-\Phi_R(k)|\le\varepsilon_R(k),                      \tag{6.13}
\]
where
\[
\begin{split}
 \varepsilon_R={}&
 \frac12a_{Q_2}\rho_1
 +|c_e|\bigl\{
 a_{P_2}a_{M_0}\rho_2
 +2a_{P_0}a_{H_1}^2a_{M_0}\rho_4\\
 &+a_{P_0}a_{H_2}a_{M_0}\rho_3
 +a_{P_0}a_{M_2}\rho_2
 +2a_{P_1}a_{H_1}a_{M_0}\rho_3\\
 &+2a_{P_1}a_{M_1}\rho_2
 +2a_{P_0}a_{H_1}a_{M_1}\rho_3\bigr\}.
\end{split}                                                   \tag{6.14}
\]
Here every \(\rho_m\) is evaluated at the same residual bound
\(\eta\).
\end{corollary}

\begin{proof}
Expand the two covariance derivatives by (6.11).  This gives exactly
the seven bubble words in (6.12), with respectively
\(2,4,3,2,3,2,3\) inverse slots.  Apply
Theorem~\ref{thm:v2-word-residual} to each word and the triangle
inequality only after the complete trial trace has been formed.  The
one-factor case gives the tadpole term.  This is (6.14).
\end{proof}

\subsection{Exact quarter-turn execution}

Define the full quarter-turn grid and its normalized mean by
\[
 \Gamma_4:=
 \left\{\frac\pi2n:n\in\{0,1,2,3\}^4\right\},
 \qquad
 \langle f\rangle_4:=\frac1{256}\sum_{k\in\Gamma_4}f(k).    \tag{6.15}
\]
Conjugate reflection pairs \(k\) and \(-k\), leaving \(136\)
representatives when the sixteen fixed points are counted once.

\begin{theorem}[Certified negative \(M=4\) endpoint-difference card]
\label{thm:v2-negative-card}
At each of the \(136\) reflection representatives, form the exact
Hermitian Gaussian-rational trial inverse of denominator \(10^8\)
used by the V92 quarter-turn certificate.  Evaluate (6.12) with exact
integer matrix products and average (6.14) with the reflection
weights.  The maximal residual bounds are
\[
 \eta_{\rm W}\le\frac{111971}{250000000000},
 \qquad
 \eta_*\le\frac{7926629}{500000000000}.                      \tag{6.16}
\]
The exact trial centres are
\[
 c_{\rm W}=
 \frac{76967350158114568540832543838838898581}
 {345600000000000000000000000000000000},                    \tag{6.17}
\]
\[
 c_*=-
 \frac{3063277481200150123558963791819240686231}
 {33246720000000000000000000000000000000}.                  \tag{6.18}
\]
Their mean error radii obey
\[
 r_{\rm W}<0.443356940,\qquad
 r_*<18.437658713.                                           \tag{6.19}
\]
Consequently
\[
 c_*-c_{\rm W}=
 -\frac{6542210354006732260741909068197214206077}
 {20779200000000000000000000000000000000}
 \approx-314.844188130762,                                  \tag{6.20}
\]
and the true finite-grid mean satisfies the completely rational
enclosure
\[
\boxed{
 -\frac{41715650473}{125000000}
 \le
 \langle F_*-F_{\rm W}\rangle_4
 \le
 -\frac{147981586239}{500000000}<0.}                         \tag{6.21}
\]
Numerically, the interval is
\[
 [-333.725203784,\,-295.963172478].                          \tag{6.22}
\]
\end{theorem}

\begin{proof}
All Laurent coefficients are rational, and every phase on
\(\Gamma_4\) belongs to \(\{1,i,-1,-i\}\).  Thus the Hessian jets,
vertices, trial inverses, residual products, and traces in (6.12) all
lie in \(\mathbb Q(i)\).  The verifier forms them by integer matrix
multiplication with explicit common denominators.  The response
reflection theorem identifies the contribution of a two-point orbit
with twice the real part at its representative.  No floating value is
used in the centre or radius.

The residual calculation gives (6.16).  Exact weighted summation gives
(6.17)--(6.20).  Corollary~\ref{cor:v2-response-radius} gives the two
radii in (6.19), whose exact rational values are retained by the
executable.  Adding the action radii, subtracting the trial centres,
and rounding outward to denominator \(10^9\) gives (6.21).  Its upper
endpoint is a negative rational, proving the sign.
\end{proof}

\begin{record}[V2 endpoint-scalar verifier]
\label{rec:v2-verifier}
The executable is
\[
 \texttt{scripts/verify\_ym\_postarchive\_v2\_endpoint\_scalar.py}.
                                                                  \tag{6.23}
\]
It has \(414\) lines, \(14310\) bytes, and SHA--256
\[
 \texttt{F6D75E19B3650EBCEDDA87F67A48B3EDE963CDA495EA5CC8B480BDE977FD0007}.
                                                                  \tag{6.24}
\]
The canonical output-payload SHA--256 is
\[
 \boxed{
 \texttt{BEB9DE24D1996A1B835EFE70CCB3238EFCC7C4F0A6D0F41D5B81DC32811BFBCB}.}
                                                                  \tag{6.25}
\]
The run checks all \(136\) reflection representatives, all residual-one
gates, the action-radius sum, strict negativity of (6.21), and inclusion
of the archived floating V61 value in the new rational interval.
\end{record}

\subsection{What this decides and what it does not}

\begin{firewall}[A finite-grid sign is not the Haar mean]
\label{fire:v2-grid-not-mean}
Theorem~\ref{thm:v2-negative-card} certifies the complete literal
V60 response on the \(M=4\) tensor grid.  It does not identify that
grid mean with the torus Haar mean (6.4).  Modes whose four Fourier
indices are multiples of four alias into (6.15), and their total has
not been bounded sharply enough.  Therefore V2 neither proves that
\(\mathfrak m_{\rm G}\ne0\) nor closes or refutes SCIC.
\end{firewall}

\begin{corollary}[Exact remaining quadrature threshold]
\label{cor:v2-threshold}
Let
\[
 \delta_4:=\frac{147981586239}{500000000}.                   \tag{6.26}
\]
Any rigorous estimate
\[
 \left|
 \langle F_*-F_{\rm W}\rangle
 -\langle F_*-F_{\rm W}\rangle_4
 \right|<\delta_4                                           \tag{6.27}
\]
would imply
\(\mathfrak m_{\rm G}<0\).  Conversely, if the full Ward/scale
comparison is to have zero scalar class while (6.27) holds, its
reset/sector boundary must supply a compensating non-Gaussian or
adjacent-scale constant class; it cannot be hidden in a periodic
divergence or sublinear stationary coboundary.
\end{corollary}

\begin{proof}
The upper endpoint in (6.21) is \(-\delta_4\).  Adding an error whose
absolute value is strictly less than \(\delta_4\) leaves the torus
mean strictly negative.  The second statement follows from
Theorem~\ref{thm:v2-endpoint-reduction} and the zero-mean classes in
Record~\ref{rec:v1-edge}.
\end{proof}

\begin{assessment}[Progress at V2]
\label{ass:v2-progress}
V2 removes two ambiguities.  First, the Gaussian endpoint scalar is
already a complete V60 object on the commuting-toron background;
Haar, stationary, and router terms are not absent from that card.
Second, its \(M=4\) value is no longer a floating diagnostic: a
floor-free multiplicative residual theorem and exact
Gaussian-rational execution prove the strict negative interval
(6.21).  The remaining Gaussian question is now the explicit alias
tail (6.27), not construction of a homotopy derivative bank and not
another evaluation of the same grid.
\end{assessment}

\begin{decision}[V3 target]
\label{dec:v2-v3}
Attack (6.27) without returning to the global entrywise inverse
recursion rejected in archived V62.  Keep the exact V60 trace sum and
use one of the following structure-preserving routes:
\begin{enumerate}[label=\textup{(\roman*)},leftmargin=3.3em]
\item a residual-retaining interval subdivision of the torus, with
complete \(RX\) products and local Fourier-tail control;
\item a coefficientwise constant-mode computation of the rational
matrix resolvent; or
\item a Ward-derived cancellation which removes the high alias modes
before any absolute value is taken.
\end{enumerate}
The target error is the exact rational \(\delta_4\) in (6.26).  If a
global bound again loses the finite-grid margin, go upstream to the
alias blocks or local Schur object rather than increasing the tensor
grid under an inoperative estimate.
\end{decision}

\begin{firewall}[Claim boundary after V2]
\label{fire:v2-final}
V2 proves an endpoint reduction, a floor-free trace-word residual
theorem, and a certified negative finite-grid response card.  It does
not certify the continuous torus mean, populate the compensating
scale boundary, prove SCIC or ADMC, construct the cofinal interacting
state or continuum measure, establish reflection positivity or
physical exponential clustering, reconstruct a Hamiltonian, or prove
the Yang--Mills mass gap.
\end{firewall}

% =============================================================================
% END APPENDED POST-TENTH-ARCHIVE V2 MATERIAL
% =============================================================================

% =============================================================================
% BEGIN APPENDED POST-TENTH-ARCHIVE V3 MATERIAL
% =============================================================================

\section{V3: cancellation-retaining cell transport and an exact jet obstruction}
\label{sec:v3-cell-transport}

\subsection{Context and the precise purpose of this iteration}

V2 reduced the Gaussian homotopy scalar to the endpoint difference
\(\langle F_*-F_{\rm W}\rangle\) and certified the strictly negative
quarter-turn-grid interval
\[
 -\frac{41715650473}{125000000}
 \le \langle F_*-F_{\rm W}\rangle_4
 \le -\frac{147981586239}{500000000}<0.                    \tag{7.1}
\]
The unresolved step is not another finite-grid evaluation.  It is a
rigorous passage from (7.1) to the normalized Haar integral on the
continuous momentum torus.  A direct global second-derivative
majorant was tested upstream and is too large by many orders of
magnitude.  This iteration therefore keeps the trial inverse and each
complete trace word intact while the momentum moves inside a cell.

Two further duties are discharged below.  First, the Laurent support
is combined with the trial inverse before norms are taken; this gives
a residual-weighted transport stock and a sharper first-order Taylor
stock.  Second, an exact coefficient comparison rules out the
factorwise identification of the external response jets with
momentum derivatives.  That comparison prevents a false periodic
integration-by-parts shortcut from being recycled.

\subsection{A local enclosure for a transported complete trace word}

Let \({\cal C}\) be a lifted torus cell with anchor \(a\).  All
matrices in the next theorem are \(n\times n\).  The notation is
deliberately local: a new trial inverse may be chosen independently
on every cell.

\begin{theorem}[Cancellation-retaining cell enclosure]
\label{thm:v3-cell-word}
Let \(H(k)=H(k)^*>0\) on \({\cal C}\), fix \(R_a=R_a^*\), and put
\[
 E_a(k):=I-R_aH(k),\qquad
 \sup_{k\in{\cal C}}\|E_a(k)\|_{\rm F}\le\eta<1,\qquad
 s:=(1-\eta)^{-1}.                                         \tag{7.2}
\]
For matrix functions \(X_i(k)\), define
\[
 A_i(k):=R_aX_i(k),\quad
 a_i:=\|A_i(a)\|_{\rm F},\quad
 d_i\ge\sup_{k\in{\cal C}}\|A_i(k)-A_i(a)\|_{\rm F},\quad
 b_i:=a_i+d_i.                                             \tag{7.3}
\]
For \(m\ge2\) one has, uniformly on \({\cal C}\),
\[
\boxed{\;
 \left|
 \Tr\!\left(H^{-1}X_1\cdots H^{-1}X_m\right)(k)
 -\Tr\!\left(A_1(a)\cdots A_m(a)\right)
 \right|
 \le s^m\prod_{i=1}^m b_i-\prod_{i=1}^m a_i .
 \;}                                                       \tag{7.4}
\]
For one factor, and any fixed \(\chi_n\ge\sqrt n\),
\[
\boxed{\;
 \left|\Tr(H^{-1}X)(k)-\Tr(A(a))\right|
 \le (s-1)b+\chi_n d .
 \;}                                                       \tag{7.5}
\]
In particular, neither estimate divides by a separately estimated
spectral floor for \(H\).
\end{theorem}

\begin{proof}
At a fixed \(k\), write \(E=E_a(k)\), \(S=(I-E)^{-1}\), and
\(A_i=A_i(k)\).  As in Theorem~\ref{thm:v2-word-residual},
\(H^{-1}=SR_a\), \(\|S\|_{\rm op}\le s\), and
\(\|S-I\|_{\rm F}\le s-1\).  Replacing the \(m\) copies of \(S\)
in
\[
 \Tr(SA_1SA_2\cdots SA_m)
\]
one at a time gives, for \(m\ge2\),
\[
 \left|\Tr(SA_1\cdots SA_m)-\Tr(A_1\cdots A_m)\right|
 \le (s^m-1)\prod_i\|A_i(k)\|_{\rm F}
 \le (s^m-1)\prod_i b_i.                                  \tag{7.6}
\]
Here two matrix factors are put in Hilbert--Schmidt norm and every
remaining matrix factor in operator norm, bounded by Frobenius norm.

Expand
\(\prod_i(A_i(a)+[A_i(k)-A_i(a)])-\prod_iA_i(a)\)
without commuting the factors.  Schatten--H\"older bounds every
nonconstant term by the corresponding product of \(a_i\)'s and
\(d_i\)'s.  Summing those terms yields
\[
 \left|\Tr(A_1(k)\cdots A_m(k))
       -\Tr(A_1(a)\cdots A_m(a))\right|
 \le\prod_i b_i-\prod_i a_i.                              \tag{7.7}
\]
Adding (7.6) and (7.7) proves (7.4).

For \(m=1\), Hilbert--Schmidt duality gives
\[
 |\Tr((S-I)A(k))|\le(s-1)b .
\]
Also
\[
 |\Tr(A(k)-A(a))|
 \le\|I\|_{\rm F}\|A(k)-A(a)\|_{\rm F}
 \le\sqrt n\,d\le\chi_n d .
\]
Their sum is (7.5).
\end{proof}

\begin{remark}[Where cancellation is retained]
\label{rem:v3-cancellation}
The centre of (7.4) is the signed trace of the complete word.
No entrywise absolute value and no absolute value of a partial trace
is introduced.  Absolute values enter only in the error radius after
the exact signed centre has been formed.
\end{remark}

\subsection{The complete V60 response on one cell}

For a fixed action \(e\in\{{\rm W},*\}\), let
\(\Phi_{e,a}\) be the signed trial response (6.12), with the single
matrix \(R_a\) and every factor evaluated at \(a\).  The seven bubble
words, including their signed multiplicities, are
\[
\begin{array}{c|c|c}
q&\gamma_q&(X_{q,1},\ldots,X_{q,m_q})\\ \hline
1& 1 &(P_2,M_0)\\
2& 2 &(P_0,H_1,H_1,M_0)\\
3&-1 &(P_0,H_2,M_0)\\
4& 1 &(P_0,M_2)\\
5&-2 &(P_1,H_1,M_0)\\
6& 2 &(P_1,M_1)\\
7&-2 &(P_0,H_1,M_1).
\end{array}                                                \tag{7.8}
\]
For every named factor \(X\), choose \(a_X,d_X,b_X\) as in (7.3).
Set
\[
 {\cal T}_1(X):=(s-1)b_X+\chi_n d_X,\qquad
 {\cal T}_{m_q}(q):=
 s^{m_q}\prod_{i=1}^{m_q}b_{X_{q,i}}
 -\prod_{i=1}^{m_q}a_{X_{q,i}}.                            \tag{7.9}
\]

\begin{corollary}[Cell radius for the complete endpoint response]
\label{cor:v3-cell-response}
Assume (7.2) for the action \(e\).  With
\(c_{\rm W}=1/3\), \(c_*=3\), define
\[
 {\cal R}_{e}({\cal C};a):=
 \frac12{\cal T}_1(Q_2)
 +|c_e|\sum_{q=1}^7|\gamma_q|{\cal T}_{m_q}(q).             \tag{7.10}
\]
Then
\[
\boxed{\quad
 |F_e(k)-\Phi_{e,a}|
 \le {\cal R}_{e}({\cal C};a)
 \quad(k\in{\cal C}).
\quad}                                                     \tag{7.11}
\]
At zero cell radius, (7.10) is exactly the point-residual radius
(6.14).
\end{corollary}

\begin{proof}
Apply (7.5) to the tadpole and (7.4) to the seven complete words
listed in (7.8), retaining their signed centres.  Only after those
centres have been combined is the triangle inequality applied to
the eight error radii.  If every \(d_X=0\), equations (7.9)--(7.10)
reduce to \(\rho_m(\eta)\prod_i a_i\) and hence to (6.14).
\end{proof}

\begin{corollary}[An exact sufficient condition for passing from
\(\Gamma_4\) to Haar measure]
\label{cor:v3-grid-to-haar}
Let \({\cal C}_a\), \(a\in\Gamma_4\), be the periodic Voronoi cubes
of the quarter-turn grid.  Each has normalized Haar weight \(1/256\)
and lifted \(\ell^\infty\)-radius \(\pi/4\).  Let
\(\varepsilon^0_{e,a}\) be the point radius used in the exact V2
certificate, and suppose (7.2)--(7.3) certify radii
\({\cal R}_{e,a}:={\cal R}_e({\cal C}_a;a)\).  If
\[
 \frac1{256}\sum_{a\in\Gamma_4}
 \left[
  ({\cal R}_{*,a}-\varepsilon^0_{*,a})
 +({\cal R}_{{\rm W},a}-\varepsilon^0_{{\rm W},a})
 \right]
 <
 \delta_4:=\frac{147981586239}{500000000},                 \tag{7.12}
\]
then
\[
 \langle F_*-F_{\rm W}\rangle<0.                           \tag{7.13}
\]
\end{corollary}

\begin{proof}
Integrate (7.11) on each cell and sum with weight \(1/256\).
The signed trial centres are the same ones used by the V2 exact
grid certificate.  Its upper endpoint is \(-\delta_4\); its point
radii are \(\varepsilon^0_{*,a}+\varepsilon^0_{{\rm W},a}\).
Replacing those radii by the cell radii increases the upper endpoint
by precisely the left side of (7.12).  The strict inequality leaves
that endpoint negative.
\end{proof}

\subsection{Residual-weighted Laurent and Taylor transport}

The preceding result becomes executable once \(\eta\) and the
factor variations are bounded on a cell.  The following estimates
use the finite Laurent support already constructed in the archive.
They improve the crude product
\(\|R_a\|_{\rm op}\sum_\alpha|\alpha|_1\|Y_\alpha\|_{\rm F}\)
by multiplying every coefficient by \(R_a\) before taking its norm.

\begin{theorem}[Residual-weighted transport stocks]
\label{thm:v3-weighted-stock}
Let
\[
 Y(k)=\sum_{\alpha\in{\cal A}}Y_\alpha e^{i\alpha\cdot k},
 \qquad {\cal A}\Subset\mathbb Z^4,
 \qquad |\alpha|_1:=\sum_{\mu=1}^4|\alpha_\mu|.              \tag{7.14}
\]
For a fixed matrix \(R_a\), define
\[
\begin{split}
 L^{(1)}_{R_aY}
   &:=\sum_{\alpha\in{\cal A}}|\alpha|_1
       \|R_aY_\alpha\|_{\rm F},\\
 \ell_{R_aY}(a)
   &:=\sum_{\mu=1}^4\|R_a\partial_\mu Y(a)\|_{\rm F},\\
 L^{(2)}_{R_aY}
   &:=\frac12\sum_{\alpha\in{\cal A}}|\alpha|_1^2
       \|R_aY_\alpha\|_{\rm F}.
\end{split}                                                 \tag{7.15}
\]
If \(k\) and \(a\) have lifts satisfying
\(\|k-a\|_\infty\le r\), then
\[
 \|R_a[Y(k)-Y(a)]\|_{\rm F}
 \le rL^{(1)}_{R_aY},                                      \tag{7.16}
\]
and the cancellation-retaining Taylor estimate
\[
 \|R_a[Y(k)-Y(a)]\|_{\rm F}
 \le r\ell_{R_aY}(a)+r^2L^{(2)}_{R_aY}.                    \tag{7.17}
\]
In particular, for
\(\eta_0(a):=\|I-R_aH(a)\|_{\rm F}\),
\begin{align}
 \|I-R_aH(k)\|_{\rm F}
 &\le\eta_0(a)+rL^{(1)}_{R_aH},                            \tag{7.18}\\
 \|I-R_aH(k)\|_{\rm F}
 &\le\eta_0(a)+r\ell_{R_aH}(a)+r^2L^{(2)}_{R_aH}.          \tag{7.19}
\end{align}
The same two bounds with \(H\) replaced by any response factor
\(X\) give admissible \(d_X\)'s for (7.3).
\end{theorem}

\begin{proof}
Put \(h=k-a\).  Since
\[
 |e^{i\alpha\cdot h}-1|
 \le|\alpha\cdot h|
 \le|\alpha|_1r,
\]
the triangle inequality applied after multiplying each Laurent
coefficient by \(R_a\) gives (7.16).  For the second estimate, use
\[
 |e^{ix}-1-ix|\le\frac{x^2}{2}
\]
and expand
\[
 R_a[Y(k)-Y(a)]
 =\sum_{\mu=1}^4h_\mu R_a\partial_\mu Y(a)
 +R_a{\cal E}_2(a,h).
\]
The linear term is bounded by \(r\ell_{R_aY}(a)\), while the
remainder is at most
\[
 \frac{r^2}{2}\sum_\alpha|\alpha|_1^2
       \|R_aY_\alpha\|_{\rm F}.
\]
This proves (7.17).  Finally,
\[
 I-R_aH(k)=I-R_aH(a)-R_a[H(k)-H(a)]
\]
gives (7.18)--(7.19).
\end{proof}

\begin{record}[Exact toron-anchor stock]
\label{rec:v3-toron-stock}
The exact verifier evaluated Theorem~\ref{thm:v3-weighted-stock} at
\[
 a=(0,0,0,0).                                               \tag{7.20}
\]
For Wilson, the point residual and trial-inverse operator bound are
\[
 \eta_0=0,\qquad \|R_a\|_{\rm op}\le\frac{471}{16}.         \tag{7.21}
\]
The unweighted Laurent stock, residual-weighted stock, and their
crude-product improvement ratio are
\[
\begin{split}
 \widehat L^{(1)}_H
 &\le\frac{3646905023083}{125000000000},\\
 L^{(1)}_{R_aH}
 &\le\frac{53709527030729}{500000000000},\\
 \frac{\|R_a\|_{\rm op}\widehat L^{(1)}_H}
      {L^{(1)}_{R_aH}}
 &=\frac{1717692265872093}{214838108122916}
   >7.995.
\end{split}                                                 \tag{7.22}
\]
Thus the first-order sufficient Neumann radius is
\[
 r<
 \frac{500000000000}{53709527030729}
 >0.0093093.                                                \tag{7.23}
\]
The cancellation-retaining stocks are
\[
 \ell_{R_aH}(a)\le\frac{21900262682627}{500000000000},
 \qquad
 L^{(2)}_{R_aH}\le\frac{15718119490079}{200000000000}.      \tag{7.24}
\]
At \(r=1/64\), (7.19) is at most
\[
 \frac{703570366413}{1000000000000}<1.                     \tag{7.25}
\]

For the endpoint action, the corresponding exact values are
\[
\begin{split}
 \eta_0&\le\frac{11205513}{1000000000000},\qquad
 \|R_a\|_{\rm op}\le\frac{151225919}{100000000},\\
 \widehat L^{(1)}_H
 &\le\frac{159654323446791}{125000000000},\qquad
 L^{(1)}_{R_aH}
 \le\frac{27340210537099}{125000000000},\\
 \frac{\|R_a\|_{\rm op}\widehat L^{(1)}_H}
      {L^{(1)}_{R_aH}}
 &=\frac{24143871785564216575929}
         {2734021053709900000000}
   >8.830.
\end{split}                                                 \tag{7.26}
\]
The first-order sufficient radius is
\[
 r<
 \frac{999988794487}{218721684296792}
 >0.0045719.                                                \tag{7.27}
\]
The Taylor stocks satisfy
\[
 \ell_{R_aH}(a)\le\frac{62926726824303}{1000000000000},
 \qquad
 L^{(2)}_{R_aH}\le\frac{11707924794243}{62500000000}.       \tag{7.28}
\]
At \(r=1/128\), (7.19) is at most
\[
 \frac{100611955827}{200000000000}<1.                      \tag{7.29}
\]
The nonconstant Hessian supports have respectively \(20\) and \(56\)
frequencies.  Every displayed comparison was performed over exact
rationals.
\end{record}

\begin{assessment}[What the one-anchor result does and does not buy]
\label{ass:v3-one-anchor}
The Taylor cancellations enlarge the certified dyadic radius from
the first-order thresholds to \(1/64\) and \(1/128\).  They do not
cover a quarter-turn Voronoi cube, whose radius is \(\pi/4\).
Indeed \(333/424<\pi/4\), while the particular rational majorant
obtained by inserting the certified coefficient stocks into the
right side of (7.19) at \(r=333/424\) has the values
\[
 \frac{10545561616293}{125000000000}>84
 \quad\hbox{and}\quad
 \frac{5368095933223}{31250000000}>171                    \tag{7.30}
\]
for Wilson and the endpoint action.  Thus this recorded
single-anchor majorant does not certify the whole cell.  The numbers
in (7.30) are not lower bounds on the actual residual and therefore
give no no-go theorem for local trial inverses, subdivision, grouped
phase intervals, or a Schur-level argument.
\end{assessment}

\subsection{Exact rejection of a factorwise derivative shortcut}

One might try to identify every external first or second response jet
with a periodic momentum derivative and then integrate it away.  Such
an identification would be useful only if it held for the actual
Laurent coefficient banks, before traces and cancellations are
formed.  The next theorem tests precisely that necessary coefficient
identity.

\begin{theorem}[The archived response jets are not factorwise
momentum derivatives]
\label{thm:v3-jet-obstruction}
Write the order-\(q\) external jet of one quartic or cubic response
factor in the archived first-axis convention as
\[
 X^{[q]}(k)=\sum_{\alpha\in\mathbb Z^4}
              X^{[q]}_\alpha e^{i\alpha\cdot k}.
\]
Factorwise first- and second-momentum differentiation would require
\[
 X^{[1]}_\alpha=\alpha_1X^{[0]}_\alpha,\qquad
 X^{[2]}_\alpha=-\alpha_1^2X^{[0]}_\alpha                 \tag{7.31}
\]
for every entry and every exponent.  Neither the Wilson nor endpoint
quartic bank satisfies (7.31), and none of the four signed cubic
banks satisfies it.  The exact numbers of failing
entry--exponent--order records are
\[
\begin{array}{c|rrr}
 &\text{quartic}&\text{cubic \(+\)}&\text{cubic \(-\)}\\ \hline
 {\rm Wilson}&270&664&664\\
 *&7153&10557&10603.
\end{array}                                                \tag{7.32}
\]
In particular, exact witnesses \((i,j;\alpha;q;
X^{[q]}_\alpha,\text{target})\), with zero-based matrix indices, are
\[
\begin{aligned}
 (0,1;(0,0,0,0);2;\;7/4,\;0)
       &\quad\text{in the Wilson quartic bank},\\
 (0,24;(1,0,0,0);1;\;-1/2,\;0)
       &\quad\text{in the Wilson \(+\) cubic bank},\\
 (0,0;(-2,-2,0,0);2;\;-12,\;0)
       &\quad\text{in the endpoint quartic bank}.
\end{aligned}                                               \tag{7.33}
\]
\end{theorem}

\begin{proof}
For a finite Laurent polynomial, equality of two matrix-valued
functions is equivalent to equality of every matrix Laurent
coefficient.  Momentum differentiation acts diagonally on those
coefficients, giving the two targets in (7.31) in the archived jet
convention.  The verifier loads the exact rational V60 quartic and
cubic banks, forms both targets coefficientwise, and compares them
with the stored first and second external jets using rational
arithmetic.  Its exhaustive defect ledger has the six positive
counts in (7.32).  The three displayed records are members of that
ledger and already contradict (7.31) in both actions.  Hence the
proposed factorwise identities fail.
\end{proof}

\begin{firewall}[Scope of the jet obstruction]
\label{fire:v3-jet-scope}
Theorem~\ref{thm:v3-jet-obstruction} rules out only a factor-by-factor
replacement of external jets by momentum derivatives.  It does not
rule out a separate Ward identity after all tadpole and bubble words
have been assembled into a closed trace.  Any such identity must be
derived for that complete signed scalar; it cannot be inferred by
applying (7.31) to its individual factors.
\end{firewall}

\subsection{A higher-grid diagnostic, kept outside the proof chain}

\begin{record}[The \(M=5\) floating diagnostic]
\label{rec:v3-m5}
An independent floating evaluation on the \(5^4=625\) tensor grid
gave
\[
\begin{split}
 \langle F_{\rm W}\rangle_5&\doteq 229.22438737241905,\\
 \langle F_*\rangle_5&\doteq -95.06225223097269,\\
 \langle F_*-F_{\rm W}\rangle_5
   &\doteq -324.28663960339173.
\end{split}                                                 \tag{7.34}
\]
The computation used \(313\) reflection representatives and had
maximum reported inverse residual \(2.63\times10^{-12}\).  Together
with the archived \(M=3\) value
\(-318.121568758848\) and \(M=4\) central value
\(-314.844179178182\), it supports continued pursuit of a negative
Haar mean.  It is not an interval computation and is not used in
Corollary~\ref{cor:v3-grid-to-haar}.  Moreover, the sampled \(M=4\)
pointwise difference assumes both signs (approximately
\(-1056.87\) to \(108.33\)); therefore pointwise negativity is not a
valid replacement for the average argument.
\end{record}

\subsection{Exact provenance}

\begin{record}[V3 verifier identities]
\label{rec:v3-provenance}
The residual-stock certificate is reproduced by
\begin{center}
\path{scripts/verify_ym_postarchive_v3_weighted_residual_stock.py}.
\end{center}
Its source has \(215\) lines, \(7447\) bytes, and SHA--256
\[
 \texttt{7BC12999878444C1DE5B98A11DA2F59345078032F00E2FFB7CDA88F939113886}.
                                                                  \tag{7.35}
\]
The canonical JSON payload has SHA--256
\[
 \texttt{798AB0D3537D7F5D7DD5509FF8281F8AF902F547A96BD517CD1B52B589CC903A}.
                                                                  \tag{7.36}
\]
The factorwise-jet comparison is reproduced by
\begin{center}
\path{scripts/verify_ym_postarchive_v3_jet_obstruction.py}.
\end{center}
Its source has \(152\) lines, \(5002\) bytes, and SHA--256
\[
 \texttt{F2CFAADCFFE76AFF8285DCC78BB9DF8D59ADC1843F3FE1BFB9FEC0B74F9E4ED6}.
                                                                  \tag{7.37}
\]
Its canonical JSON payload has SHA--256
\[
 \texttt{9E302F73C1461B6269CE5F14F1755128F5FE771F79D0D42AA500C8701E924521}.
                                                                  \tag{7.38}
\]
Both verifiers import the V60 exact response oracle with SHA--256
\[
 \texttt{7D5BFEE7FB834DFB5132FA2E1A9C5BC3CED7D57D366C557C8F370B6AD24D6FC8}.
                                                                  \tag{7.39}
\]
\end{record}

\subsection{Upstream decision for V4}

\begin{decision}[Do not enlarge a failed global derivative bound]
\label{dec:v3-next}
The operative target is the average transport increment (7.12), not
a uniform pointwise sign.  V4 should proceed in the following order.
\begin{enumerate}[label=\arabic*.]
\item Build residual-weighted \(H,Q_2,P_i,M_i\) Laurent stocks at all
      \(256\) quarter-turn anchors and test the complete response
      radius (7.10), retaining the seven signed centres.
\item Subdivide only cells whose contribution prevents (7.12), and
      use grouped phase intervals or local Taylor cancellation before
      taking coefficient norms.
\item If the certified radii remain of order \(10^{-2}\), go upstream
      to the archived alias decomposition: the \(45\)-dimensional
      vertical block is inverted first and the constrained retained
      four-direction Schur object is transported with its moving
      kernel.  One must not replace that constraint by an independent
      \(4\times4\) positivity assertion.
\end{enumerate}
The rejected alternatives are a global second-derivative majorant,
which numerically demands an inoperative grid scale, and the
factorwise derivative shortcut excluded by
Theorem~\ref{thm:v3-jet-obstruction}.
\end{decision}

\begin{assessment}[V3 progress]
\label{ass:v3-progress}
V3 supplies the missing local theorem that turns exact trial
inverses, residuals, and factor transports into a complete
cell-response interval.  It gives the exact sufficient inequality
(7.12) whose discharge would promote the V2 grid sign to a Haar-mean
sign.  It also proves that coefficientwise residual weighting yields
roughly eightfold improvement over the crude operator-product stock
at the toron anchor, and that local Taylor cancellation enlarges the
certified Neumann radii.  The quarter-turn cells are nevertheless far
larger than the one-anchor certificate, so (7.12) has not yet been
verified.
\end{assessment}

\begin{firewall}[Claim boundary after V3]
\label{fire:v3-final}
V3 does not prove the sign of the continuous Haar mean.  It therefore
does not yet prove a nonzero Gaussian constant edge class.  It does
not populate the compensating scale boundary, prove SCIC or ADMC,
construct the cofinal interacting state or continuum measure,
establish reflection positivity or physical exponential clustering,
reconstruct a Hamiltonian, or prove the Yang--Mills mass gap.
\end{firewall}

% =============================================================================
% END APPENDED POST-TENTH-ARCHIVE V3 MATERIAL
% =============================================================================

% =============================================================================
% BEGIN APPENDED POST-TENTH-ARCHIVE V4 MATERIAL
% =============================================================================

\section{V4: constraint-preserving alias inversion and a stable
rank-three compiler}
\label{sec:v4-alias-compiler}

\subsection{Why the V3 cell architecture must be moved upstream}

V3 supplied a rigorous cell theorem and the exact sufficient
Haar-sign condition (7.12).  Before building a large rational
subdivision, the complete response radius was evaluated at every
quarter-turn anchor with the same residual-weighted Taylor stocks.
This diagnostic is decisive about architecture, although not about
the sign of the true scalar.

\begin{record}[All-anchor transport reconnaissance]
\label{rec:v4-all-anchor}
At the full quarter-cell radius \(r=\pi/4\), every one of the \(256\)
weighted points fails the sufficient Neumann condition for both
actions.  The ranges of the computed Hessian residual majorants are
\[
\begin{array}{c|cc}
 &\min&\max\\ \hline
 {\rm Wilson}&63.0621&95.1889\\
 *&124.6242&174.6364.
\end{array}                                                \tag{8.1}
\]
The local radii at which the Hessian majorant first reaches one have
the ranges
\[
\begin{array}{c|cc}
 &\min&\max\\ \hline
 {\rm Wilson}&0.0181524&0.0282462\\
 *&0.0151127&0.0215843.
\end{array}                                                \tag{8.2}
\]
More severely, after all factor-variation radii in the eight complete
response words are included, the pointwise radii that consume the
entire exact margin
\(\delta_4=147981586239/500000000\) range from
\[
 7.9088\times10^{-7}
 \quad\hbox{to}\quad
 6.7901\times10^{-6},                                      \tag{8.3}
\]
with median \(2.1590\times10^{-6}\).
\end{record}

\begin{firewall}[Meaning of the reconnaissance]
\label{fire:v4-recon-scope}
The numbers in (8.1)--(8.3) are floating evaluations of rigorous
upper-bound formulae.  They do not lower-bound the true residual or
quadrature error and hence do not prove that subdivision is
impossible.  They do prove that blindly rationalizing this particular
triangle-majorant implementation would target millions of cells per
axis before it even addresses the Haar error.  The correct response
is to expose the alias constraint and cancel the moving-kernel terms
before taking norms.
\end{firewall}

\subsection{The sixteen-alias transverse geometry}

Work in the fundamental coarse cube
\({\cal K}=[-\pi,\pi]^4\).  For
\(\epsilon\in\{0,1\}^4\), put
\[
\begin{split}
 p_\epsilon(k)&:=\frac{k+2\pi\epsilon}{2},\qquad
 q_\epsilon:=e^{ip_\epsilon}-{\bf1},\\
 b_\epsilon&:={\bf1}+e^{ip_\epsilon},\qquad
 \lambda_\epsilon:=q_\epsilon^*q_\epsilon,\\
 Q&:=e^{ik}-{\bf1},\qquad
 D_\epsilon:=\operatorname{diag}(b_\epsilon).
\end{split}                                                 \tag{8.4}
\]
All exponentials in (8.4) are componentwise.  The elementary
half-angle identity gives
\[
 D_\epsilon q_\epsilon=Q.                                 \tag{8.5}
\]
Let
\[
 {\cal C}_\epsilon
 :=\lambda_\epsilon I_4-q_\epsilon q_\epsilon^*
 =\lambda_\epsilon P_\epsilon,                            \tag{8.6}
\]
where \(P_\epsilon\) is the orthogonal projection onto
\(q_\epsilon^\perp\) when \(\lambda_\epsilon>0\), and is set equal
to zero when \(\lambda_\epsilon=0\).  On
\({\cal E}=\bigoplus_\epsilon\mathbb C^4\), define
\[
\begin{split}
 {\cal M}_{{\rm W},\epsilon}
   &:=\frac13{\cal C}_\epsilon,\\
 {\cal M}_{*,\epsilon}
   &:=\left(\frac13+\lambda_\epsilon\right){\cal C}_\epsilon,
 \qquad
 {\cal M}_e:=\bigoplus_\epsilon{\cal M}_{e,\epsilon},\\
 P&:=\bigoplus_\epsilon P_\epsilon.
\end{split}                                                 \tag{8.7}
\]

Let \(V(k):\mathbb C^{45}\to{\cal E}\) be the normalized parity
Fourier frame restricted to the V21 retained edges.  The archived
alias reduction gives
\[
 V^*V=I_{45},\qquad H_e=V^*{\cal M}_eV.                    \tag{8.8}
\]
Set
\[
 T:=PV.                                                     \tag{8.9}
\]
Since \({\cal M}_e=P{\cal M}_eP\), (8.8) is equivalently
\[
 H_e=T^*{\cal M}_eT.                                      \tag{8.10}
\]

Define the coarse-link map
\[
 {\cal B}(a_\epsilon)_\epsilon
 :=\sum_\epsilon D_\epsilon a_\epsilon.                   \tag{8.11}
\]

\begin{lemma}[Exact retained-edge constraint]
\label{lem:v4-coarse-link}
For every \(k\in{\cal K}\),
\[
 {\cal B}V=0.                                              \tag{8.12}
\]
If \(k\ne0\), let \(J:\mathbb C^3\to Q^\perp\) be any isometry
onto \(Q^\perp\), and define
\[
 C:=J^*{\cal B}P:{\cal E}\longrightarrow\mathbb C^3.       \tag{8.13}
\]
Then
\[
 CT=0,\qquad \operatorname{rank}C=3,\qquad
 \operatorname{Ran}T=\ker C\cap\operatorname{Ran}P.        \tag{8.14}
\]
At \(k=0\), \(\operatorname{Ran}P\) has dimension \(45\) and
\[
 \operatorname{Ran}T=\operatorname{Ran}P.                 \tag{8.15}
\]
\end{lemma}

\begin{proof}
For a retained parity edge \((s,\mu)\), its contribution to
\({\cal B}V\) is confined to component \(\mu\) and, apart from the
nonzero common factor \(e^{-ik\cdot s/2}/4\), has the two Fourier
coefficients
\[
 \sum_\epsilon(-1)^{\epsilon\cdot s},
 \qquad
 \sum_\epsilon(-1)^{\epsilon\cdot s+\epsilon_\mu}.         \tag{8.16}
\]
Every retained V21 edge has \(s_\nu=1\) for at least one
\(\nu\ne\mu\).  Summing first over \(\epsilon_\nu\) makes both
coefficients in (8.16) zero.  This proves (8.12) coefficientwise.

Suppose \(k\ne0\).  All \(\lambda_\epsilon\) are then positive.
For \(a=PVx\), equations (8.5) and (8.12) give
\[
 {\cal B}a
 =-Q\sum_\epsilon
       \frac{q_\epsilon^*(Vx)_\epsilon}{\lambda_\epsilon}
 \in\operatorname{span}\{Q\}.
                                                                  \tag{8.17}
\]
Therefore \(CT=0\).

For the low alias \(\epsilon=0\), \(D_0\) is invertible on
\({\cal K}\).  Let \(J_0:\mathbb C^3\to q_0^\perp\) be an
isometry.  If
\(J^*D_0J_0y=0\), then \(D_0J_0y\in\operatorname{span}\{Q\}\).
By (8.5) and invertibility of \(D_0\),
\(J_0y\in\operatorname{span}\{q_0\}\).  It is also orthogonal to
\(q_0\), so \(y=0\).  Thus the \(3\times3\) map
\(J^*D_0J_0\) is invertible and \(C\) has rank three.

The global floors (1.2) and (8.10) imply that \(T\) is injective,
so \(\dim\operatorname{Ran}T=45\).  The active transverse space
\(\operatorname{Ran}P\) has dimension \(48\) for \(k\ne0\).
Together with \(CT=0\) and \(\operatorname{rank}C=3\), dimension
counting proves the last equality in (8.14).

At \(k=0\), the low curl block vanishes and the other fifteen curl
blocks have rank three.  Hence \(\dim\operatorname{Ran}P=45\).
Injectivity of \(T\) then proves (8.15).
\end{proof}

\subsection{The constrained inverse and its scalar Schur collapse}

The Moore--Penrose inverse of the direct-sum Hessian is elementary.
On the active transverse space write
\[
 {\cal D}_e:={\cal M}_e^+
 =\bigoplus_\epsilon d_e(\lambda_\epsilon)P_\epsilon,       \tag{8.18}
\]
where a zero-curl block contributes zero and, for \(\lambda>0\),
\[
 d_{\rm W}(\lambda)=\frac3\lambda,\qquad
 d_*(\lambda)=\frac1{\lambda(1/3+\lambda)}.                \tag{8.19}
\]

\begin{theorem}[Constraint-preserving alias inverse]
\label{thm:v4-constrained-inverse}
Let
\[
 K_e:=T H_e^{-1}T^*.                                      \tag{8.20}
\]
For \(k\ne0\), set
\[
 S_e:=C{\cal D}_eC^*.
\]
Then \(S_e\) is positive definite on \(\mathbb C^3\) and
\[
\boxed{\quad
 K_e={\cal D}_e
 -{\cal D}_eC^*S_e^{-1}C{\cal D}_e .
\quad}                                                     \tag{8.21}
\]
At \(k=0\), the correction disappears:
\[
 K_e={\cal D}_e.                                           \tag{8.22}
\]
Thus the action-dependent inversion of the \(45\times45\) Hessian
is replaced exactly by sixteen scalar transverse inverses and one
\(3\times3\) Schur inverse.
\end{theorem}

\begin{proof}
For \(k\ne0\), Lemma~\ref{lem:v4-coarse-link} identifies
\(\operatorname{Ran}T\) with the kernel of \(C\) in the active
space.  Since \({\cal D}_e\) is positive definite on that space and
\(C\) is onto, \(S_e=C{\cal D}_eC^*>0\).

Call the right side of (8.21) \(\widetilde K_e\).  Direct
multiplication gives
\[
 C\widetilde K_e=0.                                       \tag{8.23}
\]
If \(y\in\ker C\cap\operatorname{Ran}P\), then
\({\cal D}_e{\cal M}_ey=y\), and hence
\[
 \widetilde K_e{\cal M}_ey
 =y-{\cal D}_eC^*S_e^{-1}Cy=y.                            \tag{8.24}
\]
The matrix \(\widetilde K_e\) is Hermitian and its range is contained
in \(\operatorname{Ran}T\).  Therefore
\(\widetilde K_e=TAT^*\) for a Hermitian \(45\times45\) matrix
\(A\).  Substitute \(y=Tx\) in (8.24), use (8.10), and obtain
\[
 TA H_e x=Tx.
\]
Injectivity of \(T\) gives \(A=H_e^{-1}\), proving (8.21).
At the toron, (8.15) says that the active space itself is
\(\operatorname{Ran}T\); \({\cal D}_e\) is its inverse, which proves
(8.22).
\end{proof}

The \(3\times3\) Schur matrix is simpler than it first appears.  Put
\[
 \sigma_{e,\mu}(k)
 :=\sum_\epsilon
 d_e(\lambda_\epsilon)|b_{\epsilon,\mu}|^2,\qquad
 \Sigma_e:=\operatorname{diag}
  (\sigma_{e,1},\ldots,\sigma_{e,4}).                       \tag{8.25}
\]

\begin{theorem}[Four scalar sums determine the Schur inverse]
\label{thm:v4-scalar-schur}
For \(k\ne0\),
\[
 S_e=J^*\Sigma_eJ.                                        \tag{8.26}
\]
Moreover the inverse lifted back to \(Q^\perp\) is
\[
\boxed{\quad
 {\cal R}_e:=
 J(J^*\Sigma_eJ)^{-1}J^*
 =
 \Sigma_e^{-1}
 -\frac{\Sigma_e^{-1}QQ^*\Sigma_e^{-1}}
        {Q^*\Sigma_e^{-1}Q}.
\quad}                                                     \tag{8.27}
\]
Consequently (8.21) is independent of the chosen frame \(J\) and can
be written
\[
 K_e={\cal D}_e
 -{\cal D}_e({\cal B}P)^*{\cal R}_e
       ({\cal B}P){\cal D}_e.                              \tag{8.28}
\]
\end{theorem}

\begin{proof}
On the \(\epsilon\)-block,
\[
\begin{split}
 D_\epsilon P_\epsilon D_\epsilon^*
 &=D_\epsilon D_\epsilon^*
   -\frac{(D_\epsilon q_\epsilon)
          (D_\epsilon q_\epsilon)^*}{\lambda_\epsilon}\\
 &=\operatorname{diag}(|b_{\epsilon,1}|^2,\ldots,
                       |b_{\epsilon,4}|^2)
   -\frac{QQ^*}{\lambda_\epsilon}.
\end{split}                                                 \tag{8.29}
\]
Insert (8.18), sum over aliases, and sandwich by \(J^*\) and \(J\).
Because \(J^*Q=0\), all rank-one terms in (8.29) vanish and (8.26)
follows.

Let the right side of (8.27) be \(R_e\).  It is Hermitian,
\(Q^*R_e=0\), and for every \(y\perp Q\),
\[
 R_e\Sigma_ey=y.                                          \tag{8.30}
\]
Thus \(R_e\) is precisely the inverse of \(\Sigma_e\) restricted to
\(Q^\perp\), lifted to \(\mathbb C^4\).  This proves (8.27).
Substitution of \(C=J^*{\cal B}P\) into (8.21) gives (8.28).
\end{proof}

\subsection{A toron-stable low-alias Woodbury form}

Formula (8.28) is exact, but near \(k=0\) its low-alias block contains
a large pseudoinverse which is cancelled by the correction.  The
next parameterization performs that cancellation algebraically.

For \(k\ne0\), choose isometries
\[
 J_\epsilon:\mathbb C^3\longrightarrow q_\epsilon^\perp,
 \qquad
 J:\mathbb C^3\longrightarrow Q^\perp,
\]
and put
\[
 A_\epsilon:=J^*D_\epsilon J_\epsilon,\qquad
 A_h:=(A_\epsilon)_{\epsilon\ne0},\qquad
 L:=A_0^{-1}A_h.                                          \tag{8.31}
\]
Define the transverse energy weights
\[
 w_{\rm W}(\lambda):=\frac{\lambda}{3},\qquad
 w_*(\lambda):=\lambda\left(\frac13+\lambda\right),        \tag{8.32}
\]
and
\[
 D_{e,h}:=\bigoplus_{\epsilon\ne0}
             w_e(\lambda_\epsilon)I_3.                    \tag{8.33}
\]

\begin{theorem}[Stable high-alias covariance]
\label{thm:v4-stable-woodbury}
The matrix \(A_0\) in (8.31) is invertible.  Every constrained
transverse vector has the unique form
\[
 y_0=-Ly_h,\qquad
 (a_0,a_h)=(-J_0Ly_h,(J_\epsilon y_\epsilon)_{\epsilon\ne0}).
                                                                  \tag{8.34}
\]
In these \(45\) high-alias coordinates, the Hessian is
\[
 N_e=D_{e,h}+w_e(\lambda_0)L^*L.                           \tag{8.35}
\]
For \(k\ne0\),
\[
\boxed{\quad
 N_e^{-1}
 =D_{e,h}^{-1}
 -D_{e,h}^{-1}L^*
  \left[
   w_e(\lambda_0)^{-1}I_3
   +LD_{e,h}^{-1}L^*
  \right]^{-1}
  LD_{e,h}^{-1}.
\quad}                                                     \tag{8.36}
\]
At \(k=0\), the low block has zero energy and
\[
 N_e^{-1}=D_{e,h}^{-1}.                                   \tag{8.37}
\]
Hence the toron limit is represented without subtracting two terms
of order \(\lambda_0^{-1}\).
\end{theorem}

\begin{proof}
The proof of Lemma~\ref{lem:v4-coarse-link} already shows that
\(A_0=J^*D_0J_0\) is invertible.  In transverse coordinates the
constraint is
\[
 A_0y_0+A_hy_h=0,
\]
which is equivalent to (8.34).  The direct-sum action on
\(q_\epsilon^\perp\) is scalar multiplication by
\(w_e(\lambda_\epsilon)\).  Substituting (8.34) into the sum of the
sixteen block energies gives (8.35).  The Woodbury identity gives
(8.36).  At \(k=0\), the low block vanishes, the fifteen high
transverse blocks already have total dimension \(45\), and the
rank-three update in (8.35) is zero.  This proves (8.37).
\end{proof}

\subsection{A complete trace-word compiler}

The covariance in (8.20) acts on the active alias space, whereas the
response factors in (6.12) are stored in the \(45\) retained
coordinates.  The following basis-free identity connects the two
without inserting another action-dependent inverse.

\begin{theorem}[Constraint-first trace-word identity]
\label{thm:v4-trace-compiler}
Define the action-independent left inverse
\[
 \Lambda:=(T^*T)^{-1}T^*,\qquad \Lambda T=I_{45}.          \tag{8.38}
\]
For every reduced \(45\times45\) factor \(X\), put
\[
 \widehat X:=\Lambda^*X\Lambda\quad\text{on }{\cal E}.      \tag{8.39}
\]
Then, for arbitrary reduced factors \(X_1,\ldots,X_m\),
\[
\boxed{\quad
 \Tr_{\mathbb C^{45}}
   (H_e^{-1}X_1\cdots H_e^{-1}X_m)
 =
 \Tr_{\cal E}
   (K_e\widehat X_1\cdots K_e\widehat X_m).
\quad}                                                     \tag{8.40}
\]
In particular, replacing every complete word in (6.12) by its
right side in (8.40), and the tadpole by the \(m=1\) case, reproduces
the full order-two response \(F_e(k)\).
\end{theorem}

\begin{proof}
Write \(G_e=H_e^{-1}\).  By definition \(K_e=TG_eT^*\), and
\(\Lambda T=I\) implies
\[
 G_e=\Lambda K_e\Lambda^*.                                \tag{8.41}
\]
Substitute (8.41) in the left side of (8.40).  Repeated use of
\(\Lambda T=I\), \(T^*\Lambda^*=I\), and cyclicity of the
rectangular trace gives
\[
\begin{split}
 \Tr(G_eX_1\cdots G_eX_m)
 &=\Tr\!\left(
   \Lambda K_e\Lambda^*X_1\cdots
   \Lambda K_e\Lambda^*X_m\right)\\
 &=\Tr_{\cal E}\!\left(
   K_e\Lambda^*X_1\Lambda\cdots
   K_e\Lambda^*X_m\Lambda\right),
\end{split}                                                 \tag{8.42}
\]
which is (8.40).  Formula (6.12) is a signed sum of one such
one-factor word and seven such multi-factor words, so the last claim
follows term by term without any triangle inequality.
\end{proof}

\begin{corollary}[Action-dependent inversion has rank three]
\label{cor:v4-rank-three-response}
After the transformations (8.38)--(8.39), every covariance occurrence
in the complete response may be evaluated by either (8.28) or the
stable form (8.36).  The only non-diagonal action-dependent inverse
has size \(3\times3\); the remaining change of coordinates is the
common geometric map \(\Lambda(k)\).
\end{corollary}

\begin{proof}
Combine Theorems~\ref{thm:v4-constrained-inverse},
\ref{thm:v4-stable-woodbury}, and
\ref{thm:v4-trace-compiler}.
\end{proof}

\begin{firewall}[The compiler is not yet an interval theorem]
\label{fire:v4-compiler-scope}
The trace identity (8.40) is exact for arbitrary reduced factors, so
it does not assume that the cubic or quartic jets separately preserve
the moving gauge kernel.  Those effects remain inside
\(\widehat X\).  However, V4 has not yet produced coefficient-exact
Laurent or interval enclosures for \(\Lambda\), \(L\), or the lifted
response factors.  The rank-three reduction therefore identifies
the correct cancellation level but does not by itself prove the Haar
sign.
\end{firewall}

\subsection{Executable checks and provenance}

\begin{record}[Constraint-preserving alias compiler]
\label{rec:v4-alias-verifier}
The executable
\begin{center}
\path{scripts/verify_ym_postarchive_v4_constrained_alias_inverse.py}
\end{center}
checks the coarse-link coefficients of all \(45\) retained edges
exactly.  It then cross-checks all \(256\) weighted quarter-turn
points for both actions.  The maximum floating relative residuals
are
\[
\begin{array}{c|cc}
 &{\rm Wilson}&*\\ \hline
 H_{\rm alias}-H_{\rm reduced}
   &1.015\times10^{-15}&1.020\times10^{-15}\\
 K_{\rm Schur}-K_{\rm direct}
   &7.86\times10^{-15}&5.16\times10^{-15}\\
 K_{\rm stable}-K_{\rm direct}
   &8.10\times10^{-15}&5.30\times10^{-15}\\
 F_{\rm alias}-F_{\rm direct}
   &1.72\times10^{-13}&2.74\times10^{-13}.
\end{array}                                                \tag{8.43}
\]
At every non-toron grid point, \(T\) has rank \(45\), \(C\) and
\(S_e\) have rank \(3\), and the Woodbury middle matrix has rank
\(3\).  At the toron the three latter ranks are zero and the
correction vanishes.  The least sampled singular value of \(A_0\)
is \(0.7653668647\).

The source has \(544\) lines, \(18467\) bytes, and SHA--256
\[
 \texttt{309FDFB53F02708FE82720E6D4491732C824F41685F98AD69C89CF5490401626}.
                                                                  \tag{8.44}
\]
Its canonical payload has SHA--256
\[
 \texttt{F955ADCC17178712B0835B2A9092B182A08D9B6D5F8A478C2F07BBDE393DED44}.
                                                                  \tag{8.45}
\]
The exact inputs include the V21 incidence source with SHA--256
\[
 \texttt{DE678E5CAE7671B90C8788AF9A489444D8A63B4F32EA78F08BC68AA0131B0AE0}
                                                                  \tag{8.46}
\]
and the V60 response oracle with hash (7.39).
\end{record}

\begin{record}[All-anchor diagnostic provenance]
\label{rec:v4-transport-provenance}
The reconnaissance in Record~\ref{rec:v4-all-anchor} is reproduced
by
\begin{center}
\path{scripts/explore_ym_postarchive_v4_all_anchor_transport.py}.
\end{center}
It has \(357\) lines, \(10604\) bytes, and SHA--256
\[
 \texttt{6FE9B6B60CFE9E92B4CC6CEBE2208D6CB716B5AD88C7789A5D9CF37F1EB895A7}.
                                                                  \tag{8.47}
\]
Its output is explicitly diagnostic and is not imported into any
theorem above.
\end{record}

\subsection{V5 programme}

\begin{decision}[Compile before bounding]
\label{dec:v4-v5}
V5 must begin with the scheduled five-version audit of the current
Standard-Model and Navier--Stokes notes in \path{latex_notes}.
Any transferable constrained-inverse, endpoint-cancellation, or
source-concentration device should be recorded before the YM route
is fixed.

The YM calculation should then use the stable high-alias coordinates
(8.34), transform the actual \(Q_2,P_i,M_i,H_i\) banks at the level
of complete signed words, and simplify the four scalar sums
\(\sigma_{e,\mu}\) before interval norms are taken.  The first gate
is an exact replay of the V2 quarter-turn scalar from (8.36) and
(8.40).  The second is a cancellation-retaining interval or Fourier
tail bound for the resulting scalar \(3\times3\) compiler.

If the geometric factor \(\Lambda\) reintroduces a large
\(45\times45\) norm loss, replace it by the explicit high-alias
extension (8.34); do not return to the V3 complete-word triangle
radius and do not split the sixteen moving-kernel blocks
independently.
\end{decision}

\begin{assessment}[V4 progress]
\label{ass:v4-progress}
V4 closes the principal logical defect in a naive alias reduction.
It proves the exact constraint, derives the constrained covariance,
collapses its Schur matrix to four scalar alias sums, removes the
toron pseudoinverse cancellation by a high-alias Woodbury formula,
and compiles every complete response trace word into that
constraint-preserving representation.  The remaining finite-stage
bottleneck is no longer a \(45\times45\) action-dependent inverse;
it is the coefficient-exact control of a common geometric extension,
four scalar sums, and one rank-three correction.
\end{assessment}

\begin{firewall}[Claim boundary after V4]
\label{fire:v4-final}
V4 does not prove the continuous Haar-mean sign or a nonzero Gaussian
edge class.  It does not close the compensating scale boundary,
SCIC, ADMC, the interacting cofinal state, regulator removal,
reflection-positive continuum reconstruction, physical exponential
clustering, or the Yang--Mills mass gap.
\end{firewall}

% =============================================================================
% END APPENDED POST-TENTH-ARCHIVE V4 MATERIAL
% =============================================================================

% =============================================================================
% BEGIN APPENDED POST-TENTH-ARCHIVE V5 MATERIAL
% =============================================================================

\section{V5: companion audit and rank-three complete-response compiler}
\label{sec:v5-rank-three}

V4 reduced the action-dependent inverse in every response word to a
diagonal transverse alias inverse minus one rank-three Schur correction.
The present note performs the scheduled companion audit, expands that
correction without destroying signs, and reassembles the seven bubble
words as one exact second derivative.  The purpose is to expose the
smallest scalar object on which a coefficient-exact enclosure should be
attempted.  No continuum or mass-gap conclusion is drawn from the finite
grid diagnostics below.

\subsection{Scheduled audit of the current companion programmes}

\begin{audit}[Sources inspected at V5]
\label{aud:v5-companions}
The current companion sources inspected for this continuation were
\path{Renewal_SM_Einstein_Continuum_Research_Notes_V32.tex} and
\path{Renewal_NS_Stability_Research_Notes_V26.tex}.  Their respective
SHA--256 values are
\[
\begin{gathered}
\texttt{0724889C70F0E181431F6D91CDEEC558}\\
\texttt{F3EA2CD5DF44760A036806F4369CBA00},\\[3pt]
\texttt{82C887F8C2C69E4F28FAD7C3DC8DE5B9}\\
\texttt{099CB5A4BF431EBA1FFF3E91935978AD}.
\end{gathered}                                                 \tag{9.1}
\]
The SM source advanced from V29 to V32 while this audit was in
progress, so the final audit used V32, which contains the earlier
material append-only.

Three lessons transfer at the level of proof strategy.  First, the SM
additive tree--Schur ledger retains all positive pieces and all actual
low-rank Schur losses before performing the final comparison; it does
not multiply unrelated local relative slacks.  Second, its V31
collision-core compiler forms the complete signed local response before
using a triangle inequality.  Third, V28--V32 keep a connected
three-arm object intact until after its common propagation has been
resolved.  In the present finite-alias problem these principles say to
retain the complete signed response and its common rank-three Schur
core, rather than bounding the seven bubble words or sixteen aliases
independently.

The NS V26 audit supplies a necessary warning.  Its exact Oseen
calculation shows that a rank-one input restriction gives no gain for
the first derivative and only a modest gain for the second derivative.
A low-rank norm is legitimate only when the actual input has that rank.
Here the Schur correction really factors through \(\mathbb C^3\), so a
rank-three trace reduction is legitimate; no such restriction will be
imposed on the unrestricted lifted vertex factors.

These are methodological transfers only.  No SM or NS estimate is
imported as a Yang--Mills theorem, and both companion continuum
programmes remain open.
\end{audit}

\subsection{Exact expansion of the constrained covariance}

Fix an action \(e\) and a non-toron momentum.  Use the notation of
Theorem~\ref{thm:v4-constrained-inverse} and set
\[
 D:=\mathcal D_e,\qquad
 S:=CD C^*,\qquad R:=S^{-1},\qquad U:=DC^*.
                                                               \tag{9.2}
\]
Thus \(U:\mathbb C^3\to\mathcal E\), and (8.20) becomes
\[
 K=D-URU^*.                                                   \tag{9.3}
\]
At the toron we use the same notation with \(U\) the unique map from
the zero-dimensional space and hence \(URU^*=0\).

Let \(X_1,\ldots,X_m\) be arbitrary endomorphisms of \(\mathcal E\).
For a nonempty subset \(I\subseteq\{1,\ldots,m\}\), list its elements
in cyclic order as \(i_1,\ldots,i_p\), put \(i_{p+1}=i_1\), and define
\[
 A_\ell(I):=
 U^*X_{i_\ell}
 \left(\prod_{j=i_\ell+1}^{\,i_{\ell+1}-1}DX_j\right)U.
                                                               \tag{9.4}
\]
The product follows the positive cyclic order modulo \(m\).  It is the
identity when the selected positions are consecutive.  In particular,
when \(p=1\), the product traverses every unselected position before
returning to \(i_1\).

\begin{theorem}[Rank-three cyclic subset expansion]
\label{thm:v5-subset-expansion}
For arbitrary \(X_1,\ldots,X_m\),
\[
\boxed{
\begin{aligned}
 \Tr_{\mathcal E}\prod_{j=1}^m(KX_j)
 ={}&
 \Tr_{\mathcal E}\prod_{j=1}^m(DX_j)\\
 &+\sum_{\varnothing\ne I\subseteq\{1,\ldots,m\}}
 (-1)^{|I|}
 \Tr_{\mathbb C^3}
 \bigl(RA_1(I)\cdots RA_{|I|}(I)\bigr).
\end{aligned}}                                                \tag{9.5}
\]
At the toron the sum on the second line is empty.  No commutation,
positivity, self-adjointness, or invariant-subspace assumption on the
\(X_j\) is required.
\end{theorem}

\begin{proof}
Expand each occurrence of \(K=D-URU^*\) in the left side of
(9.5).  The term indexed by \(I\) chooses \(URU^*\) precisely at the
positions in \(I\) and has sign \((-1)^{|I|}\).  The empty choice is
the first term on the right.

For a nonempty choice, cyclically start the ambient trace at the factor
\(U\) belonging to \(i_1\).  Between the correction at \(i_\ell\) and
the next correction there occurs exactly
\[
 U R U^*X_{i_\ell}
 \left(\prod_{j=i_\ell+1}^{\,i_{\ell+1}-1}DX_j\right)U .
\]
Move the initial rectangular \(U\) through the trace.  The identity
\(\Tr_{\mathcal E}(UV)=\Tr_{\mathbb C^3}(VU)\) then turns the selected
term into
\[
 \Tr_{\mathbb C^3}
 \bigl(RA_1(I)\cdots RA_p(I)\bigr).
\]
Summing all subsets proves (9.5).  At the toron \(U=0\), so every
nonempty term vanishes.
\end{proof}

\begin{corollary}[Exact size of the response correction]
\label{cor:v5-forty-six}
After applying the lifting (8.39), a response word with \(m\)
covariance occurrences has one direct-alias trace and \(2^m-1\)
rank-three traces.  The complete order-two response has word lengths
\[
 1;\quad 2,4,3,2,3,2,3,
\]
and therefore has eight direct-alias traces and exactly
\[
 (2^1-1)+(2^2-1)+(2^4-1)+3(2^3-1)+2(2^2-1)
 =46                                                       \tag{9.6}
\]
nonempty rank-three subset traces at every non-toron momentum.
\end{corollary}

\begin{proof}
Apply Theorem~\ref{thm:v5-subset-expansion} to the tadpole and the
seven word lengths in (6.12).  The arithmetic in (9.6) is
\(1+3+15+7+3+7+3+7=46\).
\end{proof}

\begin{remark}[Correct cancellation hierarchy]
\label{rem:v5-hierarchy}
Formula (9.5) is an identity, not a prescription to take the absolute
value of its 46 correction terms.  A rigorous implementation should
accumulate the signed \(3\times3\) traces coefficient-exactly or by
outward-rounded interval arithmetic and enclose their sum only
afterward.  The SM audit supports precisely this order of operations.
\end{remark}

\subsection{The seven bubble words are one second derivative}

The signs and multiplicities in (6.12) are not accidental.  They are
the product rule for one complete response jet.

\begin{lemma}[Inverse jets]
\label{lem:v5-inverse-jets}
Let \(H(t)\) be twice differentiable and invertible near \(t=0\), put
\(G(t)=H(t)^{-1}\), and write \(H_j=H^{(j)}(0)\),
\(G_j=G^{(j)}(0)\).  Then
\[
 G_1=-G_0H_1G_0,\qquad
 G_2=2G_0H_1G_0H_1G_0-G_0H_2G_0.                            \tag{9.7}
\]
\end{lemma}

\begin{proof}
Differentiate \(H(t)G(t)=I\) once and solve for \(G_1\).
Differentiate it twice:
\[
 H_2G_0+2H_1G_1+H_0G_2=0.
\]
Multiplication by \(G_0\) on the left and substitution of the first
identity give (9.7).
\end{proof}

\begin{theorem}[Frozen-outer-propagator reconstruction]
\label{thm:v5-bubble-jet}
Let \(P(t)\) and \(M(t)\) be twice differentiable matrix curves and
write \(P_j=P^{(j)}(0)\), \(M_j=M^{(j)}(0)\).  Hold the outer
propagator \(G_0\) fixed and define
\[
 \Psi(t):=\Tr\bigl(G_0P(t)G(t)M(t)\bigr).                    \tag{9.8}
\]
Then
\[
\begin{split}
 \Psi''(0)={}&
 \Tr(G_0P_2G_0M_0)
 +2\Tr(G_0P_0G_0H_1G_0H_1G_0M_0)\\
 &-\Tr(G_0P_0G_0H_2G_0M_0)
 +\Tr(G_0P_0G_0M_2)\\
 &-2\Tr(G_0P_1G_0H_1G_0M_0)
 +2\Tr(G_0P_1G_0M_1)\\
 &-2\Tr(G_0P_0G_0H_1G_0M_1).
\end{split}                                                   \tag{9.9}
\]
Consequently the seven bubble words of (6.12), including their
coefficients and order, are exactly \(\Psi''(0)\), multiplied by the
declared action factor.
\end{theorem}

\begin{proof}
Twice applying the product rule to (9.8), while keeping its first
\(G_0\) fixed, gives
\[
\begin{split}
 \Psi''(0)=\Tr\bigl(G_0[
 &P_2G_0M_0+2P_1G_1M_0+2P_1G_0M_1\\
 &+P_0G_2M_0+2P_0G_1M_1+P_0G_0M_2]\bigr).
\end{split}                                                   \tag{9.10}
\]
Substitute (9.7) and distribute the trace.  The resulting seven
terms are (9.9), in the same order as the seven entries of (6.12).
\end{proof}

\begin{corollary}[Smallest complete scalar target]
\label{cor:v5-smallest-target}
For each action and momentum, the full order-two scalar can be
evaluated without splitting its bubble part as
\[
 F_e(k)=\frac12\Tr\bigl(K_e\widehat Q_2\bigr)
       +a_e\,\Psi_e''(0),                                   \tag{9.11}
\]
where \(a_{\rm W}=1/3\), \(a_*=3\), and every occurrence of \(G_0\)
in (9.8)--(9.10) is represented by the exact constrained covariance
\(K_e\).  Thus a future interval proof may differentiate the complete
\(3\times3\)-corrected expression (9.8) as a unit.
\end{corollary}

\begin{proof}
The tadpole is the first term of (6.12).
Theorem~\ref{thm:v5-bubble-jet} identifies the remaining seven terms,
and Theorem~\ref{thm:v4-trace-compiler} replaces every reduced
covariance word by its constrained alias representation.
\end{proof}

\subsection{Exhaustive finite-grid diagnostic}

\begin{record}[Rank-three verifier]
\label{rec:v5-verifier}
The executable
\begin{center}
\path{scripts/verify_ym_postarchive_v5_rank_three_response_expansion.py}
\end{center}
checks all nonempty subsets in all eight words at every one of the
\(256\) weighted quarter-turn momenta for both actions.  For each
subset it compares the ambient trace with the \(3\times3\) cyclic
trace in (9.5), reconstructs the complete response, and compares it
with the independent reduced-coordinate oracle.

The largest floating residuals were
\[
\begin{array}{c|cc}
 &{\rm Wilson}&*\\ \hline
\text{subset compression}
 &1.7245\times10^{-14}&7.3619\times10^{-15}\\
\text{complete reconstruction}
 &1.6212\times10^{-14}&1.1539\times10^{-13}\\
\text{independent oracle}
 &1.7442\times10^{-13}&2.6114\times10^{-13}.
\end{array}                                                   \tag{9.12}
\]
The observed non-toron correction count was always \(46\), and the
toron correction was zero.

Pairing each non-self-opposite momentum with its conjugate gives the
following diagnostic means:
\[
\begin{array}{c|rrr}
 &\text{direct alias}&\text{rank-three correction}&\text{total}\\ \hline
{\rm Wilson}
 & 232.6560414953&-9.9495895633& 222.7064519320\\
*
 &-141.9935799984&49.8558527523&-92.1377272462\\
*-{\rm Wilson}
 &-374.6496214938&59.8054423156&-314.8441791782.
\end{array}                                                   \tag{9.13}
\]
These numbers are floating reconnaissance, not enclosures.  They show
why the exact compilation matters.  For the endpoint action the
direct and rank-three bubble means are respectively
\[
 -40.8452116313,\qquad 37.9281822832,
\]
leaving only \(-2.9170293481\).  Within that bubble, the second and
third word means are
\[
 -204.3080835972,\qquad 199.8865856619.                       \tag{9.14}
\]
Bounding those words independently would discard the dominant
cancellation that Theorem~\ref{thm:v5-bubble-jet} preserves.

The source has \(399\) lines, \(15031\) bytes, and SHA--256
\[
\texttt{4F086F6B075AD55847D990BBABE6262871E7F5C51B015AA30807AB9482AAC9E0}.
                                                               \tag{9.15}
\]
Its canonical payload has SHA--256
\[
\texttt{8EC89A46BA672C1BB434B2B27FB8F0CC8131FC7ADA8B9B3B7C2F2B9F0E1DD5D3}.
                                                               \tag{9.16}
\]
The V4 inverse compiler and V60 response-oracle input hashes are,
respectively,
\[
\begin{gathered}
\texttt{309FDFB53F02708FE82720E6D4491732C}\\
\texttt{824F41685F98AD69C89CF5490401626},\\[3pt]
\texttt{7D5BFEE7FB834DFB5132FA2E1A9C5BC3}\\
\texttt{CED7D57D366C557C8F370B6AD24D6FC8}.
\end{gathered}                                                 \tag{9.17}
\]
\end{record}

\begin{firewall}[What the grid does not prove]
\label{fire:v5-grid}
The quarter-turn calculation validates identities and diagnoses
cancellation at a finite set.  It supplies neither an outward-rounded
continuum enclosure nor the continuous Haar-mean sign.  In particular,
the values in (9.13) cannot be used as a positive-measure edge class
or as any part of a continuum spectral-gap proof.
\end{firewall}

\subsection{V6 programme and current boundary}

\begin{decision}[Differentiate the compiled scalar before enclosing]
\label{dec:v5-v6}
V6 should implement (9.11) directly.  The primary object is the
single bubble jet \(\Psi_e''(0)\), with the diagonal alias part and
the signed \(3\times3\) Schur correction accumulated together.
The first certificate target is a coefficient-exact or
outward-rounded cell enclosure of
\[
 \frac12\Tr(K_e\widehat Q_2)+a_e\Psi_e''(0),
                                                               \tag{9.18}
\]
followed by the endpoint-minus-Wilson difference.  The endpoint
second/third-word cancellation in (9.13) must remain internal to
\(\Psi_e''(0)\).

If direct interval differentiation of \(\Lambda(k)\) is still too
wide, go upstream to the high-alias extension (8.34) and
differentiate its \(3\times3\) Woodbury middle matrix.  Do not return
to a triangle bound over seven words, 46 subsets, or sixteen moving
kernel blocks.
\end{decision}

\begin{assessment}[V5 progress]
\label{ass:v5-progress}
V5 completes the scheduled companion audit, proves an exact
rank-three subset compiler for every constrained response word, and
proves that the entire seven-word bubble is one second derivative.
The result lowers the next rigorous numerical target from a
\(45\times45\) inverse and seven separately bounded words to one
complete scalar jet whose only action-dependent non-diagonal inverse
has size \(3\times3\).  The exhaustive grid check validates that
compiler and reveals the cancellation which any successful enclosure
must preserve.
\end{assessment}

\begin{firewall}[Claim boundary after V5]
\label{fire:v5-final}
V5 does not prove the continuous Haar-mean sign, a nonzero Gaussian
edge class, or any interacting scale estimate.  The compensating
scale boundary, SCIC, ADMC, cofinal state, regulator removal,
Osterwalder--Schrader reconstruction, physical clustering, and the
four-dimensional Yang--Mills mass gap all remain open.
\end{firewall}

% =============================================================================
% END APPENDED POST-TENTH-ARCHIVE V5 MATERIAL
% =============================================================================

% =============================================================================
% BEGIN APPENDED POST-TENTH-ARCHIVE V6 MATERIAL
% =============================================================================

\section{V6: toron-uniform high-alias bounds and a single jet algebra}
\label{sec:v6-high-alias-jet}

V5 showed that the seven bubble words form one second derivative but
did not yet give a uniformly conditioned coordinate system in which
to differentiate that scalar.  This continuation moves to the
fundamental coarse cube, proves explicit continuum bounds for the
high-alias extension and both propagators, and evaluates the complete
bubble inside the second-order truncated jet algebra.  The resulting
representation has no massless low-alias inverse.

\subsection{Elementary bounds on the fundamental cube}

Work on \({\cal K}=[-\pi,\pi]^4\).  Put
\[
 z_\mu=e^{ik_\mu/2},\qquad x_\mu=\cos(k_\mu/2)\ge0,\qquad
 s_{\epsilon,\mu}=(-1)^{\epsilon_\mu}.                       \tag{10.1}
\]
Then the quantities in (8.4) obey
\[
 q_{\epsilon,\mu}=s_{\epsilon,\mu}z_\mu-1,\qquad
 b_{\epsilon,\mu}=1+s_{\epsilon,\mu}z_\mu,\qquad
 \lambda_\epsilon
 =2\sum_{\mu=1}^4(1-s_{\epsilon,\mu}x_\mu).                  \tag{10.2}
\]

\begin{lemma}[Low-link floor and collective high-alias bound]
\label{lem:v6-elementary-alias}
For \(k\in{\cal K}\),
\[
 \lambda_\epsilon\ge2\quad(\epsilon\ne0),\qquad
 |b_{0,\mu}|^2=2+2x_\mu\ge2.                                \tag{10.3}
\]
Moreover, for each coordinate \(\mu\),
\[
 \sum_{\epsilon\ne0}|b_{\epsilon,\mu}|^2
 =32-|b_{0,\mu}|^2\le30.                                   \tag{10.4}
\]
\end{lemma}

\begin{proof}
If \(\epsilon\ne0\), choose \(\mu\) with
\(s_{\epsilon,\mu}=-1\).  Its summand in (10.2) is
\(2(1+x_\mu)\ge2\), while all other summands are nonnegative.
The second assertion in (10.3) follows directly from
\(|1+z_\mu|^2=2+2x_\mu\).

For fixed \(\mu\), eight aliases have \(s_{\epsilon,\mu}=1\)
and eight have \(s_{\epsilon,\mu}=-1\).  Therefore
\[
 \sum_{\epsilon}|b_{\epsilon,\mu}|^2
 =8\bigl(|1+z_\mu|^2+|1-z_\mu|^2\bigr)=32.
\]
Remove the \(\epsilon=0\) term and use (10.3) to obtain (10.4).
\end{proof}

\subsection{A uniform extension theorem}

Retain the bases \(J_\epsilon,J\), matrices \(A_0,A_h\), and
\(L=A_0^{-1}A_h\) from (8.30)--(8.31).  All these bases are
isometries onto the stated transverse spaces.

\begin{lemma}[Uniform low-link inverse]
\label{lem:v6-low-link}
For \(k\ne0\),
\[
 s_{\min}(A_0)\ge\min_\mu|b_{0,\mu}|\ge\sqrt2,\qquad
 \|A_0^{-1}\|_{\rm op}\le\frac1{\sqrt2}.                    \tag{10.5}
\]
\end{lemma}

\begin{proof}
Let \(y\in\mathbb C^3\), put \(a=J_0y\), and recall
\(Q=D_0q_0\).  Since \(a\perp q_0\),
\[
\begin{split}
 \|A_0y\|
 &=\|P_{Q^\perp}D_0a\|
   =\inf_{\alpha\in\mathbb C}\|D_0a-\alpha Q\|\\
 &=\inf_{\alpha\in\mathbb C}\|D_0(a-\alpha q_0)\|\\
 &\ge\min_\mu|b_{0,\mu}|
       \inf_{\alpha\in\mathbb C}\|a-\alpha q_0\|
  =\min_\mu|b_{0,\mu}|\,\|a\|.
\end{split}                                                   \tag{10.6}
\]
Because \(J_0\) is isometric, \(\|a\|=\|y\|\).  Apply (10.3).
\end{proof}

\begin{theorem}[Collective extension bound]
\label{thm:v6-extension}
For \(k\ne0\),
\[
 \|A_h\|_{\rm op}\le\sqrt{30},\qquad
 \|L\|_{\rm op}\le\sqrt{15}.                                \tag{10.7}
\]
The high-alias extension
\[
 {\cal E}_h y_h:=
 \bigl(-J_0Ly_h,\,(J_\epsilon y_\epsilon)_{\epsilon\ne0}\bigr)
                                                               \tag{10.8}
\]
satisfies
\[
 {\cal E}_h^*{\cal E}_h=I+L^*L,\qquad
 \|{\cal E}_h\|_{\rm op}^2\le16.                            \tag{10.9}
\]
At the toron, define
\({\cal E}_h y_h=(0,(J_\epsilon y_\epsilon)_{\epsilon\ne0})\);
then \(\|{\cal E}_h\|_{\rm op}=1\).
\end{theorem}

\begin{proof}
Write \(A_h=[A_\epsilon]_{\epsilon\ne0}\), where
\(A_\epsilon=J^*D_\epsilon J_\epsilon\).  Since
\(J_\epsilon J_\epsilon^*=P_\epsilon\),
\[
\begin{split}
 A_hA_h^*
 &=\sum_{\epsilon\ne0}
 J^*D_\epsilon P_\epsilon D_\epsilon^*J\\
 &\preceq
 J^*\left(\sum_{\epsilon\ne0}D_\epsilon D_\epsilon^*\right)J
 \preceq30I_3.
\end{split}                                                   \tag{10.10}
\]
The last inequality is exactly the diagonal identity (10.4).
This proves the first bound in (10.7); Lemma~\ref{lem:v6-low-link}
gives the second.  The low and high alias blocks in (10.8) are
orthogonal, so direct multiplication gives the first identity in
(10.9), and (10.7) gives the second.  The toron statement follows
from the declared isometric high block.
\end{proof}

\subsection{Uniform propagators and endpoint ordering}

In the coordinates (10.8), Theorem~\ref{thm:v4-stable-woodbury}
gives
\[
\begin{split}
 N_{\rm W}
 &=\bigoplus_{\epsilon\ne0}\frac{\lambda_\epsilon}{3}I_3
   +\frac{\lambda_0}{3}L^*L,\\
 N_*
 &=\bigoplus_{\epsilon\ne0}
   \lambda_\epsilon\left(\frac13+\lambda_\epsilon\right)I_3
   +\lambda_0\left(\frac13+\lambda_0\right)L^*L.
\end{split}                                                   \tag{10.11}
\]

\begin{theorem}[Toron-uniform inverse bounds]
\label{thm:v6-uniform-inverse}
On the entire fundamental cube, with the toron interpreted as in
(8.37),
\[
 N_{\rm W}\succeq\frac23I,\qquad
 N_*\succeq\frac{14}{3}I,                                   \tag{10.12}
\]
and hence
\[
 \|N_{\rm W}^{-1}\|_{\rm op}\le\frac32,\qquad
 \|N_*^{-1}\|_{\rm op}\le\frac3{14}.                         \tag{10.13}
\]
The ambient constrained covariances are
\[
 K_e={\cal E}_hN_e^{-1}{\cal E}_h^*,                         \tag{10.14}
\]
and obey the explicit global bounds
\[
 \boxed{\quad
 \|K_{\rm W}\|_{\rm op}\le24,\qquad
 \|K_*\|_{\rm op}\le\frac{24}{7}.
 \quad}                                                       \tag{10.15}
\]
\end{theorem}

\begin{proof}
Every high alias has \(\lambda_\epsilon\ge2\) by
Lemma~\ref{lem:v6-elementary-alias}.  The diagonal parts of
(10.11) are therefore bounded below by \(2/3\) and \(14/3\);
the low-alias updates are positive semidefinite.  This proves
(10.12)--(10.13).

Equation (10.14) is the covariance associated with the unique
constraint parametrization (8.34), and also follows by substituting
that parametrization into the quadratic form.  Combine (10.13)
with \(\|{\cal E}_h\|^2\le16\) from
Theorem~\ref{thm:v6-extension}.  At the toron the extension norm is
one, so the same global bounds remain valid.
\end{proof}

\begin{theorem}[Exact endpoint covariance order]
\label{thm:v6-endpoint-order}
For every \(k\in{\cal K}\),
\begin{align}
 N_*-N_{\rm W}
 &=
 \bigoplus_{\epsilon\ne0}\lambda_\epsilon^2I_3
 +\lambda_0^2L^*L\succeq0,                                  \tag{10.16}\\
 N_{\rm W}^{-1}-N_*^{-1}
 &=N_{\rm W}^{-1}(N_*-N_{\rm W})N_*^{-1}\succeq0.            \tag{10.17}
\end{align}
Consequently,
\[
 0\preceq K_{\rm W}-K_*
 ={\cal E}_h(N_{\rm W}^{-1}-N_*^{-1}){\cal E}_h^*.           \tag{10.18}
\]
\end{theorem}

\begin{proof}
Subtract the two lines of (10.11), using
\(\lambda(1/3+\lambda)-\lambda/3=\lambda^2\).
The order-reversing property of inversion on positive definite
matrices proves the last inequality in (10.17); the displayed
resolvent identity proves its equality.  Congruence by
\({\cal E}_h\) gives (10.18).  The toron is included because the
\(\lambda_0^2L^*L\) term is absent there.
\end{proof}

\begin{remark}[The common massless pole cancels algebraically]
\label{rem:v6-pole}
At the direct-alias level,
\[
 d_*(\lambda)-d_{\rm W}(\lambda)
 =\frac1{\lambda(1/3+\lambda)}-\frac3\lambda
 =-\frac9{1+3\lambda}.                                      \tag{10.19}
\]
Thus the endpoint difference has a continuous direct-alias remainder
at \(\lambda=0\).  Equations (10.11)--(10.18) are the
constraint-preserving version of the same cancellation.
The covariance order alone does not determine the response sign,
because the two actions have different cubic and quartic banks.
\end{remark}

\subsection{Congruence into high-alias coordinates}

Let \(T\) and \(\Lambda\) be as in (8.38), and define the square map
\[
 Z:=\Lambda{\cal E}_h:\mathbb C^{45}\longrightarrow
 \mathbb C^{45}.                                             \tag{10.20}
\]

\begin{theorem}[High-alias trace compiler]
\label{thm:v6-high-trace}
The map \(Z\) is invertible, \(TZ={\cal E}_h\), and
\[
 Z^*H_eZ=N_e.                                                \tag{10.21}
\]
For every reduced factor \(X\), put
\(\widetilde X=Z^*XZ\).  Then
\[
 \Tr(H_e^{-1}X_1\cdots H_e^{-1}X_m)
 =
 \Tr(N_e^{-1}\widetilde X_1\cdots
          N_e^{-1}\widetilde X_m).                           \tag{10.22}
\]
\end{theorem}

\begin{proof}
The operator \(T\Lambda\) is the orthogonal projector onto
\(\operatorname{Ran}T\), and
\(\operatorname{Ran}{\cal E}_h=\operatorname{Ran}T\).
Hence \(TZ=T\Lambda{\cal E}_h={\cal E}_h\).
Both \(T\) and \({\cal E}_h\) are injective maps from
\(\mathbb C^{45}\) onto that range, so \(Z\) is invertible.
Now
\[
 Z^*H_eZ
 =Z^*T^*{\cal M}_eTZ
 ={\cal E}_h^*{\cal M}_e{\cal E}_h=N_e,
\]
which proves (10.21).  It implies
\(H_e^{-1}=ZN_e^{-1}Z^*\).  Substitute this identity into the
left side of (10.22), move the initial \(Z\) through the trace,
and collect \(Z^*X_jZ=\widetilde X_j\).
\end{proof}

\subsection{One truncated jet instead of seven words}

For an associative matrix algebra \({\cal A}\), define
\[
 {\mathfrak J}_2({\cal A})={\cal A}[\tau]/(\tau^3).
\]
Represent an element by its ordinary coefficients
\({\bf A}=(A_0,A_1,A_2)\), and multiply by convolution:
\[
 {\bf A}\odot{\bf B}
 =\bigl(A_0B_0,\,
 A_0B_1+A_1B_0,\,
 A_0B_2+A_1B_1+A_2B_0\bigr).                                \tag{10.23}
\]

\begin{lemma}[Inverse in the second-jet algebra]
\label{lem:v6-jet-inverse}
For
\[
 {\bf H}=(N_e,\widetilde H_1,\tfrac12\widetilde H_2)
\]
the unique inverse
\({\bf G}={\bf H}^{-1}=(G_0,G_1,G_2)\) is
\[
\begin{split}
 G_0&=N_e^{-1},\\
 G_1&=-G_0\widetilde H_1G_0,\\
 G_2&=G_0\widetilde H_1G_0\widetilde H_1G_0
      -\frac12G_0\widetilde H_2G_0.
\end{split}                                                   \tag{10.24}
\]
\end{lemma}

\begin{proof}
Equate the coefficients of \(1,\tau,\tau^2\) in
\({\bf H}\odot{\bf G}=(I,0,0)\) and solve successively.
\end{proof}

\begin{theorem}[Single-jet response formula]
\label{thm:v6-single-jet}
Set
\[
\begin{split}
 {\bf P}&=(\widetilde P_0,\widetilde P_1,
                    \tfrac12\widetilde P_2),\\
 {\bf M}&=(\widetilde M_0,\widetilde M_1,
                    \tfrac12\widetilde M_2),\\
 {\bf G}_{\rm out}&=(G_0,0,0),
\end{split}                                                   \tag{10.25}
\]
and extend the trace coefficientwise to
\({\mathfrak J}_2\).  The complete endpoint response is
\[
\boxed{\quad
 F_e(k)=\frac12\Tr(G_0\widetilde Q_2)
 +2c_e[\tau^2]\Tr\!\left(
 {\bf G}_{\rm out}\odot{\bf P}\odot{\bf G}\odot{\bf M}
 \right),
\quad}                                                        \tag{10.26}
\]
where \(c_{\rm W}=1/3\), \(c_*=3\).
\end{theorem}

\begin{proof}
The second coefficient in (10.24) is one half of the second
derivative of the inverse.  Therefore twice the \(\tau^2\)
coefficient in (10.26) is precisely the product-rule expansion
in (9.10).  Substitution of (10.24) gives all seven terms of
(9.9) with their signs and multiplicities.  Theorem
\ref{thm:v6-high-trace} transfers those terms without changing
their traces, and the first term in (10.26) is the tadpole.
\end{proof}

\begin{remark}[What has improved]
\label{rem:v6-improvement}
Equation (10.26) is evaluated as one coefficient in an associative
algebra; no list of seven response words is needed.  Its only base
inverse is \(N_e^{-1}\), with the uniform bounds (10.13).  The
low-alias massless pole and the arbitrary ambient transverse-frame
conditioning have disappeared from the inversion step.
\end{remark}

\subsection{Executable interface check}

\begin{record}[High-alias jet verifier]
\label{rec:v6-verifier}
The executable
\begin{center}
\path{scripts/verify_ym_postarchive_v6_high_alias_jet.py}
\end{center}
checks the coordinate identity (10.21), the single-jet response
(10.26), both uniform inverse gates, and the endpoint orders
(10.16)--(10.18) on all \(256\) weighted quarter-turn momenta,
wrapped into \({\cal K}\).  The largest floating residuals are
\[
\begin{array}{c|cc}
 &{\rm Wilson}&*\\ \hline
 Z^*H_eZ-N_e
 &1.0510\times10^{-14}&1.2088\times10^{-14}\\
 F_e^{\rm jet}-F_e^{\rm oracle}
 &4.3779\times10^{-14}&3.1059\times10^{-14}.
\end{array}                                                   \tag{10.27}
\]
The extension reconstruction residual is
\(6.94\times10^{-15}\).  The sampled extrema attain
\[
 \min_{\epsilon\ne0}\lambda_\epsilon=2,\quad
 \min s_{\min}(A_0)=\sqrt2,\quad
 \max\|L\|=\sqrt{15},\quad
 \max\|{\cal E}_h\|^2=16,                                   \tag{10.28}
\]
up to floating roundoff, thereby exercising the sharp faces of the
proved elementary bounds.  The response means reproduce
\[
 \langle F_{\rm W}\rangle_4\doteq222.7064519320,\qquad
 \langle F_*\rangle_4\doteq-92.1377272462.                   \tag{10.29}
\]
These last values remain diagnostics; the exact rational V2 card is
the rigorous finite-grid sign result.

The source has \(417\) lines, \(15196\) bytes, and SHA--256
\[
\texttt{7D7C63161E1AC9260A5F798EDE3BDBE092A2BCD8692658A9CAA426BF46B2C26D}.
                                                               \tag{10.30}
\]
Its canonical payload has SHA--256
\[
\texttt{5D6B22A979918E704E22D9EB988F16282A3BD7E8842968C3CEE823DBD9C71ADE}.
                                                               \tag{10.31}
\]
\end{record}

\subsection{V7 programme and claim boundary}

\begin{decision}[Certificate the transformed banks, not the inverse]
\label{dec:v6-v7}
The inverse part of the continuum cell problem is now uniformly
conditioned by (10.12)--(10.15).  V7 should construct a finite,
coefficient-exact chart atlas for \(Z(k)\) and the transformed banks
\[
 \widetilde H_j,\quad\widetilde P_j,\quad
 \widetilde M_j,\quad\widetilde Q_2,                          \tag{10.32}
\]
then propagate one outward-rounded \({\mathfrak J}_2\) enclosure
through (10.26).  The certificate must combine endpoint and Wilson
responses before its final absolute value and must target the exact
quadrature margin \(\delta_4\) in (6.26).

If orthonormal-frame charts make (10.32) too wide, derive the
transformed banks directly from the retained-edge incidence and the
explicit extension equation \(TZ={\cal E}_h\).  Do not weaken
(10.13) by returning to the global reduced Hessian floor, and do not
split the jet coefficient back into seven words.
\end{decision}

\begin{assessment}[V6 progress]
\label{ass:v6-progress}
V6 proves global, toron-safe inverse bounds in the correct constrained
coordinates, proves the endpoint covariance order, cancels the common
direct massless pole explicitly, and replaces the seven-word bubble by
one truncated-jet coefficient.  These results remove inversion
conditioning as the immediate Haar-enclosure bottleneck.  The live
finite-stage obstruction is now the coefficient-exact control of the
common coordinate map and transformed vertex banks over finitely many
momentum charts.
\end{assessment}

\begin{record}[V6 append record]
\label{rec:v6-append}
V6 is appended to the validated V5 whose installed SHA--256 was
\[
\texttt{53D900C8A36175A6702FA4D1AC2D65DD748D7C7988BC389967CD1E16E008D89A}.
                                                               \tag{10.33}
\]
After structural validation and compilation only V6 is to remain the
active numeric YM note; the frozen archive files are unchanged.
\end{record}

\begin{firewall}[Claim boundary after V6]
\label{fire:v6-final}
V6 does not prove the continuous Haar-mean sign, a nonzero Gaussian
edge class, SCIC, ADMC, an interacting cofinal state, regulator
removal, reflection-positive continuum reconstruction, physical
clustering, or a Hamiltonian spectral gap.  In particular, covariance
order (10.18) is not a response-sign theorem because the endpoint
vertex banks differ.
\end{firewall}

% =============================================================================
% END APPENDED POST-TENTH-ARCHIVE V6 MATERIAL
% =============================================================================

% =============================================================================
% BEGIN APPENDED POST-TENTH-ARCHIVE V7 MATERIAL
% =============================================================================

\section{V7: a rational pivot atlas and an exact toron-collar budget}
\label{sec:v7-pivot-collar}

V6 removed the massless inverse from the response but still used
numerically chosen orthonormal transverse frames.  Such frames are
poor objects for a coefficient-exact interval certificate: their
phases may jump and their normalization introduces unnecessary square
roots.  This continuation replaces them by four rational pivot charts
on the punctured fundamental cube.  It also isolates a small toron
collar and pays for its entire integral with an exact rational bound.
The complement is consequently compact, finitely charted, and has all
chart denominators separated from zero.

\subsection{Rational elimination frames}

Let \(v\in\mathbb C^4\) and suppose \(v_r\ne0\).  Enumerate the three
indices different from \(r\) as \(j_1,j_2,j_3\), and define
\(F_r(v):\mathbb C^3\to\mathbb C^4\) by
\[
 [F_r(v)y]_{j_a}=y_a,\qquad
 [F_r(v)y]_r
 =-\sum_{a=1}^3\frac{\overline{v_{j_a}}}{\overline{v_r}}y_a.
                                                               \tag{11.1}
\]

\begin{lemma}[Elimination-frame identities]
\label{lem:v7-elimination}
The range of \(F_r(v)\) is \(v^\perp\), and
\[
 F_r(v)^*F_r(v)=I_3+\eta^*\eta,\qquad
 \eta_a=\frac{\overline{v_{j_a}}}{\overline{v_r}}.            \tag{11.2}
\]
Consequently
\[
 s_{\min}(F_r(v))=1,\qquad
 \|F_r(v)\|_{\rm op}^2
 =1+\sum_{j\ne r}\frac{|v_j|^2}{|v_r|^2},                   \tag{11.3}
\]
and the orthogonal projector onto \(v^\perp\) is
\[
 P_v=F_r(v)\bigl(F_r(v)^*F_r(v)\bigr)^{-1}F_r(v)^*.          \tag{11.4}
\]
\end{lemma}

\begin{proof}
Substitution in (11.1) gives \(v^*F_r(v)y=0\).  The map is
injective and both its range and \(v^\perp\) have dimension three,
so the ranges agree.  Direct multiplication gives (11.2).
The rank-one matrix \(\eta^*\eta\) has eigenvalues
\(0,0,\|\eta\|^2\), which proves (11.3).  Formula (11.4) is the
standard projector associated with a full-column-rank frame.
\end{proof}

On the unit torus, complex conjugation sends \(z_\mu\) to
\(z_\mu^{-1}\).  Thus, when \(v\) has Laurent-polynomial entries,
all entries in (11.1)--(11.4) are rational Laurent functions.
No square root is introduced.

\subsection{Four cones cover the punctured cube}

For \(0<\rho<\pi\), put
\[
 {\cal C}_\rho:=\{k\in{\cal K}:\|k\|_\infty<\rho\},\qquad
 {\cal U}_{r,\rho}:=
 \{k\in{\cal K}:|k_r|=\|k\|_\infty\ge\rho\}.                 \tag{11.5}
\]
Ties may be assigned to the smallest eligible index.

\begin{theorem}[Four-chart rational transverse atlas]
\label{thm:v7-four-chart}
The four sets \({\cal U}_{r,\rho}\) cover
\({\cal K}\setminus{\cal C}_\rho\).  On
\({\cal U}_{r,\rho}\), both
\[
 F_{0,r}:=F_r(q_0),\qquad C_r:=F_r(Q)                       \tag{11.6}
\]
are defined and obey
\[
 \|F_{0,r}\|_{\rm op}\le2,\qquad
 \|C_r\|_{\rm op}\le2.                                     \tag{11.7}
\]
Their pivot denominators satisfy the explicit bounds
\[
 |q_{0,r}|\ge\frac{\rho}{\pi}>\frac{\rho}{4},\qquad
 |Q_r|\ge\frac{2\rho}{\pi}>\frac{\rho}{2}.                  \tag{11.8}
\]

For each high alias \(\epsilon\ne0\), fix once and for all the
smallest index \(r_\epsilon\) with \(\epsilon_{r_\epsilon}=1\)
and set \(F_\epsilon=F_{r_\epsilon}(q_\epsilon)\).  Then
\[
 |q_{\epsilon,r_\epsilon}|\ge\sqrt2,\qquad
 \|F_\epsilon\|_{\rm op}\le\sqrt7                          \tag{11.9}
\]
throughout \({\cal K}\).
\end{theorem}

\begin{proof}
The covering assertion is immediate from the definition of the
\(\ell^\infty\) norm.  On \([0,\pi]\), both
\(\sin(t/4)\) and \(\sin(t/2)\) are increasing.  Hence a coordinate
maximizing \(|k_\mu|\) also maximizes
\[
 |q_{0,\mu}|=2\sin(|k_\mu|/4),\qquad
 |Q_\mu|=2\sin(|k_\mu|/2).
\]
Every ratio in (11.3) therefore has modulus at most one, proving
(11.7).  The elementary inequality
\(\sin x\ge2x/\pi\) for \(0\le x\le\pi/2\) gives the first
inequalities in (11.8); the second use \(\pi<4\).

For the fixed high pivot,
\[
 q_{\epsilon,r_\epsilon}=-z_{r_\epsilon}-1,
\]
whose modulus is at least \(\sqrt2\) on \({\cal K}\).
Every other component of \(q_\epsilon\) has modulus at most two.
Thus each squared ratio in (11.3) is at most two, and there are
three of them.  This proves (11.9).
\end{proof}

\subsection{A coefficient-exact nonorthonormal compiler}

On \({\cal U}_{r,\rho}\), define
\[
\begin{split}
 A_0^{(r)}&:=C_r^*D_0F_{0,r},\\
 A_h^{(r)}&:=
 [\,C_r^*D_\epsilon F_\epsilon\,]_{\epsilon\ne0},\\
 L_r&:=(A_0^{(r)})^{-1}A_h^{(r)}.
\end{split}                                                   \tag{11.10}
\]

\begin{lemma}[The rational low link is uniformly invertible]
\label{lem:v7-rational-low-link}
For every \(k\in{\cal U}_{r,\rho}\),
\[
 s_{\min}(A_0^{(r)})\ge\min_\mu|b_{0,\mu}|\ge\sqrt2.         \tag{11.11}
\]
\end{lemma}

\begin{proof}
Let \(a=F_{0,r}y\in q_0^\perp\).  The proof of
Lemma~\ref{lem:v6-low-link} gives
\[
 \|P_{Q^\perp}D_0a\|
 \ge\min_\mu|b_{0,\mu}|\,\|a\|.
\]
The nonzero singular values of \(C_r^*:Q^\perp\to\mathbb C^3\)
are those of \(C_r\), all at least one by
Lemma~\ref{lem:v7-elimination}.  The same lemma gives
\(\|a\|\ge\|y\|\).  Combining these inequalities with (10.3)
proves (11.11).
\end{proof}

Let \(F_h=\bigoplus_{\epsilon\ne0}F_\epsilon\) and define
\[
 {\cal E}_{r}y_h=
 \bigl(-F_{0,r}L_ry_h,\,F_hy_h\bigr).                        \tag{11.12}
\]
The corresponding high-coordinate Hessian is
\[
 N_e^{(r)}
 =
 \bigoplus_{\epsilon\ne0}
 w_e(\lambda_\epsilon)F_\epsilon^*F_\epsilon
 +w_e(\lambda_0)L_r^*F_{0,r}^*F_{0,r}L_r.                   \tag{11.13}
\]
Because every frame Gram matrix is at least the identity,
\[
 N_{\rm W}^{(r)}\succeq\frac23I,\qquad
 N_*^{(r)}\succeq\frac{14}{3}I.                             \tag{11.14}
\]

\begin{theorem}[Uniform retained-coordinate Gram floor]
\label{thm:v7-transverse-gram}
The transverse retained-edge frame \(T=PV\) satisfies
\[
 T^*T\succeq\frac{147}{20000}I_{45},\qquad
 \|\Lambda\|_{\rm op}
 \le\sqrt{\frac{20000}{147}}.                               \tag{11.15}
\]
\end{theorem}

\begin{proof}
The already proved full-torus Wilson floor, recorded in (1.2), is
\[
 H_{\rm W}\succeq\frac{49}{1250}I.
\]
Every alias obeys \(\lambda_\epsilon\le16\), so on the transverse
alias space
\[
 {\cal M}_{\rm W}
 =\bigoplus_\epsilon\frac{\lambda_\epsilon}{3}P_\epsilon
 \preceq\frac{16}{3}P.
\]
Consequently
\[
 \frac{49}{1250}I
 \preceq H_{\rm W}
 =T^*{\cal M}_{\rm W}T
 \preceq\frac{16}{3}T^*T,
\]
which rearranges to the first inequality in (11.15).
The second follows from
\(\|\Lambda\|=s_{\min}(T)^{-1}\).
\end{proof}

\begin{theorem}[Rational chart realization of the complete jet]
\label{thm:v7-rational-jet}
On every \({\cal U}_{r,\rho}\), put
\[
 Z_r:=\Lambda{\cal E}_r,\qquad
 \widetilde X^{(r)}:=Z_r^*XZ_r.                             \tag{11.16}
\]
Then \(Z_r\) is invertible,
\[
 Z_r^*H_eZ_r=N_e^{(r)},                                     \tag{11.17}
\]
The single-jet response formula (10.26) holds in this chart after
the substitutions
\[
 N_e\longmapsto N_e^{(r)},\qquad
 \widetilde X\longmapsto\widetilde X^{(r)}.
\]

Every entry of \(Z_r,N_e^{(r)}\), and each transformed vertex bank
is a rational Laurent function of
\(z_1,\ldots,z_4\), with rational coefficients inherited from the
V60 banks.  Its only low and coarse pivot denominators satisfy
(11.8), and all required finite matrix inverses have the positive
floors (11.11), (11.14), and (11.15).
\end{theorem}

\begin{proof}
Equations (11.10)--(11.12) solve the exact coarse-link constraint,
so \(\operatorname{Ran}{\cal E}_r=\operatorname{Ran}T\).
The proof of Theorem~\ref{thm:v6-high-trace} applies verbatim to
the nonorthonormal injective frame and gives invertibility of \(Z_r\),
(11.17), and trace invariance.  The convolution proof of
Theorem~\ref{thm:v6-single-jet} is basis independent, so it gives
the last response assertion.

By Lemma~\ref{lem:v7-elimination}, every projector can be written
as \(F(F^*F)^{-1}F^*\).  On the torus,
\(\overline z_\mu=z_\mu^{-1}\), so these projectors, \(T\),
\(\Lambda=(T^*T)^{-1}T^*\), (11.10), and (11.12) are obtained by
finite rational matrix operations from Laurent functions.  The V60
banks have rational Laurent coefficients.  Congruence by \(Z_r\)
therefore preserves rational Laurent dependence.  The denominator
and inverse floors are exactly the cited estimates.
\end{proof}

\begin{firewall}[Rational does not yet mean interval-certified]
\label{fire:v7-rational-scope}
Theorem~\ref{thm:v7-rational-jet} constructs a finite exact atlas and
proves that all its inversions are separated from singularity on the
complement of the collar.  It does not yet supply outward-rounded
ranges for the rational functions on a subdivision.  In particular,
naive interval evaluation may still suffer dependency loss even
though every true denominator is safe.
\end{firewall}

\subsection{Paying for the toron collar exactly}

For an action \(e\), let
\[
 \gamma_{\rm W}=\frac{49}{1250},\qquad
 \gamma_*=\frac{343}{1250},\qquad g_e=\gamma_e^{-1}.         \tag{11.18}
\]
Let \(h_1,h_2,p_0,p_1,p_2,m_0,m_1,m_2,q_2\) denote rational
upper bounds on the full-torus Frobenius norms of the corresponding
V60 Laurent banks.  They are obtained by summing, over Laurent
exponents, rational upper bounds on the Frobenius norm of each
coefficient matrix.  Define
\[
\begin{split}
 {\cal M}_e:={}&
 \frac{\sqrt{45}^{\,\uparrow}}2\,g_eq_2
 +|c_e|\bigl[
 g_e^2p_2m_0
 +2g_e^4p_0h_1^2m_0
 +g_e^3p_0h_2m_0\\
 &\quad
 +g_e^2p_0m_2
 +2g_e^3p_1h_1m_0
 +2g_e^2p_1m_1
 +2g_e^3p_0h_1m_1
 \bigr].
\end{split}                                                   \tag{11.19}
\]
Here \(\sqrt{45}^{\,\uparrow}\) is the rational upper square root
used by the exact coefficient-stock compiler.

\begin{lemma}[Global pointwise response envelope]
\label{lem:v7-global-envelope}
For every \(k\in{\cal K}\),
\[
 |F_e(k)|\le{\cal M}_e.                                     \tag{11.20}
\]
\end{lemma}

\begin{proof}
The floors (11.18) give \(\|H_e^{-1}\|_{\rm op}\le g_e\).
For the tadpole use
\[
 |\Tr(GQ_2)|\le\|G\|_{\rm F}\|Q_2\|_{\rm F}
 \le\sqrt{45}\,g_eq_2.
\]
For a bubble word, use the Frobenius norm on its first and last
vertex factors and the operator norm, bounded by the Frobenius norm,
on every intermediate vertex factor.  Every propagator contributes
\(g_e\).  Applying this rule to the seven signed words in (6.12)
and then taking the triangle inequality gives exactly (11.19).
The Laurent coefficient stocks dominate the corresponding
full-torus Frobenius norms by the triangle inequality.
\end{proof}

\begin{theorem}[Exact toron-collar budget]
\label{thm:v7-collar}
The coefficient-exact evaluation of (11.19) gives
\[
\begin{split}
 {\cal M}_{\rm W}&\doteq9.2170356413\times10^9,\\
 {\cal M}_*&\doteq9.4055993533\times10^{12},
\end{split}                                                   \tag{11.21}
\]
and proves the rational inequality
\[
 {\cal M}_{\rm W}+{\cal M}_*
 <9415000000000.                                            \tag{11.22}
\]
For \(\rho=10^{-3}\),
\[
\boxed{\quad
 \left|
 \frac1{(2\pi)^4}\int_{{\cal C}_\rho}
 (F_*(k)-F_{\rm W}(k))\,dk
 \right|
 <\frac{1883}{16200}
 <1<\frac{\delta_4}{4}.
\quad}                                                       \tag{11.23}
\]
\end{theorem}

\begin{proof}
All Laurent coefficients and both Hessian floors in (11.19) are
rational.  The rational upper-square-root routine squares every
declared upper bound before accepting it.  Exact fraction arithmetic
then proves (11.22); the decimals in (11.21) only report its scale.

The normalized volume of \({\cal C}_\rho\) is
\((\rho/\pi)^4<\rho^4/81\), since \(\pi>3\).  Lemma
\ref{lem:v7-global-envelope}, (11.22), and \(\rho=10^{-3}\)
give
\[
 \left|\langle(F_*-F_{\rm W})1_{{\cal C}_\rho}\rangle\right|
 <
 \frac{9415000000000}{81\cdot10^{12}}
 =\frac{1883}{16200}.
\]
Finally \(\delta_4/4=147981586239/2000000000>1\), proving
(11.23).
\end{proof}

\begin{corollary}[Reduced complement target]
\label{cor:v7-complement-target}
Put \(f=F_*-F_{\rm W}\) and define the complement remainder
\[
 E_{\rm comp}:=
 \langle f1_{{\cal K}\setminus{\cal C}_\rho}\rangle
 -\langle f\rangle_4.                                      \tag{11.27}
\]
To prove the Haar-sign condition (6.27), it is sufficient to certify
\(|E_{\rm comp}|<3\delta_4/4\), while using (11.23) for the
toron collar.  On that complement the four rational charts have the
denominator floors (11.8).
\end{corollary}

\begin{proof}
The Haar error is
\[
 \langle f\rangle-\langle f\rangle_4
 =\langle f1_{{\cal C}_\rho}\rangle+E_{\rm comp}.
\]
The collar consumes less than \(\delta_4/4\); a complement error
strictly below \(3\delta_4/4\) leaves total error below
\(\delta_4\).  The chart assertion is Theorem
\ref{thm:v7-four-chart}.
\end{proof}

\subsection{Executable checks and V8 programme}

\begin{record}[Pivot-atlas and collar verifier]
\label{rec:v7-verifier}
The executable
\begin{center}
\path{scripts/verify_ym_postarchive_v7_pivot_atlas_collar.py}
\end{center}
constructs the exact rational stocks in (11.19), checks the strict
fractions (11.22)--(11.23), and cross-checks the nonorthonormal
compiler at all \(256\) weighted quarter-turn momenta.  The largest
floating response residuals are
\[
 5.1080\times10^{-14}\quad({\rm Wilson}),\qquad
 2.4673\times10^{-14}\quad(*),                              \tag{11.24}
\]
and the largest projector and coordinate-reconstruction residuals
are \(3.60\times10^{-16}\) and \(6.96\times10^{-15}\).
The sampled frame norms reach the declared values \(2,2,\sqrt7\),
so the boundary cases of the atlas are exercised.

The source has \(478\) lines, \(16753\) bytes, and SHA--256
\[
\texttt{122A873318D06129211C5E4E8AF2C8326CC0973382A70AAE8AF8F3C341FCE3AF}.
                                                               \tag{11.25}
\]
Its canonical payload has SHA--256
\[
\texttt{1464E410A2CDC468BC138E1BCCBF6BB9F5A24B7AD1A8814E30114F035780AADA}.
                                                               \tag{11.26}
\]
\end{record}

\begin{decision}[V8: certify the charted complement]
\label{dec:v7-v8}
V8 should partition
\({\cal K}\setminus{\cal C}_{10^{-3}}\) by the deterministic pivot
rule in (11.5), subdivide each chart into boxes, and propagate
outward-rounded rational/trigonometric intervals through the single
jet (10.26).  Endpoint and Wilson values must be subtracted before
the final absolute value.  The exact complement budget is
\[
 \frac{3\delta_4}{4}
 =\frac{443944758717}{2000000000}.                           \tag{11.28}
\]
The implementation should first measure dependency loss on a small
cover.  If direct interval evaluation of
\((T^*T)^{-1}\) is too wide despite (11.15), use a midpoint
preconditioned residual inverse on each box.  If transformed vertex
congruences dominate, derive their ambient alias banks before
increasing the subdivision depth.
\end{decision}

\begin{assessment}[V7 progress]
\label{ass:v7-progress}
V7 removes the moving-frame ambiguity by constructing a four-chart
rational atlas with explicit denominator and inverse floors.  It also
proves an exact collar estimate which spends less than \(0.04\%\) of
the available finite-grid margin and leaves a compact,
uniformly nonsingular complement.  The live finite-stage bottleneck
is now an outward-rounded evaluation of the complete jet on that
charted complement, not toron singularity, gauge-frame choice, or
existence of a finite rational representation.
\end{assessment}

\begin{record}[V7 append record]
\label{rec:v7-append}
V7 is appended to the validated V6 whose installed SHA--256 was
\[
\texttt{26CF0AC3E4BDA30E3C7710BACDD49D2D1FDEDBC7A48205478E9245008B33EA2A}.
                                                               \tag{11.29}
\]
After validation only V7 is to remain the active numeric YM note.
\end{record}

\begin{firewall}[Claim boundary after V7]
\label{fire:v7-final}
V7 does not yet prove the complement quadrature bound, the continuous
Haar-mean sign, a nonzero Gaussian edge class, SCIC, ADMC, regulator
removal, reconstruction, clustering, or a Hamiltonian mass gap.  The
collar estimate is one rigorously paid part of the finite-stage
quadrature problem and no more.
\end{firewall}

% =============================================================================
% END APPENDED POST-TENTH-ARCHIVE V7 MATERIAL
% =============================================================================

% =============================================================================
% BEGIN APPENDED POST-TENTH-ARCHIVE V8 MATERIAL
% =============================================================================

\section{V8: the coupled difference jet and a convergent
live-phase cell certificate}
\label{sec:v8-coupled-cell}

\subsection{Why the global complement bound must again move upstream}

V7 separated the toron collar at exact cost less than
\(1883/16200\) and reduced the finite-stage problem to the compact
complement estimate
\[
 |E_{\rm comp}|<
 \frac{443944758717}{2000000000}.                            \tag{12.1}
\]
The first obligation of this iteration is to implement endpoint-minus-
Wilson subtraction before the last absolute value, rather than merely
to state it.  We do this in the second-order jet algebra.  The
resulting identity is exact and considerably cleaner than a difference
of two seven-word expansions.

The coefficient-stock computation then gives a useful negative result:
global Frobenius/operator stocks erase essentially all of that algebraic
cancellation.  Their tensor-trapezoid threshold remains above
\(3.9\times10^9\) points per axis.  Increasing a global grid would
therefore be circular.  The second part of V8 moves upstream to a
different finite certificate: freeze only the inverse on a cell, keep
every Laurent phase live, integrate the complete signed trial scalar,
and take a norm only on the inverse defect.

\subsection{One noncommutative jet for the endpoint difference}

Retain the algebra
\(\mathfrak J_2({\cal A})={\cal A}[\tau]/(\tau^3)\) and its product
\(\odot\) from (10.23).  For \(e\in\{{\rm W},*\}\), put
\[
\begin{aligned}
 {\bf H}_e&=(H_{e,0},H_{e,1},\tfrac12H_{e,2}),&
 {\bf G}_e&={\bf H}_e^{-1},\\
 {\bf O}_e&=(G_{e,0},0,0),&
 {\bf Q}_e&=(0,0,\tfrac12Q_{e,2}),\\
 {\bf P}_e&=(P_{e,0},P_{e,1},\tfrac12P_{e,2}),&
 {\bf M}_e&=(M_{e,0},M_{e,1},\tfrac12M_{e,2}).
\end{aligned}                                                 \tag{12.2}
\]
Thus (10.26) is equivalently
\[
 F_e=[\tau^2]\Tr\!\left(
 {\bf O}_e\odot{\bf Q}_e+
 {\bf O}_e\odot(2c_e{\bf P}_e)\odot
 {\bf G}_e\odot{\bf M}_e\right).                             \tag{12.3}
\]
For any object \(A\), write \(\Delta A=A_*-A_{\rm W}\).

\begin{lemma}[Jet resolvent identity]
\label{lem:v8-jet-resolvent}
In \(\mathfrak J_2\),
\[
 \boxed{\quad
 \Delta{\bf G}
 =-{\bf G}_*\odot\Delta{\bf H}\odot{\bf G}_{\rm W}.
 \quad}                                                       \tag{12.4}
\]
In particular,
\[
 \Delta G_0=-G_{*,0}(H_{*,0}-H_{{\rm W},0})G_{{\rm W},0}.
                                                                  \tag{12.5}
\]
\end{lemma}

\begin{proof}
The ordinary identity
\(A^{-1}-B^{-1}=-A^{-1}(A-B)B^{-1}\) holds in every unital
associative algebra whenever \(A\) and \(B\) are invertible.  Apply it
to \(A={\bf H}_*\) and \(B={\bf H}_{\rm W}\).  The constant
coefficient gives (12.5).
\end{proof}

There is a harmless scaling freedom in the bubble.  Choose nonzero
scalars \(s_e\), and define the four ordered factors
\[
 {\bf A}_{e,1}={\bf O}_e,\qquad
 {\bf A}_{e,2}=s_e{\bf P}_e,\qquad
 {\bf A}_{e,3}={\bf G}_e,\qquad
 {\bf A}_{e,4}=\frac{2c_e}{s_e}{\bf M}_e.                    \tag{12.6}
\]
Their product is independent of \(s_e\).

\begin{theorem}[Permutation telescoping of the complete bubble]
\label{thm:v8-permutation-telescope}
Let \(\pi\) be any permutation of \(\{1,\ldots,4\}\).
For \(1\le\ell\le4\), define the factor in position \(j\) by
\[
 {\bf B}_{\pi,\ell,j}:=
 \begin{cases}
  {\bf A}_{*,j},&
    j\in\{\pi(1),\ldots,\pi(\ell-1)\},\\
  \Delta{\bf A}_j,&j=\pi(\ell),\\
  {\bf A}_{{\rm W},j},&
    j\notin\{\pi(1),\ldots,\pi(\ell)\}.
 \end{cases}                                                  \tag{12.7}
\]
Then
\[
\boxed{\quad
 \prod_{j=1}^4{\bf A}_{*,j}
 -\prod_{j=1}^4{\bf A}_{{\rm W},j}
 =
 \sum_{\ell=1}^4\prod_{j=1}^4{\bf B}_{\pi,\ell,j}.
\quad}                                                        \tag{12.8}
\]
The product order in every term remains \(1,2,3,4\).  The same
statement with two factors applies to
\({\bf O}\odot{\bf Q}\).  Combining (12.4), (12.8), and (12.3)
therefore expresses \(F_*-F_{\rm W}\) as one coupled
second-order jet without first enclosing either action response.
\end{theorem}

\begin{proof}
Let \({\bf C}_\ell\) be the ordered product in which the positions
\(\pi(1),\ldots,\pi(\ell)\) use endpoint factors and all other
positions use Wilson factors.  Thus
\({\bf C}_0=\prod_j{\bf A}_{{\rm W},j}\) and
\({\bf C}_4=\prod_j{\bf A}_{*,j}\).  The difference
\({\bf C}_\ell-{\bf C}_{\ell-1}\) changes exactly the factor in
position \(\pi(\ell)\), and hence equals the \(\ell\)-th summand in
(12.8).  Summing these four differences telescopes.  No factors are
commuted.  Apply the same proof to the tadpole and then take
\([\tau^2]\Tr\).
\end{proof}

\subsection{The exact stock semiring and its diagnostic verdict}

For a Laurent factor \(A\), an external-jet order \(r\), an internal
coordinate \(\mu\), and \(0\le m\le2\), define the normalized stock
\[
 {\cal S}^{(\mu)}_{r,m}(A)
 :=\frac1{r!\,m!}
   \sum_{\alpha\in\mathbb Z^4}
   |\alpha_\mu|^m\|A^{[r]}_\alpha\|_{\rm F}.                 \tag{12.9}
\]
For covariance factors the same array denotes the operator-norm
majorant obtained recursively from the Hessian stocks and the proved
floors \(49/1250\) and \(343/1250\).  Products are bounded by the
finite convolution
\[
 (a\star b)_{r,m}
 =\sum_{\substack{r_1+r_2=r\\m_1+m_2=m}}
   a_{r_1,m_1}b_{r_2,m_2}.                                  \tag{12.10}
\]
All Frobenius norms in (12.9) are replaced by rational upper square
roots, so the compiled arrays are rational.

\begin{proposition}[Exact coupled-stock compiler]
\label{prop:v8-coupled-stock}
Apply (12.10) to the telescope (12.8), and at every bidegree bound
\(\Delta{\bf G}\) by the smaller of
\[
 {\cal S}({\bf G}_*)\star{\cal S}(\Delta{\bf H})
       \star{\cal S}({\bf G}_{\rm W})
 \quad\hbox{and}\quad
 {\cal S}({\bf G}_*)+{\cal S}({\bf G}_{\rm W}).              \tag{12.11}
\]
For the tadpole use
\(|\Tr(GQ)|\le\sqrt{45}\|G\|_{\rm op}\|Q\|_{\rm F}\);
for the bubble use the Frobenius norms of the two vertex factors and
operator norms of the covariance factors.  The resulting number at
\((r,m)=(2,0)\) is a pointwise majorant for
\(|F_*-F_{\rm W}|\), and twice the number at \((2,2)\) majorizes
\(|\partial_\mu^2(F_*-F_{\rm W})|\).
\end{proposition}

\begin{proof}
Equation (12.9) dominates the indicated derivative norm by the
Laurent triangle inequality.  The normalized Leibniz rule is exactly
the convolution (12.10).  Both entries of (12.11) are valid bounds:
the first uses Lemma~\ref{lem:v8-jet-resolvent}, and the second is the
triangle inequality.  Taking their minimum at a fixed derivative
degree is therefore valid.  Apply the Schatten inequalities stated in
the proposition to every mixed product in (12.8), sum its four
radii, and finally extract external degree two.  Since normalized
internal degree two is one half of the second derivative, the last
claim follows.
\end{proof}

The executable scan used \(s_{\rm W}=1\) and
\[
 s_*\in\left\{1,\frac32,2,3,\frac92,6,9\right\}.             \tag{12.12}
\]
For the pointwise stock and for all four coordinate second-derivative
stocks, the minimum in this family is \(s_*=9\).  Equivalently, the
minus-vertex multiplier is the common value \(2/3\) in both actions.
The optimal replacement order, with factors named as in (12.6), is
\[
 (1,3,4,2),                                                  \tag{12.13}
\]
or zero-based order \((0,2,3,1)\) in the verifier.

\begin{record}[Exact rejection of the global coupled stock]
\label{rec:v8-global-rejection}
The coefficient-exact results are
\[
\begin{split}
 |F_*-F_{\rm W}|
 &\le 9.440592848099\times10^{12},\\
 \sum_{\mu=1}^4\sup_k
 |\partial_\mu^2(F_*-F_{\rm W})(k)|
 &\le1.021267242779\times10^{21}.
\end{split}                                                   \tag{12.14}
\]
The first number is larger than the V7 independent-action envelope
\(9.415\times10^{12}\).  The second is larger than the V66
trace-ideal value \(1.018713982630\times10^{21}\) by the exact
factor whose decimal value is
\[
 1.002506356242.                                             \tag{12.15}
\]
Thus the coupled algebra does not survive these global coefficient
norms.

Using the periodic tensor-trapezoid estimate
\[
 |\langle f\rangle-\langle f\rangle_M|
 \le\frac{\pi^2}{3M^2}
       \sum_\mu\sup|\partial_\mu^2f|
 <\frac{10}{3M^2}
       \sum_\mu\sup|\partial_\mu^2f|,                         \tag{12.16}
\]
the smallest integer certified by (12.14) to beat (12.1) is
\[
 M=3916152587.                                               \tag{12.17}
\]
The corresponding V66 threshold is \(3911254156\).
These integers are certificate outputs, not proposed computations.
\end{record}

\begin{proof}
Proposition~\ref{prop:v8-coupled-stock} proves that the exact rational
outputs are upper bounds.  The verifier squares every rational
upper-square-root before accepting it, enumerates all \(24\)
replacement orders and all seven scale ratios in (12.12), and compares
the exact fractions.  Inequality (12.16) is the one-dimensional
periodic trapezoid remainder applied successively in the four
coordinates.  Substitution of the exact stock and (12.1), using
\(\pi^2<10\), gives (12.17) by integer square-root arithmetic.
\end{proof}

\begin{firewall}[What the rejection says]
\label{fire:v8-rejection-scope}
Record~\ref{rec:v8-global-rejection} is a no-go result only for the
specified global Laurent/Frobenius/operator stock calculus and the
finite scaling family (12.12).  It is not a lower bound on the actual
quadrature error and does not reject the coupled identity, local
preconditioning, affine or Taylor models, coefficientwise integration,
or a sharper Ward identity.  Its force is architectural: further
refinement of the same global triangle sums cannot discharge (12.1).
\end{firewall}

\subsection{Freeze the inverse, not the Laurent phases}

We now extract the part of the V3 cell argument which was lost when
every live factor was replaced by its centre.  Let \(B\) be a momentum
box and let \(R_{e,B}\) be a fixed Hermitian Gaussian-rational trial
inverse.  Write
\[
 E_{e,B}(k)=I-R_{e,B}H_{e,0}(k),\qquad
 \sup_{k\in B}\|E_{e,B}(k)\|_{\rm F}\le\eta_{e,B}<1,\qquad
 s_{e,B}=(1-\eta_{e,B})^{-1}.                                \tag{12.18}
\]
In every word of (6.12), replace each occurrence of
\(H_{e,0}^{-1}\) by \(R_{e,B}\), but leave
\(Q_2,P_j,M_j,H_1,H_2\) at their actual momentum \(k\).
Denote the resulting complete signed scalar by \(\Phi_{e,B}(k)\).
It is essential that \(\Phi_{e,B}\) is not evaluated at the box
centre.

\begin{theorem}[Live-phase frozen-inverse cell bound]
\label{thm:v8-live-phase-cell}
For every response factor \(X\), suppose
\[
 b_{e,B}(X)\ge
 \sup_{k\in B}\|R_{e,B}X_e(k)\|_{\rm F}.                     \tag{12.19}
\]
Use the seven word coefficients \(\gamma_q\), factor lists
\((X_{q,1},\ldots,X_{q,m_q})\), and inverse-slot counts \(m_q\)
from (7.8).  Then
\[
\boxed{\begin{aligned}
 |F_e(k)-\Phi_{e,B}(k)|
 \le\varepsilon_{e,B}:={}&
 \frac12(s_{e,B}-1)b_{e,B}(Q_2)\\
 &+|c_e|\sum_{q=1}^7|\gamma_q|
   (s_{e,B}^{m_q}-1)
   \prod_{j=1}^{m_q}b_{e,B}(X_{q,j})
\end{aligned}}                                               \tag{12.20}
\]
for every \(k\in B\).  Moreover, \(\Phi_{e,B}\) is a finite scalar
Laurent polynomial whose coefficients are Gaussian rational whenever
\(R_{e,B}\) is Gaussian rational.
\end{theorem}

\begin{proof}
From (12.18),
\[
 H_{e,0}^{-1}=(I-E_{e,B})^{-1}R_{e,B}
 =S_{e,B}R_{e,B},\qquad
 \|S_{e,B}\|_{\rm op}\le s_{e,B},\qquad
 \|S_{e,B}-I\|_{\rm F}\le s_{e,B}-1.                         \tag{12.21}
\]
For the tadpole, Hilbert--Schmidt duality and (12.19) give the
first term of (12.20).  For a bubble word, set
\(A_j=R_{e,B}X_{q,j}(k)\).  Replace the \(m_q\) copies of \(S_{e,B}\)
in
\(\Tr(SA_1\cdots SA_{m_q})\) one at a time.  Schatten--H\"older
gives
\[
 \left|\Tr(SA_1\cdots SA_{m_q})
 -\Tr(A_1\cdots A_{m_q})\right|
 \le(s_{e,B}^{m_q}-1)\prod_j b_{e,B}(X_{q,j}).               \tag{12.22}
\]
Multiply by the signed word coefficients and take the triangle
inequality only for the error radii.  This proves (12.20).

Every live factor is a finite matrix Laurent polynomial; multiplication
by a constant Gaussian-rational matrix, finite products, traces, and
finite sums preserve that property.  Hence \(\Phi_{e,B}\) is the
claimed scalar Laurent polynomial.
\end{proof}

\begin{remark}[The removed loss]
\label{rem:v8-removed-loss}
The V3 point-centred radius (7.10) also paid for
\(\|R[X(k)-X(a)]\|\) in every response factor.  Formula (12.20)
does not: all those phases remain in \(\Phi_{e,B}(k)\) and are
integrated with their signs.  Only the inverse defect is normed.
Thus Theorem~\ref{thm:v8-live-phase-cell} is not another refinement
of the failed pointwise radius; it changes the object being enclosed.
\end{remark}

\subsection{A finite coefficient-integral certificate}

If
\(\Phi_{*,B}-\Phi_{{\rm W},B}
 =\sum_\alpha d_{B,\alpha}e^{i\alpha\cdot k}\)
and \(B=\prod_{\mu=1}^4[a_\mu,b_\mu]\), then
\[
 \frac1{(2\pi)^4}\int_B
 (\Phi_{*,B}-\Phi_{{\rm W},B})\,dk
 =\frac1{(2\pi)^4}
  \sum_\alpha d_{B,\alpha}\prod_{\mu=1}^4
  J_{\alpha_\mu}(a_\mu,b_\mu),                              \tag{12.23}
\]
where
\[
 J_0(a,b)=b-a,\qquad
 J_n(a,b)=\frac{e^{inb}-e^{ina}}{in}\quad(n\ne0).            \tag{12.24}
\]
Equations (12.23)--(12.24) retain cancellation between all response
words and both actions before outward rounding.
We use boxes whose endpoints lie in \(\mathbb Q+\pi\mathbb Q\);
this class contains \(\pm\pi\) and \(\pm10^{-3}\) and is closed
under dyadic subdivision.

\begin{theorem}[Sound finite complement certificate]
\label{thm:v8-finite-certificate}
Let the charted complement
\({\cal K}\setminus{\cal C}_{10^{-3}}\) be partitioned, up to
measure-zero faces, into finitely many boxes \(B\), grouped in
\(k\mapsto-k\) pairs.  Suppose (12.18)--(12.20) are certified and a
rational real interval \(I_B\) contains the real part of the paired
value in (12.23).  Put
\[
 {\cal E}_{\rm inv}
 :=\sum_B\frac{|B|}{(2\pi)^4}
       (\varepsilon_{*,B}+\varepsilon_{{\rm W},B}),\qquad
 I_{\rm trial}:=\sum_B I_B.                                 \tag{12.25}
\]
Then
\[
 \left\langle(F_*-F_{\rm W})
 1_{{\cal K}\setminus{\cal C}_{10^{-3}}}\right\rangle
 \in I_{\rm trial}
       +[-{\cal E}_{\rm inv},{\cal E}_{\rm inv}].            \tag{12.26}
\]
Consequently, (12.1) follows if the interval obtained by subtracting
the exact V2 quarter-turn mean from the right side of (12.26) is
strictly contained in
\[
 \left(
 -\frac{443944758717}{2000000000},
  \frac{443944758717}{2000000000}\right).                   \tag{12.27}
\]
\end{theorem}

\begin{proof}
Integrate (12.20) on each box for both actions.  The normalized error
on \(B\) is at most its normalized volume times the sum of the two
pointwise inverse-defect radii.  Add the signed trial integrals as
intervals, then add all error radii.  This proves (12.26).
The definition (11.27) of \(E_{\rm comp}\) gives (12.27).
\end{proof}

\begin{theorem}[Completeness of the finitely described box architecture]
\label{thm:v8-certificate-completeness}
For every rational \(\varepsilon>0\), there exists a finite rational
box certificate of the form (12.18)--(12.26) for which
\[
 {\cal E}_{\rm inv}<\varepsilon                              \tag{12.28}
\]
and the total width of the coefficient-integral intervals
\(\sum_B I_B\) is less than \(\varepsilon\).
\end{theorem}

\begin{proof}
On the compact complement, both reduced Hessians are continuous and
have the positive global floors (1.2).  Their inverses and all finite
Laurent response factors are therefore uniformly continuous and
bounded.  At a point \(a\), Gaussian-rational Hermitian matrices are
dense, so choose \(R_{e,a}\) arbitrarily close to
\(H_{e,0}(a)^{-1}\).  By continuity, (12.18) then holds with
\(\eta_{e,B}\) arbitrarily small on a sufficiently small rational box
about \(a\), with endpoints in \(\mathbb Q+\pi\mathbb Q\).
Compactness supplies a finite box refinement.  The
quantities in (12.19) remain uniformly bounded, and the right side of
(12.20) tends uniformly to zero with the residual.  Refining until
its supremum is below \(\varepsilon\) proves (12.28), since the
normalized box volumes sum to at most one.

For fixed \(B\), Theorem~\ref{thm:v8-live-phase-cell} gives only
finitely many coefficients in (12.23).  Rational Taylor bounds for
sine and cosine enclose every endpoint value in (12.24) to arbitrary
width.  There are finitely many boxes and coefficients, so choose
their widths with total below \(\varepsilon\).  This constructs the
claimed finite rational certificate.
\end{proof}

\begin{remark}[Existence is not the sign computation]
\label{rem:v8-completeness-scope}
Theorem~\ref{thm:v8-certificate-completeness} proves that no limiting
or undecidable analytic step remains inside the finite-stage Haar
enclosure: a finite rational certificate exists at every requested
accuracy.  It gives no feasible box count and does not itself show
that the interval (12.27) passes.  Computational conditioning and
coefficient growth remain genuine tasks.
\end{remark}

\subsection{Executable provenance and next iteration}

\begin{record}[Coupled-difference verifier]
\label{rec:v8-verifier}
The executable
\begin{center}
\path{scripts/verify_ym_postarchive_v8_coupled_difference_jet.py}
\end{center}
constructs all stocks in exact fraction arithmetic, scans all
\(24\) bubble replacement orders and the seven scale ratios (12.12),
and proves the comparisons (12.14)--(12.17).  It separately evaluates
the coupled identity and the jet resolvent identity on all \(256\)
weighted quarter-turn momenta.  The largest relative residuals are
\[
 4.953\times10^{-15}
 \quad\hbox{and}\quad
 9.690\times10^{-14},                                       \tag{12.29}
\]
respectively.

The source has \(567\) lines, \(21426\) bytes, and SHA--256
\[
\texttt{342969DB2F32716012E92744DA012EA61B4F43BD4914D6F9F1E869ECB5E47A8D}.
                                                                  \tag{12.30}
\]
Its canonical payload has SHA--256
\[
\texttt{1F31D476724245435A6D5CCA6E760581A12BCEAF46E1B527887309265D006E84}.
                                                                  \tag{12.31}
\]
\end{record}

\begin{decision}[V9: integrate the live trial scalar]
\label{dec:v8-v9}
V9 should implement Theorem~\ref{thm:v8-finite-certificate}, beginning
with a shallow dyadic refinement of the quarter-turn cells intersected
with the deterministic V7 pivot charts.  On each box it should:
\begin{enumerate}[label=\arabic*.]
\item construct Gaussian-rational midpoint trial inverses and certify
      (12.18) by outward Laurent phase bounds;
\item form the complete coefficient dictionary of
      \(\Phi_{*,B}-\Phi_{{\rm W},B}\) before taking any absolute value;
\item enclose its integral by (12.23)--(12.24), add only the inverse
      defect (12.20), and split the box with largest normalized error;
\item use the high-alias V7 coordinates only as a preconditioner if
      the \(45\times45\) rational products dominate runtime.
\end{enumerate}
The stopping test is the exact interval inclusion (12.27), not a
pointwise sign and not a global derivative threshold.  V10 must begin
with the scheduled audit of the current SM and NS notes.
\end{decision}

\begin{assessment}[V8 progress]
\label{ass:v8-progress}
V8 implements endpoint-minus-Wilson cancellation as an exact
noncommutative jet telescope and proves, rather than guesses, that
global coefficient norms destroy its gain.  It then replaces the
failed point-centred cell object by a live-phase frozen-inverse
Laurent scalar, proves a sound finite complement certificate, and
proves that rational certificates of this form converge to arbitrary
accuracy.  The live finite-stage bottleneck is now constructive and
specific: produce a feasible adaptive ledger satisfying (12.27).
\end{assessment}

\begin{record}[V8 append record]
\label{rec:v8-append}
V8 is appended to the validated V7 whose installed SHA--256 was
\[
\texttt{46E8AA5EFAAB03E4141A6D7F5AB16D6C5FEE0D0E1F7530CB54CBAED4EEB9DE43}.
                                                                  \tag{12.32}
\]
After validation only V8 is to remain the active numeric YM note.
\end{record}

\begin{firewall}[Claim boundary after V8]
\label{fire:v8-final}
V8 does not yet supply the adaptive box ledger (12.27), so the
continuous Haar-mean sign and the nonzero Gaussian edge class remain
open.  It also does not close the compensating scale boundary, SCIC,
ADMC, the interacting cofinal state, regulator removal,
Osterwalder--Schrader reconstruction, physical exponential
clustering, or a Hamiltonian Yang--Mills mass gap.
\end{firewall}

% =============================================================================
% END APPENDED POST-TENTH-ARCHIVE V8 MATERIAL
% =============================================================================

% =============================================================================
% BEGIN APPENDED POST-TENTH-ARCHIVE V9 MATERIAL
% =============================================================================

\section{V9: exact phase balls and live Neumann tails}
\label{sec:v9-neumann-tail}

\subsection{Purpose of this iteration}

V8 proved that a finite rational live-phase certificate converges, but
did not determine a usable local scale.  This iteration implements the
first exact box.  It begins with the frozen-inverse error of
Theorem~\ref{thm:v8-live-phase-cell}; the result is again too narrow,
now for a precisely identified reason.  It then keeps finitely many
powers of the live inverse defect inside the Laurent scalar and bounds
only the Neumann tail.

At a generic quarter-turn anchor, first order needs half-width
\(2^{-22}\) to fit the complement budget.  The live Neumann polynomial
raises the half-width to \(2^{-8}\) at order \(47\), a factor \(2^{14}\)
in each coordinate.  This does not yet give a complement cover:
forming and integrating the resulting scalar coefficient dictionaries
remains open.  It does show that the catastrophic first-order scale is
not intrinsic to the response.

\subsection{Exact phase transport without centring the response}

Let
\[
 Y(k)=\sum_{\alpha\in{\cal A}}Y_\alpha e^{i\alpha\cdot k},
 \qquad
 \ell(Y):=\sum_\alpha|\alpha|_1\|Y_\alpha\|_{\rm F},\qquad
 \ell_R(Y):=\sum_\alpha|\alpha|_1\|RY_\alpha\|_{\rm F}.       \tag{13.1}
\]
For a periodic box
\[
 B(a,r):=\{k:\|k-a\|_\infty\le r\},                          \tag{13.2}
\]
all coordinate differences are taken in a fixed lift.  Put
\[
 a_R(Y):=\|RY(a)\|_{\rm F},\qquad
 \eta_0:=\|I-RH_0(a)\|_{\rm F}.                              \tag{13.3}
\]

\begin{lemma}[Exact phase-ball inputs for the V8 cell theorem]
\label{lem:v9-phase-ball}
On \(B(a,r)\),
\[
\begin{split}
 \|I-RH_0(k)\|_{\rm F}
 &\le\eta_0+r\ell_R(H_0),\\
 \|RX(k)\|_{\rm F}
 &\le a_R(X)+r\|R\|_{\rm op}\ell(X).
\end{split}                                                   \tag{13.4}
\]
If the Laurent coefficients and \(R\) are Gaussian rational and the
anchor phases are fourth roots of unity, every quantity on the
right side of (13.4) has a rational upper bound obtained by exact
integer arithmetic and rational upper square roots.
\end{lemma}

\begin{proof}
The first inequality is (7.16), with \(Y=H_0\), added to the anchor
residual.  For the second,
\[
\begin{split}
 \|RX(k)\|_{\rm F}
 &\le\|RX(a)\|_{\rm F}
      +\|R\|_{\rm op}\|X(k)-X(a)\|_{\rm F}\\
 &\le a_R(X)+r\|R\|_{\rm op}
      \sum_\alpha|\alpha|_1\|X_\alpha\|_{\rm F},
\end{split}
\]
again using
\(|e^{i\alpha\cdot(k-a)}-1|
 \le|\alpha|_1r\).
At a quarter-turn anchor all phases lie in
\(\{1,i,-1,-i\}\).  Matrix products, squared Frobenius norms, and
the row-sum operator majorant of a Hermitian trial inverse are then
rational operations; the declared square-root routine rounds upward
and verifies its square.
\end{proof}

\begin{remark}[No duplicated V3 transport]
\label{rem:v9-v3-use}
Lemma~\ref{lem:v9-phase-ball} deliberately reuses the V3 Laurent
transport theorem.  Its new role is to provide only the inverse
residual and the uniform live-factor sizes required by (12.20);
the signed response itself is not transported to the anchor.
\end{remark}

\subsection{Reflection pairing is part of the certificate}

The single-momentum response need not be real.  The correct real
object is the sum of a box and its reflected box.

\begin{lemma}[External-order reflection rule]
\label{lem:v9-reflection}
For every archived order-\(r\) Hessian, cubic, or quartic response
bank,
\[
 X_r(-k)=(-1)^r\overline{X_r(k)}.                            \tag{13.5}
\]
The inverse jets inherit the same rule, and hence the complete
order-two response satisfies
\[
 F_e(-k)=\overline{F_e(k)},\qquad
 (F_*-F_{\rm W})(-k)
 =\overline{(F_*-F_{\rm W})(k)}.                             \tag{13.6}
\]
Thus a reflected pair contributes twice the real part of either
member.
\end{lemma}

\begin{proof}
The Laurent coefficients before the stored jet factor are real.
Momentum reflection conjugates the phase.  An odd external derivative
or stored first-order jet contributes one factor \(i\), which changes
sign under conjugation; this proves (13.5).  The inverse recursion
(10.24) preserves total external order and therefore preserves the
same parity.  Every monomial contributing to
\([\tau^2]\) in (10.26) has total external order two.  Multiplying the
signs in (13.5) consequently gives \(+1\), proving (13.6).
\end{proof}

\subsection{Keeping finitely many live inverse-defect powers}

Fix an action and one box, suppress their subscripts, and use
\[
 E(k):=I-RH_0(k),\qquad
 \sup_{k\in B}\|E(k)\|_{\rm F}\le\eta<1,\qquad
 s:=(1-\eta)^{-1}.                                          \tag{13.7}
\]
Define
\[
 S_N(k):=\sum_{j=0}^{N-1}E(k)^j,\qquad
 G_N(k):=S_N(k)R.                                            \tag{13.8}
\]
In the seven-word formula (6.12), replace every covariance \(G\)
by \(G_N\), leaving all \(H_1,H_2,P_j,M_j,Q_2\) live.  Call the
complete signed result \(\Phi_B^{(N)}(k)\).  Since \(E\) and every
response factor are finite Laurent matrices,
\(\Phi_B^{(N)}\) is a finite scalar Laurent polynomial.

\begin{theorem}[Live Neumann-tail response bound]
\label{thm:v9-neumann-tail}
Let \(b_B(X)\) be any valid bounds from (13.4), and retain the word
data \(\gamma_q,(X_{q,j})_{j=1}^{m_q}\) from (7.8).  Then
\[
\boxed{\begin{aligned}
 |F_e(k)-\Phi_{e,B}^{(N)}(k)|
 \le\varepsilon_{e,B}^{(N)}:={}&
 \frac12s_{e,B}\eta_{e,B}^{N}b_{e,B}(Q_2)\\
 &+|c_e|\sum_{q=1}^7|\gamma_q|m_q
      s_{e,B}^{m_q}\eta_{e,B}^{N}
      \prod_{j=1}^{m_q}b_{e,B}(X_{q,j})
\end{aligned}}                                               \tag{13.9}
\]
for every \(k\in B\).  The right side tends geometrically to zero
with \(N\).
\end{theorem}

\begin{proof}
Let \(S=(I-E)^{-1}\).  Since \(S\) is a convergent power series in
\(E\), it commutes with \(E\), and
\[
 S-S_N=E^NS,\qquad
 \|S\|_{\rm op},\|S_N\|_{\rm op}\le s,\qquad
 \|S-S_N\|_{\rm op}\le s\eta^N.                             \tag{13.10}
\]
Also
\(\|S-S_N\|_{\rm F}\le s\eta^N\), because
\(\|E^N\|_{\rm F}\le\|E\|_{\rm F}^N\).

The tadpole difference is bounded by Hilbert--Schmidt duality,
giving the first term of (13.9).  In a word with \(m_q\) inverse
slots, write \(A_j=RX_{q,j}(k)\).  The exact and truncated words are
\(\Tr(SA_1\cdots SA_{m_q})\) and
\(\Tr(S_NA_1\cdots S_NA_{m_q})\).  Replace the \(m_q\) occurrences
one at a time.  Each mixed term is bounded by
\(s^{m_q}\eta^N\prod_j b_B(X_{q,j})\), so their sum gives the second
line of (13.9).  Because \(\eta<1\), geometric convergence follows.
\end{proof}

\begin{remark}[Relation to V8]
\label{rem:v9-order-one}
At \(N=1\), \(G_N=R\) and the trial scalar is exactly the V8
frozen-inverse scalar.  For that case (12.20) is sharper than (13.9),
because it retains \(s^{m_q}-1\) rather than bounding it by
\(m_qs^{m_q}\eta\).  The executable therefore uses (12.20) for the
first-order radius and (13.9) for \(N\ge2\).
\end{remark}

\begin{corollary}[Finite coefficient integral at every order]
\label{cor:v9-neumann-integral}
For every fixed \(N\), the reflection-paired box integral of
\(\Phi_{*,B}^{(N)}-\Phi_{{\rm W},B}^{(N)}\) is a finite sum of the
coefficient integrals (12.23)--(12.24).  Adding the normalized radius
\[
 \frac{|B|}{(2\pi)^4}
 \left(\varepsilon_{*,B}^{(N)}
       +\varepsilon_{{\rm W},B}^{(N)}\right)                 \tag{13.11}
\]
gives a sound real interval for the exact paired contribution.
\end{corollary}

\begin{proof}
Finite Laurent matrices are closed under the finite sum and products
in (13.8), under the seven response words, and under traces.
Apply Theorem~\ref{thm:v9-neumann-tail}, integrate its uniform error,
and use Lemma~\ref{lem:v9-reflection}.
\end{proof}

\subsection{An exact generic pilot box}

Take the quarter-turn index and momentum
\[
 p=(1,2,3,0),\qquad
 a=\frac{\pi}{2}p
   =(\pi/2,\pi,-\pi/2,0)\quad\hbox{on the periodic torus}.    \tag{13.12}
\]
The second representative in (13.12) uses \(-\pi/2\) instead of
\(3\pi/2\).  The V2 Gaussian-rational trial inverses give
\[
\begin{array}{c|ccc}
 &\eta_0&\|R\|_{\rm op}&\ell_R(H_0)\\ \hline
 {\rm W}
 &4.23399\times10^{-7}&46.156334083&109.116686040\\
 *
 &1.3300247\times10^{-5}&2.105629958&185.306737872 .
\end{array}                                                   \tag{13.13}
\]
Every displayed decimal in this subsection is the image of an exact
rational upper bound recorded by the verifier.

\begin{record}[First-order pilot obstruction]
\label{rec:v9-first-order}
For both actions, the largest dyadic half-width accepted by the
Neumann test in (13.4) is
\[
 r=2^{-8}=\frac1{256};                                      \tag{13.14}
\]
the endpoint residual there is approximately \(0.723867746\), while
the Wilson residual is approximately \(0.426237479\).

In contrast, the sharper first-order error (12.20) first falls below
the full complement budget at
\[
\begin{split}
 r&=2^{-22}=\frac1{4194304},\\
 \varepsilon_*+\varepsilon_{\rm W}
 &\doteq177.447699215
 <\frac{443944758717}{2000000000}.
\end{split}                                                   \tag{13.15}
\]
At the preceding radius \(2^{-21}\), the exact bound is approximately
\(326.330430372\) and fails.  Thus live phases alone do not cure the
product-of-factor-norm loss.
\end{record}

The exact V2 real-part interval at the anchor is
\[
 -811.672416293
 <\Re(F_*-F_{\rm W})(a)
 <-754.464063403<0.                                         \tag{13.16}
\]
The floating oracle value is
\[
 -783.068344338+73.748851896\,i.                            \tag{13.17}
\]
Its reflected value is its complex conjugate to zero reported
floating residual, illustrating why Lemma~\ref{lem:v9-reflection}
is part of the certificate rather than a cosmetic symmetry.

\begin{theorem}[Exact higher-order recovery on the pilot box]
\label{thm:v9-pilot-recovery}
The exact rational evaluation of (13.9) gives the following least
orders whose combined action tail is strictly below the complement
budget:
\[
\begin{array}{c|c|c|c}
r&N&\varepsilon_*^{(N)}+\varepsilon_{\rm W}^{(N)}
 &\hbox{tail at }N-1\\ \hline
2^{-8} &47&183.378041&253.330863\\
2^{-9} &12& 91.854779&253.897353\\
2^{-10}& 6&201.826099&1136.203821\\
2^{-11}& 5& 21.534138&247.348789 .
\end{array}                                                   \tag{13.18}
\]
All four passing inequalities and all four preceding-order failures
are strict comparisons of exact fractions.
\end{theorem}

\begin{proof}
Use (13.4) with the exact data (13.13), substitute the resulting
\(\eta\)'s and live-factor bounds into (13.9), and compare fractions
with (12.1).  The executable performs this calculation without
floating decisions.  It scans \(N=1,2,\ldots\), records the first
passing order, and separately verifies that \(N-1\) fails.
\end{proof}

\begin{assessment}[What the pilot changes]
\label{ass:v9-pilot}
The first-order radius \(2^{-22}\) would make a four-dimensional
cover hopeless.  Theorem~\ref{thm:v9-pilot-recovery} proves that this
scale is an artefact of discarding every live defect power: a finite
Neumann polynomial reaches \(2^{-8}\).  However, a uniform
\(2^{-8}\) cover is still far too large, and order \(47\) coefficient
growth may itself be prohibitive.  The correct next move is therefore
to combine moderate Neumann order with the V7 high-alias coordinates,
whose Wilson and endpoint inverse bounds are \(3/2\) and \(3/14\)
instead of the pilot values in the middle column of (13.13).
\end{assessment}

\subsection{Executable provenance and the V10 audit gate}

\begin{record}[Live-phase pilot verifier]
\label{rec:v9-verifier}
The executable
\begin{center}
\path{scripts/verify_ym_postarchive_v9_live_phase_pilot_box.py}
\end{center}
uses exact Gaussian-rational anchor matrices, exact residual-weighted
Hessian stocks, rational upper square roots, and exact fraction
comparisons.  It checks (13.15), all rows of (13.18), the V2
zero-radius identity, the negative anchor interval, and reflection
conjugacy.  Its source has \(387\) lines, \(15142\) bytes, and
SHA--256
\[
\texttt{40A1F65BF952A378B6F20C6656BFC27DB2D0004E6AFD64697A1D341E3735247B}.
                                                                  \tag{13.19}
\]
The canonical payload has SHA--256
\[
\texttt{BB32504B88A9F06374F62663A35F429014A341D698A81712F91680A9112A586F}.
                                                                  \tag{13.20}
\]
\end{record}

\begin{decision}[V10: audit, then precondition the Neumann ledger]
\label{dec:v9-v10}
As required by the five-version rule, V10 must first inspect the
current SM and NS active notes and record any transferable device.
It should then repeat the exact pilot in the rational V7 high-alias
charts:
\begin{enumerate}[label=\arabic*.]
\item use \(N_e^{(r)}\), not the poorly scaled reduced Hessian, in
      the local defect \(I-RN_e^{(r)}\);
\item compile residual-weighted chart Laurent stocks with the
      denominator floors (11.8), (11.14), and (11.15);
\item compare the pairs \((r,N)\) against (13.18) before forming a
      large cover;
\item form a coefficient dictionary for the complete reflection-paired
      scalar only after a chart shows a genuinely feasible scale.
\end{enumerate}
If the chart congruences erase the inverse gain, the programme should
derive the live response factors directly in ambient alias coordinates
instead of returning to the \(45\times45\) trial inverse.
\end{decision}

\begin{assessment}[V9 progress]
\label{ass:v9-progress}
V9 supplies the first exact live-phase box, corrects the pointwise
reality assumption by proving the reflection rule, and replaces the
failed first-order inverse by a finite live Neumann polynomial with a
rigorous geometric tail.  The certified \(2^{14}\) radius improvement
shows that higher inverse order is mathematically effective.  The
remaining finite-stage bottleneck is now the joint optimization of
chart conditioning, Neumann order, coefficient growth, and adaptive
box count.
\end{assessment}

\begin{record}[V9 append record]
\label{rec:v9-append}
V9 is appended to the validated V8 whose installed SHA--256 was
\[
\texttt{23EC89B8E57800E424713FE450882682DF6009D431C38B25A6D6E0F281A9B0C0}.
                                                                  \tag{13.21}
\]
After validation only V9 is to remain the active numeric YM note.
\end{record}

\begin{firewall}[Claim boundary after V9]
\label{fire:v9-final}
V9 certifies one pilot box and a convergent higher-order inverse tail;
it does not form the full complement coefficient-integral ledger or
prove the continuous Haar-mean sign.  The nonzero Gaussian edge class,
compensating scale boundary, SCIC, ADMC, interacting cofinal state,
regulator removal, Osterwalder--Schrader reconstruction, physical
clustering, and Hamiltonian Yang--Mills mass gap remain open.
\end{firewall}

% =============================================================================
% END APPENDED POST-TENTH-ARCHIVE V9 MATERIAL
% =============================================================================

% =============================================================================
% BEGIN APPENDED POST-TENTH-ARCHIVE V10 MATERIAL
% =============================================================================

\section{V10: a frozen rational high-alias preconditioner and an exact scale gain}
\label{sec:v10-frozen-high-alias}

V9 proved that retaining finitely many live inverse-defect powers can
enlarge a single certified box from the first-order scale \(2^{-22}\) to
\(2^{-8}\).  Its generic anchor still used the poorly scaled reduced
coordinates.  This continuation performs the required companion audit and
then goes upstream to the V7 high-alias chart.  The decisive adjustment is
to freeze the chart at the centre of each box.  This retains the improved
centre conditioning while making every live transformed factor a finite
Laurent polynomial; no moving-pivot denominator is differentiated inside
the box.

\subsection{Required V10 companion audit}

The current active companion notes inspected read-only were
\begin{center}
\path{latex_notes/Renewal_SM_Einstein_Continuum_Research_Notes_V56.tex},\\
\path{latex_notes/Renewal_NS_Stability_Research_Notes_V26.tex}.
\end{center}
Their byte counts and SHA--256 hashes were respectively
\[
\begin{split}
 &(592553,
 \texttt{C4C284AB2C89D9241A3239040D76FEBC6EAC3755EC4E10E819BAC880A08DECFF}),\\
 &(278283,
 \texttt{82C887F8C2C69E4F28FAD7C3DC8DE5B9099CB5A4BF431EBA1FFF3E91935978AD}).
\end{split}                                                     \tag{14.1}
\]

SM V56 proves a compile-then-estimate principle for cancelling chiral
phases: signed sectors are combined before any absolute value is taken.
The valid YM transfer is that the finite coefficient dictionary for
\(\Phi_*-\Phi_{\rm W}\) must be formed before coefficient norms.  It does
not allow the two independent inverse-defect error radii to cancel.

NS V26 computes norms only on the rank-one inputs actually used by its
equation.  Its first-derivative gain is zero and its second-derivative gain
is modest, so no numerical NS estimate transfers.  The sound methodological
transfer is to norm the image actually entering the defect.  Below this
means compiling \(RZ^*C_\nu Z\) coefficient by coefficient for the Hessian,
instead of multiplying three unrelated ambient norms.  No Standard-Model
or Navier--Stokes conclusion is used as a Yang--Mills premise.

\begin{assessment}[V10 audit decision]
\label{ass:v10-audit}
The SM rule governs the signed trial polynomial; the NS rule governs the
directional defect stock.  The former is postponed until a feasible pair
\((r,N)\) has been selected, while the latter is implemented exactly in the
present certificate.
\end{assessment}

\subsection{Freeze the centre chart}

Fix an anchor \(a\) and an invertible matrix \(Z\), independent of the live
momentum \(k\) in its box.  For either action put
\[
 N_e(k)=Z^*H_e(k)Z,
 \qquad \widetilde X(k)=Z^*X(k)Z.                            \tag{14.2}
\]

\begin{theorem}[Frozen-coordinate response compiler]
\label{thm:v10-frozen-compiler}
Suppose a matrix \(R_e\) satisfies
\[
 \|I-R_eN_e(a)\|_{\rm F}<1.
\]
Then \(N_e(a)\), \(Z\), and \(H_e(a)\)
are invertible.  Wherever \(H_e(k)\) is invertible,
\[
 H_e(k)^{-1}=Z N_e(k)^{-1}Z^*,                              \tag{14.3}
\]
and every cyclic response word obeys
\[
 \Tr\bigl(H_e^{-1}X_1\cdots H_e^{-1}X_m\bigr)
 =\Tr\bigl(N_e^{-1}\widetilde X_1\cdots
                 N_e^{-1}\widetilde X_m\bigr).             \tag{14.4}
\]
If \(H_e\) and the \(X_j\) are finite Laurent matrices, so are \(N_e\)
and all \(\widetilde X_j\).
\end{theorem}

\begin{proof}
The strict residual inequality makes \(R_eN_e(a)\) invertible by the
Neumann lemma, hence \(N_e(a)\) is invertible.  If \(Zx=0\), then
\(N_e(a)x=Z^*H_e(a)Zx=0\), so \(x=0\); square \(Z\) is therefore
invertible.  Equation (14.3) follows by multiplying
\(N_e=Z^*H_eZ\), and cyclicity of trace gives (14.4).  Multiplication by
constant matrices preserves finite Laurent support.
\end{proof}

This theorem is deliberately stronger than a moving rational atlas for the
local task: once \(Z\) is frozen, the V7 pivot denominators become fixed
rational data and create no derivative loss.

\subsection{An exact residual-weighted frozen-box bound}

Write
\[
 H_e(k)=\sum_\nu C_{e,\nu}e^{i\nu\cdot k},
 \qquad X(k)=\sum_\nu D_{X,\nu}e^{i\nu\cdot k},
 \qquad A_e=R_eZ^*.                                         \tag{14.5}
\]
Define
\[
\begin{split}
 \eta_{e,0}&=\|I-R_eN_e(a)\|_{\rm F},\\
 L_{e,H}^{\sharp}
   &=\sum_{\nu\ne0}|\nu|_1\,
       \|A_eC_{e,\nu}Z\|_{\rm F},\\
 L_X^{\rm raw}
   &=\sum_{\nu\ne0}|\nu|_1\|D_{X,\nu}\|_{\rm F},\\
 \chi_e&=\|A_e\|_{\rm op}\|Z\|_{\rm op},\\
 b_{e,X}(r)&=\|A_eX(a)Z\|_{\rm F}+r\chi_eL_X^{\rm raw}.
\end{split}                                                     \tag{14.6}
\]

\begin{theorem}[Frozen-box transport]
\label{thm:v10-frozen-transport}
On the periodic sup-norm box
\(B(a,r)=\{k:|k_\mu-a_\mu|\le r\}\),
\[
 \|I-R_eN_e(k)\|_{\rm F}
 \le \eta_e(r):=\eta_{e,0}+rL_{e,H}^{\sharp},
 \qquad
 \|R_e\widetilde X(k)\|_{\rm F}\le b_{e,X}(r).             \tag{14.7}
\]
Consequently, whenever \(\eta_e(r)<1\), the V9 finite-Neumann trial and
tail theorem apply with \(\eta_{e,B}=\eta_e(r)\) and
\(b_{e,B}(X)=b_{e,X}(r)\).
\end{theorem}

\begin{proof}
For \(k=a+\delta\),
\[
 R_e\bigl(N_e(k)-N_e(a)\bigr)
 =\sum_\nu A_eC_{e,\nu}Z e^{i\nu\cdot a}
       \bigl(e^{i\nu\cdot\delta}-1\bigr).
\]
The elementary inequalities
\(|e^{it}-1|\le|t|\) and
\(|\nu\cdot\delta|\le r|\nu|_1\), followed only now by the triangle
inequality, give the first assertion.  The same expansion for \(X\), with
\(\|A_eD_{X,\nu}Z\|_{\rm F}
 \le\|A_e\|_{\rm op}\|D_{X,\nu}\|_{\rm F}\|Z\|_{\rm op}\), gives the
second.  The last assertion is Theorem~\ref{thm:v9-neumann-tail}.
\end{proof}

\begin{remark}[Where the NS transfer is and is not used]
\label{rem:v10-directional-stock}
The Hessian stock in (14.6) retains the exact image
\(A_eC_{e,\nu}Z\), the analogue of restricting the NS tensor to its used
input.  The other factors presently use the coarser product
\(\chi_eL_X^{\rm raw}\).  Thus the remaining opportunity for V11 is visible and
quantified rather than silently assumed.
\end{remark}

\subsection{The exact Gaussian pilot}

Take
\[
 p=(2,2,0,0),\qquad a=(\pi,\pi,0,0).                       \tag{14.8}
\]
At this point the V7 pivot chart has a small-denominator Gaussian-rational
shadow.  The verifier defines \(Z_\sharp\in M_{45}(\mathbb Q(i))\) by
entrywise rational reconstruction with maximum denominator \(16\); every
reconstructed denominator is in fact at most \(4\).  The largest entrywise
difference from the discovery chart is \(2.318\times10^{-14}\).  This
floating comparison is not used as proof: the exact residual tests in the
next table independently prove that \(Z_\sharp\) is an invertible valid
preconditioner.

The exact identity-chart and frozen-\(Z_\sharp\) stocks are
\[
\begin{array}{c|cc|cc}
 &\eta_{{\rm W},0}&\eta_{*,0}
 &L_{{\rm W},H}^{\sharp}&L_{*,H}^{\sharp}\\ \hline
 I
 &2.56472\!\times10^{-7}&1.0255126\!\times10^{-5}
 &113.735336368&183.075547730\\
 Z_\sharp
 &1.534667\!\times10^{-6}&2.8237744\!\times10^{-5}
 &77.103612636&133.329872724
\end{array}                                                   \tag{14.9}
\]
The corresponding exact operator-norm upper bounds are
\[
\begin{array}{c|cc}
 &\|R_{\rm W}\|_{\rm op}&\|R_*\|_{\rm op}\\ \hline
 I&51.64285718&2.32937089\\
 Z_\sharp&1.125&0.08104396.
\end{array}
\]
All entries are outward images of exact fractions.  The identity row uses
the physical reduced-Hessian floors; the second row may also use the V7
floors \(2/3\) and \(14/3\), though invertibility already follows from the
displayed strict residuals.  The high-alias residuals themselves are a
little larger, but both relevant Hessian transport stocks and both inverse
bounds are sharply smaller.

At the anchor, the exact high-alias trial disk gives
\[
\begin{split}
 \Phi_*^{\rm tr}-\Phi_{\rm W}^{\rm tr}
   &=-178.245950627+2.911911\times10^{-6}i,\\
 \varepsilon_*+\varepsilon_{\rm W}&=1.444546908,\\
 -179.690497535
   &<\Re(F_*-F_{\rm W})(a)<-176.801403719<0.
\end{split}                                                     \tag{14.10}
\]
The independent floating oracle is
\(-178.246221541-1.49\times10^{-14}i\) and lies in this exact interval.
The small nonzero imaginary part of the independently rounded trial is
strictly below its residual disk, as required by self-reflection of the
exact response.  For comparison, the identity trial has radius
\(41.615988147\); the frozen chart therefore also makes the point enclosure
substantially sharper.

\begin{theorem}[Exact first-order scale gain]
\label{thm:v10-first-order-gain}
Using (14.7) and the sharp V8 first-order multiplier, the first dyadic
half-width at which the two-action inverse-defect error is strictly below
the complement budget is
\[
\begin{array}{c|c|c|c}
\hbox{coordinates}&r&\varepsilon_*+\varepsilon_{\rm W}
 &\hbox{error at the preceding dyadic radius}\\ \hline
I&2^{-23}&164.860841&288.122901\\
Z_\sharp&2^{-16}&145.267787&300.533499.
\end{array}                                                   \tag{14.11}
\]
Every comparison is exact.  Thus the frozen high-alias preconditioner gains
seven dyadic powers in half-width, a local four-dimensional volume factor
\(2^{28}\), before any higher Neumann powers are used.
\end{theorem}

\begin{proof}
Insert the exact data defining (14.9) into (14.7) and then into (12.20).
The executable compares the resulting fractions with
\(443944758717/2000000000\).  It verifies strict passage at the displayed
radii and strict failure at their immediate predecessors.
\end{proof}

\begin{theorem}[Exact higher-order comparison]
\label{thm:v10-order-gain}
For four useful radii, the least Neumann orders and their passing tails are
\[
\begin{array}{c|cc|cc}
 &\multicolumn{2}{c|}{I}&\multicolumn{2}{c}{Z_\sharp}\\
r&N&\hbox{tail}&N&\hbox{tail}\\ \hline
2^{-8} &47&178.071942&21&163.493705\\
2^{-9} &12&149.594618& 8& 94.725626\\
2^{-10}& 7& 63.250990& 4&187.090508\\
2^{-11}& 5& 38.449411& 3& 58.777677.
\end{array}                                                   \tag{14.12}
\]
For the \(Z_\sharp\) rows, the preceding-order tails are
\[
\begin{array}{c|c@{\qquad}c|c}
2^{-8}&313.899076&2^{-9}&363.947929\\
2^{-10}&1455.935704&2^{-11}&938.365350,
\end{array}
\]
and hence all fail.
The executable also verifies the four corresponding identity-chart
preceding-order failures.  All sixteen pass/fail comparisons are exact.
\end{theorem}

\begin{proof}
Apply Theorem~\ref{thm:v9-neumann-tail} with the bounds from
Theorem~\ref{thm:v10-frozen-transport}.  For each radius scan positive
integer \(N\), compare exact fractions with the budget, and separately
check \(N-1\).  This is precisely the verifier's minimality gate.
\end{proof}

\begin{assessment}[What V10 changes]
\label{ass:v10-change}
The V7 conditioning gain survives the transformed response factors and a
conservative rigorous transport bound; it is not cancelled by chart
congruence.  Among the tested choices, \(r=2^{-9},N=8\) is the first
plausible coefficient-compilation target: halving the \(2^{-8}\) half-width
costs a factor \(16\) in a uniform four-dimensional count but reduces the
inverse polynomial from order \(21\) to order \(8\).  Even this uniform
count is far too large, so V11 must combine exact factor-image stocks,
signed coefficient compilation, and adaptive boxes rather than instantiate
a uniform atlas.
\end{assessment}

\subsection{Executable provenance and the V11 route}

\begin{record}[Frozen high-alias verifier]
\label{rec:v10-verifier}
The executable
\begin{center}
\path{scripts/verify_ym_postarchive_v10_frozen_high_alias.py}
\end{center}
has \(582\) physical lines, \(20931\) bytes, and SHA--256
\[
\texttt{EDBF857C2A7638B3F44DCE8BFD922BFB94A0B0DAD608CC45DC2ACF8A5C594E47}.
                                                                  \tag{14.13}
\]
Its canonical payload has SHA--256
\[
\texttt{119E4E4E3B99E6B1B27FAF923E5C0BC9E790BC6B14F0682457F0493BDC44A45B}.
                                                                  \tag{14.14}
\]
The candidate is discovered numerically, but every invertibility,
residual, stock, interval, budget, and minimal-order decision is made with
exact fractions and outward rational norm bounds.
\end{record}

\begin{decision}[V11: compile the used factor images]
\label{dec:v10-v11}
Keep the frozen \(Z_\sharp\) construction and the target
\((r,N)=(2^{-9},8)\).  Replace each coarse factor transport bound
\(\chi_eL_X^{\rm raw}\) by the exact used-image stock
\[
 L_{e,X}^{\sharp}
 =\sum_{\nu\ne0}|\nu|_1\|A_eD_{X,\nu}Z_\sharp\|_{\rm F}.   \tag{14.15}
\]
If those stocks preserve a moderate order, form the complete finite
Laurent dictionary of
\(\Phi_{*,B}^{(8)}-\Phi_{{\rm W},B}^{(8)}\), pair its reflected box, and
take coefficient absolute values only after the signed difference has been
compiled, as required by the SM audit.  If dictionary growth is excessive,
retain the same mathematics but use sparse convolution and exact frequency
pruning; do not retreat to seven separately bounded response words.
\end{decision}

\begin{record}[V10 append record]
\label{rec:v10-append}
V10 is appended to the validated V9 whose installed SHA--256 was
\[
\texttt{C97496D124E981A1BFC9A23C13B346F82A368D7E46FDB1B5560D0711104F5D86}.
                                                                  \tag{14.16}
\]
After validation only V10 is to remain the active numeric YM note.
\end{record}

\begin{firewall}[Claim boundary after V10]
\label{fire:v10-final}
V10 proves a one-anchor frozen-preconditioner gain and selects a concrete
next coefficient target.  It does not cover the complement, compile its
signed coefficient integral, or prove the continuous Haar-mean sign.  The
nonzero Gaussian edge class, compensating scale boundary, SCIC, ADMC,
interacting cofinal state, regulator removal, Osterwalder--Schrader
reconstruction, physical clustering, and Hamiltonian Yang--Mills mass gap
remain open.
\end{firewall}

% =============================================================================
% END APPENDED POST-TENTH-ARCHIVE V10 MATERIAL
% =============================================================================

% =============================================================================
% BEGIN APPENDED POST-TENTH-ARCHIVE V11 MATERIAL
% =============================================================================

\section{V11: exact used-image stocks and the sparse signed-compiler gate}
\label{sec:v11-used-image}

V10 showed that a frozen rational shadow of the V7 high-alias chart gives
a genuine conditioning gain, but it still transported every non-Hessian
factor through the product
\(\|RZ^*\|_{\rm op}\|Z\|_{\rm op}L_X^{\rm raw}\).
This continuation carries out the NS-inspired refinement announced there:
each Laurent coefficient is first mapped through the actual used image
\(RZ^*D_{X,\nu}Z\), and only that image is normed.  It then tests the
provisional pair \((r,N)=(2^{-9},8)\), finds the true minimum \(N=5\), and
computes the frequency envelope which a signed coefficient compiler must
handle.

\subsection{An exact cubic-bank symmetry}

For action \(e\), external derivative order \(j\), and Laurent frequency
\(\nu\), denote the stored plus and minus cubic coefficient matrices by
\(P_{e,j,\nu}\) and \(M_{e,j,\nu}\).

\begin{lemma}[Coefficientwise plus/minus identity]
\label{lem:v11-plus-minus}
For both the Wilson and endpoint banks and every frequency,
\[
 M_{e,0,\nu}=P_{e,0,\nu},\qquad
 M_{e,1,\nu}=-P_{e,1,\nu},\qquad
 M_{e,2,\nu}=P_{e,2,\nu}.                                  \tag{15.1}
\]
Consequently the corresponding centre Frobenius bounds, used-image
Lipschitz stocks, and frequency supports agree exactly.
\end{lemma}

\begin{proof}
The verifier takes the union of the two dictionary supports in every one
of the \(45^2\) matrix entries and compares the integer coefficient at
each key.  It finds the signs \(+1,-1,+1\), with no absent key on either
side.  Left and right multiplication by fixed matrices preserves these
scalar signs, and the Frobenius norm removes them.  Support equality is
coefficientwise.
\end{proof}

This identity removes three redundant exact matrix compilations; it is not
an equality inferred from matching decimal norms.

\subsection{The sharp used-image transport}

Retain the V10 notation
\[
 A_e=R_eZ_\sharp^*,\qquad
 X(k)=\sum_\nu D_{X,\nu}e^{i\nu\cdot k}.
\]
Define
\[
 \boxed{\quad
 L_{e,X}^{\sharp}
 :=\sum_{\nu\ne0}|\nu|_1
       \|A_eD_{X,\nu}Z_\sharp\|_{\rm F}.
 \quad}                                                       \tag{15.2}
\]

\begin{theorem}[Used-image factor bound]
\label{thm:v11-used-image}
For \(k\in B(a,r)\),
\[
 \|R_e\widetilde X(k)\|_{\rm F}
 \le
 \|A_eX(a)Z_\sharp\|_{\rm F}+rL_{e,X}^{\sharp}.              \tag{15.3}
\]
Thus every occurrence of
\(\chi_eL_X^{\rm raw}\) in the V10 box theorem may be replaced by
\(L_{e,X}^{\sharp}\).  The V9 Neumann-tail theorem remains valid with
these smaller factor bounds.
\end{theorem}

\begin{proof}
Expand the difference from the anchor:
\[
 A_e\bigl(X(k)-X(a)\bigr)Z_\sharp
 =\sum_\nu A_eD_{X,\nu}Z_\sharp e^{i\nu\cdot a}
      \bigl(e^{i\nu\cdot(k-a)}-1\bigr).
\]
Apply \(|e^{it}-1|\le |t|\) and
\(|\nu\cdot(k-a)|\le r|\nu|_1\), then take the triangle inequality after
the used image has been formed.  This proves (15.3).  Since only the
factor bounds in Theorem~\ref{thm:v9-neumann-tail} have changed, its proof
applies verbatim.
\end{proof}

The exact stocks at \(a=(\pi,\pi,0,0)\) are
\[
\begin{array}{c|ccc}
e&L_{e,H_1}^{\sharp}=L_{e,H_2}^{\sharp}
 &L_{e,Q_2}^{\sharp}&L_{e,P_0}^{\sharp}\\ \hline
{\rm W}&30.015779081&183.414107947&71.302507744\\
*&56.360087204&1192.246820519&45.550509751
\end{array}                                                     \tag{15.4}
\]
and
\[
\begin{array}{c|cc}
e&L_{e,P_1}^{\sharp}&L_{e,P_2}^{\sharp}\\ \hline
{\rm W}&45.549965358&79.536090091\\
*&61.365069235&72.973602579.
\end{array}
\]
The \(M_j\) entries equal the \(P_j\) entries by
Lemma~\ref{lem:v11-plus-minus}.  Every displayed decimal is an outward
image of an exact fraction.  Relative to the V10 factored stocks, the
improvement factors range from \(44.4116\) to \(107.364\); the verifier
checks every factor separately and proves all inequalities strict.

\subsection{The revised exact box ledger}

The sharper factor transport does not change the Hessian defect
\[
\eta_e(r)=\eta_{e,0}+rL_{e,H}^{\sharp},
\]
but reduces every product multiplying its Neumann powers.

\begin{theorem}[First-order scale is stable but sharper]
\label{thm:v11-first-order}
The least dyadic first-order half-width below the complement budget remains
\(2^{-16}\), but its exact error improves from the V10 value \(145.267787\)
to
\[
 \varepsilon_*+\varepsilon_{\rm W}=140.327544.              \tag{15.5}
\]
At the preceding half-width \(2^{-15}\), the refined error is
\(280.518820\) and still fails.
\end{theorem}

\begin{proof}
Insert (15.3) into the V8 first-order multiplier (12.20) and compare exact
fractions with \(443944758717/2000000000\).  The stated radius passes
strictly and its predecessor fails strictly.
\end{proof}

\begin{theorem}[Exact Neumann-order reduction]
\label{thm:v11-order-reduction}
The least orders after used-image refinement are
\[
\begin{array}{c|c|c|c|c}
r&N_{\rm V10}&N_{\rm V11}
 &\hbox{passing tail}&\hbox{tail at }N_{\rm V11}-1\\ \hline
2^{-8} &21&13&192.088320&368.829802\\
2^{-9} & 8& 5&201.341968&786.243435\\
2^{-10}& 4& 3&208.682982&1706.400723\\
2^{-11}& 3& 3& 19.830558&326.731166.
\end{array}                                                     \tag{15.6}
\]
All passing tails are strictly below the complement budget and every
preceding tail is strictly above it.  The first three orders improve
strictly; the fourth is unchanged and no order worsens.
\end{theorem}

\begin{proof}
Use Theorem~\ref{thm:v9-neumann-tail} with (15.3), scan positive integer
orders, and compare fractions.  The executable independently verifies
passage at each displayed \(N\) and failure at \(N-1\).
\end{proof}

\begin{corollary}[Revised coefficient target]
\label{cor:v11-target}
At \(r=2^{-9}\), order \(5\) is minimal.  Its combined inverse-defect tail
is \(201.341968\), whereas order \(4\) gives \(786.243435\).  Hence the V10
provisional order \(8\) may be replaced by
\[
 \boxed{(r,N)=(2^{-9},5).}                                  \tag{15.7}
\]
\end{corollary}

\subsection{Support growth before coefficient arithmetic}

For finite frequency sets \(S,T\subset\mathbb Z^4\), write
\(S+T=\{s+t:s\in S,t\in T\}\).  The support of a matrix-product
coefficient dictionary is contained in this Minkowski sum.  This statement
is deliberately about a support envelope: algebraic cancellation may
remove frequencies, but cannot create a frequency outside it.

\begin{proposition}[Order-five Neumann support envelope]
\label{prop:v11-neumann-support}
Let \({\cal S}_{e,j}\) be the \(j\)-fold Minkowski sum of the active
Hessian support, with \({\cal S}_{e,0}=\{0\}\).  For
\(S_5=I+E+\cdots+E^4\), the successive envelope counts are
\[
\begin{array}{c|ccccc|c}
e&|{\cal S}_{e,0}|&|{\cal S}_{e,1}|&|{\cal S}_{e,2}|
 &|{\cal S}_{e,3}|&|{\cal S}_{e,4}|
 &|\bigcup_{j=0}^4{\cal S}_{e,j}|\\ \hline
{\rm W}&1&21&131&471&1251&1251\\
*&1&57&469&1859&5171&5171.
\end{array}                                                     \tag{15.8}
\]
Both unions lie in \([-4,4]^4\).
\end{proposition}

\begin{proof}
The verifier starts from \(\{0\}\) and performs exact integer-tuple
Minkowski addition against the nonzero coefficient supports.  The sets are
nested here, so the final union count equals the fourth-power count.
Each base frequency lies in \([-1,1]^4\), which also proves the range
inclusion.
\end{proof}

\begin{proposition}[Dense response-box obstruction]
\label{prop:v11-response-support}
After adding the factor supports in every tadpole and bubble word, the
order-five response envelopes lie in
\[
\begin{array}{c|c|r}
e&\hbox{frequency box}&\hbox{number of lattice slots}\\ \hline
{\rm W}&[-20,20]^4&2\,825\,761\\
*&[-22,22]\times[-22,22]\times[-20,20]\times[-22,22]
 &3\,736\,125.
\end{array}                                                     \tag{15.9}
\]
Therefore an implementation which allocates every matrix coefficient in
the bounding box is not a feasible exact compiler.
\end{proposition}

\begin{proof}
Each inverse slot contributes one \(S_5\) support and one response-factor
support.  Add their coordinate ranges for each of the seven words and take
the coordinatewise union with the tadpole range.  Multiplication of the
four integer side lengths gives the displayed slot counts.  Even one
double-complex \(45\times45\) matrix per slot would require more than
\(85\) GiB and \(112\) GiB respectively, before dictionary, exact-integer,
or intermediate-product overhead.  The actual rational objects are much
larger.
\end{proof}

\begin{firewall}[Envelope size is not nonzero-coefficient count]
\label{fire:v11-support-envelope}
Proposition~\ref{prop:v11-response-support} does not assert that every slot
is nonzero, nor does it obstruct a sparse scalar compiler.  It rejects only
the dense matrix-valued bounding-box strategy.
\end{firewall}

\subsection{The correct sparse signed algebra}

Let \(\mathcal A=M_{45}(\mathbb Q(i))[\mathbb Z^4]\) be the finite-support
convolution algebra.  For dictionaries \(U,V\in\mathcal A\), define
\[
 (U*V)_\nu=\sum_{\alpha+\beta=\nu}U_\alpha V_\beta.
                                                                  \tag{15.10}
\]

\begin{theorem}[Sparse trace-level compiler suffices]
\label{thm:v11-sparse-compiler}
Form \(S_5=\sum_{j=0}^4E^{*j}\) in \(\mathcal A\), but contract and trace
each complete response word as soon as its last matrix factor is inserted.
Then:
\begin{enumerate}[label=\textup{(\roman*)}]
\item the resulting scalar dictionary evaluates exactly to the V9
      order-five trial at every momentum;
\item zero matrix and zero scalar coefficients may be deleted after every
      convolution without changing that evaluation;
\item the endpoint-minus-Wilson scalar dictionary may be subtracted
      coefficientwise before any absolute value or box integral bound is
      taken.
\end{enumerate}
\end{theorem}

\begin{proof}
Evaluation
\[
 {\rm ev}_k(U)=\sum_\nu U_\nu e^{i\nu\cdot k}
\]
is an algebra homomorphism from convolution to matrix multiplication.
Hence the convolution powers reproduce the finite Neumann polynomial and
the seven complete words.  Deleting an exactly zero coefficient leaves
the represented Laurent polynomial unchanged.  Trace is linear, as is
endpoint-minus-Wilson subtraction, which proves the last assertion and
implements the SM compile-then-estimate rule from the V10 audit.
\end{proof}

\begin{assessment}[What V11 acquires]
\label{ass:v11-acquisition}
The factor-bound bottleneck is closed at this anchor: every used-image
stock is exact, all are much smaller than their factored predecessors, and
the target order falls from eight to five.  The newly exposed bottleneck is
representation, not analysis.  A dense matrix-valued frequency box is
prohibitively large, while Theorem~\ref{thm:v11-sparse-compiler} shows that
only an exact sparse scalar output is logically required.
\end{assessment}

\subsection{Executable provenance and the V12 upstream move}

\begin{record}[Used-image verifier]
\label{rec:v11-verifier}
The executable
\begin{center}
\path{scripts/verify_ym_postarchive_v11_used_image_stocks.py}
\end{center}
has \(490\) physical lines, \(16058\) bytes, and SHA--256
\[
\texttt{C3100394849AFF3898838A1EEBE40FC81A772AB1835A33D5F0DBE413AB7C77E4}.
                                                                  \tag{15.11}
\]
Its canonical payload has SHA--256
\[
\texttt{2854FA1EFB8E8E6326E511261BDBACFE1986AAE8D8747665E6752E5C6AA93943}.
                                                                  \tag{15.12}
\]
It transforms \(436\) active coefficient images with exact rational
matrix products, verifies the cubic-bank symmetry, proves every stock
comparison and order minimality statement, and computes the integer support
envelopes.
\end{record}

\begin{decision}[V12: rank-factor before sparse convolution]
\label{dec:v11-v12}
Before building a large polynomial, compute the exact or certified ranks of
the coefficient images
\[
 E_{e,\nu}=-A_eC_{e,\nu}Z_\sharp,\qquad
 B_{e,X,\nu}=A_eD_{X,\nu}Z_\sharp.
\]
If the images have low-rank factorizations, carry those factors through
the cyclic trace and frequency convolution at the target
\((2^{-9},5)\).  Use exact modular arithmetic and Chinese remaindering only
after an integer coefficient-height bound has been recorded.  If the
rational-shadow images are generically full rank, go upstream to the
block-diagonal-plus-rank-three V7 Woodbury coordinates rather than allocate
the dense boxes of Proposition~\ref{prop:v11-response-support}.  The first
executed scalar output must be the signed, reflection-paired
endpoint-minus-Wilson dictionary, not two separately normed responses.
\end{decision}

\begin{record}[V11 append record]
\label{rec:v11-append}
V11 is appended to the validated V10 whose installed SHA--256 was
\[
\texttt{341694A2E6AE928C7AF8AA30B335B9F2FFB85EBD6A614ADE0F6EFAF2D01D5B9D}.
                                                                  \tag{15.13}
\]
After validation only V11 is to remain the active numeric YM note.
\end{record}

\begin{firewall}[Claim boundary after V11]
\label{fire:v11-final}
V11 certifies one anchor's sharp factor stocks and support architecture; it
does not yet produce a signed coefficient dictionary, cover the complement,
or prove the continuous Haar-mean sign.  The nonzero Gaussian edge class,
compensating scale boundary, SCIC, ADMC, interacting cofinal state,
regulator removal, Osterwalder--Schrader reconstruction, physical
clustering, and Hamiltonian Yang--Mills mass gap remain open.
\end{firewall}

% =============================================================================
% END APPENDED POST-TENTH-ARCHIVE V11 MATERIAL
% =============================================================================

% =============================================================================
% BEGIN APPENDED POST-TENTH-ARCHIVE V12 MATERIAL
% =============================================================================

\section{V12: exact coefficient ranks and the hybrid cyclic compiler}
\label{sec:v12-rank-factor}

V11 reduced the local Neumann order from eight to five and showed that a
dense matrix-valued frequency box is the wrong representation.  Its next
decision was conditional: use rank factors if the transported Laurent
coefficients are genuinely low rank, but return to the V7 Woodbury
coordinates if they are generically full rank.  This continuation resolves
that condition exactly.  The answer is mixed.  Most coefficients are low
rank, a small family is full rank, and the storage-optimal input
representation is therefore hybrid rather than uniformly factorized.

\subsection{The actual Neumann bank and an exact rank certificate}

Write the Hessian Laurent expansion as
\[
 H_e(k)=\sum_{\nu\in\mathbb Z^4}C_{e,\nu}e^{i\nu\cdot k},
 \qquad A_e=R_eZ_\sharp^* .
\]
The frozen-chart Neumann defect is
\[
 E_e(k)=I-A_eH_e(k)Z_\sharp
       =\sum_\nu E_{e,\nu}e^{i\nu\cdot k},
\quad
 E_{e,\nu}=
 \begin{cases}
 I-A_eC_{e,0}Z_\sharp,&\nu=0,\\
 -A_eC_{e,\nu}Z_\sharp,&\nu\ne0.
 \end{cases}                                                   \tag{16.1}
\]
The identity in the zero mode is essential.  In particular,
\(\operatorname{rank}E_{e,0}\) cannot be inferred from the rank of
\(C_{e,0}\); V12 computes it directly.

Let \(\sigma(M)\) be the maximum matching number of the bipartite graph of
the exact nonzero pattern of a matrix \(M\).  Put
\[
 p=2^{31}-1.
\]
The verifier proves by trial division that \(p\) is prime and checks
\(p\equiv3\pmod4\).  Hence \(x^2+1\) is irreducible over
\(\mathbb F_p\), and
\(\mathbb F_p[i]=\mathbb F_p[x]/(x^2+1)\) is a field.

\begin{lemma}[Modular--structural rank sandwich]
\label{lem:v12-rank-sandwich}
Let \(M\in M_d(\mathbb Q(i))\), and suppose no denominator occurring in
\(M\) is divisible by \(p\).  If \(\overline M\) is its reduction in
\(M_d(\mathbb F_p[i])\), then
\[
 \operatorname{rank}_{\mathbb F_p[i]}\overline M
 \le \operatorname{rank}_{\mathbb Q(i)}M
 \le \sigma(M).                                                \tag{16.2}
\]
If the two outside quantities agree, they give the exact
\(\mathbb Q(i)\)-rank.
\end{lemma}

\begin{proof}
If an \(r\times r\) minor survives reduction modulo \(p\), its determinant
has numerator nonzero modulo \(p\), and therefore was nonzero in
\(\mathbb Q(i)\).  This proves the first inequality.  Every nonzero term
in the determinant of a minor selects nonzero entries in distinct rows and
columns.  Thus a nonzero \(r\times r\) minor supplies a matching of size
\(r\), proving the second inequality.
\end{proof}

For each coefficient the executable first applies
Lemma~\ref{lem:v12-rank-sandwich}.  In \(422\) of the \(436\) cases its two
bounds coincide.  In the remaining \(14\), genuine cancellation makes the
structural bound too large, so the executable performs complete Gaussian
elimination over \(\mathbb Q(i)\).  The exact rank in every such case
equals the modular lower bound.  Thus no rank below is a floating
singular-value decision.

\begin{theorem}[Rank preservation through the frozen chart]
\label{thm:v12-rank-preservation}
The rational matrices used by V11 satisfy
\[
 \operatorname{rank}Z_\sharp
 =\operatorname{rank}R_{\rm W}
 =\operatorname{rank}R_*
 =\operatorname{rank}A_{\rm W}
 =\operatorname{rank}A_*=45.                                  \tag{16.3}
\]
Consequently, for every response-factor coefficient \(D_{e,X,\nu}\),
\[
 \operatorname{rank}(A_eD_{e,X,\nu}Z_\sharp)
 =\operatorname{rank}D_{e,X,\nu},                              \tag{16.4}
\]
and for every nonzero Hessian frequency,
\(\operatorname{rank}E_{e,\nu}=\operatorname{rank}C_{e,\nu}\).
The two zero-frequency matrices \(E_{e,0}\) are both full rank.
\end{theorem}

\begin{proof}
The five rank equalities in (16.3) are separately certified by
Lemma~\ref{lem:v12-rank-sandwich}; each modular rank is \(45\), so every
matrix is invertible over \(\mathbb Q(i)\).  Multiplication on either side
by an invertible matrix preserves rank, which proves (16.4) and the
nonzero-frequency assertion from (16.1).  The verifier constructs each
zero mode in (16.1), including the identity, and certifies rank \(45\)
directly.
\end{proof}

\subsection{The complete mixed-rank ledger}

The base banks in the next table are \(E,Q_2,P_0,P_1,P_2\).
The minus banks need no second audit because V11 proved
\(M_0=P_0\), \(M_1=-P_1\), and \(M_2=P_2\) coefficientwise.
Likewise \(H_1\) and \(H_2\) are scalar frequency multiples of the same
nonzero Hessian coefficients and can share their factors.  In each row,
\(n\) is the frequency count, \(r_{\min},r_{\max}\) are the extreme exact
ranks, \emph{fac} counts ranks \(r\le22\), \emph{den} counts ranks
\(r\ge23\), \emph{full} counts rank \(45\), and \(\rho\) is the
hybrid-to-dense storage ratio.

\begin{theorem}[Exact coefficient-rank ledger]
\label{thm:v12-rank-ledger}
The Wilson ranks are
\[
\begin{array}{c|rrrrrr|c}
\text{bank}&n&r_{\min}&r_{\max}&\text{fac}&\text{den}&\text{full}&\rho\\ \hline
E   &21&3&45&20&1&1&341/945\\
Q_2 &11&3&39&10&1&0&197/495\\
P_0 & 9&3&14& 9&0&0&8/27\\
P_1 &11&3&36&10&1&0&59/165\\
P_2 &15&3&44&14&1&0&29/75
\end{array}                                                     \tag{16.5}
\]
and the endpoint ranks are
\[
\begin{array}{c|rrrrrr|c}
\text{bank}&n&r_{\min}&r_{\max}&\text{fac}&\text{den}&\text{full}&\rho\\ \hline
E   &57&1&45&48& 9&1&287/855\\
Q_2 &99&2&45&72&27&5&2107/4455\\
P_0 &51&2&38&45& 6&0&334/765\\
P_1 &81&1&44&60&21&0&1729/3645\\
P_2 &81&1&45&56&25&8&2017/3645 .
\end{array}                                                     \tag{16.6}
\]
Over all \(436\) audited coefficients, \(421\) are strictly low rank,
\(344\) select factors, \(92\) select retained coefficients, and exactly
\(15\) are full rank.
\end{theorem}

\begin{proof}
For every frequency the verifier records the modular pivot columns, the
matching upper bound, the exact nonzero count, and, in a cancellation case,
the rational-elimination pivots.  Grouping the \(436\) resulting exact
ranks gives (16.5)--(16.6).  The factor/dense split follows from
\(2(45)r<45^2\), which is equivalent to \(r\le22\).
\end{proof}

The full-rank family is small enough to state explicitly:
\[
\begin{split}
 &(e,E,0),\qquad e\in\{{\rm W},*\};\\
 &(*,Q_2,\nu),\qquad
   \nu\in\{0,\ \pm(1,0,0,0),\ \pm(0,0,1,0)\};\\
 &(*,P_2,\nu),\qquad
   \nu\in\{\pm e_j:1\le j\le4\}.
\end{split}                                                     \tag{16.7}
\]
There are no other rank-\(45\) base coefficients.

\subsection{Canonical factors and cyclic contraction}

\begin{lemma}[Exact pivot-column factorization]
\label{lem:v12-pivot-factor}
Let \(D\in M_{45}(\mathbb Q(i))\) have certified rank \(r\), and let
\(\mathcal P\) be the \(r\) pivot columns recorded by V12.  Then
\[
 D=UV,\qquad U=D[:,\mathcal P]\in M_{45\times r}(\mathbb Q(i)),
 \quad V\in M_{r\times45}(\mathbb Q(i)),                        \tag{16.8}
\]
with \(V\) unique.  For a transported response coefficient,
\[
 A_eDZ_\sharp=(A_eU)(VZ_\sharp).                               \tag{16.9}
\]
\end{lemma}

\begin{proof}
In a modular--structural certificate the recorded pivot-column minor is
nonzero modulo \(p\), hence nonzero over \(\mathbb Q(i)\).  In a
cancellation case the columns are exact rational-elimination pivots.
Thus the columns of \(U\) are independent.  Since their number equals
\(\operatorname{rank}D\), they form a basis of the column space of \(D\);
every column of \(D\) consequently has a unique coordinate column in that
basis.  Collecting those coordinates gives \(V\) and proves (16.8).
Associativity gives (16.9).
\end{proof}

\begin{theorem}[Rank-sized cyclic trace identity]
\label{thm:v12-cyclic-trace}
Suppose a frequency path in a complete response word contributes matrices
\(B_j=U_jV_j\), where \(U_j\) is \(45\times r_j\) and \(V_j\) is
\(r_j\times45\).  Then
\[
 \Tr(B_1B_2\cdots B_m)
 =
 \Tr_{r_1}\!\left[
 (V_1U_2)(V_2U_3)\cdots(V_{m-1}U_m)(V_mU_1)
 \right].                                                       \tag{16.10}
\]
The trace may be cyclically rotated so that the outside trace has the
smallest rank occurring on the path.  Formula (16.10) remains valid for a
dense-selected coefficient by taking \(r_j=45\), \(U_j=B_j\), and
\(V_j=I\).
\end{theorem}

\begin{proof}
Expand \(B_j=U_jV_j\), cyclically move \(U_1\) from the left end of the
trace to the right end, and use
\(\Tr_{45}(U_1T)=\Tr_{r_1}(TU_1)\).  All intermediate dimensions match.
Repeated cyclic rotation proves the final assertion.
\end{proof}

\begin{proposition}[Certified input-storage reduction]
\label{prop:v12-storage}
Storing all \(436\) base coefficients densely requires
\[
 436\cdot45^2=882\,900
\]
Gaussian-rational scalar entries.  Choosing two factors precisely when
\(r\le22\), and a dense coefficient otherwise, requires
\[
 396\,540,\qquad
 \frac{396\,540}{882\,900}=\frac{2203}{4905}
 =0.449133538\ldots .                                          \tag{16.11}
\]
Thus the certified hybrid input uses \(486\,360\) fewer scalar entries.
\end{proposition}

\begin{proof}
The cost of a rank-\(r\) pair in (16.8) is \(90r\), while the dense cost
is \(2025\).  Sum \(\min(2025,90r)\) over the exact ranks in
Theorem~\ref{thm:v12-rank-ledger}.  The verifier performs this integer sum
and reduces the resulting ratio.
\end{proof}

\begin{corollary}[No wholesale upstream retreat]
\label{cor:v12-no-retreat}
The V11 factor-or-return alternative is resolved in favour of a hybrid
factor compiler.  The \(15\) full-rank modes do not justify discarding the
\(344\) profitable factorizations; they are carried as explicit dense
exceptions.
\end{corollary}

\begin{firewall}[Storage is not convolution complexity]
\label{fire:v12-storage-only}
Proposition~\ref{prop:v12-storage} is an exact input-representation result,
not a bound on the number of frequency paths or on the rank after adding
several path products at one output frequency.  Such a sum can become full
rank.  The safe next compiler must therefore contract each complete path
to a scalar by Theorem~\ref{thm:v12-cyclic-trace}, accumulate only scalar
coefficients, and form endpoint-minus-Wilson and reflection pairs before
taking absolute values.
\end{firewall}

\subsection{Executable provenance and next compiler gate}

\begin{record}[Exact rank-factor verifier]
\label{rec:v12-verifier}
The executable
\begin{center}
\path{scripts/verify_ym_postarchive_v12_rank_factor_gate.py}
\end{center}
has \(536\) physical lines, \(17561\) bytes, and SHA--256
\[
\texttt{433B63470FDFB10D5334FC08210DA29106EB8CC53F3C07EB0E271FBE4B815577}.
                                                                  \tag{16.12}
\]
Its canonical payload has SHA--256
\[
\texttt{26428B537E058FFD44D1973E2525BC3D072780E8587B5672F42F593AA9BFA02C}.
                                                                  \tag{16.13}
\]
It certifies the primality and field condition, all five transport ranks,
all \(436\) coefficient ranks, every pivot set, the \(14\) exact
cancellation fallbacks, and the hybrid storage ledger.
\end{record}

\begin{decision}[V13: streamed signed derivative-free sector]
\label{dec:v12-v13}
Construct the exact factors (16.8), verify their reconstructions, and use
(16.10) at Neumann order five to compile the signed
endpoint-minus-Wilson scalar dictionary for the tadpole and the three
derivative-free bubble words
\[
 (P_2,M_0),\qquad(P_0,M_2),\qquad(P_1,M_1).
                                                                  \tag{16.14}
\]
Pair frequencies \(\nu\) and \(-\nu\) in that signed dictionary before
any absolute value.  Record actual scalar support, exact cancellation
counts, coefficient-height growth, and peak live state.  If path
enumeration is still prohibitive, go upstream to a meet-in-the-middle
cyclic contraction keyed by partial frequency and boundary rank indices;
do not revert to dense matrix boxes.
\end{decision}

\begin{record}[V12 append record]
\label{rec:v12-append}
V12 is appended to the validated V11 whose installed SHA--256 was
\[
\texttt{CA84CC493C55AAA8E4D931F0DA3D0B115E8F169619A9098F5A8F164E1291A565}.
                                                                  \tag{16.15}
\]
After validation only V12 is to remain the active numeric YM note.
\end{record}

\begin{firewall}[Claim boundary after V12]
\label{fire:v12-final}
V12 certifies the exact mixed-rank architecture at one frozen Gaussian
anchor.  It does not yet compile the signed order-five scalar response,
cover the complement, or prove the continuous Haar-mean sign.  The nonzero
Gaussian edge class, compensating scale boundary, SCIC, ADMC, interacting
cofinal state, regulator removal, Osterwalder--Schrader reconstruction,
physical clustering, and Hamiltonian Yang--Mills mass gap remain open.
\end{firewall}

% =============================================================================
% END APPENDED POST-TENTH-ARCHIVE V12 MATERIAL
% =============================================================================

% =============================================================================
% BEGIN APPENDED POST-TENTH-ARCHIVE V13 MATERIAL
% =============================================================================

\section{V13: constructed factors, cyclic collapse, and the path gate}
\label{sec:v13-factor-path}

V12 proved that \(344\) of the \(436\) base coefficients should be
represented by rank factors and that \(92\) should remain dense.  That was
an existence and rank result.  V13 constructs every selected factor,
checks every reconstruction over \(\mathbb Q(i)\), and then applies the
construction to the order-five derivative-free sector selected in V12.
The resulting path census is exact and forces a further algorithmic
change: complete-path streaming is combinatorially impossible even after
all presently available cyclic cancellations.

\subsection{The derivative-free sector collapses before compilation}

Put
\[
 S=S_{e,5}:=I+E_e+E_e^2+E_e^3+E_e^4,
\]
where products are matrix multiplication in the Laurent convolution
algebra.  Let \(B_{e,X}=A_eX_eZ_\sharp\), and suppress \(e\) and \(B\) in
the displayed notation.  The derivative-free part of the order-five trial
response is
\[
\begin{split}
 \Psi_e^{\rm df}
 ={}&\frac12\Tr(SQ_2)\\
 &+c_e\{\Tr(SP_2SM_0)+\Tr(SP_0SM_2)
                  +2\Tr(SP_1SM_1)\}.
\end{split}                                                     \tag{17.1}
\]
Here \(c_{\rm W}=1/3\) and \(c_*=3\).

\begin{theorem}[Exact cyclic collapse]
\label{thm:v13-cyclic-collapse}
The V11 coefficient identities imply the Laurent-polynomial identity
\[
 \boxed{\quad
 \Psi_e^{\rm df}
 =\frac12\Tr(SQ_2)
 +2c_e\{\Tr(SP_2SP_0)-\Tr(SP_1SP_1)\}.
 \quad}                                                         \tag{17.2}
\]
\end{theorem}

\begin{proof}
Substitute \(M_0=P_0\), \(M_2=P_2\), and \(M_1=-P_1\) in (17.1).
Cyclicity of the trace gives
\[
 \Tr(SP_0SP_2)=\Tr(SP_2SP_0).
\]
The first two bubble words therefore coincide and the third acquires a
minus sign.  This gives (17.2) coefficientwise, before evaluation,
integration, or any absolute value.
\end{proof}

The symbol \(P_1\) in (17.1)--(17.2) denotes the physically normalized
coefficient, including its fixed factor \(i\).  The canonical stored
factor bank below uses the raw rational \(P_1\) bank and attaches this
fixed Gaussian scalar as metadata; this does not change any rank,
frequency support, or reconstruction statement.

\subsection{Every selected factor is now explicit}

\begin{theorem}[Canonical factor-bank construction]
\label{thm:v13-factor-bank}
For each of the \(344\) coefficients selected by V12, let
\(\mathcal P\) be its certified pivot-column set and put
\[
 U=D[:,\mathcal P].
\]
There is an algorithm over \(\mathbb Q(i)\) which selects pivot rows
\(\mathcal I\), inverts the square minor
\(U[\mathcal I,:]\), and forms
\[
 V=(U[\mathcal I,:])^{-1}D[\mathcal I,:].
                                                                  \tag{17.3}
\]
For every selected coefficient the executable verifies exactly that
\[
 UV=D.                                                           \tag{17.4}
\]
It retains the other \(92\) coefficients without factorization, so the
partition exhausts all \(436\) base coefficients.
\end{theorem}

\begin{proof}
V12 proves that the columns indexed by \(\mathcal P\) form a basis of the
column space of \(D\).  Hence \(U\) has full column rank and contains an
invertible square row minor.  Exact Gaussian elimination on \(U^T\)
selects such a row set \(\mathcal I\).  Restricting a putative equality
\(UV=D\) to these rows forces (17.3), while the column-space property
proves the resulting equality on all rows.  The executable nevertheless
multiplies the two constructed factors and compares all \(45^2\) Gaussian
rational entries with \(D\), thereby checking (17.4) independently for
every coefficient.
\end{proof}

The construction absorbs the minus sign into each nonzero \(E\)-mode and
constructs the exceptional
\(E_{e,0}=I-A_eC_{e,0}Z_\sharp\) directly.  Its exact ledger is
\[
\begin{array}{c|rrrrr}
e&E&Q_2&P_0&P_1&P_2\\ \hline
{\rm W}\ {\rm factored}&20&10&9&10&14\\
{\rm W}\ {\rm dense}   & 1& 1&0& 1& 1\\
*\ {\rm factored}      &48&72&45&60&56\\
*\ {\rm dense}         & 9&27& 6&21&25 .
\end{array}                                                     \tag{17.5}
\]

\begin{corollary}[Small exact factor height]
\label{cor:v13-factor-height}
Across the \(344\) pairs \((U,V)\), the largest numerator has \(31\)
bits, the largest denominator has \(24\) bits, and there are \(7855\)
nonzero Gaussian-rational factor entries.  The canonical factor data have
SHA--256
\[
\texttt{0C5661C1FC2794CC210DF771B03B63B9FA048F87D2D94C64F12867AD1599E78A}.
                                                                  \tag{17.6}
\]
\end{corollary}

\begin{proof}
After each reconstruction check, serialize action, bank, frequency, rank,
pivot rows, and every real and imaginary numerator and denominator in a
fixed order.  Direct integer maxima give the height and sparsity data;
hashing that serialization gives (17.6).
\end{proof}

The factors need not be dressed separately.  For two nonexceptional
coefficients \(D_j=U_jV_j\), define the common bridge
\[
 K_e:=Z_\sharp A_e=Z_\sharp R_eZ_\sharp^* .
\]
Then the cyclic contraction from V12 uses the rank-sized bridge
\[
 (V_jZ_\sharp)(A_eU_{j+1})=V_jK_eU_{j+1}.                       \tag{17.7}
\]
Thus a single exact \(K_e\) suffices for all such adjacent factors.
Dense exceptions and \(E_{e,0}\) are inserted as explicit boundary
matrices.

\subsection{Exact support and path census}

Let \(m_e=|\operatorname{supp}E_e|\).  Before matrix cancellation, the
number of paths represented by \(S\) is
\[
 \Pi_e=\sum_{j=0}^4m_e^j.                                      \tag{17.8}
\]
Because \(0\in\operatorname{supp}E_e\), successive Minkowski powers are
nested.  Consequently
\[
 \operatorname{supp}S+\operatorname{supp}S
 =\bigcup_{j=0}^8j\,\operatorname{supp}E_e.                     \tag{17.9}
\]

\begin{theorem}[Order-five derivative-free candidate supports]
\label{thm:v13-support-census}
The exact combinatorial support census is
\[
\begin{array}{c|r|r|r|r|r|r|r}
e&m_e&|S|&|S+S|&|SQ_2|&|SP_2SP_0|&|SP_1SP_1|&|\mathcal U_e|\\ \hline
{\rm W}&21&1251&15421&2621&34635&34347&35069\\
*&57&5171&67789&18141&238795&260967&261797 .
\end{array}                                                     \tag{17.10}
\]
Here \(\mathcal U_e\) is the union of the three candidate supports in
(17.2).  It is reflection invariant and has respectively \(17535\) and
\(130899\) orbits under \(\nu\mapsto-\nu\).  The signed
endpoint-minus-Wilson candidate union has \(261797\) frequencies and
\(130899\) reflection orbits.
\end{theorem}

\begin{proof}
Starting with \(\{0\}\), the verifier forms exact integer-tuple Minkowski
sums with the \(E\)-support through powers four and eight.  It then adds
the exact \(Q_2,P_0,P_1,P_2\) supports as prescribed by (17.2).
Set union gives \(\mathcal U_e\).  Direct membership testing proves
\(-\mathcal U_e=\mathcal U_e\).  Only \(0\) is fixed by negation in
\(\mathbb Z^4\), so the orbit count is
\((|\mathcal U_e|+1)/2\).
\end{proof}

\begin{theorem}[Exact raw path counts]
\label{thm:v13-path-count}
The order-five path counts before matrix cancellation are
\[
\begin{array}{c|r|r}
e&{\rm W}&*\\ \hline
\Pi_e&204205&10744501\\
\text{tadpole}&2246255&1063705599\\
\text{three original bubbles}
 &16304575671775&1711230884677211823\\
\text{two collapsed bubbles}
 &10675118598400&1234330474193398692\\
\text{paths removed by (17.2)}
 &5629457073375&476900410483813131 .
\end{array}                                                     \tag{17.11}
\]
\end{theorem}

\begin{proof}
The tadpole count is \(\Pi_e|\operatorname{supp}Q_2|\).  Before the cyclic
collapse, the three bubble words contribute
\[
 \Pi_e^2\left(
 2|\operatorname{supp}P_2||\operatorname{supp}P_0|
 +|\operatorname{supp}P_1|^2\right).
\]
After (17.2), only one \(P_2P_0\) path family and one \(P_1P_1\) path
family need be evaluated, giving
\[
 \Pi_e^2\left(
 |\operatorname{supp}P_2||\operatorname{supp}P_0|
 +|\operatorname{supp}P_1|^2\right).                           \tag{17.12}
\]
Substitution of the exact support cardinalities gives (17.11).
\end{proof}

\begin{countertheorem}[Literal complete-path streaming is rejected]
\label{counter:v13-literal-stream}
Even after the exact cyclic collapse, literal enumeration of every
complete derivative-free path is not a viable compiler: it requires more
than \(10^{13}\) Wilson contractions and more than \(10^{18}\) endpoint
contractions.
\end{countertheorem}

\begin{proof}
This is the fourth row of (17.11), compared with the stated integer
thresholds.  The conclusion does not depend on matrix-multiplication
timing estimates.
\end{proof}

\begin{firewall}[Candidate support is not scalar support]
\label{fire:v13-candidate}
The counts in (17.10) are exact Minkowski candidate supports, not claims
that all corresponding matrix or scalar coefficients survive.  The path
counts in (17.11) precede cancellations among paths with the same total
frequency.  V13 therefore neither reports an actual signed scalar support
nor uses these counts as an \(L^1\) bound.
\end{firewall}

\subsection{The frequency-aggregated upstream move}

For
\[
 L_{e,X}=SB_{e,X}=\sum_\mu L_{e,X,\mu}e^{i\mu\cdot k},
\]
the two bubble coefficients in (17.2) obey
\[
\begin{split}
 [\Tr(L_{e,P_2}L_{e,P_0})]_\nu
   &=\sum_\mu\Tr(L_{e,P_2,\mu}L_{e,P_0,\nu-\mu}),\\
 [\Tr(L_{e,P_1}L_{e,P_1})]_\nu
   &=\sum_\mu\Tr(L_{e,P_1,\mu}L_{e,P_1,\nu-\mu}).
\end{split}                                                     \tag{17.13}
\]
This is the correct meet-in-the-middle identity: partial paths are
aggregated by frequency before the two halves are paired.  The V12 cyclic
rank contraction and (17.7) remain available inside each aggregated
update.

\begin{decision}[V14: certified frequency aggregation]
\label{dec:v13-v14}
Do not enumerate the paths in (17.11).  Cache the canonical factor bank
under digest (17.6), derive a common-denominator and numerator-height bound
for every coefficient in (17.13), and only then select enough
Gaussian finite-field primes for exact Chinese reconstruction.  Implement
the order-five recurrence frequency by frequency, compile the two
endpoint-minus-Wilson combinations in (17.2), and pair \(\nu\) with
\(-\nu\) before any absolute value.  The implementation must report peak
live states and distinguish modular zeros from exact reconstructed zeros.
If the aggregated boundary matrices become effectively dense, return
upstream to a block/Schur contraction of the common bridge \(K_e\).
\end{decision}

\subsection{Executable provenance and claim boundary}

\begin{record}[V13 factor-and-path verifier]
\label{rec:v13-verifier}
The executable
\begin{center}
\path{scripts/verify_ym_postarchive_v13_factor_path_gate.py}
\end{center}
has \(436\) physical lines, \(15568\) bytes, and SHA--256
\[
\texttt{A3BEC0A9068ACD53BFD768937F4051762AD9C1BB19D161C68C7BA67A0EE816C3}.
                                                                  \tag{17.14}
\]
Its canonical payload has SHA--256
\[
\texttt{2C9AE57EFC534080B0B069984DE9B626DE92E96FD5506FA2944A0C7067019B90}.
                                                                  \tag{17.15}
\]
It reconstructs every selected factor over \(\mathbb Q(i)\), certifies
the factor digest and height, proves the cubic/cyclic collapse, and
computes the support and path ledgers solely with exact arithmetic.
\end{record}

\begin{record}[V13 append record]
\label{rec:v13-append}
V13 is appended to the validated V12 whose installed SHA--256 was
\[
\texttt{1A1D41A289DA206CCB88023306DD33743C46058F731345900765A724BA1DDCBE}.
                                                                  \tag{17.16}
\]
After validation only V13 is to remain the active numeric YM note.
\end{record}

\begin{firewall}[Claim boundary after V13]
\label{fire:v13-final}
V13 closes the factor-construction question and rejects literal path
streaming, but it does not yet produce the signed order-five scalar
dictionary, cover the complement, or prove the continuous Haar-mean sign.
The nonzero Gaussian edge class, compensating scale boundary, SCIC, ADMC,
interacting cofinal state, regulator removal,
Osterwalder--Schrader reconstruction, physical clustering, and
Hamiltonian Yang--Mills mass gap remain open.
\end{firewall}

% =============================================================================
% END APPENDED POST-TENTH-ARCHIVE V13 MATERIAL
% =============================================================================

% =============================================================================
% BEGIN APPENDED POST-TENTH-ARCHIVE V14 MATERIAL
% =============================================================================

\section{V14: exact coefficient height and the finite CRT gate}
\label{sec:v14-height-crt}

V13 proved that literal enumeration of complete order-five paths is
impossible and selected frequency aggregation.  Exact aggregation over
\(\mathbb Q(i)\) risks coefficient swell, while modular aggregation is
rigorous only after one knows how many primes guarantee unique
reconstruction.  This continuation supplies that missing height theorem.
It compiles the actual frozen-chart coefficient banks, clears exact common
denominators, derives a uniform numerator bound for every signed
endpoint-minus-Wilson coefficient, and constructs a finite minimal-prefix
list of reconstruction primes.

\subsection{A submultiplicative Laurent height}

For a finite matrix Laurent polynomial
\[
 F(z)=\sum_{\nu}F_\nu z^\nu,\qquad
 F_\nu\in M_d(\mathbb Z[i]),
\]
define
\[
 {\cal A}(F):=
 \sum_\nu\max_{1\le i,j\le d}
 \bigl(|\Re(F_{\nu,ij})|+|\Im(F_{\nu,ij})|\bigr).              \tag{18.1}
\]
For a scalar Laurent polynomial \(f\), let
\(\|f\|_{1,i}\) be the same sum without the matrix maximum.

\begin{lemma}[Integer Laurent-height calculus]
\label{lem:v14-height-calculus}
For finite \(F,G\in M_d(\mathbb Z[i])[\mathbb Z^4]\),
\[
 {\cal A}(FG)\le d\,{\cal A}(F){\cal A}(G),\qquad
 \|\Tr(FG)\|_{1,i}\le d^2{\cal A}(F){\cal A}(G).               \tag{18.2}
\]
In particular, every real and imaginary component of every coefficient
of \(\Tr(FG)\) is bounded in absolute value by the right side of the
second inequality.
\end{lemma}

\begin{proof}
The Gaussian integer norm
\(|a+ib|_1=|a|+|b|\) is submultiplicative.  In each entry of \(FG\),
sum over the \(d\) matrix indices and over all frequency decompositions
\(\alpha+\beta=\nu\), then sum over \(\nu\).  Replacing every entry by the
maximum at its frequency gives the first inequality.  The trace of \(FG\)
is a sum over the \(d^2\) pairs \((i,j)\), which gives the second.
\end{proof}

\subsection{Actual frozen-chart bank heights}

For each action \(e\) and base bank \(X\in\{E,Q_2,P_0,P_1,P_2\}\), V14
chooses the least positive common denominator \(D_{e,X}\) encountered by
the exact compiler and writes
\[
 B_{e,X}=D_{e,X}^{-1}\widehat B_{e,X},
 \qquad \widehat B_{e,X}\in
 M_{45}(\mathbb Z[i])[\mathbb Z^4],\qquad
 A_{e,X}:={\cal A}(\widehat B_{e,X}).                          \tag{18.3}
\]
These are the actual coefficients \(A_eD_{X,\nu}Z_\sharp\), not raw
vertex coefficients.  The bank \(P_1\) includes its physical factor \(i\);
the zero defect mode is
\(E_{e,0}=I-A_eC_{e,0}Z_\sharp\).

\begin{theorem}[Exact bank denominator and height ledger]
\label{thm:v14-bank-ledger}
The Wilson data are
\[
\begin{array}{c|r|r}
X&D_{{\rm W},X}&A_{{\rm W},X}\\ \hline
E   &4\,800\,000\,000&53\,262\,501\,420\\
Q_2 &57\,600\,000\,000&1\,769\,396\,444\,344\\
P_0 &1\,600\,000\,000&19\,848\,214\,416\\
P_1 &12\,800\,000\,000&99\,353\,570\,640\\
P_2 &76\,800\,000\,000&941\,669\,632\,056
\end{array}                                                     \tag{18.4}
\]
and the endpoint data are
\[
\begin{array}{c|r|r}
X&D_{*,X}&A_{*,X}\\ \hline
E   &4\,800\,000\,000&107\,623\,215\,368\\
Q_2 &1\,385\,280\,000\,000\,000&190\,661\,226\,249\,146\,240\\
P_0 &9\,600\,000\,000&68\,991\,664\,712\\
P_1 &38\,400\,000\,000&288\,359\,156\,884\\
P_2 &5\,541\,120\,000\,000\,000&54\,642\,974\,339\,593\,144 .
\end{array}                                                     \tag{18.5}
\]
\end{theorem}

\begin{proof}
For every one of the \(436\) frequencies, the verifier performs the exact
rational products defining the frozen-chart image.  It takes the least
common multiple of all real and imaginary entry denominators in a bank.
Multiplication by this denominator is checked to make
\(D_{e,X}(|\Re b|+|\Im b|)\) integral for every entry.  It then takes the
entry maximum at each frequency and sums, exactly implementing (18.1).
\end{proof}

\subsection{Height of the order-five signed response}

Put \(d=45\), \(D_e=D_{e,E}\), and \(A_e=A_{e,E}\).  Clearing
\[
 S_e=I+E_e+E_e^2+E_e^3+E_e^4
\]
to denominator
\[
 D_{e,S}=D_e^4                                               \tag{18.6}
\]
and applying Lemma~\ref{lem:v14-height-calculus} gives
\[
 A_{e,S}:=
 D_e^4+\sum_{j=1}^4D_e^{4-j}d^{j-1}A_e^j
 \ \ge\ {\cal A}(D_{e,S}S_e).                                 \tag{18.7}
\]
For \(L_{e,X}=S_eB_{e,X}\), one may therefore take
\[
 D_{e,L_X}=D_{e,S}D_{e,X},\qquad
 A_{e,L_X}=dA_{e,S}A_{e,X}.                                   \tag{18.8}
\]

\begin{theorem}[Uniform signed coefficient bound]
\label{thm:v14-signed-height}
Let \(\psi_\nu\in\mathbb Q(i)\) be the coefficient at frequency \(\nu\)
of the endpoint-minus-Wilson derivative-free scalar in (17.2).  There
are Gaussian integers \(N_\nu\) such that
\[
 \psi_\nu=\frac{N_\nu}{\Delta},\qquad
 |\Re N_\nu|+|\Im N_\nu|\le B,                                \tag{18.9}
\]
where
\[
 \Delta
 =2^{126}3^{11}5^{82}\cdot13\cdot37                           \tag{18.10}
\]
has \(343\) bits and
\[
\begin{split}
 B={}&
24\,959\,593\,000\,859\,949\,653\,031\,106\,227\,650\,636\,323\,739\,452
 \cdot10^{90}\\
&+
918\,818\,268\,888\,098\,979\,461\,819\,918\,386\,592\,460\,769\,601\,597
 \cdot10^{45}\\
&+
612\,427\,838\,226\,993\,950\,792\,852\,785\,626\,808\,320\,000\,000\,000 .
\end{split}                                                     \tag{18.11}
\]
Equivalently, \(B\) is the \(444\)-bit integer whose decimal SHA--256 is
\[
\texttt{144B73BC8DB36A043EB6A7F4F0016011B9FB61D0549A3321302DF82B4005D420}.
                                                                  \tag{18.12}
\]
\end{theorem}

\begin{proof}
The tadpole in (17.2), after denominator clearing, has numerator
\(\ell^1\)-height at most
\[
 d^2A_{e,S}A_{e,Q_2}
\]
and denominator \(2D_{e,S}D_{e,Q_2}\).  If
\(2c_e=n_e/q_e\) in lowest terms, the \(XY\) bubble has numerator height
at most
\[
 |n_e|d^2A_{e,L_X}A_{e,L_Y}
\]
and denominator
\[
 q_eD_{e,L_X}D_{e,L_Y}.
\]
Apply this to \(XY=P_2P_0\) and \(P_1P_1\), take the least common
denominator within each action, and then the least common denominator
between endpoint and Wilson.  Scaling and adding the three action bounds
and then the two action bounds yields (18.9).  Direct integer arithmetic
gives the factorization (18.10) and the exact decimal integer (18.11).
The executable additionally hashes its ungrouped decimal representation,
giving (18.12).
\end{proof}

\begin{remark}[Reading the grouped integer]
\label{rem:v14-B-reading}
Equation (18.11) is a base-ten block decomposition, not a floating
approximation.  The exact unbroken decimal value is retained in the
verifier payload.  The Wilson and endpoint action bounds have respectively
\((418,444)\) numerator bits and \((328,343)\) denominator bits.
\end{remark}

\subsection{A finite exact reconstruction list}

Let \(p_1,\ldots,p_{15}\) be
\[
\begin{array}{rrrrr}
2147483647&2147483587&2147483579&2147483563&2147483543\\
2147483423&2147483399&2147483323&2147483179&2147483171\\
2147483123&2147483059&2147482951&2147482943&2147482867 .
\end{array}                                                     \tag{18.13}
\]

\begin{theorem}[Certified CRT sufficiency and minimal prefix]
\label{thm:v14-crt}
Every number in (18.13) is a distinct prime congruent to \(3\bmod4\) and
is coprime to \(\Delta\).  If
\[
 M=\prod_{\ell=1}^{15}p_\ell,
\]
then
\[
 M>2B,\qquad \operatorname{bits}(M)=465.                       \tag{18.14}
\]
The product of the first fourteen primes is at most \(2B\).  Hence the
fifteen-prime prefix is sufficient and minimal within this deterministic
descending list.
\end{theorem}

\begin{proof}
The verifier trial-divides each candidate by every integer up to its
square root, checks its residue modulo \(4\), and checks its gcd with
\(\Delta\).  Exact multiplication gives (18.14), with \(20\) bits of
excess over \(2B\).  Exact division of \(M\) by \(p_{15}\) proves failure
of the preceding prefix.  The prime-list SHA--256 is
\[
\texttt{9F69CC2F4FEB1C26BA37E0DB82BD00CA36B4414E60AADF6B3CA83CC87B89DA2F},
                                                                  \tag{18.15}
\]
and the ungrouped decimal SHA--256 of \(M\) is
\[
\texttt{2658F1E3EDC2BD671113BBAD98E64A86C456FC6A77744906D18253BEB6221420}.
                                                                  \tag{18.16}
\]
\end{proof}

\begin{corollary}[Exact reconstruction of every signed coefficient]
\label{cor:v14-reconstruction}
Suppose the real and imaginary components of every numerator \(N_\nu\)
in (18.9) are computed modulo all fifteen primes in (18.13).  Centered
integer CRT reconstructs every \(N_\nu\) uniquely and exactly.
\end{corollary}

\begin{proof}
CRT determines each component modulo \(M\).  Two integers in
\([-B,B]\) with the same residues differ by a multiple of \(M\), but
their difference has absolute value at most \(2B<M\).  They are therefore
equal.
\end{proof}

\begin{assessment}[The arithmetic-height bottleneck is closed]
\label{ass:v14-height-closed}
The coefficient-height question is no longer an obstacle: exact
reconstruction requires only fifteen \(31\)-bit prime runs.  The live
bottleneck is the frequency-aggregated matrix contraction within one
prime.  V14 does not infer feasibility of that contraction merely from
the modest CRT count.
\end{assessment}

\begin{firewall}[Height is not execution]
\label{fire:v14-no-execution}
Theorem~\ref{thm:v14-crt} proves uniqueness conditional on correct modular
residues.  V14 has not yet computed those residues, actual scalar support,
reflection-paired box integrals, or the Haar-mean sign.  The bound \(B\)
is deliberately global and cancellation-free; it certifies reconstruction
but is not itself an analytic error estimate.
\end{firewall}

\subsection{Executable provenance and the scheduled V15 audit}

\begin{record}[V14 height-and-CRT verifier]
\label{rec:v14-verifier}
The executable
\begin{center}
\path{scripts/verify_ym_postarchive_v14_height_crt_gate.py}
\end{center}
has \(348\) physical lines, \(12095\) bytes, and SHA--256
\[
\texttt{B48CDFE6BC95E25322319615E57C07F05504516735E5AA0839C3B27346549707}.
                                                                  \tag{18.17}
\]
Its canonical payload has SHA--256
\[
\texttt{850A99369956AB10C94A0E1D847D99055D277EDE4B93088B0952D873AF6086E5}.
                                                                  \tag{18.18}
\]
It compiles all actual coefficient images, checks every denominator
clearance, derives all integer bounds, proves the fifteen primes by trial
division, and verifies the strict CRT inequalities.
\end{record}

\begin{decision}[V15: cross-programme audit, then one-prime aggregation]
\label{dec:v14-v15}
In accordance with the five-version rule, V15 must first inspect the
current SM and NS notes and import only mathematically applicable
techniques.  It should then cache the V13 factor bank and execute the
frequency-aggregated recurrence for one prime in (18.13), measuring
support, cancellation, boundary-matrix rank, and peak live memory for the
two terms in (17.13).  Only after that one-prime result is checked against
direct evaluation at several finite-field phase points should the
remaining fourteen primes be run.  Endpoint-minus-Wilson subtraction and
reflection pairing remain prior to every absolute value.
\end{decision}

\begin{record}[V14 append record]
\label{rec:v14-append}
V14 is appended to the validated V13 whose installed SHA--256 was
\[
\texttt{8FE2899FEB248D8398E1A18B0A73D2C135C035C86B4249753FA12330DDBD9F27}.
                                                                  \tag{18.19}
\]
After validation only V14 is to remain the active numeric YM note.
\end{record}

\begin{firewall}[Claim boundary after V14]
\label{fire:v14-final}
V14 closes the arithmetic-height and CRT-uniqueness gate for the
derivative-free order-five sector.  It does not yet compute its signed
dictionary, cover the complement, or prove the continuous Haar-mean sign.
The derivative-containing sector, nonzero Gaussian edge class,
compensating scale boundary, SCIC, ADMC, interacting cofinal state,
regulator removal, Osterwalder--Schrader reconstruction, physical
clustering, and Hamiltonian Yang--Mills mass gap remain open.
\end{firewall}

% =============================================================================
% END APPENDED POST-TENTH-ARCHIVE V14 MATERIAL
% =============================================================================

% =============================================================================
% BEGIN APPENDED POST-TENTH-ARCHIVE V15 MATERIAL
% =============================================================================

\section{V15: companion audit and a one-prime aggregate-rank gate}
\label{sec:v15}

V14 made exact reconstruction finite but did not establish that one modular
run fits a mathematically honest representation.  The present version first
performs the required five-version companion audit.  That audit changes one
important detail of the computation: low rank may be used only for the
\emph{aggregate coefficient after all paths to its frequency have been
summed}.  V15 then executes this rule on the first nontrivial Neumann
product \(E^2\), for both frozen actions and at the first V14 prime.  This is
an early state-growth gate, not yet the order-five scalar response.

\subsection{Compulsory audit of the current SM and NS notes}

The SM file advanced while this audit was in progress, so the final
read-only snapshot was
\begin{center}
\path{latex_notes/Renewal_SM_Einstein_Continuum_Research_Notes_V78.tex}\\
\(24313\) lines,\quad \(860587\) bytes,\\
\(\operatorname{SHA256}=\)
\texttt{59B7A9F2B3310807C706340D663BABC9939F484995A50C88C962D8E7D206DAA2}.
\end{center}
The current NS snapshot was
\begin{center}
\path{latex_notes/Renewal_NS_Stability_Research_Notes_V26.tex}\\
\(5319\) lines,\quad \(278283\) bytes,\\
\(\operatorname{SHA256}=\)
\texttt{82C887F8C2C69E4F28FAD7C3DC8DE5B9099CB5A4BF431EBA1FFF3E91935978AD}.
\end{center}

Three finite-algebra rules in the cumulative SM V78 file apply to the
present task.  SM V71 permits hybrid factor/dense storage only after an
exact rank certificate and warns that subsequent operations may densify.
SM V73 keeps shared occurrences in one sparse polynomial object and
performs signed subtraction before interval evaluation.  SM V75, itself
responding to YM V13, replaces literal path enumeration by aggregation
through the relevant finite operator before taking a norm.  SM V78 proves
a local Banach contraction for a different coexact nonabelian problem.  Its
field theorem is not imported here; its cofinal-uniformity firewall remains
relevant later, after the finite response is known.

NS V26 gives the complementary warning.  Its rank-one norm improves one
kernel derivative but not another, and only because the input is proved to
belong to the stated rank-one class.  A sum of structured inputs need not
remain in that class.  Translated to the present finite algebra, V12 ranks
of the individual coefficient images do not certify ranks of their
frequency sums.

\begin{theorem}[V15 companion-audit decision]
\label{thm:v15-audit}
No SM or NS field theorem is used as a Yang--Mills theorem.  The applicable
certificate rule is the following:
\begin{enumerate}[leftmargin=2.2em]
\item compile every ordered matrix path with its sign in the common Laurent
      dictionary;
\item sum all paths having the same frequency;
\item only then compute the rank of the resulting boundary matrix and
      choose factor or dense storage; and
\item postpone endpoint-minus-Wilson subtraction, reflection pairing, and
      every absolute value until the common scalar dictionary exists.
\end{enumerate}
\end{theorem}

\begin{proof}
Items one and two are the evaluation-homomorphism and aggregation lessons
of SM V73 and V75.  Item three is forced by the SM V71 and NS V26
used-class restrictions: rank is not additive and the sum of rank-\(r\)
matrices can have larger rank.  Item four is the V13--V14 signed-order rule
and is unchanged by the audit.  These are statements about the design of
the finite certificate, not imports of conclusions from the companion
equations.
\end{proof}

\subsection{The one-prime coefficient object}

Set
\[
 p=2147483647=2^{31}-1,\qquad
 k=\mathbb F_p[\iota]/(\iota^2+1).                         \tag{19.3}
\]
V14 certified that \(p\) is prime, \(p\equiv3\bmod4\), and \(p\) divides
none of the cleared denominators.  Since \(-1\) is a nonsquare in
\(\mathbb F_p\), \(k\) is the field \(\mathbb F_{p^2}\).

For action \(e\in\{W,\partial\}\), reduce its actual V14 residual bank
coefficientwise:
\[
 E_e(z)=\sum_{\alpha\in A_e}E_{e,\alpha}z^\alpha
 \in M_{45}\bigl(k[z_1^{\pm1},\ldots,z_4^{\pm1}]\bigr).     \tag{19.4}
\]
The aggregate square coefficient is
\[
 G^{(2)}_{e,\nu}
   :=\sum_{\substack{\alpha,\beta\in A_e\\\alpha+\beta=\nu}}
       E_{e,\alpha}E_{e,\beta}.                            \tag{19.5}
\]

\begin{lemma}[Frequency aggregation is the exact square]
\label{lem:v15-convolution}
One has
\[
 E_e(z)^2=\sum_\nu G^{(2)}_{e,\nu}z^\nu.                   \tag{19.6}
\]
Thus a rank computed from \(G^{(2)}_{e,\nu}\) is a rank of the used
frequency image, whereas ranks of the individual path products in
(19.5) are not substitutes for it.
\end{lemma}

\begin{proof}
Distribute the finite product in (19.4).  Laurent monomials multiply by
\(z^\alpha z^\beta=z^{\alpha+\beta}\).  Collecting equal exponents gives
(19.5)--(19.6).  No rearrangement of an infinite sum is involved.
\end{proof}

\subsection{Exact modular matrix multiplication}

The executable uses binary64 matrix multiplication only as an exact
integer dot-product engine.  Write each residue uniquely as
\[
 a=a_0+2^{16}a_1,\qquad
 0\le a_0<2^{16},\quad0\le a_1<2^{15}.                    \tag{19.7}
\]

\begin{lemma}[Exact two-limb product]
\label{lem:v15-limb}
Every length-\(45\) dot product required by the two-limb Karatsuba formula
is an exactly represented binary64 integer.  The four relevant upper
bounds are
\[
\begin{array}{c|r}
\hbox{dot type}&\hbox{strict upper bound}\\ \hline
00&193267630125\\
\hbox{two mixed dots}&193264681050\\
11&48315433005\\
(0+1)(0+1)&434847744180 ,
\end{array}                                                 \tag{19.8}
\]
all strictly smaller than \(2^{53}=9007199254740992\).
Moreover
\[
 (2^{16})^2\equiv2\pmod p.                                \tag{19.9}
\]
Consequently the executable's real and Gaussian matrix products equal
matrix multiplication over \(k\), entry for entry.
\end{lemma}

\begin{proof}
The largest possible summands and the length \(45\) give the four integer
bounds in (19.8).  Binary64 represents every integer of magnitude at most
\(2^{53}\) exactly; every partial sum is nonnegative and bounded by its
corresponding total.  Therefore the three real matrix products used to
recover the low, high, and mixed limbs are exact, including a fused
multiply-add implementation.  Recombination uses
\[
 c_{00}+2^{16}(c_{01}+c_{10})+2c_{11}\pmod p,
\]
by (19.9).  Finally, for Gaussian matrices the identity
\[
 (A+\iota B)(C+\iota D)
 =(AC-BD)+\iota\bigl((A+B)(C+D)-AC-BD\bigr)
\]
uses three already exact real products.  Hence the returned residues are
the exact product in \(k\).
\end{proof}

\begin{lemma}[Aggregate rank]
\label{lem:v15-rank}
For each frequency, the V15 elimination returns exactly the rank over \(k\)
of the completed aggregate matrix \(G^{(2)}_{e,\nu}\).
\end{lemma}

\begin{proof}
For a nonzero pivot \(a+\iota b\), its norm
\(a^2+b^2\) is nonzero in \(\mathbb F_p\), since \(k\) is a field, and
\[
 (a+\iota b)^{-1}=(a-\iota b)(a^2+b^2)^{-1}.
\]
The algorithm normalizes that pivot and eliminates entries below it using
only the exact field operations of Lemma~\ref{lem:v15-limb}.  Standard
row-echelon induction therefore makes the returned pivot count equal to
the rank over \(k\).
\end{proof}

\subsection{Executed aggregate census}

All \(21^2=441\) Wilson paths and all \(57^2=3249\) endpoint paths in
(19.5) were accumulated.  Whole-matrix zero testing and rank elimination
were performed only after the sums were complete.

\begin{theorem}[Executed one-prime \(E^2\) gate]
\label{thm:v15-executed}
At the prime (19.3), exact execution gives
\[
\begin{array}{c|rr}
 &W&\partial\\ \hline
|A_e|&21&57\\
\hbox{ordered paths}&441&3249\\
\hbox{candidate frequencies}&131&469\\
\hbox{nonzero aggregate matrices}&131&467\\
\hbox{whole-matrix modular cancellations}&0&2\\
\hbox{maximum paths merged at one frequency}&21&57\\
\hbox{factor-profitable aggregate matrices}&110&372\\
\hbox{dense aggregate matrices}&21&95\\
\hbox{full-rank aggregate matrices}&1&41 .
\end{array}                                                  \tag{19.10}
\]
The two zero endpoint frequencies are
\[
 (-2,-2,2,0),\qquad(2,2,-2,0).                             \tag{19.11}
\]
The complete nonzero rank distributions \(r:n_r\) are
\[
\begin{split}
W:\;&2{:}4,\ 3{:}8,\ 4{:}6,\ 5{:}22,\ 6{:}20,\ 8{:}6,\
10{:}8,\ 11{:}20,\ 14{:}6,\ 19{:}8,\ 20{:}2,\\
&25{:}2,\ 26{:}8,\ 28{:}2,\ 34{:}8,\ 45{:}1.
\end{split}                                                  \tag{19.12}
\]
\[
\begin{split}
\partial:\;&1{:}36,\ 2{:}120,\ 4{:}24,\ 5{:}20,\ 6{:}36,\
7{:}8,\ 8{:}10,\ 9{:}4,\ 11{:}12,\ 12{:}22,\\
&14{:}40,\ 17{:}4,\ 18{:}8,\ 19{:}6,\ 20{:}8,\ 21{:}2,\
22{:}12,\ 23{:}12,\ 26{:}8,\ 28{:}4,\\
&30{:}2,\ 35{:}2,\ 36{:}2,\ 37{:}2,\ 38{:}10,\ 40{:}6,\
41{:}2,\ 42{:}4,\ 45{:}41.
\end{split}
                                                               \tag{19.13}
\]
Here a factor is retained precisely when \(90r<2025\), i.e.\ \(r\le22\).
\end{theorem}

\begin{proof}
Lemma~\ref{lem:v15-convolution} identifies the dictionary being computed.
Lemma~\ref{lem:v15-limb} proves every path product and accumulation exact
modulo \(p\), and Lemma~\ref{lem:v15-rank} proves the rank counts.  The
program independently checks twelve scalar entries against a direct
length-\(45\) Gaussian loop for each action.  At each of six deterministic
quarter-turn points \(t\in(\mathbb Z/4)^4\), it also verifies
\[
 \sum_\nu G^{(2)}_{e,\nu}\iota^{\nu\cdot t}
 =\left(\sum_\alpha E_{e,\alpha}\iota^{\alpha\cdot t}\right)^2.
                                                               \tag{19.14}
\]
These checks validate the implementation of the coefficient identity;
the identity itself follows for all Laurent variables from
Lemma~\ref{lem:v15-convolution}.  Counting the resulting ranks gives
(19.10)--(19.13).
\end{proof}

\begin{corollary}[The aggregate state is still compressible]
\label{cor:v15-storage}
Dense storage of every nonzero \(45\)-square coefficient would use
\[
265275\quad\hbox{and}\quad945675
\]
Gaussian entries for Wilson and endpoint.  The post-aggregation hybrid
gate uses respectively
\[
124425=\frac{553}{1179}(265275),\qquad
427095=\frac{9491}{21015}(945675).                         \tag{19.15}
\]
It therefore saves \(140850\) and \(518580\) Gaussian entries.
The largest retained input-plus-aggregate array payloads in this gate are
\[
4924800\ {\rm bytes}\quad\hbox{and}\quad17042400\ {\rm bytes}. \tag{19.16}
\]
\end{corollary}

\begin{proof}
A dense coefficient costs \(45^2=2025\) entries and a rank-\(r\) factor
pair costs \(45r+r45=90r\).  Sum the smaller cost over the distributions
(19.12)--(19.13) to obtain (19.15).  Equation (19.16) counts the two
signed-int64 component arrays retained for every input and candidate
frequency.  It excludes interpreter dictionaries and transient BLAS
scratch, and is therefore an exact coefficient-array payload, not a claim
about process resident memory.
\end{proof}

\begin{remark}[Meaning of the two cancellations]
\label{rem:v15-modular-zero}
Equation (19.11) is an exact statement in \(k\).  A zero at one prime does
not by itself prove that the corresponding rational coefficient vanishes.
That conclusion is reserved for the full CRT reconstruction.
\end{remark}

\subsection{Provenance and next acquisition}

\begin{record}[V15 verifier]
\label{rec:v15-verifier}
The executable
\begin{center}
\path{scripts/verify_ym_postarchive_v15_one_prime_densification_gate.py}
\end{center}
has \(450\) physical lines, \(16094\) bytes, and SHA--256
\[
\texttt{23570455B6805DA73A36C0E13440D8C8A682F4E646E00E6BB30419B850618F2E}.
                                                               \tag{19.17}
\]
The canonical payload SHA--256 is
\[
\texttt{7C819F5776C36BB7F3B521D15855F85C3E7DD59F4132FE3BBF60F16AC22F0AD7}.
                                                               \tag{19.18}
\]
The Wilson and endpoint aggregate-dictionary hashes are respectively
\[
\begin{split}
&\texttt{943785A1206DDAE81195FBAAEAE9EBCEE3B91893B62BBA3AD630C50AD540042E},\\
&\texttt{53DD1C12A572E08A39CA0A42987AD93DFCB7B497C3E207C21FB177090DE10FE6}.
\end{split}
                                                               \tag{19.19}
\]
\end{record}

\begin{decision}[V16: factor-aware \(E^3\) recurrence]
\label{dec:v15-v16}
The early-abort condition is not met: most aggregate coefficients remain
factor-profitable and the coefficient arrays remain modest.  V16 should
construct canonical modular factors of the matrices certified in
(19.12)--(19.13), multiply the factored \(E^2\) boundary by the actual
\(E\) bank, sum by frequency, and re-rank the completed \(E^3\)
coefficients.  It must compare factor-aware products with dense products
and repeat the phase identity.  If \(E^3\) loses material compression, the
programme returns upstream to a scalar/block contraction before forming
\(E^4\); otherwise it completes \(S=I+E+E^2+E^3+E^4\) at this prime.
\end{decision}

\begin{assessment}[V15 acquisition and boundary]
\label{ass:v15}
V15 completes the compulsory companion audit and the first exact modular
frequency aggregation.  It proves that the used aggregate ranks, not the
input ranks, still support hybrid storage through \(E^2\).  It does not yet
compute \(E^3,E^4\), either scalar term in (17.13), endpoint-minus-Wilson
subtraction, or any Haar integral.
\end{assessment}

\begin{record}[V15 append record]
\label{rec:v15-append}
V15 is appended to the validated V14 whose installed SHA--256 was
\[
\texttt{E8A71C2B09204402F33E1006F0CE0D11ABF6E4D4D4DD1EFE12CFD71FD228275E}.
                                                               \tag{19.20}
\]
After validation only V15 is to remain the active numeric YM note.  The
next compulsory companion audit is V20.
\end{record}

\begin{firewall}[Claim boundary after V15]
\label{fire:v15-final}
The modular \(E^2\) certificate is a finite algebraic sub-result.  The full
order-five signed response, its exact CRT lift, complement control,
continuous Haar-mean sign, derivative-containing sector, nonzero Gaussian
edge class, compensating scale boundary, SCIC, ADMC, interacting cofinal
state, regulator removal, Osterwalder--Schrader reconstruction, physical
clustering, and Hamiltonian Yang--Mills mass gap remain open.
\end{firewall}

% =============================================================================
% END APPENDED POST-TENTH-ARCHIVE V15 MATERIAL
% =============================================================================

% =============================================================================
% BEGIN APPENDED POST-TENTH-ARCHIVE V16 MATERIAL
% =============================================================================

\section{V16: canonical modular factors and the E-cubed aggregate}
\label{sec:v16}

V15 showed that low rank survives the first aggregate multiplication, but
it retained the \(E^2\) matrices densely during execution.  V16 constructs
and verifies canonical factors over the same finite field, uses those
factors in the next recurrence, and then discards all inherited rank
claims: every completed \(E^3\) coefficient is ranked afresh.

\subsection{Canonical factors over the CRT field}

Continue with \(k=\mathbb F_p[\iota]/(\iota^2+1)\) and \(p=2^{31}-1\).
For \(M\in M_{45}(k)\) of rank \(r>0\), exact elimination selects the
lexicographically first pivot-column set \(J\).  Put
\[
 U=M_{[:,J]}\in M_{45\times r}(k).                         \tag{20.1}
\]
Elimination of \(U^\top\) selects pivot rows \(I\), so
\(U_I\in M_r(k)\) is invertible.  Define
\[
 V=(U_I)^{-1}M_I\in M_{r\times45}(k).                     \tag{20.2}
\]

\begin{lemma}[Canonical factor reconstruction]
\label{lem:v16-factor}
The matrices (20.1)--(20.2) satisfy \(M=UV\);
\(\operatorname{rank}U=\operatorname{rank}V=r\).  The construction is
deterministic for the fixed row order and first-pivot convention.
\end{lemma}

\begin{proof}
The columns of \(U\) are pivot columns of \(M\), hence are a basis of the
column space of \(M\).  Every column \(m\) of \(M\) has a unique coordinate
vector \(c\) with \(m=Uc\).  Restriction to the rows \(I\) gives
\(m_I=U_Ic\), whence \(c=(U_I)^{-1}m_I\).  These coordinate vectors are
exactly the columns of \(V\), proving \(M=UV\).  The invertibility of
\(U_I\) proves both factor ranks.  First available pivots uniquely determine
\(J\) and \(I\).
\end{proof}

As before, the factor is stored only when
\[
 90r<2025,\quad\hbox{equivalently}\quad r\le22.             \tag{20.3}
\]
Every selected factor is multiplied back and compared entrywise with its
dense source before it may enter the recurrence.

\subsection{Factor-aware multiplication}

Let a boundary coefficient and an \(E\)-coefficient have factorizations
\[
 A=U_A V_A,\qquad B=U_B V_B,
 \qquad \operatorname{rank}A=r_A,\quad
 \operatorname{rank}B=r_B.                                \tag{20.4}
\]
When both factors are profitable, V16 evaluates
\[
 AB=U_A(V_AU_B)V_B,                                       \tag{20.5}
\]
associating the last two products so that the final dense expansion uses
\(\min(r_A,r_B)\).  If only \(A\) or \(B\) is factored, it uses respectively
\[
 AB=U_A(V_AB),\qquad AB=(AU_B)V_B.                         \tag{20.6}
\]

\begin{proposition}[Exact cost gate]
\label{prop:v16-cost}
Every route (20.5)--(20.6) equals the dense product in \(M_{45}(k)\).
Its charged number of Gaussian scalar multiply-accumulates is the sum of
\(mkn\) over the rectangular products.  The route is used only with
profitable factors; the total charge is compared with \(45^3\) for every
dense boundary product.
\end{proposition}

\begin{proof}
Equations (20.5)--(20.6) follow from Lemma~\ref{lem:v16-factor} and
associativity of matrix multiplication.  An \(m\times k\) by \(k\times n\)
classical product has \(mkn\) scalar multiply-accumulates.  Summing that
quantity over the displayed associations gives the charged cost.  The
executable records each route and verifies that the aggregate charge does
not exceed the dense baseline; it also performs periodic full dense
comparisons of the returned matrices.
\end{proof}

\subsection{The grouped cubic recurrence}

With the V15 coefficients \(G^{(2)}_{e,\rho}\), define
\[
 G^{(3)}_{e,\nu}
  =\sum_{\substack{\rho+\alpha=\nu\\
                   G^{(2)}_{e,\rho}\ne0}}
     G^{(2)}_{e,\rho}E_{e,\alpha}.                         \tag{20.7}
\]

\begin{lemma}[Grouped paths are complete]
\label{lem:v16-paths}
Over \(k\),
\[
 E_e(z)^3=\sum_\nu G^{(3)}_{e,\nu}z^\nu.                  \tag{20.8}
\]
Skipping a zero \(G^{(2)}_{e,\rho}\) loses no length-three path
contribution, even if it removes a frequency from the intermediate
envelope.
\end{lemma}

\begin{proof}
Multiply the exact identity (19.6) by (19.4) on the right and collect equal
frequencies.  A zero intermediate coefficient contributes the zero matrix
to every product, so deletion is exact in \(k\).  Expanding (20.7) recovers
the sum of all ordered triples \((\alpha,\beta,\gamma)\) with
\(\alpha+\beta+\gamma=\nu\).
\end{proof}

\subsection{Executed E-cubed certificate}

\begin{theorem}[Factor-aware one-prime \(E^3\) gate]
\label{thm:v16-executed}
Exact execution at \(p=2147483647\) gives
\[
\begin{array}{c|rr}
 &W&\partial\\ \hline
E\ \hbox{factor/dense coefficients}&20/1&48/9\\
E^2\ \hbox{factor/dense coefficients}&110/21&372/95\\
\hbox{grouped boundary products}&2751&26619\\
\hbox{represented ordered length-three paths}&9261&185193\\
\hbox{raw cubic frequency envelope}&471&1859\\
\hbox{frequencies actually accumulated}&471&1857\\
\hbox{nonzero completed coefficients}&471&1847\\
\hbox{whole-matrix modular cancellations}&0&10\\
E^3\ \hbox{factor/dense coefficients}&342/129&1176/671\\
\hbox{full-rank coefficients}&43&361.
\end{array}                                                  \tag{20.9}
\]
The endpoint zero frequencies are
\[
\begin{gathered}
(-3,-3,2,1),\quad(-3,-2,3,-1),\quad(-3,-2,3,1),\\
(-2,-3,3,-1),\quad(-2,-3,3,1),\\
(2,3,-3,-1),\quad(2,3,-3,1),\\
(3,2,-3,-1),\quad(3,2,-3,1),\quad(3,3,-2,-1).
\end{gathered}                                               \tag{20.10}
\]
The Wilson ranks range from \(2\) to \(45\), and the endpoint ranks from
\(1\) to \(45\).  The complete rank histograms are part of the canonical
payload.
\end{theorem}

\begin{proof}
The V15 square is reconstructed first.  Lemma~\ref{lem:v16-factor}
certifies every factor, Proposition~\ref{prop:v16-cost} certifies every
factor-aware multiplication route, and Lemma~\ref{lem:v16-paths} proves
that their frequency sums equal the cubic coefficients.  Exact finite-field
elimination is applied only after each sum is complete.  The executable
compares \(3\) Wilson and \(27\) endpoint routed products with independent
dense multiplication.  At the six V15 quarter-turn points it verifies
\[
 \sum_\nu G^{(3)}_{e,\nu}\iota^{\nu\cdot t}
 =\left(\sum_\alpha E_{e,\alpha}\iota^{\alpha\cdot t}\right)^3.
                                                               \tag{20.11}
\]
All comparisons pass.  Counting the completed dictionaries and their
fresh ranks gives (20.9)--(20.10).
\end{proof}

\begin{corollary}[Measured savings through \(E^3\)]
\label{cor:v16-savings}
The factor-aware Gaussian scalar charges are
\[
\begin{array}{c|rr}
 &W&\partial\\ \hline
\hbox{dense baseline}&250684875&2425656375\\
\hbox{factor-aware charge}&52129125&418354335\\
\hbox{saved charge}&198555750&2007302040\\
\hbox{exact charge ratio}
 &\dfrac{46337}{222831}&\dfrac{3098921}{17967825}.
\end{array}                                                  \tag{20.12}
\]
For the completed \(E^3\) dictionary the dense/hybrid entry counts are
\[
\begin{array}{c|rr}
 &W&\partial\\ \hline
\hbox{dense entries}&953775&3740175\\
\hbox{hybrid entries}&610965&2230875\\
\hbox{saved entries}&342810&1509300\\
\hbox{exact hybrid ratio}
 &\dfrac{13577}{21195}&\dfrac{3305}{5541}.
\end{array}                                                  \tag{20.13}
\]
The peak retained \(E,E^2,E^3\) coefficient-array payloads are
\(20185200\) and \(77209200\) bytes.
\end{corollary}

\begin{proof}
Sum the route costs from Proposition~\ref{prop:v16-cost} over all grouped
products to obtain (20.12).  Sum
\(\min(2025,90r)\) over the fresh \(E^3\) rank histograms to obtain
(20.13).  The byte count retains two signed-int64 arrays per dense
coefficient in the three recurrence levels and, as in V15, excludes
interpreter and BLAS workspace.
\end{proof}

\begin{remark}[Modular cancellation firewall]
\label{rem:v16-cancellation}
The ten zeros in (20.10) are exact only modulo the first CRT prime.  Their
frequency symmetry is a consistency check, not a rational-zero theorem.
They remain subject to the other fourteen prime runs.
\end{remark}

\subsection{Executable provenance and continuation}

\begin{record}[V16 verifier]
\label{rec:v16-verifier}
The executable
\begin{center}
\path{scripts/verify_ym_postarchive_v16_factor_aware_e3_gate.py}
\end{center}
has \(612\) physical lines, \(21304\) bytes, and SHA--256
\[
\texttt{E229D8B013FD61049DEA7E74AD06B165B91D2E91E2280682A4E92D23476EAEEA}.
                                                               \tag{20.14}
\]
Its canonical payload SHA--256 is
\[
\texttt{4C465D28A62A9DC70E2FC6F96224D5ADCAE905EAEAB269C0D60C2E5FD5AA6D73}.
                                                               \tag{20.15}
\]
The Wilson and endpoint \(E^3\) dictionary hashes are
\[
\begin{split}
&\texttt{F869FBC72310B73D59D7A352B904FE17320CD7217031E1921CF7E8EB9C46572A},\\
&\texttt{B46BED1C1365819D36E83A88B8CE5C20791D65ADF909B3939A4AC9BEF4196882}.
\end{split}                                                    \tag{20.16}
\]
\end{record}

\begin{decision}[V17: complete \(E^4\) and the one-prime Neumann sum]
\label{dec:v16-v17}
The cubic aggregate still saves about \(79\%\) and \(83\%\) of the dense
field-multiplication charge, and its hybrid storage saves \(342810\) and
\(1509300\) entries.  The upstream-retreat condition is therefore not met.
V17 should factor the completed \(E^3\) bank, compute \(E^4=E^3E\), and
form the coefficientwise sum
\[
 S=I+E+E^2+E^3+E^4
\]
at this prime.  It must separately rank the \(E^4\) and \(S\) aggregates,
verify both phase identities, and decide whether response multiplication
should use the \(S\) factors or a block/trace contraction.
\end{decision}

\begin{assessment}[V16 acquisition and boundary]
\label{ass:v16}
V16 proves and executes the first factor-aware aggregate recurrence.
It represents every length-three path, verifies the routed products, and
shows that material compression survives through \(E^3\).  It has not yet
formed \(E^4\), \(S\), the scalar response, any CRT lift, or a Haar sign.
\end{assessment}

\begin{record}[V16 append record]
\label{rec:v16-append}
V16 is appended to the validated V15 whose installed SHA--256 was
\[
\texttt{6DAEA47E6EBDE32C4C2D41AE6C98C531327EDC67377D5CBFDC851A8B02DF3D15}.
                                                               \tag{20.17}
\]
After validation only V16 is to remain the active numeric YM note.
\end{record}

\begin{firewall}[Claim boundary after V16]
\label{fire:v16-final}
The exact one-prime \(E^3\) certificate does not prove the order-five
response, rational coefficients, Haar-mean sign, complement bound,
interacting cofinal state, regulator removal, Osterwalder--Schrader
reconstruction, physical clustering, or Hamiltonian Yang--Mills mass gap.
\end{firewall}

% =============================================================================
% END APPENDED POST-TENTH-ARCHIVE V16 MATERIAL
% =============================================================================

% =============================================================================
% BEGIN APPENDED POST-TENTH-ARCHIVE V17 MATERIAL
% =============================================================================

\section{V17: the fourth aggregate power and the one-prime Neumann bank}
\label{sec:v17}

V16 left one residual multiplication before the order-five Neumann
polynomial could be used in the response.  V17 performs that multiplication
with the same post-aggregation factor rule, freezes an independent
\(E^4\) certificate, and then reuses its dense arrays to form
\[
 S=I+E+E^2+E^3+E^4                                      \tag{21.1}
\]
coefficientwise.  The \(E^4\) and \(S\) dictionaries are separately hashed,
ranked, and checked at finite-field phase points.

\subsection{Exact fourth-power recurrence}

For each action \(e\), let
\[
 G^{(4)}_{e,\nu}
 =\sum_{\substack{\rho+\alpha=\nu\\
                  G^{(3)}_{e,\rho}\ne0}}
     G^{(3)}_{e,\rho}E_{e,\alpha}.                         \tag{21.2}
\]

\begin{lemma}[Fourth-power completeness]
\label{lem:v17-e4}
In the Laurent matrix algebra over \(k\),
\[
 E_e(z)^4=\sum_\nu G^{(4)}_{e,\nu}z^\nu.                  \tag{21.3}
\]
The grouped recurrence represents all \(|A_e|^4\) ordered length-four
paths, including those that cancel in an intermediate frequency sum.
\end{lemma}

\begin{proof}
Multiply (20.8) by (19.4) on the right and collect equal frequencies.
Expanding (21.2) gives the ordered fourfold convolution.  If the complete
cubic coefficient at \(\rho\) is zero, the sum of every path arriving
there is already zero in \(k\); multiplying that sum by any
\(E_{e,\alpha}\) remains zero by distributivity.  Removing the zero
aggregate therefore does not remove a nonzero contribution.
\end{proof}

V17 constructs a canonical V16 representation of every nonzero cubic
coefficient and uses (20.5)--(20.6) in (21.2).  The completed fourth-power
coefficient is again ranked only after every product at its frequency has
been added.

\subsection{Freezing the fourth power before in-place summation}

Let \(\mathcal G_e^{(j)}\) denote the dense coefficient dictionary of
\(E_e^j\), with \(\mathcal G_e^{(0)}\) equal to the identity at zero
frequency.  After (21.2) is complete, V17 first computes:
\begin{enumerate}[leftmargin=2.2em]
\item the canonical byte digest of \(\mathcal G_e^{(4)}\);
\item its zero-frequency-set digest;
\item its complete rank histogram; and
\item its hybrid storage census.
\end{enumerate}
Only then does it mutate the same arrays by
\[
 \mathcal G_e^{(4)}
 \mathrel{+}=\mathcal G_e^{(3)}
             +\mathcal G_e^{(2)}
             +\mathcal G_e^{(1)}
             +\mathcal G_e^{(0)}.                          \tag{21.4}
\]

\begin{lemma}[Sound in-place Neumann construction]
\label{lem:v17-inplace}
After (21.4), the mutated dictionary is exactly the coefficient dictionary
of \(S_e\) in (21.1).  The previously frozen \(E^4\) digest and census
remain valid records of the pre-mutation object.
\end{lemma}

\begin{proof}
For every frequency, (21.4) is precisely the coefficient rule for the
finite sum (21.1); a missing key represents the zero matrix and is inserted
before addition.  The digest and census are finite values computed from the
completed \(E^4\) arrays before their first mutation.  Subsequent mutation
cannot retroactively change those recorded values.
\end{proof}

\subsection{Executed fourth-power and Neumann certificate}

\begin{theorem}[Complete one-prime \(E^4\) and \(S\) banks]
\label{thm:v17-executed}
At \(p=2147483647\), exact execution gives
\[
\begin{array}{c|rr}
 &W&\partial\\ \hline
\hbox{nonzero represented \(E^3\) coefficients}&471&1847\\
\hbox{factor/dense \(E^3\) representations}&342/129&1176/671\\
\hbox{grouped \(E^3E\) products}&9891&105279\\
\hbox{represented ordered length-four paths}&194481&10556001\\
E^4\ \hbox{candidate frequencies}&1251&5159\\
E^4\ \hbox{nonzero frequencies}&1251&5135\\
E^4\ \hbox{whole-matrix modular cancellations}&0&24\\
E^4\ \hbox{factor/dense coefficients}&716/535&2558/2577\\
E^4\ \hbox{full-rank coefficients}&237&1567.
\end{array}                                                  \tag{21.5}
\]
For the Neumann bank \(S\),
\[
\begin{array}{c|rr}
 &W&\partial\\ \hline
\hbox{candidate frequencies}&1251&5159\\
\hbox{nonzero frequencies}&1251&5135\\
\hbox{whole-matrix modular zeros}&0&24\\
\hbox{factor/dense coefficients}&716/535&2558/2577\\
\hbox{full-rank coefficients}&237&1567\\
\hbox{minimum nonzero rank}&2&1\\
\hbox{maximum rank}&45&45.
\end{array}                                                  \tag{21.6}
\]
Although the two rank censuses happen to coincide action by action, the
\(E^4\) and \(S\) coefficient hashes are different.
\end{theorem}

\begin{proof}
Lemma~\ref{lem:v17-e4} proves the recurrence object.  Every cubic factor is
reconstructed exactly by Lemma~\ref{lem:v16-factor}, and every routed
product is exact by Proposition~\ref{prop:v16-cost}.  V17 compares \(10\)
Wilson and \(106\) endpoint routed products with dense multiplication.
After complete frequency summation it performs field elimination on every
nonzero \(E^4\) coefficient, giving (21.5).

Lemma~\ref{lem:v17-inplace} proves the subsequent coefficientwise sum.
All \(S\) coefficients are independently re-ranked, giving (21.6).  At six
deterministic quarter-turn points, the executable verifies
\[
\begin{split}
 \operatorname{ev}_t(E^4)&=\operatorname{ev}_t(E)^4,\\
 \operatorname{ev}_t(S)&=
 I+\operatorname{ev}_t(E)+\operatorname{ev}_t(E)^2
   +\operatorname{ev}_t(E)^3+\operatorname{ev}_t(E)^4.
\end{split}
                                                               \tag{21.7}
\]
Because the \(E^4\) arrays are reused, the first equality is checked by
subtracting the four lower evaluated powers from the stored \(S\)
evaluation.  All twelve action--point pairs pass.
\end{proof}

\begin{corollary}[Fourth-power operation and storage savings]
\label{cor:v17-savings}
The fourth-power factor-aware charge is
\[
\begin{array}{c|rr}
 &W&\partial\\ \hline
\hbox{dense Gaussian baseline}&901317375&9593548875\\
\hbox{factor-aware Gaussian charge}&236265345&2096901675\\
\hbox{saved charge}&665052030&7496647200\\
\hbox{exact charge ratio}
 &\dfrac{5250341}{20029275}&\dfrac{345169}{1579185}.
\end{array}                                                  \tag{21.8}
\]
The \(E^4\), and separately the \(S\), storage census is
\[
\begin{array}{c|rr}
 &W&\partial\\ \hline
\hbox{dense Gaussian entries}&2533275&10398375\\
\hbox{hybrid Gaussian entries}&1916235&7372485\\
\hbox{saved entries}&617040&3025890\\
\hbox{exact hybrid ratio}
 &\dfrac{42583}{56295}&\dfrac{54611}{77025}.
\end{array}                                                  \tag{21.9}
\]
The peak retained coefficient-array payload through the in-place \(S\)
stage is \(60717600\) bytes for Wilson and \(244360800\) bytes for endpoint.
\end{corollary}

\begin{proof}
Sum the rectangular \(mkn\) charges of the recorded V16 product routes and
compare with \(45^3\) per grouped product to obtain (21.8).  Sum
\(\min(2025,90r)\) over the independently computed rank histogram to obtain
(21.9).  The peak count includes the retained dense arrays for
\(E,E^2,E^3\), and one shared \(E^4/S\) bank, but excludes factor copies,
interpreter objects, and BLAS scratch.
\end{proof}

\begin{record}[V17 immutable dictionary identities]
\label{rec:v17-digests}
The Wilson \(E^4\) and \(S\) SHA--256 values are
\[
\begin{split}
&\texttt{6037555BB23BAB27EFBBB119352DBD8B917F319350383DA0B04FA3662799AEB8},\\
&\texttt{E532AD9253B3E7BA25260152517D55EBD60DF0E2BF037EC01D1DA9964900179B}.
\end{split}                                                    \tag{21.10}
\]
The endpoint values are
\[
\begin{split}
&\texttt{E5CEAA6974438CD810742D366A9B29E454C7FA59E6420A0A44411740D823A9E6},\\
&\texttt{95C5BEA0F1BE0DF756EE53456A308DB725DEE9C776B31A141C970FE7789C284E}.
\end{split}                                                    \tag{21.11}
\]
The endpoint set of \(24\) modular-zero frequencies has SHA--256
\[
\texttt{94E92386EF80CB6F3B2080F7FAB9EF80F93853E09DE05EB07B88A9AC6C6BC80E}.
                                                               \tag{21.12}
\]
The same zero set remains zero in \(S\); these outer frequencies receive no
lower-power contribution.  As in V15--V16, this is not yet a rational-zero
claim.
\end{record}

\subsection{Executable provenance and next multiplication}

\begin{record}[V17 verifier]
\label{rec:v17-verifier}
The executable
\begin{center}
\path{scripts/verify_ym_postarchive_v17_e4_neumann_sum_gate.py}
\end{center}
has \(319\) physical lines, \(11336\) bytes, and SHA--256
\[
\texttt{3B6781FBC8E6FCA2E1CD5090036F4A5DFC5E76D0E2DC7157AF93438968F8F0FF}.
                                                               \tag{21.13}
\]
Its canonical payload SHA--256 is
\[
\texttt{A2CEAC92BB93C5EC57D08D3227CFAFBC565C643C9BA583D5CC2128ADC208099C}.
                                                               \tag{21.14}
\]
\end{record}

\begin{decision}[V18: one-prime response-bank multiplication]
\label{dec:v17-v18}
The matrix Neumann stage is now complete at one CRT prime.  V18 should
construct verified hybrid representations of the \(S\) coefficients and
of the actual \(Q_2,P_0,P_1,P_2\) banks, then compute
\[
 L_X=SX,\qquad X\in\{Q_2,P_0,P_1,P_2\},                   \tag{21.15}
\]
by frequency aggregation.  It should first measure one response bank for
each action, then continue through all four if the measured state remains
tractable.  Ranks must again be recomputed after aggregation.  No trace,
endpoint-minus-Wilson subtraction, reflection pairing, or absolute value
may precede the common scalar response dictionary.
\end{decision}

\begin{assessment}[V17 acquisition and boundary]
\label{ass:v17}
V17 completes and certifies \(E^4\) and the order-five Neumann matrix bank
at the first V14 prime.  The live finite bottleneck is no longer the
Neumann recurrence but the four \(S\)-times-response-bank convolutions and
their scalar trace contractions.
\end{assessment}

\begin{record}[V17 append record]
\label{rec:v17-append}
V17 is appended to the validated V16 whose installed SHA--256 was
\[
\texttt{93514BC316AEF4B842585A434729AFB26AB21AFF34A60A694A78EF36DDCCAB59}.
                                                               \tag{21.16}
\]
After validation only V17 is to remain the active numeric YM note.
\end{record}

\begin{firewall}[Claim boundary after V17]
\label{fire:v17-final}
The one-prime matrix Neumann bank is not yet a scalar response or a
rational certificate.  The remaining fourteen primes, exact CRT lift,
signed Haar analysis, complement and derivative sectors, cofinal
interacting construction, regulator removal, Osterwalder--Schrader
reconstruction, physical clustering, and Hamiltonian Yang--Mills mass gap
remain open.
\end{firewall}

% =============================================================================
% END APPENDED POST-TENTH-ARCHIVE V17 MATERIAL
% =============================================================================

% =============================================================================
% BEGIN APPENDED POST-TENTH-ARCHIVE V18 MATERIAL
% =============================================================================

\section{V18: the SP-zero response-bank pilot and an upstream return}
\label{sec:v18}

V17 completed the matrix Neumann bank at the first CRT prime.  Before
forming all four response banks, V18 tests the smallest bubble bank,
\[
 L_{e,P_0}=S_eP_{e,0}.                                    \tag{22.1}
\]
The test is exact and successful, but its state census changes the next
algorithmic direction: endpoint frequency aggregation is now predominantly
dense.  V18 therefore accepts (22.1) as a certificate and rejects
simultaneous materialization of all four \(L_X\) banks.

\subsection{Actual response coefficients}

The \(P_0\) bank is compiled with the same frozen action transport as V14,
then reduced in \(k=\mathbb F_{p^2}\).  Write
\[
 P_{e,0}(z)=\sum_{\beta\in B_{e,0}}P_{e,0,\beta}z^\beta.
                                                               \tag{22.2}
\]
For the already certified V17 bank
\[
 S_e(z)=\sum_{\alpha\in A_{e,S}}S_{e,\alpha}z^\alpha,
\]
define
\[
 L_{e,P_0,\nu}
 =\sum_{\substack{\alpha+\beta=\nu\\
                  S_{e,\alpha}\ne0}}
   S_{e,\alpha}P_{e,0,\beta}.                              \tag{22.3}
\]

\begin{lemma}[Response-bank convolution]
\label{lem:v18-convolution}
Equation (22.3) is the exact coefficient dictionary of (22.1).  If
\[
 \Pi_e=\sum_{j=0}^4|A_e|^j,
\]
then it represents exactly \(\Pi_e|B_{e,0}|\) unaggregated
Neumann--response paths.
\end{lemma}

\begin{proof}
Distribute the two finite Laurent sums and use
\(z^\alpha z^\beta=z^{\alpha+\beta}\).  Each coefficient of \(S\) is
itself the sum of all residual paths of lengths zero through four, by
V15--V17.  Distributing once more yields the stated path count.  Skipping a
zero completed \(S\)-coefficient is exact by distributivity.
\end{proof}

V18 reconstructs \(S\) rather than trusting a mutable in-memory object.
Before (22.3) it requires the reconstructed hashes to equal the frozen V17
values (21.10)--(21.11).  It then constructs canonical factors with the
V16 pivot rule for \(S\) and \(P_0\), executes the factor-aware product
routes, sums by frequency, and re-ranks only the completed \(L_{P_0}\)
coefficients.

\subsection{Executed response-bank census}

\begin{theorem}[Exact one-prime \(L_{P_0}\) certificate]
\label{thm:v18-executed}
At \(p=2147483647\), exact execution gives
\[
\begin{array}{c|rr}
 &W&\partial\\ \hline
S\ \hbox{nonzero factor/dense representations}&716/535&2558/2577\\
P_0\ \hbox{factor/dense representations}&9/0&45/6\\
\hbox{grouped \(SP_0\) products}&11259&261885\\
\hbox{represented raw paths}&1837845&547969551\\
L_{P_0}\ \hbox{candidate frequencies}&2499&14981\\
L_{P_0}\ \hbox{nonzero frequencies}&2499&14951\\
\hbox{whole-matrix modular cancellations}&0&30\\
L_{P_0}\ \hbox{factor/dense coefficients}&1493/1006&6624/8327\\
\hbox{full-rank coefficients}&0&4987\\
\hbox{minimum/maximum rank}&2/38&1/45.
\end{array}                                                  \tag{22.4}
\]
The canonical Wilson and endpoint dictionaries have SHA--256
\[
\begin{split}
&\texttt{10833A2B7A224FDE6E63FBA480C99D234724FA25BBB07CC300083AAA183BE754},\\
&\texttt{EE51C467A5768EEBB1A36DA812266C6FFE00C363323E442EC6779B86EB1AAE80}.
\end{split}                                                    \tag{22.5}
\]
\end{theorem}

\begin{proof}
The frozen \(S\) digest comparison passes for both actions.
Lemma~\ref{lem:v18-convolution} identifies the computed dictionary.
The V16 factor reconstruction is checked entrywise for every accepted
factor.  Twelve Wilson and \(263\) endpoint routed products are compared
with dense multiplication.  At each of six deterministic quarter-turn
points the executable verifies
\[
 \sum_\nu L_{e,P_0,\nu}\iota^{\nu\cdot t}
 =\left(\sum_\alpha S_{e,\alpha}\iota^{\alpha\cdot t}\right)
  \left(\sum_\beta P_{e,0,\beta}\iota^{\beta\cdot t}\right).
                                                               \tag{22.6}
\]
It finally performs exact field elimination on the completed frequency
sums.  The resulting counts give (22.4), while canonical serialization
gives (22.5).
\end{proof}

\begin{corollary}[Measured cost and state growth]
\label{cor:v18-cost}
The factor-aware and storage censuses are
\[
\begin{array}{c|rr}
 &W&\partial\\ \hline
\hbox{dense multiplication baseline}&1025976375&23864270625\\
\hbox{factor-aware charge}&253289700&7874176860\\
\hbox{saved charge}&772686675&15990093765\\
\hbox{exact charge ratio}
 &\dfrac{375244}{1519965}&\dfrac{2160268}{6547125}\\ \hline
\hbox{dense output entries}&5060475&30275775\\
\hbox{hybrid output entries}&3312360&22495185\\
\hbox{saved output entries}&1748115&7780590\\
\hbox{exact hybrid ratio}
 &\dfrac{24536}{37485}&\dfrac{166631}{224265}.
\end{array}                                                  \tag{22.7}
\]
The retained \(S,P_0,L_{P_0}\) dense-array payloads peak at
\[
121791600\ {\rm bytes}\quad\hbox{and}\quad654188400\ {\rm bytes}.
                                                               \tag{22.8}
\]
\end{corollary}

\begin{proof}
Sum the exact rectangular route charges as in V16 and compare with
\(45^3\) per grouped product.  Sum \(\min(2025,90r)\) over the fresh output
rank histograms for the storage counts.  Equation (22.8) counts two
signed-int64 component arrays for each retained dense coefficient of the
three banks; factor copies, dictionary overhead, and BLAS scratch are not
included.
\end{proof}

\subsection{Why the next step moves upstream}

\begin{proposition}[The four-bank materialization plan fails its gate]
\label{prop:v18-return}
The V17 plan to materialize \(L_{Q_2},L_{P_0},L_{P_1},L_{P_2}\)
simultaneously is rejected for the endpoint action.
\end{proposition}

\begin{proof}
The smallest bubble response bank already has \(14951\) nonzero matrix
coefficients.  Of these, \(8327\) fail the factor-storage inequality and
\(4987\) are full rank.  Its measured retained-array payload is
\(654188400\) bytes before interpreter, factor, and workspace overhead.
Thus the premise that the response boundary remains predominantly
factor-compressible is false.  Materializing three additional banks would
compound this measured state without yet producing the scalar quantity in
(17.13).  The certificate remains feasible, but the simultaneous
matrix-bank representation fails the declared state gate.
\end{proof}

\begin{remark}[This is a representation rejection, not a theorem of
impossibility]
\label{rem:v18-not-impossible}
V18 does not prove that the other three matrix banks cannot be computed.
It proves that the chosen reason for computing them together---continued
low-rank compression---has ceased to hold on the actual endpoint output.
The programme therefore follows the user's upstream rule and changes the
contraction order before spending that state.
\end{remark}

\subsection{Executable provenance and upstream acquisition}

\begin{record}[V18 verifier]
\label{rec:v18-verifier}
The executable
\begin{center}
\path{scripts/verify_ym_postarchive_v18_p0_response_bank_gate.py}
\end{center}
has \(326\) physical lines, \(11508\) bytes, and SHA--256
\[
\texttt{137E901CF8A2EE5B5811EDE836F24EAB1C346208E68965DAEFC540898BB91243}.
                                                               \tag{22.9}
\]
Its canonical payload SHA--256 is
\[
\texttt{37D6171F83A2B2821EDDD4E083EFE3F34F7999C8FC743EE46D16E645B6B49C0C}.
                                                               \tag{22.10}
\]
\end{record}

\begin{decision}[V19: trace-first frequency tiling]
\label{dec:v18-v19}
Return upstream to the scalar objects actually required:
\[
\begin{split}
 T_{20,\nu}
  &=\sum_\mu\operatorname{Tr}
       (L_{P_2,\mu}L_{P_0,\nu-\mu}),\\
 T_{11,\nu}
  &=\sum_\mu\operatorname{Tr}
       (L_{P_1,\mu}L_{P_1,\nu-\mu}),\\
 T_{Q,\nu}&=\operatorname{Tr}(L_{Q_2,\nu}).
\end{split}
                                                               \tag{22.11}
\]
V19 should compute the exact coordinate spans and admissible transform
lengths at the V14 prime, compare sparse pair contraction with a
memory-bounded multidimensional number-theoretic transform, and prove a
tiling/streaming theorem that never retains all four \(L_X\) banks.  The
signed endpoint-minus-Wilson scalar dictionary remains the first object on
which reflection pairing or absolute values may act.
\end{decision}

\begin{assessment}[V18 acquisition and boundary]
\label{ass:v18}
V18 computes the first complete response bank at one CRT prime and locates
the densification boundary with actual ranks and state.  It does not
compute the two scalar bubbles, the tadpole, the other fourteen primes, or
any Haar sign.
\end{assessment}

\begin{record}[V18 append record]
\label{rec:v18-append}
V18 is appended to the validated V17 whose installed SHA--256 was
\[
\texttt{C0D6AC6B4F4F527220A10F2590DB9254B792C0539B1ACAB104C3271CF4983260}.
                                                               \tag{22.12}
\]
After validation only V18 is to remain the active numeric YM note.
\end{record}

\begin{firewall}[Claim boundary after V18]
\label{fire:v18-final}
The one-prime \(L_{P_0}\) bank is not the signed scalar response.
Trace contraction, exact CRT lifting, Haar-mean sign, complement and
derivative sectors, cofinal interacting construction, regulator removal,
Osterwalder--Schrader reconstruction, physical clustering, and the
Hamiltonian Yang--Mills mass gap remain open.
\end{firewall}

% =============================================================================
% END APPENDED POST-TENTH-ARCHIVE V18 MATERIAL
% =============================================================================

% =============================================================================
% BEGIN APPENDED POST-TENTH-ARCHIVE V19 MATERIAL
% =============================================================================

\section{V19: alias-free trace-first transform geometry}
\label{sec:v19}

V18 showed that the matrix bank \(L_{P_0}\) is already the wrong retained
object for the endpoint action.  V19 moves one step upstream.  It computes
the exact coordinate spans of the three scalar terms, proves that the first
CRT field contains roots of the required orders, and derives a pointwise
spectral formula that never materializes any \(L_X\) bank.  This closes the
transform-geometry question; execution cost inside each spectral point
remains to be measured.

\subsection{Exact scalar support geometry}

Let \(A_S,A_{Q_2},A_{P_0},A_{P_1},A_{P_2}\subset\mathbb Z^4\) be the
candidate supports already fixed in V13.  Put
\[
\begin{split}
 A_{L_X}&=A_S+A_X,\\
 A_{20}&=A_{L_{P_2}}+A_{L_{P_0}},\\
 A_{11}&=A_{L_{P_1}}+A_{L_{P_1}},\\
 A_{\rm sc}&=A_{L_{Q_2}}\cup A_{20}\cup A_{11}.
\end{split}
                                                               \tag{23.1}
\]
The executable does not form the two large sums in (23.1) literally.
Since \(0\in A_E\) and \(A_S=\bigcup_{j=0}^4jA_E\),
\[
 A_S+A_S=\bigcup_{j=0}^8jA_E,                             \tag{23.2}
\]
and associativity gives
\[
\begin{split}
 A_{20}&=(A_S+A_S)+(A_{P_2}+A_{P_0}),\\
 A_{11}&=(A_S+A_S)+(A_{P_1}+A_{P_1}).
\end{split}
                                                               \tag{23.3}
\]
Thus (23.3) is an exact reduction of enumeration work, not a support
majorant.

\begin{theorem}[Executed coordinate-span census]
\label{thm:v19-spans}
The Wilson supports have
\[
\begin{array}{c|r|c|c}
\hbox{object}&|A|&[\min,\max]^4&\hbox{widths}\\ \hline
S&1251&[-4,4]^4&(9,9,9,9)\\
L_{P_0}&2499&[-5,5]^4&(11,11,11,11)\\
L_{P_1}=L_{Q_2}&2621&[-5,5]^4&(11,11,11,11)\\
L_{P_2}&2721&[-5,5]^4&(11,11,11,11)\\
A_{20}&34635&[-10,10]^4&(21,21,21,21)\\
A_{11}&34347&[-10,10]^4&(21,21,21,21)\\
A_{\rm sc}&35069&[-10,10]^4&(21,21,21,21).
\end{array}                                                  \tag{23.4}
\]
The endpoint supports have
\[
\begin{array}{c|r|c|c}
\hbox{object}&|A|&(\min;\max)&\hbox{widths}\\ \hline
S&5171&((-4)^4;(4)^4)&(9,9,9,9)\\
L_{P_0}&15075&((-6,-6,-5,-6);(6,6,5,6))&(13,13,11,13)\\
L_{P_1}&17991&((-6)^4;(6)^4)&(13,13,13,13)\\
L_{P_2}&17771&((-6)^4;(6)^4)&(13,13,13,13)\\
L_{Q_2}&18141&((-6)^4;(6)^4)&(13,13,13,13)\\
A_{20}&238795&
((-12,-12,-11,-12);(12,12,11,12))&(25,25,23,25)\\
A_{11}&260967&((-12)^4;(12)^4)&(25,25,25,25)\\
A_{\rm sc}&261797&((-12)^4;(12)^4)&(25,25,25,25).
\end{array}                                                  \tag{23.5}
\]
The signed endpoint-minus-Wilson candidate union equals the endpoint
\(A_{\rm sc}\), is invariant under \(\nu\mapsto-\nu\), and has \(130899\)
reflection orbits.
\end{theorem}

\begin{proof}
The raw supports are selected by the same nonzero coefficient predicates
as V12--V14.  Repeated integer-tuple Minkowski sums give \(A_S\) and
(23.2).  Equations (23.3) then give the two bubble supports.  Direct minima,
maxima, set cardinalities, and reflection membership tests give
(23.4)--(23.5).  The endpoint scalar-union SHA--256 is
\[
\texttt{F6C17DCD47F124FD4EDEBF8CC92C3ED23329B4D5D6EBC39953667A114E9088C0}.
                                                               \tag{23.6}
\]
This matches the V13 cardinality and orbit count while adding the new
coordinate-span certificate.
\end{proof}

\subsection{Roots in the first CRT field}

The first modulus satisfies
\[
 p-1=2147483646
 =2\cdot3^2\cdot7\cdot11\cdot31\cdot151\cdot331.           \tag{23.7}
\]
All seven displayed bases and \(p\) are prime.  The factorization has
\(192\) positive divisors.

\begin{theorem}[Primitive generator and required roots]
\label{thm:v19-roots}
The residue \(g=7\) generates \(\mathbb F_p^\times\).  For each required
length \(N\), let \(\omega_N=g^{(p-1)/N}\).  The exact roots are
\[
\begin{array}{c|r}
N&\omega_N\\ \hline
9&309107220\\
11&1969212174\\
14&610312904\\
21&1295308677\\
31&512.
\end{array}                                                  \tag{23.8}
\]
Each \(\omega_N\) has order exactly \(N\), and \(N^{-1}\) exists modulo
\(p\).
\end{theorem}

\begin{proof}
For every prime \(q\) in (23.7), direct modular exponentiation gives
\[
 7^{(p-1)/q}\ne1\pmod p.
\]
The standard cyclic-group order criterion therefore makes \(7\) primitive.
For each row of (23.8), the executable checks
\[
 \omega_N^N=1,\qquad
 \omega_N^{N/q}\ne1
 \quad\hbox{for every prime }q\mid N.                      \tag{23.9}
\]
Thus the order is \(N\).  Since \(N<p\), it is invertible in
\(\mathbb F_p\subset k\).
\end{proof}

\subsection{Smallest alias-free tensor grids}

For a finite support with coordinate widths \(w_j\), shift its minimum to
zero.  A tensor cyclic transform of lengths \(N_j\) recovers the linear
coefficient array whenever \(N_j\ge w_j\) for every coordinate.

\begin{lemma}[No-wrap recovery]
\label{lem:v19-no-wrap}
Let a Laurent polynomial over \(k\) be supported in
\(\prod_j[a_j,b_j]\), and let \(w_j=b_j-a_j+1\).  If
\(N_j\ge w_j\) and \(N_j\mid p-1\), then its values on the tensor roots
\((\omega_{N_1}^{t_1},\ldots,\omega_{N_4}^{t_4})\) determine every
coefficient uniquely by the inverse tensor DFT.
\end{lemma}

\begin{proof}
After multiplication by \(z^{-a}\), every exponent lies in
\(\prod_j[0,N_j-1]\); reduction modulo \(N_j\) is therefore injective on
each coordinate.  Orthogonality
\[
 \frac1{N_j}\sum_{t=0}^{N_j-1}
 \omega_{N_j}^{t(r-s)}=\delta_{rs}
\]
holds because \(\omega_{N_j}\) has exact order \(N_j\).  Applying it
successively in four coordinates gives the inverse formula and no cyclic
aliasing.
\end{proof}

\begin{theorem}[Minimal base-field grids]
\label{thm:v19-grids}
Among transform lengths dividing \(p-1\), the smallest componentwise
choices are
\[
\begin{array}{c|c|r|r}
\hbox{action/object}&(N_1,N_2,N_3,N_4)&\prod_jN_j&
\hbox{bytes for one packed \(k\)-scalar grid}\\ \hline
W:S&(9,9,9,9)&6561&52488\\
W:L_X&(11,11,11,11)&14641&117128\\
W:\hbox{scalar bubbles}&(21,21,21,21)&194481&1555848\\
\partial:S&(9,9,9,9)&6561&52488\\
\partial:L_{P_0}&(14,14,11,14)&30184&241472\\
\partial:L_{P_1},L_{P_2},L_{Q_2}&(14,14,14,14)&38416&307328\\
\partial:\hbox{scalar bubbles}&(31,31,31,31)&923521&7388168.
\end{array}                                                  \tag{23.10}
\]
The last row is also the common signed grid.
\end{theorem}

\begin{proof}
For each width in (23.4)--(23.5), enumerate all \(192\) divisors of
\(p-1\) and select the first divisor not below it.  This yields
(23.10); Theorem~\ref{thm:v19-roots} supplies every root and
Lemma~\ref{lem:v19-no-wrap} supplies recovery.  Two unsigned 32-bit
residues store one element of \(k\), giving eight times the grid cardinality
in the last column.
\end{proof}

\subsection{Trace-first spectral identity}

Let \(\mathcal F\) be the common tensor DFT, applied entrywise to matrix
Laurent dictionaries.  Evaluation is an algebra homomorphism, so
\[
 \mathcal F(AB)(t)=\mathcal F(A)(t)\mathcal F(B)(t).        \tag{23.11}
\]

\begin{theorem}[No-\(L_X\) pointwise contraction]
\label{thm:v19-trace-first}
The scalar spectral arrays required by (22.11) are
\[
\begin{split}
\widehat T_{20}(t)
 &=\operatorname{Tr}\bigl(
   \widehat S(t)\widehat P_2(t)
   \widehat S(t)\widehat P_0(t)\bigr),\\
\widehat T_{11}(t)
 &=\operatorname{Tr}\bigl(
   \widehat S(t)\widehat P_1(t)
   \widehat S(t)\widehat P_1(t)\bigr),\\
\widehat T_Q(t)
 &=\operatorname{Tr}\bigl(\widehat S(t)\widehat Q_2(t)\bigr).
\end{split}
                                                               \tag{23.12}
\]
Inverse transformation on the grid (23.10) returns the exact scalar
coefficient dictionaries.  At each spectral point,
\[
\widehat S(t)=I+\widehat E(t)+\widehat E(t)^2
                +\widehat E(t)^3+\widehat E(t)^4,           \tag{23.13}
\]
so no matrix \(L_X\) coefficient bank need be retained.
\end{theorem}

\begin{proof}
Apply (23.11) to \(L_X=SX\), then apply it once more to the two scalar
convolutions in (22.11).  Trace commutes with entrywise evaluation, giving
(23.12).  Apply (23.11) to each power of \(E\) to obtain (23.13).
Lemma~\ref{lem:v19-no-wrap} then recovers the exact Laurent coefficients
from the scalar arrays.
\end{proof}

\begin{corollary}[State reduction at the scalar boundary]
\label{cor:v19-state}
The endpoint sparse pair route exposes
\[
 |A_{L_{P_2}}||A_{L_{P_0}}|=267897825,\qquad
 |A_{L_{P_1}}|^2=323676081                              \tag{23.14}
\]
candidate matrix pairs.  The trace-first route instead has \(923521\)
independent spectral points and retains only scalar spectral outputs,
\(7388168\) bytes per packed array, plus a fixed number of temporary
\(45\times45\) matrices per point.
\end{corollary}

\begin{proof}
The products in (23.14) use the exact cardinalities in (23.5).
Equation (23.12) is pointwise, so after its three scalar values have been
written the temporary matrices at that point can be discarded.  The scalar
array size is the last row of (23.10).  This is a peak-state theorem, not
yet a runtime theorem: producing the matrix values
\(\widehat E,\widehat P_j,\widehat Q_2\) efficiently remains the next
executable question.
\end{proof}

\subsection{Signed order and streaming semantics}

Both actions use the same \(31^4\) grid.  Their pointwise scalar responses
may therefore be multiplied by their rational action constants in \(k\)
and subtracted on the grid.  Linearity of the inverse DFT makes this
identical to coefficientwise endpoint-minus-Wilson subtraction.  Only
after inverse transformation may the resulting coefficient dictionary be
paired under \(\nu\mapsto-\nu\) or bounded in absolute value.

\begin{firewall}[A transform architecture is not its execution]
\label{fire:v19-runtime}
The scalar grid is small, but each of its \(923521\) points still requires
the transformed \(45\)-dimensional base symbols and noncommutative products
in (23.12)--(23.13).  V19 proves exactness and memory feasibility of the
output representation; it does not claim that a naive per-point matrix
evaluation is fast.  The next verifier must measure and exploit the base
bank structure rather than hide this cost.
\end{firewall}

\subsection{Executable provenance and V20 checkpoint}

\begin{record}[V19 verifier]
\label{rec:v19-verifier}
The executable
\begin{center}
\path{scripts/verify_ym_postarchive_v19_trace_ntt_geometry.py}
\end{center}
has \(315\) physical lines, \(10532\) bytes, and SHA--256
\[
\texttt{BE3469200EDE8E1BA6D9100036F9092532733609838F17592FD64A585BB512D3}.
                                                               \tag{23.15}
\]
Its canonical payload SHA--256 is
\[
\texttt{F6BCBB606A580411EA0203B277A4F911324CAF9EF93AE6906FC5FE3D64F2C6E2}.
                                                               \tag{23.16}
\]
The verifier certifies (23.7), primitive generator \(7\), every exact root
order, every support digest and span, minimal divisor selection, reflection
closure, and equality of the signed and endpoint scalar unions.
\end{record}

\begin{decision}[V20: compulsory audit, then spectral producer gate]
\label{dec:v19-v20}
V20 must first inspect the then-current SM and NS notes.  It should import
only a genuinely applicable technique, if one exists.  It must then test
how to produce the matrices in (23.12) on the \(31^4\) grid: compare
direct sparse phase evaluation, separable transforms in bounded entry
batches, and any certified block/factor evaluation inherited from V13.
The output of the pilot is a measured operation/state gate and an exact
pointwise comparison with direct polynomial evaluation.  The signed scalar
array, rather than four \(L_X\) banks, remains the target.
\end{decision}

\begin{assessment}[V19 acquisition and boundary]
\label{ass:v19}
V19 turns the V18 upstream return into an exact alias-free transform
theorem.  It reduces the retained endpoint output to scalar grids and
certifies all roots and spans.  It has not executed the \(31^4\) matrix
producer, inverse transform, CRT lift, or Haar sign theorem.
\end{assessment}

\begin{record}[V19 append record]
\label{rec:v19-append}
V19 is appended to the validated V18 whose installed SHA--256 was
\[
\texttt{8D0EB367F8C5798EC4F650D2B6E30349722ED8E3F0541976E2EBA17D86EF61DA}.
                                                               \tag{23.17}
\]
After validation only V19 is to remain the active numeric YM note.  The
next compulsory companion audit is V20.
\end{record}

\begin{firewall}[Claim boundary after V19]
\label{fire:v19-final}
The exact transform geometry is a finite algebraic sub-result.  The scalar
response residues, their CRT reconstruction, continuous Haar-mean sign,
complement and derivative sectors, cofinal interacting construction,
regulator removal, Osterwalder--Schrader reconstruction, physical
clustering, and Hamiltonian Yang--Mills mass gap remain open.
\end{firewall}

% =============================================================================
% END APPENDED POST-TENTH-ARCHIVE V19 MATERIAL
% =============================================================================

% =============================================================================
% BEGIN APPENDED POST-TENTH-ARCHIVE V20 MATERIAL
% =============================================================================

\section{V20: companion audit, coefficientwise producer gate, and raw push-through}
\label{sec:v20}

V19 reduced the desired response certificate to one scalar array on an
alias-free \(31^4\) grid, but left open the cost of producing the scalar
spectral values.  V20 first performs the compulsory fifth-version audit of
the contemporaneous Standard Model and Navier--Stokes programmes.  It then
tests, on every coefficient of the actual frozen Yang--Mills banks, whether
the \(45\)-dimensional spectral producer has a simultaneous coordinate-block
or common-hull reduction.  The answer is negative.  This rules out an
attractive but inapplicable shortcut and forces the next step upstream.  The
final theorem pushes the trace through the fixed maps \(A\) and \(Z\),
exposing the raw physical-sector matrices and the single bridge \(K=ZA\).

\subsection{The compulsory V20 companion audit}

\begin{record}[Exact files inspected at the V20 audit]
\label{rec:v20-audit-files}
The audit used the then-current files
\begin{center}
\path{latex_notes/Renewal_SM_Einstein_Continuum_Research_Notes_V86.tex},\\
\path{latex_notes/Renewal_NS_Stability_Research_Notes_V26.tex}.
\end{center}
The SM snapshot has \(26197\) physical lines, \(932700\) bytes, and
SHA--256
\[
\texttt{0ABBD654D47D2EABFE1DBB3F67B2531F32704701D81392D5E89EA8C602B288DB}.
                                                               \tag{24.1}
\]
The NS snapshot has \(5319\) physical lines, \(278283\) bytes, and
SHA--256
\[
\texttt{82C887F8C2C69E4F28FAD7C3DC8DE5B9099CB5A4BF431EBA1FFF3E91935978AD}.
                                                               \tag{24.2}
\]
Thus the audit is reproducible even if either companion programme later
advances.
\end{record}

\begin{assessment}[What transfers from the companion programmes]
\label{ass:v20-audit-transfer}
The following two proof disciplines transfer.
\begin{enumerate}
\item The SM note constructs exact invariant orientation blocks before
      estimating transformed output, and aggregates all modal contributions
      in the physical coordinate basis before applying a norm or sign
      decision.  Accordingly, the present producer gate tests simultaneous
      structure on the full coefficient family and keeps the trace sum
      intact.
\item The NS note permits a rank-one or reduced-coordinate shortcut only
      after it has been proved on the used class; its computed gain may be
      zero.  Accordingly, isolated low ranks at sample points below are
      diagnostics only.  Uniform reduction is inferred solely from the
      coefficient hull.
\end{enumerate}
The SM product--Hodge cancellation and positive massive Green
representation do not transfer: the Yang--Mills response banks here are
complex, noncommutative matrices, and no corresponding invariant
orientation decomposition or positivity theorem has been proved.  Moreover,
the audited SM conclusion is strip-uniform rather than a full
four-dimensional cofinal theorem.  The NS field theorem likewise is not
imported.  The audit changes the order of the Yang--Mills tests, not their
mathematical conclusion.
\end{assessment}

\subsection{Coefficientwise block and hull criteria}

Let
\[
 {\cal X}(z)=\sum_{\alpha\in F}X_\alpha z^\alpha,
 \qquad X_\alpha\in M_d(\mathbb K),                         \tag{24.3}
\]
where \(F\subset\mathbb Z^4\) is finite and \(\mathbb K\) is any field.
Define the undirected coefficient graph \(G_{\cal X}\) on
\(\{1,\ldots,d\}\) by joining distinct \(i,j\) when
\((X_\alpha)_{ij}\ne0\) or \((X_\alpha)_{ji}\ne0\) for at least one
\(\alpha\).  Define also
\[
 U_{\cal X}:=\sum_{\alpha\in F}\operatorname{im}X_\alpha,
 \qquad
 W_{\cal X}:=\sum_{\alpha\in F}\operatorname{im}X_\alpha^{\mathsf T}.
                                                               \tag{24.4}
\]

\begin{theorem}[Exact simultaneous coordinate blocks]
\label{thm:v20-block-graph}
The connected components of \(G_{\cal X}\) give a coordinate partition
for which every \(X_\alpha\), and hence every \({\cal X}(z)\), is block
diagonal.  This is the finest such coordinate partition: every coordinate
partition for which all \(X_\alpha\) are block diagonal is a regrouping of
those connected components.
\end{theorem}

\begin{proof}
If \(i\) and \(j\) lie in different connected components, no coefficient
has a nonzero \(ij\) or \(ji\) entry; hence all entries crossing two
components vanish.  Conversely, if a coordinate partition makes every
coefficient block diagonal, the endpoints of every edge must lie in one
block.  The endpoints of every path, hence every connected component, must
therefore lie in one block.  Linearity proves the assertion for
\({\cal X}(z)\).
\end{proof}

\begin{theorem}[Coefficient hulls are uniform pointwise rank certificates]
\label{thm:v20-hull}
For every \(z\) in every extension field of \(\mathbb K\),
\[
 \operatorname{im}{\cal X}(z)\subseteq U_{\cal X},\qquad
 \operatorname{im}{\cal X}(z)^{\mathsf T}\subseteq W_{\cal X},
                                                               \tag{24.5}
\]
and therefore
\[
 \operatorname{rank}{\cal X}(z)
 \leq \min\{\dim U_{\cal X},\dim W_{\cal X}\}.                \tag{24.6}
\]
If \(\dim U_{\cal X}=r<d\) and \(B\in M_{d\times r}(\mathbb K)\)
has columns forming a basis of \(U_{\cal X}\), then there are unique
\(Y_\alpha\in M_{r\times d}(\mathbb K)\) with
\[
 X_\alpha=BY_\alpha,\qquad
 {\cal X}(z)=B\sum_\alpha Y_\alpha z^\alpha.                  \tag{24.7}
\]
There is an analogous common-row factorization when
\(\dim W_{\cal X}<d\).
\end{theorem}

\begin{proof}
Each column of a linear combination of the \(X_\alpha\) belongs to the sum
of their column spaces; applying this observation to the transposes proves
(24.5).  The rank bound follows.  Since every column of \(X_\alpha\) lies
in the column space of the full-rank matrix \(B\), its coordinates in that
basis exist and are unique, giving \(Y_\alpha\).  Summing the coefficient
identities gives (24.7).  Transposition proves the row statement.
\end{proof}

\begin{remark}[The exact profitability threshold used here]
\label{rem:v20-profit}
Replacing one dense \(d\times d\) application by the two rectangular
applications associated with a rank-\(r\) common factor has the elementary
entry-operation proxy \(2dr\), versus \(d^2\).  The gate therefore calls
the factor profitable only if
\[
                 2dr<d^2.                                   \tag{24.8}
\]
For \(d=45\), this requires \(r\leq22\).  This deliberately conservative
gate ignores possible hardware-dependent constants; failure of (24.8)
does not say that the factor is algebraically false, only that it does not
solve the present producer bottleneck.
\end{remark}

\subsection{The actual frozen-bank computation}

The verifier reconstructs the rational chart with maximum denominator
\(4\), reduces the exact coefficients in
\[
 \mathbb F_{p^2}=\mathbb F_p[\iota]/(\iota^2+1),
 \qquad p=2147483647,                                        \tag{24.9}
\]
and uses the V19 order-\(31\) root \(\omega=512\).  The two reconstruction
errors are
\[
 \epsilon_{\max}=2.3173136731078448\cdot10^{-14},\qquad
 \epsilon_{\rm F}=2.293018683569624\cdot10^{-13}.
\]
The eight diagnostic grid exponents are
\[
\begin{split}
 &(0,0,0,0),(1,0,0,0),(0,1,2,3),(1,1,1,1),\\
 &(1,2,3,5),(3,5,7,11),(7,11,13,17),(19,23,29,30).
\end{split}                                                   \tag{24.10}
\]
All uniform assertions below are computed from all coefficients, not from
these eight points.

\begin{theorem}[Negative producer gate for the endpoint action]
\label{thm:v20-endpoint-gate}
For each endpoint family
\[
              E,\quad P_0,\quad P_1,\quad P_2,\quad Q_2,      \tag{24.11}
\]
the exact coefficient graph over (24.9) is connected on all \(45\)
coordinates.  Their combined coefficient graph is also connected.  For
each family,
\[
       \dim U_{\cal X}=\dim W_{\cal X}=45.                    \tag{24.12}
\]
Consequently neither a simultaneous coordinate block nor a proper common
row/column factor can reduce the endpoint \(45\)-dimensional producer.
The frequency counts and diagnostic ranks are
\[
\begin{array}{c|c|c|c|c}
{\cal X}&|F|&\#\text{ active positions}&
(\dim U_{\cal X},\dim W_{\cal X})&
\text{ranks at the points in (24.10)}\\ \hline
E   &57&2025&(45,45)&42,45,45,45,45,45,45,45\\
P_0 &51&1983&(45,45)&42,44,45,45,45,45,45,45\\
P_1 &81&2025&(45,45)&43,43,45,45,45,45,45,45\\
P_2 &81&2025&(45,45)&45,45,45,45,45,45,45,45\\
Q_2 &99&2025&(45,45)&45,45,45,45,45,45,45,45
\end{array}                                                    \tag{24.13}
\]
The combined family occupies all \(2025\) matrix positions, and its sole
component has size \(45\).
\end{theorem}

\begin{proof}
For every nonzero coefficient entry the verifier inserts the corresponding
edge and computes graph components.  For each family and for their union
it obtains one component containing indices \(0,\ldots,44\), with no
off-component nonzero entry.  Theorem~\ref{thm:v20-block-graph} gives the
block claim.

For the column hull it horizontally concatenates every coefficient matrix
and performs exact Gaussian elimination over \(\mathbb F_{p^2}\); for the
row hull it vertically stacks them and performs the same calculation.
Every reported rank is \(45\).  Thus Theorem~\ref{thm:v20-hull} gives
(24.12) and excludes every proper factor of the two certified types.  The
last column of (24.13) is obtained by exact root evaluation followed by
exact Gaussian elimination and is only a reproducibility check.
\end{proof}

\begin{theorem}[Wilson action gate]
\label{thm:v20-wilson-gate}
Each Wilson family and their combined family again have one coordinate
component of size \(45\).  The exact hull and sample-rank census is
\[
\begin{array}{c|c|c|c|c}
{\cal X}&|F|&(\dim U_{\cal X},\dim W_{\cal X})&
\text{sample ranks}&\text{factor profitable}\\ \hline
E   &21&(45,45)&33,38,45,45,45,45,45,45&\text{no}\\
P_0 & 9&(38,38)&34,34,34,34,34,34,34,34&\text{no}\\
P_1 &11&(42,42)&36,37,38,40,40,40,40,40&\text{no}\\
P_2 &15&(45,45)&44,44,45,45,45,45,45,45&\text{no}\\
Q_2 &11&(45,45)&45,45,45,45,45,45,45,45&\text{no}
\end{array}                                                    \tag{24.14}
\]
The constructed common-column factors for \(P_0\) and \(P_1\) reconstruct
every coefficient exactly and have respective SHA--256 digests
\[
\begin{split}
&\texttt{F3D73AB24629A5F4E9D5C95CD8EB135285C9E678A46AE9C56185DE7561A587B5},\\
&\texttt{730E7F2715D2E5653F7EBC0D9EF53D8AD1A9A409A5BFA5511E311D4EE64B99C0}.
\end{split}                                                    \tag{24.15}
\]
Their ranks \(38\) and \(42\) fail (24.8), so neither is admitted as the
spectral-producer architecture.
\end{theorem}

\begin{proof}
The graph and hull computations are exactly those in the proof of
Theorem~\ref{thm:v20-endpoint-gate}.  For each proper column hull the
verifier selects pivot columns of the concatenated bank to form \(B\),
selects an invertible pivot-row minor, solves for every \(Y_\alpha\), and
checks \(BY_\alpha=X_\alpha\) entrywise.  Since
\[
 2\cdot45\cdot38>45^2,\qquad
 2\cdot45\cdot42>45^2,                                      \tag{24.16}
\]
the final assertion follows from the declared gate (24.8).
\end{proof}

\begin{firewall}[What the rank samples do not prove]
\label{fire:v20-samples}
The ranks in (24.13)--(24.14) neither bound all \(31^4\) points nor prove
a characteristic-zero generic-rank statement.  Only the coefficientwise
graph and hull computations support uniform claims.  Conversely, reduction
modulo one prime can lose a characteristic-zero nonzero minor, but that
does not weaken the negative result used here: a determinant minor nonzero
modulo \(p\) is the reduction of a nonzero rational minor.  Hence full hull
rank modulo \(p\) certifies full rational hull rank.
\end{firewall}

\subsection{The upstream trace push-through}

Let \(A,Z\in M_d(\mathbb K)\) be the fixed reconstruction and coordinate
maps of the frozen chart.  Let \(C(z)\) be the raw Hessian bank and let
\(R_j(z)\) denote any raw response bank.  Put
\[
\begin{split}
 E(z)&=I-A C(z)Z,\qquad K:=ZA,\\
 \widetilde E(z)&=I-C(z)K,\qquad
 {\cal S}(X):=I+X+X^2+X^3+X^4.
\end{split}                                                    \tag{24.17}
\]

\begin{theorem}[Exact push-through identities]
\label{thm:v20-push-through}
For every spectral point \(z\),
\[
             {\cal S}(E(z))A=A{\cal S}(\widetilde E(z)).       \tag{24.18}
\]
Consequently the linear trace term satisfies
\[
 \operatorname{Tr}\!\left({\cal S}(E)A R_2 Z\right)
 =\operatorname{Tr}\!\left(K{\cal S}(\widetilde E)R_2\right), \tag{24.19}
\]
and the two quadratic trace terms satisfy
\[
\begin{split}
 &\operatorname{Tr}\!\left(
 {\cal S}(E)A R_2 Z\,{\cal S}(E)A R_0 Z\right)\\
 &\hspace{23mm}=
 \operatorname{Tr}\!\left(
 K{\cal S}(\widetilde E)R_2
 K{\cal S}(\widetilde E)R_0\right).
\end{split}                                                    \tag{24.20}
\]
\[
\begin{split}
 &\operatorname{Tr}\!\left(
 {\cal S}(E)A R_1 Z\,{\cal S}(E)A R_1 Z\right)\\
 &\hspace{23mm}=
 \operatorname{Tr}\!\left(
 K{\cal S}(\widetilde E)R_1
 K{\cal S}(\widetilde E)R_1\right).
\end{split}                                                    \tag{24.21}
\]
Thus the scalar spectral producer can be evaluated without forming a
coefficient or pointwise matrix of the dense image \(A(\,\cdot\,)Z\).
\end{theorem}

\begin{proof}
The elementary identity
\[
             (I-ACZ)A=A(I-CZA)=A(I-CK)                       \tag{24.22}
\]
gives \(EA=A\widetilde E\).  Induction gives
\[
                  E^nA=A\widetilde E^n\qquad(n\geq0),         \tag{24.23}
\]
and summation for \(0\leq n\leq4\) proves (24.18).  Substitute
(24.18) in the left side of (24.19) and cyclically move \(Z\) to the
front:
\[
 \operatorname{Tr}(A{\cal S}(\widetilde E)R_2Z)
 =\operatorname{Tr}(ZA{\cal S}(\widetilde E)R_2),
\]
which is (24.19).  After substituting (24.18) twice in the first
quadratic trace, cyclically move the final \(Z\) to the front:
\[
\begin{split}
 &\operatorname{Tr}\!\left(
 A{\cal S}(\widetilde E)R_2ZA
 {\cal S}(\widetilde E)R_0Z\right)\\
 &\quad=
 \operatorname{Tr}\!\left(
 ZA{\cal S}(\widetilde E)R_2ZA
 {\cal S}(\widetilde E)R_0\right).
\end{split}
\]
This is (24.20) because \(ZA=K\).  The same argument with \(R_1\)
in both slots proves (24.21).
\end{proof}

\begin{remark}[Why this is the correct upstream move]
\label{rem:v20-upstream}
V13 already introduced the bridge \(K=ZA\) coefficientwise.  The new
content here is the pointwise polynomial identity (24.18) and its use
inside the V19 trace-first producer.  It replaces five dense frozen banks
by the raw \(C,R_0,R_1,R_2\) banks and one fixed bridge.  Whether this is
computationally decisive depends on their simultaneous raw-sector
structure; V20 does not assume the answer.
\end{remark}

\subsection{Executable provenance and V21 checkpoint}

\begin{record}[V20 verifier]
\label{rec:v20-verifier}
The executable
\begin{center}
\path{scripts/verify_ym_postarchive_v20_shared_hull_producer_gate.py}
\end{center}
has \(350\) physical lines, \(12040\) bytes, and SHA--256
\[
\texttt{97680D08EEA4DB1640B333C208F7605BF063F150CC201AD78EC7E3A54ADF8949}.
                                                               \tag{24.24}
\]
Its canonical payload SHA--256 is
\[
\texttt{5A76D7B7363B9CC5579D2698C14EE2F5899DBD85858B46C9E185D727F4FDD189}.
                                                               \tag{24.25}
\]
It certifies all coefficient graphs, all hull ranks, both proper factor
reconstructions, the sample rank bounds, reuse of the V19 order-\(31\)
root, and absence of off-block coefficients.
\end{record}

\begin{decision}[V21: audit the raw bridge architecture]
\label{dec:v20-v21}
V21 should compile the raw coefficient families \(C,R_0,R_1,R_2\) and
the exact bridge \(K=ZA\), verify (24.18)--(24.21) pointwise against the
frozen formulas, and compute:
\begin{enumerate}
\item the common coordinate blocks and hulls of every raw family;
\item the way \(K\) couples those blocks;
\item an exact arithmetic-cost comparison with the dense frozen producer.
\end{enumerate}
If \(K\) couples every raw sector and destroys the saving, the programme
must move further upstream to a tensor, Schur-complement, or representation
structure actually present in \(C\) and \(K\); it must not fall back to an
unmeasured naive \(31^4\) dense run.
\end{decision}

\begin{assessment}[V20 acquisition and boundary]
\label{ass:v20}
V20 eliminates simultaneous frozen-bank blocks and common hulls as the
endpoint producer mechanism and proves an exact raw-sector reformulation.
It does not yet establish that the raw reformulation is cheaper, nor does it
execute the \(31^4\) scalar transform.
\end{assessment}

\begin{record}[V20 append record]
\label{rec:v20-append}
V20 is appended to the validated V19 whose installed SHA--256 was
\[
\texttt{1425DAD95958093EA027B6D48AC0AD57D9BBFCEAE6041CEDF7129FB43B83F5D6}.
                                                               \tag{24.26}
\]
After validation only V20 is to remain the active numeric YM note.  The
next compulsory companion audit is V25.  At V100 the current lineage is to
be archived with suffix \path{_fX} and the active notes restarted at V1
with context and a rigorous summary, as required by the standing rollover
rule.
\end{record}

\begin{firewall}[Claim boundary after V20]
\label{fire:v20-final}
Neither the companion audit, the negative structural gate, nor the
push-through identities prove a response sign or a Yang--Mills mass gap.
The endpoint scalar response residues, inverse transform and CRT
reconstruction, continuous Haar-mean sign, complement and derivative
sectors, cofinal interacting construction, regulator removal,
Osterwalder--Schrader reconstruction, physical clustering, and Hamiltonian
mass gap remain open.
\end{firewall}

% =============================================================================
% END APPENDED POST-TENTH-ARCHIVE V20 MATERIAL
% =============================================================================

% =============================================================================
% BEGIN APPENDED POST-TENTH-ARCHIVE V21 MATERIAL
% =============================================================================

\section{V21: exact raw-bridge architecture and complete operation gate}
\label{sec:v21}

V20 proved that the frozen \(45\times45\) coefficient banks have neither a
simultaneous block decomposition nor a profitable common hull.  It also
proved that their scalar traces can be pushed through the fixed chart maps
to the raw Hessian and response banks with bridge \(K=ZA\).  V21 now audits
that raw architecture coefficientwise, verifies its implementation against
the frozen formulas at exact spectral points, and counts the full declared
phase-production and matrix workload.  The raw architecture wins this
gate, although its absolute endpoint cost remains too large to execute
naively.

\subsection{Raw banks and exact implementation checks}

For each action write
\[
\begin{split}
 C(z)&=\sum_\alpha C_\alpha z^\alpha,\\
 Q_2(z)&=\sum_\alpha Q_{2,\alpha}z^\alpha,\qquad
 P_j(z)=\sum_\alpha P_{j,\alpha}z^\alpha\quad(0\leq j\leq2).
\end{split}                                                    \tag{25.1}
\]
These are the physical-sector coefficients before applying \(A\) and
\(Z\).  The factor \(i\) belonging to the first external derivative is
included in \(P_1\).  Over the first CRT field the verifier constructs
\[
 E=I-ACZ,\qquad \widetilde E=I-CK,\qquad K=ZA,                \tag{25.2}
\]
directly from the same rational chart used in V10--V20.

\begin{proposition}[Exact spectral implementation check]
\label{prop:v21-pointwise}
At each of the eight order-\(31\) spectral exponents in (24.10), for both
the Wilson and endpoint actions, the verifier certifies entrywise over
\(\mathbb F_{p^2}\) that
\[
\begin{split}
 EA&=A\widetilde E,\\
 {\cal S}(E)A&=A{\cal S}(\widetilde E),\\
 A R Z&=R^{\rm frozen}
       \quad\text{for }R\in\{Q_2,P_0,P_1,P_2\},
\end{split}                                                    \tag{25.3}
\]
and that the two sides of each trace identity
(24.19)--(24.21) agree.  The Wilson and endpoint trace transcripts have
respective SHA--256 digests
\[
\begin{split}
 &\texttt{CFCE1E5CF707D5215053264DF53E4855B12A49447465DCB7F9DD1B89F76081F2},\\
 &\texttt{C4708A1577F8ACC6B9381A207390C6B790E39F4B105143B1E2223E59D35AA43F}.
\end{split}                                                    \tag{25.4}
\]
\end{proposition}

\begin{proof}
For a Laurent bank \(X\) and exponent \(k\), the executable evaluates
\[
 X(\omega^k)=
 \sum_\alpha X_\alpha\,
 \omega^{\langle\alpha,k\rangle},\qquad \omega=512,           \tag{25.5}
\]
using all coefficients modulo \(p=2147483647\).  It independently compiles
the frozen banks by the V14 route and the raw banks by (25.1).  It then
performs exact pair-array matrix multiplication in
\(\mathbb F_p[\iota]/(\iota^2+1)\), compares every entry in (25.3), and
compares both components of all three traces.  Any mismatch raises an
exception before the transcript is emitted.  All \(16\) action--point
tests pass.  This proposition checks the executable realization; the
all-\(z\) algebraic proof remains Theorem~\ref{thm:v20-push-through}.
\end{proof}

\subsection{Coefficientwise raw structure and the bridge}

\begin{theorem}[Raw coefficient census]
\label{thm:v21-raw-census}
For the Wilson action the exact raw census is
\[
\begin{array}{c|r|r|r|c|c}
\text{bank}&|F|&\text{coefficient nonzeros}&
\text{active positions}&\text{component sizes}&
(\dim U,\dim W)\\ \hline
C   &21&673&561&45&(45,45)\\
P_0 & 9& 92& 92&39,1,1,1,1,1,1&(38,38)\\
P_1 &11&220&188&42,1,1,1&(42,42)\\
P_2 &15&444&374&45&(45,45)\\
Q_2 &11&264&232&45&(45,45)
\end{array}                                                    \tag{25.6}
\]
For the endpoint action it is
\[
\begin{array}{c|r|r|r|c|c}
\text{bank}&|F|&\text{coefficient nonzeros}&
\text{active positions}&\text{component sizes}&
(\dim U,\dim W)\\ \hline
C   &57&4067&1523&45&(45,45)\\
P_0 &51&1912& 664&45&(45,45)\\
P_1 &81&4132&1037&45&(45,45)\\
P_2 &81&6425&1374&45&(45,45)\\
Q_2 &99&6773&1407&45&(45,45)
\end{array}                                                    \tag{25.7}
\]
For each action the combined raw coefficient graph has one component of
size \(45\).  Its active-position count is \(561\) for Wilson and \(1523\)
for endpoint.  Hence the endpoint raw family has no simultaneous coordinate
block or proper common row/column hull.
\end{theorem}

\begin{proof}
The graph and hull definitions are those of (24.3)--(24.4).  The verifier
visits every exact nonzero coefficient, constructs the union pattern,
computes connected components, and performs exact Gaussian elimination on
the horizontally concatenated and vertically stacked coefficient matrices.
Theorems~\ref{thm:v20-block-graph} and \ref{thm:v20-hull} convert the
reported graph and hull data into the final assertion.
\end{proof}

\begin{proposition}[The chart bridge is almost dense and invertible]
\label{prop:v21-bridge}
For each action, \(K=ZA\) has rank \(45\) modulo \(p\) and exactly \(1937\)
nonzero entries out of \(2025\).  Its Wilson and endpoint SHA--256 digests
are respectively
\[
\begin{split}
 &\texttt{BB4403D1655B45AAD6AC2954185F264C06E220CFE84D65242CFE565BCBE6F0D4},\\
 &\texttt{CCD7FB8FCE5764F02AB4DB0DE42C9BF410FE3F112C87BFC0C3642961C1B4A760}.
\end{split}                                                    \tag{25.8}
\]
Thus sparsity of the raw coefficients, rather than a bridge block or bridge
rank defect, is the only reduction certified here.
\end{proposition}

\begin{proof}
The verifier reduces the exact rational representation product \(ZA\) and
independently multiplies the reduced \(Z\) and \(A\), checking equality.
It counts the nonzero pair entries and obtains rank \(45\) by exact
elimination.  Since the combined raw graph already has one component, no
nontrivial raw block survives to be protected by \(K\).
\end{proof}

\subsection{A complete declared-operation comparison}

We count multiply-accumulates in \(\mathbb F_{p^2}\).  This gives a
machine-independent ledger for two explicit algorithms, not a wall-clock
model.  Put \(d=45\), \(N=31\), and let \(M\) be the
number of nonzero coefficient entries across the required input banks.
Direct sparse evaluation of every grid point uses
\[
                        N^4M                              \tag{25.9}
\]
phase multiply-accumulates.  If \(q\) matrix positions carry at least one
coefficient, a separable direct \(N\)-point DFT on all four axes uses
\[
                        4N^5q.                            \tag{25.10}
\]
The gate takes the lesser of (25.9) and (25.10) separately for each
architecture.

\begin{lemma}[Pointwise matrix-kernel counts]
\label{lem:v21-kernel-count}
The frozen trace-first producer has scheduled count
\[
             c_{\rm fr}=6d^3+3d^2=552825.                    \tag{25.11}
\]
Let \(s_X\) be the number of active matrix positions in a raw bank \(X\).
The raw producer has the uniform scheduled upper count
\[
\begin{split}
 c_{\rm raw}={}&4d^3+
 d(s_C+s_{P_0}+s_{P_1}+s_{P_2})\\
 &\quad+2d^2+s_{Q_2}.                                     
\end{split}                                                   \tag{25.12}
\]
It equals \(423457\) for Wilson and \(576867\) for endpoint.
\end{lemma}

\begin{proof}
For the frozen producer, three dense products form
\[
                         E^2,\quad E^3,\quad E^4.
\]
Three more form
\[
 {\cal S}(E)P_0,\qquad{\cal S}(E)P_1,\qquad{\cal S}(E)P_2.
\]
The linear trace and two bubble traces each use \(d^2\) scalar
contractions, proving (25.11).

For the raw producer, multiplication of sparse \(C\) by dense \(K\) uses
at most \(d s_C\) multiply-accumulates.  Three dense products form the
powers in \({\cal S}(\widetilde E)\), and one forms
\(K{\cal S}(\widetilde E)\).  Right multiplication of this dense matrix by
the three sparse \(P_j\) uses at most
\(d(s_{P_0}+s_{P_1}+s_{P_2})\).  The two bubble contractions use
\(2d^2\), while the linear trace contracts only the \(s_{Q_2}\) possible
nonzeros.  This proves (25.12).  Substitution of (25.6)--(25.7) gives the
two numerical values.
\end{proof}

\begin{theorem}[The complete raw architecture wins the V21 gate]
\label{thm:v21-cost-gate}
The phase inputs are
\[
\begin{array}{c|r|r|r|r}
&M_{\rm fr}&q_{\rm fr}&M_{\rm raw}&q_{\rm raw}\\ \hline
\text{Wilson}&80452&9672&1693&1447\\
\text{endpoint}&517823&10083&23309&6005
\end{array}                                                    \tag{25.13}
\]
For every entry in (25.13), direct sparse evaluation (25.9) is cheaper
than the separable direct DFT (25.10).  The resulting complete ledgers are
\[
\begin{array}{c|r|r}
&\text{Wilson}&\text{endpoint}\\ \hline
\text{frozen phase}&74299111492&478220414783\\
\text{raw phase}&1563521053&21526350989\\
\text{frozen matrix}&510545496825&510545496825\\
\text{raw matrix}&391071432097&532748788707\\ \hline
\text{frozen total}&584844608317&988765911608\\
\text{raw total}&392634953150&554275139696
\end{array}                                                    \tag{25.14}
\]
Therefore the raw architecture has strictly lower declared complete cost
for both actions, even though its endpoint matrix kernel alone is more
expensive.
\end{theorem}

\begin{proof}
The values \(M\) are exact counts of all nonzero coefficient entries; the
values \(q\) are exact counts of matrix positions active in each bank,
summed over the five banks.  Substituting \(N=31\) in (25.9)--(25.10)
selects (25.9) in all four cases.  Multiplying the two values in
Lemma~\ref{lem:v21-kernel-count} by
\[
                         N^4=923521                           \tag{25.15}
\]
gives the matrix row of (25.14).  Adding the selected phase rows gives the
totals, whose strict inequalities prove the claim.
\end{proof}

\begin{firewall}[Meaning and limitation of the operation ledger]
\label{fire:v21-cost}
The ledger counts a specified uniform schedule; it does not assert that
every possible algorithm requires this many operations, and it does not
translate a Gaussian-field operation into a hardware-independent running
time.  It omits lower-order additions used to assemble
\({\cal S}(\widetilde E)\) in both architectures.  More importantly,
\(5.54\cdot10^{11}\) endpoint operations is still not a practical
certificate run.  V21 selects the correct side of the V20 fork but does
not authorize the full grid execution.
\end{firewall}

\subsection{Polynomial form of the next upstream target}

\begin{lemma}[Alternating bridge polynomial]
\label{lem:v21-bridge-polynomial}
For every square matrix \(X\),
\[
 {\cal S}(I-X)=5I-10X+10X^2-5X^3+X^4.                       \tag{25.16}
\]
In the raw producer \(X=CK\).  Hence every scalar response term is a
finite linear combination of cyclic traces of alternating words in
\(C,K,Q_2,P_0,P_1,P_2\), of \(C\)-degree at most eight.
\end{lemma}

\begin{proof}
Expand \((I-X)^n\) for \(0\leq n\leq4\); \(I\) commutes with \(X\), so the
ordinary binomial theorem applies.  Summing the coefficients of each
power of \(X\) gives \(5,-10,10,-5,1\).  Substitution in (24.19)--(24.21)
gives the word statement.  Each of the two factors in a bubble contributes
at most four copies of \(C\).
\end{proof}

\subsection{Executable provenance and V22 checkpoint}

\begin{record}[V21 verifier]
\label{rec:v21-verifier}
The executable
\begin{center}
\path{scripts/verify_ym_postarchive_v21_raw_bridge_gate.py}
\end{center}
has \(555\) physical lines, \(18085\) bytes, and SHA--256
\[
\texttt{EF7718C8C7DE0C928BE4D373FAAD3CC645C8FA6C29F67480CA2598B986AE8ACD}.
                                                               \tag{25.17}
\]
Its canonical payload SHA--256 is
\[
\texttt{1AB68E3A1FC394A355823E5AAE85217F1CB93C5A8580844546C54E15FF090807}.
                                                               \tag{25.18}
\]
\end{record}

\begin{decision}[V22: find structure inside the alternating words]
\label{dec:v21-v22}
The raw producer is now selected, but its endpoint count remains too large.
V22 should therefore work upstream on the word family in
Lemma~\ref{lem:v21-bridge-polynomial}.  Before any full-grid run it must:
\begin{enumerate}
\item test exact reductions under a common change of basis, including the
      algebra generated by the matrices; coordinate blocks alone are
      insufficient;
\item test tensor, displacement, or Schur structure of the exact bridge
      \(K\) in the natural physical-sector labeling;
\item if those gates are negative, construct a measured batched kernel on
      a small exact grid tile and derive a defensible whole-grid
      time/state bound.
\end{enumerate}
This is the upstream response to the remaining \(5.54\cdot10^{11}\)
operation count.  V22 must not expand the coefficient paths literally:
V13 already proved that the corresponding collapsed path stream exceeds
one trillion before cancellations.
\end{decision}

\begin{assessment}[V21 acquisition and boundary]
\label{ass:v21}
V21 proves the raw and frozen spectral formulas agree in the exact
implementation, rules out a raw coordinate-block explanation, and selects
the raw architecture by a complete declared-operation gate.  It does not
yet make the endpoint computation practical or produce one scalar residue
array.
\end{assessment}

\begin{record}[V21 append record]
\label{rec:v21-append}
V21 is appended to the clean V20 whose installed SHA--256 was
\[
\texttt{86E816946E39A997F78EE5FB84F02890935B4964C779E844172D18D2106C327E}.
                                                               \tag{25.19}
\]
After validation only V21 is to remain the active numeric YM note.  The
next compulsory companion audit remains V25.
\end{record}

\begin{firewall}[Claim boundary after V21]
\label{fire:v21-final}
No finite-field residue array or characteristic-zero response sign has
been proved.  The Haar inequality, interacting cofinal construction,
regulator removal, OS reconstruction, physical clustering, and Hamiltonian
Yang--Mills mass gap also remain open.  The exact finite producer programme
is only a local algebraic component of that larger chain.
\end{firewall}

% =============================================================================
% END APPENDED POST-TENTH-ARCHIVE V21 MATERIAL
% =============================================================================

% =============================================================================
% BEGIN APPENDED POST-TENTH-ARCHIVE V22 MATERIAL
% =============================================================================

\section{V22: full generated-algebra obstruction to invariant sectors}
\label{sec:v22}

V21 showed that the raw endpoint producer is cheaper than the frozen one,
but still requires more than \(5.5\cdot10^{11}\) operations in its declared
schedule.  Coordinate connectivity alone does not exclude a block
decomposition after a non-coordinate change of basis.  V22 closes that
loophole.  A family of sixteen rank-one endpoint Hessian coefficients,
together with the bridge \(K\), generates the entire matrix algebra
\(M_{45}\).  Therefore no common invariant subspace, block triangularization,
or exact simultaneous similarity reduction exists for the endpoint family.
The bridge also has maximal Kronecker rearrangement rank in every standard
nontrivial factorization of \(45\).

\subsection{A rank-one family criterion}

Let \(\mathbb K\) be a field, let \(K\in M_d(\mathbb K)\), and let
\[
                   D_s=u_sv_s^{\mathsf T}\qquad(s\in S)       \tag{26.1}
\]
be nonzero rank-one matrices.  Define the two Krylov spans
\[
\begin{split}
 {\cal U}&=\operatorname{span}\{K^a u_s:s\in S,\ 0\leq a<d\},\\
 {\cal V}&=\operatorname{span}\{(K^{\mathsf T})^b v_s:
                         s\in S,\ 0\leq b<d\}.
\end{split}                                                    \tag{26.2}
\]
Give \(S\) the directed interaction graph in which \(s\to t\) when
\[
             v_s^{\mathsf T}K^m u_t\ne0
             \quad\hbox{for some }0\leq m<d.                 \tag{26.3}
\]

\begin{theorem}[Krylov--interaction full-algebra criterion]
\label{thm:v22-family-criterion}
If
\[
             {\cal U}=\mathbb K^d,\qquad
             {\cal V}=\mathbb K^d,                           \tag{26.4}
\]
and the graph (26.3) is strongly connected, then the unital algebra
generated by \(K\) and the \(D_s\) is all of \(M_d(\mathbb K)\).
\end{theorem}

\begin{proof}
An edge \(s\to t\), witnessed by \(m\), gives
\[
 D_sK^mD_t
 =u_s(v_s^{\mathsf T}K^m u_t)v_t^{\mathsf T},                \tag{26.5}
\]
a nonzero scalar multiple of \(u_sv_t^{\mathsf T}\).  Multiplying identities
of the form (26.5) along a directed path shows that the generated algebra
contains \(u_sv_t^{\mathsf T}\) for every ordered pair \(s,t\): strong
connectivity supplies a path, and the product of its nonzero scalar
couplings is nonzero.  Left and right multiplication by powers of \(K\)
then supplies
\[
 K^a u_s v_t^{\mathsf T}K^b
 =(K^a u_s)\bigl((K^{\mathsf T})^b v_t\bigr)^{\mathsf T}.
                                                               \tag{26.6}
\]
By (26.4), the first factors in (26.6) span \(\mathbb K^d\) and the
second factors span its dual.  Their outer products span
\(\mathbb K^d\otimes(\mathbb K^d)^*=M_d(\mathbb K)\).
\end{proof}

\begin{remark}[Why the family criterion is needed]
\label{rem:v22-single-failure}
The endpoint raw banks contain \(40\) rank-one coefficients:
\[
                    24\text{ in }C,\qquad
                     8\text{ in }P_1,\qquad
                     8\text{ in }P_2.                        \tag{26.7}
\]
None is individually cyclic on both sides.  The exact pairs
\((\dim\operatorname{Krylov}_K(u_s),
\dim\operatorname{Krylov}_{K^{\mathsf T}}(v_s))\), with multiplicity, are
\[
\begin{array}{c|rrrrrrrrr}
(r_\ell,r_r)&(23,23)&(23,24)&(24,23)&(24,24)&(24,26)&
(24,27)&(26,24)&(27,24)&(27,27)\\ \hline
\#&4&1&3&6&2&6&2&4&12 .
\end{array}                                                    \tag{26.8}
\]
Thus the tempting one-coefficient cyclic proof is false.  Passing to the
actual family is the required upstream correction.
\end{remark}

\subsection{The endpoint certificate}

Order all rank-one endpoint coefficients first by bank name and then by
frequency.  The deterministic greedy witness selects the canonical indices
\[
\begin{split}
 2,14,3,13,0,4,5,7,\;&9,11,12,20,\\
                    &10,16,17,23 .
\end{split}                                                    \tag{26.9}
\]
All sixteen are coefficients of \(C\).  The executable records their
frequencies, factor pivots, individual Krylov ranks, and vector digests.

\begin{theorem}[The endpoint raw family generates \(M_{45}\)]
\label{thm:v22-full-algebra}
For the sixteen rank-one matrices in (26.9), over
\(\mathbb F_{p^2}\),
\[
                  \dim{\cal U}=\dim{\cal V}=45.               \tag{26.10}
\]
Every one of the \(16^2=256\) directed pairs has a witness (26.3), so its
interaction graph is complete and hence strongly connected.  Therefore
\[
 \operatorname{alg}_{\mathbb F_{p^2}}
     (K,D_s:s\text{ in (26.9)})
          =M_{45}(\mathbb F_{p^2}).                           \tag{26.11}
\]
The family certificate has SHA--256
\[
\texttt{03625747E892D6CAE03B35FB1F88E85751DB845333AE6C2E038031B0016D8889}.
                                                               \tag{26.12}
\]
\end{theorem}

\begin{proof}
Each rank-one factorization is reconstructed entrywise.  The verifier forms
the \(45\) Krylov columns on each side, concatenates them for the selected
family, and performs exact Gaussian elimination over
\(\mathbb F_{p^2}\), giving (26.10).  For each ordered pair it tests
(26.3) successively for \(0\leq m<45\), recording the first nonzero
witness.  It obtains all \(256\) edges.  Theorem
\ref{thm:v22-family-criterion} gives (26.11).
\end{proof}

\begin{corollary}[No simultaneous invariant-sector reduction]
\label{cor:v22-no-similarity}
The rational endpoint matrices \(K,D_s\in M_{45}(\mathbb Q(i))\) generate
all of \(M_{45}(\mathbb Q(i))\).  After any field extension they have no
common nonzero proper invariant subspace.  In particular, no common change
of basis can put the endpoint raw producer into a nontrivial simultaneous
block-diagonal or block-upper-triangular form.
\end{corollary}

\begin{proof}
The nonzero pivots certifying (26.10) and the nonzero scalars certifying
the \(256\) interactions are reductions of rational Gaussian numbers whose
denominators are invertible modulo \(p\).  Nonvanishing after reduction
implies nonvanishing before reduction.  Repeating the proof of
Theorem~\ref{thm:v22-family-criterion} over \(\mathbb Q(i)\) gives the full
rational matrix algebra.  Scalar extension preserves the equality with
the full matrix algebra, whose natural module has no nonzero proper
invariant submodule.
\end{proof}

\subsection{Standard tensor reshapes of the bridge}

For \(d=mn\), identify the standard coordinate index with a pair
\((a,b)\in\{0,\ldots,m-1\}\times\{0,\ldots,n-1\}\).  Define the
Kronecker rearrangement
\[
 {\cal R}_{m,n}(K)_{(a,c),(b,e)}
              :=K_{(a,b),(c,e)}.                             \tag{26.13}
\]

\begin{lemma}[Rearrangement rank equals Kronecker rank]
\label{lem:v22-kronecker}
The least \(r\) for which
\[
                         K=\sum_{j=1}^r B_j\otimes C_j        \tag{26.14}
\]
in the standard \(m\times n\) reshape equals
\(\operatorname{rank}{\cal R}_{m,n}(K)\).
\end{lemma}

\begin{proof}
Rearrangement sends a simple tensor \(B\otimes C\) to
\(\operatorname{vec}(B)\operatorname{vec}(C)^{\mathsf T}\), a rank-one
matrix.  Hence any representation (26.14) bounds the rearrangement rank
by \(r\).  Conversely, a rank decomposition of the rearranged matrix into
\(r\) outer products can be unvectorized term by term, and inverse
rearrangement gives (26.14).
\end{proof}

\begin{theorem}[Maximal standard Kronecker ranks]
\label{thm:v22-kronecker-gate}
For both the Wilson and endpoint bridges, exact elimination gives
\[
\begin{array}{c|c|c}
(m,n)&\operatorname{shape}{\cal R}_{m,n}(K)&
\operatorname{rank}{\cal R}_{m,n}(K)\\ \hline
(3,15)&9\times225&9\\
(15,3)&225\times9&9\\
(5,9)&25\times81&25\\
(9,5)&81\times25&25 .
\end{array}                                                    \tag{26.15}
\]
Every rank in (26.15) is the maximum possible.  Thus no low-Kronecker-rank
compression exists in any standard ordered nontrivial factorization of
\(45\).
\end{theorem}

\begin{proof}
The executable forms (26.13) entry by entry over \(\mathbb F_{p^2}\) and
computes the ranks in (26.15).  A nonzero maximal minor after reduction is
a nonzero rational maximal minor, so the same ranks hold over
\(\mathbb Q(i)\).  Lemma~\ref{lem:v22-kronecker} gives the conclusion.
\end{proof}

\begin{firewall}[Scope of the algebra obstruction]
\label{fire:v22-algebra}
Corollary~\ref{cor:v22-no-similarity} rules out exact common invariant
sectors for the finite endpoint matrices, and
Theorem~\ref{thm:v22-kronecker-gate} rules out the stated standard
Kronecker compressions.  These facts do not exclude a trace-specific
identity, a nonstandard data-sparse representation, or a faster
matrix-multiplication implementation.  Full matrix algebra is therefore
an obstruction to a class of producer designs, not a lower bound for every
algorithm computing the three scalar traces.
\end{firewall}

\subsection{Executable provenance and V23 checkpoint}

\begin{record}[V22 verifier]
\label{rec:v22-verifier}
The executable
\begin{center}
\path{scripts/verify_ym_postarchive_v22_full_algebra_gate.py}
\end{center}
has \(464\) physical lines, \(16502\) bytes, and SHA--256
\[
\texttt{F365C9ED1F711B5350881E76CBDB1C905F8FFA347ACA84E8577A3998DF564150}.
                                                               \tag{26.16}
\]
Its canonical payload SHA--256 is
\[
\texttt{5418222A1C616A7EF86064C69F94C801E04C81C8976CE8C0B8792AE25F9F52D7}.
                                                               \tag{26.17}
\]
It certifies all forty rank-one reconstructions, the deterministic
sixteen-member selection, both full Krylov spans, every directed coupling,
and every rearrangement rank.
\end{record}

\begin{decision}[V23: trace-specific compression and measured kernel]
\label{dec:v22-v23}
General simultaneous matrix reduction is now closed for the endpoint
producer.  V23 must exploit the fact that only the three cyclic traces in
(24.19)--(24.21), with \(C\)-degree at most eight, are required.  It should:
\begin{enumerate}
\item compile the exact cyclic-word classes after the polynomial expansion
      (25.16), without enumerating the rejected V13 coefficient paths;
\item test whether the needed trace functionals span a proper quotient of
      the full matrix algebra on the actual response families;
\item implement and time an exact batched raw kernel on a reproducible grid
      tile, then derive a whole-grid time and memory gate.
\end{enumerate}
If the trace quotient is also full and the measured run is impractical,
the programme must return upstream to the choice of finite certificate,
not disguise the cost with an unexecuted \(31^4\) prescription.
\end{decision}

\begin{assessment}[V22 acquisition and boundary]
\label{ass:v22}
V22 upgrades the V20 coordinate-connectivity observation to a
basis-independent full-algebra theorem and closes the standard tensor
shortcut.  It is a negative but decisive structural result.  No scalar
response residue has yet been produced.
\end{assessment}

\begin{record}[V22 append record]
\label{rec:v22-append}
V22 is appended to the clean V21 whose installed SHA--256 was
\[
\texttt{B020585A4AB1BD97B9EBF22C4AFE801D65B040D1DAFA71802D2109996968EFF6}.
                                                               \tag{26.18}
\]
After validation only V22 is to remain the active numeric YM note.  The
next compulsory companion audit remains V25.
\end{record}

\begin{firewall}[Claim boundary after V22]
\label{fire:v22-final}
The full-algebra obstruction proves neither a sign nor a gap.  The scalar
residue array, CRT lift, Haar inequality, complement and derivative
sectors, cofinal interacting measure, regulator removal, OS reconstruction,
physical clustering, and Hamiltonian Yang--Mills mass gap remain open.
\end{firewall}

% =============================================================================
% END APPENDED POST-TENTH-ARCHIVE V22 MATERIAL
% =============================================================================

% =============================================================================
% BEGIN APPENDED POST-TENTH-ARCHIVE V23 MATERIAL
% =============================================================================

\section{V23: trace-influence obstruction and exact batched endpoint kernel}
\label{sec:v23}

V22 rules out every common invariant-sector reduction of the endpoint raw
matrices.  V23 asks the narrower question that full matrix algebra leaves
open: can the three required traces ignore a fixed subset of the entries of
\[
                    H(z):=K{\cal S}(I-C(z)K)?                 \tag{27.1}
\]
The exact support answer is again negative.  Nevertheless, batching the
spectral points, evaluating each raw Laurent bank by one exact limb product,
forming the degree-four polynomial with three products, and contracting the
traces directly gives a working exact kernel.  A reproducible host
benchmark moves one prime from an abstract \(5.5\cdot10^{11}\)-operation
ledger to a projected run of about \(1.36\) hours.

\subsection{Which entries of the trace kernel can matter?}

For a raw response bank \(R\), let \(A_R\subset\{1,\ldots,45\}^2\) be the
union of the nonzero patterns of all its Laurent coefficients.  Let
\(\operatorname{row}(A_R)\) and \(\operatorname{col}(A_R)\) be its active
row and column sets.

\begin{lemma}[Coefficientwise trace-influence masks]
\label{lem:v23-influence}
For an otherwise arbitrary matrix \(H\), the linear term
\(\operatorname{Tr}(HQ_2)\) can use precisely the mask
\[
                         A_{Q_2}^{\mathsf T}.                 \tag{27.2}
\]
In
\[
                 \operatorname{Tr}(HP_2HP_0),                \tag{27.3}
\]
the first copy of \(H\) can use only
\[
 \operatorname{col}(A_{P_0})\times
 \operatorname{row}(A_{P_2}),                                \tag{27.4}
\]
and the second can use only
\[
 \operatorname{col}(A_{P_2})\times
 \operatorname{row}(A_{P_0}).                                \tag{27.5}
\]
The two copies in \(\operatorname{Tr}(HP_1HP_1)\) use
\[
 \operatorname{col}(A_{P_1})\times
 \operatorname{row}(A_{P_1}).                                \tag{27.6}
\]
These are uniform coefficient-pattern masks; cancellation at an individual
spectral point may make them smaller.
\end{lemma}

\begin{proof}
The linear identity
\[
                  \operatorname{Tr}(HQ_2)
                  =\sum_{i,j}H_{ij}(Q_2)_{ji}                \tag{27.7}
\]
gives (27.2).  Expanding (27.3) gives
\[
 \operatorname{Tr}(HP_2HP_0)
 =\sum_{i,j,k,l}H_{ij}(P_2)_{jk}H_{kl}(P_0)_{li}.             \tag{27.8}
\]
The first \(H_{ij}\) can occur only if row \(j\) of \(P_2\) and column
\(i\) of \(P_0\) are active, proving (27.4).  Reading the same monomial at
the second copy gives (27.5).  Replacing both responses by \(P_1\) gives
(27.6).
\end{proof}

\begin{theorem}[No uniform entry projection]
\label{thm:v23-mask-gate}
The exact influence counts are
\[
\begin{array}{c|r|r}
\text{mask}&\text{Wilson}&\text{endpoint}\\ \hline
A_{Q_2}^{\mathsf T}&232&1407\\
\text{first mask in (27.3)}&1755&2025\\
\text{second mask in (27.3)}&1755&2025\\
\text{mask in (27.6)}&1764&2025\\ \hline
\text{union}&2025&2025 .
\end{array}                                                    \tag{27.9}
\]
For endpoint each of \(P_0,P_1,P_2,Q_2\) has all \(45\) rows and all
\(45\) columns active.  Hence no fixed proper entry projection of \(H\)
can retain all three endpoint trace functionals by coefficient-pattern
support alone.
\end{theorem}

\begin{proof}
The verifier forms every union pattern from all exact raw coefficients,
computes the row and column sets, constructs (27.2), (27.4)--(27.6) by
Boolean outer products, and counts their union.  Lemma
\ref{lem:v23-influence} turns the resulting count \(45^2\) into the claimed
negative gate.
\end{proof}

\begin{firewall}[A full mask is not a universal trace lower bound]
\label{fire:v23-mask}
Theorem~\ref{thm:v23-mask-gate} excludes only a uniform shortcut obtained
by discarding entries according to coefficient support.  It does not prove
that every scalar algorithm must explicitly form all of \(H\), nor does it
exclude cancellations after the Laurent phases are inserted.
\end{firewall}

\subsection{Exact batched arithmetic}

Write each residue in two \(16\)-bit limbs.  The existing first-prime
matrix product uses three real products and the identity
\[
 (a+ib)(c+id)=(ac-bd)+i\bigl((a+b)(c+d)-ac-bd\bigr).          \tag{27.10}
\]
V23 applies the same construction to batches and to rectangular phase
panels.

\begin{lemma}[Binary64 limb exactness for every V23 contraction]
\label{lem:v23-limb}
Every integer dot product performed through binary64 in the V23 kernel is
exact before modular reduction.
\end{lemma}

\begin{proof}
The longest dot product is a final trace contraction of length
\(45^2=2025\).  Each sum formed by the three-product limb identity is
strictly bounded by
\[
                   2025(2^{17})^2<2^{46}<2^{53}.             \tag{27.11}
\]
The matrix products have length \(45\), and the phase-panel products have
length at most \(99\), so they satisfy the same bound with room to spare.
All limb factors and all partial sums are integers below the consecutive
integer range of binary64.  Rounding therefore returns the exact integer,
after which the Mersenne reduction modulo \(2^{31}-1\) is exact.
\end{proof}

\begin{theorem}[Batched trace-only endpoint kernel]
\label{thm:v23-kernel}
For a batch of spectral points the executable:
\begin{enumerate}
\item evaluates all five raw Laurent banks by rectangular exact limb
      products;
\item forms \(\widetilde E=I-CK\), then computes
      \(\widetilde E^2,\widetilde E^3,\widetilde E^4\) with three products;
\item forms \(H=K{\cal S}(\widetilde E)\) and
      \(HP_0,HP_1,HP_2\);
\item contracts the linear trace and both bubble traces directly, without
      forming the final trace products.
\end{enumerate}
At the eight points (24.10), all three outputs agree exactly with the
independent V21 reference.  Their transcript SHA--256 is
\[
\texttt{C4708A1577F8ACC6B9381A207390C6B790E39F4B105143B1E2223E59D35AA43F}.
                                                               \tag{27.12}
\]
The retained trace output occupies \(3\cdot2\cdot8=48\) bytes per point.
\end{theorem}

\begin{proof}
Lemma~\ref{lem:v23-limb} proves exactness of every batched product.
The three-product power schedule is algebraically identical to (25.16).
Direct index contraction is the definition of the trace.  For validation,
the executable separately evaluates each Laurent bank coefficient by
coefficient, uses the unbatched V21 matrix path, and compares both residue
components of all three traces pointwise.  Equality at all eight points
and the transcript digest give (27.12).
\end{proof}

\subsection{Measured host gate}

The timing record is deliberately outside the canonical algebraic payload.
After one eight-point warm-up, the finalized \(64\)-point endpoint tile was
run five times on Windows 11, Python \(3.12.10\), and NumPy \(2.5.1\).

\begin{record}[V23 endpoint tile benchmark]
\label{rec:v23-benchmark}
The five durations in seconds were
\[
 0.339520200,\quad0.319882900,\quad0.335220000,\quad
 0.351704700,\quad0.345062800.                               \tag{27.13}
\]
The median was \(0.339520200\) seconds, or
\(188.501302731\) points per second.  Linear projection to all
\(923521\) points gives
\[
             4899.281790730\ {\rm s}
             =1.360911609\ {\rm h}.                          \tag{27.14}
\]
The \(64\)-point output SHA--256 was
\[
\texttt{195F5738756E92C56A2E3536A9C8EC35C6BC3B06A8D3A4E66077FD2214037588}.
                                                               \tag{27.15}
\]
\end{record}

\begin{assessment}[What the benchmark changes]
\label{ass:v23-benchmark}
The V21 operation ledger made a full prime look prohibitive if interpreted
as scalar Python work.  V23 batches the same exact arithmetic into vectorized
limb products and places one endpoint prime within an hours-scale run on
the present host.  This is sufficient to design a checkpointed producer.
It is not yet sufficient to launch all CRT work: the number of required
primes and the cost of Wilson, inverse transforms, and reconstruction must
be included first.
\end{assessment}

\begin{firewall}[Timing boundary]
\label{fire:v23-timing}
Equation (27.14) is an empirical linear projection after setup, not a
theorem about elapsed time.  Thermal throttling, competing processes,
tile-edge overhead, and storage can change it.  Only (27.2)--(27.12) and
the deterministic payload below are algebraic certificates.
\end{firewall}

\subsection{Executable provenance and V24 checkpoint}

\begin{record}[V23 verifier]
\label{rec:v23-verifier}
The executable
\begin{center}
\path{scripts/verify_ym_postarchive_v23_trace_kernel_gate.py}
\end{center}
has \(431\) physical lines, \(13576\) bytes, and SHA--256
\[
\texttt{CA5941E10CE7E7FF8A0E41FBAFFE0996837C380D57B9A8218414B4F8E28A0A48}.
                                                               \tag{27.16}
\]
Its deterministic canonical payload SHA--256 is
\[
\texttt{E8424CD237698CDB046DCEF5A45F4A46D440FBCC458AD03D71EBDB8E5F208E9C}.
                                                               \tag{27.17}
\]
The empirical timing fields are excluded from that canonical hash.
\end{record}

\begin{decision}[V24: whole-certificate execution gate]
\label{dec:v23-v24}
Before starting an hours-long calculation, V24 must return to the V14
height and CRT ledger and combine it with the measured V23 kernel.  It
should determine:
\begin{enumerate}
\item the exact number of primes needed for each scalar trace or for the
      signed aggregate;
\item whether reconstruction can be certified from a smaller adaptive
      prime set using independent height bounds;
\item the Wilson and endpoint tile rates, inverse-transform cost, retained
      state, and checkpoint format.
\end{enumerate}
If this whole-certificate gate is practical, V24 should implement a
resumable first-prime producer and validate two disjoint tiles against the
reference path.  If the prime multiplier makes the calculation
impractical, it must reduce the height bound upstream before launching
the grid.
\end{decision}

\begin{assessment}[V23 acquisition and boundary]
\label{ass:v23}
V23 closes fixed entry projection, supplies an exact batched trace-only
kernel, and measures an hours-scale first-prime endpoint run.  It does not
yet compute the complete first-prime array, inverse transform it, or lift
an integer coefficient.
\end{assessment}

\begin{record}[V23 append record]
\label{rec:v23-append}
V23 is appended to the clean V22 whose installed SHA--256 was
\[
\texttt{455387C7F28001CD4A1EF118931D19C7F65EDD3D39C03ABE9AC2BC215317C93E}.
                                                               \tag{27.18}
\]
After validation only V23 is to remain the active numeric YM note.  The
next compulsory companion audit remains V25.
\end{record}

\begin{firewall}[Claim boundary after V23]
\label{fire:v23-final}
No full scalar grid, CRT reconstruction, Haar sign, complement-sector
estimate, cofinal interacting measure, regulator removal, OS
reconstruction, physical clustering, or Hamiltonian Yang--Mills mass gap
has been proved.
\end{firewall}

% =============================================================================
% END APPENDED POST-TENTH-ARCHIVE V23 MATERIAL
% =============================================================================

% =============================================================================
% BEGIN APPENDED POST-TENTH-ARCHIVE V24 MATERIAL
% =============================================================================

\section{Transform-compatible CRT execution gate (V24)}
\label{sec:v24}

\subsection{Purpose and upstream correction}

V23 made one exact spectral prime computationally plausible.  Before
launching a long calculation, V24 returns upstream to the V14 height
certificate and asks whether its CRT moduli actually support the
length-\(31\) Fourier grid used from V19 onward.  This audit exposes a real
interface defect: the V14 uniqueness bound is correct, but thirteen of its
fifteen moduli admit no element of order \(31\), even after passing to the
quadratic Gaussian extension.  Thus the old list cannot execute the later
transform.  The present section replaces only that list, proves the
replacement sufficient, validates the signed spectral assembly on a tile,
and specifies a resumable whole-prime checkpoint.

\begin{definition}[Requirements on an executable modulus]
\label{def:v24-modulus}
Let \(N=31\).  A prime \(p<2^{31}\) is called \emph{V24-compatible} when
\[
 p\equiv3\pmod4,\qquad N\mid p-1,\qquad \gcd(p,\Delta)=1,       \tag{28.1}
\]
where the V14 signed common denominator is
\[
 \Delta=2^{126}3^{11}5^{82}\cdot13\cdot37.                    \tag{28.2}
\]
The first condition makes \(X^2+1\) irreducible over \(\mathbb F_p\), so
the pair arithmetic used by the kernel realizes
\(\mathbb F_{p^2}=\mathbb F_p[i]\).  The second places an order-\(31\)
Fourier root already in the scalar subfield.
\end{definition}

\begin{lemma}[Order-\(31\) criterion]
\label{lem:v24-root-criterion}
For a prime \(p\ne31\), the field \(\mathbb F_{p^2}\) contains an element
of order \(31\) if and only if \(31\mid p^2-1\).  It contains such an
element in \(\mathbb F_p\) if and only if \(31\mid p-1\).
\end{lemma}

\begin{proof}
The multiplicative groups of the two finite fields are cyclic of orders
\(p^2-1\) and \(p-1\), respectively.  A finite cyclic group has an element
of prime order \(31\) exactly when \(31\) divides its order.
\end{proof}

\begin{theorem}[The V14 modulus list is transform-incompatible]
\label{thm:v24-old-list}
Number the fifteen V14 primes in their recorded order.  Exactly entries
\(1\) and \(7\) satisfy \(31\mid p-1\), and exactly the same two satisfy
\(31\mid p^2-1\).  Consequently thirteen V14 moduli cannot carry the
length-\(31\) Fourier transform even in \(\mathbb F_{p^2}\).
\end{theorem}

\begin{proof}
The verifier reduces each of the fifteen recorded integers modulo \(31\).
The residues are
\[
 1,3,26,10,21,25,1,18,29,21,4,2,18,10,27.                   \tag{28.3}
\]
Since \(31\mid p^2-1\) precisely when \(p\equiv1\) or \(-1\pmod{31}\),
only the two displayed residues equal to \(1\) qualify; none equals \(30\).
Lemma~\ref{lem:v24-root-criterion} gives the conclusion.
\end{proof}

\begin{firewall}[What is and is not corrected]
\label{fire:v24-v14}
Theorem~\ref{thm:v24-old-list} does not invalidate the V14 height bound or
its abstract CRT uniqueness implication.  It invalidates the use of that
particular prime list with the later length-\(31\) spectral algorithm.
The correction below retains the V14 bound and replaces the moduli.
\end{firewall}

\subsection{Certified replacement moduli}

The simultaneous congruence
\[
                         p\equiv63\pmod{124}                 \tag{28.4}
\]
implies \(p\equiv3\pmod4\) and \(p-1\equiv62\pmod{124}\), hence
\(31\mid p-1\).  Searching downward from \(2^{31}-1\) in this progression
gives the following list.  The table also records
\(\omega=2^{(p-1)/31}\bmod p\) and its inverse.
\[
\begin{array}{r|r|r}
p&\omega&\omega^{-1}\\ \hline
2147483647&4096&524288\\
2147483399&450851339&464861500\\
2147480299&1165718630&1649503587\\
2147479307&812766432&592660951\\
2147478563&799606820&784802703\\
2147477323&286518496&2003217633\\
2147476951&1664145400&519555813\\
2147475587&2064147492&561159384\\
2147474843&1504384037&1023829032\\
2147473231&704944775&1565482360\\
2147472611&1472976815&680800178\\
2147470751&1695997229&1913953696\\
2147470627&1893255531&472353335\\
2147470007&659340190&1016799973\\
2147469263&1893760090&1489689141
\end{array}                                                   \tag{28.5}
\]

\begin{theorem}[Primality, roots, and denominator compatibility]
\label{thm:v24-new-list}
Every integer in (28.5) is prime, is V24-compatible, and the displayed
\(\omega\) has exact order \(31\) in \(\mathbb F_p^\times\).
\end{theorem}

\begin{proof}
For each candidate, the verifier divides by every one of the \(4792\)
primes at most
\(\lfloor\sqrt{2^{31}-1}\rfloor=46340\).  No divisor is found for any
accepted integer, which is a complete trial-division primality proof.
Congruence (28.4), checked for every entry, gives the first two conditions
of (28.1).  None divides (28.2), as verified by the Euclidean algorithm.
Finally, modular exponentiation gives \(\omega^{31}=1\) and
\(\omega\ne1\) for every row.  Because \(31\) is prime, the order is
exactly \(31\).  The inverse column is checked by
\(\omega\omega^{-1}\equiv1\pmod p\).
\end{proof}

\begin{record}[Search and list transcript]
\label{rec:v24-primes}
The descending search tests \(117\) candidates.  The SHA--256 of the
ordered prime list is
\[
\texttt{0423622DBA6E34296584AD78DAFADA339D307137585753E4624D4C4DCC396423}.
                                                               \tag{28.6}
\]
\end{record}

\subsection{Height sufficiency and signed reconstruction}

The V24 executable independently recompiles the actual Wilson and endpoint
banks and reproduces the V14 signed numerator bound
\[
\begin{split}
B={}&2495959300085994965303110622765063632373945291881826888809897946\\
   &181991838659246076960159761242783822699395079285278562680832000000000,
\end{split}                                                    \tag{28.7}
\]
which has \(444\) bits and SHA--256
\[
\texttt{144B73BC8DB36A043EB6A7F4F0016011B9FB61D0549A3321302DF82B4005D420}.
                                                               \tag{28.8}
\]
The common denominator has \(343\) bits and SHA--256
\[
\texttt{91FD90C6F1BA3C9AEBD08235C84FD1EFA15D7EAC41240BCCE1CF362DD7855A03}.
                                                               \tag{28.9}
\]

\begin{theorem}[Minimal 31-bit CRT cardinality]
\label{thm:v24-crt-height}
Let \(M\) be the product of the fifteen primes in (28.5).  Then
\[
\begin{split}
M={}&9526295680543063045474653463761972593388397514444528338906763762\\
   &4770248780487195240422597713771999621047331141083919587802037888499805335207
\end{split}                                                    \tag{28.10}
\]
has \(465\) bits and satisfies \(M>2B\).  The product of the first fourteen
entries does not exceed \(2B\).  Moreover, no collection of fourteen
positive moduli below \(2^{31}\) can exceed \(2B\).  Thus fifteen is the
minimal cardinality in the stipulated 31-bit architecture.
\end{theorem}

\begin{proof}
Exact integer multiplication gives (28.10) and direct comparison gives
\[
 \prod_{j=1}^{14}p_j\le 2B<\prod_{j=1}^{15}p_j.               \tag{28.11}
\]
For the global cardinality statement, the product of any fourteen positive
integers below \(2^{31}\) is less than \(2^{434}\), whereas \(B\) has
\(444\) bits, so \(2B\ge2^{444}>2^{434}\).
\end{proof}

\begin{corollary}[Unique signed lift]
\label{cor:v24-signed-lift}
If an integer coefficient \(a\) satisfies \(|a|\le B\), its residues
modulo the primes in (28.5) determine \(a\) uniquely.
\end{corollary}

\begin{proof}
Two possible lifts differ by a multiple of \(M\), while their difference
has absolute value at most \(2B<M\).
\end{proof}

\subsection{Signed spectral observable}

At a spectral point let
\[
\Psi_e=\frac{1}{2}L_e+
 2c_e\bigl(B_{20,e}-B_{11,e}\bigr),\qquad
 c_{\rm W}=\frac{1}{3},\quad c_{\rm end}=3,                  \tag{28.12}
\]
where \(L_e\), \(B_{20,e}\), and \(B_{11,e}\) are the exact V23 trace
outputs for action \(e\).  The producer needs only
\[
             F=\Delta\bigl(\Psi_{\rm end}-\Psi_{\rm W}\bigr). \tag{28.13}
\]

\begin{lemma}[Integral signed assembly]
\label{lem:v24-signed}
Every Laurent coefficient of \(F\) is a Gaussian integer.  For every
prime in (28.5), pointwise reduction of (28.13) is well-defined and the
inverse length-\(31\) tensor transform returns the corresponding coefficient
residue of \(F\).
\end{lemma}

\begin{proof}
The V14 denominator certificate states that \(\Delta\) clears all rational
denominators in both terms of (28.13), including those in (28.12);
therefore \(F\in\mathbb Z[i][z_1^{\pm1},\ldots,z_4^{\pm1}]\).
Theorem~\ref{thm:v24-new-list} gives \(\gcd(p,\Delta)=1\), so all rational
operations and multiplication by \(\Delta\) reduce consistently modulo
\(p\).  It also supplies an exact order-\(31\) root.  Orthogonality,
\[
 \frac{1}{31}\sum_{r=0}^{30}\omega^{r(k-\ell)}
   =\begin{cases}1,&k\equiv\ell\pmod{31},\\0,&k\not\equiv\ell\pmod{31},
     \end{cases}                                               \tag{28.14}
\]
follows by the finite geometric-series identity.  Applying (28.14) in
each coordinate proves the tensor interpolation statement.
\end{proof}

\begin{theorem}[Independent signed-tile gate]
\label{thm:v24-tile}
For \(p=2147483647\), the final \(64\)-point V23 Wilson and endpoint
kernels have output hashes
\[
\begin{split}
{\rm Wilson}:&\quad
\texttt{400FA2F08941DCB18661596CDFA5E946206037A6CA272D02E316FC9182C0BE0C},\\
{\rm endpoint}:&\quad
\texttt{195F5738756E92C56A2E3536A9C8EC35C6BC3B06A8D3A4E66077FD2214037588}.
\end{split}                                                    \tag{28.15}
\]
After exact assembly of (28.13), the signed tile SHA--256 is
\[
\texttt{302B75C315019ABC934FB6EE11D786CD5DCB1E9003DD7770ABC85E2E45ABD48E}.
                                                               \tag{28.16}
\]
\end{theorem}

\begin{proof}
The action kernels are the independently checked V23 kernels.  V24
multiplies their residue outputs by the reductions of
\(\Delta/2\), \(2\Delta c_{\rm W}\), and
\(2\Delta c_{\rm end}\), combines them according to (28.12)--(28.13), and
hashes both residue limbs in fixed point order.  Exact modular arithmetic
and Lemma~\ref{lem:v24-signed} prove that the resulting tile is the
reduction of \(F\).
\end{proof}

\begin{firewall}[A tile is not a coefficient certificate]
\label{fire:v24-tile}
Equation (28.16) validates the signed assembly and data convention on one
tile.  It neither evaluates the remaining \(923457\) spectral points nor
performs an inverse transform.  No Laurent coefficient or Haar sign is
therefore certified here.
\end{firewall}

\subsection{Checkpoint and whole-certificate cost}

\begin{proposition}[Lossless per-prime checkpoint]
\label{prop:v24-checkpoint}
A complete signed spectral grid for one prime consists of two
\(\mathrm{uint32}\) arrays of length \(31^4=923521\), one for each
\(\mathbb F_p\) component.  It occupies
\[
                     2\cdot4\cdot923521=7388168\ {\rm bytes}. \tag{28.17}
\]
With tiles of \(64\) points there are \(14431\) tiles.  A sufficient
resumption record contains the prime, \(\omega\), the next flat point,
the producer digest, and a rolling output digest.  A direct separable
four-dimensional inverse DFT requires
\[
                    4\cdot31^5=114516604                    \tag{28.18}
\]
\(\mathbb F_{p^2}\) multiply--accumulates.
\end{proposition}

\begin{proof}
There are \(31^4\) points and each field element is represented by two
four-byte residues, proving (28.17).  Ceiling division gives
\(\lceil923521/64\rceil=14431\).  Along one coordinate there are \(31^3\)
lines, each with \(31\) outputs and \(31\) summands, hence \(31^5\)
field multiply--accumulates.  There are four coordinates, proving
(28.18).  The listed metadata identifies the arithmetic instance and an
unambiguous completed prefix, while the rolling digest detects corruption.
\end{proof}

\begin{record}[V24 host timing]
\label{rec:v24-timing}
Five finalized \(64\)-point runs give Wilson durations
\[
0.287415000,\ 0.299039800,\ 0.304080300,\ 0.287397700,\
0.293984800\ {\rm s},                                        \tag{28.19}
\]
with median \(0.293984800\) seconds and projected full-grid time
\(1.178390349\) hours.  Endpoint durations are
\[
0.304937300,\ 0.312293000,\ 0.290507600,\ 0.278887700,\
0.276282000\ {\rm s},                                        \tag{28.20}
\]
with median \(0.290507600\) seconds and projected full-grid time
\(1.164452557\) hours.  Linear projection for both actions at all fifteen
primes is \(35.142643590\) kernel-hours, excluding transforms, checkpoint
I/O, and reconstruction.
\end{record}

\begin{firewall}[Timing remains empirical]
\label{fire:v24-timing}
The measured durations and their projections are not included in the
canonical certificate hash.  They authorize only a checkpointed first-prime
trial, not an unattended thirty-five-hour calculation.
\end{firewall}

\subsection{Executable provenance and V25 decision}

\begin{record}[V24 verifier]
\label{rec:v24-verifier}
The executable
\begin{center}
\path{scripts/verify_ym_postarchive_v24_transform_crt_execution_gate.py}
\end{center}
has \(374\) physical lines, \(12918\) bytes, and SHA--256
\[
\texttt{E43B00508466141CC968D98CFBC10934CEB2F3E94E11C7AFE9C218D483DB620D}.
                                                               \tag{28.21}
\]
Its deterministic canonical payload SHA--256 is
\[
\texttt{F5C92543D4B726A7612E2B2D8754E4BB703648DAE751736877A43ED240A52AE3}.
                                                               \tag{28.22}
\]
The product digest is
\[
\texttt{B59685B6DF3874BEF781650F2095DE7D2B0767B638C9BFE345D8C25793F9FB45}.
                                                               \tag{28.23}
\]
\end{record}

\begin{decision}[V25: audit, then resumable first prime]
\label{dec:v24-v25}
V25 is a compulsory fifth-version companion audit.  After reading the
then-current SM and NS notes, it should implement a prime-parameterized,
atomic, resumable producer for (28.13), validate two disjoint tiles against
the scalar reference, and complete the \(p=2147483647\) grid before
authorizing the other fourteen primes.  It must also implement and verify
the separable inverse transform, so the first actual coefficient residues
rather than spectral samples are obtained.
\end{decision}

\begin{assessment}[V24 acquisition and boundary]
\label{ass:v24}
V24 repairs the CRT/transform interface, proves that fifteen compatible
31-bit primes are necessary and sufficient for the present height bound,
and closes the signed tile and checkpoint design gates.  It does not
complete one prime, invert a grid, reconstruct an integer coefficient, or
prove the Haar sign.
\end{assessment}

\begin{record}[V24 append record]
\label{rec:v24-append}
V24 is appended to the clean V23 whose installed SHA--256 was
\[
\texttt{5DA58A569F82DF22D7F070266BCFBA353241F3E1653595E343AC396375551878}.
                                                               \tag{28.24}
\]
After validation only V24 is to remain the active numeric YM note.
\end{record}

\begin{firewall}[Claim boundary after V24]
\label{fire:v24-final}
No complete spectral prime, inverse-transform coefficient array, CRT
reconstruction, Haar sign, complement-sector estimate, cofinal interacting
measure, regulator removal, Osterwalder--Schrader reconstruction, physical
clustering, or Hamiltonian Yang--Mills mass gap has been proved.
\end{firewall}

% =============================================================================
% END APPENDED POST-TENTH-ARCHIVE V24 MATERIAL
% =============================================================================

% =============================================================================
% BEGIN APPENDED POST-TENTH-ARCHIVE V25 MATERIAL
% =============================================================================

\section{Mandatory companion audit at V25}
\label{sec:v25}

\subsection{Frozen sources}

Before further execution work, the current active companion notes in
\path{latex_notes} were inspected read-only:
\[
\begin{array}{c|r|l}
\text{source}&\text{lines}&\text{SHA--256}\\ \hline
\text{SM V97}&28879&
\texttt{D5B2414339CE0DD358CFA116ABA1CB59C3C82E4228B012069382867D22D084E1}\\
\text{NS V26}&5319&
\texttt{82C887F8C2C69E4F28FAD7C3DC8DE5B9099CB5A4BF431EBA1FFF3E91935978AD}
\end{array}                                                    \tag{29.1}
\]
Here SM V97 is
\path{Renewal_SM_Einstein_Continuum_Research_Notes_V97.tex}, and NS V26 is
\path{Renewal_NS_Stability_Research_Notes_V26.tex}.  The latter has not
changed since the V20 audit.

\subsection{SM V96--V97}

SM V96 defines the dyadic coarse link by the ordered product of its two
fine links.  It proves associativity, gauge covariance, and an exact
decomposition of a coarse plaquette into four fine plaquettes transported
to a common base point.  For the continuum-scaled curvature its derivative
is the four-face average \(B_2=\frac14S_2\).  SM V97 then proves, on its
rooted small-field chart,
\[
 \left\|{\cal G}^{\rm c}_N-B_2{\cal G}_{2N}\right\|_\infty
 <\frac{1}{2000000N}.                                      \tag{29.2}
\]
Its own firewall states that this is a one-step kinematic estimate, not
control of accumulated refinement errors, compatible probability laws, or
an interacting continuum measure.

\begin{proposition}[Consequence for the YM acquisition order]
\label{prop:v25-sm-consequence}
The SM audit supplies a justified later continuum architecture:
coarsen gauge-covariant link words before comparing curvature charts.
It does not alter the immediate V24 decision to compute the finite
signed Haar certificate.
\end{proposition}

\begin{proof}
The V24 obstruction is an unevaluated finite Laurent coefficient problem
for a fixed local response producer.  Equation (29.2) presupposes a
small-field chart and controls a distinct map between lattice scales; it
contains neither the signed response coefficient nor its Haar average.
Conversely, once the local sign problem is settled, the exact associative
link block and its \(O(N^{-1})\) nonlinear estimate are the appropriate
objects for the later cofinal compatibility stage.  Importing them before
the finite sign is known would change no present premise or conclusion.
\end{proof}

\subsection{NS V26}

NS V26 contracts first and second Oseen-kernel derivatives against
rank-one inputs.  Its first-derivative norm gives no improvement over the
full symmetric-tracefree norm; its second-derivative norm improves only on
a bounded radial window, with a maximal displayed pointwise gain of about
eleven percent.  Those facts sharpen a Navier--Stokes ancient-form ledger
but contain no Yang--Mills group integral, Laurent support, or finite-field
transform.

\begin{theorem}[V25 audit decision]
\label{thm:v25-audit-decision}
No estimate is imported from SM V97 or NS V26 into the finite V24
certificate.  V26 should implement and validate the resumable first-prime
producer of Decision~\ref{dec:v24-v25}.  The SM link-block result is retained
as a routing constraint for the later continuum stage.
\end{theorem}

\begin{proof}
Proposition~\ref{prop:v25-sm-consequence} proves the SM statement.  The
NS quantities act on Oseen kernels and rank-one velocity tensors, whereas
the V24 task is exact arithmetic in a finite Laurent response algebra.
There is no common operator or hypothesis from which a valid numerical
constant could be transferred.
\end{proof}

\begin{firewall}[Audit boundary]
\label{fire:v25-audit}
Reading a companion theorem does not prove its hypotheses in the present
Yang--Mills construction.  In particular, (29.2) is not a YM measure-flow
estimate, and the NS rank-one reductions are not response-matrix
reductions.
\end{firewall}

\begin{assessment}[V25 acquisition]
\label{ass:v25}
The mandatory audit is complete.  It adds an exact-word-first constraint
to the later cofinal route and leaves the finite spectral bottleneck
unchanged.  No new Haar sign or mass-gap claim is made.
\end{assessment}

\begin{record}[V25 append record]
\label{rec:v25-append}
V25 is appended to the clean V24 whose installed SHA--256 was
\[
\texttt{3F49865A90958502A6457AEF65118AC6E80649F6F039AF6B45D710CFF0A3EF02}.
                                                               \tag{29.3}
\]
After validation only V25 is to remain the active numeric YM note.  The
next compulsory companion audit is V30.
\end{record}

\begin{firewall}[Claim boundary after V25]
\label{fire:v25-final}
No complete spectral prime, inverse-transform coefficient array, CRT
reconstruction, Haar sign, compatible interacting measure, regulator
removal, Osterwalder--Schrader reconstruction, physical clustering, or
Hamiltonian Yang--Mills mass gap has been proved.
\end{firewall}

% =============================================================================
% END APPENDED POST-TENTH-ARCHIVE V25 MATERIAL
% =============================================================================

% =============================================================================
% BEGIN APPENDED POST-TENTH-ARCHIVE V26 MATERIAL
% =============================================================================

\section{V26: completion of the first signed spectral prime}
\label{sec:v26}

V24 reduced the finite signed-response problem to a length-\(31\)
four-dimensional transform at fifteen compatible primes.  V25 performed
the compulsory companion audit.  V26 now executes the first complete
prime, but first repairs a root-indexing mismatch in the bounded V24
tile record.  The repair changes neither the Laurent polynomial nor the
CRT list; it makes the advertised sample ordering exact.

\subsection{Root-index correction}

Let \(F\) be the integral Gaussian Laurent polynomial of
(28.13), and for an order-\(31\) root \(\omega\) write
\[
 {\cal E}_{\omega}(r)
 =F(\omega^{r_1},\omega^{r_2},\omega^{r_3},\omega^{r_4}),
 \qquad r\in(\mathbb Z/31\mathbb Z)^4.                       \tag{30.1}
\]
The V23 kernel used \(\omega_0=512\) at \(p=2147483647\), whereas the
first row of (28.5) lists \(\omega_1=4096\).

\begin{theorem}[Exact root automorphism and corrected tile]
\label{thm:v26-root-correction}
In \(\mathbb F_{2147483647}\),
\[
                    4096=512^{22}.                          \tag{30.2}
\]
Consequently
\[
              {\cal E}_{4096}(r)={\cal E}_{512}(22r)         \tag{30.3}
\]
with multiplication by \(22\) in every coordinate modulo \(31\).
For the first \(64\) points in the declared C-order at root \(4096\),
the signed-residue SHA--256 is
\[
\texttt{2652DAD06370CDF207A0EA91DA5886B9BFC9FFE6FF15F46C9AEC956B80C6AC33}.
                                                               \tag{30.4}
\]
\end{theorem}

\begin{proof}
Modular exponentiation gives (30.2).  Substitution into every Laurent
monomial \(z^\nu\) gives
\[
 (\omega_1^r)^\nu=(\omega_0^{22r})^\nu,
\]
and linearity gives (30.3).  The V26 verifier evaluates two disjoint
\(64\)-point tiles, beginning at flat indices \(0\) and \(461760\).
After applying \(r\mapsto22r\), every Wilson and endpoint value of each
of the three traces \(L,B_{20},B_{11}\) agrees with the V23 kernel.
It then assembles (28.13) at the unpermuted root-\(4096\) points and
hashes the first tile, giving (30.4).
\end{proof}

\begin{remark}[Precise status of the V24 hash]
\label{rem:v26-v24-hash}
The V24 hashes (28.15)--(28.16) remain valid witnesses for the primitive
root \(512\).  They are not hashes of the same fixed flat indices at the
listed root \(4096\).  Equation (30.4) supersedes (28.16) only when the
latter is read as a root-\(4096\) indexed-tile claim.  Since (30.3) is a
permutation of the complete grid, the interpolation and CRT theorems of
V24 are unchanged.
\end{remark}

\subsection{Exact arithmetic for every listed prime}

Write a residue \(a<p<2^{31}\) as
\[
             a=a_0+2^{16}a_1,\qquad 0\le a_0,a_1<2^{16}.     \tag{30.5}
\]
For a real dot product the executable computes the low, high, and
mixed-limb sums in binary64 and then reconstructs
\[
 L+2^{16}(M-L-H)+2^{32}H\pmod p.                             \tag{30.6}
\]

\begin{theorem}[Prime-parameterized limb exactness]
\label{thm:v26-limb}
Every residue returned by the V26 batched kernel is the exact residue of
the corresponding Gaussian-rational expression for every prime in
(28.5).
\end{theorem}

\begin{proof}
The longest real dot product is a \(45^2=2025\)-entry trace
contraction.  The phase panels have inner length at most \(99\), and
matrix products have inner length \(45\).  In the worst mixed product,
both limb factors are below \(2^{17}\), so every nonnegative partial sum
is bounded by
\[
                  2025(2^{17})^2<2^{46}<2^{53}.              \tag{30.7}
\]
Binary64 represents every integer in this range exactly, independently
of summation order; hence \(L,H,M\) and \(M-L-H\) are exact integers.
In (30.6), reduction of each limb precedes multiplication.  The largest
integer product is below \(2^{31}2^{31}=2^{62}<2^{63}\), so signed
64-bit recombination is also exact.  This argument uses only
\(p<2^{31}\), not the Mersenne identity used by the original V23
specialization.

Gaussian multiplication uses three such exact real products and
\[
 (a+ib)(c+id)=(ac-bd)
   +i\bigl((a+b)(c+d)-ac-bd\bigr).                           \tag{30.8}
\]
All additions and reductions are exact in
\(\mathbb F_p[i]\).  Finally, Theorem~\ref{thm:v24-new-list} states that
none of the listed primes divides a required denominator, so exact
rational reduction is defined.  These observations cover every
operation in the producer.
\end{proof}

\subsection{Independent kernel and resumption gates}

\begin{record}[Bounded V26 gates]
\label{rec:v26-gates}
The deterministic verification payload has SHA--256
\[
\texttt{BF75F8462A7BB949D28140813CF069B150ABF4999FBBA714DA36D333B89B8BC5}.
                                                               \tag{30.9}
\]
The two root-permuted tile traces have transcript digest
\[
\texttt{592A22A4503E0CE607DC7C606A0D29ACB940F9CA709BEAF68EC457D65267803B}.
                                                               \tag{30.10}
\]
At the second prime \(2147483399\), root \(450851339\), and point
\((30,30,30,30)\), the fast kernel agrees in both Gaussian limbs with an
independent scalar path for both actions and all three traces.  That
transcript has digest
\[
\texttt{43D41FDA9DF2FC1075CC396209C69B705B74470220D8B806D2FD63938BE9ED70}.
                                                               \tag{30.11}
\]
An uninterrupted four-tile prefix and a two-tile plus resumed two-tile
prefix are byte-identical; their common rolling digest is
\[
\texttt{4BF0DD6E988CFC36795BFEF7EA52B61352A7800C45BB949D97D46F9C69E16CB8}.
                                                               \tag{30.12}
\]
\end{record}

For a committed prefix of length \(n\), define its byte stream by
visiting consecutive \(64\)-point frames, placing that frame's
big-endian real \(\mathrm{uint32}\) residues before its imaginary
residues, and hashing the concatenation.

\begin{proposition}[Process-interruption-safe checkpoint invariant]
\label{prop:v26-resume}
Assume the host filesystem honors successful file flush, file
\texttt{fsync}, and atomic same-directory \texttt{replace}.  After every
completed metadata replacement, the recorded prefix is sufficient to
resume without changing any committed residue.
\end{proposition}

\begin{proof}
Before replacing metadata, the producer flushes and synchronizes both
arrays.  It then synchronizes a temporary metadata file and atomically
replaces the prior metadata.  Thus metadata never declares a prefix
before its two arrays have been synchronized.  On reopening, the
producer checks the schema, prime, root, grid geometry, tile size, source
digest, array shapes and dtypes, and prefix alignment.  It recomputes the
entire prefix digest and refuses to proceed on a mismatch.  If an
interruption occurred after array writes but before metadata replacement,
those extra tail entries are outside the declared prefix and are
overwritten deterministically.  The four-tile split experiment in
Record~\ref{rec:v26-gates} tests this invariant on executed data.
The proposition is an application-level process-interruption statement;
it does not assert durability under a filesystem that violates the
stated host contract.
\end{proof}

\subsection{Completed first-prime artifact}

\begin{theorem}[Complete signed spectral grid at \(p=2^{31}-1\)]
\label{thm:v26-complete-prime}
The V26 producer evaluates all
\[
                       31^4=923521                            \tag{30.13}
\]
points for \(p=2147483647\) and root \(4096\).  Its terminal metadata reads
\[
 \texttt{complete=true},\qquad
 \texttt{next\_flat\_point}=923521.
\]
The canonical full-stream digest is
\[
\texttt{875178AF3DD67A4745BE9F0F2110A3419573D0889BBE308864693DE5E488ABDB}.
                                                               \tag{30.14}
\]
The real and imaginary arrays each have shape \((923521,)\), dtype
\(\mathrm{uint32}\), and \(3694212\) stored bytes.  Their complete
\(\texttt{.npy}\) file hashes are respectively
\[
\begin{split}
\texttt{69D8EC75482E61C5768B79D8BA5ED21BD7FFB5D0919C6DA1C41B790BD66CC7FC},\\
\texttt{CEC419D013D931D9DA241C1CE574DEB8FC300E72E2FDD2A3F0C6010A2D001CD6}.
\end{split}                                                    \tag{30.15}
\]
\end{theorem}

\begin{proof}
Theorem~\ref{thm:v26-limb} proves every produced residue exact.
The process exits with code zero after writing the terminal checkpoint.
An independent read-only pass reloads both arrays, visits all \(14430\)
full \(64\)-point frames and the final one-point frame in the declared
order, and recomputes (30.14) exactly.  The recomputed value equals the
metadata value.  The same pass checks the endpoint, shapes, and dtypes;
independent whole-file SHA--256 calculations give (30.15).
\end{proof}

\begin{record}[Execution and source provenance]
\label{rec:v26-provenance}
The final invocation resumes at point \(1024\), evaluates \(14415\)
tiles, and takes \(6574.290507600\) seconds on the host.  This timing is
empirical and is not part of the canonical payload.  The executable
\begin{center}
\path{scripts/verify_ym_postarchive_v26_resumable_prime_producer.py}
\end{center}
has \(939\) physical lines, \(33487\) bytes, and SHA--256
\[
\texttt{80BA1AE9FF855EE3DC8DB2138870DB3D7D5D720CCC69F01FEE001CA9ECC391D9}.
                                                               \tag{30.16}
\]
The terminal checkpoint file has SHA--256
\[
\texttt{25ECA2D55FF19E293957DD86B53B28E0012C309E329A76F8551AD6B24800D304}.
                                                               \tag{30.17}
\]
\end{record}

\begin{decision}[V27: invert before authorizing more primes]
\label{dec:v26-v27}
V27 should reload and rehash the completed artifact, perform the exact
separable inverse tensor transform, prove zero outside the V19 candidate
support, and transform forward again over the complete grid.  Reflection
pairing must remain after coefficient reconstruction.
\end{decision}

\begin{assessment}[V26 acquisition]
\label{ass:v26}
V26 repairs the root-indexed tile witness and completes the first actual
signed spectral prime with independent on-disk verification.  This closes
the first-prime evaluation bottleneck but not coefficient reconstruction.
\end{assessment}

\begin{record}[V26 append record]
\label{rec:v26-append}
V26 is appended to the clean V25 whose installed SHA--256 was
\[
\texttt{5D6759B911E076304A17D09955E6777C3ACD6C99D3B47CD42BBA75EB0EA96BE6}.
                                                               \tag{30.18}
\]
After validation only V26 is to remain the active numeric YM note.
\end{record}

\begin{firewall}[Claim boundary after V26]
\label{fire:v26-final}
One complete spectral prime does not determine a Laurent coefficient over
\(\mathbb Z[i]\), its sign, or its Haar average.  No inverse coefficient
array, fifteen-prime CRT lift, complement-sector estimate, compatible
interacting measure, regulator removal, Osterwalder--Schrader
reconstruction, physical clustering, or Hamiltonian Yang--Mills mass gap
has been proved.
\end{firewall}

% =============================================================================
% END APPENDED POST-TENTH-ARCHIVE V26 MATERIAL
% =============================================================================
% =============================================================================
% BEGIN APPENDED POST-TENTH-ARCHIVE V27 MATERIAL
% =============================================================================

\section{V27: exact first-prime coefficient reconstruction}
\label{sec:v27}

V26 completed the first signed spectral grid.  V27 turns that grid into
Laurent coefficient residues.  It first reloads and independently hashes
the complete input, then applies an exact separable inverse transform,
checks the independently enumerated support, transforms forward over the
entire grid, and reloads the atomically stored output.

\subsection{Transform convention and no-alias theorem}

Put \(p=2147483647\), \(\omega=4096\), and \(N=31\).  If
\[
 F(z)=\sum_{\nu\in\mathbb Z^4}c_\nu z^\nu,\qquad
 S(r)=F(\omega^{r_1},\ldots,\omega^{r_4}),                   \tag{31.1}
\]
then for \(t\in(\mathbb Z/N\mathbb Z)^4\) define
\[
 {\cal I}S(t)=N^{-4}\sum_{r\in(\mathbb Z/N\mathbb Z)^4}
                  S(r)\omega^{-r\cdot t}.                   \tag{31.2}
\]

\begin{theorem}[Exact inverse tensor transform]
\label{thm:v27-inverse}
For every residue index \(t\),
\[
                 {\cal I}S(t)=\sum_{\nu\equiv t\ ({\rm mod}\ 31)}c_\nu
                 \quad\hbox{in }\mathbb F_p[i].              \tag{31.3}
\]
For the signed V26 polynomial there is at most one supported
\(\nu\) in each residue class, so (31.2) returns its exact coefficient
residue.
\end{theorem}

\begin{proof}
Insert (31.1) into (31.2), exchange the finite sums, and factor the
four coordinates.  For one coordinate,
\[
 N^{-1}\sum_{r=0}^{N-1}\omega^{r(\nu-t)}
 =\begin{cases}1,&\nu\equiv t\pmod N,\\0,&\text{otherwise}.
 \end{cases}                                                  \tag{31.4}
\]
The denominator \(N\) is invertible modulo \(p\), and \(\omega\) has
exact order \(N\), so the geometric-series proof of (31.4) applies.
V19 proves that the signed candidate support is contained in
\([-12,12]^4\).  Two integers in \([-12,12]\) with the same residue
modulo \(31\) differ by a multiple of \(31\) of absolute value at most
\(24\), hence are equal.  This proves injectivity in every coordinate.

The executable applies the \(31\times31\) transform successively along
the four axes, using \(\omega^{-1}=524288\), and finally multiplies by
\((31^4)^{-1}\).  The longest new dot product has length \(31\), so
Theorem~\ref{thm:v26-limb} proves every modular operation exact.
\end{proof}

\begin{record}[Independent support and orientation gates]
\label{rec:v27-gates}
The V19 signed candidate support has \(261797\) frequencies and SHA--256
\[
\texttt{F6C17DCD47F124FD4EDEBF8CC92C3ED23329B4D5D6EBC39953667A114E9088C0}.
                                                               \tag{31.5}
\]
The V27 synthetic gate places four prescribed Gaussian residues at four
distinct supported indices.  Its forward transform agrees with direct
Laurent evaluation at four separately chosen grid points; the inverse
round trip is exact.  The synthetic spectral digest is
\[
\texttt{442E8B745C2643EC315B998CC8DCED4DD0F67940961849D70B90C48EA9255AFF}.
                                                               \tag{31.6}
\]
The direct samples fix the sign and orientation independently of the
forward/inverse round trip.
\end{record}

\subsection{Executed coefficient certificate}

\begin{theorem}[First-prime coefficient residues]
\label{thm:v27-coefficients}
The inverse transform of the complete V26 grid has:
\[
\begin{array}{c|r}
\text{candidate-support cardinality}&261797\\
\text{nonzero Gaussian residue count}&258729\\
\text{nonzero imaginary-limb count}&258681\\
\text{nonzero residues outside candidate support}&0.
\end{array}                                                    \tag{31.7}
\]
In C-order, with index \(t_j=\nu_j\bmod31\), the constant coefficient
residue is
\[
                    c_0=(1452792677,\ 367879794)
                    \quad\hbox{in }\mathbb F_p[i].            \tag{31.8}
\]
The complete coefficient-stream SHA--256, with the full real array
followed by the full imaginary array as big-endian \(\mathrm{uint32}\),
is
\[
\texttt{ABCCFC24307DBFB9352575A0D5AE1FDE5DF5596DD1861A671EB26F5FF9427DC1}.
                                                               \tag{31.9}
\]
Transforming these coefficients forward reproduces both V26 spectral
arrays at all \(923521\) points.
\end{theorem}

\begin{proof}
The input reader requires terminal V26 metadata and recomputes the entire
digest (30.14) before transformation.  Theorem~\ref{thm:v27-inverse}
then identifies every inverse-transform entry with the corresponding
coefficient residue.  A Boolean mask is constructed directly from the
V19 support whose digest is (31.5); exhaustive comparison gives the four
counts in (31.7).  A second complete tensor transform with root \(4096\)
is compared entrywise against both input arrays and has no mismatch.
After atomic writes, an independent reload gives the same coefficient
digest (31.9), the same counts, and the same constant entry (31.8).
\end{proof}

\begin{remark}[Zeros at one prime]
\label{rem:v27-modular-zeros}
The \(3068\) candidate-support positions with zero Gaussian residue need
not have zero integer coefficient: a nonzero coefficient may be divisible
by \(p\).  By contrast, positions outside the candidate support are zero
already by the exact support algebra of V19.  V27 records a modular
census, not an integer sparsity theorem.
\end{remark}

\begin{record}[Stored coefficient files]
\label{rec:v27-files}
Both coefficient arrays have shape \((923521,)\), dtype
\(\mathrm{uint32}\), and \(3694212\) stored bytes.  Their complete file
hashes are
\[
\begin{split}
\text{real:}\quad&
\texttt{029F0B618D941DFB7DBDA3CCF841E5CA25657968DA1315ACD596C2319592F146},\\
\text{imaginary:}\quad&
\texttt{BFF89ADF76AD4BFD04395B06D0A64D3077C451DC5DC66D832F13E9C55CBE6DB9}.
\end{split}                                                    \tag{31.10}
\]
The JSON certificate has SHA--256
\[
\texttt{1C40599143E760A6430A6D9F8A990971DE06EB29090C596E266E548957F15AEF},
                                                               \tag{31.11}
\]
and its deterministic canonical payload has SHA--256
\[
\texttt{9973BF7A9003634ED1B6D022A23BC7AD37E548936DF31DF9167D61402F4211D3}.
                                                               \tag{31.12}
\]
\end{record}

\subsection{Why reflection is deliberately not tested spectrally}

Let \(\sigma(a+ib)=a-ib\) be Gaussian conjugation on
\(\mathbb F_p[i]\).  Since the transform root lies in the base field,
\(\sigma(\omega)=\omega\), not \(\omega^{-1}\).

\begin{proposition}[Finite-field conjugation does not reverse momentum]
\label{prop:v27-conjugation}
For a general Gaussian Laurent polynomial,
\[
\begin{aligned}
\sigma(S(r))&=\sum_\nu \sigma(c_\nu)\omega^{r\cdot\nu},\\
S(-r)&=\sum_\nu c_\nu\omega^{-r\cdot\nu}.
\end{aligned}                                                 \tag{31.13}
\]
These are equal only under an additional coefficient reflection identity
such as \(\sigma(c_\nu)=c_{-\nu}\).
\end{proposition}

\begin{proof}
Apply \(\sigma\) termwise to (31.1) and use
\(\sigma(\omega)=\omega\), which gives the first line.  Replacing \(r\)
by \(-r\) gives the second.  Relabeling \(\nu\mapsto-\nu\) shows the
stated sufficient identity and also displays the missing hypothesis in
the generic case.
\end{proof}

\begin{firewall}[Reflection pairing remains downstream]
\label{fire:v27-reflection}
The current frozen high-alias chart is not itself an already
reflection-paired global polynomial.  Therefore the \(258681\) nonzero
imaginary coefficient residues in (31.7) are not a failed reality test.
The physical reflection rule of Lemma~\ref{lem:v9-reflection} relates the
appropriate reflected external order after reconstruction.  Applying
absolute values, discarding imaginary limbs, or imposing
\(S(-r)=\sigma(S(r))\) at V27 would be an unjustified alteration of the
signed response.
\end{firewall}

\begin{record}[V27 verifier provenance]
\label{rec:v27-provenance}
The executable
\begin{center}
\path{scripts/verify_ym_postarchive_v27_inverse_ntt_first_prime.py}
\end{center}
has \(406\) physical lines, \(14614\) bytes, and SHA--256
\[
\texttt{80009D865E45D9B95551DB857DE60CBC62F6551830A046B1F733F243884BF068}.
                                                               \tag{31.14}
\]
\end{record}

\begin{decision}[V28: exact batch acceleration before fourteen primes]
\label{dec:v27-v28}
The first-prime coefficient route is correct, but repeating the
\(64\)-point execution partition at fourteen further primes would spend
substantial host time without adding mathematical information.  V28
should go upstream to the partition choice: prove that larger row batches
preserve every exact residue, retain the canonical \(64\)-point digest
framing, benchmark several safe batch sizes on identical point sets, and
select the fastest verified architecture.  It should then authorize the
remaining primes only under source-bound, disjoint-output checkpoints.
\end{decision}

\begin{assessment}[V27 acquisition]
\label{ass:v27}
V27 obtains the first complete Laurent coefficient-residue array and
closes transform orientation, aliasing, candidate-support containment,
full round trip, and stored-output integrity at the first prime.  The
observed constant entry is still only one residue pair.
\end{assessment}

\begin{record}[V27 append record]
\label{rec:v27-append}
V27 is appended to the clean V26 whose installed SHA--256 was
\[
\texttt{13DD7799A838F519C71B4C30A927761333A12E781546E7F72BD608C5A4019BDD}.
                                                               \tag{31.15}
\]
After validation only V27 is to remain the active numeric YM note.  The
next compulsory companion audit remains V30.
\end{record}

\begin{firewall}[Claim boundary after V27]
\label{fire:v27-final}
One-prime coefficient residues do not determine signed Gaussian integers
under the \(444\)-bit height bound.  No fifteen-prime CRT reconstruction,
reflection-paired box inequality, continuous Haar sign,
complement-sector estimate, compatible interacting measure, regulator
removal, Osterwalder--Schrader reconstruction, physical clustering, or
Hamiltonian Yang--Mills mass gap has been proved.
\end{firewall}

% =============================================================================
% END APPENDED POST-TENTH-ARCHIVE V27 MATERIAL
% =============================================================================
% =============================================================================
% BEGIN APPENDED POST-TENTH-ARCHIVE V28 MATERIAL
% =============================================================================

\section{V28: batch invariance and a negative acceleration gate}
\label{sec:v28}

V27 proposed increasing the number of spectral rows evaluated by each
kernel call before committing fourteen further primes.  V28 proves that
this change is algebraically safe, implements source-bound checkpoints
for it, and tests the performance premise.  The premise is false on the
present host: the original \(64\)-point partition is faster.

\subsection{Partition independence}

For a finite ordered point set \(P=(r_0,\ldots,r_{m-1})\), let
\({\cal K}_p(P)\) denote the two residue arrays returned by the exact V26
signed kernel.  A batch partition is a decomposition of \(P\) into
consecutive sublists.

\begin{theorem}[Exact row-batch invariance]
\label{thm:v28-batch-invariance}
For every listed prime and every two consecutive batch partitions of the
same ordered point set, concatenating the V26 outputs gives the same
array \({\cal K}_p(P)\).
\end{theorem}

\begin{proof}
Every phase-panel row depends only on its own grid point.  Rectangular
matrix multiplication contracts that row against fixed Laurent-bank
columns; subsequent \(45\times45\) Gaussian products and trace
contractions preserve the leading row index and never combine distinct
grid points.  Thus the mathematical output is rowwise.
Theorem~\ref{thm:v26-limb} proves that every binary64 limb contraction is
the exact integer contraction before modular reduction.  Hence a BLAS
partition cannot introduce a rounding dependence: each row in every
partition is the same exact field computation.  Concatenation in the
original order proves the claim.
\end{proof}

The V28 producer retains the V26 hash convention even for a larger
execution batch: it subdivides each output into consecutive \(64\)-point
microframes, hashing the real frame and then its imaginary frame.

\begin{corollary}[Canonical digest invariance]
\label{cor:v28-digest}
If two executions visit the same ordered points but use different batch
partitions whose boundaries are multiples of \(64\), their rolling
digests agree at every common \(64\)-point boundary.
\end{corollary}

\begin{proof}
Theorem~\ref{thm:v28-batch-invariance} gives identical residues in every
position.  The microframe rule presents the same bytes to SHA--256 in
the same order, independently of execution-batch boundaries.
\end{proof}

\begin{proposition}[V28 checkpoint binding]
\label{prop:v28-checkpoint}
A V28 checkpoint binds the schema, prime, root, grid geometry,
\(64\)-point microframe, execution-batch size, checkpoint frequency,
V28 source digest, and frozen V26 dependency digest.  Subject to the
filesystem assumptions of Proposition~\ref{prop:v26-resume}, a resumed
prefix preserves all committed residues and the canonical digest.
\end{proposition}

\begin{proof}
The write order and reopen checks are those of
Proposition~\ref{prop:v26-resume}, with the stated fields added to the
equality gate.  The endpoint must be execution-batch aligned unless it is
the terminal point.  The full prefix digest is recomputed in canonical
microframes before resumption.  Corollary~\ref{cor:v28-digest} shows that
this digest is the V26 convention rather than a new batch-dependent one.
\end{proof}

\subsection{Executed equivalence and resumption gates}

\begin{theorem}[V28 executed batch gate]
\label{thm:v28-gate}
On the first \(512\) points at \(p=2147483647\), partitions of sizes
\[
                         64,\ 128,\ 256,\ 512                \tag{32.1}
\]
all agree entrywise with the completed V26 artifact.  Their common
canonical prefix digest is
\[
\texttt{D3DA90A201DBAABC08A747C8F9D4DC531C222CB3D89AFA8AA1AE836B792BB65A}.
                                                               \tag{32.2}
\]
At \(p=2147483399\), one \(128\)-point evaluation agrees entrywise with
two consecutive \(64\)-point evaluations.  Finally, an uninterrupted
two-batch run of batch size \(256\) is byte-identical to a one-batch plus
resumed one-batch run through point \(512\), with digest (32.2).
\end{theorem}

\begin{proof}
The executable reloads the frozen V26 source and requires SHA--256
\[
\texttt{80BA1AE9FF855EE3DC8DB2138870DB3D7D5D720CCC69F01FEE001CA9ECC391D9}.
                                                               \tag{32.3}
\]
It also requires the complete first-prime artifact and its stream digest
(30.14).  It evaluates each partition independently and compares both
limbs entrywise with the first \(512\) stored residues.  The second-prime
comparison repeats the partition test after recompiling every rational
bank modulo the second prime.  The resume test uses two disjoint temporary
checkpoint directories and compares both stored prefixes and metadata.
All comparisons are exhaustive on their declared point sets.
\end{proof}

\begin{record}[Deterministic V28 certificate]
\label{rec:v28-certificate}
All six exact checks pass.  The deterministic canonical payload has
SHA--256
\[
\texttt{710EBC9E24561F64B885981F7BD8E1D5277C0DD29EB54181072B1ABB28BD49C8}.
                                                               \tag{32.4}
\]
Floating chart reconstruction diagnostics are labelled non-hashed and
are excluded from (32.4); this prevents platform-dependent last-bit
diagnostics from masquerading as canonical data.
\end{record}

\subsection{The performance premise fails}

\begin{record}[Equal-work host timing]
\label{rec:v28-timing}
On identical \(512\)-point inputs, one execution of each verified
partition gives
\[
\begin{array}{c|c|c}
\text{batch size}&\text{seconds}&\text{ratio to batch 64}\\ \hline
64 &7.915398700&1.000000\\
128&9.144364100&1.155262\\
256&9.254803300&1.169214\\
512&9.208491700&1.163362
\end{array}                                                    \tag{32.5}
\]
The timings are empirical, excluded from the canonical payload, and do
not claim a machine-independent ordering.
\end{record}

\begin{decision}[Reject larger batches on this host]
\label{dec:v28-reject}
No tested larger batch is authorized for the remaining-prime production.
The \(64\)-point partition is retained.  Since the per-prime row kernel
has not accelerated, V29 should move one level upstream and test
inter-prime concurrency: distinct primes have no shared mutable
mathematical state and can use disjoint source-bound output directories.
The gate must compare serial and concurrent fixed-prefix outputs and
measure aggregate rather than single-process throughput before launching
the fourteen remaining primes.
\end{decision}

\begin{record}[V28 verifier provenance]
\label{rec:v28-provenance}
The executable
\begin{center}
\path{scripts/verify_ym_postarchive_v28_batched_prime_producer.py}
\end{center}
has \(517\) physical lines, \(18171\) bytes, and SHA--256
\[
\texttt{EE5E3536B7CB0C251508005C7E1ACB1DB70218572034FC76B7A6F85FBACA2335}.
                                                               \tag{32.6}
\]
\end{record}

\begin{assessment}[V28 acquisition]
\label{ass:v28}
V28 proves execution-partition invariance and supplies a compatible
large-batch checkpoint engine, but its measured acceleration hypothesis
fails.  This negative result prevents a slower fourteen-prime run and
routes the programme upstream to independent-prime scheduling.
\end{assessment}

\begin{record}[V28 append record]
\label{rec:v28-append}
V28 is appended to the clean V27 whose installed SHA--256 was
\[
\texttt{3D80B53EDC2C565509BED36604346980C6A86F4BCC42A3A560C1368C7B685EE5}.
                                                               \tag{32.7}
\]
After validation only V28 is to remain the active numeric YM note.
\end{record}

\begin{firewall}[Claim boundary after V28]
\label{fire:v28-final}
Batch invariance and host timing do not compute a new prime.  Fourteen
spectral grids, their inverse transforms, the fifteen-prime signed CRT
lift, reflection-paired box inequality, continuous Haar sign,
complement-sector estimate, compatible interacting measure, regulator
removal, Osterwalder--Schrader reconstruction, physical clustering, and
Hamiltonian Yang--Mills mass gap remain open.
\end{firewall}

% =============================================================================
% END APPENDED POST-TENTH-ARCHIVE V28 MATERIAL
% =============================================================================
% =============================================================================
% BEGIN APPENDED POST-TENTH-ARCHIVE V29 MATERIAL
% =============================================================================

\section{V29: exact inter-prime concurrency gate}
\label{sec:v29}

V28 showed that increasing the row-batch size does not accelerate the
present host.  V29 therefore moves upstream from row partitioning to prime
scheduling.  It tests whether two complete standalone producer processes
can run concurrently without changing either prime's exact prefix.

\subsection{Mathematical independence of prime jobs}

For prime index \(j\), let \({\cal P}_j(n)\) be the two V26 residue arrays
through flat point \(n\), together with their source-bound metadata and
canonical rolling digest.

\begin{theorem}[Disjoint prime-job independence]
\label{thm:v29-independence}
Suppose two V26 producer processes use distinct listed primes and disjoint
output directories.  For each process, the committed object
\({\cal P}_j(n)\) is independent of whether the other process is scheduled
serially or concurrently.
\end{theorem}

\begin{proof}
Each process recompiles the immutable exact rational banks modulo its own
prime and constructs its own root-power table.  Its field operations,
spectral points, arrays, metadata, digest state, and temporary metadata
path are process-local.  The only common inputs are read-only source
files.  There is no random input, shared cache that enters a residue
formula, or shared writable output.  By Theorem~\ref{thm:v26-limb}, every
binary64 contraction returns its exact integer value before modular
reduction; thread or process scheduling therefore cannot change a rounded
residue.  The checkpoint digest is a deterministic function of those
residues.  Hence concurrency can affect elapsed time but not
\({\cal P}_j(n)\).
\end{proof}

\subsection{Executed serial-versus-concurrent comparison}

The gate launches standalone V26 processes at prime indices \(2\) and
\(3\), each through \(64\) tiles:
\[
                         n=64\cdot64=4096.                   \tag{33.1}
\]
It first runs the two jobs serially in separate temporary directories and
then runs identical jobs concurrently in two further directories.

\begin{theorem}[Byte-identical concurrent prefixes]
\label{thm:v29-gate}
Every child exits with code zero.  At prime \(2147483399\), root
\(450851339\), the serial and concurrent real and imaginary prefixes are
entrywise identical and have common rolling digest
\[
\texttt{BF97E435EDD98189D5F9BBD0C8ACC0B551631DE41439A924CC3E192AE2AE7D99}.
                                                               \tag{33.2}
\]
At prime \(2147480299\), root \(1165718630\), the corresponding prefixes
are entrywise identical and have common rolling digest
\[
\texttt{33D45C03BDFADF54E453D3855DAA8CCE2D8474E1C7D701B2FFD68FE523EBF36D}.
                                                               \tag{33.3}
\]
\end{theorem}

\begin{proof}
For each index the verifier reloads both serial arrays and both concurrent
arrays, checks shape, dtype, prime, root, and endpoint \(4096\), and
compares all \(8192\) stored limb residues.  It separately compares the
rolling digests.  All comparisons pass.  Theorem~\ref{thm:v29-independence}
explains why this executed equality is the required scheduling invariant
rather than an accidental analytic approximation.
\end{proof}

\begin{record}[Canonical concurrency certificate]
\label{rec:v29-certificate}
The deterministic payload binds the frozen V26 dependency hash (32.3),
the two prime indices, prefix geometry, disjoint-directory condition,
both prime/root pairs, both digests (33.2)--(33.3), and all exact checks.
Its SHA--256 is
\[
\texttt{676FD50878694C0422F787EAF0390728B11FEB50321F41AF57FDEA237F02D8CD}.
                                                               \tag{33.4}
\]
\end{record}

\subsection{Aggregate timing and scheduling decision}

\begin{record}[Serial and concurrent host timing]
\label{rec:v29-timing}
The two serial child durations are
\[
           145.944283000\ {\rm s},\qquad141.336580500\ {\rm s},
\]
with aggregate wall time \(287.281489700\) seconds.  The two concurrent
child durations are
\[
           144.883670300\ {\rm s},\qquad145.224777500\ {\rm s},
\]
with aggregate wall time \(145.226377600\) seconds.  Thus the measured
aggregate speedup is
\[
             \frac{287.281489700}{145.226377600}
             =1.978163295.                                  \tag{33.5}
\]
All timing fields and the decision derived from them are excluded from
the canonical payload.
\end{record}

\begin{decision}[Authorize concurrency two]
\label{dec:v29-concurrency}
The exact equality gate passes and the measured speedup exceeds the
declared \(1.10\) threshold.  The remaining prime producer may therefore
run two standalone processes concurrently, using the unchanged
\(64\)-point V26 kernel and one disjoint source-bound directory per prime.
No higher concurrency is authorized by V29.
\end{decision}

\begin{record}[V29 verifier provenance]
\label{rec:v29-provenance}
The executable
\begin{center}
\path{scripts/verify_ym_postarchive_v29_prime_concurrency_gate.py}
\end{center}
has \(220\) physical lines, \(7077\) bytes, and SHA--256
\[
\texttt{2DD10D17A64A19B39E7B15B21950365D46959EFF1CA16C538E344CCAC7885729}.
                                                               \tag{33.6}
\]
\end{record}

\begin{decision}[V30 audit before production]
\label{dec:v29-v30}
V30 is the compulsory fifth-version audit.  It must inspect the then-current
SM and NS notes before launching long-lived prime jobs.  If those notes do
not supply an exact shortcut for the finite coefficient problem, V31
should start source-bound prime indices \(2\) and \(3\) concurrently under
Decision~\ref{dec:v29-concurrency}.
\end{decision}

\begin{assessment}[V29 acquisition]
\label{ass:v29}
V29 proves and executes scheduling invariance for two independent primes
and nearly halves aggregate fixed-prefix wall time.  It changes the
authorized acquisition schedule, not the underlying polynomial.
\end{assessment}

\begin{record}[V29 append record]
\label{rec:v29-append}
V29 is appended to the clean V28 whose installed SHA--256 was
\[
\texttt{D049E0925E055ADBE0F05A815187404F3962EA865618C566B3BAA20BF4018B58}.
                                                               \tag{33.7}
\]
After validation only V29 is to remain the active numeric YM note.
\end{record}

\begin{firewall}[Claim boundary after V29]
\label{fire:v29-final}
The scheduling gate does not complete either tested prime.  Fourteen
spectral grids, their inverse transforms, the fifteen-prime signed CRT
lift, reflection-paired box inequality, continuous Haar sign,
complement-sector estimate, compatible interacting measure, regulator
removal, Osterwalder--Schrader reconstruction, physical clustering, and
Hamiltonian Yang--Mills mass gap remain open.
\end{firewall}

% =============================================================================
% END APPENDED POST-TENTH-ARCHIVE V29 MATERIAL
% =============================================================================
% =============================================================================
% BEGIN APPENDED POST-TENTH-ARCHIVE V30 MATERIAL
% =============================================================================

\section{Mandatory companion audit at V30}
\label{sec:v30}

\subsection{Active sources inspected}

The active numeric companion notes in \path{latex_notes} were inspected
read-only after the V29 concurrency gate:
\[
\begin{array}{c|r|r}
\text{source}&\text{lines}&\text{bytes}\\ \hline
\text{SM V29}&6700&245282\\
\text{NS V26}&5319&278283
\end{array}                                                     \tag{34.1}
\]
Their respective SHA--256 values are
\[
\texttt{02D53AB3BD8EBA36768EC83D1DA6740207D565714F1B9CA22D101628882C9870}
\]
and
\[
\texttt{82C887F8C2C69E4F28FAD7C3DC8DE5B9099CB5A4BF431EBA1FFF3E91935978AD}.
\]
The SM programme has rolled over since the V25 YM audit; archived
sources, including its \(\texttt{V157\_f1}\), were not mistaken for the
active compact source.  The NS source is unchanged.

\subsection{SM V29}

SM V29 proves existence and uniqueness of a positive solution to the
flat-torus CMC Lichnerowicz equation
\[
 -8\Delta\psi-|\sigma|^2\psi^{-7}
       +\frac23\tau^2\psi^5=0                               \tag{34.2}
\]
for every smooth nonzero transverse-traceless seed \(\sigma\) and
nonzero constant mean curvature \(\tau\).  A positive solution of a
linear elliptic equation supplies a lower barrier even where
\(|\sigma|\) vanishes.  Its own firewall states that this is a classical
constraint solution, not a quantum measure, an interacting transfer
operator, or reflection-positive continuum construction.

\begin{proposition}[SM consequence for the YM route]
\label{prop:v30-sm}
Equation (34.2) supplies no coefficient identity, modulus, transform, or
runtime shortcut for the finite Yang--Mills CRT problem.  It is retained
only as a possible later gravitational constraint interface.
\end{proposition}

\begin{proof}
The current YM objects are residues of a fixed Gaussian Laurent
polynomial at finite-field roots.  The variables and operators in (34.2)
are a positive real conformal factor, a smooth TT tensor, and an elliptic
Laplacian.  There is no map in SM V29 from these data to the V24 prime
list or the V27 coefficient array.  Conversely the one-prime residue
array proves no hypothesis of the CMC theorem.  Therefore importing any
constant or conclusion would be type-incorrect.
\end{proof}

SM V25 independently read the completed YM V26--V27 steps and reached
the same no-import decision.  It correctly records the V27 source hash
\[
\texttt{3D80B53EDC2C565509BED36604346980C6A86F4BCC42A3A560C1368C7B685EE5}
                                                               \tag{34.3}
\]
and byte count \(385758\), but lists \(9119\) lines.  The source with
(34.3) has \(10020\) physical lines.  This companion metadata discrepancy
does not enter any YM theorem; the hash and mathematical reading are
consistent.

\subsection{NS V26}

NS V26 still gives closed pointwise rank-one norms for the first and
second derivatives of the Oseen kernel.  Its first-derivative
restriction gives no improvement, and its second-derivative improvement
holds only over a bounded radial window.

\begin{proposition}[No NS import]
\label{prop:v30-ns}
The NS rank-one norm identities neither accelerate nor alter the V29
inter-prime producer schedule.
\end{proposition}

\begin{proof}
Those identities contract spatial derivatives of a real Oseen kernel
against rank-one velocity tensors.  The YM producer performs exact
Gaussian finite-field matrix and trace contractions at independent
primes.  The two computations share no operator, norm, variable, or
hypothesis.  The unchanged NS source therefore supplies no new interface
to the current bottleneck.
\end{proof}

\begin{theorem}[V30 audit decision]
\label{thm:v30-decision}
No theorem, numerical constant, or executable artifact is imported from
SM V29 or NS V26.  The exact concurrency-two acquisition order of
Decision~\ref{dec:v29-concurrency} is unchanged.  V31 should launch
complete prime indices \(2\) and \(3\) concurrently, with the frozen V26
source and disjoint checkpoint directories.
\end{theorem}

\begin{proof}
Propositions~\ref{prop:v30-sm} and \ref{prop:v30-ns} exclude a valid
finite-certificate import.  The V29 gate already proves byte identity
under concurrent scheduling and measures a \(1.978163295\) aggregate
speedup.  No audited result invalidates its source, arithmetic, or
disjoint-output hypotheses.  Hence the stated acquisition remains the
strongest certified next action.
\end{proof}

\begin{assessment}[V30 acquisition]
\label{ass:v30}
The compulsory audit is complete.  It preserves the SM CMC theorem for
a later coupled constraint stage, records one harmless companion
metadata defect, and leaves the finite exact CRT route unchanged.
\end{assessment}

\begin{record}[V30 append record]
\label{rec:v30-append}
V30 is appended to the clean V29 whose installed SHA--256 was
\[
\texttt{F10A2A9A585B2C5DFDD33522B4722DC3D36BD70A5078391D42A2CC90A3EC595B}.
                                                               \tag{34.4}
\]
After validation only V30 is to remain the active numeric YM note.  The
next compulsory companion audit is V35.
\end{record}

\begin{firewall}[Claim boundary after V30]
\label{fire:v30-final}
The companion audit does not compute a new prime or import continuum
closure.  Fourteen spectral grids, their inverse transforms, the
fifteen-prime signed CRT lift, reflection-paired box inequality,
continuous Haar sign, complement-sector estimate, compatible
interacting measure, regulator removal, Osterwalder--Schrader
reconstruction, physical clustering, and Hamiltonian Yang--Mills mass
gap remain open.
\end{firewall}

% =============================================================================
% END APPENDED POST-TENTH-ARCHIVE V30 MATERIAL
% =============================================================================

% =============================================================================
% BEGIN APPENDED POST-TENTH-ARCHIVE V31 MATERIAL
% =============================================================================

\section{V31: exact scope of the frozen-chart artifact, the constant-class
refutation theorem, and a second-order complex strip}
\label{sec:v31}

V30 authorized concurrent production of the spectral grids at prime
indices \(2\) and \(3\).  Before the remaining fourteen prime runs are
spent, V31 goes upstream to the question that the V11--V30 chain has not
asked explicitly: what would the completed fifteen-prime constant
coefficient certify?  The answer, proved exactly below, is that the
frozen order-five polynomial satisfies its own Neumann hypothesis at
exactly one of the \(256\) quarter-turn momenta, so its Haar mean has no
certified relation to \(\mathfrak m_{\rm G}\), and the box-uniform V8
architecture at the V11 parameters needs more than \(6.69\times10^{12}\)
boxes.  V31 therefore withdraws the production route as a Haar-sign
instrument.  It then proves, from the archived SCIC definitions alone,
that a certified nonzero constant class would \emph{refute} SCIC for the
Wilson--endpoint pair rather than close it, and that the archived reset
ledger cannot absorb such a class under ADMC as stated.  Finally it
proves a second-order Hermitian-part strip theorem which sharpens the
archived V53 polystrip by a certified factor above \(190\), and shows
that even this gain leaves the absolute alias route infeasible; the
binding obstruction is the supremum stock, not the strip.

\subsection{Partial production record}

\begin{record}[Frozen prime-two and prime-three prefixes]
\label{rec:v31-prefixes}
When V31 began, the directories
\path{artifacts/ym_v31_prime_02} and \path{artifacts/ym_v31_prime_03}
contained resumable V26 prefixes for the primes \(2147483399\) (root
\(450851339\)) and \(2147480299\) (root \(1165718630\)), with
\[
 \texttt{next\_flat\_point}=438272
 \quad\hbox{and}\quad
 \texttt{next\_flat\_point}=424960,                          \tag{35.1}
\]
producer digest (32.3), and rolling digests
\begin{center}
\scriptsize
\texttt{ACE0998DFC327038271DDDD6B9B98D59B5CD5BF7C0DEFEA245E656B24C2D1364},
\par\medskip
\texttt{E3963713AF183378259D1989CD75740E0CFE6F538DFA2BF549C867C0C181AE4E}.
\end{center}
Both jobs were started at 20:57 and last checkpointed at 22:28 local
time on 11 September 2026; no producer process was running when V31
began.  The prefixes are retained under the V26 checkpoint invariant
and are not resumed by this note.  This record is empirical and enters
no theorem.
\end{record}

\subsection{Exact global census of the frozen chart}

Recall the frozen data of V10 and V12: the rational chart
\(Z_\sharp\) reconstructed at the anchor \(a=(\pi,\pi,0,0)\), the
Hermitian Gaussian-rational trial inverses \(R_e\) of
\(N_e(a)=Z_\sharp^*H_e(a)Z_\sharp\), and the defect
\[
 E_e(k):=I-R_eZ_\sharp^*H_e(k)Z_\sharp.                       \tag{35.2}
\]
The polynomial whose coefficients V26--V27 reconstruct replaces
\(N_e(k)^{-1}\) by \(\sum_{n=0}^4E_e(k)^nR_e\).  Every theorem of
V9--V30 which relates that polynomial to the response
\(F_e\) assumes \(\|E_e(k)\|_{\rm F}<1\) on the region considered.

\begin{theorem}[Exact quarter-turn defect census]
\label{thm:v31-census}
For both actions and every \(k\in\Gamma_4\), the quantity
\(\|E_e(k)\|_{\rm F}^2\) is a nonnegative rational number computable in
Gaussian-rational arithmetic.  Its exact values satisfy
\[
\begin{array}{c|c|c|c|c}
 e&\hbox{anchor}&\hbox{least other point}&\hbox{toron }k=0&\hbox{largest}\\ \hline
 {\rm W}&2.3552\times10^{-12}&92.1516&423.0852&1559.3810\\
 *&7.9737\times10^{-10}&115.9973&515.1875&2683.1331 .
\end{array}                                                    \tag{35.3}
\]
Consequently the Neumann condition \(\|E_e(k)\|_{\rm F}<1\) holds at
exactly one of the \(256\) quarter-turn momenta, the anchor, and fails
at the other \(255\) with
\[
 \|E_{\rm W}(k)\|_{\rm F}\ge9.599,\qquad
 \|E_*(k)\|_{\rm F}\ge10.770\qquad(k\in\Gamma_4\setminus\{a\}).
                                                                 \tag{35.4}
\]
\end{theorem}

\begin{proof}
At a quarter-turn point every Laurent phase lies in
\(\{1,i,-1,-i\}\), so \(H_e(k)\) is a Gaussian-rational matrix.  The
verifier forms \(Z_\sharp^*H_e(k)Z_\sharp\) and \(R_e(\cdot)\) by exact
integer matrix products with common denominators, subtracts the
identity, and sums the squares of all real and imaginary numerators.
The anchor values are the squares of the outward residuals recorded in
(14.9): their square roots are \(1.5346661\times10^{-6}\) and
\(2.8237743\times10^{-5}\), consistent with the outward V10 values
\(1.534667\times10^{-6}\) and \(2.8237744\times10^{-5}\).  Exact
comparison with \(1\) at all \(256\) points gives the counts; the
decimals in (35.3) are images of the stored fractions, and (35.4)
follows from \(92<9.599^2\) and \(115<10.770^2\).
\end{proof}

\begin{corollary}[The V27 constant coefficient is not a Haar-mean
approximant]
\label{cor:v31-not-approximant}
The identity \(\langle\Phi^{(5)}_*-\Phi^{(5)}_{\rm W}\rangle=c_0/\Delta\)
is exact, but no theorem of V9--V30 relates
\(\langle\Phi^{(5)}_*-\Phi^{(5)}_{\rm W}\rangle\) to
\(\mathfrak m_{\rm G}\): the tail theorem
Theorem~\ref{thm:v9-neumann-tail} is inapplicable at every quarter-turn
momentum except the anchor, and its transported form
Theorem~\ref{thm:v10-frozen-transport} certifies only the box of
half-width \(2^{-9}\) selected in (15.7).
\end{corollary}

\begin{proof}
Theorem~\ref{thm:v31-census} shows that the hypothesis
\(\eta_{e,B}<1\) of Theorem~\ref{thm:v9-neumann-tail} fails on every
cell containing a non-anchor quarter-turn point.  The only certified
region is the V11 box.
\end{proof}

\begin{theorem}[Exact uniform box count]
\label{thm:v31-box-count}
Any family of closed cubes of side \(2^{-8}\) covering
\([-\pi,\pi]^4\) has at least \((512\pi)^4\) members.  With
\(\pi>3.1415\),
\[
 (512\pi)^4>\left(\frac{201056}{125}\right)^4
 >6\,693\,112\,119\,092.                                        \tag{35.5}
\]
At the executed first-prime cost of Record~\ref{rec:v26-provenance},
namely \(15\times6574.290507600\) seconds per box, the box-uniform
certificate of Theorem~\ref{thm:v8-finite-certificate} at the V11
parameters requires more than
\[
 \frac{82504619007841980620811}{125000}\ {\rm s}
 >1.83\times10^{14}\ {\rm host\ hours}.                           \tag{35.6}
\]
\end{theorem}

\begin{proof}
The cubes have total volume at least the volume \((2\pi)^4\) of the
cube they cover, so their number is at least
\((2\pi)^4/2^{-32}=(512\pi)^4\).  Insert \(\pi>3.1415\) and multiply
exactly; the verifier records the integer part of the rational in
(35.5) and the product (35.6).
\end{proof}

\begin{decision}[Withdrawal of the production route]
\label{dec:v31-withdraw}
The acquisition authorized in Theorem~\ref{thm:v30-decision} is
withdrawn as a Haar-sign instrument.  The completed first-prime
coefficient array of Theorem~\ref{thm:v27-coefficients} remains a
verified datum for the single V11 box; the prefixes of
Record~\ref{rec:v31-prefixes} are frozen and not resumed.  No further
prime is to be computed unless a theorem first relates a frozen-chart
polynomial to \(\mathfrak m_{\rm G}\) on a set of positive Haar measure
at feasible box count.
\end{decision}

\subsection{What a certified sign would decide}

The ninth archive defines the measured comparison cochain and SCIC by
\[
 \gamma_j^{\rm cmp}-\gamma_j^{\rm cov}
 =A_{j+1}-A_j+E_j+C_j,\qquad
 \mathfrak b(C_j)=0,\qquad
 \mathfrak b(A_N)=o(N),\qquad
 \frac1N\sum_{j<N}\mathfrak b(E_j)\longrightarrow0,             \tag{35.7}
\]
proves that SCIC together with the covariant hypothesis
\(\mathfrak b\gamma_j^{\rm cov}=b_{\rm YM}+e_j^{\rm cov}\),
\(N^{-1}\sum_{j<N}e_j^{\rm cov}\to0\), implies
\(N^{-1}\sum_{j<N}\mathfrak b\gamma_j^{\rm cmp}\to b_{\rm YM}\), and
proves that the edge row of SCIC holds if and only if the Cesaro mean
\(\overline m_N\) of the homotopy scalars \(\int_0^1m_{j,t}\,dt\) tends
to zero.  V2 proved that, with scale-independent rescaled endpoint
banks, \(\int_0^1m_{j,t}\,dt\) is the same number at every scale; call
it \(\sigma\mathfrak m_{\rm G}\), where \(\sigma\in\{1,-1\}\) is fixed by
the orientation of the archived homotopy relative to the
endpoint-minus-Wilson convention (6.4).

\begin{theorem}[A constant class refutes SCIC]
\label{thm:v31-refutation}
Assume the hypotheses of Theorem~\ref{thm:v2-endpoint-reduction} and
the archived uniform analytic stock (17.14).  Then:
\begin{enumerate}[label=\textup{(\roman*)},leftmargin=3.3em]
\item \(\overline m_N=\sigma\mathfrak m_{\rm G}\) for every \(N\);
\item the edge row of SCIC holds if and only if
      \(\mathfrak m_{\rm G}=0\);
\item if \(\mathfrak m_{\rm G}\ne0\) and the covariant hypothesis holds,
      then
      \[
       \frac1N\sum_{j<N}\mathfrak b\gamma_j^{\rm cmp}
       \longrightarrow b_{\rm YM}+\sigma\mathfrak m_{\rm G},         \tag{35.8}
      \]
      so the compact Gaussian half (14.16) of ADMC fails by the linear
      term \(\sigma\mathfrak m_{\rm G}N\);
\item any ledger rows absorbing that discrepancy in (14.10) must total
      \(-\sigma\mathfrak m_{\rm G}N+o(N)\), which violates the sublinear
      reset row (14.17).  Hence ADMC cannot hold for this tower: either
      the clock limit is \(b_{\rm YM}+\sigma\mathfrak m_{\rm G}\) or
      (14.17) fails.
\end{enumerate}
\end{theorem}

\begin{proof}
(i) Theorem~\ref{thm:v2-endpoint-reduction} gives
\(\int_0^1m_{j,t}\,dt=\sigma\mathfrak m_{\rm G}\) for each \(j\); the
Cesaro mean of a constant sequence is that constant.
(ii) This is the archived edge criterion applied to (i).
(iii) Apply \(\mathfrak b\) to (35.7) and sum.  The contacts vanish.
By the archived analytic alias cohomology, the nonconstant part of each
rescaled edge scalar is an alias coboundary with uniformly bounded
potential, hence contributes to \(A_j\) with
\(\sup_N|\mathfrak b(A_N)|<\infty\), and the constant part contributes
\(\sigma\mathfrak m_{\rm G}\) per scale.  Therefore
\(N^{-1}\sum_{j<N}\mathfrak b(\gamma_j^{\rm cmp}-\gamma_j^{\rm cov})
\to\sigma\mathfrak m_{\rm G}\), and the covariant hypothesis gives
(35.8).  The identification of \(\sum_{j<N}\mathfrak b\gamma_j^{\rm cmp}\)
with \(\mathfrak b\Gamma_{\rm G}^{[0,N)}\) is the archived determinant
cocycle (14.8).
(iv) In the ledger (14.10) the nonlinear rows are Cesaro-null by the
archived estimate (14.14)--(14.15), so
\(\beta_N/N\to b_{\rm YM}\) forces
\(R_N^{\rm reset}+R_N^{\rm bdry}=-\sigma\mathfrak m_{\rm G}N+o(N)\),
which is not \(o(N)\).
\end{proof}

\begin{corollary}[A sign certificate is a refutation instrument]
\label{cor:v31-refutation-instrument}
Under the hypotheses of Theorem~\ref{thm:v31-refutation}, a rigorous
certificate \(\langle F_*-F_{\rm W}\rangle<0\) would prove that the
Wilson--endpoint pair does not satisfy SCIC and that the self-similar
reset ledger does not satisfy ADMC.  It cannot close either card.
\end{corollary}

\begin{proof}
Negativity implies \(\mathfrak m_{\rm G}\ne0\); apply (ii) and (iv).
\end{proof}

\begin{remark}[No conflict with one-loop universality]
\label{rem:v31-universality}
Universality of the one-loop coefficient is a statement about the
per-shell coefficients \(b_{{\rm G},j}\) of a \emph{nested} Gaussian
tower, the successive Schur complements of one fine quadratic form in
(14.7), whose \(j\)-th shell uses the \(j\)-times blocked kernel.  The
scale-independence hypothesis of Theorem~\ref{thm:v2-endpoint-reduction}
describes instead a tower reset to the bare banks at every scale.  In
such a tower the first-shell finite constant of each action recurs at
every scale, and the difference of first-shell constants between two
actions is a finite scheme-type number which need not vanish.  The
numerical diagnostics of V2--V3, namely the certified \(M=4\) interval
(6.21) and the floating values \(-318.12\), \(-314.84\), \(-324.29\) at
\(M=3,4,5\), make \(\mathfrak m_{\rm G}=0\) implausible, although this
is not a proof.
\end{remark}

\begin{proposition}[Nested-tower coboundary]
\label{prop:v31-nested}
Suppose the per-shell Gaussian coefficients of the nested tower have
the form
\[
 b_{{\rm G},j}=b_{\rm YM}+c_j-c_{j+1}+e_j,
 \qquad \sup_j|c_j|<\infty,\qquad
 \frac1N\sum_{j<N}e_j\longrightarrow0.                         \tag{35.9}
\]
Then (14.16) holds with remainder \(c_0-c_N+o(N)=o(N)\).  The
hypotheses (35.9) and the scale-independence hypothesis of
Theorem~\ref{thm:v2-endpoint-reduction} are compatible only if
\(\mathfrak m_{\rm G}=0\).
\end{proposition}

\begin{proof}
Summation telescopes the coboundary and the Cesaro assumption removes
\(\sum_{j<N}e_j\).  If both hypotheses held, the constant class per
scale would be simultaneously \(\sigma\mathfrak m_{\rm G}\) and the
Cesaro-null quantity \(c_j-c_{j+1}+e_j\), which forces
\(\mathfrak m_{\rm G}=0\) by (i) of Theorem~\ref{thm:v31-refutation}.
\end{proof}

\begin{assessment}[Consequence for the acquisition order]
\label{ass:v31-order}
The finite sign problem is not the Gaussian anomaly.  Whether the
parity-cell tower satisfies (35.9), with \(c_j\) the finite constant of
the \(j\)-times blocked kernel, is the actual open Gaussian statement,
and it is decided at the level of nested determinants before any cubic
or quartic vertex is used.  The V2--V30 quantity
\(\mathfrak m_{\rm G}\) is the first-shell constant difference between
two actions; its exact sign is of no use to (14.16) and, if nonzero,
only certifies that the self-similar comparison must be abandoned.
\end{assessment}

\subsection{A second-order Hermitian-part strip}

The archived strips (18.8) and (53.11) perturb from the real floor by
the first-order quantity \(e^{Rr}-1\).  The following identity shows
that the Hermitian part of the complexified Hessian moves only at
second order.

\begin{lemma}[Hermitian part of a complexified Laurent Hessian]
\label{lem:v31-cosh}
Let \(H(k)=\sum_{\nu}C_\nu e^{i\nu\cdot k}\) be a finite Laurent matrix
which is Hermitian for all real \(k\), so that \(C_{-\nu}=C_\nu^*\).
For \(z=k+iy\) with real \(k,y\), put
\(H(z)=\sum_\nu C_\nu e^{i\nu\cdot k}e^{-\nu\cdot y}\).  Then
\[
 \frac12\bigl(H(z)+H(z)^*\bigr)
 =\sum_\nu C_\nu e^{i\nu\cdot k}\cosh(\nu\cdot y).               \tag{35.10}
\]
\end{lemma}

\begin{proof}
The adjoint of the \(\nu\)-th term is
\(C_\nu^*e^{-i\nu\cdot k}e^{-\nu\cdot y}=C_{-\nu}e^{-i\nu\cdot k}
e^{-\nu\cdot y}\).  Relabel \(\nu\mapsto-\nu\) in the adjoint sum and
add: the \(\nu\)-th term becomes
\(C_\nu e^{i\nu\cdot k}(e^{-\nu\cdot y}+e^{\nu\cdot y})/2\).
\end{proof}

For a Laurent matrix \(P\) recall the induced stocks (53.4):
\({\cal R}(P)\), \({\cal C}(P)\), and the range \(R(P)\).

\begin{theorem}[Second-order strip]
\label{thm:v31-strip}
Let \(H\) be as in Lemma~\ref{lem:v31-cosh} with \(H(k)\succeq\theta I\)
on the real torus.  If \(\|y\|_\infty\le r\), then
\[
 \left\|\tfrac12\bigl(H(z)+H(z)^*\bigr)-H(k)\right\|_2
 \le\sqrt{{\cal R}(H){\cal C}(H)}\,\bigl(\cosh(R(H)r)-1\bigr).   \tag{35.11}
\]
If \(\gamma:=\theta-\sqrt{{\cal R}{\cal C}}(\cosh(Rr)-1)>0\), then
\(H(z)\) is invertible on the strip and
\(\|H(z)^{-1}\|_2\le\gamma^{-1}\).
\end{theorem}

\begin{proof}
By (35.10) the difference is
\(\sum_\nu C_\nu e^{i\nu\cdot k}(\cosh(\nu\cdot y)-1)\).  Since
\(\cosh\) is even and increasing on \([0,\infty)\) and
\(|\nu\cdot y|\le R r\), every weight lies in \([0,\cosh(Rr)-1]\).
Row and column absolute sums are therefore bounded by
\({\cal R}(\cosh(Rr)-1)\) and \({\cal C}(\cosh(Rr)-1)\), and
\(\|M\|_2\le\sqrt{\|M\|_1\|M\|_\infty}\) gives (35.11).  If the
Hermitian part of \(H(z)\) is at least \(\gamma I\), then for every
vector \(v\),
\(\|H(z)v\|\,\|v\|\ge\operatorname{Re}\langle v,H(z)v\rangle
\ge\gamma\|v\|^2\), which proves invertibility and the inverse bound.
\end{proof}

\begin{lemma}[Rational majorant of \(\cosh t-1\)]
\label{lem:v31-cosh-bound}
For \(t^2<56\),
\[
 \cosh t-1\le\frac{t^2}{2}+\frac{t^4}{24}
 +\frac{t^6}{720\,(1-t^2/56)}.                                   \tag{35.12}
\]
\end{lemma}

\begin{proof}
For \(n\ge4\), \((2n)!/(2n-2)!=2n(2n-1)\ge56\), so the tail
\(\sum_{n\ge3}t^{2n}/(2n)!\) is at most
\((t^6/720)\sum_{m\ge0}(t^2/56)^m\).
\end{proof}

\begin{theorem}[Certified second-order radii]
\label{thm:v31-radii}
For the unscaled family \(B_t=D_R^*D_R+3tK_R^*K_R\) of (53.2), with the
exact stocks \({\cal R}={\cal C}\le1066\), \(R\le3\) and real floor
\(147/1250\), the radius
\[
 \boxed{\ \rho_{\rm c}:=\frac{3501}{1000000}\ }                  \tag{35.13}
\]
satisfies \(1066(\cosh(3\rho_{\rm c})-1)\le147/2500\).  Hence for all
\(0\le t\le1\) and \(\|y\|_\infty\le\rho_{\rm c}\),
\(A_t(k+iy)=B_t(k+iy)/3\) is invertible with
\(\|A_t(k+iy)^{-1}\|_2\le2500/49\).  For the Wilson endpoint alone,
with stocks \(23\) and range \(2\), the same conclusion holds on
\[
 \boxed{\ \rho_{\rm c}^{\rm W}:=\frac{7149}{200000}.\ }          \tag{35.14}
\]
Relative to the archived radii (53.11) and (53.22),
\[
 \frac{\rho_{\rm c}}{\rho_{53}}\ge\frac{1866033}{9800}>190.41,
 \qquad
 \frac{\rho_{\rm c}^{\rm W}}{\rho_{53}^{\rm W}}
 \ge\frac{54809}{1960}>27.96.                                    \tag{35.15}
\]
\end{theorem}

\begin{proof}
The verifier regenerates \(B_0,B_1\) from the V21 incidence rules,
reproduces the stocks \(23,1066\) and ranges \(2,3\) of (53.8), and
uses the archived endpoint convexity (53.5)--(53.7) for
intermediate \(t\).  It selects the largest \(r\) with denominator
\(10^6\) for which the rational majorant (35.12) certifies the
perturbation budget \(147/2500\), and rechecks that inequality
exactly.  Theorem~\ref{thm:v31-strip} with
\(\theta=147/1250\) gives \(\gamma=147/2500\) and hence
\(\|B_t^{-1}\|_2\le2500/147\).  For (35.15) use
\(\log(1+x)<x\): \(\rho_{53}<49/2665000\) and
\(\rho_{53}^{\rm W}<147/115000\); the displayed fractions are the exact
quotients of (35.13)--(35.14) by these bounds.
\end{proof}

\begin{countertheorem}[The sharper strip does not make the absolute
alias route feasible]
\label{counter:v31-mesh}
Under the archived allocation \(\rho_W=\rho/8\) and the alias
multiplier \(C_M(\rho)=\coth^4(\rho M/2)-1\), the necessary condition
\(C_M(\rho_W)<1\) requires, with \(x_*=\operatorname{artanh}(2^{-1/4})
>6121/5000\),
\[
 M>\frac{16x_*}{\rho_{\rm c}}>\frac{19587200}{3501},
 \qquad\hbox{hence}\qquad M\ge5595,                              \tag{35.16}
\]
for the common family, and \(M\ge548\) for the Wilson-only radius.
The common threshold alone exceeds \(9.79\times10^{14}\) grid points,
before the unpopulated Wiener stock \(N_W\) of (49.4) enters the
enclosure.
\end{countertheorem}

\begin{proof}
The archived threshold lemma (54.4) gives
\(C_M<1\iff\rho_WM/2>x_*\).  Substitute \(\rho_W=\rho_{\rm c}/8\) and the
rational lower bound for \(x_*\); the verifier records the exact
quotients and their integer parts, and \(5595^4\) exceeds the stated
count.
\end{proof}

\begin{record}[Sampled Hermitian-part floors, non-hashed]
\label{rec:v31-diagnostic}
On the \(10^4\) grid of the real torus, the smallest eigenvalue of the
Hermitian part (35.10) evaluated at \(y=r(1,1,1,1)\) has sampled minima
\[
\begin{array}{c|cccccccc}
 r&0&0.01&0.02&0.05&0.1&0.2&0.3&0.5\\ \hline
 B_0&0.1176&0.1176&0.1175&0.1169&0.1149&0.1043&0.0661&-0.156\\
 B_1&2.678&2.677&2.674&2.647&2.532&2.006&0.953&-4.83 ,
\end{array}                                                    \tag{35.17}
\]
and at \(y=r(1,0,0,0)\) the \(B_1\) minima remain positive through
\(r=0.5\), where the value is \(0.108\).  The sampled real floor of
\(B_0\) equals \(147/1250\) to four digits and is attained near
\((\pi,\pi,\pi,\cdot)\), not at the toron; the sampled real floor of
\(B_1\) is \(2.678\), about \(23\) times the inherited floor.  These
are floating diagnostics excluded from the canonical payload.
\end{record}

\begin{assessment}[Where the quadrature obstruction actually lies]
\label{ass:v31-obstruction}
The diagnostics indicate a true Hermitian-part strip of the order
\(0.3\) for the Wilson Gram matrix, two orders of magnitude beyond the
certified (35.14) and four beyond (53.22).  With \(\rho\approx0.3\) the
necessary mesh of Countertheorem~\ref{counter:v31-mesh} would be
\(M\approx66\), a feasible \(1.9\times10^7\)-point grid.  The entire
remaining gap between the archived certificate and a feasible one is
therefore the supremum stock \(N_W\), for which every archived bound is
a product of factor norms and loses many orders of magnitude.  A
feasible sign certificate would require an exact alias-level proof of
the wide strip and a cancellation-preserving supremum bound; by
Corollary~\ref{cor:v31-refutation-instrument} it would then only refute
SCIC for the present pair.
\end{assessment}

\subsection{Executable provenance}

\begin{record}[V31 verifiers]
\label{rec:v31-verifiers}
The census and box-count verifier
\begin{center}
\path{scripts/verify_ym_postarchive_v31_frozen_chart_global_residual.py}
\end{center}
has \(154\) physical lines, \(6031\) bytes, and SHA--256
\[
\texttt{CD0C8D166CA30BDE712B315222FEC96AE9F4E53704099127EFD5F88148CC4D37}.
                                                               \tag{35.18}
\]
Its canonical payload SHA--256 is
\[
\texttt{A3D3B2EBE02AB8A74C9A9F4EE8BE1F3BF1891261AAE89AFD593361E44D30E1B4},
                                                               \tag{35.19}
\]
and its stored artifact
\path{artifacts/ym_postarchive_v31_frozen_chart_global_residual.json},
which lists all \(512\) exact defect squares, has SHA--256
\[
\texttt{4389816E3B4571941C31FA11F96276220DA9F82FA5B33A60C505169309E4AB75}.
                                                               \tag{35.20}
\]
The strip verifier
\begin{center}
\path{scripts/verify_ym_postarchive_v31_cosh_hermitian_strip.py}
\end{center}
has \(205\) physical lines, \(7942\) bytes, and SHA--256
\[
\texttt{5750F96A86D3AEB678FE5C2A56094659E0EEF07313EB3EA42B9AD790417A4A23}.
                                                               \tag{35.21}
\]
Its canonical payload SHA--256, which excludes the diagnostics of
Record~\ref{rec:v31-diagnostic}, is
\[
\texttt{8D20234995C1E53506097E91B36335726F149E1AD0395152E22120CA2DA80A21}.
                                                               \tag{35.22}
\]
Both executables import only the frozen exact sources already recorded
in V2, V10, V12, and V53.
\end{record}

\subsection{Decision for V32}

\begin{decision}[Upstream to the nested Gaussian tower]
\label{dec:v31-v32}
V32 must not resume prime production.  It should:
\begin{enumerate}[label=\arabic*.]
\item record the exact normalization relating the V60 response
      \(F_e\) to the extractor \(\mathfrak b\), so that
      \(\langle F_{\rm W}\rangle\) can be compared with
      \(b_{\rm YM}=\frac{22C_A}{3(4\pi)^2}\log2\) in one convention;
\item derive the exact Schur-complement recursion of the parity-cell
      Hessian across two nested dyadic shells, prove the induced
      second-shell response formula at the determinant level, and
      evaluate the first two per-shell Gaussian coefficients as
      floating diagnostics, to test which of the two mutually exclusive
      hypotheses of Proposition~\ref{prop:v31-nested} the actual tower
      satisfies;
\item if the second-shell coefficient moves toward a fixed value,
      formulate (35.9) as the live Gaussian card and route the
      compensating class into \(c_j\); if it does not move, record the
      self-similar obstruction as a theorem about the parity-cell
      tower.
\end{enumerate}
An alias-level proof of the wide strip suggested by (35.17) is a
secondary target, useful only after item 2 shows that a sign
certificate is needed at all.  The compulsory companion audit remains
V35, and the ten-iteration sweep of newly generated companion notes is
due at V40.
\end{decision}

\begin{assessment}[V31 acquisition]
\label{ass:v31}
V31 proves exactly that the V26--V27 artifact certifies a single
anchor box and that the box-uniform route needs more than
\(6.69\times10^{12}\) boxes; withdraws the fourteen-prime production;
proves that a certified nonzero constant class refutes SCIC and ADMC
for the self-similar Wilson--endpoint tower rather than closing them;
and supplies a second-order strip theorem with a certified
\(190\)-fold gain whose propagated mesh threshold shows the supremum
stock, not the strip, to be the binding obstruction.  It redirects the
programme to the nested Gaussian tower.
\end{assessment}

\begin{record}[V31 append record]
\label{rec:v31-append}
V31 is appended to the clean V30 whose installed record was
\[
\begin{array}{c|c}
\text{lines}&10535\\
\text{bytes}&407339\\
\text{SHA--256}&
\texttt{C777BD838A65CB56FFE00EAF0283A1C00044D6FC83E47DBA65E49E8C2CC51790}.
\end{array}                                                    \tag{35.23}
\]
After validation only V31 is to remain the active numeric YM note.
\end{record}

\begin{firewall}[Claim boundary after V31]
\label{fire:v31-final}
V31 proves no Haar sign, no nested-tower coefficient, no strip wider
than (35.13)--(35.14), and no continuum statement.  The Gaussian
anomaly (14.16), the reset row (14.17), the compatible interacting
measure, regulator removal, Osterwalder--Schrader reconstruction,
physical clustering, and the Hamiltonian Yang--Mills mass gap remain
open.
\end{firewall}

% =============================================================================
% END APPENDED POST-TENTH-ARCHIVE V31 MATERIAL
% =============================================================================
% =============================================================================
% BEGIN APPENDED POST-TENTH-ARCHIVE V32 MATERIAL
% =============================================================================

\section{V32: the exact nested Gaussian kernel recursion and the failure of
self-similarity}
\label{sec:v32}

V31 redirected the programme from the single-anchor sign problem to the
nested Gaussian tower and asked which of two mutually exclusive
hypotheses that tower satisfies: scale independence of the shell banks,
or a convergent coboundary of blocked constants.  This note settles the
question at the level of the Gaussian kernel, where it is decidable
exactly.  The archived level-one Schur matrix of the router equation is
recognized as the first step of an exact renormalization recursion for
the four-by-four coarse-link kernel.  The recursion is evaluated in
Gaussian-rational arithmetic at all \(256\) quarter-turn momenta at
level one, for both actions, and at the sixteen parity corners at level
two.  The level-one kernel is not a scalar multiple of the bare kernel:
at the corner \((\pi,0,0,0)\) the transverse eigenvalue drops from
\(4\) to \(15/7\), and at level two to \(25823/24660\), while the
small-momentum Maxwell coefficient is preserved along an exact rational
sequence.  Hence the hypothesis of Theorem~\ref{thm:v2-endpoint-reduction}
fails for the nested tower and the archived reference reset has nonzero
memory already at tree level.  Floating diagnostics show the corner value
and the vertical floor halving at every level, so the sharp dyadic
coarse-link rule does not produce a uniformly conditioned Gaussian fixed
point.  The note ends with the normalization audit requested in
Decision~\ref{dec:v31-v32} and with the V33 target.

\subsection{The recursion}

Let \(\mathcal E=\bigoplus_{\epsilon\in\{0,1\}^4}\mathbb C^4\) be the
alias space of V4 and let \(V_R(k)\in M_{64\times45}\) and
\(V_C(k)\in M_{64\times4}\) be the parity Fourier frames of the \(45\)
retained and the \(4\) second-constituent edges: row \((\epsilon,\mu)\)
of the column belonging to the edge \((s,\mu)\) carries
\(e^{-ip_\epsilon\cdot s}/4\), \(p_\epsilon=(k+2\pi\epsilon)/2\).  For a
Hermitian positive semidefinite four-by-four kernel function
\(\widehat S(p)\), define the level map
\[
 \mathcal T[\widehat S](k)
 :=C-B^*A^{-1}B,\qquad
 \begin{pmatrix}A&B\\B^*&C\end{pmatrix}
 =\begin{pmatrix}V_R^*\\V_C^*\end{pmatrix}
 \Bigl(\bigoplus_\epsilon\widehat S(p_\epsilon)\Bigr)
 \begin{pmatrix}V_R&V_C\end{pmatrix},                           \tag{36.1}
\]
whenever \(A\) is invertible, and put
\[
 \widehat S^{(0)}(p):=|q(p)|^2I-q(p)q(p)^*,\qquad
 q_\mu(p)=e^{ip_\mu}-1,\qquad
 \widehat S^{(j+1)}:=\mathcal T[\widehat S^{(j)}].                \tag{36.2}
\]
The cleared kernels carry the factor \(3\): the physical Wilson kernels
are \(\widehat S^{(j)}/3\).  For the endpoint form, \(\widehat S^{(0)}\)
is replaced by \((1+3\lambda_\epsilon)\widehat S^{(0)}\) at alias
\(\epsilon\), with \(\lambda_\epsilon=|q(p_\epsilon)|^2\).

\begin{proposition}[Level one is the archived Schur matrix]
\label{prop:v32-level-one}
With \(D=(D_R\ D_C\ D_T)\) the V21 incidence at coarse momentum \(k\),
\[
 V_R^*\Bigl(\bigoplus_\epsilon\widehat S^{(0)}(p_\epsilon)\Bigr)V_R
 =D_R^*D_R,\qquad
 V_R^*\Bigl(\bigoplus_\epsilon\widehat S^{(0)}(p_\epsilon)\Bigr)V_C
 =D_R^*D_C,                                                       \tag{36.3}
\]
and \(\widehat S^{(1)}(k)\) equals three times the archived Wilson
Schur matrix (23.9).  Both identities in (36.3) are verified exactly
at \(k=(\pi,0,0,0)\) by the V32 executable, and the archived corner
properties (rank zero at the toron, rank three elsewhere, exact
symmetry) are reproduced.
\end{proposition}

\begin{proof}
The alias energy identity (63.6) of the ninth archive states that the
fine curl energy of a retained configuration equals
\(16\sum_\epsilon(a^\epsilon)^*C_\epsilon a^\epsilon\) with
\(a^\epsilon\) the normalized parity amplitudes; the frame normalization
\(1/4\) reproduces the factor \(16\), and the same computation with one
second-constituent column gives the mixed block.  The executable checks
both identities entrywise in Gaussian-rational arithmetic.  The Schur
complement of the blocks \((D_R^*D_R,D_R^*D_C,D_C^*D_C)\) is by
definition \(C^c_{\rm W}+(B^c_{\rm W})^*L_{\rm W}\) with the router
\(L_{\rm W}=-(A^c_{\rm W})^{-1}B^c_{\rm W}\), which is three times (23.9).
\end{proof}

\begin{theorem}[Transversality is preserved by the recursion]
\label{thm:v32-transverse}
Suppose \(\widehat S^{(j)}(p)\succeq0\) and
\(\widehat S^{(j)}(p)q(p)=0\) for all \(p\), and suppose the vertical
block \(A\) in (36.1) is invertible at \(k\).  Then
\(\widehat S^{(j+1)}(k)\succeq0\) and
\[
 \widehat S^{(j+1)}(k)\,q(k)=0,\qquad q_\mu(k)=e^{ik_\mu}-1.       \tag{36.4}
\]
In particular \(\operatorname{rank}\widehat S^{(j+1)}(k)\le3\).
\end{theorem}

\begin{proof}
A Schur complement of a positive semidefinite Hermitian form with
positive definite pivot block is positive semidefinite, and equals the
minimum of the form over the pivot variables:
\(c^*\widehat S^{(j+1)}c=\min_v (v,c)^*\mathcal H(v,c)\).  Let
\(\phi\) be a level-\((j+1)\) gauge function at momentum \(k\) and
\(c=q(k)\phi\) its coarse gradient.  Extend \(\phi\) to level \(j\) by
the piecewise constant rule \(\phi_j(2n+s)=\phi(2n)\).  The level-\(j\)
pure gauge field \(d\phi_j\) has coarse links
\(\phi(2n+2e_\mu)-\phi(2n)\), namely \(c\), and its alias amplitudes are
\(q(p_\epsilon)\widehat\phi_j(p_\epsilon)\), which the level-\(j\) form
annihilates by the hypothesis \(\widehat S^{(j)}q=0\).  Its energy is
therefore zero.  The tree gauge is a section of the orbits of the
level-\(j\) gauge transformations vanishing at base points; such
transformations preserve the coarse links and the energy, so a
configuration in the slice with coarse links \(c\) and zero energy
exists.  The minimum is thus zero, i.e.\ \(q(k)^*\widehat S^{(j+1)}(k)q(k)=0\),
and positive semidefiniteness turns this into (36.4).  Rank at most
three follows.
\end{proof}

\subsection{Exact level-one and level-two computations}

\begin{theorem}[Exact level-one census]
\label{thm:v32-level-one-census}
For both the Wilson and the endpoint forms and all \(256\) quarter-turn
momenta \(k\in\Gamma_4\), the Gaussian-rational Schur complement
\(\widehat S^{(1)}(k)\) is Hermitian, satisfies (36.4) exactly, and has
rank three at every \(k\ne0\) and rank zero at the toron.  At the
parity corners the Wilson kernels have the closed forms
\[
\begin{array}{c|c}
 k&\widehat S^{(1)}_{\rm W}(k)\\ \hline
 (\pi,0,0,0)&\operatorname{diag}(0,\ 15/7,\ 15/7,\ 15/7)\\[2pt]
 (\pi,\pi,0,0)&
 \begin{pmatrix}12/7&-12/7\\-12/7&12/7\end{pmatrix}
 \oplus\operatorname{diag}(8/3,\ 8/3)\\[6pt]
 (\pi,\pi,\pi,0)&
 \bigl(\tfrac52\delta_{\mu\nu}-\tfrac54(1-\delta_{\mu\nu})\bigr)_{\mu,\nu\le3}
 \oplus(3)\\[4pt]
 (\pi,\pi,\pi,\pi)&3I-J,
\end{array}                                                      \tag{36.5}
\]
where \(J\) is the all-ones matrix, whereas the bare kernels at the same
corners are
\(\operatorname{diag}(0,4,4,4)\),
\(\begin{pmatrix}4&-4\\-4&4\end{pmatrix}\oplus\operatorname{diag}(8,8)\),
\((8\delta_{\mu\nu}-4(1-\delta_{\mu\nu}))\oplus(12)\), and \(16I-4J\).
For the endpoint form at \((\pi,0,0,0)\),
\[
 \widehat S^{(1)}_*=\operatorname{diag}(0,\ 61845/2851,\ 61845/2851,\ 61845/2851),
 \qquad
 \widehat S^{(0)}_*=\operatorname{diag}(0,52,52,52).               \tag{36.6}
\]
\end{theorem}

\begin{proof}
At a quarter-turn momentum every phase \(e^{ik_\mu}\) lies in
\(\{1,i,-1,-i\}\), so the incidence blocks are Gaussian-rational
matrices; the executable performs Gauss--Jordan elimination over
\(\mathbb Q(i)\) for the router, forms the Schur complement, and tests
Hermitian symmetry, the defect (36.4), and the rank by exact row
reduction.  The displayed matrices are the stored exact entries.  The
bare values follow from (36.2) with \(\lambda=4,8,12,16\).
\end{proof}

\begin{corollary}[Self-similarity fails at tree level]
\label{cor:v32-not-self-similar}
\(\widehat S^{(1)}_{\rm W}\ne\widehat S^{(0)}\) and
\(\widehat S^{(1)}_*\ne\widehat S^{(0)}_*\).  The ratio of transverse
eigenvalues at \((\pi,0,0,0)\) is \(15/28\) for Wilson and
\(61845/148252\) for the endpoint; at \((\pi,\pi,\pi,\pi)\) the Wilson
ratio is \(1/4\).  Consequently the scale-independence hypothesis of
Theorem~\ref{thm:v2-endpoint-reduction} is false for the nested tower,
and the archived reference reset, which replaces the blocked kernel by
the bare one at the next scale, changes the quadratic form by a factor
between \(1/4\) and \(15/28\) at the corners.
\end{corollary}

\begin{proof}
Compare (36.5)--(36.6) with the bare values.
\end{proof}

\begin{theorem}[Exact level two at the parity corners]
\label{thm:v32-level-two}
For the Wilson form, \(\widehat S^{(2)}(k)\) at the sixteen parity
corners \(k\in\{0,\pi\}^4\) is obtained exactly from the \(256\)
level-one kernels, because the sixteen aliases \(p_\epsilon\) of a
corner are quarter-turn momenta.  All sixteen results are Hermitian,
satisfy (36.4) exactly, have rank three except at the toron where the
kernel vanishes, and
\[
 \widehat S^{(2)}_{\rm W}(\pi,0,0,0)
 =\operatorname{diag}\Bigl(0,\ \tfrac{25823}{24660},\ \tfrac{25823}{24660},\ \tfrac{25823}{24660}\Bigr).
                                                                 \tag{36.7}
\]
Thus the transverse corner eigenvalue evolves exactly as
\(4\to15/7\to25823/24660\), with successive ratios \(15/28\) and
\(25823/52830\).
\end{theorem}

\begin{proof}
The executable builds the \(64\times64\) block-diagonal form from the
stored level-one kernels at \(\pm i\) phases, applies the exact frames,
and repeats the elimination.  All checks are exact.
\end{proof}

\begin{theorem}[Exact preservation of the Maxwell coefficient along a
rational sequence]
\label{thm:v32-moment}
Let \(z_n=\frac{n^2-1}{n^2+1}+\frac{2n}{n^2+1}i\), a Gaussian rational
on the unit circle, and \(k_n=(\theta_n,0,0,0)\) with
\(e^{i\theta_n}=z_n\), so that \(|q_1(k_n)|^2=4/(n^2+1)\).  The exact
Wilson level-one transverse entries satisfy
\[
\begin{array}{c|c|c|c}
 n&\theta_n&\widehat S^{(1)}_{22}(k_n)\big/|q_1(k_n)|^2&\hbox{defect from }1\\ \hline
 4&0.48996&56033/59172&-5.30\times10^{-2}\\
 8&0.24871&813281/825252&-1.45\times10^{-2}\\
 16&0.12484&12689633/12736932&-3.71\times10^{-3}\\
 32&0.06248&201752801/201941412&-9.34\times10^{-4}\\
 64&0.03125&3222929633/3223683492&-2.34\times10^{-4}.
\end{array}                                                      \tag{36.8}
\]
The defects decrease by the factors \(3.66,3.91,3.98,3.99\), i.e.\ as
\(\theta_n^2\), with \(\widehat S^{(1)}_{22}/|q_1|^2=1-0.24\,\theta^2+O(\theta^4)\)
along the sequence.
\end{theorem}

\begin{proof}
The phases \(e^{ik_n}\) are Gaussian rational, so the level-one
elimination is exact; the executable records the reduced fractions in
(36.8).  The ratios of consecutive defects are exact rational numbers
whose decimals are displayed.
\end{proof}

\begin{remark}[Two exact facts, one conclusion]
\label{rem:v32-two-facts}
The corner ratio \(15/28\) and the small-momentum ratio above
\(0.9997\) are both exact.  Together they prove that the blocked kernel
is the bare kernel neither up to a scalar nor up to a small
perturbation: the sharp dyadic coarse-link rule preserves the Maxwell
coefficient at long wavelength and roughly halves the kernel at the
Brillouin corner.  This is the tree-level content of the reset memory
whose one-loop shadow appears in the ledger row (14.17).
\end{remark}

\subsection{Deeper levels: floating diagnostics}

\begin{record}[Recursion to level four, non-hashed]
\label{rec:v32-deep}
Evaluating (36.1)--(36.2) in floating point along the same recursion at
the corner \((\pi,0,0,0)\) gives the transverse eigenvalue
\[
 4,\quad2.1428571,\quad1.0471614,\quad0.50981238,\quad0.2511057
 \qquad(j=0,\ldots,4),                                           \tag{36.9}
\]
with successive ratios \(0.5357,\ 0.4887,\ 0.4868,\ 0.4925\).  The
levels one and two agree with the exact values \(15/7\) and
\(25823/24660\) to the displayed digits.  On the axis
\(k=(t,0,0,0)\) the ratio \(\widehat S^{(j)}_{22}/|q_1|^2\) is
\[
\begin{array}{c|cccc}
 j\backslash t&0.2&0.5&1.0&2.0\\ \hline
 1&0.99055&0.94493&0.82304&0.61306\\
 2&0.97240&0.85164&0.60493&0.33307\\
 3&0.93601&0.70433&0.38822&0.17093\\
 4&0.87039&0.52235&0.22556&0.08640 ,
\end{array}                                                      \tag{36.10}
\]
so the defect at fixed \(t=0.2\) grows as
\(0.0095,\ 0.0276,\ 0.064,\ 0.130\), roughly doubling per level.  The
smallest eigenvalue of the cleared vertical block \(A\) in (36.1),
sampled at eleven momenta, is
\[
 0.1269,\quad0.0457,\quad0.0188,\quad0.00844\qquad(j=1,\ldots,4),
                                                                 \tag{36.11}
\]
the first being consistent with the proved level-one floor
\(147/1250=0.1176\).  These values are diagnostics and are excluded
from the canonical payload.
\end{record}

\begin{assessment}[No uniformly conditioned Gaussian fixed point]
\label{ass:v32-fixed-point}
The diagnostics indicate that, under the sharp dyadic coarse-link rule,
the coarse kernel keeps its Maxwell coefficient at long wavelength while
its corner value and the vertical floor decay geometrically, and the
momentum range on which it is Maxwell-like shrinks with the level.  A
rescaling by \(2^{j}\) stabilizes the corner but not the long-wavelength
part, so no single normalization yields a convergent local kernel.  The
archived reference reset is therefore not a small correction: it is the
device that keeps the tower uniformly conditioned, at the price of a
per-scale memory which V31 proved cannot be sublinear when constant.
The Gaussian tower of the programme is thus caught between a nested
form with degenerating conditioning and a reset form with linear memory.
\end{assessment}

\subsection{Normalization audit}

\begin{record}[What the archive fixes and what it does not]
\label{rec:v32-normalization}
The archived Ward density is
\(W_{j,t}=\frac32\partial_q^2\Omega_{j,t}|_{q=0}\) by (16.6), and the
V60 oracle returns the second external-momentum derivative
\(\partial_q^2\) of the tadpole and bubble words, so
\[
 \mathfrak b\gamma_j^{(e)}=\langle W\rangle=\tfrac32\langle F_e\rangle
                                                                 \tag{36.12}
\]
in the archive's own convention, giving the diagnostic values
\(\tfrac32\cdot222.706=334.06\) and \(\tfrac32\cdot(-92.138)=-138.21\)
from (10.29).  The archive expresses the universal candidate as
\(b_{\rm YM}=\frac{22C_A}{3(4\pi)^2}\log2\), i.e.\ \(0.0644\) for
\(C_A=2\) and \(0.0966\) for \(C_A=3\), but the colour, background, and
cell-volume factors relating the one-colour commuting-toron response to
\(b_{\rm YM}\) are not recorded in any compact source consulted here.
No comparison between (36.12) and \(b_{\rm YM}\) is therefore asserted;
the missing factor is an open bookkeeping item to be recovered from the
tenth archive before any universality claim.
\end{record}

\subsection{Executable provenance}

\begin{record}[V32 verifier]
\label{rec:v32-verifier}
The executable
\begin{center}
\path{scripts/verify_ym_postarchive_v32_nested_gaussian_kernel.py}
\end{center}
has \(545\) physical lines, \(20234\) bytes, and SHA--256
\[
\texttt{CC27C0AE6CC15AE6EB378A993A9660CC60A18AF670DF55679469DE9922DA902F}.
                                                               \tag{36.13}
\]
Its canonical payload, which contains all exact kernels, hashes, ranks,
and the ratios (36.8) but not the diagnostics of
Record~\ref{rec:v32-deep}, has SHA--256
\[
\texttt{9F15813288A7600ABABC5545C75342F9461C46FD6F5D49543C6F4074ABA48B17},
                                                               \tag{36.14}
\]
and the stored artifact
\path{artifacts/ym_postarchive_v32_nested_gaussian_kernel.json} has
SHA--256
\[
\texttt{9B81CE614675974B8A5D8ECC95A631C23B54F6B8B46922B833EEE48DF21FAD65}.
                                                               \tag{36.15}
\]
The exact part ran in \(400\) seconds and the diagnostics in \(92\)
seconds; the only imported source is the frozen V21 incidence grammar.
\end{record}

\subsection{Decision for V33}

\begin{decision}[A smeared dyadic blocking at the Gaussian level]
\label{dec:v32-v33}
Since the sharp rule has nonzero tree-level memory
(Corollary~\ref{cor:v32-not-self-similar}) and degenerating
conditioning (Record~\ref{rec:v32-deep}), V33 should test the
standard remedy at the Gaussian level: replace the exact coarse link
by a coarse link fluctuating about the blocked product with a
gauge-covariant Gaussian weight of stiffness \(\kappa\).  For the
abelian one-colour kernel this amounts to the recursion
\(\widehat S\mapsto\bigl[(\mathcal T[\widehat S])^{-1}+\kappa^{-1}P\bigr]^{-1}\)
on the transverse three-space.  V33 must:
\begin{enumerate}[label=\arabic*.]
\item prove that the smeared level map preserves transversality and
      positivity and reduces to (36.1) as \(\kappa\to\infty\);
\item evaluate the smeared recursion exactly at the corners for
      rational \(\kappa\) and in floating point to level six, and
      determine whether the corner value, the axis defect, and the
      vertical floor converge for some \(\kappa\);
\item if a convergent kernel exists, record it as the candidate
      Gaussian fixed point and reformulate the reset row (14.17) as a
      contraction toward it; if none exists in the tested range, record
      the obstruction and return to the interacting reset ledger.
\end{enumerate}
The bookkeeping item of Record~\ref{rec:v32-normalization} is to be
recovered from the tenth archive at the V35 audit.
\end{decision}

\begin{assessment}[V32 acquisition]
\label{ass:v32}
V32 turns the archived router Schur matrix into an exact Gaussian
renormalization recursion, proves that it preserves transversality,
computes it exactly at \(256\) level-one momenta for both actions and at
sixteen level-two corners for Wilson, and proves that the blocked kernel
differs from the bare kernel by corner factors between \(1/4\) and
\(15/28\) while preserving the long-wavelength Maxwell coefficient.  It
thereby refutes the self-similar hypothesis for the nested tower at tree
level and locates the tower's conditioning problem.
\end{assessment}

\begin{record}[V32 append record]
\label{rec:v32-append}
V32 is appended to the clean V31 whose installed record was
\[
\begin{array}{c|c}
\text{lines}&11118\\
\text{bytes}&432885\\
\text{SHA--256}&
\texttt{73712A413552FD14065BA67CA448B57F3742DED3CC0EFF1A74512BCF2B83DDB}.
\end{array}                                                    \tag{36.16}
\]
After validation only V32 is to remain the active numeric YM note.
\end{record}

\begin{firewall}[Claim boundary after V32]
\label{fire:v32-final}
V32 proves exact tree-level statements about Gaussian kernels.  It
proves no one-loop shell coefficient, no fixed point, no bound on the
reset memory at one loop, no Haar sign, and no continuum statement.
The Gaussian anomaly (14.16), the reset row (14.17), the compatible
interacting measure, regulator removal, Osterwalder--Schrader
reconstruction, physical clustering, and the Hamiltonian Yang--Mills
mass gap remain open.
\end{firewall}

% =============================================================================
% END APPENDED POST-TENTH-ARCHIVE V32 MATERIAL
% =============================================================================
% =============================================================================
% BEGIN APPENDED POST-TENTH-ARCHIVE V33 MATERIAL
% =============================================================================

\section{V33: a transversely averaged dyadic blocking and its exact
Gaussian fixed point on the coordinate axes}
\label{sec:v33}

V32 proved that the sharp straight coarse-link rule has nonzero
tree-level memory and observed that its nested Gaussian kernel halves at
the Brillouin corner at every level.  This note replaces the straight
rule by the average of the eight staircase paths through the transverse
offsets of the dyadic cell, which is the abelian Gaussian limit of a
gauge-covariant block-link rule.  The blocked kernel is again an exact
constrained minimum and is evaluated in Gaussian-rational arithmetic at
all \(256\) quarter-turn momenta at level one and at the sixteen parity
corners at level two.  The transverse aliases decouple exactly on every
coordinate axis, where the recursion reduces to a scalar Riccati-type
map with the closed-form solution
\(s_j(t)=4\sin^2(t/2)/(1-\tfrac23(1-4^{-j})\sin^2(t/2))\).  The
level-\(j\) axis kernel therefore converges geometrically, with ratio
\(1/4\), to the local fixed point \(12\sin^2(t/2)/(2+\cos t)\), which is
analytic in the strip \(|\operatorname{Im}t|<\log(2+\sqrt3)\) and
preserves the Maxwell coefficient.  Floating diagnostics indicate the
same convergence off the axes.  The reset memory of V31--V32 thereby
acquires a candidate replacement: convergence to a Gaussian fixed point.

\subsection{The averaged rule and its normalization}

\begin{definition}[Staircase eight-path average]
\label{def:v33-average}
For a direction \(\mu\) and a transverse offset \(t\in\{0,1\}^4\) with
\(t_\mu=0\), let \(\Pi_{\mu,t}\) be the path from the base point \(0\) of
the cell to the base point \(2e_\mu\) of the next cell which first takes
the unit steps \(e_\nu\), \(t_\nu=1\), in increasing order of \(\nu\),
then the two links \((t,\mu)\) and \((t+e_\mu,\mu)\), and finally the
same transverse steps in reverse.  The averaged coarse link is
\[
 c_\mu:=\frac18\sum_{t:\,t_\mu=0}\sum_{\ell\in\Pi_{\mu,t}}\pm a_\ell,
                                                                 \tag{37.1}
\]
with the sign of the orientation.  In alias amplitudes, with
\(p_\epsilon=(k+2\pi\epsilon)/2\) and \(z_\mu=e^{ik_\mu}\),
\[
 c_\mu=\frac18\sum_{t:\,t_\mu=0}\sum_\epsilon\Bigl[
 e^{ip_\epsilon\cdot t}\bigl(1+e^{ip_{\epsilon,\mu}}\bigr)a^\epsilon_\mu
 +(1-z_\mu)\sum_{(\nu,\sigma)\in{\rm steps}(t)}
   e^{ip_\epsilon\cdot\sigma}a^\epsilon_\nu\Bigr],                 \tag{37.2}
\]
where \({\rm steps}(t)\) lists the outgoing transverse steps
\((\nu,\sigma)\) with their starting sites \(\sigma\).  The straight rule
of V32 is the single term \(t=0\).
\end{definition}

\begin{lemma}[Ward property and slice normalization]
\label{lem:v33-ward}
For every fine gauge function \(\phi\), the averaged coarse link of the
pure gauge \(d\phi\) is \(\phi(2e_\mu)-\phi(0)\), i.e.\ the coarse
gradient of the restriction of \(\phi\) to base points.  The map (37.2)
is invariant under fine gauge transformations vanishing at base points.
If \(V\) is the parity frame with entries \(e^{-ip_\epsilon\cdot s}/4\)
and \(a\) is a slice configuration, then \(Va=4a^\epsilon\), so in slice
coordinates the constraint reads \(c=\tfrac14{\cal B}Va\) with
\({\cal B}\) the matrix of (37.2).
\end{lemma}

\begin{proof}
Every path \(\Pi_{\mu,t}\) runs from \(0\) to \(2e_\mu\), so the sum of
\(d\phi\) along it telescopes to \(\phi(2e_\mu)-\phi(0)\) independently
of \(t\); averaging over \(t\) preserves this value.  A gauge
transformation vanishing at base points changes each path sum by the
difference of its values at the two endpoints, which is zero.  For the
normalization, the ninth archive's identity (63.6) uses
\(a^\epsilon=\frac1{16}\sum_se^{-ip_\epsilon\cdot s}a(s)\), whereas the
frame sums with weight \(1/4\); hence \(Va=4a^\epsilon\).  The V33
executable checks the Ward identity exactly on the base-point
concentrated gauge function \(\phi=16\delta_{s,0}\), whose alias
amplitudes are \(q(p_\epsilon)\), obtaining \(16(z_\mu-1)\) in every
component.
\end{proof}

\begin{theorem}[Constrained level map]
\label{thm:v33-level-map}
Let \(\mathcal H=\begin{pmatrix}A&B\\B^*&C\end{pmatrix}\) be the
level-\(j\) form in slice coordinates \((a_R,a_C)\) as in (36.1), and
let the constraint be \(Xa_R+Ya_C=c\) with \(Y\in M_4\) invertible.
Put \(W=Y^{-1}X\), \(\widetilde A=A-BW-W^*B^*+W^*CW\), and
\(\widetilde B=B-W^*C\).  If \(\widetilde A\) is invertible, the
constrained minimum of the form is \(c^*\widehat S'c\) with
\[
 \widehat S'=Y^{-*}\bigl(C-\widetilde B^*\widetilde A^{-1}\widetilde B\bigr)Y^{-1}.
                                                                 \tag{37.3}
\]
For the straight rule \(X=0\), \(Y=I\), and (37.3) is (36.1).  If the
level-\(j\) kernel is positive semidefinite and transverse, then
\(\widehat S'\) is positive semidefinite, transverse, and of rank at
most three.
\end{theorem}

\begin{proof}
Substitute \(a_C=Y^{-1}c-Wa_R\) into the form; the result is
\(a_R^*\widetilde Aa_R+2\operatorname{Re}(a_R^*\widetilde Bu)+u^*Cu\)
with \(u=Y^{-1}c\), whose minimum over \(a_R\) is
\(u^*(C-\widetilde B^*\widetilde A^{-1}\widetilde B)u\).  Positivity is
inherited from the form.  For transversality repeat the proof of
Theorem~\ref{thm:v32-transverse}: the piecewise constant extension of a
coarse gauge function is a level-\(j\) pure gauge with zero energy, and by
Lemma~\ref{lem:v33-ward} its averaged coarse link is the coarse gradient;
the slice argument is unchanged.  The V33 executable verifies exact
transversality at every computed momentum and reproduces \(15/7\) and
\(25823/24660\) for the straight rule through (37.3).
\end{proof}

\subsection{Exact computations}

\begin{theorem}[Exact level one and level two for the averaged rule]
\label{thm:v33-exact}
For the averaged rule the Gaussian-rational kernels
\(\widehat S^{(1)}(k)\) at all \(256\) quarter-turn momenta, computed
through the incidence route, and \(\widehat S^{(2)}(k)\) at the sixteen
parity corners, computed through the alias route from the stored
level-one kernels, are Hermitian, exactly transverse, of rank three off
the toron and zero at the toron.  At the corner \((\pi,0,0,0)\),
\[
 \widehat S^{(0)}=\operatorname{diag}(0,4,4,4),\qquad
 \widehat S^{(1)}=\operatorname{diag}(0,8,8,8),\qquad
 \widehat S^{(2)}=\operatorname{diag}\bigl(0,\tfrac{32}{3},\tfrac{32}{3},\tfrac{32}{3}\bigr),
                                                                 \tag{37.4}
\]
whereas the straight rule gives \(15/7\) and \(25823/24660\).  At the
other corners the level-one averaged kernels are
\[
 \widehat S^{(1)}(\pi,\pi,0,0)=\begin{pmatrix}\tfrac{192}{79} & -\tfrac{192}{79} & 0 & 0 \\ -\tfrac{192}{79} & \tfrac{192}{79} & 0 & 0 \\ 0 & 0 & 16 & 0 \\ 0 & 0 & 0 & 16\end{pmatrix},\qquad
 \widehat S^{(1)}(\pi,\pi,\pi,\pi)=\begin{pmatrix}\tfrac{992}{243} & -\tfrac{32}{81} & -\tfrac{224}{243} & -\tfrac{224}{81} \\ -\tfrac{32}{81} & \tfrac{224}{27} & -\tfrac{160}{81} & -\tfrac{160}{27} \\ -\tfrac{224}{243} & -\tfrac{160}{81} & \tfrac{4064}{243} & -\tfrac{1120}{81} \\ -\tfrac{224}{81} & -\tfrac{160}{27} & -\tfrac{1120}{81} & \tfrac{608}{27}\end{pmatrix}.                \tag{37.5}
\]
Along the rational sequence \(k_n=(\theta_n,0,0,0)\) of
Theorem~\ref{thm:v32-moment}, the exact level-one transverse ratio
\(\widehat S^{(1)}_{22}/|q_1|^2\) equals
\[
 \tfrac{34}{33},\quad \tfrac{130}{129},\quad \tfrac{514}{513},\quad \tfrac{2050}{2049},\quad \tfrac{8194}{8193}                                                      \tag{37.6}
\]
for \(n=4,8,16,32,64\), in agreement with (37.8) below.
\end{theorem}

\begin{proof}
The incidence route evaluates the coarse-link map (37.1) directly on the
cell edges, with the returning connectors carrying the phase \(z_\mu\) of
the neighbouring cell, restricts it to the slice columns, and applies
(37.3) to the exact Gram blocks of V32.  The alias route applies (37.2)
with the normalization of Lemma~\ref{lem:v33-ward}.  Both routes agree at
the corner \((\pi,0,0,0)\), where the executable also checks that the
straight rule returns \(15/7\).  All displayed entries are stored exact
fractions.
\end{proof}

\subsection{Exact reduction on the coordinate axes}

\begin{theorem}[Axis decoupling]
\label{thm:v33-axis}
Let \(k=(t,0,0,0)\) with \(0<t<2\pi\) and suppose the level-\(j\)
averaged kernel satisfies
\(\widehat S^{(j)}(\tau,0,0,0)=\operatorname{diag}(0,s_j(\tau),s_j(\tau),s_j(\tau))\)
for all \(\tau\), with \(s_j>0\) off \(2\pi\mathbb Z\).  Then the same
holds at level \(j+1\) with
\[
 s_{j+1}(t)=\frac{4}{\dfrac{\cos^2(t/4)}{s_j(t/2)}+\dfrac{\sin^2(t/4)}{s_j(t/2+\pi)}}.
                                                                 \tag{37.7}
\]
The hypothesis holds at level zero with \(s_0(t)=4\sin^2(t/2)\).
\end{theorem}

\begin{proof}
At \(k=(t,0,0,0)\) the aliases have \(p_\epsilon=(t/2+\pi\epsilon_1,
\pi\epsilon_2,\pi\epsilon_3,\pi\epsilon_4)\).  For a transverse
direction \(\mu\in\{2,3,4\}\) the connector terms in (37.2) carry
\(1-z_\mu=0\).  In the straight term the sum over offsets factorizes,
\(\sum_{t:\,t_\mu=0}e^{ip_\epsilon\cdot t}
 =\prod_{\nu\ne\mu}(1+e^{ip_{\epsilon,\nu}})\), and the factor
\(1+e^{ip_{\epsilon,\mu}}\) is present as well; hence every alias with a
transverse flip \(\epsilon_\nu=1\), \(\nu\ge2\), has weight zero, and the
two aliases \(\epsilon=(\epsilon_1,0,0,0)\) have weight
\(b_{\epsilon,\mu}=1+e^{i(t/2+\pi\epsilon_1)}\), with
\(|b|^2=4\cos^2(t/4)\) and \(4\sin^2(t/4)\).  At these two aliases the
level-\(j\) kernel is \(\operatorname{diag}(0,s_j,s_j,s_j)\) by
hypothesis, so the transverse amplitudes \(a^\epsilon_\mu\),
\(\mu=2,3,4\), have the decoupled energies
\(16s_j(p_{\epsilon,1})|a^\epsilon_\mu|^2\), while the longitudinal
amplitudes \(a^\epsilon_1\) cost nothing.  The constraint for a
transverse coarse component \(c_\mu\) involves only the amplitudes
\(a^\epsilon_\mu\), because its connectors vanish.  The longitudinal
constraint \(c_1=0\) does involve the transverse amplitudes through the
connectors of the paths with a step in direction \(\mu\), but it also
involves the longitudinal amplitudes through the straight weights
\(1+e^{ip_{\epsilon,1}}\), at least one of which is nonzero; since those
amplitudes have zero energy, \(c_1=0\) is enforced at no cost for every
choice of transverse amplitudes.  All amplitudes of weight zero are set
to zero, which by Lemma~\ref{lem:v33-ward} is permitted because the
minimum may be taken over all configurations rather than over the
slice.  Hence the minimum for \(c=e_\mu\) is the scalar constrained
minimum of \(16\sum_\epsilon s_j(p_{\epsilon,1})|a^\epsilon_\mu|^2\)
subject to \(\sum_\epsilon b_\epsilon a^\epsilon_\mu=1\), namely
\(16/\sum_\epsilon|b_\epsilon|^2/s_j(p_{\epsilon,1})\), which is (37.7).
For general \(c\) the componentwise minimizers add, so the kernel is
diagonal with equal transverse entries, and its longitudinal row and
column vanish by transversality.
\end{proof}

\begin{theorem}[Closed form and fixed point on the axes]
\label{thm:v33-closed-form}
For every level \(j\ge0\) and \(t\in\mathbb R\),
\[
 \boxed{\
 s_j(t)=\frac{4\sin^2(t/2)}{1-a_j\sin^2(t/2)},\qquad
 a_j=\frac23\bigl(1-4^{-j}\bigr).\ }                              \tag{37.8}
\]
In particular \(s_j(\pi)=12\cdot4^j/(4^j+2)\), so
\(s_1(\pi)=8\) and \(s_2(\pi)=32/3\) in agreement with (37.4), and
\[
 s_j\longrightarrow s^*(t)=\frac{4\sin^2(t/2)}{1-\frac23\sin^2(t/2)}
 =\frac{12\sin^2(t/2)}{2+\cos t}                                  \tag{37.9}
\]
uniformly on \(\mathbb R\), with \(\sup_t|s_j(t)-s^*(t)|\le24\cdot4^{-j}\).
\end{theorem}

\begin{proof}
Write \(S=\sin^2(t/2)\) and \(\sigma=\sin^2(t/4)\), so that
\(\sin^2(t/4+\pi/2)=\cos^2(t/4)=1-\sigma\) and \(\sigma(1-\sigma)=S/4\).
Assume (37.8) at level \(j\).  Then
\(s_j(t/2)=4\sigma/(1-a_j\sigma)\) and
\(s_j(t/2+\pi)=4(1-\sigma)/(1-a_j(1-\sigma))\), and the denominator of
(37.7) equals
\[
 \frac{(1-\sigma)(1-a_j\sigma)}{4\sigma}+\frac{\sigma(1-a_j+a_j\sigma)}{4(1-\sigma)}
 =\frac{1-(2+a_j)\sigma(1-\sigma)}{4\sigma(1-\sigma)}
 =\frac{1-\frac{2+a_j}{4}S}{S}.
\]
Hence \(s_{j+1}(t)=4S/(1-a_{j+1}S)\) with \(a_{j+1}=(2+a_j)/4\).  From
\(a_0=0\) this recursion gives \(a_j=\frac23(1-4^{-j})\), since
\(a_{j+1}-\frac23=\frac14(a_j-\frac23)\).  The values at \(t=\pi\) follow
from \(S=1\).  Since \(0\le S\le1\) and \(a_j\le2/3\), the denominators
are at least \(1/3\); the difference
\(s^*-s_j=4S^2(a^*-a_j)/((1-a^*S)(1-a_jS))\) is bounded by
\(4\cdot\frac23 4^{-j}\cdot9=24\cdot4^{-j}\), which proves uniform
geometric convergence.
\end{proof}

\begin{corollary}[Locality and the Maxwell coefficient]
\label{cor:v33-local}
The fixed point \(s^*\) is a rational function of \(\cos t\), analytic
in the strip \(|\operatorname{Im}t|<\operatorname{arccosh}2=\log(2+\sqrt3)
=1.3169\ldots\), with double zeros exactly at \(2\pi\mathbb Z\), and
\(s^*(t)=t^2+t^4/12+O(t^6)\).  Every finite-level kernel \(s_j\) is
analytic in the wider strip \(|\operatorname{Im}t|<\operatorname{arccosh}(2/a_j-1)\).
\end{corollary}

\begin{proof}
The denominator \(2+\cos t\) vanishes only where \(\cos t=-2\), i.e.\ at
\(t=\pi\pm i\operatorname{arccosh}2\) modulo \(2\pi\).  The numerator
\(12\sin^2(t/2)=6(1-\cos t)\) has double zeros at \(2\pi\mathbb Z\).
Expanding \(6(1-\cos t)/(2+\cos t)\) at \(t=0\) gives the stated series.
For finite \(j\) the denominator \(1-a_j\sin^2(t/2)=1-\frac{a_j}{2}(1-\cos t)\)
vanishes where \(\cos t=1-2/a_j\), which lies below \(-1\) since
\(a_j<1\); this gives the wider strip.
\end{proof}

\begin{remark}[Comparison with the straight rule]
\label{rem:v33-straight}
Under the straight rule, the aliases with transverse flips keep the full
weight \(|1+e^{ip_{\epsilon,\mu}}|^2=4\) at \(p_{\epsilon,\mu}=0\), so
(37.7) is replaced by a sum over all eight aliases with
\(\epsilon_\mu=0\), and the level-\(j\) axis kernel is no longer
determined by axis data.  The number of contributing aliases grows like
\(8^j\) and the sum of their inverse energies diverges like a
three-dimensional lattice sum of \(1/|n|^2\), which is the mechanism
behind the halving observed in (36.9).  The transverse average kills
every flipped alias exactly, which is why (37.7) closes.
\end{remark}

\subsection{Deeper levels off the axes: floating diagnostics}

\begin{record}[Averaged versus straight rule to level six, non-hashed]
\label{rec:v33-deep}
At the corner \((\pi,0,0,0)\) the transverse eigenvalues for levels
\(0,\ldots,6\) are
\[
\begin{array}{c|ccccccc}
 \hbox{rule}&0&1&2&3&4&5&6\\ \hline
 \hbox{straight}&4 & 2.1428571 & 1.0471614 & 0.50981238 & 0.2511057 & 0.12459744 & \hbox{---}\\
 \hbox{averaged}&4 & 8 & 10.666667 & 11.636364 & 11.906977 & 11.976608 & \hbox{---}
\end{array}                                                      \tag{37.10}
\]
The averaged values agree with (37.8) to all displayed digits, the
successive differences shrinking by the factor \(1/4\).  On the axis at
\(t=0.2,0.5,1,2\) the averaged ratio \(s_j/|q_1|^2\) at level four is
\(1.00666,\ 1.04237,\ 1.18013,\ 1.88752\), against the closed form
\(1/(1-a_4\sin^2(t/2))\) of (37.8), which gives
\(1.006663,\ 1.042369,\ 1.180128,\ 1.887523\).  For five pseudo-random momenta and the
corner, the largest relative change of the averaged kernel between
consecutive levels is
\[
 1,\ 0.436626,\ 0.176141,\ 0.0594146\qquad(j=1,\ldots,4),                              \tag{37.11}
\]
against \(0.691118,\ 0.607232,\ 0.557908,\ 0.532554\) for the straight rule, and the sampled
minimum of the cleared vertical block is
\[
 \hbox{averaged: }0.0838332,\ 0.127452,\ 0.151715,\ 0.194082,\qquad
 \hbox{straight: }0.117607,\ 0.0400899,\ 0.0162337,\ 0.00799183\qquad(j=1,\ldots,4).             \tag{37.12}
\]
These are floating diagnostics excluded from the canonical payload.
\end{record}

\begin{assessment}[A convergent Gaussian tower exists]
\label{ass:v33-fixed-point}
On every coordinate axis the averaged rule has an exactly solvable,
geometrically convergent, local Gaussian fixed point which preserves the
Maxwell coefficient.  The diagnostics indicate the same behaviour off
the axes with the same rate \(1/4\) and a vertical floor which does not
degenerate.  For the programme this means that the nested Gaussian tower
can be built on a blocking whose kernels converge, so that the archived
reference reset is unnecessary at the Gaussian level and the reset row
(14.17) can be replaced by a contraction toward the fixed point with
geometrically summable memory.  Whether this survives the interacting
vertices is the next question, not a Gaussian one.
\end{assessment}

\subsection{Executable provenance}

\begin{record}[V33 verifier]
\label{rec:v33-verifier}
The executable
\begin{center}
\path{scripts/verify_ym_postarchive_v33_averaged_blocking_kernel.py}
\end{center}
has \(407\) physical lines, \(18436\) bytes, and SHA--256
\[
\texttt{FFD55F0C5580182EC0CC2EAC3780F1571855EA3111DA0C053BCA65D3DECC2A79}.
                                                               \tag{37.13}
\]
Its canonical payload, which contains the Ward check, the level-one
census for both rules at all quarter-turn momenta, the corner kernels,
the rational-sequence ratios, the exact level-two corner kernels of
both rules, and the two regression identities \(15/7\) and
\(25823/24660\), has SHA--256
\[
\texttt{D0FB503B1E13A6976D200E18E59B3849CE0CE07AEB5F24B9DB46CC4A8722DF7B},
                                                               \tag{37.14}
\]
and the artifact
\path{artifacts/ym_postarchive_v33_averaged_blocking_kernel.json} has
SHA--256
\[
\texttt{5F509C26070EED8B046F574684E1DBA44179BE4E49F563AB5B170067272ECD3E}.
                                                               \tag{37.15}
\]
The exact part ran in \(530\) seconds and the diagnostics in
\(208\) seconds.
\end{record}

\subsection{Decision for V34}

\begin{decision}[Build the Gaussian fixed-point tower]
\label{dec:v33-v34}
V34 should:
\begin{enumerate}[label=\arabic*.]
\item symmetrize Definition~\ref{def:v33-average} over all orderings of
      the transverse steps, prove hypercubic covariance of the resulting
      kernel, and verify that the axis theory (37.7)--(37.9) is
      unchanged;
\item prove convergence of the averaged recursion off the axes: express
      the level-\(j\) propagator as an alias sum over \(\mathbb Z^4\)
      with the squared blocking weights and show that the transverse
      weights decay in every direction, so that the limit exists and is
      analytic in a strip;
\item state the Gaussian fixed-point tower as the replacement for the
      reset ledger and record which archived rows change; and
\item begin the vertex level: derive the tree-level blocking of the
      cubic and quartic banks through the constrained minimizer of
      Theorem~\ref{thm:v33-level-map}, since the one-loop shell
      coefficients of the fixed-point tower need the blocked vertices.
\end{enumerate}
The compulsory companion audit is due at V35, together with the
normalization item of Record~\ref{rec:v32-normalization}.
\end{decision}

\begin{assessment}[V33 acquisition]
\label{ass:v33}
V33 defines a transversely averaged dyadic blocking, proves its Ward
property and constrained level map, computes it exactly at level one on
all quarter-turn momenta and at level two on the corners, and proves
that on the coordinate axes the Gaussian kernel recursion is exactly
solvable with a local fixed point reached geometrically at rate
\(1/4\).  It thereby replaces the negative diagnosis of V32 by a
constructive Gaussian tower.
\end{assessment}

\begin{record}[V33 append record]
\label{rec:v33-append}
V33 is appended to the clean V32 whose installed record was
\[
\begin{array}{c|c}
\text{lines}&11540\\
\text{bytes}&452376\\
\text{SHA--256}&
\texttt{F1DE1DC4AA157BC4DF069FEDB4F34C3032565655348CB9AE35AC95F5695744D3}.
\end{array}                                                    \tag{37.16}
\]
After validation only V33 is to remain the active numeric YM note.
\end{record}

\begin{firewall}[Claim boundary after V33]
\label{fire:v33-final}
V33 proves Gaussian, tree-level statements: exact kernels at finitely
many momenta and an exact fixed point on the coordinate axes.  It proves
no off-axis fixed point, no one-loop shell coefficient, no interacting
contraction, no Haar sign, and no continuum statement.  The Gaussian
anomaly (14.16), the reset row (14.17) or its replacement, the compatible
interacting measure, regulator removal, Osterwalder--Schrader
reconstruction, physical clustering, and the Hamiltonian Yang--Mills
mass gap remain open.
\end{firewall}

% =============================================================================
% END APPENDED POST-TENTH-ARCHIVE V33 MATERIAL
% =============================================================================
% =============================================================================
% BEGIN APPENDED POST-TENTH-ARCHIVE V34 MATERIAL
% =============================================================================

\section{V34: symmetrized blocking, the exact propagator recursion, and the
fixed-point form of the Gaussian anomaly card}
\label{sec:v34}

V33 established an exactly solvable Gaussian fixed point on the coordinate
axes for the staircase eight-path average and noted that the staircase
ordering breaks permutation covariance off the axes.  This note first
symmetrizes the rule over all orderings of the transverse steps and proves
exact permutation covariance at the parity corners while leaving the axis
theory unchanged.  It then proves that the transverse pseudo-inverse of the
blocked kernel obeys a \emph{linear} alias recursion, the propagator
recursion, and verifies it exactly at levels one and two.  The linear form
is what makes the off-axis fixed point a convergent lattice sum.  Finally
the averaged Gaussian anomaly card (14.16)--(14.17) is restated on the
fixed-point tower, where the reset row is replaced by a geometrically
convergent memory.

\subsection{Symmetrized rule}

\begin{definition}[Symmetrized sixteen-path average]
\label{def:v34-symmetrized}
For a direction \(\mu\) and a transverse offset \(t\) with \(m=|t|_1\)
steps, let the paths \(\Pi_{\mu,t,\sigma}\) run through the \(m!\)
orderings \(\sigma\) of the transverse steps, each with weight
\(1/(8\,m!)\), so that every offset keeps total weight \(1/8\) and the
straight parts of (37.2) are unchanged.  There are \(1+3+6+6=16\) paths per
direction.  The straight rule and the staircase rule are the special cases
with one and eight paths.
\end{definition}

\begin{theorem}[Exact permutation covariance and axis invariance]
\label{thm:v34-covariance}
Let \(\widehat S^{(1)}_{\rm sym}(k)\) be the level-one kernel of the
symmetrized rule.  For every permutation \(\pi\) of the four directions and
every parity corner \(k\in\{0,\pi\}^4\),
\[
 \widehat S^{(1)}_{\rm sym}(\pi k)=P_\pi\widehat S^{(1)}_{\rm sym}(k)P_\pi^{\mathsf T},
                                                                 \tag{38.1}
\]
whereas the staircase kernel violates (38.1); an exact witness is the
corner \((\pi,\pi,\pi,\pi)\), where
\[
 \widehat S^{(1)}_{\rm sym}=\frac{72}{17}(3I-J),\qquad
 \widehat S^{(1)}_{\rm stair}=\begin{pmatrix}\tfrac{992}{243} & -\tfrac{32}{81} & -\tfrac{224}{243} & -\tfrac{224}{81} \\ -\tfrac{32}{81} & \tfrac{224}{27} & -\tfrac{160}{81} & -\tfrac{160}{27} \\ -\tfrac{224}{243} & -\tfrac{160}{81} & \tfrac{4064}{243} & -\tfrac{1120}{81} \\ -\tfrac{224}{81} & -\tfrac{160}{27} & -\tfrac{1120}{81} & \tfrac{608}{27}\end{pmatrix}.                 \tag{38.2}
\]
At every quarter-turn momentum with at most one nonzero component the
symmetrized and staircase kernels coincide; in particular
Theorems~\ref{thm:v33-axis}--\ref{thm:v33-closed-form} hold verbatim for
the symmetrized rule.  The other symmetrized corner kernels are
\[
 \widehat S^{(1)}_{\rm sym}(\pi,\pi,0,0)=\begin{pmatrix}\tfrac{576}{107} & -\tfrac{576}{107} & 0 & 0 \\ -\tfrac{576}{107} & \tfrac{576}{107} & 0 & 0 \\ 0 & 0 & 16 & 0 \\ 0 & 0 & 0 & 16\end{pmatrix},\qquad
 \widehat S^{(1)}_{\rm sym}(\pi,\pi,\pi,0)=\begin{pmatrix}\tfrac{720}{77} & -\tfrac{360}{77} & -\tfrac{360}{77} & 0 \\ -\tfrac{360}{77} & \tfrac{720}{77} & -\tfrac{360}{77} & 0 \\ -\tfrac{360}{77} & -\tfrac{360}{77} & \tfrac{720}{77} & 0 \\ 0 & 0 & 0 & 24\end{pmatrix}.       \tag{38.3}
\]
\end{theorem}

\begin{proof}
The set of symmetrized paths for direction \(\pi(\mu)\) at the permuted
cell is the image of the set for \(\mu\) under \(\pi\), with equal weights,
and the incidence grammar is permutation covariant; hence the constrained
minimum transforms by conjugation.  The executable verifies (38.1) exactly
for all \(24\) permutations at all \(16\) corners for the symmetrized rule
and records the first failing pair for the staircase rule.  On the axis
every connector carries the factor \(1-z_\mu=0\) for the three transverse
directions, and the straight parts agree for all three rules, so the
level-one kernels coincide there; the executable confirms equality at all
\(13\) quarter-turn momenta with at most one nonzero component.  The
recursion (37.7) then applies unchanged.
\end{proof}

\subsection{The exact propagator recursion}

For a Hermitian positive semidefinite transverse kernel \(\widehat S(k)\)
of rank three with \(\widehat Sq(k)=0\), write \(\Pi(k)=I-qq^*/|q|^2\) for
the orthogonal projector onto \(q(k)^\perp\) (and \(\Pi=I\) at the toron),
and let \(\widehat S^+\) be the Moore--Penrose inverse, which satisfies
\(\widehat S^+=\Pi\widehat S^+\Pi\).  Let \({\cal B}(k)\) be the
\(4\times64\) alias matrix of the blocking rule acting on normalized alias
amplitudes, and \(B_\epsilon(k)\) its \(4\times4\) alias blocks.

\begin{theorem}[Propagator recursion]
\label{thm:v34-propagator}
Let \(\widehat S'=\mathcal T[\widehat S]\) be the constrained level map of
Theorem~\ref{thm:v33-level-map} for any of the three rules, and let
\(k\ne0\).  Then
\[
 \boxed{\
 (\widehat S')^+(k)=\frac1{16}\,\Pi(k)\Bigl[\sum_\epsilon
 B_\epsilon(k)\,\widehat S^+(p_\epsilon)\,B_\epsilon(k)^*\Bigr]\Pi(k).\ }
                                                                 \tag{38.4}
\]
Consequently the transverse propagators \(G^{(j)}:=(\widehat S^{(j)})^+\)
obey the linear recursion \(G^{(j+1)}=\mathcal L[G^{(j)}]\) with
\(\mathcal L\) the map on the right of (38.4), and
\[
 G^{(j)}(k)=16^{-j}\,\Pi(k)\sum_{\epsilon^{(1)},\ldots,\epsilon^{(j)}}
 W_{\epsilon}(k)\,G^{(0)}(p_{\epsilon})\,W_{\epsilon}(k)^*\,\Pi(k),
                                                                 \tag{38.5}
\]
where the sum runs over the \(16^j\) alias chains, \(p_\epsilon\) is the
level-zero momentum of the chain, and \(W_\epsilon\) is the ordered product
of the alias blocks with the intermediate projectors.
\end{theorem}

\begin{proof}
In normalized amplitudes the level-\(j\) energy of a configuration is
\(16\sum_\epsilon(a^\epsilon)^*\widehat S(p_\epsilon)a^\epsilon\) and the
constraint is \(c=\sum_\epsilon B_\epsilon a^\epsilon\).  Decompose each
\(a^\epsilon\) into its component in \(\ker\widehat S(p_\epsilon)\), which
costs nothing, and its transverse part.  By the Ward property the kernel
components are mapped by the blocking onto multiples of \(q(k)\), and by
Lemma~\ref{lem:v33-ward} they can realize any multiple of \(q(k)\) at zero
cost.  Hence the constrained minimum equals the minimum of the transverse
energy subject to \(\Pi\sum_\epsilon B_\epsilon a^\epsilon_\perp=\Pi c\).
For a positive definite form \(16\widehat S_\perp\) on the transverse
amplitudes and a surjective constraint map \(\Pi{\cal B}_\perp\) onto
\(q^\perp\), the minimum is
\((\Pi c)^*\bigl[\Pi{\cal B}_\perp(16\widehat S_\perp)^{-1}{\cal B}_\perp^*\Pi\bigr]^{-1}(\Pi c)\),
where the inverse is taken on \(q^\perp\); this is the standard
Lagrange-multiplier formula.  Since \(\widehat S_\perp^{-1}\) on the
transverse space is \(\widehat S^+\), the bracket equals the right side of
(38.4), and the minimum is \(c^*\widehat S'c\) with \(\widehat S'\) the
inverse of the bracket on \(q^\perp\), i.e.\ the bracket is
\((\widehat S')^+\).  Iterating (38.4) gives (38.5).
\end{proof}

\begin{theorem}[Exact verification]
\label{thm:v34-propagator-check}
For all three rules and all fifteen nonzero corners, the level-one propagator
computed from (38.4) with \(G^{(0)}(p)=\Pi(p)/\lambda(p)\) and the
level-one kernel computed by the constrained minimum satisfy the
Moore--Penrose identities \(\widehat SG\widehat S=\widehat S\) and
\(G\widehat SG=G\) exactly.  For the symmetrized rule the same identities
hold exactly at level two at all fifteen nonzero corners, with \(G^{(1)}\) obtained
as the exact transverse pseudo-inverse \((\widehat S^{(1)}+P_q)^{-1}-P_q\)
of the stored level-one kernels.
\end{theorem}

\begin{proof}
All quantities are Gaussian rational at the corners.  The executable forms
both sides independently, the kernel through the incidence and alias
routes and the propagator through (38.4), and compares the two products
entrywise.  The formula \((\widehat S+P_q)^{-1}-P_q\) is the pseudo-inverse
of a rank-three Hermitian matrix with kernel \(\mathbb Cq\), since
\(\widehat S+P_q\) is invertible and block diagonal with respect to
\(\mathbb Cq\oplus q^\perp\).
\end{proof}

\begin{remark}[The toron]
\label{rem:v34-toron}
At \(k=0\) the coarse form vanishes identically, \(\Pi=I\), and the
constraint map is not surjective, so (38.4) does not apply: its right side
is a nonzero matrix while \((\widehat S')^+=0\).  The toron never occurs as
an alias of a nonzero momentum, so the recursion (38.5) is unaffected for
\(k\ne0\).
\end{remark}

\begin{remark}[Why the linear form matters]
\label{rem:v34-linear}
Equation (38.5) writes the level-\(j\) propagator as a positive lattice sum
over \(16^j\) aliases.  On the axis the weights reduce to the scalar
products of Theorem~\ref{thm:v33-axis} and the sum is the finite cosecant
sum behind (37.8).  Off the axis, the straight weights of the averaged
rules contain the factor \(\prod_{\nu\ne\mu}(1+e^{ip_{\epsilon,\nu}})/2\)
at every level, whose composed modulus is
\(\prod_{\nu\ne\mu}|\sin(k_\nu/2)/(2^j\sin(p_{\epsilon,\nu}/2))|\), of order
\(\prod_{\nu\ne\mu}|k_\nu+2\pi n_\nu|^{-1}\) for the alias \(n\in\mathbb Z^4\).
Together with \(G^{(0)}\sim4^{j}/|k+2\pi n|^2\) and the prefactor
\(16^{-j}\), the straight contribution to (38.5) is dominated by
\(\sum_{n\in\mathbb Z^4}\prod_{\nu\ne\mu}(k_\nu+2\pi n_\nu)^{-2}|k+2\pi n|^{-2}\),
which converges absolutely.  The connector contributions carry an extra
factor \(1-z_\mu\) and one fewer transverse cosine; their absolute
convergence is the remaining estimate, recorded as the first V35 target.
The straight rule lacks the transverse cosines and its sum diverges like
\(\sum_{n\in\mathbb Z^3}|n|^{-2}\), which is the halving of (36.9).
\end{remark}

\subsection{Exact level two for the symmetrized rule}

\begin{record}[Symmetrized corner kernels at level two]
\label{rec:v34-level-two}
The symmetrized level-two kernels at the sixteen corners are Hermitian,
exactly transverse, of rank three off the toron, and
\[
 \widehat S^{(2)}_{\rm sym}(\pi,0,0,0)=\operatorname{diag}(0,\tfrac{32}{3},\tfrac{32}{3},\tfrac{32}{3}),\qquad
 \widehat S^{(2)}_{\rm sym}(\pi,\pi,\pi,\pi)=\begin{pmatrix}\tfrac{193536}{11621} & -\tfrac{64512}{11621} & -\tfrac{64512}{11621} & -\tfrac{64512}{11621} \\ -\tfrac{64512}{11621} & \tfrac{193536}{11621} & -\tfrac{64512}{11621} & -\tfrac{64512}{11621} \\ -\tfrac{64512}{11621} & -\tfrac{64512}{11621} & \tfrac{193536}{11621} & -\tfrac{64512}{11621} \\ -\tfrac{64512}{11621} & -\tfrac{64512}{11621} & -\tfrac{64512}{11621} & \tfrac{193536}{11621}\end{pmatrix},   \tag{38.6}
\]
\[
 \widehat S^{(2)}_{\rm sym}(\pi,\pi,0,0)=\begin{pmatrix}\tfrac{94590720}{13496023} & -\tfrac{94590720}{13496023} & 0 & 0 \\ -\tfrac{94590720}{13496023} & \tfrac{94590720}{13496023} & 0 & 0 \\ 0 & 0 & \tfrac{256}{11} & 0 \\ 0 & 0 & 0 & \tfrac{256}{11}\end{pmatrix}.          \tag{38.7}
\]
\end{record}

\subsection{The Gaussian anomaly card on the fixed-point tower}

The archived ledger (14.10) attributes the inverse coupling to the nested
Gaussian determinant, the reset memory, the nonlinear rows, and the
boundary row, and ADMC demands (14.16) with (14.17).  With a convergent
Gaussian tower the reset is never performed.

\begin{theorem}[Fixed-point form of the Gaussian card]
\label{thm:v34-card}
Let \(\widehat S^{(j)}\) be the kernels of a blocking rule whose Gaussian
recursion converges: \(\sup_k\|\widehat S^{(j)}(k)-\widehat S^*(k)\|\le
C\rho^j\) with \(\rho<1\), and suppose the per-shell Gaussian coefficient
\(b_{{\rm G},j}\) of the nested tower is a functional of the level-\(j\)
kernel and vertices which is Lipschitz in the kernel, uniformly along the
tower.  Then \(b_{{\rm G},j}=b^*+O(\rho^j)\) for the fixed-point value
\(b^*\), the sum \(\sum_j(b_{{\rm G},j}-b^*)\) converges absolutely, and
(14.16) holds with \(b_{\rm YM}\) replaced by \(b^*\) and a bounded
remainder; the reset row (14.17) is vacuous.  The Gaussian half of ADMC
then reduces to the single identity \(b^*=b_{\rm YM}\).
\end{theorem}

\begin{proof}
Lipschitz continuity gives \(|b_{{\rm G},j}-b^*|\le LC\rho^j\); summation
of the geometric series gives the bounded remainder in (14.16).  No
reference reset is inserted, so \(R_N^{\rm reset}=0\).
\end{proof}

\begin{remark}[Status of the hypotheses]
\label{rem:v34-hypotheses}
For the symmetrized rule the kernel convergence is proved on the axes with
\(\rho=1/4\) (Theorem~\ref{thm:v33-closed-form}) and observed off the axes
with the same rate (Record~\ref{rec:v33-deep}); its proof off the axes is
the alias-sum estimate of Remark~\ref{rem:v34-linear}.  The Lipschitz
hypothesis concerns the blocked vertices, which have not yet been
constructed; V32's Record~\ref{rec:v32-normalization} also leaves the
identification of \(b^*\) with \(b_{\rm YM}\) as a bookkeeping item.  The
theorem therefore replaces two archived obstructions, an unabsorbable
constant class and a linear reset row, by two concrete tasks: an alias-sum
estimate and a one-loop computation at a fixed local kernel.
\end{remark}

\subsection{Executable provenance}

\begin{record}[V34 verifier]
\label{rec:v34-verifier}
The executable
\begin{center}
\path{scripts/verify_ym_postarchive_v34_symmetrized_propagator_recursion.py}
\end{center}
has \(325\) physical lines, \(13802\) bytes, and SHA--256
\[
\texttt{AEEDA53CC20B8AB0CBA51CD0E5A2177CAFE7300A6AF592A24A27BAC34082AF09}.
                                                               \tag{38.8}
\]
Its canonical payload has SHA--256
\[
\texttt{F280F43FEB5E97F2A69B3B5473E9CDC8AC98EE37367FAE2EC32E84EF8D38C69A},
                                                               \tag{38.9}
\]
and the artifact
\path{artifacts/ym_postarchive_v34_symmetrized_propagator_recursion.json}
has SHA--256
\[
\texttt{29529A94C6601A0E18B7E05E8B6707DA3EE40F9727872733F6E40D21F64E63A6}.
                                                               \tag{38.10}
\]
The run took \(341\) seconds; every check is exact.
\end{record}

\subsection{Decision for V35}

\begin{decision}[Compulsory audit and the alias-sum estimate]
\label{dec:v34-v35}
V35 is the compulsory fifth-version companion audit of the current SM and
NS notes.  It must also recover, from the tenth archive, the colour and
background normalization relating \(\langle F_e\rangle\) to
\(b_{\rm YM}\).  Its Yang--Mills content should be the absolute
convergence of the connector part of (38.5) for the symmetrized rule,
which would prove the off-axis fixed point and the rate \(1/4\); if the
connector sum resists a direct bound, the estimate may be organized by the
number of connector insertions along the chain, each of which carries a
factor \(|1-z_\mu|\le2\) and removes one transverse cosine.
\end{decision}

\begin{assessment}[V34 acquisition]
\label{ass:v34}
V34 symmetrizes the averaged blocking and proves its exact permutation
covariance, proves that the blocked propagator obeys a linear alias
recursion and verifies it exactly at levels one and two, and restates the
Gaussian anomaly card on the fixed-point tower, where the archived reset
row disappears.  The remaining Gaussian task is an explicit lattice-sum
estimate; the remaining one-loop task is the construction of blocked
vertices.
\end{assessment}

\begin{record}[V34 append record]
\label{rec:v34-append}
V34 is appended to the clean V33 whose installed record was
\[
\begin{array}{c|c}
\text{lines}&11963\\
\text{bytes}&473284\\
\text{SHA--256}&
\texttt{B417E58E5CF6BFBE9F3FB0DD21106355674D525AE9BF8464EFD7AA580D0134F9}.
\end{array}                                                    \tag{38.11}
\]
After validation only V34 is to remain the active numeric YM note.
\end{record}

\begin{firewall}[Claim boundary after V34]
\label{fire:v34-final}
V34 proves exact tree-level statements and a conditional reformulation.
It proves no off-axis fixed point, no one-loop shell coefficient, no
Lipschitz bound, no identification \(b^*=b_{\rm YM}\), no Haar sign, and no
continuum statement.  The Gaussian anomaly, the interacting measure,
regulator removal, Osterwalder--Schrader reconstruction, physical
clustering, and the Hamiltonian Yang--Mills mass gap remain open.
\end{firewall}

% =============================================================================
% END APPENDED POST-TENTH-ARCHIVE V34 MATERIAL
% =============================================================================
% =============================================================================
% BEGIN APPENDED POST-TENTH-ARCHIVE V35 MATERIAL
% =============================================================================

\section{V35: companion audit, the alias Ward identity, and the off-axis
Gaussian fixed point}
\label{sec:v35}

V34 left two Gaussian tasks: the absolute convergence of the connector part
of the alias sum (38.5), which was to prove the off-axis fixed point, and
the recovery of the colour and background normalization relating the
compact response to \(b_{\rm YM}\).  This note first performs the compulsory
fifth-version companion audit.  It then proves that the alias blocks of the
symmetrized rule satisfy an exact Ward identity, \(B_\epsilon(k)q(p_\epsilon)=q(k)\),
and that their Gram sum has a closed form of four terms.  These two facts
make the connector estimate unnecessary: the linear propagator map
\(\mathcal L\) of Theorem~\ref{thm:v34-propagator} contracts the transverse
projector by the exact factor \(\rho(k)=\tfrac14+\tfrac1{18}\max_\mu|q_\mu(k)|^2\le\tfrac{17}{36}\),
uniformly on the torus, and the first blocking step changes the propagator
by a bounded kernel.  Together these give the off-axis Gaussian fixed point
of the symmetrized rule with geometric convergence, which is the first
hypothesis of Theorem~\ref{thm:v34-card}.  The normalization item is reduced
to one explicit exact identity and one remaining question.

\subsection{Mandatory companion audit}

The active companion notes in \path{latex_notes} were inspected read-only:
\[
\begin{array}{c|r|r|l}
\text{source}&\text{lines}&\text{bytes}&\text{SHA--256}\\ \hline
\text{SM V93}&19339&795468&\texttt{F907BADAAFA9EA70378F0DAA941B03690F76578815ABE12058E01323750CB518}\\
\text{NS V26}&5319&278283&\texttt{82C887F8C2C69E4F28FAD7C3DC8DE5B9099CB5A4BF431EBA1FFF3E91935978AD}\\
\text{GRH V60}&17282&580602&\texttt{E090E736AE1C43D60756DB3E9C7929E0F2A179905B74630ADBE85266036737F4}\\
\text{YM-par V5}&1351&54174&\texttt{306F1BB84CFA8A8D47EA569EC17E9545F9867C8D32B548EC9D9C444B4F3C0512}
\end{array}                                                    \tag{39.1}
\]
Here SM V93 is \path{Renewal_SM_Einstein_Continuum_Research_Notes_V93.tex}
of the nineteenth SM cycle, GRH V60 is
\path{Renewal_GRH_Cofinal_Visibility_Attack_V60.tex}, and YM-par V5 is
\path{Renewal_YM_Parallel_Research_Notes_V5.tex}.  The NS source is
unchanged since the V20 audit.

\begin{audit}[SM V88--V93]
\label{aud:v35-sm}
The nineteenth SM cycle has moved to exact classical dynamics: V88--V89
solve the plane-symmetric Einstein equations for a time-symmetric standing
wave to second order, V91 confirms V89 by the full plane-symmetric system,
V92 proves a focusing theorem at fourth order and refutes a homogeneous
two-scale model, and V93 treats finite-mode balanced data and the Friedmann
limit.  These are classical solutions of the vacuum and null-dust equations
in synchronous gauge; they contain no lattice gauge field, no blocking
kernel, and no Gaussian measure.  Nothing is importable into the present
Gaussian tower.  The SM eighteenth cycle's link-blocking results, audited at
V25, remain the relevant SM inputs; the V32--V34 findings sharpen the
routing constraint recorded there.  The ordered-product link law of SM
V96 and its exact projective family of SM V99 linearize to the straight
rule of V32, whose Gaussian tower has no uniformly conditioned fixed point
(the kernel halves per level, (36.9)), whereas the transversely averaged
rules of V33--V34 converge.  A later projective Gaussian construction must
therefore introduce the transverse average at the Gaussian level even if
the nonlinear link law stays the ordered product.
\end{audit}

\begin{audit}[NS V26, GRH V60, and the parallel YM note]
\label{aud:v35-others}
NS V26 is unchanged.  GRH V60 proves a degree-free exclusion schema for
Pick expressions and closes a two-pair band; its objects are polynomial
inequalities in gamma-function variables and share no operator with the
present programme.  The parallel Yang--Mills note V5 reached its declared
failure stop: the energy--entropy tightness route for global integer charge
fails under tensorization, and its terminal gate recommends retaining
compact flux and toron labels exactly and assembling the continuum through
quasi-local observables.  It constructs no blocking kernel and no Gaussian
tower, so there is no duplication with V32--V35; its conclusion is
consistent with the present direction, which keeps the one-colour
commuting-toron background exact and builds the Gaussian tower locally.
\end{audit}

\begin{decision}[V35 audit decision]
\label{dec:v35-audit}
No estimate is imported.  The direction fixed by Decision~\ref{dec:v34-v35}
stands, with the connector estimate replaced by the exact Ward and Gram
identities proved below.  The next compulsory audit is V40, which is also
the ten-version sweep of the other programmes' notes.
\end{decision}

\subsection{What the archives fix about the normalization}

\begin{record}[The archived chain from the response density to \(b_{\rm YM}\)]
\label{rec:v35-normalization}
The eighth archive normalizes the adjoint Casimir by
\(-\operatorname{tr}_{\rm ad}(\operatorname{ad}X\operatorname{ad}Y)=C_A\langle X,Y\rangle\)
(its (37.1)) and the background action by \(S_{\rm bg}=\frac1{4u}\int|\overline F|^2\)
(its (37.2)), and proves that one upward dyadic shell shifts \(1/u\) by
\(b_{\rm YM}=\frac{22C_A}{3(4\pi)^2}\log2\) (its (37.11)).  On the lattice
side it uses one colour component and the parity-cell incidence blocks
\(A=\frac13D_R^*D_R\) (its (42.7)); the one-loop response of the quotient
along a unit direction \(X\in\mathfrak{su}(3)\) is
\(D^2\Gamma_M(0)[aX,aX]=\frac12\operatorname{Tr}(GQ)+\frac13\operatorname{Tr}(GRGR)\)
(its (42.12)), where the coefficient \(\frac13=-\frac12\cdot\frac19\operatorname{Tr}_{\rm colour}(\operatorname{ad}X)^2\)
uses \(\operatorname{Tr}_{\rm colour}(\operatorname{ad}X)^2=-6\) (its (42.11)).
The transverse coarse kernel of this response has the low-angle expansion
\(\gamma_\infty(q)=\kappa_{\rm tree}q^2+O(q^4)\) (its (42.19)--(42.20)),
the flat coarse Schur reference is
\[
 s_0(q)=\frac{t(t+8)(t+24)}{9(t^2+18t+16)},\qquad t=\sin^2(q/2),\qquad
 s_0(q)=\tfrac13q^2+O(q^4)                                    \tag{39.2}
\]
(its (42.22)--(42.23)), and the downward pure-Wilson shift of the
unit-Wilson inverse coupling is
\(\delta\beta^\downarrow_{\rm tree}=\lim_{q\to0}\gamma_\infty(q)/s_0(q)=3\kappa_{\rm tree}\),
with the matching condition \(\delta\beta^\downarrow_{\rm tree}=-b_{\rm YM}\)
(its (42.24)--(42.25)).  The ninth archive defines the Ward-recombined
density \(W_{j,t}(k)=\frac32\partial_q^2|_{q=0}\Omega_{j,t}(k,q)\) (its
(16.5)--(16.6)).  Since \(\int\partial_q^2\Omega\,dk=\partial_q^2\gamma_\infty|_0=2\kappa_{\rm tree}\),
\[
 \int W\,dk=\tfrac32\cdot2\kappa_{\rm tree}=3\kappa_{\rm tree}
 =\delta\beta^\downarrow_{\rm tree},                          \tag{39.3}
\]
i.e.\ the factor \(\frac32\) in the Ward density is exactly the tree
normalization \(1/(\frac12s_0''(0))=3\), and \(\mathfrak b\gamma_j=\langle W_j\rangle\)
is the shift of the unit-Wilson \(\beta=1/u\) per downward block in the
one-colour normalization \(A=\frac13D_R^*D_R\).  The identification of
\(\beta\) with \(1/u\) is the eighth archive's (37.2) with the unit Wilson
action; no further colour or background factor enters between (42.12) and
(37.11).
\end{record}

\begin{theorem}[The compact kernels are the colour-traced archived kernels]
\label{thm:v35-colour-factor}
At the quarter-turn axis momenta \(k=(\pi/2,0,0,0)\) and \((\pi,0,0,0)\),
the compact straight-rule level-one kernel of V32 equals exactly three
times the archived flat coarse Schur reference (39.2):
\[
 \widehat S^{(1)}_{\rm straight,22}(\pi/2,0,0,0)=\frac{833}{606}=3\cdot\frac{833}{1818},\qquad
 \widehat S^{(1)}_{\rm straight,22}(\pi,0,0,0)=\frac{15}{7}=3\cdot\frac57.
                                                                 \tag{39.4}
\]
Hence the compact V32--V34 kernels are the one-colour archived kernels
multiplied by \(3=-\frac12\operatorname{Tr}_{\rm colour}(\operatorname{ad}X)^2/3\),
the compact propagators are one third of the archived ones, and any one-loop
coefficient formed from compact kernels with archived vertices must carry the
factor \(3\) explicitly.
\end{theorem}

\begin{proof}
The V35 executable evaluates the straight-rule incidence route of V32 at
\(z=(i,1,1,1)\), \((-1,1,1,1)\) and \((-i,1,1,1)\) in exact Gaussian
rationals and compares with \(3s_0\) at \(t=\frac12,1,\frac12\); all three
transverse entries agree and the longitudinal entry vanishes.  The archived
blocks (42.7) carry the explicit factor \(\frac13\) in front of the
colour-traced incidence Gram matrix \(D_R^*D_R\), which is the compact
reduced block; the factor \(3\) is therefore the colour trace, not a
background or cell-volume convention.
\end{proof}

\begin{firewall}[What remains open in the normalization]
\label{fire:v35-normalization}
The compact response \(F_e(k)\) of (9.11) has the same two-term structure
as the archived (42.12), with colour weights \(a_{\rm W}=\frac13\), \(a_*=3\),
and the V60 oracle returns its second external-momentum derivative.  If the
oracle's \(\Omega\) carries the archived normalization, then
\(\mathfrak b\gamma^{(e)}=\frac32\langle F_e\rangle=\delta\beta^\downarrow_e\)
and the diagnostic values \(334.06\) and \(-138.21\) of (36.12) would be
direct shifts of \(\beta\).  Their size against \(b_{\rm YM}=0.0966\)
(\(C_A=3\), \(L=2\)) shows that at least one of two things is not yet
accounted for: the quadrature error of the \(256\)-point quarter-turn average
against the torus integral for a sharply peaked integrand, or a frame factor
in the oracle (the alias normalization \(Va=4a^\epsilon\) of
Lemma~\ref{lem:v33-ward} contributes \(16\) per quadratic form).  No
identification of \(\langle F_e\rangle\) with \(b_{\rm YM}\) is asserted; the
open bookkeeping item is now this single explicit question.
\end{firewall}

\subsection{Closed form of the alias blocks and the alias Ward identity}

Throughout, \(k\) is the coarse momentum, \(p_\epsilon=(k+2\pi\epsilon)/2\)
with \(\epsilon\in\{0,1\}^4\) are its aliases, \(w=e^{ip_\epsilon}\)
componentwise so that \(w_\mu^2=e^{ik_\mu}=z_\mu\), and
\(q(p)_\mu=e^{ip_\mu}-1\).  The alias block \(B_\epsilon(k)\) is the
\(4\times4\) matrix of (37.2) for the symmetrized rule of
Definition~\ref{def:v34-symmetrized}, acting on normalized alias amplitudes.

\begin{lemma}[Closed form]
\label{lem:v35-closed-form}
For the symmetrized rule,
\[
 B_\epsilon(k)_{\mu\mu}=\frac18\prod_{\nu=1}^4(1+w_\nu),\qquad
 B_\epsilon(k)_{\mu\nu}=(1-w_\mu^2)\,c_{\mu\nu}(w),\quad
 c_{\mu\nu}(w)=\frac{7+2(w_\rho+w_{\rho'})+w_\rho w_{\rho'}}{24}\quad(\mu\ne\nu),
                                                                 \tag{39.5}
\]
where \(\{\rho,\rho'\}\) is the complement of \(\{\mu,\nu\}\).  In
particular \(B_\epsilon(k)=2I\) at \(p_\epsilon=0\), \(c_{\mu\nu}=\frac12\) there,
and the linear part of \(B-2I\) in the variables \(u_\nu=w_\nu-1\) is
\(\operatorname{diag}(\sum_\nu u_\nu)\) on the diagonal and \(-u_\mu\) in
position \((\mu,\nu)\).
\end{lemma}

\begin{proof}
The straight term of (37.2) sums \(e^{ip\cdot t}(1+e^{ip_\mu})\) over the
eight offsets \(t\) with \(t_\mu=0\), each of total weight \(\frac18\) after
the orderings are summed, and factorizes as stated.  For the connector
term fix \(\nu\ne\mu\) and an offset set \(D\ni\nu\), \(D\not\ni\mu\), of
size \(m\).  The step in direction \(\nu\) starts at the site \(1_P\) where
\(P\subset D\setminus\{\nu\}\) is the set of steps preceding it; among the
\(m!\) orderings, \(|P|!\,(m-1-|P|)!\) have this preceding set.  Summing
the phases \(e^{ip\cdot1_P}\) with weight \(1/(8m!)\) over the four sets
\(D\) gives
\(\frac18\bigl[1+\frac12(1+w_\rho)+\frac12(1+w_{\rho'})+\frac13(1+w_\rho w_{\rho'})+\frac16(w_\rho+w_{\rho'})\bigr]\),
which is \(c_{\mu\nu}\).  The V35 executable constructs \(B\) from the
sixteen-path enumeration of the V34 executable as a matrix of Laurent
polynomials in \(w\) with rational coefficients and verifies (39.5) as a
polynomial identity; it also verifies \(B(1)=2I\) and the stated linear
part after the substitution \(w=1+u\).
\end{proof}

\begin{lemma}[Alias Ward identity]
\label{lem:v35-ward}
For every coarse momentum and every alias,
\[
 B_\epsilon(k)\,q(p_\epsilon)=q(k).                              \tag{39.6}
\]
\end{lemma}

\begin{proof}
Write \(a,b,c\) for the three \(w_\nu\) with \(\nu\ne\mu\).  By (39.5),
the \(\mu\)-th component of \(Bq(p_\epsilon)\) is
\(\frac18(1+w_\mu)(1+a)(1+b)(1+c)(w_\mu-1)+(1-w_\mu^2)\sum_\nu c_{\mu\nu}(w_\nu-1)\),
i.e.\ \((w_\mu^2-1)f(a,b,c)\) with
\[
 f=\tfrac18(1+a)(1+b)(1+c)-\tfrac1{24}\bigl[((b+2)(c+2)+3)(a-1)+((a+2)(c+2)+3)(b-1)+((a+2)(b+2)+3)(c-1)\bigr].
\]
Expanding, the bracket equals \(3abc+3(ab+bc+ca)+3(a+b+c)-21\), so
\(f=\frac18[1+s_1+s_2+s_3]-\frac18[s_3+s_2+s_1-7]=1\) with \(s_j\) the
elementary symmetric functions.  Hence the component is \(w_\mu^2-1=q(k)_\mu\).
The executable verifies (39.6) as a polynomial identity.
\end{proof}

Lemma~\ref{lem:v35-ward} is the exact form of the Ward property used in the
proof of Theorem~\ref{thm:v34-propagator}: the blocking maps the kernel
vector \(q(p_\epsilon)\) of the fine transverse form onto \(q(k)\) itself, with
coefficient one for every alias.

\begin{lemma}[Alias Gram matrix]
\label{lem:v35-gram}
Let \(S(k)=\sum_\epsilon B_\epsilon(k)B_\epsilon(k)^*\).  Then, with
\(q=q(k)\), \(\mathbb 1=(1,1,1,1)^{\mathsf T}\), and \(D_q=\operatorname{diag}(|q_\mu|^2)\),
\[
 S(k)=4I+\tfrac89D_q-q\mathbb 1^*-\mathbb 1q^*+\tfrac{53}{18}\,qq^*.  \tag{39.7}
\]
Equivalently \(S_{\mu\mu}=4+\frac{29}6|q_\mu|^2=\frac{41}3-\frac{29}3\cos k_\mu\)
and \(S_{\mu\nu}=\frac{89}{18}-\frac{71}{18}z_\mu-\frac{71}{18}\bar z_\nu+\frac{53}{18}z_\mu\bar z_\nu\)
for \(\mu\ne\nu\).
\end{lemma}

\begin{proof}
The sixteen aliases are the sign patterns \(w\mapsto(\pm w_1,\ldots,\pm w_4)\)
of one square root \(w\) of \(z\).  Each entry of \(B_\epsilon B_\epsilon^*\)
is a Laurent polynomial in \(w\) (with \(\bar w=w^{-1}\) on the torus), and
the sum over the sign patterns retains exactly the monomials of even degree
in every \(w_\nu\), which are Laurent polynomials in \(z=w^2\).  The
executable performs this sum exactly and compares with (39.7) as a
polynomial identity.  The two displayed forms agree because
\(-q_\mu-\bar q_\mu=|q_\mu|^2\) on the diagonal.
\end{proof}

\subsection{The contraction}

Recall from Theorem~\ref{thm:v34-propagator} the linear map on transverse
kernel-valued functions,
\[
 \mathcal L[X](k)=\frac1{16}\,\Pi(k)\Bigl[\sum_\epsilon B_\epsilon(k)X(p_\epsilon)B_\epsilon(k)^*\Bigr]\Pi(k),
                                                                 \tag{39.8}
\]
so that \(G^{(j+1)}=\mathcal L[G^{(j)}]\) for \(k\ne0\).  It is positive:
\(X\le Y\) pointwise in the semidefinite order implies \(\mathcal L[X]\le\mathcal L[Y]\).

\begin{lemma}[Inner projectors are invisible]
\label{lem:v35-inner}
For every transverse Hermitian \(X\) (i.e.\ \(X(p)q(p)=0\)) and \(k\ne0\),
\[
 \Pi(k)B_\epsilon\Pi(p_\epsilon)X(p_\epsilon)\Pi(p_\epsilon)B_\epsilon^*\Pi(k)
 =\Pi(k)B_\epsilon X(p_\epsilon)B_\epsilon^*\Pi(k),\qquad
 \mathcal L[\Pi](k)=\frac1{16}\bigl[4\Pi(k)+\tfrac89\Pi(k)D_q\Pi(k)\bigr].   \tag{39.9}
\]
\end{lemma}

\begin{proof}
\(\Pi(p_\epsilon)=I-q_\epsilon q_\epsilon^*/\lambda(p_\epsilon)\) with
\(q_\epsilon=q(p_\epsilon)\), and \(\Pi(k)B_\epsilon q_\epsilon=\Pi(k)q(k)=0\) by
Lemma~\ref{lem:v35-ward}.  Hence the extra terms vanish.  With \(X=\Pi\),
\(\mathcal L[\Pi]=\frac1{16}\Pi S\Pi\), and by (39.7) the three terms of
\(S\) containing \(q\) or \(q^*\) are annihilated by the outer projectors.
\end{proof}

\begin{theorem}[Exact contraction of the transverse projector]
\label{thm:v35-contraction}
For every \(k\ne0\),
\[
 \mathcal L[\Pi](k)\le\rho(k)\,\Pi(k),\qquad
 \rho(k):=\frac14+\frac1{18}\max_\mu|q_\mu(k)|^2\le\rho:=\frac{17}{36}.   \tag{39.10}
\]
On the coordinate axes \(\mathcal L[\Pi]=\frac14\Pi\) exactly.  Consequently
\(\mathcal L^j[\Pi]\le\rho^j\Pi\) for all \(j\), and for every transverse
Hermitian \(X\) with \(\|X(p)\|\le c\) for all \(p\ne0\),
\(\|\mathcal L^j[X](k)\|\le c\rho^j\).
\end{theorem}

\begin{proof}
For a unit vector \(v\perp q(k)\), (39.9) gives
\(v^*\mathcal L[\Pi]v=\frac14+\frac1{18}\sum_\mu|q_\mu|^2|v_\mu|^2\le\frac14+\frac1{18}\max_\mu|q_\mu|^2\),
and \(|q_\mu|^2=4\sin^2(k_\mu/2)\le4\).  On the axis \(k=(t,0,0,0)\) the
transverse vectors have \(v_1=0\) and the sum vanishes, reproducing the
alias weights \(4\cos^2(t/4)+4\sin^2(t/4)=4\) of (37.7).  Positivity of
\(\mathcal L\) and induction give \(\mathcal L^j[\Pi]\le\rho^j\Pi\); for \(X\)
with \(-c\Pi\le X\le c\Pi\) pointwise, \(-c\rho^j\Pi\le\mathcal L^j[X]\le c\rho^j\Pi\).
The executable checks (39.9) exactly at the fifteen nonzero parity corners
against the V34 propagator step applied to \(\Pi\), and samples (39.10) in
floating arithmetic over the torus.
\end{proof}

\begin{remark}[Why the connector estimate is not needed]
\label{rem:v35-no-connector}
Remark~\ref{rem:v34-linear} organized the alias sum (38.5) by the number of
connector insertions and asked for absolute convergence.  The identities
(39.6)--(39.7) show that the connector parts are not small individually:
the raw Gram sum \(\frac1{16}\|S\|\) exceeds \(1\) on most of the torus, and
its excess is concentrated on the longitudinal direction \(q(k)\), where
the Ward identity places it.  The outer projector removes it exactly,
leaving the straight-rule weight \(4\) plus the bounded transverse remainder
\(\frac89\Pi D_q\Pi\).  The sharp rate on the axis is \(\frac14\); the
uniform rate \(\frac{17}{36}\) is attained in the limit of momenta with a
transverse component at \(\pi\).
\end{remark}

\subsection{The first step is a bounded correction}

\begin{lemma}[Bounded first step]
\label{lem:v35-first-step}
Let \(D_1(p):=G^{(1)}(p)-G^{(0)}(p)\) for \(p\ne0\), where
\(G^{(0)}(p)=\Pi(p)/\lambda(p)\) and \(G^{(1)}=\mathcal L[G^{(0)}]\).  Then
\[
 d_1:=\sup_{p\in\mathbb T^4\setminus\{0\}}\|D_1(p)\|<\infty.        \tag{39.11}
\]
\end{lemma}

\begin{proof}
By Lemma~\ref{lem:v35-inner} with \(X=G^{(0)}\),
\(G^{(1)}(p)=\frac1{16}\sum_\epsilon\Pi(p)B_\epsilon B_\epsilon^*\Pi(p)/\lambda(p_\epsilon)\).
For \(\epsilon\ne0\) some component of \(p_\epsilon\) lies in
\([\pi/2,3\pi/2]\), so \(\lambda(p_\epsilon)\ge2\), and
\(0\le\frac1{16}\sum_{\epsilon\ne0}\Pi B_\epsilon B_\epsilon^*\Pi/\lambda(p_\epsilon)\le\frac1{32}\Pi S\Pi\le\frac{17}{72}\Pi\)
by (39.9)--(39.10).  It remains to bound the alias \(\epsilon=0\), whose
block we write \(B_0=2I+E\) with \(E=E(p)\) analytic and \(E(0)=0\)
(Lemma~\ref{lem:v35-closed-form}).  Put \(\lambda_{1/2}=\lambda(p/2)\).  Then
\[
 T(p):=\Pi\Bigl[\frac{B_0B_0^*}{16\lambda_{1/2}}-\frac{I}{\lambda(p)}\Bigr]\Pi
 =\frac{\Pi N\Pi}{16\lambda_{1/2}\lambda(p)},\qquad
 N=4\bigl(\lambda(p)-4\lambda_{1/2}\bigr)I+\lambda(p)\bigl[2(E+E^*)+EE^*\bigr].
\]
The function \(T\) is continuous on \(\mathbb T^4\setminus\{0\}\), so it
suffices to show that it is bounded near \(p=0\), where
\(\lambda(p)=|p|^2+O(|p|^4)\), \(\lambda_{1/2}=\frac14|p|^2+O(|p|^4)\), and
\(\lambda(p)\ge\frac4{\pi^2}|p|^2\).  Three facts control the numerator.
(i) Exactly, \(\lambda(p)-4\lambda(p/2)=-16\sum_\mu\sin^4(p_\mu/4)=O(|p|^4)\);
the executable verifies this as a Laurent-polynomial identity in \(e^{ip/4}\).
(ii) \(\|E\|=O(|p|)\), so \(\lambda(p)\,EE^*=O(|p|^4)\), and the part of
\(E+E^*\) of order \(|p|^2\) contributes \(O(|p|^4)\) after multiplication by
\(\lambda(p)\).  (iii) The linear part of \(E+E^*\): by
Lemma~\ref{lem:v35-closed-form} with \(u_\nu=e^{ip_\nu/2}-1=\frac i2p_\nu+O(|p|^2)\),
the diagonal of \(E+E^*\) is \(\sum_\nu(u_\nu+\bar u_\nu)=-4\sum_\nu\sin^2(p_\nu/4)=O(|p|^2)\)
to linear order zero, and the off-diagonal entry \((\mu,\nu)\) is
\(-u_\mu-\bar u_\nu=\frac i2(p_\nu-p_\mu)+O(|p|^2)\).  Hence
\(E+E^*=\frac i2(\mathbb 1p^{\mathsf T}-p\mathbb 1^{\mathsf T})+O(|p|^2)\).
Now \(q(p)=ip+O(|p|^2)\), \(q^*p=-i|p|^2+O(|p|^3)\), so
\(\Pi(p)p=p-q(q^*p)/\lambda(p)=O(|p|^2)\) and likewise \(p^{\mathsf T}\Pi(p)=O(|p|^2)\).
Therefore \(\Pi(E+E^*)\Pi=O(|p|^2)\) and \(\lambda(p)\,\Pi(E+E^*)\Pi=O(|p|^4)\).
Altogether \(\Pi N\Pi=O(|p|^4)\) while the denominator
\(16\lambda_{1/2}\lambda(p)\ge\frac{16}{\pi^4}|p|^4\), so \(T\) is bounded near
\(0\).  Since \(D_1=T+\frac1{16}\sum_{\epsilon\ne0}(\cdots)\), (39.11) follows.
\end{proof}

\begin{remark}[The value of \(d_1\)]
\label{rem:v35-d1}
The lemma proves finiteness; the number is a floating diagnostic.  Sampling
\(3000\) random momenta and \(1000\) small momenta down to \(|p|=0.003\), the
executable finds \(\|D_1\|\le 0.1954\), approached as \(p\to0\) (the values
0.1932,\ 0.1954,\ 0.1954,\ 0.1951,\ 0.1946 at \(|p|=0.3,0.1,0.03,0.01,0.003\)).  The ratio
\(\|\widehat S^{(1)}-\widehat S^{(0)}\|/\lambda(p)^2\) stays at most
0.1965 over the same scan, which is the kernel-level form of the same
statement.  With \(d_1\approx0.20\) the region \(\lambda(k)<(1-\rho)/d_1\) of
Corollary~\ref{cor:v35-kernels} covers \(\lambda(k)\lesssim2.7\), i.e.\ well
beyond the long-wavelength regime.
\end{remark}

\subsection{The off-axis Gaussian fixed point}

\begin{theorem}[Off-axis Gaussian fixed point of the symmetrized rule]
\label{thm:v35-fixed-point}
Let \(G^{(j)}=\mathcal L^j[G^{(0)}]\) be the transverse propagators of the
symmetrized blocking tower and \(\rho=\frac{17}{36}\), \(d_1\) as in (39.11).
Then for all \(j,m\ge0\) and all \(k\ne0\),
\[
 \|G^{(j+m)}(k)-G^{(j)}(k)\|\le\frac{d_1}{1-\rho}\,\rho^j,             \tag{39.12}
\]
so \(G^{(j)}\) converges uniformly on \(\mathbb T^4\setminus\{0\}\) to a
transverse Hermitian limit \(G^*\) with
\(\|G^{(j)}-G^*\|\le d_1\rho^j/(1-\rho)\).  The limit satisfies
\(G^*=\mathcal L[G^*]\), the bounds
\[
 \Bigl(\frac1{\lambda(k)}-\frac{d_1}{1-\rho}\Bigr)\Pi(k)\le G^*(k)\le\Bigl(\frac1{\lambda(k)}+\frac{d_1}{1-\rho}\Bigr)\Pi(k),
                                                                 \tag{39.13}
\]
and it is the unique fixed point of \(\mathcal L\) in the class
\(G^{(0)}+\{\text{bounded transverse Hermitian functions on }\mathbb T^4\setminus\{0\}\}\).
\end{theorem}

\begin{proof}
Since \(\mathcal L\) is linear, \(G^{(i+1)}-G^{(i)}=\mathcal L^i[D_1]\).  By
Lemma~\ref{lem:v35-first-step}, \(-d_1\Pi\le D_1\le d_1\Pi\) pointwise, and
by Theorem~\ref{thm:v35-contraction},
\(-d_1\rho^i\Pi\le\mathcal L^i[D_1]\le d_1\rho^i\Pi\).  Summing the geometric
series from \(i=j\) to \(j+m-1\) gives (39.12); summing from \(0\) gives
(39.13) after adding \(G^{(0)}=\Pi/\lambda\).  Uniform convergence gives a
limit \(G^*\); it is transverse and Hermitian as a limit of such matrices.
Fixed-point property: \(\mathcal L\) applied pointwise to the sixteen aliases
of \(k\) is continuous under uniform convergence, so
\(\mathcal L[G^*]=\lim\mathcal L[G^{(j)}]=\lim G^{(j+1)}=G^*\).  Uniqueness:
if \(G'=\mathcal L[G']\) with \(G'-G^{(0)}\) bounded, then \(G'-G^*\) is
bounded, say by \(c\), and \(G'-G^*=\mathcal L^j[G'-G^*]\) has norm at most
\(c\rho^j\) for every \(j\), hence vanishes.
\end{proof}

\begin{corollary}[Convergence of the kernels]
\label{cor:v35-kernels}
Let \(\widehat S^{(j)}=(G^{(j)})^+\) and \(\widehat S^*=(G^*)^+\) be the
transverse pseudo-inverses.  On the set
\(K:=\{k\ne0:\ \lambda(k)<(1-\rho)/d_1\}\) the limit is nondegenerate,
\(G^*\ge\delta(k)\Pi\) with \(\delta(k)=1/\lambda(k)-d_1/(1-\rho)>0\), and
\[
 \|\widehat S^{(j)}(k)-\widehat S^*(k)\|\le\frac{d_1\rho^j}{(1-\rho)\,\delta(k)^2}
 \qquad(k\in K).                                                   \tag{39.14}
\]
The same holds on any set where \(\inf_j\lambda_{\min}^\perp(G^{(j)})>0\).
\end{corollary}

\begin{proof}
On \(K\), (39.13) and the same bound for every \(G^{(j)}\) give
\(G^{(j)},G^*\ge\delta\Pi\).  For Hermitian positive semidefinite \(A,B\)
with the common kernel \(\mathbb Cq\), \(A^+-B^+=A^+(B-A)B^+\) on \(q^\perp\),
and \(\|A^+\|\le1/\lambda^\perp_{\min}(A)\).
\end{proof}

\begin{remark}[Observed rates]
\label{rem:v35-rates}
At three pseudo-random momenta the propagator increments
\(\|G^{(j+1)}-G^{(j)}\|\), \(j=0,\ldots,3\), are at most
\(0.076,\ 0.027,\ 0.0094,\ 0.0031\), each below the proven bound
\(d_1\rho^j=0.195,\ 0.092,\ 0.044,\ 0.021\); their successive ratios are
\(0.52,0.47,0.37\), \(0.39,0.39,0.34\), and \(0.35,0.35,0.33\), which move from
the uniform bound \(\frac{17}{36}\) towards the axis rate \(\frac14\) as
(39.10) predicts.  The kernel increments start large (the kernels are of
size ten at these momenta) with ratios \(1.06,0.66,0.44\), \(0.96,0.61,0.41\),
\(0.92,0.56,0.39\), and the smallest transverse eigenvalue of \(G^{(4)}\) at
these points is at least \(0.030\).  The sampled nondegeneracy suggests that
Corollary~\ref{cor:v35-kernels} holds on the whole torus, which is the first
V36 target.
\end{remark}

\begin{theorem}[First hypothesis of the fixed-point card]
\label{thm:v35-card-hypothesis}
For the symmetrized rule the kernel-convergence hypothesis of
Theorem~\ref{thm:v34-card} holds on \(K\) with \(\rho=\frac{17}{36}\), and
the propagator-convergence statement holds on all of \(\mathbb T^4\setminus\{0\}\).
\end{theorem}

\begin{proof}
Theorem~\ref{thm:v35-fixed-point} and Corollary~\ref{cor:v35-kernels}.
\end{proof}

\subsection{Executable provenance}

\begin{record}[V35 verifier]
\label{rec:v35-verifier}
The executable
\begin{center}
\path{scripts/verify_ym_postarchive_v35_alias_gram_contraction.py}
\end{center}
has 514 physical lines, 19429 bytes, and SHA--256
\[
\texttt{02869F8722822F778E0D784865E01D8334B7D400959BD19B6B7B38FC0DCBCE51}.
                                                               \tag{39.15}
\]
Its canonical payload SHA--256 is
\[
\texttt{7FE24D9B1954559DE8AFEDCD9CA14851C089EEE7C9D11716E62B0D7435610A5A}.
                                                               \tag{39.16}
\]
The hashed payload records: the closed form (39.5) against the sixteen-path
enumeration, the Ward identity (39.6), and the Gram form (39.7), all as
exact identities of Laurent polynomials in \(w\) over \(\mathbb Q\); the
corner identity (39.9) against the V34 propagator step at the fifteen
nonzero corners; exact momentum-inversion and permutation covariance of the
symmetrized level-one kernel at all \(256\) quarter-turn momenta, together
with the exact count of failures of axis-reflection covariance
(\(924\) of the \(1024\) point--axis pairs fail, the other \(100\) being trivial reflections: a base-point blocking has no reflection
symmetry, which is why Lemma~\ref{lem:v35-first-step} uses the expansion
rather than a symmetry argument); full rank \(45\) of the reduced
fluctuation block at the toron; the identity of Lemma~\ref{lem:v35-first-step}(i)
and the linear part of Lemma~\ref{lem:v35-closed-form}; and (39.4).  The
floating diagnostics of Remarks~\ref{rem:v35-d1} and~\ref{rem:v35-rates}
are stored under a separate non-hashed key.
\end{record}

\subsection{Decision for V36}

\begin{decision}[Nondegeneracy, blocked vertices, and the oracle frame]
\label{dec:v35-v36}
V36 should first extend Corollary~\ref{cor:v35-kernels} to the whole torus
by a lower bound on \(G^{(j)}\) from the straight-chain minorant
\(G^{(j+1)}(k)\ge\frac1{16}\Pi B_0G^{(j)}(k/2)B_0^*\Pi\), or by an exact
positivity argument for \(\Pi B_0B_0^*\Pi\) on \(q^\perp\).  It should then
begin the blocked vertices: the one-colour commuting-toron cubic vertex of
the tenth archive transported through one symmetrized blocking step, whose
Lipschitz dependence on the kernel is the second hypothesis of
Theorem~\ref{thm:v34-card}.  In parallel it should settle the single open
normalization question of Firewall~\ref{fire:v35-normalization} by
evaluating the V60 oracle's frame convention on the exactly solvable axis of
Theorem~\ref{thm:v33-closed-form}.
\end{decision}

\begin{assessment}[V35 acquisition]
\label{ass:v35}
V35 completes the compulsory audit, proves the alias Ward identity and the
four-term alias Gram matrix in closed form, derives from them the exact
contraction \(\rho\le\frac{17}{36}\) of the propagator map, proves that the
first blocking step is a bounded correction, and concludes that the
symmetrized Gaussian tower converges geometrically to a unique fixed point
off the axes.  It fixes the colour factor \(3\) between the compact and the
archived kernels exactly.  The remaining Gaussian task is nondegeneracy on
the whole torus; the remaining one-loop tasks are the blocked vertices and
the oracle frame convention.
\end{assessment}

\begin{record}[V35 append record]
\label{rec:v35-append}
V35 is appended to the reconstructed V34 whose installed record was
\[
\begin{array}{c|c}
\text{lines}&12282\\
\text{bytes}&489911\\
\text{SHA--256}&
\texttt{3CF829F57455C43D3CE54387782C3C60EB3497F1B680A11DDCAED13306D8C469}.
\end{array}                                                    \tag{39.17}
\]
The V34 file was rebuilt from the session transcript after an accidental
deletion of the installed V33; its V1--V33 text was recovered line by line
from the recorded reads, and the four unrecorded lines were restored from
the same session's re-emission.  After validation only V35 is to remain the
active numeric YM note.  The next compulsory companion audit and the
ten-version sweep are both V40.
\end{record}

\begin{firewall}[Claim boundary after V35]
\label{fire:v35-final}
V35 proves a tree-level fixed point of a linear Gaussian recursion for one
blocking rule.  It proves no one-loop shell coefficient, no Lipschitz bound
for blocked vertices, no identification \(b^*=b_{\rm YM}\), no Haar sign,
and no continuum statement.  The Gaussian anomaly, the interacting measure,
regulator removal, Osterwalder--Schrader reconstruction, physical
clustering, and the Hamiltonian Yang--Mills mass gap remain open.
\end{firewall}

% =============================================================================
% END APPENDED POST-TENTH-ARCHIVE V35 MATERIAL
% =============================================================================

% =============================================================================
% BEGIN APPENDED POST-TENTH-ARCHIVE V36 MATERIAL
% =============================================================================

\section{V36: levelwise nondegeneracy, lift chains, the vertical covariance,
and the single remaining Gaussian inequality}
\label{sec:v36}

V35 proved that the symmetrized Gaussian tower has a unique fixed-point
propagator \(G^*\) and that the kernels converge on the region
\(\lambda(k)<(1-\rho)/d_1\).  The one-loop shell functional needs the coarse
kernel \(\widehat S^{(j+1)}=(G^{(j+1)})^+\) explicitly, through the
fine-given-coarse (vertical) covariance, so uniform control along the tower
requires \(G^*\) to be nondegenerate on the whole torus.  This note settles
what can be settled exactly: every finite level is nondegenerate, the full
alias Gram matrix has the uniform floor \(\frac{140}{53}\), any degeneracy
of the limit would have to be carried by the chains that descend towards
zero momentum along the sixteen lifts of \(k\), and the vertical covariance
has an exact closed form of rank \(45\).  It then reduces global
nondegeneracy to one explicit finite-dimensional inequality, the three-step
lift Gram floor (H\(_U\)) of Theorem~\ref{thm:v36-conditional}, whose
sampled value is about \(6.9\), and explains why the straight-chain
minorant cannot replace it: the straight alias block is exactly singular at
some momenta.

\subsection{The full alias Gram floor and levelwise nondegeneracy}

\begin{theorem}[Uniform floor of the alias Gram matrix]
\label{thm:v36-gram-floor}
For every \(k\), \(S(k)=\sum_\epsilon B_\epsilon(k)B_\epsilon(k)^*\ge\frac{140}{53}\,I\)
as a full \(4\times4\) matrix.  Consequently, for every momentum \(p\) and
every vector \(u\), \(\sum_\epsilon\|B_\epsilon(p)^*u\|^2\ge\frac{140}{53}|u|^2\), and for
every chain length \(n\) and unit \(v\perp q(k)\),
\[
 \sum_{M\bmod2^n}\|W_M(k)^*v\|^2\ \ge\ 4\Bigl(\frac{140}{53}\Bigr)^{n-1},
                                                                 \tag{40.1}
\]
where \(W_M\) is the ordered product of the alias blocks along the chain
\(M\).
\end{theorem}

\begin{proof}
By (39.7), for a unit vector \(u\), with \(a=u^*q(k)\) and \(b=\mathbb 1^*u\),
\(u^*Su=4+\frac89\sum_\mu|q_\mu|^2|u_\mu|^2-2\operatorname{Re}(a\bar b)+\frac{53}{18}|a|^2
 \ge4-4|a|+\frac{53}{18}|a|^2\), since \(|b|\le2\).  The quadratic in \(|a|\)
is minimized at \(|a|=\frac{36}{53}\) with value \(4-\frac{144}{53}+\frac{72}{53}=\frac{140}{53}\).
For (40.1) apply the transverse bound \(v^*Sv\ge4\) of (39.9) at the first
step and the full bound at every later step.  The executable verifies
\(S-\frac{140}{53}I\succeq0\) exactly at the fifteen nonzero corners and at
ten Pythagorean momenta by exact elimination.
\end{proof}

\begin{theorem}[Every level is nondegenerate]
\label{thm:v36-levelwise}
For every \(j\ge0\) and every \(k\ne0\), \(G^{(j)}(k)\) has rank three, i.e.\ it
is positive definite on \(q(k)^\perp\).  Equivalently, at every level the
constraint map from transverse fine amplitudes to transverse coarse links
is surjective, so the minimizer map and the vertical covariance below are
defined at every finite level.
\end{theorem}

\begin{proof}
Induction on \(j\); \(G^{(0)}=\Pi/\lambda\) is nondegenerate for \(p\ne0\).
For a unit \(v\perp q(k)\), (39.8) and Lemma~\ref{lem:v35-inner} give
\(v^*G^{(j+1)}(k)v=\frac1{16}\sum_\epsilon(B_\epsilon^*v)^*G^{(j)}(p_\epsilon)(B_\epsilon^*v)\),
and \(B_\epsilon^*v\perp q(p_\epsilon)\) by Lemma~\ref{lem:v35-ward}.  If the
sum vanished, every \(B_\epsilon^*v\) would lie in the kernel of
\(G^{(j)}(p_\epsilon)\) on \(q(p_\epsilon)^\perp\), which is trivial by
induction because \(p_\epsilon\ne0\) when \(k\ne0\).  Hence
\(\sum_\epsilon\|B_\epsilon^*v\|^2=0\), contradicting \(v^*Sv\ge4\).
\end{proof}

\subsection{Lift chains and the degeneracy set of the limit}

For a lift \(k'=k+2\pi M\), \(M\in\mathbb Z^4\), put \(\theta_l=k'/2^l\).  The
chain of \(k'\) is the alias chain whose momentum after \(l\) steps is
\(\theta_l\) (modulo \(2\pi\)); its blocks are \(B(e^{i\theta_l})\), \(l=1,\ldots,n\),
in this order, and \(W_{k'}^{(n)}\) denotes their product.  Chains of
length \(n\) are labelled by \(M\bmod2^n\).

\begin{theorem}[Lift-chain minorant]
\label{thm:v36-minorant}
For any set \(\mathcal M\) of lifts that are distinct modulo \(2^n\), any
\(j\), and \(k\ne0\),
\[
 G^{(j+n)}(k)\ \ge\ 16^{-n}\,\Pi(k)\sum_{k'\in\mathcal M}W^{(n)}_{k'}\,G^{(j)}(k'/2^n)\,W^{(n)*}_{k'}\,\Pi(k),
                                                                 \tag{40.2}
\]
and the same with \(G^{(j+n)},G^{(j)}\) replaced by \(G^*,G^*\).  The iterated
Ward identity \(W^{(n)}_{k'}q(k'/2^n)=q(k)\) holds, so the inner projectors
are invisible: \(\Pi(k)W_{k'}\Pi(k'/2^n)W_{k'}^*\Pi(k)=\Pi(k)W_{k'}W_{k'}^*\Pi(k)\).
In particular, if \(G^*(k)v=0\) for a unit \(v\perp q(k)\), then
\(W^{(n)*}_{k'}v=0\) for every lift \(k'\) with \(k'/2^n\) in the
nondegeneracy region \(K\) of Corollary~\ref{cor:v35-kernels}.
\end{theorem}

\begin{proof}
Iterating (39.8) \(n\) times writes \(G^{(j+n)}\) as \(16^{-n}\Pi\sum_M W_MG^{(j)}(p_M)W_M^*\Pi\)
over all \(16^n\) chains, all terms being positive semidefinite; dropping
chains gives (40.2).  The iterated Ward identity follows from
Lemma~\ref{lem:v35-ward} applied at each step, and it removes the inner
projectors exactly as in Lemma~\ref{lem:v35-inner}.  For the fixed point use
\(G^*=\mathcal L^n[G^*]\).  If \(v^*G^*(k)v=0\), every retained term vanishes;
at a chain end in \(K\), \(G^*\ge\delta\Pi\) with \(\delta>0\), so
\(\Pi(k'/2^n)W^*v=W^*v=0\).  The executable verifies the two-step iterated
Ward identity exactly at thirty Pythagorean momenta.
\end{proof}

\subsection{Why the straight chain is not enough}

\begin{proposition}[Exact singularities of the straight block]
\label{prop:v36-singular}
The alias block \(B(w)\) of (39.5) is singular at \(w=(1,-i,1,i)\) (rank
three), and vanishes identically at every sign vector \(w\in\{\pm1\}^4\)
other than \(w=1\).  At \(k=(-\pi,0,\pi,0)\) the straight-alias Gram
\(\Pi(k)B_0(k)B_0(k)^*\Pi(k)\) has rank two on the three-dimensional
transverse space.  Consequently the straight-chain minorant (the single lift
\(M=0\)) has zero transverse floor at some momenta, and no lower bound of
\(G^*\) can be obtained from it alone.
\end{proposition}

\begin{proof}
Exact ranks are computed by the executable in Gaussian rationals.  At a
sign vector with some \(w_\nu=-1\) the straight entries contain the factor
\(1+w_\nu=0\) and every connector contains a factor \(1-w_\mu^2=0\).
\end{proof}

\begin{record}[Where the straight block is and is not small, non-hashed]
\label{rec:v36-sigma-boxes}
Sampling \(\sigma_{\min}(B(e^{i\theta}))\) on the boxes
\(|\theta|_\infty\le\pi/4,\ 3\pi/8,\ \pi/2\) gives minima about
\(0.55,\ 0.057,\ 0.0003\); the torus maximum of \(\|B(w)\|\) is about \(3.1\).
Thus the singular set reaches into \(|\theta|_\infty\approx1.15\), and the
box \(|\theta|_\infty\le3\pi/8\), which the third chain step visits, is not
safe.  The fourth step and beyond stay in \(|\theta|_\infty\le3\pi/16\).
\end{record}

\begin{lemma}[Small-angle bound]
\label{lem:v36-small-angle}
With \(E(\theta)=B(e^{i\theta})-2I\),
\(\|E\|\le(4+\sqrt3)|\theta|_2+\tfrac12|\theta|_2^2\), hence
\(\sigma_{\min}(B(e^{i\theta}))\ge2-5.74|\theta|_2-\tfrac12|\theta|_2^2\).
\end{lemma}

\begin{proof}
\(D_{\mu\mu}-2=2(\prod_\nu e^{i\theta_\nu/2}\cos(\theta_\nu/2)-1)\), and
\(|\prod e^{i\theta_\nu/2}\cos(\theta_\nu/2)-1|\le|e^{i\sum\theta_\nu/2}-1|+1-\prod\cos(\theta_\nu/2)\le|\theta|_1/2+|\theta|_2^2/8\).
A connector \((1-e^{2i\theta_\mu})c_{\mu\nu}\) has modulus at most
\(2|\sin\theta_\mu|\cdot\frac12\le|\theta_\mu|\), since \(|c_{\mu\nu}|\le\frac{12}{24}\).
Hence \(\|E\|_F^2\le4(|\theta|_1+|\theta|_2^2/4)^2+3|\theta|_2^2\), and
\(|\theta|_1\le2|\theta|_2\).
\end{proof}

\begin{lemma}[Certified floor on the fourth-step box]
\label{lem:v36-certified}
\(\|\partial_{\theta_\nu}B(e^{i\theta})\|\le2.72\) for every \(\nu\) and
\(\theta\), and
\[
 \sigma_{\min}\bigl(B(e^{i\theta})\bigr)\ \ge\ 0.739\qquad\text{for all }|\theta|_\infty\le3\pi/16.
                                                                 \tag{40.3}
\]
\end{lemma}

\begin{proof}
\(\partial_{\theta_\nu}D_{\mu\mu}=\frac{iw_\nu}{1+w_\nu}D_{\mu\mu}\) has modulus
\(\prod_{\rho\ne\nu}|\cos(\theta_\rho/2)|\le1\).  For the connector in position
\((\mu,\nu')\): if \(\nu=\mu\) the derivative is \(-2iw_\mu^2c_{\mu\nu'}\), of
modulus at most \(1\); if \(\nu\) lies in the complement of \(\{\mu,\nu'\}\) it is
\((1-w_\mu^2)\,iw_\nu(w_{\nu''}+2)/24\), of modulus at most \(\frac14\); otherwise
it is zero.  Summing squares, \(\|\partial_\nu B\|_F^2\le4+3+6/16\), so
\(\|\partial_\nu B\|\le2.72\).  For \(\theta,\theta'\) with
\(|\theta-\theta'|_\infty\le h/2\), integrating along coordinate segments gives
\(\|B(e^{i\theta})-B(e^{i\theta'})\|\le5.44h\), and \(\sigma_{\min}\) is
\(1\)-Lipschitz in the operator norm.  The executable evaluates
\(\sigma_{\min}\) at the \(60^4\) points of the grid of spacing
\(h=0.019968\) on the box, finds the minimum \(0.8480\), and records the
certified floor \(0.8480-5.44h-10^{-9}=0.7394\); the singular value
decomposition of a \(4\times4\) matrix in double precision is backward
stable to well within the subtracted \(10^{-9}\).  This is a certified
floating computation with a proved Lipschitz constant, and it is included
in the canonical payload.
\end{proof}

\subsection{The conditional global theorem}

For \(k\in(-\pi,\pi]^4\) and \(\epsilon\in\{0,1\}^4\) let
\(k'_\epsilon=k+2\pi\epsilon\) be the sixteen lifts with components in
\((-\pi,3\pi]\), so that the chain angles satisfy
\(|\theta_l|_\infty=|k'_\epsilon/2^l|_\infty\le3\pi/2^l\).  The first block
of the chain of \(k'_\epsilon\) is the alias block \(B_\epsilon(k)\), so the
sixteen lift chains are the sixteen aliases continued straight.

\begin{definition}[The three-step lift Gram and hypothesis (H\(_U\))]
\label{def:v36-HU}
\(U(k):=\sum_\epsilon W^{(3)}_{k'_\epsilon}W^{(3)*}_{k'_\epsilon}
 =\sum_\epsilon B_\epsilon(k)B(e^{ik'_\epsilon/4})B(e^{ik'_\epsilon/8})\bigl[\cdots\bigr]^*\).
Hypothesis (H\(_U\)): there is \(u_*>0\) with \(v^*U(k)v\ge u_*\) for all
\(k\ne0\) and all unit \(v\perp q(k)\).
\end{definition}

\begin{theorem}[Global nondegeneracy under (H\(_U\))]
\label{thm:v36-conditional}
Assume (H\(_U\)).  Then there is \(c>0\), explicit in \(u_*\), \(d_1\) and
the constants of Lemmas~\ref{lem:v36-small-angle}--\ref{lem:v36-certified},
such that \(G^*(k)\ge c\,\Pi(k)\) for every \(k\ne0\).  Consequently
\(\sup_j\sup_k\|\widehat S^{(j)}(k)\|\le1/c\), the kernels converge uniformly
on the whole torus with
\(\|\widehat S^{(j)}-\widehat S^*\|\le d_1\rho^j/((1-\rho)c^2)\), and the first
hypothesis of Theorem~\ref{thm:v34-card} holds on all of \(\mathbb T^4\setminus\{0\}\).
\end{theorem}

\begin{proof}
Choose \(n\ge6\) with \(36\pi^2\,4^{-n}d_1/(1-\rho)\le\frac12\).  For each lift
\(k'_\epsilon\), \(\lambda(k'_\epsilon/2^n)\le|k'_\epsilon/2^n|_2^2\le36\pi^24^{-n}\),
so by (39.13) \(G^*(k'_\epsilon/2^n)\ge\frac{4^n}{72\pi^2}\Pi\).  Apply (40.2)
with the sixteen lifts: for unit \(v\perp q(k)\),
\(v^*G^*(k)v\ge16^{-n}\frac{4^n}{72\pi^2}\sum_\epsilon\|W^{(n)*}_{k'_\epsilon}v\|^2\).
Split \(W^{(n)}=W^{(3)}R\) with \(R\) the product of the blocks at steps
\(4,\ldots,n\); \(\|R^*x\|\ge\prod_{l=4}^n\sigma_l\,|x|\) with
\(\sigma_l\ge0.739\) for \(l=4,5\) by (40.3) (since \(|\theta_l|_\infty\le3\pi/16\))
and \(\sigma_l\ge2-5.74\cdot6\pi/2^l-\frac12(6\pi/2^l)^2\) for \(l\ge6\) by
Lemma~\ref{lem:v36-small-angle} (since \(|\theta_l|_2\le6\pi/2^l\)); the
latter is positive from \(l=6\) on (\(0.26,1.14,1.57,\ldots\)) and
\(\prod_{l\ge6}(\sigma_l/2)^2\) converges to a positive number.  Hence
\(v^*G^*(k)v\ge\frac{u_*}{72\pi^2}\,4^{-3}\prod_{l=4}^n(\sigma_l/2)^2=:c>0\),
uniformly in \(k\) and in \(n\).  The kernel statements follow as in
Corollary~\ref{cor:v35-kernels} with \(\delta=c\).
\end{proof}

\begin{record}[The status of (H\(_U\)), non-hashed]
\label{rec:v36-HU-status}
Sampling \(4000\) momenta, the transverse floor of \(U(k)\) is at least
\(6.90\), and that of the two-step Gram \(T(k)=\sum_\epsilon W^{(2)}W^{(2)*}\)
at least \(2.18\); \(T-I\succ0\) holds exactly at thirty Pythagorean momenta
in the fundamental domain.  The two-step Gram alone does not suffice,
because the third chain step visits the unsafe box of
Record~\ref{rec:v36-sigma-boxes}.  A direct grid certificate of (H\(_U\))
is out of reach: the proved Lipschitz constants of a three-fold product
of blocks are of order \(10^3\), which would require of order \(10^{13}\)
grid points; (H\(_U\)) must be proved by algebra.  It is a statement about
an explicit polynomial matrix in \(e^{ik/8}\) with coefficients in
\(\mathbb Q(i,\sqrt2)\) on the fundamental domain.
\end{record}

\subsection{The vertical covariance}

\begin{theorem}[Minimizer map and vertical covariance]
\label{thm:v36-vertical}
At level \(j+1\) and coarse momentum \(k\ne0\), the constrained minimizer of
the level-\(j\) energy with coarse link \(c\) has alias amplitudes
\(a^\epsilon=M_\epsilon c\), \(M_\epsilon=\frac1{16}G^{(j)}(p_\epsilon)B_\epsilon(k)^*\widehat S^{(j+1)}(k)\),
and the covariance of the transverse fine field conditioned on the coarse
link (the vertical covariance) is
\[
 \operatorname{Cov}^{\rm vert}_{\epsilon\epsilon'}(k)
 =\frac1{16}G^{(j)}(p_\epsilon)\delta_{\epsilon\epsilon'}
 -\frac1{256}G^{(j)}(p_\epsilon)B_\epsilon(k)^*\widehat S^{(j+1)}(k)B_{\epsilon'}(k)G^{(j)}(p_{\epsilon'}).
                                                                 \tag{40.4}
\]
It is Hermitian positive semidefinite, annihilated by the constraint
\(\Pi(k)\sum_\epsilon B_\epsilon(\cdot)\), and of rank \(45=16\cdot3-3\).
\end{theorem}

\begin{proof}
In normalized amplitudes the level-\(j\) Gaussian has covariance
\(C=\frac1{16}\bigoplus_\epsilon G^{(j)}(p_\epsilon)\) on the transverse
space, and the coarse link is \(Lc=\Pi(k)\sum_\epsilon B_\epsilon a^\epsilon\).
Conditioning a Gaussian on a linear image gives the mean \(CL^*(LCL^*)^+Lc\)
and the covariance \(C-CL^*(LCL^*)^+LC\).  Here \(LCL^*=\mathcal L[G^{(j)}]=G^{(j+1)}\)
by (39.8), \((G^{(j+1)})^+=\widehat S^{(j+1)}\), and \(CL^*\) has blocks
\(\frac1{16}G^{(j)}(p_\epsilon)B_\epsilon^*\Pi\); since \(\widehat S\Pi=\widehat S\)
the projector is absorbed, giving \(M_\epsilon\) and (40.4).  The rank is the
dimension of the transverse fine space minus the rank of \(L\), which is
three by Theorem~\ref{thm:v36-levelwise}.  The executable builds (40.4)
exactly at the fifteen nonzero corners from the exact level-zero propagators
and the exact level-one kernel, and verifies Hermiticity, positive
semidefiniteness by exact elimination, annihilation by the constraint, and
rank \(45\).
\end{proof}

\begin{remark}[The gate]
\label{rem:v36-gate}
The number \(45\) is the dimension of the reduced fluctuation block of V32:
the vertical covariance is the inverse of that block transported to alias
coordinates.  Its background dependence is what the one-loop shell
coefficient of the fixed-point card differentiates.  Formula (40.4) contains
the coarse kernel \(\widehat S^{(j+1)}\), which is why the uniform kernel
bound of Theorem~\ref{thm:v36-conditional}, hence (H\(_U\)), is the gate to
the one-loop programme; at every finite level the formula is unconditional
by Theorem~\ref{thm:v36-levelwise}.
\end{remark}

\subsection{Executable provenance}

\begin{record}[V36 verifier]
\label{rec:v36-verifier}
The executable
\begin{center}
\path{scripts/verify_ym_postarchive_v36_nondegeneracy_and_vertical_covariance.py}
\end{center}
has \(376\) physical lines, \(15217\) bytes, and SHA--256
\[
\texttt{FA7E8EF685EC29F7E4B6DA2E94EED56ABFAF1E90288A4EF88713631A4C0B03DA}.
                                                               \tag{40.5}
\]
Its canonical payload SHA--256 is
\[
\texttt{CF9CCC7BB30F854696B8E6C75B43BDD12AD49BDE350FE0EEBE5D353801BF3807}.
                                                               \tag{40.6}
\]
The hashed payload records the exact Gram-floor checks at \(25\) points, the
exact singular ranks, the exact vertical covariance properties at the fifteen
nonzero corners, the exact two-step lift Gram positivity and iterated Ward
identity at \(30\) Pythagorean momenta, and the certified grid of
Lemma~\ref{lem:v36-certified}.  The sampled floors of
Records~\ref{rec:v36-sigma-boxes} and~\ref{rec:v36-HU-status} are stored
under the separate non-hashed key.
\end{record}

\subsection{Decision for V37}

\begin{decision}[Prove (H\(_U\)) by algebra]
\label{dec:v36-v37}
V37 should prove the three-step lift Gram floor.  Three routes are
admissible: (i) exact expansion of \(U\) in the Laurent-polynomial class over
\(\mathbb Q(i,\sqrt2)\), followed by a decomposition
\(U-u_*I=\sum_\alpha F_\alpha F_\alpha^*\) or a Finsler certificate on
\(q(k)^\perp\); (ii) a lower bound of the two-step straight Gram
\(B(e^{ip/2})B(e^{ip/4})[\cdots]^*\) on the transverse space of \(p\) that
uses the Ward identity to exclude the singular directions of
Proposition~\ref{prop:v36-singular}, combined with the first-step bound
\(v^*Sv\ge4\); (iii) an upper bound of the kernels by explicit trial fields
built from the minimizer map of Theorem~\ref{thm:v36-vertical} at level one.
Failing all three, V37 should record the obstruction and move to the
one-loop shell functional at finite level, where
Theorem~\ref{thm:v36-levelwise} makes every object unconditional.
\end{decision}

\begin{assessment}[V36 acquisition]
\label{ass:v36}
V36 proves the uniform Gram floor \(\frac{140}{53}\), nondegeneracy of every
finite level, the lift-chain minorant with its exact consequence for the
degeneracy set of the limit, a certified singular-value floor for the deep
chain steps, and the exact rank-\(45\) vertical covariance.  It shows that
global nondegeneracy of the fixed point, hence uniform kernel convergence
and the first hypothesis of the fixed-point card on the whole torus, follows
from the single inequality (H\(_U\)), which is numerically robust and not
yet proved.
\end{assessment}

\begin{record}[V36 append record]
\label{rec:v36-append}
V36 is appended to the clean V35 whose installed record was
\[
\begin{array}{c|c}
\text{lines}&12860\\
\text{bytes}&520002\\
\text{SHA--256}&
\texttt{6C014195DC1E1A52F7943D3840A669EB813C202A77B90B80F31E251AB1C1BB9D}.
\end{array}                                                    \tag{40.7}
\]
After validation only V36 is to remain the active numeric YM note.  The
next compulsory companion audit and the ten-version sweep are both V40.
\end{record}

\begin{firewall}[Claim boundary after V36]
\label{fire:v36-final}
Global nondegeneracy of the fixed-point propagator is proved only under
(H\(_U\)); without it, nondegeneracy holds at every finite level and on the
region of Corollary~\ref{cor:v35-kernels}.  Lemma~\ref{lem:v36-certified}
rests on a floating-point grid with a proved Lipschitz constant.  No
one-loop shell coefficient, no Lipschitz bound for blocked vertices, no
identification \(b^*=b_{\rm YM}\), no Haar sign, and no continuum statement
is made; the Gaussian anomaly, the interacting measure, regulator removal,
Osterwalder--Schrader reconstruction, physical clustering, and the
Hamiltonian Yang--Mills mass gap remain open.
\end{firewall}

% =============================================================================
% END APPENDED POST-TENTH-ARCHIVE V36 MATERIAL
% =============================================================================

% =============================================================================
% BEGIN APPENDED POST-TENTH-ARCHIVE V37 MATERIAL
% =============================================================================

\section{V37: upstream to the block-average rule, whose Gaussian tower is
explicitly solvable}
\label{sec:v37}

V36 reduced global nondegeneracy of the symmetrized fixed point to the
three-step lift inequality (H\(_U\)) and showed that it cannot be certified
by grids.  Following the programme's rule to go upstream when blocked, this
note changes the linear blocking rule of the Gaussian tower.  The base-point
path rules of V32--V34 were modelled on the nonabelian parity-tree quotient;
the Gaussian half of the fixed-point card only needs a gauge-covariant
linear rule with hypercubic symmetry.  The classical block average of
lattice gauge theory, the average over all sixteen sites of a cell of the
straight two-link sums, has both properties, and its alias blocks are
\emph{diagonal}.  Every object of V32--V36 then becomes an explicit
positive lattice sum: the level-\(j\) propagators, the fixed point, the
contraction constant \(\rho=\frac14\) uniformly, the global nondegeneracy
floor \(256/\pi^{12}\), and the vertical covariance.  On the coordinate axes
the new tower coincides with the V33 tower at every level, so the axis
fixed point \(12\sin^2(t/2)/(2+\cos t)\) is unchanged.  The first hypothesis
of Theorem~\ref{thm:v34-card} thereby holds on the whole torus without any
unproved inequality.

\subsection{The rule}

\begin{definition}[Block-average rule]
\label{def:v37-rule}
For the cell \(\{0,1\}^4\) with base point \(0\) and a direction \(\mu\), the
coarse link is
\[
 c_\mu:=\frac1{16}\sum_{x\in\{0,1\}^4}\bigl(a_{(x,\mu)}+a_{(x+e_\mu,\mu)}\bigr),
                                                                 \tag{41.1}
\]
where a link leaving the cell is the corresponding link of the neighbouring
cell and carries its phase \(z_\mu=e^{ik_\mu}\).  Thus the cell edge
\((y,\mu)\) enters \(c_\mu\) with weight \(\frac1{16}(1+y_\mu+z_\mu(1-y_\mu))\).
In alias amplitudes, with \(w=e^{ip_\epsilon}\),
\[
 c_\mu=\sum_\epsilon s_\epsilon(1+w_\mu)\,a^\epsilon_\mu,\qquad
 s_\epsilon=\frac1{16}\prod_{\nu=1}^4(1+w_\nu),\qquad
 B_\epsilon(k)=s_\epsilon\,\Delta_\epsilon,\ \ \Delta_\epsilon=\operatorname{diag}(1+w_\mu).
                                                                 \tag{41.2}
\]
\end{definition}

\begin{lemma}[Ward property, transversality, normalization]
\label{lem:v37-ward}
\(B_\epsilon(k)q(p_\epsilon)=s_\epsilon\,q(k)\) for every alias.  The coarse
link of a fine pure gauge \(d\phi\) is the coarse gradient of the block
average of \(\phi\); the map (41.1) is invariant under fine gauge
transformations whose block averages vanish; the blocked kernels are
transverse and the propagator recursion (39.8) holds with the blocks (41.2)
and the slice normalization of Lemma~\ref{lem:v33-ward}.
\end{lemma}

\begin{proof}
\((1+w_\mu)(w_\mu-1)=w_\mu^2-1=q(k)_\mu\).  A pure gauge has alias
amplitudes \(a^\epsilon=q(p_\epsilon)\phi^\epsilon\), so
\(c=q(k)\sum_\epsilon s_\epsilon\phi^\epsilon\), and
\(\sum_\epsilon s_\epsilon\phi^\epsilon\) is the block average of \(\phi\)
in alias coordinates.  Pure gauges cost no energy and reach every multiple of
\(q(k)\), which gives transversality of the constrained minimum as in
Theorem~\ref{thm:v33-level-map}; the derivation of
Theorem~\ref{thm:v34-propagator} uses only the constraint form
\(c=\sum_\epsilon B_\epsilon a^\epsilon\) and applies verbatim.  The
executable checks that the incidence route with the edge weights of (41.1)
and the alias route with the blocks (41.2) give the same level-one kernels at
the sixteen corners exactly, and verifies the pseudo-inverse identities of
the propagator recursion at levels one and two at the fifteen nonzero
corners.
\end{proof}

\begin{lemma}[Hypercubic covariance]
\label{lem:v37-symmetry}
The level-one kernels satisfy, exactly at all \(256\) quarter-turn momenta,
permutation covariance (38.1), momentum inversion as complex conjugation,
and axis reflection covariance with the cell-shift phase,
\[
 \widehat S^{(1)}(R_\nu k)=D_\nu(k)\,R_\nu\widehat S^{(1)}(k)R_\nu\,D_\nu(k)^*,\qquad
 D_\nu(k)=\operatorname{diag}(1,\ldots,e^{-ik_\nu},\ldots,1),                       \tag{41.3}
\]
with \(R_\nu=\operatorname{diag}(1,\ldots,-1,\ldots,1)\).  The sixteen-segment
family of (41.1) is mapped to itself by the hypercubic group up to a cell
translation, which is the origin of the phase.
\end{lemma}

\begin{proof}
The segment from \(x\) to \(x+2e_\mu\), \(x\in\{0,1\}^4\), is sent by the
reflection \(x\mapsto1-x\) about the cell centre to the reversed segment
from \(1-x-2e_\mu\) to \(1-x\), which is a segment of the family based at the
cell shifted by \(-2e_\mu\) when \(\mu=\nu\); the shift produces
\(e^{-ik_\nu}\) on the \(\nu\)-th coarse component and the reversal the sign.
The executable verifies (41.3), permutation covariance for all \(24\)
permutations, and conjugation at every quarter-turn point.  The base-point
rules of V33--V34 have no reflection symmetry (Record~\ref{rec:v35-verifier}),
because their paths must start at the base point.
\end{proof}

\subsection{The explicit tower}

For a lift \(x=k+2\pi M\), \(M\in\mathbb Z^4\), and \(j\ge1\) put
\[
 F^{(j)}_\nu(x):=\prod_{l=1}^j\bigl|1+e^{ix_\nu/2^l}\bigr|^2=\frac{\sin^2(k_\nu/2)}{\sin^2(x_\nu/2^{j+1})},
                                                                 \tag{41.4}
\]
with the convention \(F^{(j)}_\nu=4^j\) when \(k_\nu=0\) and \(M_\nu\equiv0\pmod{2^j}\),
and \(F^{(j)}_\nu=0\) when \(k_\nu=0\) and \(M_\nu\not\equiv0\).  The identity
(41.4) is the telescoping \(\prod_{l=1}^j2\cos(x/2^{l+1})=\sin(x/2)/\sin(x/2^{j+1})\)
together with \(\sin^2(x_\nu/2)=\sin^2(k_\nu/2)\).

\begin{theorem}[Explicit level-\(j\) propagators]
\label{thm:v37-explicit}
For every \(j\ge0\) and \(k\ne0\),
\[
 G^{(j)}(k)=\Pi(k)\operatorname{diag}\bigl(\gamma^{(j)}_\mu(k)\bigr)\Pi(k),\qquad
 \gamma^{(j)}_\mu(k)=16^{-3j}\sum_{M\bmod2^j}\frac{\prod_{\nu}F^{(j)}_\nu(k+2\pi M)\;F^{(j)}_\mu(k+2\pi M)}{\lambda\bigl((k+2\pi M)/2^j\bigr)}.
                                                                 \tag{41.5}
\]
\end{theorem}

\begin{proof}
By (38.5) and Lemma~\ref{lem:v37-ward}, \(G^{(j)}=16^{-j}\Pi\sum_MW_MG^{(0)}(p_M)W_M^*\Pi\)
with \(W_M=\prod_lB(e^{ix/2^l})=(\prod_ls_l)\prod_l\Delta_l\), and the inner
projector of \(G^{(0)}(p_M)=\Pi(p_M)/\lambda(p_M)\) is removed by the iterated
Ward identity \(W_Mq(p_M)=(\prod_ls_l)q(k)\) and \(\Pi(k)q(k)=0\).  The product
of the diagonal matrices \(\Delta_l\Delta_l^*\) is \(\operatorname{diag}(F_\mu)\)
and \(\prod_l|s_l|^2=16^{-2j}\prod_\nu F_\nu\).  The executable verifies (41.5)
against the recursion exactly at the fifteen nonzero corners for \(j=1\) and
at fifteen Pythagorean momenta for \(j=2\), where all phases are Gaussian
rational.
\end{proof}

\begin{theorem}[The fixed point in closed form]
\label{thm:v37-fixed-point}
Let \(\sigma_\nu(k,M)=\sin^2(k_\nu/2)/(k_\nu+2\pi M_\nu)^2\), with
\(\sigma_\nu=\frac14\delta_{M_\nu,0}\) when \(k_\nu=0\).  Then \(G^{(j)}\to G^*\)
uniformly on \(\mathbb T^4\setminus\{0\}\), where
\[
 G^*(k)=\Pi(k)\operatorname{diag}\bigl(\Gamma_\mu(k)\bigr)\Pi(k),\qquad
 \Gamma_\mu(k)=1024\sum_{M\in\mathbb Z^4}\frac{\prod_\nu\sigma_\nu(k,M)\,\sigma_\mu(k,M)}{|k+2\pi M|^2},
                                                                 \tag{41.6}
\]
an absolutely convergent series with positive terms.  Moreover
\(\mathcal L[\Pi](k)=\frac14\Pi(k)\operatorname{diag}\bigl(1-\tfrac18|q_\mu(k)|^2\bigr)\Pi(k)\le\frac14\Pi(k)\),
so Theorem~\ref{thm:v35-fixed-point} applies with \(\rho=\frac14\) and the
finite constant \(d_1=\sup\|G^{(1)}-G^{(0)}\|\).
\end{theorem}

\begin{proof}
Choose the residues \(M\bmod2^j\) in the symmetric window
\(|M_\nu|\le2^{j-1}\), so that \(|x_\nu|\le2^j\pi\) for \(x=k+2\pi M\).  Then
\(y=x_\nu/2^{j+1}\) satisfies \(|y|\le\pi/2\), and \((2y/\pi)^2\le\sin^2y\le y^2\)
gives \(4^{j+1}\sigma_\nu\le F^{(j)}_\nu\le(\pi^2/4)4^{j+1}\sigma_\nu\) and
\(4^{-j}|x|^2\ge\lambda(x/2^j)\ge(4/\pi^2)4^{-j}|x|^2\).  Hence each term of
(41.5) is bounded by \(1024(\pi^2/4)^5\prod_\nu\sigma_\nu\sigma_\mu/|x|^2\),
which is summable over \(\mathbb Z^4\) (it decays like \(|M|^{-12}\)), and
converges termwise to the corresponding term of (41.6) as \(j\to\infty\);
dominated convergence gives the limit, uniformly in \(k\) because the
dominating series is bounded uniformly on the torus.  For the contraction
constant, (39.8) with diagonal blocks gives
\(\mathcal L[\Pi]=\frac1{16}\Pi\operatorname{diag}\bigl(\sum_\epsilon|s_\epsilon|^2|1+w_{\epsilon,\mu}|^2\bigr)\Pi\),
and with \(|1+w_\nu|^2=4\cos^2(p_{\epsilon,\nu}/2)\) the two aliases in each
direction contribute \(\cos^2(k_\nu/4)+\sin^2(k_\nu/4)=1\) for \(\nu\ne\mu\) and
\(16(\cos^4+\sin^4)(k_\mu/4)=16(1-\tfrac12\sin^2(k_\mu/2))\) for \(\nu=\mu\), so
the diagonal entry is \(4(1-\tfrac12\sin^2(k_\mu/2))=4(1-\tfrac18|q_\mu|^2)\).
The executable verifies this identity exactly at the corners and at
Pythagorean points.  Finiteness of \(d_1\) follows from (41.5) at \(j=1\):
the alias \(\epsilon=0\) gives \(\frac14\prod_\nu\cos^2(k_\nu/4)\operatorname{diag}(\cos^2(k_\mu/4))/\lambda(k/2)\),
whose difference from \(\Pi/\lambda(k)\) is bounded near \(0\) because
\(\lambda(k/2)=\lambda(k)/4+O(|k|^4)\) and the cosines are \(1+O(|k|^2)\); the
other aliases have \(\lambda(p_\epsilon)\ge2\).
\end{proof}

\begin{theorem}[Global nondegeneracy and uniform kernel convergence]
\label{thm:v37-nondegenerate}
For every \(k\ne0\), \(\Gamma_\mu(k)\ge256/\pi^{12}\) for all \(\mu\), hence
\[
 G^*(k)\ge\frac{256}{\pi^{12}}\,\Pi(k),\qquad
 \|\widehat S^*(k)\|\le\frac{\pi^{12}}{256},\qquad
 \|\widehat S^{(j)}(k)-\widehat S^*(k)\|\le\frac{d_1\,4^{-j}}{(3/4)\,c^2},\ \ c=\frac{256}{\pi^{12}}-\frac{4d_1}{3}4^{-j},
                                                                 \tag{41.7}
\]
for all \(j\) with \(c>0\).  In particular the kernels converge uniformly on
the whole torus and the first hypothesis of Theorem~\ref{thm:v34-card} holds
for the block-average rule on \(\mathbb T^4\setminus\{0\}\), with rate \(\frac14\).
\end{theorem}

\begin{proof}
Every term of (41.6) is nonnegative; keep \(M=0\).  For \(0<|k_\nu|\le\pi\),
\(\sigma_\nu(k,0)=\sin^2(k_\nu/2)/k_\nu^2\ge1/\pi^2\), and for \(k_\nu=0\) it is
\(\frac14\ge1/\pi^2\); also \(|k|^2\le4\pi^2\).  Hence
\(\Gamma_\mu\ge1024\,\pi^{-10}/(4\pi^2)=256/\pi^{12}\).  A diagonal matrix with
entries at least \(c_0\) is at least \(c_0I\), and compressing by \(\Pi\) gives
\(G^*\ge c_0\Pi\).  The kernel bounds follow from \(G^{(j)}\ge G^*-\frac{d_1\rho^j}{1-\rho}\Pi\)
and \(A^+-B^+=A^+(B-A)B^+\) on \(q^\perp\), as in Corollary~\ref{cor:v35-kernels}.
\end{proof}

\begin{theorem}[The axes are unchanged]
\label{thm:v37-axes}
On the coordinate axes \(k=(t,0,0,0)\) the block-average tower coincides with
the V33 tower at every level: the transverse kernels are the \(s_j(t)\) of
Theorem~\ref{thm:v33-closed-form}, and the fixed point is
\[
 \Gamma_2(t,0,0,0)=4\sin^2(t/2)\sum_{M\in\mathbb Z}\frac1{(t+2\pi M)^4}=\frac{2+\cos t}{12\sin^2(t/2)},
                                                                 \tag{41.8}
\]
i.e.\ \(\widehat S^*_{22}=12\sin^2(t/2)/(2+\cos t)\).
\end{theorem}

\begin{proof}
For \(k_\nu=0\), \(\nu\ge2\), the transverse alias weights
\(|s_\epsilon|^2|1+w_{\epsilon,\mu}|^2\) vanish unless \(\epsilon_\nu=0\) for
\(\nu\ge2\), and for the two remaining aliases equal
\(|1+e^{i(t/2+\pi\epsilon_1)}|^2=4\cos^2(t/4),\ 4\sin^2(t/4)\), which are the
weights of (37.7); the axis recursions therefore coincide at every level.  The
executable confirms the level-one and level-two values \(8\) and \(\frac{32}{3}\)
at \((\pi,0,0,0)\).  In (41.6) with \(k=(t,0,0,0)\) only \(M=(M_1,0,0,0)\)
contributes, with \(\prod_\nu\sigma_\nu\sigma_2=\sigma_1\cdot\frac1{64}\cdot\frac14\)
and \(|x|^2=(t+2\pi M_1)^2\), giving \(4\sin^2(t/2)\sum_{M_1}(t+2\pi M_1)^{-4}\).
Differentiating \(\sum_M(t+2\pi M)^{-2}=1/(4\sin^2(t/2))\) twice gives
\(\sum_M(t+2\pi M)^{-4}=(2+\cos t)/(48\sin^4(t/2))\), whence (41.8).
\end{proof}

\begin{record}[Corner values and the comparison with the symmetrized rule, non-hashed]
\label{rec:v37-corners}
At \((\pi,0,0,0)\) both rules have the transverse fixed-point propagator
\(\frac1{12}\) (kernel \(12\)).  At \((\pi,\pi,\pi,\pi)\) the block-average
fixed point has the transverse propagator \(0.005123\) (kernel \(195.2\)),
against \(0.03294\) (kernel \(30.4\)) for the symmetrized rule at level five;
the two rules have different fixed points off the axes.  The sampled
transverse floor of \(G^*\) over \(1500\) random momenta is \(0.0056\), near
the corner \((\pi,\pi,\pi,\pi)\), whose exact value \(0.005123\) is the
smallest found, against the proved floor \(2.8\times10^{-4}\).  At
\((\pi,\pi,0,0)\) and \((\pi,\pi,\pi,0)\) the transverse eigenvalues of the
block-average fixed point are \(0.0139,0.0351,0.0351\) and
\(0.0079,0.0079,0.0202\).  The sampled \(d_1\) is \(0.124\); the level-five
recursion agrees with (41.6) to \(1.3\times10^{-4}\).  The exact level-two
kernel at \((\pi,\pi,\pi,\pi)\) is \(\frac{14336}{443}(3I-J)\), transverse
eigenvalue \(\frac{57344}{443}\approx129.4\), on its way to the fixed-point
value \(\approx195.2\).
\end{record}

\subsection{The vertical covariance in diagonal form}

\begin{corollary}[Explicit vertical covariance]
\label{cor:v37-vertical}
For the block-average rule, (40.4) reads
\[
 \operatorname{Cov}^{\rm vert}_{\epsilon\epsilon'}(k)
 =\frac{\delta_{\epsilon\epsilon'}}{16}G^{(j)}(p_\epsilon)
 -\frac{\bar s_\epsilon s_{\epsilon'}}{256}\,G^{(j)}(p_\epsilon)\Delta_\epsilon^*\,\widehat S^{(j+1)}(k)\,\Delta_{\epsilon'}G^{(j)}(p_{\epsilon'}),
                                                                 \tag{41.9}
\]
with \(G^{(j)}\) given by (41.5) and \(\widehat S^{(j+1)}\) the transverse
inverse of the diagonal-compressed (41.5); all factors are explicit, and by
Theorem~\ref{thm:v37-nondegenerate} the family is uniformly bounded along the
tower.
\end{corollary}

\begin{proof}
Theorem~\ref{thm:v36-vertical} with \(B_\epsilon=s_\epsilon\Delta_\epsilon\).
\end{proof}

\subsection{Executable provenance}

\begin{record}[V37 verifier]
\label{rec:v37-verifier}
The executable
\begin{center}
\path{scripts/verify_ym_postarchive_v37_block_average_tower.py}
\end{center}
has 412 physical lines, 17326 bytes, and SHA--256
\[
\texttt{D1C7221E8973E1DA35BA6AB484F56D987056101D7EAC3CA5455D7220A31A9FCA}.
                                                               \tag{41.10}
\]
Its canonical payload SHA--256 is
\[
\texttt{BC00CFA8DBBB7BA569DAC8030A17EB39C7180E072F2B551B2877B5EE25E28AB7}.
                                                               \tag{41.11}
\]
The hashed payload records the exact level-one census at all \(256\)
quarter-turn momenta (Hermitian, transverse, rank three off the toron,
permutation and reflection covariant, conjugation under inversion), the
agreement of the incidence and alias routes at the corners, the
pseudo-inverse identities at levels one and two, the explicit formula (41.5)
against the recursion at the corners and at fifteen Pythagorean momenta, the
identity for \(\mathcal L[\Pi]\) at twenty points, and the axis values \(8\)
and \(\frac{32}{3}\).  The floating comparisons of Record~\ref{rec:v37-corners}
are stored under the separate non-hashed key.
\end{record}

\subsection{Decision for V38}

\begin{decision}[The Gaussian tower is the block-average tower]
\label{dec:v37-v38}
From V38 on the Gaussian half of the fixed-point card uses the block-average
rule; the symmetrized rule, its exact records, and the open inequality
(H\(_U\)) are retained as records only.  V38 should construct the first
transported vertex: the tree-level cubic vertex of the one-colour
commuting-toron background pushed through one block-average step by the
minimizer map \(M_\epsilon=\frac1{16}s_\epsilon^{\,*}G^{(j)}(p_\epsilon)\Delta_\epsilon^*\widehat S^{(j+1)}(k)\)
(a matrix now explicit by (41.5)), and state the level-\(j\) shell functional
in the form of the archived (16.5) with the vertical covariance (41.9) in
place of the fine propagator.  The oracle frame question of
Firewall~\ref{fire:v35-normalization} is to be settled on the axis, where
(41.8) is exact.
\end{decision}

\begin{assessment}[V37 acquisition]
\label{ass:v37}
V37 replaces the base-point path average by the block average and thereby
solves the Gaussian tower in closed form: explicit level-\(j\) propagators,
an explicit fixed point, contraction constant \(\frac14\) uniformly, global
nondegeneracy with the floor \(256/\pi^{12}\), uniform kernel convergence,
full hypercubic covariance, and an explicit vertical covariance.  The axis
fixed point of V33 is recovered exactly.  The first hypothesis of the
fixed-point card is now proved on the whole torus for this rule.
\end{assessment}

\begin{record}[V37 append record]
\label{rec:v37-append}
V37 is appended to the clean V36 whose installed record was
\[
\begin{array}{c|c}
\text{lines}&13242\\
\text{bytes}&539561\\
\text{SHA--256}&
\texttt{57DE9862277623D9AEE2DE7870ED86627E5620DB18DB3C0E8ACF5F1EF16817A5}.
\end{array}                                                    \tag{41.12}
\]
After validation only V37 is to remain the active numeric YM note.  The
next compulsory companion audit and the ten-version sweep are both V40.
\end{record}

\begin{firewall}[Claim boundary after V37]
\label{fire:v37-final}
The block-average rule is a linear, Gaussian-sector rule; its nonabelian
counterpart, whose linearization at the flat background it should be, is not
constructed here, and the transport of nonabelian vertices through it is the
next task.  The results are tree-level and exact for one blocking rule.  No
one-loop shell coefficient, no Lipschitz bound for blocked vertices, no
identification \(b^*=b_{\rm YM}\), no Haar sign, and no continuum statement is
made; the Gaussian anomaly, the interacting measure, regulator removal,
Osterwalder--Schrader reconstruction, physical clustering, and the
Hamiltonian Yang--Mills mass gap remain open.
\end{firewall}

% =============================================================================
% END APPENDED POST-TENTH-ARCHIVE V37 MATERIAL
% =============================================================================

% =============================================================================
% BEGIN APPENDED POST-TENTH-ARCHIVE V38 MATERIAL
% =============================================================================

\section{V38: harmonic extensions, the shell decomposition of the free
propagator, and the level-\(j\) shell functional}
\label{sec:v38}

With the block-average tower explicit (V37), the objects that the one-loop
shell coefficient of the fixed-point card needs can be written down in
closed form.  This note does three things.  It gives the composed minimizer
map (harmonic extension) from a level-\(j\) coarse field to the fine
lattice as an explicit diagonal chain formula.  It proves that the free
fine propagator decomposes exactly into positive shells
\(\Gamma_l=\mathcal C_l-\mathcal C_{l+1}\), where \(\mathcal C_l\) is the part
of the fine covariance explained by the level-\(l\) field; the shells are
the vertical covariances of V36 pushed down to the fine lattice, and they
telescope.  It then states the level-\(j\) shell functional in the archived
form (16.5) with the shell \(\Gamma_j\) in place of the vertical symbol and
the harmonically extended background in place of the background, records
that the archived response oracle is executable in the present environment
and reproduces its exact corner cards, and identifies the one missing
input for its evaluation: the fine vertex banks on the full edge set.

\subsection{Harmonic extensions}

Throughout, the rule is the block-average rule (41.1); \(G^{(j)}\) and
\(\widehat S^{(j)}=(G^{(j)})^+\) are its propagators and kernels, and for a
lift \(x=k+2\pi M\) the chain block at level \(l\) is
\(B(e^{ix/2^l})=s_l\Delta_l\) with \(s_l=\prod_\nu(1+e^{ix_\nu/2^l})/16\) and
\(\Delta_l=\operatorname{diag}(1+e^{ix_\mu/2^l})\); \(W_M=\prod_{l=1}^jB(e^{ix/2^l})\).
Level-\(l\) fields are normalized so that the level-\(l\) coarse field of a
level-zero Gaussian field with covariance \(G^{(0)}\) has covariance
\(G^{(l)}\); in these units the level-\((l-1)\) alias amplitudes of a
level-\(l\) field have covariance \(\frac1{16}G^{(l-1)}(p_\epsilon)\delta_{\epsilon\epsilon'}\).

\begin{theorem}[Composed minimizer map]
\label{thm:v38-harmonic}
For \(k\ne0\) at level \(j\), the constrained minimizer of the level-zero
energy among fine fields whose \(j\)-fold block average is the coarse field
\(c\) (transverse) has, at the lift \(M\) (fine momentum \(p_M=(k+2\pi M)/2^j\)),
the amplitude
\[
 a(p_M)=\mathcal M^{(j)}_M(k)\,c,\qquad
 \mathcal M^{(j)}_M(k)=16^{-j}\,G^{(0)}(p_M)\,W_M^*\,\widehat S^{(j)}(k).      \tag{42.1}
\]
It reproduces the constraint, \(\Pi(k)\sum_MW_M\mathcal M^{(j)}_M=\Pi(k)\), and
its energy equals the coarse energy, \(16^j\sum_M\mathcal M_M^*\widehat S^{(0)}(p_M)\mathcal M_M=\widehat S^{(j)}(k)\).
The one-step maps compose: with \(M^{(l+1)}_\epsilon=\frac1{16}G^{(l)}(p_\epsilon)B_\epsilon^*\widehat S^{(l+1)}\) of
Theorem~\ref{thm:v36-vertical}, the product of the one-step maps along a
chain equals (42.1), because every intermediate factor
\(\widehat S^{(l)}G^{(l)}=\Pi\) acts as the identity on the vectors it meets.
\end{theorem}

\begin{proof}
Composition: \(a^{(0)}=M^{(1)}_{\epsilon_1}M^{(2)}_{\epsilon_2}\cdots M^{(j)}_{\epsilon_j}c\)
contains the products \(\widehat S^{(l)}(p^{(l)})G^{(l)}(p^{(l)})=\Pi(p^{(l)})\)
applied to \(B^{(l+1)*}u\) with \(u\perp q(p^{(l+1)})\); by Lemma~\ref{lem:v37-ward},
\(q(p^{(l)})^*B^{(l+1)*}u=(B^{(l+1)}q(p^{(l)}))^*u=\bar s\,q(p^{(l+1)})^*u=0\), so the
projector is the identity there and (42.1) follows.  The constraint:
\(\Pi\sum_MW_M\mathcal M_M=16^{-j}\Pi\sum_MW_MG^{(0)}(p_M)W_M^*\widehat S^{(j)}=G^{(j)}\widehat S^{(j)}=\Pi\)
by (38.5).  The energy: the identity \(\widehat S^{(0)}(p)G^{(0)}(p)=\Pi(p)\) gives
\(16^j\sum_M\mathcal M_M^*\widehat S^{(0)}\mathcal M_M=\widehat S^{(j)}\bigl[16^{-j}\Pi\sum_MW_MG^{(0)}W_M^*\Pi\bigr]\widehat S^{(j)}=\widehat S^{(j)}G^{(j)}\widehat S^{(j)}=\widehat S^{(j)}\).
The executable verifies the constraint and energy identities exactly at the
fifteen nonzero corners for \(j=1\), and the composition identity for
\(j=2\) at six Pythagorean momenta, where all chain phases are Gaussian
rational.
\end{proof}

\subsection{The shell decomposition}

\begin{definition}[Explained covariance and shells]
\label{def:v38-shells}
On the alias class of \(k\) at level \(j\) (fine momenta \(p_M\), \(M\bmod2^j\)),
the covariance of the harmonic extension of the level-\(j\) Gaussian field
is
\[
 \mathcal C_j(M,M')=\mathcal M^{(j)}_MG^{(j)}\mathcal M^{(j)*}_{M'}
 =16^{-2j}\,G^{(0)}(p_M)\,W_M^*\,\widehat S^{(j)}(k)\,W_{M'}\,G^{(0)}(p_{M'}),          \tag{42.2}
\]
and \(\mathcal C_0(M,M')=16^{-j}G^{(0)}(p_M)\delta_{MM'}\) is the free covariance in
the same units.  For \(l<j\), \(\mathcal C_l\) is defined on the same class by
grouping its fine momenta into level-\(l\) classes.  The shells are
\(\Gamma_l:=\mathcal C_l-\mathcal C_{l+1}\).
\end{definition}

\begin{theorem}[Shell decomposition]
\label{thm:v38-shells}
For every \(k\ne0\) and \(j\ge1\), on the alias class of \(k\),
\[
 \mathcal C_0=\sum_{l=0}^{j-1}\Gamma_l+\mathcal C_j,\qquad \Gamma_l\ge0,               \tag{42.3}
\]
each \(\Gamma_l\) is the vertical covariance (41.9) of the level-\(l\) field
given the level-\((l+1)\) field, pushed down to the fine lattice by
\(\mathcal M^{(l)}\); it is annihilated by the level-\((l+1)\) composed
blocking, and it has rank \(45\cdot16^{\,l}\) per level-\((l+1)\) class.
Moreover \(\mathcal C_j\to0\) as \(j\to\infty\) at every fixed fine momentum,
so \(\mathcal C_0=\sum_{l\ge0}\Gamma_l\).
\end{theorem}

\begin{proof}
Gaussian conditioning: the level-\(l\) field equals the harmonic extension
of the level-\((l+1)\) field plus an independent vertical fluctuation whose
covariance is (40.4); pushing the level-\(l\) decomposition down to the fine
lattice by the linear map \(\mathcal M^{(l)}\) gives
\(\mathcal C_l=\mathcal C_{l+1}+\mathcal M^{(l)}\operatorname{Cov}^{\rm vert}_l\mathcal M^{(l)*}\),
and the last term is positive semidefinite and annihilated by the composed
blocking.  Summation over \(l\) gives (42.3).  For the decay,
\(\|\mathcal C_j(M,M)\|\le16^{-2j}\|G^{(0)}(p_M)\|^2\|W_M\|^2\|\widehat S^{(j)}\|\le16^{-2j}\cdot4^j\cdot\pi^{12}/256\cdot\|G^{(0)}(p_M)\|^2\)
by \(|s_l|\le1\), \(\|\Delta_l\|\le2\), and Theorem~\ref{thm:v37-nondegenerate}.
The executable verifies at six Pythagorean momenta and \(j=2\) the
telescoping identity \(\Gamma_0+\Gamma_1=\mathcal C_0-\mathcal C_2\) exactly on
the full class of \(256\) fine momenta, the exact positivity of both shells
on principal \(64\)-blocks, the exact annihilation of \(\Gamma_1\) by the
composed level-two blocking, and the full spectra of both shells in
floating arithmetic (smallest eigenvalues at rounding level, numerical
ranks \(720\) and \(45\), i.e.\ \(45\cdot16\) and \(45\)); at the fifteen
nonzero corners it verifies that \(\Gamma_0\) is Hermitian, exactly positive
semidefinite, annihilated, and of rank \(45\).
\end{proof}

\begin{remark}[Scale form]
\label{rem:v38-scale}
The shell \(\Gamma_l\) is supported on fine momenta whose level-\(l\) classes
are those of one level-\((l+1)\) momentum; in the fixed-point regime
\(G^{(l)}\approx G^*\) it is a dilation of a single shell up to the factor
\(4^{-l}\), the propagator scaling.  This is the block-spin analogue of the
momentum-shell decomposition of the continuum propagator, with the shell
\(l\) carrying the scales between \(2^{-l-1}\) and \(2^{-l}\) in lattice
units.
\end{remark}

\subsection{The level-\(j\) shell functional}

The ninth archive's density (16.5) is
\(\Omega_{j}(k,q)=\frac12\operatorname{Tr}[G_jQ_j]-\frac13\operatorname{Tr}[G_j(k+q\hat e_1)R_jG_j(k)R_j^*]\)
with \(G_j=A_j^{-1}\) the inverse of the \(45\times45\) vertical symbol and
\(Q_j,R_j\) the quadratic and linear background vertices of the level-\(j\)
action.  In the tower, at tree level in the interaction, the level-\(j\)
action is the fine action evaluated on the harmonic extension, so its
vertices are the fine vertices with every leg pushed down by
\(\mathcal M^{(j)}\), and the vertical symbol's inverse pushed down is the
shell \(\Gamma_j\).

\begin{definition}[Tree-level shell functional]
\label{def:v38-shell-functional}
Let \(Q_0(b)\) and \(R_0(b)\) be the fine quadratic and linear background
vertices (bilinear forms on fine fields, depending on a fine background
\(b\), with the colour trace already taken as in (42.12) of the eighth
archive).  For a coarse background \(\beta\) at level \(j+1\) with momentum
\(q\), put \(b=\mathcal M^{(j+1)}\beta\).  The level-\(j\) tree-level shell
functional is
\[
 \Omega^{\rm tree}_j(k,q):=\frac12\operatorname{Tr}\bigl[\Gamma_j(k)\,Q_0(b)\bigr]
 -\frac13\operatorname{Tr}\bigl[\Gamma_j(k+q\hat e_1)\,R_0(b)\,\Gamma_j(k)\,R_0(b)^*\bigr],       \tag{42.4}
\]
where the traces run over the alias class of \(k\) at level \(j+1\), and the
shell coefficient is \(b_{{\rm G},j}=\frac32\int\partial_q^2|_{q=0}\Omega^{\rm tree}_j(k,q)\,dk\)
as in (39.3).
\end{definition}

\begin{proposition}[What is and is not explicit]
\label{prop:v38-status}
In (42.4) the shells \(\Gamma_j\) and the extension \(\mathcal M^{(j+1)}\) are
explicit finite lattice sums by (42.1)--(42.2) and Theorem~\ref{thm:v37-explicit},
and they are uniformly bounded along the tower by
Theorem~\ref{thm:v37-nondegenerate}.  Consequently the map from the
propagator data to \(\Omega^{\rm tree}_j\) is Lipschitz with a constant
controlled by the norms of \(Q_0,R_0\) on the fine alias class, which is the
second hypothesis of Theorem~\ref{thm:v34-card} at tree level.  The fine
vertices \(Q_0,R_0\) are the archived response banks: the executable
\path{scripts/verify_ym_v60_response_oracle.py} runs in the present
environment, reproduces its exact sixteen-corner cards from the cached
V57 and V59 certificates, and returns the Wilson corner means
\[
 \bigl(\tfrac{176943}{627200},\ 0,\ -\tfrac{2822949104411}{265531392000}\bigr)
                                                                 \tag{42.5}
\]
for the three external-momentum jets.  Its banks, however, are stored on the
\(45\) retained edges of the parity-tree slice, on which the coarse
coordinates of the straight rule vanish; the block-average vertical space is
not a coordinate slice, so (42.4) requires the banks on the full edge set,
which the archived chain constructs before restriction
(\path{verify_ym_v21_incidence.py}, \path{verify_ym_v31_cubic_bubble_jets.py},
\path{verify_ym_v56_exponent_lift.py}).
\end{proposition}

\begin{proof}
Boundedness: \(\|\mathcal M^{(j)}_M\|\le16^{-j}\|G^{(0)}(p_M)\|\,2^j\,\pi^{12}/256\) and
\(\|\Gamma_j\|\le\|\mathcal C_0\|\) by positivity of the shells.  The Lipschitz
statement is the multilinearity of (42.4) in \(\Gamma_j\) and \(\mathcal M^{(j+1)}\).
The oracle statements are the executable's own checks
(\texttt{exact\_corner\_means\_reproduce\_V48}, \texttt{all\_sixteen\_exact\_total\_cards\_formed}),
re-run here; the restriction to \(45\) retained edges is the
\texttt{retained} filter of \path{verify_ym_postarchive_v4_constrained_alias_inverse.py}.
\end{proof}

\subsection{Executable provenance}

\begin{record}[V38 verifier]
\label{rec:v38-verifier}
The executable
\begin{center}
\path{scripts/verify_ym_postarchive_v38_shell_decomposition.py}
\end{center}
has 299 physical lines, 14487 bytes, and SHA--256
\[
\texttt{E0492D8FCAE95DE7058BAD5478CC332B30CFAF9B0489364217CE260ABB4A22ED}.
                                                               \tag{42.6}
\]
Its canonical payload SHA--256 is
\[
\texttt{B31855183BABF002A113033A0D90787A7747F026E5998C8122B1C0C40040CB36}.
                                                               \tag{42.7}
\]
The hashed payload records the exact constraint and energy identities and
the exact rank-\(45\) shell \(\Gamma_0\) at the fifteen nonzero corners, and
at six Pythagorean momenta the exact composition identity (42.1) for
\(j=2\), the exact telescoping \(\Gamma_0+\Gamma_1=\mathcal C_0-\mathcal C_2\) on
the \(256\)-momentum class, exact positivity on principal blocks, and exact
annihilation of \(\Gamma_1\).  The floating spectra and the explained-fraction
decay are stored under the separate non-hashed key.
\end{record}

\subsection{Decision for V39}

\begin{decision}[Fine vertex banks on the full edge set]
\label{dec:v38-v39}
V39 should rebuild the Wilson background vertices \(Q_0,R_0\) on the full
\(64\)-edge alias space from the archived incidence and word-jet modules,
verify them against the \(45\)-edge banks by restriction and against the
exact corner cards (42.5), and then evaluate (42.4) at level \(j=0\) for the
block-average rule at the sixteen corners exactly.  V40 is the compulsory
audit and ten-version sweep; the fixed-point limit of the shell coefficient,
\(b^*=\lim_jb_{{\rm G},j}\), and its comparison with \(b_{\rm YM}\) of
Record~\ref{rec:v35-normalization} are the targets after it.
\end{decision}

\begin{assessment}[V38 acquisition]
\label{ass:v38}
V38 makes the harmonic extension and the shell decomposition of the free
propagator explicit for the block-average tower, proves the telescoping and
positivity of the shells, states the tree-level shell functional in the
archived form with explicit propagator data, and locates the single missing
input, the fine vertex banks on the full edge set.  The second hypothesis of
the fixed-point card holds at tree level once those banks are bounded on
the fine alias class.
\end{assessment}

\begin{record}[V38 append record]
\label{rec:v38-append}
V38 is appended to the clean V37 whose installed record was
\[
\begin{array}{c|c}
\text{lines}&13597\\
\text{bytes}&557920\\
\text{SHA--256}&
\texttt{3E1BDBDD5446DC1D65C392200FA16E8DCB58C57D009963EBCE68E9C23D8388C0}.
\end{array}                                                    \tag{42.8}
\]
After validation only V38 is to remain the active numeric YM note.  The
next compulsory companion audit and the ten-version sweep are both V40.
\end{record}

\begin{firewall}[Claim boundary after V38]
\label{fire:v38-final}
The shell functional (42.4) is tree level in the interaction and is not yet
evaluated.  No one-loop shell coefficient, no value of \(b^*\), no
identification \(b^*=b_{\rm YM}\), no Haar sign, and no continuum statement
is made; the Gaussian anomaly, the interacting measure, regulator removal,
Osterwalder--Schrader reconstruction, physical clustering, and the
Hamiltonian Yang--Mills mass gap remain open.
\end{firewall}

% =============================================================================
% END APPENDED POST-TENTH-ARCHIVE V38 MATERIAL
% =============================================================================

% =============================================================================
% BEGIN APPENDED POST-TENTH-ARCHIVE V39 MATERIAL
% =============================================================================

\section{V39: the level-zero shell response of the block-average rule at the
sixteen corners}
\label{sec:v39}

V38 identified the missing input for the shell functional (42.4) as the fine
vertex banks on the full edge set.  This note obtains them, validates the
whole level-zero pipeline against the archived exact corner cards, and
then evaluates the block-average rule's level-zero Wilson response at the
sixteen parity corners in exact arithmetic.  The result is the first
number of the fixed-point programme: the corner mean of the block-average
shell response, to be compared with the archived straight-rule mean.

\subsection{The archived pipeline, rebuilt}

\begin{lemma}[Archived conventions]
\label{lem:v39-conventions}
In the archived response oracle
(\path{scripts/verify_ym_v60_response_oracle.py}): (i) the Hessian bank is
the plaquette Gram \(D_R^*D_R\) on the \(45\) retained edges as a Laurent
polynomial in \(z=e^{ik}\), and the vertical covariance used in the cards is
\(3(D_R^*D_R)^{-1}\), the archived normalization \(\frac13D_R^*D_R\) of
Record~\ref{rec:v35-normalization}; (ii) the background is the straight-rule
harmonic extension of the coarse polarization plane wave, with jets in the
external momentum along axis \(0\) at \(k=0\), the polarization being the
second link in direction \(2\) (the archived index \(1\) counts the second
links in edge-index order); (iii) external-momentum jets are true
derivatives with the imaginary unit factored out of odd orders, so the
bubble cross terms carry the coefficient \(-2\); (iv) the corner cards are
\(\Omega=\frac12\operatorname{Tr}[GQ(q)]+\frac13\operatorname{Tr}[G(k)P(q)G(k+q)M(q)]\)
as jets of order \(0,1,2\) in \(q\).
\end{lemma}

\begin{proof}
The V39 executable rebuilds \(D_R^*D_R\) from the V32 incidence data and
finds it equal to the archived bank entry by entry; it rebuilds the
straight-rule background jets \(a_R=-A^{-1}Be_{\rm pol}\), \(a_C=e_{\rm pol}\)
and finds them equal to the archived router fields on all \(64\) edges at
orders \(0,1,2\); and with the covariance \(3A^{-1}\) and true complex jets it
reproduces the archived exact Wilson cards (tadpole, bubble, total, all
three orders) at all sixteen corners.  With \(A^{-1}\) in place of
\(3A^{-1}\) the tadpole and bubble come out divided by \(3\) and \(9\).
\end{proof}

\begin{theorem}[Full-edge vertex banks]
\label{thm:v39-banks}
The archived assemblers accept the full column set.  The Wilson cubic
symbols assembled on all \(64\) edges and converted to Laurent jets restrict
on the \(45\) retained columns exactly to the archived V59 certificate, and
the lifted Wilson quartic assembled on \(64\) edges restricts exactly to the
archived \(45\)-column bank; the \(19\) further columns (the fifteen tree
edges and the four second links) carry \(378\) nonzero cubic entries over
the three orders.  Consequently the vertex forms are available on the
\(49\)-dimensional slice space (retained and second links).
\end{theorem}

\begin{proof}
Restriction is column selection; the executable compares the restricted
Laurent dictionaries entry by entry.
\end{proof}

\subsection{The block-average rule in the slice}

In tree gauge the fine field has the \(45\) retained coordinates \(a_R\), the
four second links \(a_C\), and zero tree links.  For the block-average rule
(41.1) the coarse link reads \(Xa_R+Ya_C\) with \(Y=\frac18I\) and
\(X_{\mu,(y,\nu)}=\delta_{\mu\nu}\frac1{16}(1+y_\mu+z_\mu(1-y_\mu))\), so
\(W=Y^{-1}X=8X\) is a Laurent polynomial of degree one.  By
Theorem~\ref{thm:v33-level-map}, with \(\widetilde A=A-BW-W^*B^*+W^*CW\) and
\(\widetilde B=B-W^*C\):

\begin{proposition}[Vertical covariance and background of the block rule]
\label{prop:v39-block-slice}
The vertical fluctuation of the block-average rule is \((a_R,-Wa_R)\) with
\(a_R\) of covariance \(3\widetilde A^{-1}\) in the archived normalization;
the harmonic extension of the coarse polarization \(e_{\rm pol}\) is
\(a_R=-8\widetilde A^{-1}\widetilde Be_{\rm pol}\), \(a_C=8e_{\rm pol}-Wa_R\).
Its external-momentum jets at \(k=0\) have the archived reality structure
(real, imaginary, real), and the vertex banks assembled with this
background on all \(64\) edges, restricted to the slice and compressed by
\(T=\begin{pmatrix}I\\-W\end{pmatrix}\) (with the jets of \(W\)), are the
effective vertices \(T^*QT\), \(T^*PT\), \(T^*MT\) of the block-average
level-zero shell functional (42.4).
\end{proposition}

\begin{proof}
Minimizing the slice form with \(a_C=Y^{-1}c-Wa_R\) gives the stated
minimizer and the reduced form \(\widetilde A\); the conditional covariance
of \(a_R\) given \(c\) is \(\widetilde A^{-1}\) in the \(D^*D\) normalization.
The fluctuation's second-link components are determined by the constraint,
which is the map \(T\).  The executable verifies the reality structure of the
background jets exactly.
\end{proof}

\subsection{The exact corner cards}

\begin{theorem}[Block-average level-zero Wilson response at the corners]
\label{thm:v39-cards}
With the archived conventions, the block-average rule's level-zero Wilson
response cards at the sixteen corners are exact rationals, recorded in the
payload, and their means over the corners are
\[
 \bigl(\overline\Omega_0,\ \overline\Omega_1,\ \overline\Omega_2\bigr)_{\rm block}
 =\bigl(\tfrac{5587}{1568},\ 0,\ -\tfrac{254367096889}{47416320000}\bigr),
                                                                 \tag{43.1}
\]
against the archived straight-rule means
\[
 \bigl(\tfrac{176943}{627200},\ 0,\ -\tfrac{2822949104411}{265531392000}\bigr)
 \approx(0.2821,\ 0,\ -10.632).                                    \tag{43.2}
\]
Numerically the block-average means are \((3.5631,\ 0,\ -5.3645)\): the order-two mean is about half of the straight-rule value and the order-zero mean is about twelve times larger, so the two rules differ substantially already at level zero, as their different vertical covariances (Record~\ref{rec:v37-corners}) suggest.  At the toron corner the block-average total card is \((-\tfrac{11}{16},0,-\tfrac{4823}{3072})\) and at \((\pi,\pi,\pi,\pi)\) it is \((\tfrac{123}{16},0,-\tfrac{7727}{1024})\).
\end{theorem}

\begin{proof}
Exact arithmetic throughout: Laurent jets of \(\widetilde A\) at each corner,
exact inversion and jet inversion, effective vertices by exact products,
traces.  The straight-rule branch of the same code reproduces (43.2)
exactly, which validates every convention used in the block branch.
\end{proof}

\begin{remark}[What the number means]
\label{rem:v39-meaning}
The order-two mean is the corner average of \(\partial_q^2\Omega\), i.e.\ of
the level-zero shell functional's second external-momentum derivative in the
polarization direction; by (39.3) its torus average times \(\frac32\) is the
level-zero shell coefficient in the unit-Wilson normalization, with the
oracle frame question of Firewall~\ref{fire:v35-normalization} still open.
The corner average is a sixteen-point quadrature of the torus integral and
is only a first diagnostic; the tower's shell coefficients at level \(j\)
and their limit \(b^*\) are the target.
\end{remark}

\subsection{Executable provenance}

\begin{record}[V39 verifier]
\label{rec:v39-verifier}
The executable
\begin{center}
\path{scripts/verify_ym_postarchive_v39_block_rule_response.py}
\end{center}
has 435 physical lines, 18242 bytes, and SHA--256
\[
\texttt{F88FE539A76FC3C5897AE73300AF100494BA65461EE06E86B36469959703BDA3}.
                                                               \tag{43.3}
\]
Its canonical payload SHA--256 is
\[
\texttt{A8FB591A06AEF799E2238FD6EAFE349E0D162ADE3E85DC5E9A601778EF29E325}.
                                                               \tag{43.4}
\]
It depends on the archived modules and on the cached certificates
\path{certificates/ym_v57_completion_lift_checkpoint.json} and
\path{certificates/ym_v59_cubic_vertex_oracle.json}, whose SHA--256 values
the archived oracle checks.  All recorded quantities are exact rationals or
Gaussian rationals.
\end{record}

\subsection{Decision for V40}

\begin{decision}[Audit, sweep, and the level-\(j\) shell coefficients]
\label{dec:v39-v40}
V40 is the compulsory companion audit and the ten-version sweep of the
other programmes' notes.  After it, the shell functional (42.4) is to be
evaluated at levels \(j\ge1\) of the block-average tower: the shell
\(\Gamma_j\) of V38 replaces \(3\widetilde A^{-1}\), the background is the
level-\((j+1)\) harmonic extension, and the fine vertex banks of
Theorem~\ref{thm:v39-banks} are contracted on the alias class of \(k\).  The
corner quadrature is to be replaced by the torus integral, and the limit
\(b^*\) compared with \(b_{\rm YM}\) once the frame factor is settled.
\end{decision}

\begin{assessment}[V39 acquisition]
\label{ass:v39}
V39 rebuilds the archived one-loop response pipeline from the slice data,
proves its conventions by exact reproduction of the sixteen archived
corner cards, extends the vertex banks to the full edge set, and computes
the block-average rule's level-zero Wilson response at the corners exactly.
The programme now has every ingredient of the level-zero shell functional
for the rule whose Gaussian tower is explicitly solvable.
\end{assessment}

\begin{record}[V39 append record]
\label{rec:v39-append}
V39 is appended to the clean V38 whose installed record was
\[
\begin{array}{c|c}
\text{lines}&13875\\
\text{bytes}&572590\\
\text{SHA--256}&
\texttt{CDC39872DC193FADBBB04B923A0370825E8EE9C1A3B08E498B2D4516F67801DE}.
\end{array}                                                    \tag{43.5}
\]
After validation only V39 is to remain the active numeric YM note.  The
next compulsory companion audit and the ten-version sweep are both V40.
\end{record}

\begin{firewall}[Claim boundary after V39]
\label{fire:v39-final}
The corner cards are level-zero, tree-level-in-the-blocking quantities in
the archived normalization; they are not a shell coefficient of the tower,
not a torus integral, and not compared with \(b_{\rm YM}\).  No Haar sign, no
continuum statement, and no mass-gap statement is made; the Gaussian
anomaly, the interacting measure, regulator removal, Osterwalder--Schrader
reconstruction, physical clustering, and the Hamiltonian Yang--Mills mass
gap remain open.
\end{firewall}

% =============================================================================
% END APPENDED POST-TENTH-ARCHIVE V39 MATERIAL
% =============================================================================

% =============================================================================
% BEGIN APPENDED POST-TENTH-ARCHIVE V40 MATERIAL
% =============================================================================

\section{V40: compulsory companion audit, ten-version sweep, and the
programme for the level-\(j\) shell coefficients}
\label{sec:v40}

V40 is the compulsory fifth-version companion audit and the ten-version
sweep of the other programmes' notes.  It records the state of every
active companion note, extracts what could bear on the Gaussian tower and
its shell functional, and fixes the direction for V41--V49.

\subsection{Active sources inspected}

The active companion notes in \path{latex_notes} on 12 September 2026 were
\[
\begin{array}{c|r|r|l}
\text{source}&\text{lines}&\text{bytes}&\text{SHA--256}\\ \hline
\text{SM V21 (20th cycle)}&4093&179562&\texttt{00A6756579FD884034339D93BD41159EC12065B6FAD0216E6A21206C02291CD5}\\
\text{NS V26}&5319&278283&\texttt{82C887F8C2C69E4F28FAD7C3DC8DE5B9099CB5A4BF431EBA1FFF3E91935978AD}\\
\text{GRH V67}&18652&640849&\texttt{4A050D2FE5D6B29E11A878BD2EBC071A1FB1227F2636E98D7844D86C5F80AA58}\\
\text{YM-par V5}&1351&54174&\texttt{306F1BB84CFA8A8D47EA569EC17E9545F9867C8D32B548EC9D9C444B4F3C0512}
\end{array}                                                    \tag{44.1}
\]
Since the V35 audit the SM programme has closed its nineteenth cycle
(archived as \texttt{V100\_f19}, whose V94--V100 treat the fourth-order
cosmological constant, radiation universes, the massless and massive scalar,
and the electromagnetic standing wave in the plane-symmetric cosmology) and
started a twentieth cycle, now at V21; the GRH note advanced from V60 to
V67; the NS note and the parallel YM note are unchanged.  No new programme
appeared in the folder.

\begin{audit}[SM twentieth cycle, V1--V21]
\label{aud:v40-sm}
The twentieth SM cycle is a classical study of the plane-symmetric
\(SU(2)\) Einstein--Yang--Mills system: relative colour phases and
non-abelian generation of gravitational waves (V2--V4), an integrable
Hamiltonian reduction of the averaged colour dynamics (V6--V7), the
polarization sphere and the critical coupling \(e^2=2/3\) (V8--V9), the
Higgs sector (V11--V14), and the fermionic sector with a spin--isospin
locked doublet, its colour condensate, and the Yukawa coupling (V16--V21).
Its own audits (V5, V10, V15, V20) fingerprint the present YM note at V35
and state that the plane-symmetric system has no object in common with
the lattice tower; the same holds in the other direction.  The cycle's
Yang--Mills content is the classical field equations of two colour modes;
it contains no lattice, no Gaussian measure, no blocking rule, and no
one-loop quantity.  Nothing is imported.
\end{audit}

\begin{audit}[GRH V61--V67, NS V26, and the parallel YM note]
\label{aud:v40-others}
GRH V61--V67 treat the three-pair determinant with at least seven real
padding zeros, an Appell reformulation of the real-axis problem, a signed
weight and its Fourier symbol, and shifted parabolic operators; its
companion sweep at V65 fingerprints earlier YM versions and records a
warning taken from the YM trace work against replacing a global
cancellation by termwise bounds, which the present programme already
follows (Record~\ref{rec:v35-normalization} keeps the tadpole and bubble in
one card).  Nothing else is transferable.  NS V26 is unchanged since the
V20 audit.  The parallel YM note remains at its failure stop of V5, whose
terminal recommendation to keep compact flux and toron labels exact and to
assemble the continuum through quasi-local observables remains compatible
with the Gaussian-tower route.
\end{audit}

\begin{decision}[V40 audit and sweep decision]
\label{dec:v40-audit}
No estimate, constant, or executable is imported.  The direction of
Decision~\ref{dec:v39-v40} stands.  The next compulsory audit is V45; the
next sweep is V50.
\end{decision}

\subsection{State of the Gaussian tower and its shell functional}

\begin{record}[What V31--V39 established]
\label{rec:v40-state}
\begin{enumerate}[leftmargin=2.2em]
\item The frozen-chart certificate route of the archives cannot close the
programme's card (V31); the compact tower is the object (V32).
\item The base-point path rules have an exactly solvable axis
(V33), permutation covariance when symmetrized (V34), a unique off-axis
fixed point with rate \(\frac{17}{36}\) (V35), and an open global
nondegeneracy inequality (V36).
\item The block-average rule has diagonal alias blocks, a hypercubic
covariant tower that is an explicit positive lattice sum, rate \(\frac14\),
the global floor \(256/\pi^{12}\), uniform kernel convergence, the same axis
fixed point as V33, and an explicit shell decomposition
\(\mathcal C_0=\sum_l\Gamma_l\) of the free propagator (V37--V38).
\item The archived one-loop response pipeline is rebuilt and exactly
reproduced; the fine vertex banks exist on the full edge set; the
block-average level-zero Wilson corner means are
\((\frac{5587}{1568},0,-\frac{254367096889}{47416320000})\) (V39).
\end{enumerate}
\end{record}

\begin{decision}[Programme for V41--V49]
\label{dec:v40-programme}
\begin{enumerate}[leftmargin=2.2em]
\item V41: the level-one shell functional of the block-average tower at
the corners, exactly: \(\Gamma_1\) on the level-two alias class, the
level-two harmonic extension of the background, and the full-edge banks
contracted on the class; comparison with the level-zero cards of V39.
\item V42: the torus average of the level-zero and level-one functionals
by the quarter-turn quadrature and by an explicit lattice sum, and the
first estimate of the shell coefficient's level dependence.
\item V43--V44: the fixed-point form of the shell functional, with the
explicit \(\Gamma_j\) of V38 in the self-similar regime, and a bound on
\(|b_{{\rm G},j}-b^*|\) in terms of the rate \(\frac14\).
\item V45: companion audit; settlement of the frame factor of
Firewall~\ref{fire:v35-normalization} on the exact axis, where the
block-average tower and the archived quotient coincide at level one.
\item V46--V49: the comparison \(b^*\) against \(b_{\rm YM}\) of
Record~\ref{rec:v35-normalization}, and the nonlinear rows of the card.
\end{enumerate}
\end{decision}

\begin{assessment}[V40 acquisition]
\label{ass:v40}
The compulsory audit and sweep are complete; nothing is imported, and no
companion result conflicts with V31--V39.  The programme for the next
cycle of ten versions is fixed on the explicit block-average tower and its
shell functional.
\end{assessment}

\begin{record}[V40 append record]
\label{rec:v40-append}
V40 is appended to the clean V39 whose installed record was
\[
\begin{array}{c|c}
\text{lines}&14090\\
\text{bytes}&583282\\
\text{SHA--256}&
\texttt{28BA834376B95AEE1D6DBEEB31A1308ED09065013E641AAEE5D5291642D90022}.
\end{array}                                                    \tag{44.2}
\]
After validation only V40 is to remain the active numeric YM note.  The
next compulsory companion audit is V45; the next ten-version sweep is V50.
\end{record}

\begin{firewall}[Claim boundary after V40]
\label{fire:v40-final}
Reading a companion theorem does not prove its hypotheses in the present
construction, and none was imported.  No shell coefficient beyond the
level-zero corner cards, no value of \(b^*\), no identification
\(b^*=b_{\rm YM}\), no Haar sign, and no continuum statement is made; the
Gaussian anomaly, the interacting measure, regulator removal,
Osterwalder--Schrader reconstruction, physical clustering, and the
Hamiltonian Yang--Mills mass gap remain open.
\end{firewall}

% =============================================================================
% END APPENDED POST-TENTH-ARCHIVE V40 MATERIAL
% =============================================================================

% =============================================================================
% BEGIN APPENDED POST-TENTH-ARCHIVE V41 MATERIAL
% =============================================================================

\section{V41: the tree-gauge representative of the tower's shells and the
fixed-background convention}
\label{sec:v41}

The shell functional (42.4) contracts a vertical covariance with the fine
vertex banks.  V39 evaluated it at level zero in the slice (tree-gauge)
coordinates of the archive.  The tower's shells \(\Gamma_j\) of V38 are
transverse alias-space objects; to use them at levels \(j\ge1\) one must
know whether the card depends on the gauge representative of the
fluctuation and, if so, how to map the transverse representative to the
archived one.  This note answers both questions exactly.  The card is
representative dependent; the archived representative is reached by an
explicit linear map (zero tree links per cell, zero longitudinal coarse
link); with this map the alias route reproduces the archived cards and the
slice-route cards exactly.  The note also fixes the background convention
of the tower to the archived fine background, so that no new vertex banks
are ever needed, and records the level-zero block-rule cards in that
convention.

\subsection{Representative dependence}

\begin{proposition}[The card is not gauge invariant in the fluctuation]
\label{prop:v41-gauge-dependence}
Let \(\Gamma_0^{\perp}\) be the level-zero vertical covariance of a rule in
transverse alias coordinates, transported to edge coordinates by the frame
\(E_{(s,\mu),(\epsilon,\nu)}=e^{ip_\epsilon\cdot s}\delta_{\mu\nu}\) with
\(\operatorname{Cov}_{\rm edge}=\frac1{16}E\operatorname{Cov}(\hat a)E^*\) and
\(\operatorname{Cov}(\hat a)=3[G^{(0)}\delta-\frac1{16}G^{(0)}B^*\widehat S^{(1)}BG^{(0)}]\)
in the archived normalization.  The cards
\(\frac12\operatorname{Tr}[\operatorname{Cov}_{\rm edge}Q]+\frac13\operatorname{Tr}[\operatorname{Cov}_{\rm edge}(k)P\operatorname{Cov}_{\rm edge}(k+q)M]\)
differ from the archived (tree-gauge) cards already at order zero; at the
corner \((\pi,0,0,0)\) the straight rule gives
\((-\tfrac{118777}{51450},0,-\tfrac{3401839261}{1512630000})\) against the
archived \((\tfrac{636}{245},0,-\tfrac{1110886837}{43218000})\).
\end{proposition}

\begin{proof}
Direct exact evaluation (recorded in the session's probe; the executable of
this note evaluates the corrected route below).  The transverse and
tree-gauge fields are gauge equivalent by a fine gauge function vanishing at
the base points, so the discrepancy is the fluctuation-gauge dependence of
the vertex banks, which are not gauge invariant in their fluctuation legs.
\end{proof}

\subsection{The gauge map}

\begin{definition}[Tree-gauge representative]
\label{def:v41-map}
For a cell field \(a\) with coarse phases \(z\), let \(\phi\) be the gauge
function with \(\phi(0)=0\) obtained by integrating \(a\) along the fifteen
tree links from the base point, extended to the neighbouring cells by the
Bloch phases.  Put \(a'=a-d\phi\), whose tree links vanish, and
\(a''=a'-c\,v\), where \(v_{(s,\mu)}=q_\mu(k)[s_\mu=1]\) is the field of the
constant-per-cell gauge function and \(c=q(k)^*L(a')/\lambda(k)\) with \(L\)
the rule's coarse-link map.  Then \(a''\) has zero tree links and zero
longitudinal coarse link, and its transverse coarse link equals that of
\(a\).  The map \(T_{\rm g}:a\mapsto a''\) is linear and Laurent in \(z\).
\end{definition}

\begin{theorem}[The alias route with \(T_{\rm g}\) equals the slice route]
\label{thm:v41-map}
With \(\operatorname{Cov}^{\rm tree}=T_{\rm g}\operatorname{Cov}_{\rm edge}T_{\rm g}^*\) and
the archived fine background, the alias-route cards reproduce, at all
fifteen nonzero corners and at all three orders, the archived exact cards
for the straight rule and the slice-route cards of the block-average rule.
\end{theorem}

\begin{proof}
The executable computes both routes in exact arithmetic with order-two
Taylor jets in the external momentum (the class \(J2\) over Gaussian
rationals: the alias phases carry \(e^{iq/2}=1+\tfrac i2q-\tfrac18q^2\), and
all inversions are series inversions).  The slice-route field of V39 is the
unique representative with zero tree links and zero coarse link in the
straight case, respectively zero tree links and block-average constraint in
the block case; \(T_{\rm g}\) produces the same representative because the
transverse coarse components of a vertical field vanish and the residual
constant-per-cell gauge freedom is fixed by the longitudinal condition.  The
toron corner is excluded from the alias route because the constant alias has
vanishing kernel there and the series in \(q\) would need a pole; the slice
route is regular at the toron and is used for it.
\end{proof}

\subsection{The fixed-background convention}

\begin{decision}[Background of the tower's shell functional]
\label{dec:v41-background}
At every level the shell functional (42.4) uses the archived fine
background: the router extension of the coarse polarization plane wave in
direction \(2\), with its external-momentum jets at zero, in tree gauge.
The level-\((j+1)\) harmonic extension of V38 is not used as background.
\end{decision}

\begin{remark}[Why]
\label{rem:v41-why}
The harmonic extension of a coarse plane wave at level \(j+1\) contains
alias components at fine momenta \(2^{-j-1}(q+2\pi M)\), \(M\ne0\), whose
amplitudes are of order \(q\); the vertex banks are jets at zero background
momentum and cannot be contracted with such components.  With a fixed fine
background the banks of Theorem~\ref{thm:v39-banks} serve all levels, the
shells \(\Gamma_j\) enter through \(T_{\rm g}\Gamma_jT_{\rm g}^*\), and the
difference between the two background choices is a lattice-scale effect: at
level zero it changes only the order-two card, and the fixed-point limit
\(b^*\), being a long-wavelength coefficient, is expected to be insensitive
to it.  This expectation is an assumption to be tested along the tower, not
a theorem.
\end{remark}

\begin{theorem}[Block-average level-zero cards with the fixed background]
\label{thm:v41-cards}
By the slice route with the archived background, the block-average rule's
level-zero Wilson corner means are
\[
 (\tfrac{5587}{1568},\ 0,\ -\tfrac{407872449589}{47416320000})\approx(3.5631,\ 0,\ -8.6019),                 \tag{45.1}
\]
against \((\tfrac{5587}{1568},0,-\tfrac{254367096889}{47416320000})\approx(3.563,0,-5.365)\)
with the block-rule background (V39) and the archived straight-rule means
\((\tfrac{176943}{627200},0,-\tfrac{2822949104411}{265531392000})\approx(0.282,0,-10.632)\).
The orders zero and one coincide with V39, as they must; the order-two mean
differs by \(-\tfrac{3480847}{1075200}\approx-3.2374\).
\end{theorem}

\begin{proof}
Exact computation, all sixteen corners, recorded in the payload.
\end{proof}

\subsection{Executable provenance}

\begin{record}[V41 verifier]
\label{rec:v41-verifier}
The executable
\begin{center}
\path{scripts/verify_ym_postarchive_v41_tree_gauge_representative.py}
\end{center}
has 474 physical lines, 18828 bytes, and SHA--256
\[
\texttt{32F70755D6634F2C469270B8DCB9204DE74E4076DDAE4ACBBFDB749B817882E8}.
                                                               \tag{45.2}
\]
Its canonical payload SHA--256 is
\[
\texttt{A899B45EFD9E19A2B8290F73A724092CBA0CC27422F0E488738E11B9EDF5981A}.
                                                               \tag{45.3}
\]
The hashed payload records the full-edge banks' assembly, the sixteen exact
block-rule cards with the fixed background, the two agreement flags of
Theorem~\ref{thm:v41-map}, and the three corner means.
\end{record}

\subsection{Decision for V42}

\begin{decision}[The level-one shell functional]
\label{dec:v41-v42}
V42 evaluates (42.4) at level one for the block-average rule at the
nonzero corners of the level-two momentum: \(\Gamma_1=\mathcal C_1-\mathcal C_2\)
of V38 on the \(256\)-momentum class, mapped by the nested tree-gauge
representative (the level-one field in tree gauge on the level-one lattice,
then each fine cell in tree gauge), and contracted with the fixed-background
banks, which are block diagonal over the sixteen level-one classes; the
rank-three structure of \(\mathcal C_2\) keeps the cost to sixteen
\(64\times64\) products per corner plus low-rank corrections.
\end{decision}

\begin{assessment}[V41 acquisition]
\label{ass:v41}
V41 proves that the one-loop card depends on the fluctuation
representative, constructs the exact linear map to the archived
representative, verifies that the transverse tower route then reproduces
the archived and slice-route cards exactly, fixes the background convention
of the tower, and records the level-zero block-rule cards in that
convention.  The tower's shells can now be contracted with the archived
banks at every level.
\end{assessment}

\begin{record}[V41 append record]
\label{rec:v41-append}
V41 is appended to the clean V40 whose installed record was
\[
\begin{array}{c|c}
\text{lines}&14245\\
\text{bytes}&591089\\
\text{SHA--256}&
\texttt{B58BD26835D492C031ADC46DE301CDEF8B89CD4A0B84EC0C27D9A9B38861CA66}.
\end{array}                                                    \tag{45.4}
\]
After validation only V41 is to remain the active numeric YM note.  The
next compulsory companion audit is V45; the next ten-version sweep is V50.
\end{record}

\begin{firewall}[Claim boundary after V41]
\label{fire:v41-final}
The fixed-background convention is a definition whose insensitivity in the
fixed-point limit is an assumption.  The cards remain level-zero quantities;
no shell coefficient of the tower, no value of \(b^*\), no identification
\(b^*=b_{\rm YM}\), no Haar sign, and no continuum statement is made; the
Gaussian anomaly, the interacting measure, regulator removal,
Osterwalder--Schrader reconstruction, physical clustering, and the
Hamiltonian Yang--Mills mass gap remain open.
\end{firewall}

% =============================================================================
% END APPENDED POST-TENTH-ARCHIVE V41 MATERIAL
% =============================================================================

% =============================================================================
\section{V42: toron regularity before certification of shell integrals}
\label{sec:v42-toron}
% =============================================================================

\begin{record}[Resumption and scope]
The latest active YM source on resumption was V41, not the obsolete
frozen-chart CRT programme.  Its SHA--256 was
\texttt{2EEC02EB12520BAA0651AB07A171970809451FB0DEB83B596B084312F9EE4961}.
V31--V41 already supply the Gaussian blocking tower and the
fixed-background/tree-representative convention; they are not rederived
here.  This appendix corrects a regularity claim in V41 and supplies an
analytic gate between finite momentum evaluations and an integrated
second-order response.  It does not change the background convention.
The preceding body is retained, apart from the displayed active version.

The existing executable named
\path{scripts/verify_ym_postarchive_v42_level_one_shell.py}
was still running when this continuation was prepared; no completed
payload from it was available.  The already completed file named
\path{artifacts/ym_postarchive_v43_sector_conditioning.json}
is an exact check at one nonzero momentum, at external momentum zero,
without gauge maps.  Its six checks pass and its recorded canonical
digest is
\texttt{9A3F8F35A078F8BE088C56315663AB80D0B2A5DC455EBD6F5388E0C4BEC95897}.
Executable numbering does not certify completion of the corresponding
note version.  These existing jobs and files were preserved.
\end{record}

\subsection{The non-Laurent projector and its precise regularity}

Write \(z_\mu=e^{ik_\mu}\), \(q_\mu=z_\mu-1\), and
\(\lambda=q^*q\).  The symbol \(q\) in this subsection is the lattice
gradient vector, not the external scalar Taylor parameter.
Let \(T_0(k)\) be V41's tree-zeroing map before the final longitudinal
subtraction, and let \(L(k)\) be its straight or block-average coarse
link map.  Introduce the constant map \(R:\mathbb C^4\to\mathbb C^{64}\),
\[
 (Ru)_{s,\mu}={\bf1}_{s_\mu=1}u_\mu,\qquad
 P(k)=\frac{q(k)q(k)^*}{\lambda(k)}.
\]
For \(k\ne0\) on the torus, the implemented representative is exactly
\[
 T(k)=(I-RP(k)L(k))T_0(k).                         \tag{46.1}
\]
Here and below the zero momentum is taken modulo \(2\pi\).

\begin{proposition}[Correction of V41's Laurent assertion]
\label{prop:v42-not-laurent}
For either rule, \(T(k)\) is bounded and smooth off the toron, but has no
continuous extension to the toron as a map on all edge fields.
In particular V41's assertion that this map is Laurent in \(z\) is false.
It is rational on the punctured torus.  This correction does not
invalidate evaluations at nonzero momenta.
\end{proposition}
\begin{proof}
The tree contains only within-cell edges, so \(T_0R=R\).
Also \(LR=I_4\).  For the straight rule this follows by selecting the
second edge of the coarse link.  For the block rule each of the eight
boundary edges in a given direction has weight \(2/16\); hence their
weights sum to one.  Thus (46.1) implies the exact identity
\[
                  T(k)R=R(I-P(k)).                 \tag{46.2}
\]
If \(k=t e_1\), \(0<|t|<\pi\), then \(P=e_1e_1^*\) and
\(T(k)Re_1=0\).  If \(k=t e_2\), then \(P=e_2e_2^*\) and
\(T(k)Re_1=Re_1\ne0\).  The limits disagree.  A finite Laurent
polynomial on the unit torus is continuous, so cannot equal \(T\).
Finally \(\|P\|=1\) off zero, \(\|R\|=\sqrt8\), and \(L,T_0\) are
finite Laurent matrices and therefore bounded.  Equation (46.1) proves
boundedness and rationality.  No assertion about a restriction to a
special covariance range is needed or inferred from this counterexample.
\end{proof}

\begin{lemma}[Derivative and weak-derivative bounds]
\label{lem:v42-derivatives}
Let \(r=|k|\) in a sufficiently small coordinate ball about zero.  For
\(j=0,1,2\), finite constants \(C_j\), depending only on the chosen
finite maps, satisfy
\[
                  \|D^jT(k)\|\le C_j r^{-j}.       \tag{46.3}
\]
Assigning any value at zero gives an essentially bounded representative
in \(W^{1,p}(\mathbb T^4)\) for \(1\le p<4\), and in
\(W^{2,p}(\mathbb T^4)\) for \(1\le p<2\).  In particular
\(DT\in L^2\) and \(D^2T\in L^1\).
\end{lemma}
\begin{proof}
On \([-\pi,\pi]^4\),
\(\lambda=4\sum_\mu\sin^2(k_\mu/2)\ge4r^2/\pi^2\).
The numerator \(qq^*\) and its derivatives of orders zero, one, two
are respectively \(O(r^2),O(r),O(1)\).  The derivatives of
\(\lambda^{-1}\) of these orders are \(O(r^{-2}),O(r^{-3}),O(r^{-4})\).
The product rule gives \(D^jP=O(r^{-j})\).  Applying it to (46.1)
gives (46.3), since \(L,T_0\) have bounded derivatives.  Radial
integration uses the volume element \(r^3\,dr\,d\omega\), so a
derivative of order \(j\) lies in \(L^p\) when \(jp<4\).
These are weak derivatives as well: integrate by parts outside a ball
of radius \(\epsilon\).  For the first derivative the boundary term
is \(O(\epsilon^3)\); for differentiating a first derivative it is
\(O(\epsilon^2)\).  Both vanish.  Smoothness away from zero completes
the argument on the torus.
\end{proof}

\subsection{An actual sufficient condition for integrating response jets}

\begin{theorem}[Finite products of translated gauge maps]
\label{thm:v42-integrated-jets}
Fix finitely many velocities \(v_i\in\mathbb R^4\).  Let
\(A_i(k,t)\) be compatible finite matrix functions jointly \(C^2\) on
\(\mathbb T^4\times[-t_0,t_0]\).  Form a finite product, followed by
trace, of these functions and factors \(T(k+t v_i)\) or their adjoints,
in any fixed order.  Call it \(F(k,t)\).  Then
\[
 t\longmapsto F(\cdot,t)\text{ is }C^2\text{ into }L^1(\mathbb T^4).
\]
Consequently its torus integral is twice differentiable and both
derivatives are the integrals of the almost-everywhere product-rule
derivatives.  At \(t=0\), for sufficiently small \(r\),
\[
 |\partial_t^j F(k,0)|\le B_j r^{-j},\quad j=0,1,2.
\]
With normalized torus measure the omitted second-derivative integral
over a ball of radius \(\delta\) is bounded in absolute value by
\[
 \frac{\pi^2 B_2}{(2\pi)^4}\,\delta^2.             \tag{46.4}
\]
\end{theorem}
\begin{proof}
Translations are continuous in \(L^p\) for finite \(p\).  A first
derivative of the product has at most one differentiated gauge factor,
in \(L^2\); a second derivative has either one second derivative in
\(L^1\), or two first derivatives in \(L^2\).  The remaining factors
are uniformly bounded.  These products lie in \(L^1\), by
Cauchy--Schwarz in the two-first-derivative case.

For completeness, continuity of the resulting second derivatives
does not require translation continuity in \(L^\infty\).
If \(f_n\to f\) in \(L^1\), and \(g_n\) are uniformly bounded and
converge in measure to \(g\), then \(f_ng_n\to fg\) in \(L^1\).
Indeed split off \((f_n-f)g_n\), and use bounded convergence in measure
against the finite measure \(|f|\,dk\) on the remaining term.
The same argument applies to finitely many bounded factors.
Products of the two differentiated factors converge in \(L^1\) by
their \(L^2\) convergence.  Thus all proposed derivative expressions
are continuous in \(L^1\).

Mollify each gauge factor in \(k\).  Mollification preserves uniform
boundedness and converges in \(W^{1,2}\) and \(W^{2,1}\).
The preceding estimates let the smooth product and its first two
derivatives converge locally uniformly in \(t\) in \(L^1\): translation
orbits over a compact \(t\)-interval are compact in the relevant
\(L^p\) spaces, and their \(L^1\) products are uniformly integrable.
Passing twice through the fundamental theorem of calculus proves the
claimed \(C^2\) statement.  Integration is a bounded linear functional
on \(L^1\), so commutes with those derivatives.
Finally (46.3) and the product rule give the pointwise bounds at zero.
The area of the unit three-sphere is \(2\pi^2\); integrating
\(B_2r^{-2}\) from zero to \(\delta\) gives (46.4).
\end{proof}

\begin{firewall}[What this criterion does not establish]
The theorem requires the remaining matrix factors to be jointly
\(C^2\), with finite norms.  No such statement about the actual shell
covariances or their harmonic lifts at every alias toron has been proved
here.  A growing covariance dimension also requires bounds uniform in
the level; the constants above are only finite-dimensional constants.
V37--V38's convergence of Gaussian kernels by itself is not a
\(C^2\) estimate for a response integrand.  Shifted singularities cannot
be handled by simply asserting a fixed pointwise dominator; the
translation argument above is the reason the stated theorem works.
\end{firewall}

\subsection{Sector normalization and the next upstream gate}

\begin{lemma}[A sector factor per additional covariance]
\label{lem:v42-sector}
Suppose a sector containing \(N\) cells has covariance
\(\Gamma=N Y C Y^*\).  For compatible vertices \(V_1,\ldots,V_r\),
\[
 \frac1N\Tr(\Gamma V_1\cdots\Gamma V_r)
 =N^{r-1}\Tr(CY^*V_1Y\cdots CY^*V_rY).             \tag{46.5}
\]
No positivity or invertibility is needed for this identity.
\end{lemma}
\begin{proof}
Substitute the displayed expression for each of the \(r\) covariance
factors and use cyclicity of trace, including the rectangular identity
\(\Tr(AB)=\Tr(BA)\).  The \(r\) factors of \(N\) and the cell average
\(1/N\) leave \(N^{r-1}\).
\end{proof}

For \(N=16\), the tadpole has no extra factor and the bubble has
factor 16.  This explains the normalization reported by the completed
conditioning check; it is not a new proof that the proposed
\(\Gamma=16YCY^*\) represents all gauged, externally shifted shells.
When the identities hold pointwise with two distinct shifted
covariances, the same substitution applies to the mixed bubble.

\begin{record}[Independent finite corroboration and companion check]
The new executable
\path{scripts/verify_ym_postarchive_v42_gauge_toron.py}
checks (46.2) for both rules at three exact Gaussian-rational phase
points.  All six checks pass.  Its source digest is
\texttt{e759aa88d90087de3158c1482f937968165fe25b1ff7ba2664c79b506af441a4};
the payload digest in
\path{artifacts/ym_postarchive_v42_gauge_toron.json} is
\texttt{6d0406d784671736ac191f91a583e6fa98a78bdd6ddead13c6c5de401f599650}.
The proof of the directional limits is analytic, not an inference from
these samples.  Floating shifted-grid averages already present in the
folder remain exploratory: neither a finite grid nor agreement of grids
is an enclosure of the integral or its level limit.

The updated SM endpoint is twentieth-cycle V42.  It supplies a
normalized reduced-model Yukawa benchmark and a distributional
large-circle density calculation, but explicitly leaves quantum
state-counting identification and the quantum continuum open.
The NS endpoint remains fourth-cycle V26, with pointwise rank-one
norms and no global regularity claim.  Neither supplies the missing
YM toron estimates.  This resumption check supplements, and does not
cancel, the mandatory companion audit at V45.
\end{record}

\begin{decision}[Next: covariance cancellation before more grids]
For each actual harmonic-lifted shell, determine the leading covariance
and its first two external derivatives near every alias toron.  Test
whether the range annihilates the directional term in (46.2), or whether
the full contracted integrand instead admits an integrable majorant.
Only after proving this with constants should a punctured-domain
quadrature and an explicit omitted-ball estimate be used as a
certificate.  If covariance factors are not \(C^2\), prove the needed
weighted replacement rather than invoking the preceding theorem.
Existing running calculations are retained as diagnostics, not as a
substitute for this gate.
\end{decision}

\begin{firewall}[Position after V42]
This appendix corrects the gauge map's regularity, proves its derivative
bounds, gives a sufficient integrated-jet theorem, and explains the
sector trace factor.  It proves no shell integral, no universal one-loop
coefficient, no interacting continuum measure, and no Yang--Mills
Hamiltonian mass gap.  The full goal remains open.
\end{firewall}

% END APPENDED POST-TENTH-ARCHIVE V42 MATERIAL

% =============================================================================
\section{V43: a quantitative level-zero covariance bound at the toron}
\label{sec:v43-cancellation}
% =============================================================================

The next step does not assume smoothness of a transverse frame through
zero.  Instead it uses the exact compatibility between the low-alias
gradient and the blocking map to bound the conditioned covariance.
The result concerns a whole punctured neighbourhood, not sampled points.

\begin{lemma}[Graph bound for a conditioned covariance]
\label{lem:v43-graph}
Let \(H=H_\ell\oplus H_h\) be finite-dimensional Hilbert spaces and
\(K=\operatorname{diag}(A,D)\ge0\).  Let
\(B=(B_\ell,B_h):H\to Z\), with \(B_\ell:H_\ell\to Z\) an
isomorphism.  Set
\[
 V=K-KB^*(BKB^*)^+BK,\quad
 S=B_\ell^{-1},\quad Jx=(-SB_hx,x).
\]
Then
\[
 0\le V\le JDJ^*,\qquad
 \|V\|\le(1+\|S\|^2\|B_h\|^2)\|D\|.             \tag{47.1}
\]
The bound does not depend on the size of \(A\).
\end{lemma}
\begin{proof}
Write \(F=BK^{1/2}\).  The matrix
\(F^*(FF^*)^+F\) is the orthogonal projector onto
\(\operatorname{Ran}F^*\), as follows immediately from a singular-value
decomposition of \(F\).  Thus \(0\le V\le K\) and \(BV=0\).
The kernel of \(B\) is exactly the graph of \(-SB_h\).
If \(P_h\) is the projection onto \(H_h\), this implies
\(V=J(P_hVP_h^*)J^*\).  Since \(P_hVP_h^*\le D\), the order bound
follows.  Finally \(\|J\|^2=1+\|SB_h\|^2\), proving (47.1).
This proof allows singular \(A,D\) and introduces no inverse at zero
momentum.
\end{proof}

\begin{theorem}[Uniform level-zero bound in the central half-cube]
\label{thm:v43-levelzero}
Let \(0\ne k\in[-\pi/2,\pi/2]^4\), \(p=k/2\), and
\(p_\epsilon=p+\pi\epsilon\) for \(\epsilon\in\{0,1\}^4\).
Use the free covariance
\(G_\epsilon=\Pi(p_\epsilon)/\lambda(p_\epsilon)\), where
\(\Pi\) is the transverse projector.  On the direct sum of the
sixteen transverse spaces set
\[
 K=\operatorname{diag}(G_\epsilon),\qquad
 B=(\Pi(k)B_\epsilon)_\epsilon,
 \qquad V=K-KB^*(BKB^*)^+BK.                       \tag{47.2}
\]
For either V41 blocking rule, with \(c=\cos(\pi/8)\),
\[
             \|V(k)\|\le\frac{1+15c^{-10}}{4c^2}. \tag{47.3}
\]
For the straight rule \(c^{-10}\) may be replaced by \(c^{-2}\).
In particular the \(|k|^{-2}\) divergence of the free low alias does
not survive in the norm of this conditioned covariance.
\end{theorem}
\begin{proof}
Write \(w_\mu=e^{ip_\mu}\), \(q_\ell=w-1\), and
\(q_c=w^2-1\).  The low-alias block is
\[
 B_0=f\,\operatorname{diag}(1+w_\mu),\qquad
 f=1\ \hbox{(straight)},\quad
 f=\frac1{16}\prod_\mu(1+w_\mu)\ \hbox{(block)}.
\]
It satisfies \(B_0q_\ell=fq_c\) exactly.  Every
\(|1+w_\mu|\ge2c\); hence \(|f|\ge c^4\) for the block rule and
\(\|B_0^{-1}\|\le(2c^5)^{-1}\), or \((2c)^{-1}\) for straight.

Take \(H_\ell=q_\ell^\perp\), \(Z=q_c^\perp\), and let
\(B_\ell=\Pi(k)B_0|_{H_\ell}\).  An explicit inverse is
\[
              S=\Pi(p)B_0^{-1}|_Z.               \tag{47.4}
\]
Indeed applying \(\Pi(k)B_0\) to this expression gives the identity
on \(Z\): the term subtracted by \(\Pi(p)\) is a multiple of
\(q_\ell\), which \(B_0\) sends into the null direction of
\(\Pi(k)\).  Both spaces have dimension three, so this right inverse
is the inverse, and its norm has the same bound as \(B_0^{-1}\).
This argument uses no choice of basis for these moving spaces.

Each high alias has some \(\epsilon_\mu=1\); in that direction
\(|e^{i(p_\mu+\pi)}-1|=|1+w_\mu|\ge2c\).
Thus its covariance has norm at most \((4c^2)^{-1}\), giving the
same bound for the block-diagonal high covariance \(D\).
For every alias, \(\|B_\epsilon\|\le2\), since the averaging scalar
has modulus at most one.  Projection cannot increase this norm.
The row of fifteen high blocks consequently has norm at most
\(2\sqrt{15}\), by Cauchy--Schwarz.  Applying (47.1) with these
three bounds yields (47.3) and its straight-rule refinement.
\end{proof}

\begin{corollary}[Bounded tree representative of the level-zero shell]
Let \(E(k)\) be the V41 alias-to-edge frame, so \(E^*E=16I\).
In that executable's colour convention the gauged edge covariance is
\[
          V_{\rm edge}=\frac3{16}T E V E^*T^*.
\]
It is uniformly bounded near the toron by
\(3\sup_k\|T(k)\|^2(1+15c^{-10})/(4c^2)\), with the
straight-rule refinement above.
\end{corollary}
\begin{proof}
The factor \(1/16\) cancels \(\|E\|^2=16\).
Use (47.3) and V42's boundedness of \(T\).
In (47.2), replacing \(B\) by \(B/4\) leaves the conditional
subtraction unchanged, so Fourier normalization does not alter this
comparison.  The prefactor three is specific to the stated code
convention, not a claim for every gauge group.
\end{proof}

\begin{assessment}[Acquisition and limitations]
V38 supplied positivity and telescoping of the Gaussian shells.
The new result here is an explicit toron-uniform bound for the
level-zero conditional covariance, independent of its singular free
low-mode variance.  Together with V42 it bounds the zeroth-order
level-zero traces with bounded vertex banks in this neighbourhood.
It does not prove continuity of the gauge representative at zero,
second-derivative bounds for the shell, or a level-uniform bound for
the fine-lifted higher shells.  The graph identity is an upstream
replacement for estimating the two separately divergent terms in
(47.2).  The next step is to differentiate an appropriately controlled
graph-covariance formula, and then address the harmonic lifts; the
V42 integrated-jet theorem cannot yet be applied to the actual response.
No extra numerical run is needed for this analytic bound.
\end{assessment}

\begin{record}[Active-version policy]
V43 retains the preceding V42 body, apart from the displayed active
version, and replaces its active numeric filename.  All archives and
the pre-existing running computations are preserved.  The next
mandatory companion audit remains V45.  Neither this bound nor the
finite checks establish the full Yang--Mills mass gap.
\end{record}

% END APPENDED POST-TENTH-ARCHIVE V43 MATERIAL

% =============================================================================
\section{V44: graph precision and integrable level-zero response jets}
\label{sec:v44-precision}
% =============================================================================

The predecessor V43 has SHA--256
\texttt{7E8245FBDAECB0AF0B9E19DDB7FC87D0A614D14701FE45FDA8CFD6476DA1E06E}.
Its body is retained apart from the displayed active version.  On
resumption the four existing YM Python jobs were confirmed live, and
no completed level-one exact payload was present.  The following
argument is analytic and independent of their eventual output.

\begin{lemma}[Exact covariance on a constraint graph]
\label{lem:v44-precision}
In Lemma~\ref{lem:v43-graph}, suppose \(A,D\) are positive definite
on their stated spaces.  Put \(U=-B_\ell^{-1}B_h\), \(Jx=(Ux,x)\).
Then
\[
 V=J WJ^*,\qquad W=(D^{-1}+U^*A^{-1}U)^{-1}.       \tag{48.1}
\]
No orthonormalization of the graph map \(J\) is implicit.
\end{lemma}
\begin{proof}
The graph support argument in V43 gives \(V=JWJ^*\) for some \(W\).
Directly from the conditional covariance formula and \(BJ=0\),
\(VK^{-1}J=J\).  Since \(J\) is injective, this implies
\(W(J^*K^{-1}J)=I\).  The matrix in parentheses is the positive
definite matrix \(D^{-1}+U^*A^{-1}U\); inversion proves the claim.
In particular there is no extra graph-volume factor in the covariance.
\end{proof}

\begin{theorem}[Symbol bounds for the level-zero shell]
\label{thm:v44-symbol}
For the covariance \(V(k)\) of Theorem~\ref{thm:v43-levelzero}, viewed
as a \(64\)-by-\(64\) ambient alias matrix with zero longitudinal
components, there are finite constants \(C_\alpha\) such that
\[
 \|\partial_k^\alpha V(k)\|\le C_\alpha |k|^{-|\alpha|},
 \qquad 0<|k|<r_0,\quad |\alpha|\le2.             \tag{48.2}
\]
The same conclusion holds for
\(V_{\rm edge}=(3/16)TEVE^*T^*\).  These are neighbourhood
bounds for either blocking rule, not finite-sample statements.
\end{theorem}
\begin{proof}
All notation is that of V43.  Let \(P_h(k)\) be the block-diagonal
projector onto the fifteen high-alias transverse spaces, embedded in
\(\mathbb C^{60}\), and let
\[
 Q_h(k)=\operatorname{diag}_{\epsilon\ne0}
                (\lambda(p_\epsilon) I_4).
\]
All high-alias gradients stay nonzero near zero.  Consequently
\(P_h,Q_h\) are smooth with bounded derivatives there, they commute,
and \(Q_h\ge4c^2I\) on the half-cube of V43.
Let \(\mathcal B_h\) be the row of the fifteen unprojected blocks
\(B_\epsilon\), and extend the graph map to the ambient high space by
\[
 U=-\Pi(p)B_0^{-1}\Pi(k)\mathcal B_h P_h,
 \qquad \mathcal Jx=(Ux,P_hx).                    \tag{48.3}
\]
Here \(UP_h=U\), its range lies in \(q_\ell^\perp\), and restriction
of \(\mathcal J\) to \(\operatorname{Ran}P_h\) is precisely the graph
map of (48.1).  Both low and coarse transverse projectors have
derivatives \(O(|k|^{-j})\), \(0\le j\le2\), by the quotient estimate
in V42, applied also at \(p=k/2\).  The matrix \(B_0^{-1}\) and its
derivatives are bounded near zero.  Thus
\(D^jU=O(|k|^{-j})\).

The low covariance on its transverse space is
\(A=\lambda(p)^{-1}I\).  Define on the full high ambient space
\[
 \mathcal H=Q_h+\lambda(p)U^*U,
 \qquad \mathcal W=P_h\mathcal H^{-1}P_h.          \tag{48.4}
\]
Since \(UP_h=U\), \(\mathcal H\) commutes with \(P_h\).  It is
uniformly positive, \(\mathcal H\ge4c^2I\).  Formula (48.1) on
\(\operatorname{Ran}P_h\), extended by zero on the longitudinal
directions, is therefore the exact identity
\[
                   V=\mathcal J\mathcal W\mathcal J^*. \tag{48.5}
\]
This also proves that (48.4) has introduced no longitudinal covariance.

The derivatives of \(\lambda(p)\) of orders zero, one, two are
\(O(|k|^2),O(|k|),O(1)\), respectively.  The product rule gives
\(D^j[\lambda(p)U^*U]=O(|k|^{2-j})\) for \(j\le2\).
Consequently \(\mathcal H\), its first two derivatives, and
\(\mathcal H^{-1}\) are bounded on the punctured neighbourhood.
For each coordinate derivative,
\[
 D\mathcal H^{-1}=-\mathcal H^{-1}(D\mathcal H)\mathcal H^{-1};
\]
differentiating again shows boundedness of every second derivative
of its inverse.  Hence \(\mathcal W\) has bounded derivatives through
order two.  Equation (48.5), with (48.3), proves (48.2).
The alias frame \(E\) is smooth on this local chart and has bounded
derivatives; V42 gives \(D^jT=O(|k|^{-j})\).  A final product-rule
application proves the asserted edge-covariance bounds.
\end{proof}

\begin{corollary}[Weak derivatives and local response integrals]
\label{cor:v44-integral}
Let \(\chi\) be a smooth cutoff supported in a sufficiently small
toron neighbourhood.  In the V41 fixed-background level-zero
convention, take smooth vertex matrices \(P(k,t),M(k,t),Q(k,t)\),
jointly \(C^2\), and a fixed momentum velocity \(v\).  Then
\[
 \begin{split}
 I_{\rm tad}(t)&=\int\chi(k)\Tr[V_{\rm edge}(k)Q(k,t)]\,dk,\\
 I_{\rm bub}(t)&=\int\chi(k)\Tr[V_{\rm edge}(k)P(k,t)
                 V_{\rm edge}(k+tv)M(k,t)]\,dk
 \end{split}                                                        \tag{48.6}
\]
are \(C^2\) for sufficiently small \(t\).  Their first two derivatives
are obtained by differentiating the integrands almost everywhere.
If \(B_2\) bounds \(|k|^2\) times the absolute second-derivative
integrand at \(t=0\), the omitted ball of radius \(\delta\) contributes
at most \(\pi^2B_2\delta^2\) with unnormalized measure, or the
bound (46.4) with normalized torus measure.
\end{corollary}
\begin{proof}
Choose a second smooth cutoff equal to one on the support of \(\chi\)
and its small translates.  Multiplication by it extends the local
edge covariance to a bounded symbol on the torus.  By (48.2), radial
integration and the vanishing boundary terms from V42, this symbol
belongs to \(W^{1,2}\cap W^{2,1}\).
The proof of Theorem~\ref{thm:v42-integrated-jets} used only these
two Sobolev properties and essential boundedness of each translated
factor; it did not use the special algebraic form of \(T\).
Apply that proof with the cutoff covariance in place of the gauge
factor.  The cutoff equals one at every occurrence in (48.6), so this
proves the assertion for the original integrals.  The pointwise
second-derivative integrand is \(O(|k|^{-2})\) by (48.2) and boundedness
of the vertices and their derivatives.  Radial integration gives
the stated omitted-ball estimate.  The same argument permits the
numerical colour and diagram prefactors in the code, absorbed in
\(B_2\).
\end{proof}

\begin{assessment}[What changed]
The toron integrability gate is now proved locally at level zero for
the stated fixed-background response and smooth banks: it is no
longer merely a conditional application requiring a smooth covariance.
The finite Laurent vertex banks meet the local smoothness hypothesis.
The result concerns integrability and differentiation, not a numerical
enclosure of the integral.  No explicit certified value of \(B_2\)
has been computed here.  It also does not identify a chosen finite
Taylor-jet implementation with every global branch convention.
\end{assessment}

\begin{decision}[Next upstream test]
Before extrapolating the observed higher-level averages, determine
whether the harmonic lifts of the level-one shell have the same
weighted derivative class.  Track every alias and the cell-volume
factors of (46.5).  Uniformity in the level is a separate problem;
neither the graph argument at one level nor a finite collection of
shell values supplies it.  The next version V45 also requires the
scheduled audit of the current SM and NS notes.
\end{decision}

\begin{firewall}[Position after V44]
An exact graph-precision formula, local level-zero derivative bounds,
and legitimacy of the local second-order response integral are proved.
The higher-shell limit, its normalization as a physical running
coupling, interacting regulator removal, reconstruction, and the
Yang--Mills Hamiltonian mass gap remain unproved.  Only V44 is retained
as the active numeric note; prior bodies and archives are preserved.
\end{firewall}

% END APPENDED POST-TENTH-ARCHIVE V44 MATERIAL

% =============================================================================
\section{V45: companion audit and fixed-level harmonic-lift regularity}
\label{sec:v45-fixedlevel}
% =============================================================================

\begin{record}[Required companion audit]
The inspected SM snapshot is twentieth-cycle V45, 355305 bytes,
SHA--256
\texttt{FA264CA71E5ED7CB8651E25216027F821D9D948F414F671650779D42682A448F}.
Its V44 derives occupation identities from the canonical
anticommutation relations and shows that equal occupations need not
determine conversion curvature.  Its V45 audit directs the next step
toward an exact one-excitation memory problem, not an assumed
Markovian closure.  These results do not construct the YM measure or
bound its Gaussian harmonic lifts.

The NS snapshot remains fourth-cycle V26, 278283 bytes,
SHA--256
\texttt{82C887F8C2C69E4F28FAD7C3DC8DE5B9099CB5A4BF431EBA1FFF3E91935978AD}.
It supplies pointwise rank-one Oseen-kernel norms and leaves its
integrated refinements and global regularity open.  No NS constant
is imported into the present finite-dimensional covariance estimates.
The audit therefore supports continuing the exact covariance
analysis, without identifying distinct theories or taking an
unproved limit.  The next mandatory companion audit is V50.
\end{record}

\subsection{A symbol class that survives a finite tower}

Write \(F\in\mathcal S^a_2\) near a puncture when \(F\) is smooth
there and, for every multi-index \(|\alpha|\le2\),
\[
 \|\partial^\alpha F(k)\|\le C_\alpha |k|^{a-|\alpha|}.
\]
Constants may depend on the blocking level.  The product rule gives
\(\mathcal S^a_2\mathcal S^b_2\subset\mathcal S^{a+b}_2\).
Smooth bounded matrices belong to \(\mathcal S^0_2\); the transverse
and longitudinal projectors belong to this class by V42.

\begin{theorem}[Propagators and inverses at each finite level]
\label{thm:v45-propagator}
For the V37 block-average Gaussian tower, at every fixed finite
level \(j\), the propagator \(G_j(k)\) is smooth and positive definite
on \(q(k)^\perp\) away from zero, with kernel spanned by \(q(k)\).
There exist \(a_j,b_j>0\) and a punctured neighbourhood such that
\[
 a_j|k|^{-2}\Pi(k)\le G_j(k)\le b_j|k|^{-2}\Pi(k),\qquad
 G_j\in\mathcal S^{-2}_2,\quad G_j^+\in\mathcal S^2_2. \tag{49.1}
\]
No uniformity of \(a_j,b_j\) or the derivative constants in \(j\)
is asserted.
\end{theorem}
\begin{proof}
For \(j=0\), \(G_0=\Pi/\lambda\), so all assertions follow from
\(\lambda\asymp|k|^2\) and the quotient differentiation in V42.
Suppose they hold at level \(j\).  The recursion is
\[
 G_{j+1}(k)=\frac1{16}\Pi(k)
       \sum_\epsilon B_\epsilon(k)G_j(k/2+\pi\epsilon)
                    B_\epsilon(k)^*\Pi(k).          \tag{49.2}
\]
In a small neighbourhood of zero, the fifteen high terms are smooth
with bounded derivatives.  The low term lies in \(\mathcal S^{-2}_2\).
For the low block the transverse map
\(B_\ell=\Pi(k)B_0|_{q(k/2)^\perp}\) has inverse
\(S=\Pi(k/2)B_0^{-1}|_{q(k)^\perp}\), by (47.4).
Both maps have bounded norm near zero.  The low term therefore
supplies a positive lower bound of order \(|k|^{-2}\) on the entire
coarse transverse space.  The low upper bound plus the bounded high
terms supplies the upper bound of that order.  This proves the
two-sided estimate and the derivative estimates for \(G_{j+1}\).

Smoothness and strict positivity off zero follow by the same low
term argument locally on the torus.  In a principal coarse chart
the low phases satisfy \(|p_\mu|\le\pi/2\), so \(B_0\) is invertible;
the other terms are positive semidefinite.  Across a chart boundary
the alias labels permute and the summed symbol agrees.  Thus no
additional zero eigenvalue has been assumed absent.

It remains to justify the inverse derivative estimate without a
smooth transverse frame through zero.  Let \(r^2=|k|^2\) in the
local chart, \(P_q=I-\Pi\), and
\(N=r^2G_{j+1}+P_q\).  The two-sided estimate makes \(N\) uniformly
positive and bounded.  It belongs to \(\mathcal S^0_2\).
Differentiating \(N^{-1}N=I\) once and twice proves
\(N^{-1}\in\mathcal S^0_2\).  On transverse and longitudinal
spaces separately one checks the exact identity
\[
                 G_{j+1}^+=r^2(N^{-1}-P_q).
\]
It proves \(G_{j+1}^+\in\mathcal S^2_2\) and completes the induction.
\end{proof}

\begin{theorem}[Finite-level conditional shells and harmonic lifts]
\label{thm:v45-lifts}
At each fixed level \(j\), form the unnormalized alias covariance
\(K_j=\operatorname{diag}_\epsilon G_j(p_\epsilon)\), the transverse
block row \(B=(\Pi(k)B_\epsilon)_\epsilon\), and
\[
 V_j=K_j-K_jB^*(BK_jB^*)^+BK_j,\qquad
 H_j=K_jB^*(BK_jB^*)^+.
\]
Then, as ambient alias matrices near zero,
\[
                 V_j\in\mathcal S^0_2,
 \qquad H_j\in\mathcal S^0_2.                      \tag{49.3}
\]
The high-alias rows of \(H_j\) in fact belong to
\(\mathcal S^2_2\).  Any fixed finite composition of these harmonic
lifts and the V41 gauge representatives, applied to a shell, retains
\(\mathcal S^0_2\) bounds locally at the coarsest toron.  The finite
Fourier and colour normalization factors do not change the orders.
\end{theorem}
\begin{proof}
For \(V_j\) use the graph argument of V44.  Its low covariance is now
\(A=G_j(k/2)|_{q(k/2)^\perp}\), whose transverse inverse, extended
by zero, belongs to \(\mathcal S^2_2\) by (49.1).
The graph map \(U=-\Pi(k/2)B_0^{-1}\Pi(k)\mathcal B_hP_h\) remains
in \(\mathcal S^0_2\).  The high covariance \(D\) is smooth and
positive on its transverse spaces, uniformly on a small compact
neighbourhood at this fixed level.  With \(D^+\) its ambient inverse
on those spaces, set
\[
 \mathcal H_j=D^++(I-P_h)+U^*A^+U.
\]
The first two terms are smooth and uniformly positive on the full
ambient high space; the last is positive semidefinite, commutes
with \(P_h\), and has derivatives of order \(m\le2\) bounded by
\(O(r^{2-m})\).  Thus \(\mathcal H_j^{-1}\) and its first two
derivatives are bounded.  The exact graph formula gives
\(V_j=\mathcal J P_h\mathcal H_j^{-1}P_h\mathcal J^*\), where
\(\mathcal Jx=(Ux,P_hx)\).  Product differentiation proves the
first assertion in (49.3).

For the harmonic lift, (49.2) gives \(BK_jB^*=16G_{j+1}\).
Its inverse on the transverse space is in \(\mathcal S^2_2\).
The low row of \(K_jB^*\) is in \(\mathcal S^{-2}_2\); the high
rows are in \(\mathcal S^0_2\).  Multiplication proves the claimed
orders, including the stronger high-row statement, without bounding
the two divergent terms in the conditional covariance separately.

For a fixed number of blocking steps, every intermediate alias
momentum is an affine function of the coarsest momentum.  Near its
toron only the all-low branch at a given intermediate stage can
approach that stage's zero; the other branches are smooth.  On the
low branch the linear rescaling changes finite constants but not
symbol orders.  There are finitely many branches at a fixed level.
V41's tree-zeroing map is smooth, and its longitudinal correction
has the \(\mathcal S^0_2\) bounds of V42.  The same bounds apply if
that correction is omitted while retaining the longitudinal coarse
component.  Products, adjoints, and finite sums now prove the
composition assertion.  No infinite-level product is taken.
\end{proof}

\begin{corollary}[Fixed-level local response differentiation]
For any fixed finite shell level, fine-lifted shell covariances in
the preceding theorem, contracted with smooth fixed-background
vertex banks in the chosen local chart, give locally integrable
second external-momentum derivatives.  With a smooth toron cutoff,
their integrated tadpole and bubble expressions are \(C^2\) in the
external momentum, and differentiation passes under the integral.
The omitted-ball second-derivative error is \(O(\delta^2)\), with a
constant depending on the level and banks.
\end{corollary}
\begin{proof}
The symbols in (49.3) are bounded, have first derivatives in
\(L^2\), and second derivatives in \(L^1\), by the radial and weak
derivative argument of V44.  Apply Corollary~\ref{cor:v44-integral}
and its translation proof to each finite trace product.
\end{proof}

\begin{decision}[Do not confuse a fixed-level theorem with a limit]
The finite-level toron obstruction is now controlled analytically.
The next acquisition must track constants with the level, or find
a rescaled norm in which the lifted response integrands converge.
In particular one needs uniform integrability of their second jets,
not just existence of each integral.  A sequence of finite certified
shell integrals or apparently stable grid averages cannot replace
that estimate.  The exact running level-one calculation is retained
as a normalization diagnostic; the present theorem gives no numerical
value for its integral or a universal running-coupling coefficient.
\end{decision}

\begin{firewall}[Position after V45]
The compulsory companion audit is complete and the Gaussian
fixed-level harmonic-lift regularity gate has an analytic proof.
No uniform-in-level response bound, nonlinear renormalization theorem,
interacting continuum measure, or Hamiltonian Yang--Mills mass gap
is established.  The full goal remains open.
\end{firewall}

\begin{record}[Append and next audit]
The predecessor V44 has SHA--256
\texttt{C91594A0AF9633A5DFD9D4D9D43BDFA795AE081D8AD05209B9FA01FA15E50B2F}.
Its body is retained apart from the displayed active version.
Only V45 remains the active numeric YM note; archives and companion
files are unchanged.  The next compulsory audit and ten-version
sweep are V50.
\end{record}

% END APPENDED POST-TENTH-ARCHIVE V45 MATERIAL

% =============================================================================
\section{V46: uniform Gaussian symbol bounds from an anisotropic majorant}
\label{sec:v46-uniform}
% =============================================================================

The predecessor V45 has SHA--256
\texttt{F1F7F2D19389BFAA45DD026CF1C432FB8B7694C386D5E020DB8C54B7D9381462}.
Its body is retained apart from the displayed active version.
The four previously identified YM computation processes were again
confirmed live; the following proof does not depend on their output.

\subsection{Correcting the summability argument}

V37's proof describes the summands in (41.6) as decaying like
\(|M|^{-12}\).  That description is not valid uniformly in direction:
if just one coordinate of \(M\) tends to infinity, the other coordinate
factors do not decay.  The conclusion can be justified instead by a
product majorant.  One must also remove the singular \(M=0\) term
before asserting a bounded majorant uniform in momentum.

Put \(n=2^j\), \(j\ge1\), take the fixed representatives
\(-n/2\le M_\nu<n/2\), and let \(x=k+2\pi M\),
\(k\in[-\pi,\pi]^4\).  Set
\[
 a_{n,\nu}(k,M)=\frac{\sin^2(k_\nu/2)}
                       {n^2\sin^2(x_\nu/(2n))},\qquad
 h_n(x)=4n^2\sum_\nu\sin^2(x_\nu/(2n)).
\]
Removable zero-over-zero values are assigned by continuity.  Formula
(41.5) becomes exactly
\[
 \gamma^{(j)}_\mu(k)=\sum_{M\bmod n} t_{n,\mu}(k,M),\qquad
 t_{n,\mu}=\frac{\prod_\nu a_{n,\nu}\,a_{n,\mu}}{h_n}. \tag{50.1}
\]

\begin{lemma}[Summable derivative majorant]
\label{lem:v46-majorant}
For \(|\alpha|\le2\), a constant independent of \(n,k,M\) bounds
\[
 |\partial_k^\alpha t_{n,\mu}(k,M)|
 \le C_\alpha\prod_{\nu=1}^4(1+|M_\nu|)^{-2},
 \qquad M\ne0.                                    \tag{50.2}
\]
For \(M=0\), on a fixed punctured neighbourhood of zero,
\(|\partial^\alpha t_{n,\mu}(k,0)|\le C_\alpha
 |k|^{-2-|\alpha|}\), again uniformly in \(n\).
\end{lemma}
\begin{proof}
Let \(s(y)=\sin y/y\), continuously extended at zero.  On the
chosen representatives \(|x_\nu/(2n)|\le3\pi/4\).  The function
\(s\) is strictly positive there, and it and its inverse have bounded
derivatives of every fixed finite order on that interval.
If \(M_\nu=0\), then
\[
 a_{n,\nu}=s(k_\nu/2)^2s(k_\nu/(2n))^{-2},
\]
which has uniformly bounded derivatives.  If \(M_\nu\ne0\), use
\[
 a_{n,\nu}=\frac{4\sin^2(k_\nu/2)}{x_\nu^2}
                         s(x_\nu/(2n))^{-2}.
\]
Here \(|x_\nu|\ge\pi|M_\nu|\), and each derivative through order two
is bounded by a constant times \((1+|M_\nu|)^{-2}\).
Thus all differentiated numerator products in (50.1) retain the
four-factor product decay in (50.2); the extra \(\mu\)-factor can
be bounded rather than needed for summability.

Uniformly on the representatives,
\(h_n\asymp|x|^2\), \(Dh_n=O(|x|)\), and \(D^2h_n=O(1)\).
Consequently \(D^m(h_n^{-1})=O(|x|^{-2-m})\) for \(m\le2\).
For \(M\ne0\), \(|x|\ge\pi\), proving (50.2) by the product rule.
For \(M=0\), the smooth bounded numerator and
\(h_n\asymp|k|^2\) give the stated singular bound.
Finally the right side of (50.2) is summable, since its sum factors
into four convergent one-dimensional sums.
\end{proof}

\begin{theorem}[Level-uniform Gaussian symbol estimates]
\label{thm:v46-symbol}
The constants in the local symbol bounds
\[
 G_j\in\mathcal S^{-2}_2,\qquad G_j^+\in\mathcal S^2_2
                                                               \tag{50.3}
\]
can be chosen independently of \(j\ge0\).  There are constants
\(a,b>0\), independent of \(j\), such that near zero
\[
       a|k|^{-2}\Pi(k)\le G_j(k)\le b|k|^{-2}\Pi(k).
                                                               \tag{50.4}
\]
On compact subsets away from zero, the matrices and their first two
derivatives are uniformly bounded.  Moreover
\(G_j\to G^*\) and \(G_j^+\to(G^*)^+\) in \(C^2\) on each such
compact subset.
\end{theorem}
\begin{proof}
Summing (50.2) gives a uniformly bounded \(C^2\) contribution from
\(M\ne0\).  The \(M=0\) contribution has the asserted singular
derivative bounds.  Multiplication by the two projectors \(\Pi(k)\)
in (41.5) preserves \(\mathcal S^{-2}_2\), uniformly in \(j\).
For the lower bound retain the nonnegative \(M=0\) term.
Since \(\sin(k_\nu/2)/(k_\nu/2)\ge2/\pi\) on this cube and
\(s(k_\nu/(2n))\le1\), every \(a_{n,\nu}(k,0)\ge4/\pi^2\).
Also \(h_n(k)\le|k|^2\), whence
\[
 \gamma^{(j)}_\mu(k)\ge\frac{(4/\pi^2)^5}{|k|^2}.
                                                               \tag{50.5}
\]
This is also a positive lower bound on every compact set away from
zero, uniformly in the level.  Include \(j=0\) separately using its
explicit free formula.  The inverse argument with
\(N=|k|^2G_j+P_q\) from V45 now has uniform constants and proves
the second assertion in (50.3).

For each fixed \(M\), all derivatives through order two in (50.1)
converge on compact sets away from the singular point to those of
the term in (41.6): \(a_{n,\nu}\to4\sigma_\nu\) and
\(h_n\to|x|^2\).  Every fixed \(M\) eventually belongs to the
representative window.  The summable derivative majorant controls
its complement uniformly; finite sums converge uniformly with their
derivatives.  Treat the single \(M=0\) term separately on the chosen
compact set.  This proves \(C^2\) convergence for \(G_j\).
The uniform transverse lower bound and ordinary inverse
differentiation give the inverse convergence.  Where a torus chart
changes, the alias sum is relabelled; the actual summed symbol is
smooth there, so these local estimates apply on a finite chart cover.
\end{proof}

\begin{corollary}[Uniform one-step estimates]
\label{cor:v46-onestep}
For the one-step alias covariances \(V_j\) and harmonic lifts \(H_j\)
defined in V45, their \(\mathcal S^0_2\) bounds can be chosen
uniformly in \(j\).  The high rows of \(H_j\) have uniform
\(\mathcal S^2_2\) bounds.  The matrices converge in \(C^2\) away
from the toron to the corresponding one-step matrices formed from
\(G^*\).  Local second-order response integrals formed from these
one-step matrices and fixed, smooth, finite-dimensional vertex
banks therefore converge to the integrals of the limiting jets.
\end{corollary}
\begin{proof}
In V45's graph-precision proof, the high covariance and its inverse
are now bounded uniformly, with two derivatives, on the fifteen
high-alias neighbourhoods.  Uniform boundedness of the high covariance
also gives a uniform positive lower bound for its inverse.  The
low inverse has the uniform \(\mathcal S^2_2\) bound of (50.3).
All geometric maps are independent of \(j\).  Thus every inverse
and product estimate in that proof has uniform constants.
The same formulas and the \(C^2\) convergence in the theorem prove
convergence away from zero.  Second-order trace integrands with fixed
smooth banks have a common bound \(C|k|^{-2}\) near zero.  Dominated
convergence proves convergence of their cutoff integrals.  This
does not require a rate of \(C^2\) convergence.
\end{proof}

\begin{firewall}[The growing lift is still a different problem]
The last corollary concerns a single blocking step with fixed-size
banks, not the programme's fine-lifted level-\(j\) Wilson banks.
Those involve \(j\) harmonic lifts, a growing alias space, derivatives
in the fixed fine-background momentum convention, and sector trace
normalizations.  A uniform bound for each factor need not bound
their growing product or its derivatives.  No limit of those actual
shell integrals, no value of \(b^*\), and no Yang--Mills mass gap
is claimed here.
\end{firewall}

\begin{decision}[Next acquisition]
Use the explicit composed harmonic extension from V38 to estimate
the normalized fine-space trace directly, rather than multiplying
one-step norm bounds.  Keep the low-alias singular term separate,
use a summable product majorant for the other aliases, and account
for derivatives of the background and the growing sector size.
This is the remaining distinction between the proved one-step
fixed point and the desired shell-response limit.
\end{decision}

\begin{record}[Version policy]
Only V46 remains the active numeric YM note.  Prior text and frozen
archives are retained; no running computation was restarted.
The next compulsory companion audit and ten-version sweep remain V50.
\end{record}

% END APPENDED POST-TENTH-ARCHIVE V46 MATERIAL

% =============================================================================
\section{V47: square-summable composed lifts and the scale factors in traces}
\label{sec:v47-composed}
\providecommand{\mathscr}{\mathcal}
% =============================================================================

The predecessor V46 has SHA--256
\texttt{74849C0791C414BCAC4098ABCCF9B19A4EEAFA03E24B3CB0817DA171079CC6CD}.
Its body is retained apart from the displayed active version.
The four existing YM computations were confirmed live on resumption.
This appendix estimates the exact transverse harmonic extension of
V38, before applying additional tree-gauge maps.

\begin{theorem}[Uniform square-summable composed-lift bound]
\label{thm:v47-lift}
Let \(n=2^j\), and stack the blocks \(\mathcal M^{(j)}_M(k)\) of
(42.1) into a map \(\mathscr M_j:\mathbb C^4\to\mathbb C^{4n^4}\).
For \(0<|k|<r_0\) and \(|\alpha|\le2\),
\[
 \|\partial_k^\alpha\mathscr M_j(k)\|_{\rm HS}
                  \le C_\alpha n^{-1}|k|^{-|\alpha|},       \tag{51.1}
\]
with constants independent of \(j\).  The same bound, without the
singular factor, holds on compact sets away from zero.
Use the fixed symmetric residue windows of V46 and extend the stacked
map by zero to \(\ell^2(\mathbb Z^4;\mathbb C^4)\).  Then
\(n\mathscr M_j\) converges in Hilbert--Schmidt norm, with its first
two derivatives locally uniformly away from zero, to the map with
blocks
\[
 \mathscr M_{\infty,M}(k)=
 \frac{\Pi_{\rm c}(x)}{|x|^2}\,
 \operatorname{diag}\!\left(\prod_\nu d_\nu(x)\,d_\mu(x)\right)^*
 (G^*(k))^+,\quad x=k+2\pi M,                         \tag{51.2}
\]
where \(\Pi_{\rm c}(x)=I-xx^{\mathsf T}/|x|^2\) and
\(d_\nu(x)=(e^{ix_\nu}-1)/(ix_\nu)\), with value one at zero.
No tree-gauge transformation is included in (51.1)--(51.2).
\end{theorem}
\begin{proof}
Define the normalized geometric sum
\[
 d_{n,\nu}(x)=\frac1n\sum_{a=0}^{n-1}e^{iax_\nu/n}
             =\frac{e^{ix_\nu}-1}{n(e^{ix_\nu/n}-1)}.
\]
The geometric-sum expression specifies all removable values.  The
diagonal chain formula in V38 gives exactly
\[
 n^{-1}W_M=\operatorname{diag}\left(
                  \prod_\nu d_{n,\nu}\,d_{n,\mu}\right),\qquad
 n\mathcal M^{(j)}_M=
 [n^{-2}G_0(x/n)]\,[n^{-1}W_M]^*\,G_j(k)^+.           \tag{51.3}
\]
For each coordinate, \(d_{n,\nu}\) and its derivatives through order
two are bounded uniformly by \(C(1+|M_\nu|)^{-1}\).
If \(M_\nu=0\), this follows from the geometric sum: a derivative
introduces \((ia/n)^r\), of modulus at most one.  Otherwise write
\[
 d_{n,\nu}=e^{i(n-1)x_\nu/(2n)}
       \frac{2\sin(x_\nu/2)}{x_\nu}
       s(x_\nu/(2n))^{-1},\qquad s(y)=\sin y/y.
\]
As in V46, \(|x_\nu|\ge\pi|M_\nu|\) and
\(|x_\nu/(2n)|\le3\pi/4\).  Differentiating this expression proves
the claimed bounds; all differentiated phase factors are bounded.

The free projector and denominator obey
\[
 \|\partial_k^\alpha[n^{-2}G_0(x/n)]\|
                  \le C_\alpha |x|^{-2-|\alpha|},\quad |\alpha|\le2.
\]
Indeed \(n^2\lambda(x/n)\asymp|x|^2\) on this window, with
first and second derivatives \(O(|x|)\) and \(O(1)\); the projector
quotient has derivatives \(O(|x|^{-|\alpha|})\).  These estimates
also follow by applying the quotient rule directly to
\(n(e^{ix/n}-1)\), whose norm is comparable to \(|x|\).
Together with the uniform \(\mathcal S^2_2\) bound for \(G_j^+\),
(51.3) bounds the \(M=0\) block by \(C_\alpha|k|^{-|\alpha|}\).
For \(M\ne0\), it gives, with a larger constant,
\[
 \|\partial_k^\alpha[n\mathcal M^{(j)}_M]\|
 \le C_\alpha |k|^{2-|\alpha|}
       \prod_\nu(1+|M_\nu|)^{-1}                  \tag{51.4}
\]
near zero.  The square of the product on the right is summable.
Summing the squared block norms, and using the fixed four-dimensional
domain to compare operator and Hilbert--Schmidt norms, proves (51.1).

For each fixed \(M\), \(d_{n,\nu}\to d_\nu\),
\(n^{-2}G_0(x/n)\to\Pi_{\rm c}(x)/|x|^2\), with two derivatives
away from \(x=0\).  V46 gives the corresponding convergence of
\(G_j^+\).  The squared summable majorant permits truncation to a
finite set of aliases uniformly in \(j\) and in a compact momentum
set.  Finite-block convergence followed by removal of that
truncation proves (51.2), including its two derivative statements.
\end{proof}

\begin{lemma}[Cell trace scaling]
\label{lem:v47-traces}
Let \(N=n^4\), let \(E_n^*E_n=NI\), and suppose
\(\|\mathscr M_j\|_{\rm HS}\le C/n\) at the momentum in question.
For a positive coarse covariance \(0\le C_j\le cI\), set
\(\Gamma_j=E_n\mathscr M_jC_j\mathscr M_j^*E_n^*\).
Then
\[
 \frac1N|\Tr(\Gamma_j Q)|\le cC^2 n^{-2}\|Q\|,
 \qquad
 \frac1N|\Tr(\Gamma_j P\Gamma_j R)|
                         \le c^2C^4\|P\|\|R\|.       \tag{51.5}
\]
These are estimates in the stated covariance and cell normalization;
they assume uniform bounds for the displayed vertex operators.
\end{lemma}
\begin{proof}
Positivity and \(E_n^*E_n=NI\) give
\(\Tr\Gamma_j\le Nc\|\mathscr M_j\|_{\rm HS}^2\).
Thus the first inequality follows from
\(|\Tr(\Gamma_jQ)|\le\Tr\Gamma_j\|Q\|\).
For the second use
\(|\Tr(\Gamma_jP\Gamma_jR)|\le
\|P\|\|R\|(\Tr\Gamma_j)^2\), and
\(Nn^{-4}=1\).  The inequality used here follows by expanding the
positive matrix \(\Gamma_j\) in an orthonormal eigenbasis and bounding
each matrix entry of \(P,R\) by the corresponding operator norm.
\end{proof}

\begin{assessment}[Why this is not yet the shell coefficient]
The composed transverse lift is now controlled without multiplying
a growing number of one-step norm bounds.  But (51.5) concerns a
bounded input covariance; \(G_j\) itself is not bounded at the toron.
For actual shells one must first use the bounded conditional
covariance and combine the finite coarse alias blocks.  Additional
gauge representatives are not unitary and must also be estimated.
The fine Wilson banks and their derivatives must satisfy the stated
operator bounds in precisely the chosen normalization before this
lemma can be applied to them.

There is also a separate momentum-scale issue.  If a fine momentum
is shifted by \(t/2\) while \(p=(k+2\pi M)/n\), the corresponding
coarse shift is \(k\mapsto k+nt/2\).  For any twice differentiable
function \(f\),
\[
 \left.\frac{d^2}{dt^2}f(k+ntv/2)\right|_{t=0}
                 =\frac{n^2}{4}(v\cdot\nabla_k)^2f(k).
                                                               \tag{51.6}
\]
Thus a coarse-momentum derivative bound cannot be substituted
unchanged for the code's fixed fine-background derivative.
Formula (51.6) alone does not determine the total response derivative,
which also differentiates vertices and other momentum arguments.
\end{assessment}

\begin{decision}[Next: normalized contracted symbols]
Estimate the fine Wilson contractions after substituting (51.3),
including gauge maps, rather than estimating each bank separately
in an unscaled space.  Track the factor in (51.6) and test which
Ward cancellations actually hold for the fixed background.  A
bounded zeroth-order bubble is not a uniform bound for its second
fine-momentum derivative.  Until those cancellations or bounds are
proved, no limit of the physical response is certified.
\end{decision}

\begin{firewall}[Position after V47]
A square-summable transverse composed-lift limit and conditional
trace estimates are proved.  They do not establish the gauged
Wilson response limit, an interacting continuum theory, or a
Yang--Mills Hamiltonian mass gap.  Only V47 remains active; previous
bodies and archives are retained.  The next companion audit is V50.
\end{firewall}

% END APPENDED POST-TENTH-ARCHIVE V47 MATERIAL

% =============================================================================
\section{V48: longitudinal loss in recursive push-down and a full-link lift}
\label{sec:v48-full-link}
% =============================================================================

The predecessor V47 has SHA--256
\texttt{A578DC1E7F3460DD530DBB7B270A39A546DDC02BACCC2272CED3A22EBE2DC0B5}.
Its body is retained apart from the active-version label.  Inspection
of the actual recursion, rather than another asymptotic estimate,
reveals a distinction between a quotient harmonic lift and a lift of
a specified gauge representative.  This distinction is material
because V41's vertex banks depend on the fluctuation representative.

\subsection{The projector that the recursion inserts}

The implemented harmonic block is
\(H_\epsilon=G_l(p_\epsilon)B_\epsilon^*G_{l+1}^+/16\).
It satisfies \(Hq_c=0\), because \(G_{l+1}^+q_c=0\).
In contrast, a tree-zeroing map generally adds a pure-gradient
component to a transverse field.  The higher-level recursion in
\path{scripts/ym_shell_float.py} converts the tree-gauged edge field
back to alias coordinates and then applies this harmonic map to
each alias block.  Thus it inserts a transverse projection on that
input.  The exact level-one script also contracts its top
tree-gauged covariance with maps containing \(G_1^+\).

\begin{proposition}[Exact vertical-vector witness]
\label{prop:v48-witness}
At top phase \(z=(-1,1,1,1)\), take fine phases
\(w_\epsilon=((-1)^{\epsilon_1}i,(-1)^{\epsilon_2},
(-1)^{\epsilon_3},(-1)^{\epsilon_4})\).
Let \(a\) have its only nonzero alias component equal to \(e_3\)
at \(\epsilon=(0,1,0,0)\).  Then \(a\) is transverse and has zero
block-average coarse link.  If \(E\) is the alias frame and \(T\)
is the implemented top tree representative, put \(y=E^{-1}TEa\).
With \(\boldsymbol\Pi=\operatorname{diag}_\epsilon\Pi(w_\epsilon)\),
\[
             \|(I-\boldsymbol\Pi)y\|^2=\frac32.       \tag{52.1}
\]
Hence the longitudinal loss occurs already on a vertical fluctuation,
not merely on an arbitrary vector outside the shell's support.
\end{proposition}
\begin{proof}
The input gradient has nonzero entries only in directions one and
two, so \(e_3\) is transverse.  The block-average scalar has the
factor \(1+w_2=0\); hence the coarse link of the input vanishes.
All phases are Gaussian integers and every nonzero gradient
denominator is rational.  The executable listed below constructs
the tree map exactly, multiplies the single input column, applies
the sixteen exact transverse projectors, and sums the squared
Gaussian-rational residual components.  There are four nonzero
components and their squared norms sum to \(3/2\).  This is a
finite exact arithmetic proof of this specified witness, not a
sampling proof of a generic identity.
\end{proof}

This does not invalidate V38's quotient minimizer formula or the
transverse estimates V43--V47.  Nor does it alone prove that two
gauge-related response integrals differ: such a conclusion would
require an identity or counterexample for the complete response.
It does show that the current recursion is not a right inverse on
the full coarse representative that enters it.  Gauge independence
of the fixed-background banks cannot be assumed to repair this.

\subsection{An explicit representative-level completion}

\begin{theorem}[Full-link harmonic lift]
\label{thm:v48-completion}
Let \(L\) be the physical four-component coarse-link map on cell
edges, \(R\) the boundary-edge injection of V42, and \(P_q\) the
coarse longitudinal projector.  Thus \(LR=I\).
Let \(H_e\) be the transverse harmonic lift in edge coordinates,
so \(\Pi_cLH_e=\Pi_c\), \(H_eq_c=0\).  Define
\[
 \widetilde H_e=H_e+R(I-LH_e),\qquad
       Y=T_{\rm keep}\widetilde H_e,                \tag{52.2}
\]
where
\(T_{\rm keep}=T_0-RP_qL(T_0-I)\) is the implemented tree map
that retains the coarse link.  Then
\[
 L\widetilde H_e=I,\qquad LY=I,\qquad
                   \widetilde H_eq_c=Rq_c.          \tag{52.3}
\]
The added term is a fine pure gauge, has zero free curl energy,
and leaves the transverse minimum energy unchanged.  At one fixed
blocking step it has \(\mathcal S^0_2\) bounds, uniformly in the
Gaussian level under V46's hypotheses.
\end{theorem}
\begin{proof}
Since \(\Pi_c(I-LH_e)=0\), the defect \(I-LH_e\) has longitudinal
range.  Thus \(R(I-LH_e)=RP_q(I-LH_e)\).  A field \(Rq_c\) is the
gradient of a Bloch scalar constant on the sites of each cell: it
vanishes on internal edges and equals \(z_\mu-1\) on boundary
edges in direction \(\mu\).  Hence the correction is pure gauge.
The identity \(LR=I\) proves the first equation in (52.3), and
\(H_eq_c=0\) proves the third.  The free curl of a gradient is zero,
so both its quadratic energy and its energy cross term vanish.

The difference \(T_0-I\) is a fine gradient map.  The Ward identity
for the block rule implies that \(L(T_0-I)\) has longitudinal range.
Consequently
\(LT_{\rm keep}=LT_0-P_qL(T_0-I)=L\), giving \(LY=I\).
The map \(T_{\rm keep}\) has zero tree edges: both its correction
and \(R(I-LH_e)\) are supported on boundary edges.  Finally \(L,R\)
and the one-cell alias frame are smooth bounded matrices;
V46 bounds \(H_e\) and V42 bounds \(P_q,T_{\rm keep}\) in
\(\mathcal S^0_2\).  The product rule proves the last assertion.
No growing composition bound is inferred from it.
\end{proof}

\begin{record}[Exact verifier]
\path{scripts/verify_ym_postarchive_v48_longitudinal_loss.py}
establishes the witness and checks \(L\widetilde H_e=I\) and
\(LT_{\rm keep}\widetilde H_e=I\) at the same nonzero phase.
All checks pass.  The source SHA--256 is
\texttt{fa2616975a8e479cea15688b6093380a9acb869231796b4a0b50b54b46e80781};
the canonical payload digest in
\path{artifacts/ym_postarchive_v48_longitudinal_loss.json} is
\texttt{a84ddbc3b5d12edd88fe1aaa542e4d5570a896a723418943031e703492af2f41}.
The general identities (52.2)--(52.3) have the analytic proof above;
the finite checks corroborate their implementation.
\end{record}

\begin{decision}[Upstream correction before interpreting higher shells]
The existing running scripts are retained unchanged as diagnostics
of their original convention.  Their output must not be called a
certificate for the full-link lift (52.2).  Next implement a separate
full-link route, verify the right-inverse identity at each push-down
stage, and compare its contracted response with the transverse-only
route.  The extra pure-gauge terms may disappear only if a proved
identity for the complete fixed-background response makes them do
so.  If they persist, the representative convention must be included
explicitly in any response-limit statement.  Changing a lift is
not itself a proof of physical gauge independence.
\end{decision}

\begin{firewall}[Position after V48]
The recursive representative mismatch has an exact witness, and a
full-coarse-link one-step completion has a rigorous construction.
Its iterated response, its gauge independence, nonlinear regulator
removal, and the Yang--Mills mass gap remain open.  Only V48 remains
active; earlier bodies, archives, and running scripts are preserved.
The next mandatory companion audit remains V50.
\end{firewall}

% END APPENDED POST-TENTH-ARCHIVE V48 MATERIAL

% =============================================================================
\section{V49: full-link composition and a representative-sensitive diagnostic}
\label{sec:v49-composition}
% =============================================================================

The predecessor V48 has SHA--256
\texttt{FABA6EBF3DFF78225DB8723BB5682CAAD94062DD7A540B1DD5C7E35CDFFFC555}.
Its body is retained apart from the active-version label.
The old level-zero/one shifted-grid process is now absent and its
artifact contains the requested grid sizes through ten; it was not
restarted.  The exact level-one and the higher-level diagnostic
processes were still live when checked.  Their original routes have
not been modified.

\begin{theorem}[Composition preserves full coarse representatives]
\label{thm:v49-rightinverse}
Let \(L_r:\mathcal H_{r-1}\to\mathcal H_r\) be successive full
coarse-link maps, and let \(Y_r:\mathcal H_r\to\mathcal H_{r-1}\)
be the completed maps (52.2), assembled on all cells or alias
blocks in compatible coordinates.  Put
\(\mathscr Y_j=Y_1\cdots Y_j\),
\(\mathscr L_j=L_j\cdots L_1\).  Then
\[
                         \mathscr L_j\mathscr Y_j=I.       \tag{53.1}
\]
The output has the finest-level zero-tree constraint, and every
intermediate reblocked field equals the corresponding intermediate
representative in the composition.  This identity also holds for
Taylor jets through any order for which the maps are differentiable
at the chosen nonzero momenta.
\end{theorem}
\begin{proof}
V48 proves \(L_rY_r=I\) for arbitrary inputs, including their
longitudinal parts.  In the product for (53.1), first cancel
\(L_1Y_1\), then \(L_2Y_2\), and continue.  Stopping the cancellation
at any intermediate stage proves the intermediate assertion.
The leftmost map \(Y_1\) has zero finest-level tree edges.
Coordinate changes cancel with their inverses and do not alter
these identities.  Differentiating (53.1), or computing in a
truncated Taylor algebra, proves the jet statement.  No assertion
about an infinite product or its operator norm is made.
\end{proof}

\begin{lemma}[What a pure-gauge change does to a quadratic bank]
\label{lem:v49-bank}
For two lift maps related by \(Y'=Y+D\) and a vertex matrix \(Q\),
\[
 Y'^*QY'-Y^*QY=Y^*QD+D^*QY+D^*QD.                 \tag{53.2}
\]
If \(D\) has pure-gradient range and both \(Q\) and \(Q^*\)
annihilate all such gradients, the right side vanishes.  The
vanishing free curl energy of \(D\) alone does not imply this
conclusion for another vertex bank.
\end{lemma}
\begin{proof}
Expand the left side.  Under the two annihilation hypotheses,
\(QD=Q^*D=0\), which kills all three terms.  Without them, even
\(Y=0\) and \(Q\) the orthogonal projector onto a nonzero gradient
direction can give \(D^*QD\ne0\), despite zero curl energy.
This finite-dimensional counterexample distinguishes the hypotheses;
it is not asserted to be a Wilson-bank counterexample.
\end{proof}

\subsection{Isolated implementation comparison}

The new diagnostic
\path{scripts/explore_ym_v49_full_link_comparison.py}
loads the existing floating shell implementation into an isolated
namespace and changes only its push-stage map to (52.2).  It checks
that its intended replacement occurs exactly once, and fails if
the upstream source no longer matches.  The original module and
running processes are unchanged.  At every completed stage it
evaluates \(L_rY_r-I\) through second-order jets.  Both original
and completed routes use the same vertex banks, momentum shift,
and diagram prefactors.  Source hashes and full complex arrays are
saved in \path{artifacts/ym_v49_full_link_comparison.json}.

At the single top momentum \(k=(0.7,-0.9,0.4,1.1)\), the first
comparison gave the following \emph{floating diagnostics}:
\[
\begin{array}{c|r|r|r|r}
\text{level}&\text{push stages}&\max\|L_rY_r-I\|_{\rm entry}
 &\Re\Omega_{\rm old}''&\Re\Omega_{\rm full}''\\\hline
1&16&3.34\,10^{-16}&-1.595415016&9.054236504\\
2&272&7.45\,10^{-16}&-0.851180413&3.455238860
\end{array}                                                    \tag{53.3}
\]
The residual maximum includes all three Taylor coefficients.
The response jets at this unsymmetrized single momentum also have
non-negligible imaginary parts, which are retained in the payload.
Only real parts are displayed in (53.3); these are not real torus
integrals, certified interval enclosures, or values of a physical
running-coupling coefficient.  Floating agreement with (53.1)
corroborates the implementation but is not its analytic proof.

\begin{assessment}[Consequence for the programme]
There is now a concrete implementation of the full-link route,
and the pointwise diagnostics do not support treating its response
as numerically interchangeable with the old route.  The rigorous
conclusion remains the right-inverse theorem and V48's exact loss
witness.  An exact nonzero response difference, or equality after
the complete torus integration and all required terms, has not
been established by this diagnostic.
\end{assessment}

\begin{decision}[Next exact test, not another extrapolated grid]
At nonzero Gaussian-rational phases, compare a corrected contracted
bank or complete response with the original route in exact
arithmetic.  Track separately the terms in (53.2), the bubble
cross terms, and the external-momentum jets.  Before interpreting
either convention as physical, determine which representative and
background Ward identity the intended effective action requires;
a full-link right inverse by itself does not select that identity.
V50 must also perform the scheduled companion audit and lineage
sweep, including this change in the response convention.
\end{decision}

\begin{firewall}[Position after V49]
Full-link composition is proved and an isolated numerical route
has been tested at two finite levels.  No representative-independent
Wilson response, nonlinear continuum construction, or Yang--Mills
mass gap is proved.  Only V49 remains active; previous bodies and
frozen archives are retained.
\end{firewall}

% END APPENDED POST-TENTH-ARCHIVE V49 MATERIAL

% =============================================================================
\section{V50: exact bank obstruction and compulsory lineage audit}
\label{sec:v50-audit}
% =============================================================================

\subsection{A finite exact comparison, beyond floating evidence}

\begin{proposition}[The quartic bank distinguishes the two lifts]
\label{prop:v50-bank}
At the coarse phase \(z=(-1,1,1,1)\), let \(Y_o=T_{\rm keep}H_e\)
be the original lift and \(Y_f\) the full-link lift (52.2), both
evaluated at external momentum zero.  Freeze these two lift matrices
while differentiating the V41 fixed-background quartic bank \(Q(t)\).
Write \(K_o^{(r)}=Y_o^*Q^{(r)}(0)Y_o\) and
\(K_f^{(r)}=Y_f^*Q^{(r)}(0)Y_f\).  Then
\[
 K_o^{(0)}=\operatorname{diag}(0,20,12,20),\qquad
 K_f^{(0)}-K_o^{(0)}=
 \begin{pmatrix}16&0&16&0\\0&0&0&0\\16&0&0&0\\0&0&0&0\end{pmatrix},
                                                               \tag{54.1}
\]
both first derivatives vanish, and
\[
 K_f^{(2)}-K_o^{(2)}=
 \begin{pmatrix}
 -49/3&-1&-49/3&-10/3\\
 -1&0&0&0\\-49/3&0&0&0\\-10/3&0&0&0
 \end{pmatrix}.                                                \tag{54.2}
\]
In particular this quartic bank does not annihilate the added
longitudinal representative in isolation.
\end{proposition}
\begin{proof}
The executable
\path{scripts/verify_ym_postarchive_v50_full_link_bank.py}
constructs the sixteen aliases \((\pm i,\pm1,\pm1,\pm1)\), their
free transverse covariances, the exact blocked inverse, and the two
lift matrices by (52.2).  It verifies \(LY_f=I\), constructs the
existing exact V41 banks, and performs both \(64\)-edge contractions
using Gaussian-rational arithmetic.  The resulting four-by-four
matrices at derivative orders zero, one, and two are recorded in
\path{artifacts/ym_postarchive_v50_full_link_bank.json} and give
(54.1)--(54.2) by entrywise subtraction.  No floating conversion
is used.  The source SHA--256 is
\texttt{0fb502768f75fa04e0c663e7b210e256f9468956398d0d7f5213027d9432a02c};
the canonical payload digest is
\texttt{8262d25b7c95c09749d6a35dfbf17447c2be53a28874cc856b00364f8762ead6}.
This is a finite exact proof for the stated phase and bank, not an
identity asserted at every momentum.
\end{proof}

\begin{remark}[Precise scope of the obstruction]
At this phase the coarse longitudinal direction is \(e_1\).
The transverse restrictions of the two matrices agree:
\(\Pi_c(K_f^{(r)}-K_o^{(r)})\Pi_c=0\) for the displayed orders.
Thus this test does not distinguish the banks on a covariance
supported entirely on the coarse transverse space.  V48 explains
why a preceding tree map need not leave such support intact.
The lift derivatives are frozen in this test; (54.2) is not the
second derivative of a complete momentum-dependent bubble or
response.  Cancellation between diagrams, or after integration,
is still possible and has not been tested exactly here.
\end{remark}

\subsection{Required companion snapshots}

The following current files were inspected read-only:
\begin{itemize}
\item SM twentieth-cycle V50, 401869 bytes, SHA--256
\texttt{49919AFC25EBAA98A48C6BC6EF2A5F8BE777DD8DCFD1E2F4467D54FF96771B5D}.
Its latest block constructs a renormalized one-parent resonant model,
its spectral measure, and fixed-coupling threshold tails.  Its own
review explicitly does not identify that model with the full
interacting SM or establish a kinetic--Einstein continuum.
\item NS fourth-cycle V26, 278283 bytes, SHA--256
\texttt{82C887F8C2C69E4F28FAD7C3DC8DE5B9099CB5A4BF431EBA1FFF3E91935978AD}.
Its endpoint remains the rank-one Oseen-kernel norm calculation;
global regularity is not proved.
\item Parallel YM V5, 54174 bytes, SHA--256
\texttt{306F1BB84CFA8A8D47EA569EC17E9545F9867C8D32B548EC9D9C444B4F3C0512}.
Its extensive-charge tail obstruction is unchanged.  It directs
the programme toward local or density estimates and retention of
compact labels, not a volume-uniform tail of total integer charge.
\end{itemize}
None supplies the background Ward identity needed to cancel
(54.1), nor a construction of the interacting YM measure.
SM's operator construction is a useful reminder to specify an
actual operator and its domain before interpreting a spectral
coefficient, but its result is not imported into the YM problem.

\subsection{Review of the preceding ten-version block}

V40 audited the Gaussian-tower direction.  V41 fixed the background
and introduced tree representatives, but incorrectly called the
representative map Laurent.  V42 corrected that assertion and
proved its punctured-torus derivative bounds.  V43 controlled the
conditioned level-zero covariance by its graph, and V44 used graph
precision to justify local second-order response differentiation.
V45 extended the regularity argument to each finite Gaussian level.
V46 corrected the anisotropic summability argument and obtained
uniform Gaussian and one-step symbol estimates.  V47 controlled the
composed transverse lift in a square-summable alias norm, while
retaining the cell-volume and momentum-rescaling factors.
V48 then found that the implemented recursive representative was
not reproduced by its transverse-only harmonic lift.  V49 proved
full-link composition and compared an isolated floating implementation.

\begin{assessment}[What remains valid and what must not be inferred]
The Gaussian transverse estimates V43--V47 remain estimates for
their precisely defined objects; the full-link correction does not
retroactively identify them with the completed Wilson response.
V48's exact longitudinal-loss witness and (54.1) rule out silently
discarding the correction at the level of an individual bank.
The old higher-shell outputs remain diagnostics of their original
convention.  They are not certificates for the completed convention,
and neither convention has been identified here with a physical
renormalized coupling.  No claimed mass-gap proof is hidden in the
completion of this Gaussian block.
\end{assessment}

\begin{decision}[Upstream priority after the audit]
Derive the finite-regulator background gauge identity for the exact
effective-action quantity that the programme intends to compute,
including the action, the chosen slice or quotient measure, and
its background dependence.  Determine the complete collection of
terms that transforms with the quartic and bubble contractions.
Then compare the two lift conventions against that identity in
exact arithmetic.  Merely accumulating more isolated shell values
does not answer this question.  Once the physical quantity is fixed,
the existing uniform Gaussian estimates can be applied to its
actual contracted symbols, followed by the still-missing nonlinear
and continuum estimates.
\end{decision}

\begin{firewall}[Position after V50]
The mandatory audit is complete and the representative sensitivity
of one Wilson bank is proved exactly at a specified phase.  The
complete gauge identity, physical response limit, interacting
continuum construction, reconstruction and Hamiltonian mass gap
remain open.  The full objective has not been replaced by a
Gaussian fixed-point or one-loop objective.
\end{firewall}

\begin{record}[Append and retention]
The predecessor V49 has SHA--256
\texttt{B3A68063B9366EF8D77D0A3618457A103200A4F7644E79218923F476EF537FB9}.
Its body is retained apart from the displayed active version.
Only V50 remains active.  Companions, archives and existing running
calculations are unchanged.  The next companion audit is V55,
the next ten-version review V60, and archival rollover remains V100.
\end{record}

% END APPENDED POST-TENTH-ARCHIVE V50 MATERIAL

% =============================================================================
\section{V51: background-dependent gauge tangents and frozen-bank curvature}
\label{sec:v51-moving-gauge}
% =============================================================================

The predecessor V50 has SHA--256
\texttt{67CEF5A0A2942A59BB27B255D2C926344A6157952DC0917C42A7764DF20A39D1}.
Its body is retained apart from the displayed active version.
This continuation goes upstream from the exact bank mismatch to
the hypotheses of the gauge identity, rather than inserting an
unspecified measure term.

\begin{record}[Reconciliation with already completed measure accounting]
Proposition~\ref{prop:v2-complete-endpoint} already states that the
stationary centre and logarithmic Haar-translation Jacobian have
vanishing background derivatives at the declared commuting-toron
Gaussian endpoints.  The ninth archive's compact V34 also records
the unit tree-gauge Faddeev--Popov determinant in its setting.
Those results are not rederived or withdrawn here.  V50's proposal
to identify the complete finite-regulator gauge identity must not
be read as authorization to add arbitrary Haar or stationary
terms to that already specified endpoint.  What remains to identify
for the new tower is the background-dependent representative and
its relation to that endpoint construction.
\end{record}

\subsection{The finite-dimensional gauge identity with its hypotheses}

\begin{theorem}[Off-shell and stationary gauge identities]
\label{thm:v51-gauge}
Work in a finite-dimensional real coordinate chart, including real
coordinates for complex matrix fields when needed.  Let \(S\) be
a sufficiently smooth gauge-invariant action, and let \(R(x)\xi\)
be the infinitesimal action of a fixed gauge parameter \(\xi\).
Write \(g=\nabla S\) and \(H=D^2S\).  Then
\[
 DS(x)[R(x)\xi]=0,\qquad
 H(x)R(x)\xi+(D_x(R(x)\xi))^*g(x)=0.             \tag{55.1}
\]
If \(b(t)\) is a stationary background family, \(g(b(t))=0\),
set \(H(t)=H(b(t))\), \(R(t)=R(b(t))\).  If these matrices are
twice differentiable, then, with subscripts denoting derivatives
at zero rather than Taylor coefficients,
\[
 H_0R_0=0,\quad H_1R_0+H_0R_1=0,\quad
 H_2R_0+2H_1R_1+H_0R_2=0.                        \tag{55.2}
\]
In particular
\[
                       R_0^*H_2R_0=2R_1^*H_0R_1. \tag{55.3}
\]
The last formula need not vanish, even though every \(H(t)\)
annihilates its actual gauge tangent \(R(t)\).
\end{theorem}
\begin{proof}
Differentiate \(S(g_s\cdot x)=S(x)\) in the gauge parameter \(s\)
at zero, obtaining the first identity of (55.1).  Differentiate
that scalar identity in an arbitrary coordinate direction \(u\):
\(D^2S[u,R\xi]+DS[D(R\xi)u]=0\).  Its vector form is the
second identity of (55.1).  At stationary backgrounds its second
term is zero, so \(H(t)R(t)=0\).  Differentiate this product zero,
one, and two times to get (55.2).  Multiply the last identity on
the left by \(R_0^*\).  Symmetry gives \(R_0^*H_0=0\) and,
from the first-derivative identity,
\(R_0^*H_1=-R_1^*H_0\).  Substitution proves (55.3).
For a nonstationary background the gradient term in (55.1) must
be retained and differentiated; (55.2) is then not justified.
\end{proof}

\begin{lemma}[Curvature of a moving lift]
\label{lem:v51-curvature}
Let \(U(t)\) be a twice differentiable lift matrix, and use the
derivative notation above.  Then
\[
\begin{split}
 (U^*HU)''(0)={}&U_0^*H_2U_0
 +2U_1^*H_1U_0+2U_0^*H_1U_1\\
 &+U_2^*H_0U_0+U_0^*H_0U_2+2U_1^*H_0U_1.
\end{split}                                                   \tag{55.4}
\]
For \(U=R\) along a stationary family, the full expression vanishes.
The single frozen-lift term \(R_0^*H_2R_0\) need not vanish.
\end{lemma}
\begin{proof}
Apply the product rule twice.  For \(U=R\),
\(R(t)^*H(t)R(t)=0\) identically by (55.2)'s undifferentiated
identity.  Thus its second derivative is zero, whereas (55.3)
gives the possible nonzero frozen contribution.
\end{proof}

\begin{proposition}[An exact elementary counterexample to the frozen inference]
Consider
\[
 S_t(x,y)=\tfrac12(x-ty)^2,\quad
 H(t)=\begin{pmatrix}1&-t\\-t&t^2\end{pmatrix},\quad
 R(t)=\begin{pmatrix}t\\1\end{pmatrix}.
\]
The translation \((x,y)\mapsto(x+t\epsilon,y+\epsilon)\) leaves
\(S_t\) invariant and preserves Lebesgue measure.  Nevertheless
\[
 R(0)^*H(t)R(0)=t^2,\qquad R(t)^*H(t)R(t)=0.       \tag{55.5}
\]
Thus the frozen longitudinal curvature equals two although the
complete gauge-direction curvature is zero, with no Jacobian term.
\end{proposition}
\begin{proof}
The expression \(x-ty\) is unchanged by the translation and its
Jacobian is one.  Matrix multiplication gives both identities.
Here \(R_1=(1,0)^\mathsf T\), so (55.3) gives
\(2R_1^*H_0R_1=2\), agreeing with (55.5).
\end{proof}

\subsection{Application boundary for the YM calculation}

V48's correction is a pure gradient at the zero-background Gaussian
level.  To compare complete background responses, one must establish
its continuation as a tangent to the background gauge orbit, or
otherwise account for the change of representative.  The assertion
that a free pure gradient remains a gauge tangent at every nonzero
background is not supplied by the zero-background Ward identity.
Consequently the nonzero quartic-bank matrix (54.1) is compatible
with gauge invariance of a complete moving-background construction.
It remains a correct obstruction to discarding the correction in
that individual frozen bank.

There are two distinct parameters to keep separate: \(t\) in
(55.2)--(55.4) differentiates the background family (for example its
amplitude), whereas the external momentum in the existing jet
scripts differentiates a Fourier label.  V50's order-two bank data
concern the latter, with the lifts frozen.  They are not an
evaluation of all six terms of (55.4).  A mixed amplitude/momentum
calculation must identify both parameter dependences before using
these formulas.

\begin{decision}[Concrete next construction]
Construct the infinitesimal gauge generator in the actual lattice
link coordinates around the declared toron background, and compute
its first two amplitude derivatives.  Compare the full-link
correction with that generator.  At each momentum used, explicitly
check whether the background is stationary for the differentiated
action; retain the off-shell term of (55.1) if it is not.  Assemble
the moving-lift terms only after this identification.  Do not add
a measure term already proved zero, and do not assume that the
fixed-background bank is the Hessian of this newly identified
moving representative without checking the coordinate map.
\end{decision}

\begin{firewall}[Position after V51]
The finite-dimensional differentiated gauge identity and a precise
counterexample to the frozen-gauge inference are proved.  The
actual YM background generator and its match to the tower remain
to be constructed.  No cancellation of the complete Wilson
response, interacting continuum theory, or mass gap is claimed.
Only V51 remains active; prior bodies and archives are retained.
The next companion audit remains V55.
\end{firewall}

% END APPENDED POST-TENTH-ARCHIVE V51 MATERIAL

% =============================================================================
\section{V52: the gauge tangent in the implemented split-link coordinates}
\label{sec:v52-generator}
% =============================================================================

The predecessor V51 has SHA--256
\texttt{3C7C545DF32A282DCAB758E0ECEC08F3547521E03D163762BCA77204442196F6}.
Its body is retained apart from the displayed active version.
The present acquisition identifies the link-coordinate convention
before applying the abstract gauge identities of V51.

\begin{record}[Active bank convention]
The cubic assembler called by V41 uses
\path{scripts/verify_ym_v31_cubic_bubble_jets.py}, whose translated
factor expression orders background colours before vertical
colours on forward links and reverses those groups on inverse
links.  The quartic routine
\path{scripts/verify_ym_v48_uncompressed_trace.py}, function
\texttt{factor\_expr}, uses the same ordering and separate factors
\(1/(n_b!n_a!)\).  They are the derivatives of
\[
 U_e=B_e\exp(a_e),\qquad
 U_e^{-1}=\exp(-a_e)B_e^{-1},\qquad B_e=\exp(tb_e), \tag{56.1}
\]
with additive coordinates inside each of the two exponentials.
They are not derivatives of a single \(\exp(tb_e+a_e)\).
The older unsplit word routine alone would not have identified
the convention of these active translated banks.
\end{record}

\begin{theorem}[Split-link generator and its amplitude jets]
\label{thm:v52-generator}
Let \(e\) be an oriented edge with source \(s\) and target \(v\).
Under \(U_e\mapsto g_sU_eg_v^{-1}\), the tangent in the coordinates
(56.1) at \(a=0\), holding \(B\) fixed while taking the gauge
variation, is
\[
 (R_B\omega)_e=\operatorname{Ad}_{B_e^{-1}}\omega_s-\omega_v.
                                                               \tag{56.2}
\]
For \(B_e(t)=\exp(tb_e)\), its first two amplitude derivatives are
\[
 (R_0\omega)_e=\omega_s-\omega_v,\quad
 (R_1\omega)_e=-[b_e,\omega_s],\quad
 (R_2\omega)_e=[b_e,[b_e,\omega_s]].               \tag{56.3}
\]
The sign of \(R_0\) is opposite to the forward difference used in
some earlier alias formulas; that is a choice of gauge-parameter
sign, not a different gradient range.
\end{theorem}
\begin{proof}
Take \(g_x(\epsilon)=\exp(\epsilon\omega_x)\).
At \(a=0\), the first-order change of the link is
\(\delta U_e=\omega_sB_e-B_e\omega_v\).
In (56.1), \(\delta U_e=B_e\delta a_e\), so multiplication by
\(B_e^{-1}\) proves (56.2).  Since
\(\operatorname{Ad}_{\exp(-tb)}=\exp(-t\operatorname{ad}_b)\),
differentiation proves (56.3), with \(R_2\) a second derivative
rather than the quadratic Taylor coefficient.
\end{proof}

\begin{remark}[Changing the fluctuation chart changes these jets]
For the alternative chart \(U_e=\exp(a_e)B_e\), the tangent is
\(\omega_s-\operatorname{Ad}_{B_e}\omega_v\), whose first two
derivatives are \(-[b_e,\omega_v]\) and
\(-[b_e,[b_e,\omega_v]]\).  One cannot combine these latter
formulas with the split ordering (56.1).  This follows directly
by multiplying \(\delta U_e\) on the other side.
\end{remark}

\begin{theorem}[When the stationary Ward specialization applies]
\label{thm:v52-stationary}
On a finite lattice with Wilson action
\(S(U)=\sum_p(d-\Re\Tr U_p)\), every flat configuration
\(U_p=I\) is stationary.  If all \(b_e\) belong to a common
commuting Lie subalgebra and their oriented sum around every
plaquette is zero, \(B_e(t)=\exp(tb_e)\) is flat for every \(t\).
For such a family, the Hessian in the coordinates (56.1) obeys
\(H(t)R(t)=0\) and all three identities (55.2), with (56.3).
\end{theorem}
\begin{proof}
For unitary matrices, \(\Re\Tr U_p\le d\), so a flat configuration
is a global minimum of the smooth finite-dimensional action and
its first derivative is zero.  Commutativity allows the product
around a plaquette to be written as the exponential of \(t\)
times the oriented sum; the assumed sum is zero.  Gauge invariance
of the Wilson action follows because each plaquette holonomy is
conjugated at its base point.  Apply (55.1) to the action in the
chart (56.1): stationarity kills the gradient term and leaves
\(H(t)R(t)=0\).  Twice differentiating proves the remaining
identities.  This uses the full link Hessian, not an arbitrary
restricted vertex matrix after a fixed projection.
\end{proof}

\begin{proposition}[A nonconstant transverse profile is not a stationary family]
\label{prop:v52-profile}
On a periodic finite lattice choose a nonzero anti-Hermitian
generator \(T\), directions \(1\ne2\), and a real nonconstant
profile \(\phi(x_1)\).  Set
\(b_{(x,2)}=\phi(x_1)T\), and all other directional components to
zero.  If \(\phi(x_1+1)-\phi(x_1)\) is nonzero somewhere, then
\(B(t)=\exp(tb)\) cannot be a stationary Wilson background for
all sufficiently small \(t\).  In contrast a constant profile
gives a flat commuting toron family.
\end{proposition}
\begin{proof}
On a \((1,2)\)-plaquette the oriented sum equals
\((\phi(x_1+1)-\phi(x_1))T\), up to an irrelevant overall
orientation sign.  Taylor expansion of its exponential gives
\[
 \left.\frac{d^2}{dt^2}S(B(t))\right|_{t=0}
 =\sum_{(1,2)\text{-plaquettes}}
   (\phi(x_1+1)-\phi(x_1))^2\,\Tr(T^*T)>0.          \tag{56.4}
\]
Other plaquette contributions vanish for this profile.  If every
\(B(t)\) in a neighbourhood were stationary, the derivative of
\(S(B(t))\) along the family would be identically zero, contradicting
(56.4).  For a constant profile every oriented sum vanishes and
the preceding theorem applies.
\end{proof}

\begin{assessment}[Consequence for momentum differentiation]
Real sine and cosine modes give examples covered by (56.4).
Their complex Fourier versions may be used as formal multilinear
coefficients but are not by themselves real unitary background
paths.  The constant commuting endpoint can legitimately use the
stationary identities, while a calculation that varies an external
transverse momentum must check its stationarity anew.  It cannot
infer \(H R=0\) at every nonzero momentum merely from flatness of
the zero-momentum toron.  The appropriate general identity is
\(HR+(DR)^*\nabla S=0\), with the gradient term retained before
the momentum derivatives are taken.  On a fixed finite periodic
lattice the allowed Fourier momenta are discrete; any continuous
momentum jet additionally uses the explicitly declared Bloch or
symbol extension, not a family of arbitrary periodic modes.
\end{assessment}

\begin{decision}[Next finite test]
Build the full, colour-resolved link Hessian jets in the actual
split convention (56.1), and verify (55.2) using (56.3) on a flat
commuting background.  Then derive the off-shell source terms for
the declared external-momentum symbol before projecting onto a
vertical slice or compressing colour indices.  The scalar divided
colour factors in the existing banks are not a substitute for this
matrix identity.  Only after that check should the new full-link
representative be compared with the complete contracted response.
\end{decision}

\begin{firewall}[Position after V52]
The infinitesimal gauge generator is now matched to the active
split-link bank convention, with explicit amplitude jets and a
proved stationarity distinction.  The colour-resolved Hessian
verification, full response cancellation, interacting continuum,
and Hamiltonian mass gap remain open.  Only V52 remains active;
previous bodies and archives are preserved.  The next companion
audit remains V55.
\end{firewall}

% END APPENDED POST-TENTH-ARCHIVE V52 MATERIAL

% =============================================================================
\section{V53: the full flat-background Wilson Hessian and exact Ward jets}
\label{sec:v53-flat-hessian}
% =============================================================================

The predecessor V52 has SHA--256
\texttt{37A55909854CE5EBE88E1BB0DF0E41C5BD2D4CBCF146926F1B1FF2DB17E8330B}.
Its body is retained apart from the active-version label.
Here the full link Hessian is computed before a vertical slice or
colour compression is chosen.  This avoids assuming that a scalar
compressed bank already satisfies the full gauge identity.

\begin{definition}[Transported plaquette differential]
For an oriented plaquette word let its background factors be
\(V_i=B_{e_i}\) on forward links and \(V_i=B_{e_i}^{-1}\) on
backward links.  Write \(P_0=I\), \(P_i=V_1\cdots V_i\).
In the split coordinates \(U_e=B_e\exp(a_e)\), define
\[
 (C_Ba)_p=
 \sum_{i:\,+}\operatorname{Ad}_{P_i}a_{e_i}
 -\sum_{i:\,-}\operatorname{Ad}_{P_{i-1}}a_{e_i}.       \tag{57.1}
\]
The edge Lie algebra uses its real Hilbert--Schmidt inner product;
the plaquette space uses the sum of those inner products.
\end{definition}

\begin{theorem}[Full Wilson Hessian at a flat background]
\label{thm:v53-gram}
For \(S(U)=\sum_p(d-\Re\Tr U_p)\) and any flat background,
\[
                         H_B=C_B^*C_B.                \tag{57.2}
\]
For arbitrary backgrounds the plaquette differential satisfies
\[
 (C_BR_B\omega)_p=(I-\operatorname{Ad}_{B_p})\omega_{x_p},
                                                               \tag{57.3}
\]
where \(x_p\) is the plaquette base point.  Hence at flat backgrounds
\(C_BR_B=0\) and \(H_BR_B=0\), with the generator of (56.2).
\end{theorem}
\begin{proof}
For a forward factor, \(\delta V_i V_i^{-1}
=\operatorname{Ad}_{B_{e_i}}a_{e_i}\); for a backward factor it is
\(-a_{e_i}\).  Differentiating the product and right-translating
its tangent by \(P_m^{-1}\) gives (57.1).
Under a gauge variation, the plaquette holonomy changes by
\(\delta B_p=\omega_{x_p}B_p-B_p\omega_{x_p}\).
Right translation proves (57.3).

At a flat background, along any real variation \(a\), write
\(U_p(s)=I+sX+s^2Y/2+o(s^2)\), with \(X=(C_Ba)_p\).
Unitarity gives \(X^*=-X\) and
\(Y+Y^*=-2X^*X\).  Thus
\(-\Re\Tr Y=\Tr(X^*X)\), so the second variation of the Wilson
action is \(\sum_p\|(C_Ba)_p\|^2\).  Polarization proves (57.2).
For flat \(B_p=I\), (57.3) vanishes, which proves the Hessian
identity without assuming any particular gauge slice.
\end{proof}

\begin{corollary}[Amplitude-jet construction]
Along a smooth flat background family, with subscripts denoting
derivatives at zero,
\[
\begin{split}
 H_0&=C_0^*C_0,\\
 H_1&=C_1^*C_0+C_0^*C_1,\\
 H_2&=C_2^*C_0+2C_1^*C_1+C_0^*C_2.
\end{split}                                                     \tag{57.4}
\]
The curl jets obey
\(C_0R_0=0\), \(C_1R_0+C_0R_1=0\), and
\(C_2R_0+2C_1R_1+C_0R_2=0\).  Consequently the Hessian jets
obey all identities (55.2).
\end{corollary}
\begin{proof}
Differentiate (57.2) and the flat specialization of (57.3).
Alternatively substitute these curl identities into (57.4).
\end{proof}

\subsection{An exact colour-resolved lattice test}

Use one square with vertices \(0,1,2,3\) in cyclic order and
forward edges \((0,1),(1,2),(3,2),(0,3)\); the latter two are
traversed backwards in the plaquette word.  Take commuting
backgrounds \(\exp(tb),\exp(2tb),\exp(tb),\exp(2tb)\),
with \(b=i\sigma_3\in\mathfrak{su}(2)\).
In the orthonormal colour basis \(i\sigma_a/\sqrt2\), its adjoint
generator is the integer matrix
\[
 J=\begin{pmatrix}0&2&0\\-2&0&0\\0&0&0\end{pmatrix}.
\]
The four curl blocks in order are
\(e^{tJ},e^{3tJ},-e^{3tJ},-e^{2tJ}\).
The gauge-generator block on an edge with source \(s\), target
\(v\), and background rate \(r\) is
\(e^{-rtJ}\omega_s-\omega_v\).  These give complete
three-colour curl matrices of size \(3\times12\), gauge matrices
of size \(12\times12\), and Hessians of size \(12\times12\).

\begin{record}[Exact integer verification]
\path{scripts/verify_ym_v53_flat_plaquette_ward.py}
forms the zero, first and second derivatives using only integer
matrix operations.  All six curl/Hessian Ward tests pass.
It also verifies
\[
 R_0^*H_2R_0=2R_1^*H_0R_1\ne0,\qquad
                    \Tr(R_0^*H_2R_0)=96.             \tag{57.5}
\]
The source SHA--256 is
\texttt{e7ec49920a4f9291547b191061ff130a5c508070df40f43edcda883ffcb9cc60};
the canonical payload digest in
\path{artifacts/ym_v53_flat_plaquette_ward.json} is
\texttt{377615a096c7eaea46014ea6ccab1d8cc82da6e1694b0bf64092b31701d3f673}.
This is a finite, full-colour check in the stated SU(2) normalization.
Its value 96 is not to be compared numerically with V50's divided
colour-bank entries, which use different matrices and conventions.
\end{record}

\begin{assessment}[Acquisition and remaining scope]
The full flat-background Wilson Hessian and its moving-gauge
identity now have an explicit transported-curl realization.
The exact test demonstrates within lattice gauge theory, rather
than only in V51's elementary model, that frozen gauge-direction
curvature can be nonzero while the moving-direction identity holds.
It does not prove a cancellation in the programme's projected
one-loop response.  At a nonflat background (57.3) is nonzero,
and the Hessian is not generally just \(C_B^*C_B\); terms involving
the background plaquette holonomy and second link variations remain.
Those terms are necessary for the off-shell identity of V51.
\end{assessment}

\begin{decision}[Next upstream calculation]
Derive the nonflat plaquette Hessian in the same split coordinates
and retain the terms omitted by the flat Gram formula.  Use it
to compute the mixed amplitude/external-momentum Ward identity
before any fibre projection.  Then match that full-colour identity
to the scalar cubic and quartic bank contractions and the
background-dependent representative.  This is the needed bridge
to a complete response, not another extrapolation of the old grids.
\end{decision}

\begin{firewall}[Position after V53]
The flat finite-regulator Hessian identity is proved and independently
checked in a small full-colour example.  The off-shell projected
response, its tower limit, interacting regulator removal, and the
Yang--Mills mass gap remain open.  Only V53 remains active; previous
bodies and archives are preserved.  The next companion audit is V55.
\end{firewall}

% END APPENDED POST-TENTH-ARCHIVE V53 MATERIAL

% =============================================================================
\section{V54: the nonflat plaquette Hessian and its gauge-source term}
\label{sec:v54-nonflat}
% =============================================================================

The predecessor V53 has SHA--256
\texttt{33795455E7C5009E80B3EF0C9B560660739D4614349ABD34F94D9BAC1577A5DE}.
Its body is retained apart from the active-version label.  This
appendix supplies the terms that the flat Gram formula omits.

\begin{theorem}[Exact nonflat plaquette gradient and Hessian]
\label{thm:v54-hessian}
Fix arbitrary background links and the split coordinates
\(U_e=B_e\exp(a_e)\).  For each occurrence in a plaquette word let
\(X_i(a)\) denote its transported signed tangent from (57.1), and
put \(C(a)=\sum_iX_i(a)\).  Define the symmetric bilinear matrix
\[
\begin{split}
 T(a,d)={}&\tfrac12\sum_i\{X_i(a),X_i(d)\}\\
 &+\sum_{i<j}\bigl(X_i(a)X_j(d)+X_i(d)X_j(a)\bigr),
\end{split}                                                     \tag{58.1}
\]
where \(\{A,D\}=AD+DA\).  For the plaquette action
\(S_p=d-\Re\Tr U_p\), at the zero fluctuation,
\[
 DS_p[a]=-\Re\Tr(C(a)B_p),\qquad
 D^2S_p[a,d]=-\Re\Tr(T(a,d)B_p).                 \tag{58.2}
\]
Equivalently
\[
 T(a,d)=\tfrac12\{C(a),C(d)\}
 +\tfrac12\sum_{i<j}\bigl([X_i(a),X_j(d)]
                         +[X_i(d),X_j(a)]\bigr).   \tag{58.3}
\]
At \(B_p=I\) the commutator traces vanish and (58.2) reduces to
the Gram formula.  For general \(B_p\), neither replacing it by
\(I\) nor dropping the commutator contribution is justified.
\end{theorem}
\begin{proof}
Move all fixed background factors in the plaquette product to
the right.  For forward factors use
\(B_e e^{a_e}=e^{\operatorname{Ad}_{B_e}a_e}B_e\), and for inverse
factors use \((B_e e^{a_e})^{-1}=e^{-a_e}B_e^{-1}\).
This gives the exact identity
\(U_p(a)=\prod_i e^{X_i(a)}B_p\), in the original word order.
Its first derivative is \(C(a)B_p\).  Its mixed second derivative
has the within-factor anticommutator and the two ordered
cross-factor products in (58.1).  Taking the real trace proves
(58.2).  Expanding (58.3) verifies it term by term.  At flat
background the trace of each commutator is zero, and
\(-\Re\Tr(C(a)C(d))\) is the real Hilbert--Schmidt pairing of
the anti-Hermitian tangents, as required.
\end{proof}

\begin{theorem}[Explicit off-shell gauge source in split coordinates]
\label{thm:v54-source}
Write \(\alpha_e=\operatorname{Ad}_{B_e^{-1}}\omega_{s(e)}\),
\(\beta_e=\omega_{v(e)}\).  The gauge generator at zero is
\((R_B\omega)_e=\alpha_e-\beta_e\), and its derivative in a
fluctuation direction \(a\) is
\[
 (D_aR_B[a]\omega)_e=-\tfrac12[a_e,\alpha_e+\beta_e]. \tag{58.4}
\]
For the full finite Wilson action, with its background held fixed,
\[
 D^2S[a,R_B\omega]+DS[D_aR_B[a]\omega]=0.          \tag{58.5}
\]
The derivatives here are fluctuation-coordinate derivatives, not
the background-amplitude derivatives \(R_1,R_2\) of V52.
\end{theorem}
\begin{proof}
At nonzero fluctuation coordinate \(u_e\), the derivative of the
exponential, right-translated by its inverse, is
\(J(u_e)\delta u_e\), where
\(J(u)=I+\tfrac12\operatorname{ad}_u+O(\|u\|^2)\).
The gauge variation in that frame is
\(\alpha_e-e^{\operatorname{ad}_{u_e}}\beta_e\).
Therefore
\[
 R(u)\omega=\alpha-\beta
             -\tfrac12[u,\alpha+\beta]+O(\|u\|^2),
\]
proving (58.4).  Differentiate the exact gauge-invariance identity
\(DS(u)[R(u)\omega]=0\) in direction \(a\) at zero to obtain
(58.5).  The plaquette formulas (58.2) can be substituted in both
terms without a stationarity assumption.
\end{proof}

\subsection{Exact nonflat full-colour verification}

On the square and edge orientations of V53 take the first
background link to be
\[
 B=\operatorname{diag}((3+4i)/5,(3-4i)/5,1),
\]
and the other three links to be the identity.  This is an SU(2)
subgroup embedded in SU(3), with nonidentity plaquette holonomy.
Use the three anti-Hermitian Pauli directions in its upper
two-by-two block as both link and site colour bases.

\begin{record}[Independent exact test of the source term]
\path{scripts/verify_ym_v54_nonflat_ward.py}
implements (58.1)--(58.2) as ordered matrix products and (58.4)
as a commutator.  All 144 edge-direction/site-gauge pairings
satisfy (58.5) in Gaussian-rational arithmetic; all twelve site
gauge basis vectors also satisfy (57.3).  Eighteen of the Hessian
terms are nonzero, with maximum absolute source value \(8/5\),
and cancel their source terms exactly.  Thus the test is not
merely checking two separately vanishing expressions.

The source SHA--256 is
\texttt{2c992477438bde5ccae333ff129b1130c27d73db770c9232e48aeefe138bfa76};
the canonical payload digest in
\path{artifacts/ym_v54_nonflat_ward.json} is
\texttt{b4e9dbaca31eadf15cf121a9399761c178722dc1d4c5bc1a65ccc9a53c635e7b}.
These finite tests corroborate the general proofs, but do not
identify their normalization with the divided-colour response banks.
\end{record}

\begin{assessment}[New control and the remaining response problem]
The full nonflat Wilson Hessian and its off-shell gauge identity
are now explicit.  In particular, flat-background Gram positivity
and the identity \(HR=0\) must not be continued into a nonflat
external-momentum calculation by assumption.  The background
holonomy terms and the source in (58.5) are part of that calculation.
This is a finite-regulator identity, not a cancellation theorem
for a projected Gaussian determinant.  It still has to be carried
through the actual background-dependent fibre map, covariance,
and colour contractions.
\end{assessment}

\begin{decision}[Next interface and scheduled audit]
Derive the Hessian of the action pulled back by a moving fibre
embedding, including its second derivative paired with \(DS\).
Only at a stationary background can that second-embedding term
be discarded automatically.  Match that pullback formula to the
implemented harmonic and full-link maps before interpreting their
response difference.  V55 must also inspect the latest SM and NS
notes under the scheduled companion audit.
\end{decision}

\begin{firewall}[Position after V54]
The nonflat Hessian/source identities are proved and tested exactly.
The projected effective-action response, uniform nonlinear flow,
interacting continuum construction, and Yang--Mills mass gap
remain open.  Only V54 remains active; prior bodies and archives
are retained.  The next companion audit is V55.
\end{firewall}

% END APPENDED POST-TENTH-ARCHIVE V54 MATERIAL

% =============================================================================
\section{V55: nonlinear fibre curvature, gauge completion, and companion audit}
\label{sec:v55-fibre}
% =============================================================================

The predecessor V54 has SHA--256
\texttt{DEF87279AC31A0DE0249A26A2C93BACA43EA1DA913309FC8420CB97F75C4D6C2}.
Its body is retained apart from the displayed active version.
The previous turn established the nonflat source identity; this
continuation identifies where that source enters a nonlinear
fibre pullback.  It does not presume that the current linear
Gaussian block rule already defines a nonlinear compact block map.

\subsection{Required companion audit}

SM advanced during this read-only audit from V56 to twentieth-cycle
V57.  The inspected V57 snapshot has 464641 bytes and SHA--256
\texttt{70A769D3E77FEF83DEEF5DAF9F6D39EB7C6F2584E5588B750167DAF5AADA6DF0}.
Its endpoint compares a fixed two-excitation CAR sector with a
bosonic model, with a sufficient joint box/weak-coupling scaling.
It explicitly leaves a growing-excitation, finite-density limit
uncontrolled.  Its embedding argument is methodological context,
not an estimate for a nonlinear YM fibre.
NS remains fourth-cycle V26, 278283 bytes, SHA--256
\texttt{82C887F8C2C69E4F28FAD7C3DC8DE5B9099CB5A4BF431EBA1FFF3E91935978AD}.
Its pointwise rank-one norms do not supply a YM constraint map or
its curvature.  No companion theorem is imported into the
following finite-dimensional proof.  The next audit is V60.

\begin{theorem}[Hessian of a nonlinear pullback]
\label{thm:v55-pullback}
Let \(S\) be a real \(C^2\) function in a finite-dimensional real
coordinate chart, and let \(F\) be a \(C^2\) embedding with
\(F(0)=b\), \(DF(0)=Y\), \(D^2F(0)=Z\).  Then
\[
 D^2(S\circ F)(0)[u,v]
       =D^2S(b)[Yu,Yv]+DS(b)[Z(u,v)].             \tag{59.1}
\]
Full stationarity \(DS(b)=0\) removes the second term.  Merely
\(DS(b)Y=0\), stationarity in the embedded directions, does not.
\end{theorem}
\begin{proof}
Differentiate \(D(S\circ F)(y)[u]=DS(F(y))[DF(y)u]\) in direction
\(v\).  The two product-rule terms give (59.1).  The weaker
stationarity condition only controls the range of \(Y\), while
\(Z(u,v)\) need not belong to that range.
\end{proof}

\begin{theorem}[Constraint curvature at a fibre stationary point]
\label{thm:v55-constraint}
Let \(\mathcal B\) be a \(C^2\) constraint map, with
\(L=D\mathcal B(b)\) surjective, and suppose
\(\mathcal B(F(y))=c\) identically.  Put
\(B_2=D^2\mathcal B(b)\).  Then
\[
                 LY=0,\qquad LZ(u,v)=-B_2[Yu,Yv]. \tag{59.2}
\]
If \(DS(b)\) annihilates \(\ker L\), there is a unique multiplier
\(\lambda\) in the dual constraint space with \(DS(b)=\lambda L\).
In this case
\[
 D^2(S\circ F)(0)[u,v]
       =(D^2S(b)-\lambda B_2)[Yu,Yv].             \tag{59.3}
\]
Thus constrained stationarity in general gives the Hessian of a
Lagrangian, not just the restriction of \(D^2S\).
\end{theorem}
\begin{proof}
Differentiate the constraint once and twice to get (59.2).
Since \(L\) is surjective, any linear functional vanishing on
\(\ker L\) factors uniquely through \(L\); this gives \(\lambda\).
Substitute \(DS(b)[Z]=\lambda LZ=-\lambda B_2[Yu,Yv]\) into
(59.1).  No choice of the component of \(Z\) in \(\ker L\)
affects this formula.
\end{proof}

\begin{remark}[Linear right inverses do not specify this data]
For a local section satisfying \(\mathcal B(F(c+\eta))=c+\eta\),
the first equation is instead \(LY=I\), but its second derivative
still satisfies \(LZ=-B_2[Y,Y]\).  A chosen right inverse \(Y\)
therefore permits the second-jet choice
\(Z(u,v)=-YB_2[Yu,Yv]+Z_{\ker L}(u,v)\).
The free kernel-valued term and the as-yet-unspecified \(B_2\)
are not determined by \(LY=I\).  In particular V48's full-link
linear completion alone is not a nonlinear fibre construction.
This algebraic jet observation does not assert global existence
of a section or a continuum measure.
\end{remark}

\subsection{Gauge-equivalent embeddings and the missing second jet}

\begin{theorem}[Second-jet cancellation for a genuine gauge change]
\label{thm:v55-gauge}
Let \(R(x)\omega\) be the infinitesimal generator of a smooth
gauge action preserving \(S\).  Let \(F(0)=b\), and let
\(\Omega\) be a linear map from the parameter space to the gauge
Lie algebra.  Define the genuinely gauge-related embedding
\(F'(y)=\exp(\Omega y)\cdot F(y)\).  Write \(R=R(b)\),
\(D R[w]\omega=D_x(R(x)\omega)|_b[w]\).  Its jets satisfy
\[
 Y'u=Yu+R\Omega u,                                \tag{59.4}
\]
\[
\begin{split}
 Z'(u,v)-Z(u,v)={}&D R[Yu]\Omega v+D R[Yv]\Omega u\\
 &+\tfrac12D R[R\Omega u]\Omega v
   +\tfrac12D R[R\Omega v]\Omega u.
\end{split}                                      \tag{59.5}
\]
With these second jets, the two pullback Hessians in (59.1)
agree even at a nonstationary background.
\end{theorem}
\begin{proof}
For a fixed small gauge parameter \(\omega\), the group action
is the time-one flow of \(x\mapsto R(x)\omega\), whose expansion is
\(x+R(x)\omega+\tfrac12D R(x)[R(x)\omega]\omega+O(\|\omega\|^3)\).
Insert \(x=F(y)\), \(\omega=\Omega y\), and differentiate through
second order to obtain (59.4)--(59.5).
The difference of the first terms of the pullback Hessians is
\[
 H[Yu,R\Omega v]+H[Yv,R\Omega u]+H[R\Omega u,R\Omega v].
\]
Use the off-shell identity
\(H[w,R\omega]+DS[D R[w]\omega]=0\) on each cross term.
Use it twice, with weights one half, on the last term, exploiting
symmetry of \(H\).  The result is precisely minus
\(DS[Z'-Z]\) from (59.5).  This proves equality directly and
agrees with the exact invariance \(S\circ F'=S\circ F\).
If the gauge parameter has a nonzero second jet, an additional
term \(R\Omega_2(u,v)\) appears; its pairing with \(DS\) vanishes
by the undifferentiated gauge identity.
\end{proof}

\begin{assessment}[Application to the present discrepancy]
V54 verifies the identity that cancels the curvature terms in
(59.5) on a nonflat plaquette.  The present theorem explains why
changing only a linear representative need not preserve a
nonstationary restricted Hessian: a genuine nonlinear gauge
change changes its second jet as well.  It does not show that
V48's two linear lifts arise from the same nonlinear fibre or
that their difference equals \(R_B\Omega\) at nonzero background.
A gauge change also stays in a fixed coarse fibre only if it
preserves the chosen nonlinear coarse constraint.  That condition
has not been established for the Gaussian block-average tower.
The already proved Haar-flat endpoint accounting is unaffected;
the term here is curvature of an embedding, not an invented Haar
Jacobian contribution.
\end{assessment}

\begin{decision}[The actual upstream input]
Specify a nonlinear compact, gauge-equivariant blocking map whose
linearization is the V37 block-average rule, or prove that the
intended effective-action construction uses a different declared
constraint.  Compute its second jet \(B_2\), determine which gauge
transformations preserve its fibres, and construct a compatible
embedding to second order.  Then use (59.1)--(59.5) to compare the
full-link and transverse-only routes.  Neither the nonlinear map
nor its fibre measure is determined by the existing Gaussian
covariance recursion, so their compatibility cannot be inferred
from the one-step Ward identity alone.
\end{decision}

\begin{firewall}[Position after V55]
The required companion audit and the nonlinear pullback/gauge
curvature identities are complete at finite-dimensional scope.
The nonlinear compact blocking map and its match to the response
remain unresolved.  No interacting continuum or Yang--Mills mass
gap is proved.  Only V55 remains active; prior bodies and archives
are retained.  The next companion audit and ten-version review
are V60.
\end{firewall}

% END APPENDED POST-TENTH-ARCHIVE V55 MATERIAL

% =============================================================================
\section{V56: obstruction to a field-independent averaged gauge action}
\label{sec:v56-gauge-average}
% =============================================================================

The predecessor V55 has SHA--256
\texttt{839721DEB7902D841FD0182280C5B8804C56C8FB14EC535B6DEAE99EFDB29565}.
Its body is retained apart from the active-version label.
Before constructing a nonlinear block map, this appendix checks
whether the proposed transformation of its gauge parameters can
come from a field-independent group action.  The answer under
that specific hypothesis is negative for a nonabelian group.

\begin{theorem}[Scalar-weight averaging is not a nonabelian group differential]
\label{thm:v56-average}
Let \(G\) be a Lie group with nonabelian Lie algebra \(\mathfrak g\).
Suppose a smooth homomorphism \(\sigma:G^m\to G\) has differential
at the identity
\[
        d\sigma(X_1,\ldots,X_m)=\sum_i w_iX_i,
                         \qquad w_i\in\mathbb R.      \tag{60.1}
\]
Then every \(w_i\) is zero or one and at most one is nonzero.
If \(\sum_iw_i=1\), the differential selects exactly one site;
it cannot be a nontrivial average.  The conclusion also holds
for a smooth local homomorphism near the identity.
\end{theorem}
\begin{proof}
A homomorphism sends left-invariant vector fields and their
commutators to the corresponding fields in the target, so its
differential preserves Lie brackets.  Choose \(X,Y\in\mathfrak g\)
with \([X,Y]\ne0\).  Put both in the same \(i\)-th factor.
Bracket preservation gives \(w_i[X,Y]=w_i^2[X,Y]\), hence
\(w_i\in\{0,1\}\).  Put \(X\) in factor \(i\) and \(Y\) in
a distinct factor \(j\).  Their bracket in the product algebra
is zero, whereas the target bracket is \(w_iw_j[X,Y]\).
Thus \(w_iw_j=0\).  These local arguments use only the
differential and apply equally to a local homomorphism.
\end{proof}

\begin{corollary}[The sixteen-site average fails the bracket test]
For \(A(X)=\frac1{16}\sum_{i=1}^{16}X_i\),
\[
 A([X^{(i)},Y^{(i)}])-[AX^{(i)},AY^{(i)}]
                         =\frac{15}{256}[X,Y],       \tag{60.2}
\]
and for distinct sites the analogous defect is
\(-[X,Y]/256\).  Both can be nonzero.  Therefore the averaged
gauge parameter of V37 is not the differential of a smooth
field-independent homomorphism \(G^{16}\to G\).
\end{corollary}
\begin{proof}
Substitute \(w_i=1/16\) into the two bracket calculations.
For example \(X=i\sigma_1\), \(Y=i\sigma_2\) in
\(\mathfrak{su}(2)\) have \([X,Y]=-2i\sigma_3\ne0\), so these
are actual nonzero defects, not merely formal coefficients.
\end{proof}

\subsection{What is and is not ruled out}

V37's linear Ward identity remains correct:
the linear block of a fine gradient is the coarse gradient of
the sixteen-site averaged parameter.  At that Gaussian order
the gradient calculation does not test nonabelian brackets.
The new result rules out the additional requirement that this
parameter average be the differential of a field-independent
homomorphism between the fine and coarse gauge groups.
It does not rule out every nonlinear blocking construction,
nor does it invalidate the Gaussian covariance estimates.

In particular, the following possibilities are not excluded:
\begin{itemize}
\item A different linear blocking rule with a chosen coarse-root
gauge parameter, such as a rule built from paths with common
endpoints.  Its linearization and Gaussian tower must then be
recomputed or related to the present tower by a proved comparison.
\item A field-dependent transformation of coarse gauge parameters,
with a proved composition law and a compatible nonlinear block map.
\item A construction directly on gauge-invariant observables or
quotient variables, with the relation to the present fluctuation
coordinates established separately.
\end{itemize}
This list states possible directions, not completed constructions.

\begin{proposition}[Consistency condition for a field-dependent alternative]
Suppose a candidate block map \(\mathcal B\) transforms as
\(\mathcal B(g\cdot U)=\sigma_U(g)\cdot\mathcal B(U)\).
A sufficient composition condition for this transformation law is
\[
 \sigma_U(gh)=\sigma_{h\cdot U}(g)\sigma_U(h),
 \qquad \sigma_U(1)=1.                              \tag{60.3}
\]
If \(\sigma_U\) is independent of \(U\), this condition reduces
to the homomorphism law obstructed above.  Equality of actions
on a particular coarse field alone may hold only modulo its
stabilizer and must not be substituted for (60.3) without analysis.
\end{proposition}
\begin{proof}
Apply first \(h\), then \(g\), and use the stated transformation
law twice.  Condition (60.3) makes the result agree with applying
\(gh\) once.  With field independence it reads
\(\sigma(gh)=\sigma(g)\sigma(h)\).  The stabilizer qualification
follows because two group elements may act identically on one
field without being equal as group elements.
\end{proof}

\begin{assessment}[Change in the live dependency]
V55 requested a nonlinear compact block map matching the Gaussian
average.  That request now has a precise restriction: it cannot
use the literal averaged gauge parameter as the differential of
a field-independent coarse gauge homomorphism.  Merely choosing
a group-valued mean of links does not resolve this compatibility
condition.  A field-dependent construction would have to supply
both (60.3) and the second constraint jet required by (59.2).
Until then, the V37 tower is a controlled Gaussian model, not an
identified nonlinear renormalization transformation for the compact
Yang--Mills measure.
\end{assessment}

\begin{decision}[Next upstream choice to investigate]
Inspect the existing common-endpoint path/router construction and
determine whether a field-dependent dressing can recover the
averaged linear rule with a valid composition law.  If not, return
to a declared root-equivariant nonlinear block map and quantify
the change of Gaussian rule.  Do not assert a nonlinear completion
merely because the linear Ward identity holds, and do not infer
nonexistence of Yang--Mills theory from this specific obstruction.
\end{decision}

\begin{firewall}[Position after V56]
A precise nonabelian bracket obstruction has been proved for the
field-independent averaged gauge action.  The permissible nonlinear
blocking alternatives, physical response, interacting continuum,
and Hamiltonian mass gap remain open.  Only V56 remains active;
prior bodies and archives are retained.  The next audit is V60.
\end{firewall}

% END APPENDED POST-TENTH-ARCHIVE V56 MATERIAL

% =============================================================================
\section{V57: a local dressed block with the averaged linearization}
\label{sec:v57-dressed}
% =============================================================================

The predecessor V56 has SHA--256
\texttt{AD85B8FEAE66B3599078EAC26CDD444F0E2ABCC936A5888309C6A52E7637BC8F}.
Its body is retained apart from the active-version label.
The field-independent obstruction is respected: the construction
below uses a field-dependent coarse gauge transformation and is
defined only in an explicit local logarithm chart.

\subsection{Construction and exact composition law}

Fix root-to-site lattice paths inside each parity cell.  For a cell
with root \(r\), let \(T_x(U)\) be the path transport from \(r\)
to \(x\in r+\{0,1\}^4\).  For a coarse edge from \(r\) to
\(r'=r+2e_\mu\), set
\[
 V_x=U_{(x,\mu)}U_{(x+e_\mu,\mu)},\qquad
 P_x=T_xV_xT_{x+2e_\mu}^{-1}.                     \tag{61.1}
\]
The last transport is rooted at \(r'\).  Thus every \(P_x\)
has the same endpoints.  Choose one \(P_*\) from these paths.
Work near the trivial configuration, where all logarithms below
are uniquely defined in the Lie algebra.  With averages over the
sixteen cell sites, define
\[
 h_r=\exp\bigl(\overline{\log T_x}\bigr),\quad
 K_{r\mu}=P_*\exp\bigl(\overline{\log(P_*^{-1}P_x)}\bigr),\quad
 \mathcal B_{r\mu}=h_r^{-1}K_{r\mu}h_{r'}.        \tag{61.2}
\]
For a compact matrix Lie group, an invariant sufficiently small
exponential neighbourhood supplies the required logarithms.
There are only finitely many path products per cell, so the
trivial field has an open neighbourhood on which (61.2) is smooth.

\begin{theorem}[Local gauge compatibility with field dependence]
\label{thm:v57-cocycle}
On the common domain of the expressions, \(K\) is root-equivariant:
\(K_{r\mu}(g\cdot U)=g_rK_{r\mu}(U)g_{r'}^{-1}\).
The dressed block obeys
\[
 \mathcal B_{r\mu}(g\cdot U)
  =\sigma_{U,r}(g)\mathcal B_{r\mu}(U)\sigma_{U,r'}(g)^{-1},
 \quad \sigma_{U,r}(g)=h_r(g\cdot U)^{-1}g_rh_r(U). \tag{61.3}
\]
Furthermore
\(\sigma_U(gh)=\sigma_{h\cdot U}(g)\sigma_U(h)\)
whenever the three configurations involved lie in the chart.
\end{theorem}
\begin{proof}
Path transport transforms by its endpoint gauge factors, so
\(P_x\mapsto g_rP_xg_{r'}^{-1}\).  Therefore
\(P_*^{-1}P_x\) is conjugated by \(g_{r'}\).
The local logarithm commutes with conjugation on its invariant
domain, giving the stated transformation of \(K\).  Substitution
in (61.2) proves (61.3).  In the product
\(\sigma_{h\cdot U,r}(g)\sigma_{U,r}(h)\), the middle factors
\(h_r(h\cdot U)h_r(h\cdot U)^{-1}\) cancel, leaving
\(h_r(gh\cdot U)^{-1}g_rh_r^{\rm gauge}h_r(U)\),
where \(h_r^{\rm gauge}\) denotes the value at site \(r\) of the
gauge transformation \(h\), not the dressing.  This is exactly
\(\sigma_{U,r}(gh)\).
\end{proof}

\begin{theorem}[The linearization is precisely the V37 average]
\label{thm:v57-linear}
For \(U_e(t)=\exp(ta_e)\), let \(\tau_x\) be the oriented sum
of \(a\) along the root path to \(x\), let \(\eta_x\) be the
corresponding sum in the target cell, and let
\(v_x=a_{(x,\mu)}+a_{(x+e_\mu,\mu)}\).  Then
\[
 \left.\frac{d}{dt}\log\mathcal B_{r\mu}(U(t))\right|_0
                            =\overline{v_x}.       \tag{61.4}
\]
At \(U=I\), the infinitesimal transformation in (61.3) is
\(\frac1{16}\sum_{x\text{ in cell }r}\omega_x\).
\end{theorem}
\begin{proof}
The first logarithmic jet of \(P_x\) is
\(p_x=\tau_x+v_x-\eta_x\).  The first jet of the path mean
\(K\) is \(\overline p\), irrespective of \(P_*\).
The two dressings subtract \(\overline\tau\) and add
\(\overline\eta\), proving (61.4).  On a pure gauge transform
of the identity, \(T_x=g_rg_x^{-1}\).  For
\(g_x=\exp(\epsilon\omega_x)\), the first jet of \(h_r(g\cdot I)\)
is \(\omega_r-\overline\omega\).  Formula (61.3) then gives
\(\overline\omega\).  Field dependence is essential, so this
does not contradict Theorem~\ref{thm:v56-average}.
\end{proof}

\subsection{An explicit second constraint jet}

For a cell quantity write \(\delta f_x=f_x-\overline f\).
\begin{theorem}[Quadratic logarithmic jet]
\label{thm:v57-quadratic}
The second derivative of the logarithmic block along \(U_e(t)=e^{ta_e}\)
is
\[
 B_2[a,a]=\overline{\,[a_{(x,\mu)},a_{(x+e_\mu,\mu)}]
 +[\delta\tau_x,\delta v_x]-[\delta\tau_x,\delta\eta_x]
 -[\delta v_x,\delta\eta_x]\,}.                   \tag{61.5}
\]
Its mixed bilinear form is recovered by polarization.  The path
choice enters through \(\tau,\eta\); the reference path \(P_*\)
does not affect this second jet.
\end{theorem}
\begin{proof}
Write
\(\log T_x=t\tau_x+t^2\kappa_x/2+O(t^3)\), and use analogous
notation \(\eta_x,\kappa'_x\) at the target.  The two-link product
has logarithm \(tv_x+t^2d_x/2+O(t^3)\), with
\(d_x=[a_{(x,\mu)},a_{(x+e_\mu,\mu)}]\).
The Baker--Campbell--Hausdorff expansion through second order
gives \(\log P_x=tp_x+t^2q_x/2+O(t^3)\), where
\[
 q_x=\kappa_x+d_x-\kappa'_x
       +[\tau_x,v_x]-[\tau_x,\eta_x]-[v_x,\eta_x].
\]
In (61.2), expansion of \(\log(P_*^{-1}P_x)\) and then of
\(P_*\exp(\overline{\log(P_*^{-1}P_x)})\) cancels the two
reference commutators.  Thus
\(\log K=t\overline p+t^2\overline q/2+O(t^3)\).
Multiplying by the dressings gives the second jet
\[
 -\overline\kappa+\overline q+\overline{\kappa'}
 -[\overline\tau,\overline p]-[\overline\tau,\overline\eta]
 +[\overline p,\overline\eta].
\]
Insert \(\overline p=\overline\tau+\overline v-\overline\eta\).
The path second jets cancel; grouping each averaged commutator
minus the commutator of its averages gives (61.5).
\end{proof}

\begin{record}[Exact noncommuting jet check]
\path{scripts/verify_ym_v57_dressed_block_jet.py}
evaluates (61.2) with truncated matrix exponential/logarithm series
and sixteen noncommuting transport/link jets in an SU(2) subgroup.
It independently compares the resulting first and second jets
with (61.4)--(61.5); both checks pass in Gaussian-rational arithmetic.
The source digest is
\texttt{ca086dd621b8dcf47cda59096b4226e9c548b46e49c112d6369af445db4cc798};
the payload digest in \path{artifacts/ym_v57_dressed_block_jet.json} is
\texttt{4f135ab33e48943226ddb41e2e77c43b0ac09a475ce57d07865cd3471924584b}.
The test treats the transport jets as algebraic inputs; the general
proof supplies their dependence on actual lattice paths.
\end{record}

\begin{firewall}[Local compatibility is not a blocking-measure theorem]
This constructs a group-valued local nonlinear map with the desired
linearization and a field-dependent composition law.  It does not
extend the logarithm chart to all compact configurations, specify
global sector transitions, prove reflection positivity or hypercubic
symmetry, or identify the pushforward measure with a conventional
coarse gauge-invariant measure.  In particular field-dependent
\(\sigma_U(g)\) does not automatically imply invariance under an
arbitrary fixed coarse gauge transformation after integration.
The compatible nonlinear fibres and their measure still require
analysis; the Yang--Mills mass gap remains open.
\end{firewall}

\begin{decision}[Next use of the acquired second jet]
Insert (61.5) into the constrained-Hessian formula (59.3) and
construct a local fibre chart by a verified right-inverse argument.
Track its measure and the background multipliers.  Separately
determine whether the dressing can be removed at the observable
level or whether a global chart/sector construction is needed.
Do not promote the local cocycle into a field-independent gauge
homomorphism or a global compact blocking theorem.
\end{decision}

\begin{record}[Retention]
Only V57 remains active; prior bodies and frozen archives are
preserved.  The next companion audit and ten-version review are V60.
\end{record}

% END APPENDED POST-TENTH-ARCHIVE V57 MATERIAL

% =============================================================================
\section{V58: a uniform local fibre chart and its actual density}
\label{sec:v58-chart}
% =============================================================================

The predecessor V57 has SHA--256
\texttt{3E1292DBFFCF6A5261B195DE24B51FC7A571FF9BBB90A3CBAD0402A6334BD1CA}.
Its body is retained apart from the displayed active version.
Fix the paths and local logarithm conventions of V57 on a finite
periodic even lattice.  All vector spaces below are finite real
Lie-algebra coordinate spaces, not gauge quotients.

\begin{lemma}[A bounded real-space right inverse]
\label{lem:v58-rightinverse}
Let \(L\) be the derivative of the V57 logarithmic block at the
identity.  Define \(Rc\) on fine links by assigning the coarse
value \(c_{r\mu}\) to each of the eight direction-\(\mu\) boundary
links in cell \(r\), and zero to its internal links.  Then
\[
 LR=I,\qquad \|R\|_\infty=1,\qquad \|L\|_\infty\le2. \tag{62.1}
\]
Here the norm is the maximum of a fixed Lie-algebra norm over
links.  Every fine field decomposes uniquely as
\(x=u+Rv\), with \(u\in\ker L\), \(v=Lx\).
\end{lemma}
\begin{proof}
Each of the sixteen two-link paths averaged in (61.4) contains
one boundary link of the source cell and one internal link,
possibly in the neighbouring cell.  On \(Rc\) its sum is exactly
\(c_{r\mu}\).  Averaging proves \(LR=I\).  The remaining norm
bounds follow from the assignment and the triangle inequality
for the average of two links.  Put \(u=x-RLx\); uniqueness follows
by applying \(L\) to a proposed decomposition.
\end{proof}

\begin{theorem}[Volume-independent small-field chart radius]
\label{thm:v58-chart}
Write the local logarithmic block as
\(F(x)=\log\mathcal B(\exp x)=Lx+N(x)\).
There exist \(r>0\) and \(K<\infty\), independent of the number
of cells for the fixed path pattern, such that
\[
 \|DN(x)\|_\infty\le K\|x\|_\infty,\quad
 \|N(x)\|_\infty\le\tfrac K2\|x\|_\infty^2,
 \quad \|x\|_\infty\le r,\quad Kr\le\tfrac12.       \tag{62.2}
\]
For \(u\in\ker L\), \(\|u\|_\infty<r/4\), and
\(\|c\|_\infty<r/4\), there is a unique \(v\) in the closed
ball \(\|v\|_\infty\le r/2\) satisfying
\[
 F(u+Rv)=c.                                      \tag{62.3}
\]
It depends smoothly on \((u,c)\), and
\(\Psi(u,c)=u+Rv(u,c)\) is a one-to-one local fibre chart.
Furthermore \(\|D_cv\|_\infty\le8/5\).
\end{theorem}
\begin{proof}
Each block component in V57 is a fixed finite composition of
matrix multiplication, inversion, exponential and local logarithm,
depending on a bounded number of fine links.  On a sufficiently
small closed coordinate ball those operations have bounded first
and second derivatives.  The maximum-link norm bounds the finite
stencil sums by constants independent of the lattice volume.
Thus \(D^2F\) has a uniform bound \(K\).  Since \(N(0)=DN(0)=0\),
integration along a line proves the first two estimates in (62.2).
Decrease \(r\) to ensure its last condition and logarithm-chart
validity.

For fixed \(u,c\), apply the contraction theorem to
\(v\mapsto c-N(u+Rv)\).  In the stated ball,
\(\|u+Rv\|_\infty\le3r/4\); the image has norm at most
\(r/4+9Kr^2/32\le r/4+9r/64<r/2\).
The Lipschitz constant is at most \(3Kr/4\le3/8\).
Existence and uniqueness follow.  The derivative in \(v\) of
(62.3) is \(I+DN(\Psi)R\), with inverse norm at most
\((1-3/8)^{-1}=8/5\).  The finite-dimensional implicit-function
theorem gives smooth dependence and the derivative bound.
From \(x=\Psi(u,c)\), one recovers
\(u=(I-RL)x\) and \(c=F(x)\), proving injectivity and the local
coordinate assertion.  No radius bound for large-field charts
has been claimed.
\end{proof}

\begin{corollary}[Fibre curvature supplied by the explicit block]
At \(c=0\), the chart has
\[
 D_u\Psi(0,0)u=u,\qquad
 D_u^2\Psi(0,0)[u,w]=-R B_2[u,w],\quad u,w\in\ker L, \tag{62.4}
\]
where \(B_2\) is the explicit commutator form (61.5).
\end{corollary}
\begin{proof}
Differentiate (62.3) at zero, where \(DN=0\), once and twice.
The second derivative gives
\(D_u^2v=-D^2F(0)[u,w]=-B_2[u,w]\).
Multiplication by \(R\) proves (62.4).
\end{proof}

\begin{theorem}[Local Haar/action density in the fibre chart]
\label{thm:v58-density}
Use Euclidean Lebesgue measure on the fine and coarse Lie-algebra
spaces and the induced measure on \(\ker L\).  Let \(J_R>0\)
be the constant absolute determinant of the linear isomorphism
\((u,v)\mapsto u+Rv\).  In exponential coordinates write product
Haar measure as \(j(x)\,dx\), with its local smooth positive
density \(j\).  Then on the chart image
\[
 e^{-S(\exp x)}j(x)\,dx
 =\frac{J_R e^{-S(\exp\Psi(u,c))}j(\Psi(u,c))}
 {|\det(I+DN(\Psi(u,c))R)|}\,du\,dc.              \tag{62.5}
\]
The determinant is nonzero.  This formula is a local coordinate
identity for the raw finite measure, not a gauge-fixed quotient
measure or a global disintegration theorem.
\end{theorem}
\begin{proof}
The linear splitting contributes \(J_R\).  In the coordinates
\((u,v)\), the transformation to \((u,c)\), with
\(c=F(u+Rv)\), has block-triangular derivative and diagonal blocks
\(I\) and \(I+DN(x)R\).  The ordinary change-of-variables theorem
therefore gives (62.5).  The inverse bound in the preceding
theorem proves nonvanishing.  Haar measure has a smooth positive
density in any sufficiently small Lie-group coordinate chart;
taking products gives \(j\).  The identity applies, for example,
to test functions supported compactly inside this chart image,
without imposing an unmentioned boundary cutoff.
\end{proof}

\begin{assessment}[What the local construction does not control]
The radius in (62.2) is uniform in volume, but the chart event
requires every participating link to be small.  No lower bound
on the probability of that event, or estimate for its complement,
has been established.  The determinant and Haar product in
(62.5) have dimensions growing with volume; a uniform inverse
operator bound is not a volume-uniform bound on their densities
or logarithmic derivatives.  These terms must be retained when
constructing an actual blocked measure.  They arise from this
new nonlinear chart, not from the Haar translation whose
Jacobian was already proved one at the archived endpoint.
Residual gauge directions have not been removed, and positivity
of a physical vertical Hessian is not inferred from submersion.
\end{assessment}

\begin{decision}[Next measure-level test]
Compute the first nontrivial logarithmic density jets of (62.5),
using (61.5) and the next derivative of the nonlinear block where
required.  Check their transformation under the field-dependent
cocycle and identify the gauge quotient before taking a Gaussian
determinant.  A local fibre chart is now available; the missing
step is compatibility of its measure and gauge reduction with
the response whose coefficient the programme seeks.
\end{decision}

\begin{firewall}[Position after V58]
A local fibre chart with volume-independent small-field radius,
its second jet, and its exact finite-dimensional density are
proved.  Global chart coverage, large-field control, quotient
measure compatibility, the interacting continuum and the
Yang--Mills mass gap remain open.  Only V58 remains active;
prior bodies and archives are retained.  The next audit is V60.
\end{firewall}

% END APPENDED POST-TENTH-ARCHIVE V58 MATERIAL

% =============================================================================
\section{V59: the quadratic local density requires a third block jet}
\label{sec:v59-density-jets}
% =============================================================================

The predecessor V58 has SHA--256
\texttt{F8A224D8DF97D4DD0B9D83FD81A134B19579363000D280BE7C833B2A594E0820}.
Its body is retained apart from the active-version label.
The following derivatives concern the raw finite-dimensional
fibre-chart density of V58, before gauge reduction.

\begin{definition}[Density and derivative notation]
Set \(B_2=D^2F(0)\), \(B_3=D^3F(0)\), and define endomorphisms
of the coarse coordinate space by
\[
 A_a v=B_2[a,Rv],\qquad A_{ab}v=B_3[a,b,Rv].
\]
The logarithm of the measure factor, omitting the constant \(J_R\)
and the action, is
\[
 \ell(x)=\log j(x)-\log|\det(DF(x)R)|.             \tag{63.1}
\]
All traces on coarse coordinates below are real matrix traces.
\end{definition}

\begin{theorem}[First and second density derivatives]
\label{thm:v59-density}
For the V57 block and V58 right inverse,
\[
 D\ell(0)[a]=0,                                  \tag{63.2}
\]
\[
 D^2\ell(0)[a,b]=
 \frac1{12}\sum_e\Tr_{\mathfrak g}
       (\operatorname{ad}_{a_e}\operatorname{ad}_{b_e})
       -\Tr A_{ab}+\Tr(A_aA_b).                  \tag{63.3}
\]
The third block derivative in this formula is not determined by
the linear right inverse or the explicit second jet (61.5).
\end{theorem}
\begin{proof}
Near zero, \(DF(x)R=I+DN(x)R\) is invertible and its real
determinant is positive by continuity from the identity.
The elementary determinant identities give
\[
 D\log\det(DFR)(0)[a]=\Tr A_a,
\quad
 D^2\log\det(DFR)(0)[a,b]=\Tr A_{ab}-\Tr(A_aA_b).
\]
For completeness, in exponential coordinates the one-link Haar
density relative to its value at zero is the absolute determinant
of the translated exponential derivative
\((1-e^{-\operatorname{ad}_x})/\operatorname{ad}_x\).
Its power series at zero gives
\[
 \log j_e(x)-\log j_e(0)
 =-\tfrac12\Tr_{\mathfrak g}\operatorname{ad}_x
   +\tfrac1{24}\Tr_{\mathfrak g}(\operatorname{ad}_x)^2
   +O(\|x\|^3).
\]
In a compact Lie algebra with an invariant inner product,
\(\operatorname{ad}_x\) is skew-adjoint, so its trace is zero.
Polarization and multiplication of the link densities give the
Haar contribution in (63.3).

Finally (61.5), after polarization, is a finite sum of brackets
of scalar-weighted edge coordinates.  Both root-path sums and
\(R\) act identically on colour components.  Therefore each
coarse-coordinate block of \(A_a\) is a linear combination of
operators \(\operatorname{ad}_z\), with real scalar coefficients.
In particular every diagonal block has zero colour trace.
Thus \(\Tr A_a=0\) for every \(a\), which proves (63.2) and
completes (63.3).  This argument does not make \(A_a\) itself zero.
\end{proof}

\begin{corollary}[The density along the identity fibre]
For \(u,w\in\ker L\), let \(\Psi\) be the V58 chart at \(c=0\).
Then
\[
 D_u^2(\ell\circ\Psi)(0,0)[u,w]=D^2\ell(0)[u,w].  \tag{63.4}
\]
For the logarithm of the full Haar/action density, the second
derivative at the trivial Wilson field is
\[
 -H_0[u,w]+\frac1{12}\sum_e\Tr_{\mathfrak g}
       (\operatorname{ad}_{u_e}\operatorname{ad}_{w_e})
       -\Tr A_{uw}+\Tr(A_uA_w).                  \tag{63.5}
\]
\end{corollary}
\begin{proof}
The second-order chain rule would add
\(D\ell(0)[D_u^2\Psi(u,w)]\), but (63.2) makes that term zero.
The first derivative of \(\Psi\) on \(\ker L\) is the inclusion.
The trivial Wilson field is fully stationary, so its action
pullback contributes only \(-H_0[u,w]\).  The constant \(J_R\)
has zero derivatives.
\end{proof}

\begin{proposition}[Second constraint jets do not determine density curvature]
Even at a regular point with a fixed right inverse, two maps
having the same first and second derivatives can have different
quadratic chart-density terms.
\end{proposition}
\begin{proof}
On \(\mathbb R\) take \(F_\kappa(x)=x+\kappa x^3/6\), \(R=1\),
and constant reference density.  All maps have \(DF(0)=1\)
and \(D^2F(0)=0\), but
\[
 \left.\frac{d^2}{dx^2}[-\log F_\kappa'(x)]\right|_0=-\kappa.
\]
The derivative stays positive in a sufficiently small interval,
so these are genuine local coordinate changes.  This example
proves the logical need for third-jet data; it is not asserted
to be a Yang--Mills block map.
\end{proof}

\begin{assessment}[What was acquired, and what was not]
The first logarithmic measure derivative vanishes for the actual
V57 bracket structure.  The complete symbolic second derivative
is (63.3), with a required \(B_3\) term.  No value, sign, or
cancellation of that third-jet contraction has been computed here.
The traces sum over all coarse or fine coordinates and can be
extensive in volume; (63.3) is not a volume-uniform density bound.
It also contains no gauge-quotient determinant or removal of
residual gauge zero modes.  These are derivatives of the new
nonlinear chart, not a contradiction to the archived unit
Haar-translation Jacobian at a different specified endpoint.
\end{assessment}

\begin{decision}[Next concrete acquisition]
Expand the explicit dressed block (61.2) to third order, retaining
the chosen reference path, which cancelled only through second
order in V57.  Compute \(\Tr(v\mapsto B_3[u,w,Rv])\) together with
\(\Tr(A_uA_w)\), and compare their sum with the Haar term before
any sign or smallness estimate.  Then determine how the local
gauge reduction changes this density.  V60 also requires the
companion audit and ten-version review.
\end{decision}

\begin{firewall}[Position after V59]
The local density's first derivative and its full second-derivative
formula are proved.  The actual third-jet contraction, quotient
measure, nonlinear continuum construction and Yang--Mills mass
gap remain open.  Only V59 remains active; earlier bodies and
archives are retained.  The next audit is V60.
\end{firewall}

% END APPENDED POST-TENTH-ARCHIVE V59 MATERIAL

% =============================================================================
\section{V60: cubic reference dependence and the compulsory review}
\label{sec:v60-review}
% =============================================================================

The predecessor V59 has SHA--256
\texttt{9ADE8D0A12DD9301F72B0E347AB012D132E05775486069C5EFB334643961BA4B}.
Its body is retained apart from the active-version label.

\begin{theorem}[Cubic term of the reference-based path mean]
\label{thm:v60-cubic}
Let \(P_i(t)\) be smooth group-valued paths through the identity,
\(X_i(t)=\log P_i(t)=tp_i+O(t^2)\), and let
\(A(t)=\log P_*(t)=ta+O(t^2)\).  For fixed averaging weights
summing to one define
\(K_*(t)=P_*(t)\exp(\overline{\log(P_*^{-1}P_i)})\).
With \(\delta p_i=p_i-\overline p\),
\[
 \log K_*(t)=\overline{X_i(t)}
     +\frac{t^3}{12}\overline{[\delta p_i,[a,\delta p_i]]}
     +O(t^4).                                      \tag{64.1}
\]
Assume sufficient smoothness for the displayed remainder, as
holds for the analytic block of V57.  For two reference paths
with first jets \(a,a'\), the difference of their third derivatives
is
\[
 \frac12\overline{[\delta p_i,[a-a',\delta p_i]]}.  \tag{64.2}
\]
The common dressings in (61.2) leave this leading cubic difference
unchanged in the logarithm of the dressed block.
\end{theorem}
\begin{proof}
Use the third-order Baker--Campbell--Hausdorff formula and put
\(B=\overline{X_i}\).  To degree three,
\[
 \overline{\log(e^{-A}e^{X_i})}
 =B-A-\tfrac12[A,B]+\tfrac1{12}[A,[A,B]]
                      +\tfrac1{12}\overline{[X_i,[A,X_i]]}.
\]
Applying the same formula once more to the product with \(e^A\)
cancels every quadratic term and all cubic terms involving only
\([A,[A,B]]\).  The remaining correction is
\(\tfrac1{12}(\overline{[X_i,[A,X_i]]}-[B,[A,B]])\).
Since the average of \(X_i-B\) is zero, this equals
\(\tfrac1{12}\overline{[X_i-B,[A,X_i-B]]}\).
Only the first jets enter its cubic term, proving (64.1).
Subtract the two expressions and multiply their cubic coefficient
by \(3!=6\) to get (64.2).  Multiplication on either side by a
common path through the identity changes a degree-three difference
only at degree four or higher; the same holds after taking its
local logarithm.  This proves the dressing assertion.
\end{proof}

\begin{proposition}[Reference dependence can be nonzero]
Take three equally weighted paths with first jets
\(p_1=X=i\sigma_1\), \(p_2=Y=i\sigma_2\), \(p_3=0\), and choose
references with first jets \(a=X\) and \(a'=0\).  Then the
quantity averaged in (64.2), before the factor one half, is
\[
             \frac89X+\frac49Y\ne0.               \tag{64.3}
\]
\end{proposition}
\begin{proof}
The average equals
\(\frac13[Y,[X,Y]]-\frac19[X+Y,[X,Y]]\).
The Pauli commutators give
\([Y,[X,Y]]=4X\) and \([X,[X,Y]]=-4Y\), proving (64.3).
These are genuine local paths, for example \(e^{tX},e^{tY},I\).
The example concerns the path-mean operation; it is not a computed
YM density or a claim that its contracted trace is nonzero.
\end{proof}

\begin{assessment}[Effect on the density calculation]
The reference path in V57 was arbitrary but fixed.  Its choice
did not affect the first two block jets; it can affect the third.
Consequently a calculation of V59's \(\Tr B_3[a,b,R\,\cdot]\)
must retain that choice.  Formula (64.2) determines its possible
change after polarization, but this note does not claim that the
coarse trace of that change is nonzero.  Proving cancellation in
that trace, or showing equivalence of the resulting full measures,
is a separate task.  No canonical reference-independent cubic
block has been established merely by matching the Gaussian rule.
\end{assessment}

\subsection{Required companion audit}

The inspected SM snapshot is twentieth-cycle V62, 514562 bytes,
SHA--256
\texttt{4BBC3AC3649C8753E2664AF138822DF894ACEDBF5FF6FF9A93125FECFF175882}.
Its endpoint directs a fixed-time continuum of spin variables
after a large-population limit in finitely many groups, and
explicitly distinguishes this sequential procedure from an
unproved simultaneous quantum limit.  It leaves ultraviolet
removal and an Einstein source unresolved.  It supplies no
third derivative or quotient-density estimate for the YM block.
NS remains fourth-cycle V26, 278283 bytes, SHA--256
\texttt{82C887F8C2C69E4F28FAD7C3DC8DE5B9099CB5A4BF431EBA1FFF3E91935978AD}.
Its rank-one kernel norms do not resolve the same dependency.
No companion theorem is imported into the cubic calculation.

\subsection{Ten-version review: V50--V59}

V50 proved a frozen quartic-bank discrepancy and audited its
scope.  V51 separated that discrepancy from a complete moving
gauge identity.  V52 matched the active split-link coordinates
and distinguished flat from nonstationary profiles.  V53 and
V54 derived and exactly tested the flat and nonflat full-colour
Hessian identities, including the nonzero off-shell source.
V55 identified nonlinear fibre curvature and the second-jet
requirements for gauge-equivalent embeddings.  V56 ruled out
a field-independent nonabelian gauge average.  V57 supplied a
local dressed alternative with a field-dependent composition
law and the desired linearization.  V58 constructed its local
fibres and raw density.  V59 identified the third block derivative
required by that density's quadratic term.

\begin{decision}[Next acquisition after review]
Fix an explicit root-path pattern and reference convention, then
compute the contracted cubic block term in (63.3), not just the
uncontracted third derivative.  Use (64.2) to test whether its
reference dependence survives the trace.  Only after combining
the Haar and Jacobian terms should the local gauge quotient be
formed and compared with the original endpoint response.  The
local construction must still be extended or patched to cover
the configurations relevant to the compact measure; the global
and continuum requirements have not been removed by this review.
\end{decision}

\begin{firewall}[Position after V60]
The compulsory audit and review are complete, and cubic reference
dependence is explicitly proved for the local path mean.
The actual contracted density, quotient compatibility, global
blocking measure, interacting continuum and Yang--Mills mass gap
remain open.  Only V60 remains active; prior bodies and archives
are retained.  The next companion audit is V65, the next review
V70, and the archival rollover remains V100.
\end{firewall}

% END APPENDED POST-TENTH-ARCHIVE V60 MATERIAL

% =============================================================================
\section{V61: an exact boundary section and reference-independent density curvature}
\label{sec:v61-reference-trace}
% =============================================================================

The predecessor V60 has SHA--256
\texttt{DC5CDD5429450F46738F141D137BA873664E722A55FB97FBF25535AFD27F003C}.
Its body is retained apart from the displayed active version.
Keep the same internal root-to-site paths as V57 and use the
boundary-edge right inverse \(R\) of V58.  Reference choices in
this appendix are choices among the sixteen paths \(P_x\) for
each coarse edge; the internal root paths themselves are not changed.

\begin{theorem}[The linear boundary injection is an exact local section]
\label{thm:v61-section}
For the logarithmic nonlinear block
\(F(x)=\log\mathcal B(\exp x)\) of V57,
\[
                         F(Rc)=c                  \tag{65.1}
\]
for all sufficiently small coarse fields \(c\).
In particular \(DF(Rc)R=I\), and at zero
\[
 D^mF(0)[Rc_1,\ldots,Rc_m]=0\quad(m\ge2).          \tag{65.2}
\]
The statements hold for every allowed reference choice.
\end{theorem}
\begin{proof}
For \(x=Rc\), all internal fine links are the identity.
Every root-to-site transport inside a cell is therefore the
identity, as is its dressing \(h_r\).  Each of the sixteen
two-link paths in direction \(\mu\) consists of one internal
identity link and the boundary link \(\exp(c_{r\mu})\).
Thus every \(P_x\) for that coarse edge equals
\(\exp(c_{r\mu})\), even when the second internal link lies in
the neighbouring cell.  All relative logarithms in (61.2) are
zero, and the dressed block equals \(\exp(c_{r\mu})\).
Taking its local logarithm proves (65.1).  Differentiation in
the arbitrary coarse coordinates proves the rest.  This section
need not be a harmonic minimizer or a stationary section.
\end{proof}

\begin{theorem}[Reference variation vanishes after one right-inverse insertion]
\label{thm:v61-reference}
Let \(F_*,F_\dagger\) use different reference paths but the same
internal transports, and write
\(\Delta B_3=D^3F_*(0)-D^3F_\dagger(0)\).  Then
\[
                  \Delta B_3[u,w,Rc]=0             \tag{65.3}
\]
for arbitrary fine fields \(u,w\) and coarse field \(c\).
Consequently the quadratic logarithmic density in (63.3) is
independent of this reference choice, although the uncontracted
third block jet can differ.
\end{theorem}
\begin{proof}
For any one coarse edge, let \(p_x(a)\) be the linearized
transported path and \(\delta p_x=p_x-\overline p\).
On the input \(Rc\), its internal transport variations vanish
and its two-link variation is \(c_{r\mu}\), independently of
\(x\).  Hence \(\delta p_x(Rc)=0\).  Also the difference
\(d=p_*-p_\dagger\) between the reference variations satisfies
\(d(Rc)=0\).

By (64.2), the diagonal cubic difference is
\(\Delta B_3[a,a,a]=\tfrac12\overline{
[\delta p_x(a),[d(a),\delta p_x(a)]]}\).
Polarize this homogeneous cubic expression.  Every mixed term
with one input \(Rc\) puts it into either a \(\delta p_x\) slot
or a \(d\) slot, both of which vanish.  This proves (65.3)
componentwise, before a trace is taken.
The first two derivatives of the blocks already agree by V57.
Thus \(A_u,A_w\) in (63.3) agree, and the sole possible difference
\(\Tr(\Delta B_3[u,w,R\,\cdot])\) is zero.  The Haar contribution
is unchanged.  This proves the density assertion without
assuming a global equivalence of the two block measures.
\end{proof}

\begin{corollary}[A density benchmark on the exact section]
For the measure-only logarithm \(\ell\) of (63.1), normalized by
its value at zero,
\[
 \ell(Rc)-\ell(0)=8\sum_{r,\mu}
       \bigl(\log j_G(c_{r\mu})-\log j_G(0)\bigr). \tag{65.4}
\]
Here \(j_G\) is the one-link Haar density in exponential
coordinates.  In particular
\[
 D^2\ell(0)[Rc,Rd]=\frac23\sum_{r,\mu}
        \Tr_{\mathfrak g}(\operatorname{ad}_{c_{r\mu}}
                           \operatorname{ad}_{d_{r\mu}}).
                                                               \tag{65.5}
\]
\end{corollary}
\begin{proof}
By (65.1), \(DF(Rc)R=I\), so the constraint Jacobian determinant
in (63.1) is exactly one along this section.  Each coarse value
appears in exactly eight boundary fine links; all internal links
remain zero in Lie-algebra coordinates.  Product Haar density
therefore gives (65.4).  The one-link second derivative from V59
is \(\frac1{12}\Tr_{\mathfrak g}(\operatorname{ad}_c
\operatorname{ad}_d)\); multiplication by eight gives (65.5).
This is a section benchmark, not a formula for generic points
of a fibre or for the integrated coarse density.
\end{proof}

\begin{assessment}[Progress on the actual contraction]
V60's cubic reference dependence does not survive the
right-inverse insertion required for V59's quadratic density.
No reference-dependent sign ambiguity at that order remains
from this particular choice.  The common value of
\(\Tr B_3[u,w,R\,\cdot]\) for general fibre directions still
has to be computed; (65.3) proves equality, not that it vanishes.
Changing the internal root paths is outside the theorem and
can change the second jet itself.  The exact section supplies
a useful check of future implementations but covers a
lower-dimensional set, not the compact measure.
\end{assessment}

\begin{decision}[Next computation]
Choose the simplest reference among the now-equivalent options
for quadratic density, while fixing the internal path pattern.
Compute the common cubic contraction for directions in \(\ker L\),
combining it with \(\Tr(A_uA_w)\) and the Haar term.  Check the
result against (65.5) in section directions.  Only afterwards
perform the required local gauge reduction and assess whether
the resulting response matches either earlier lifting route.
\end{decision}

\begin{firewall}[Position after V61]
An exact local section, cancellation of cubic reference dependence
in the quadratic density, and a section-density benchmark are
proved.  The general fibre density, gauge quotient, global
measure construction, interacting continuum and Yang--Mills
mass gap remain open.  Only V61 remains active; prior bodies
and archives are retained.  The next companion audit is V65.
\end{firewall}

% END APPENDED POST-TENTH-ARCHIVE V61 MATERIAL

% =============================================================================
\section{V62: an explicit quadratic density on boundary-supported fibres}
\label{sec:v62-boundary}
% =============================================================================

The predecessor V61 has SHA--256
\texttt{287266DB63980F79E6FACCCB695D7EC09324B66190BC32B264F212E39D6A4248}.
Its body is retained apart from the active-version label.
This appendix computes an actual third-jet contraction, rather
than only its invariance under reference changes.

Restrict fine coordinate fields to boundary links, leaving every
internal link zero.  For each coarse edge \(e\), denote the eight
distinct boundary-link values by \(a_{e,s}\), \(1\le s\le8\),
their average by \(\bar a_e\), and their centred values by
\(\delta a_{e,s}=a_{e,s}-\bar a_e\).  The sixteen paths in V57
repeat each of these eight links twice.  Internal transports and
dressings are identically the identity on this subspace.

\begin{theorem}[Boundary restriction of the Jacobian jets]
\label{thm:v62-trace}
For boundary-supported \(u,w\), the coarse endomorphisms of V59
satisfy \(A_u=A_w=0\), and
\[
 \Tr\bigl(c\mapsto B_3[u,w,Rc]\bigr)
 =-\frac16\sum_e\frac18\sum_{s=1}^8
       \Tr_{\mathfrak g}(\operatorname{ad}_{\delta u_{e,s}}
                         \operatorname{ad}_{\delta w_{e,s}}).
                                                               \tag{66.1}
\]
\end{theorem}
\begin{proof}
On this subspace, each coarse block reduces exactly to the
reference-based mean of the eight matrices \(\exp(a_{e,s})\).
Distinct coarse edges use distinct boundary-link sets.
Formula (64.1) therefore gives
\[
 F_e(a)=\bar a_e+\tfrac1{12}\frac18\sum_s
         [\delta a_{e,s},[a_{e,*},\delta a_{e,s}]]+O(\|a\|^4).
\]
There is no quadratic term on this entire subspace, so
\(B_2[u,Rc]=0\) and \(A_u=0\).
Insert \(a_{e,s}=t u_{e,s}+z c_e\).  The centred field is
\(t\delta u_{e,s}\); only the middle reference slot contributes
to the coefficient of \(t^2z\).  Differentiation yields
\[
 B_3[u,u,Rc]_e=\frac16\frac18\sum_s
 [\delta u_{e,s},[c_e,\delta u_{e,s}]]
 =-\frac16\frac18\sum_s\operatorname{ad}_{\delta u_{e,s}}^2c_e.
\]
Polarize in \(u,w\).  Taking the real colour trace, the two
operator orders have equal traces by cyclicity, giving (66.1).
The factor \(1/6\) includes the two derivatives in \(t\);
it is not the coefficient of \(t^2z\) itself.
\end{proof}

\begin{corollary}[Complete measure-only quadratic form on this subspace]
\label{cor:v62-density}
The logarithmic density \(\ell\) from (63.1) has
\[
\begin{split}
 D^2\ell(0)[u,w]={}&\frac1{12}\sum_{e,s}
       \Tr_{\mathfrak g}(\operatorname{ad}_{u_{e,s}}
                         \operatorname{ad}_{w_{e,s}})\\
 &+\frac16\sum_e\frac18\sum_s
       \Tr_{\mathfrak g}(\operatorname{ad}_{\delta u_{e,s}}
                         \operatorname{ad}_{\delta w_{e,s}}).
\end{split}                                                     \tag{66.2}
\]
For boundary-supported fibre tangents \(\bar u_e=\bar w_e=0\),
\[
 D^2\ell(0)[u,w]=\frac5{48}\sum_{e,s}
       \Tr_{\mathfrak g}(\operatorname{ad}_{u_{e,s}}
                         \operatorname{ad}_{w_{e,s}}).          \tag{66.3}
\]
On the diagonal this is nonpositive; it vanishes precisely when
all displayed Lie-algebra values are central.  For section
directions \(u=Rc,w=Rd\), (66.2) reduces to (65.5).
\end{corollary}
\begin{proof}
Substitute (66.1) and \(A_u=A_w=0\) in (63.3).
For fibre tangents the centred and original values agree, and
\(1/12+1/48=5/48\).  With an invariant positive inner product,
\(\operatorname{ad}_u\) is skew-adjoint, so
\(\Tr\operatorname{ad}_u^2=-\|\operatorname{ad}_u\|_{\rm HS}^2\).
Its kernel as a function of \(u\) is the centre of the Lie algebra.
For section directions the centred terms vanish and the eight
equal boundary contributions give the factor \(8/12=2/3\).
\end{proof}

\begin{record}[Exact noncommuting third-jet check]
\path{scripts/verify_ym_v62_boundary_density.py}
computes the reference mean through third order using matrix
exponential and logarithm series, rather than inserting (66.1).
On one coarse edge choose eight values
\((X,-X,Y,-Y,0,0,0,0)\), with \(X=i\sigma_1\), \(Y=i\sigma_2\).
Polarization against all three SU(2) colour directions gives
\[
 \Tr B_3[u,u,R\,\cdot]=\frac23,\quad
 D^2\log j(0)[u,u]=-\frac83,\quad
                         D^2\ell(0)[u,u]=-\frac{10}{3}.         \tag{66.4}
\]
All operations used for the third-jet trace are Gaussian-rational;
the Haar value uses the stated adjoint trace formula.  The source
digest is
\texttt{57f74b7d931d06bf60b43e89b1d752ae4749566708a93c6cfe3cde900c2bfd71};
the payload digest in \path{artifacts/ym_v62_boundary_density.json} is
\texttt{1d92769de487dc3a71657b647624faae6a445efdb9514c85f088f8f08303b0d6}.
The test is in an SU(2) subgroup and does not set a universal
normalization for other gauge groups.
\end{record}

\begin{assessment}[Interpretation and limits]
The new nonlinear fibre chart has a nonzero local measure
correction on this explicit subspace.  Its negative logarithmic
curvature is not a physical mass term: the action, the remaining
fibre directions, the gauge quotient and the continuum scaling
have not been combined.  It is also not an extra correction to
the previously proved Haar-translation endpoint; it belongs to
the different nonlinear coordinate construction of V57--V58.
For general fine fields with internal components, transport and
dressing derivatives still contribute.  Their mixed terms are
not bounded or evaluated by (66.2).
\end{assessment}

\begin{decision}[Next nontrivial sector]
Compute the mixed internal/boundary terms of the same density,
using the fixed internal paths.  Test the complete local density
against gauge directions and a specified quotient chart before
inferring a determinant or a physical response.  The exactly
computed boundary sector and the section benchmark should be
retained as regression tests for that larger calculation.
\end{decision}

\begin{firewall}[Position after V62]
A nontrivial restricted fibre-density quadratic form and its
third-jet contraction are proved and exactly checked.  The full
fibre density, gauge reduction, global blocking measure,
interacting continuum and Yang--Mills mass gap remain open.
Only V62 remains active; prior bodies and archives are retained.
The next companion audit is V65.
\end{firewall}

% END APPENDED POST-TENTH-ARCHIVE V62 MATERIAL

% =============================================================================
\section{V63: exact local factorization of the nonlinear fibre Jacobian}
\label{sec:v63-local}
% =============================================================================

The predecessor V62 has SHA--256
\texttt{53AEE781E054E57266C89B48AAC31DBAA03123DF68579A567F526CFF0E12E68B}.
Its body is retained apart from the active-version label.
The boundary calculation suggests examining the full dependency
structure before expanding all internal/boundary derivatives.
Keep the fixed internal root paths and reference convention of V57.

\begin{theorem}[Block-diagonal coarse Jacobian]
\label{thm:v63-jacobian}
For the logarithmic block \(F\), let \(e=(r,\mu)\) index a
coarse link and let \(\partial_e\) be its family of eight boundary
fine links.  The component \(F_e(x)\) depends only on:
the boundary coordinates in \(\partial_e\), and a fixed finite
set of internal coordinates in the source and target cells.
It is independent of every other boundary family.
Consequently
\[
 (DF(x)R)_{ef}=0\quad(e\ne f),\qquad
 DF(x)R=\bigoplus_e M_e(x),                        \tag{67.1}
\]
where \(M_e(x):\mathfrak g\to\mathfrak g\) is the derivative
when the same Lie-algebra increment is added to all eight
coordinates of \(\partial_e\).  In the V58 chart,
\[
 \det(DF(x)R)=\prod_e\det M_e(x),\qquad
 \ell(x)=\sum_l\log j_G(x_l)-\sum_e\log\det M_e(x).
                                                               \tag{67.2}
\]
Each determinant is positive near zero.
\end{theorem}
\begin{proof}
Both root-to-site transports in (61.1), and both dressings in
(61.2), use only internal links.  Each two-link path \(V_x\)
has exactly one boundary link, belonging to \(\partial_e\),
and one internal link in the source or target cell.  All relative
path products in the mean retain this same dependency set.
Thus varying a distinct boundary family cannot affect \(F_e\).
The definition of \(R\) now proves (67.1).  The determinant of a
block-diagonal matrix is the product of its block determinants.
Each block is the identity at zero and remains invertible in
the V58 chart, hence has positive determinant there.  Taking
logarithms gives (67.2).  This factorization concerns \(DFR\),
not the generally rectangular full derivative \(DF\).
\end{proof}

\begin{corollary}[The implicit fibre solve is local in the coarse variables]
For fixed \(u\in\ker L\), the V58 equation
\(F(u+Rv)=c\) splits into independent equations
\[
 F_e(u+R_ev_e)=c_e,                               \tag{67.3}
\]
where the notation in each equation retains all internal entries
of \(u\) and only its relevant boundary family.
The solution \(v_e\) depends only on those local entries and
\(c_e\), not on \(c_f\) for \(f\ne e\).
\end{corollary}
\begin{proof}
The dependency assertion in the theorem applies throughout the
small-field neighbourhood, not just at zero.  Thus the equations
decouple in \(v\).  V58's contraction proof applies to each block
and gives its unique local solution and stated dependence.
\end{proof}

\begin{proposition}[Local independent coordinates on the linear fibre]
\label{prop:v63-coordinates}
For each \(\partial_e\), choose one boundary link as a pivot.
All internal fine coordinates and the other seven boundary
coordinates per coarse link are independent coordinates on
\(\ker L\).  The pivot is determined by
\[
 u_{e,0}=-\sum_{s=1}^7u_{e,s}-8(L_{
 \rm int}u)_e,                                   \tag{67.4}
\]
where \(L_{\rm int}\) is the contribution of internal links
to \(L\).  This is a fixed local linear parametrization with
60 Lie-algebra coordinates per parity cell in four dimensions.
\end{proposition}
\begin{proof}
Each boundary value appears twice in the sixteen two-link paths,
so \((Lu)_e=\frac18\sum_{s=0}^7u_{e,s}+(L_{\rm int}u)_e\).
Solving \((Lu)_e=0\) gives (67.4).  No pivot belongs to two
different boundary families, so there is no recursive elimination.
There are 32 internal links and \(4\cdot7=28\) free boundary
links per cell.  The parametrization is an injective linear map
onto \(\ker L\), with a constant positive Euclidean volume
Jacobian on each finite lattice.  That constant need not itself
factor cellwise; it is independent of the field.
\end{proof}

\begin{theorem}[Uniform finite-range density bounds]
\label{thm:v63-density-bound}
On a sufficiently small closed subchart, the logarithmic density
\(\ell\) is a sum of finite-range smooth functions with derivative
bounds of any fixed order independent of volume.  In particular
there is a finite constant \(C\), independent of the number of
cells, such that
\[
 |D^2\ell(x)[a,b]|\le C\|a\|_2\|b\|_2             \tag{67.5}
\]
for fine-coordinate Euclidean norms and \(x\) in that subchart.
This is a bound on the Hessian as an operator, not a bound on
the total logarithmic density independent of volume.
\end{theorem}
\begin{proof}
Every term in (67.2) has a fixed finite stencil.  The smooth
finite-dimensional functions defining each \(M_e\) have uniformly
bounded derivatives on the smaller closed chart, and its inverse
has the V58 bound.  Differentiating \(\log\det M_e\) therefore
gives uniformly bounded local derivatives.  The same holds for
the one-link Haar terms.  If \(S_e\) is the stencil of a local
term, its Hessian is bounded by
\(C_0\|a|_{S_e}\|_2\|b|_{S_e}\|_2\).
Each fine link belongs to at most a fixed number \(D\) of these
stencils.  Summation and Cauchy--Schwarz yield
\(C_0D\|a\|_2\|b\|_2\), plus the one-link contribution,
proving (67.5).  Finite-order higher derivatives follow by the
same local differentiation, without a claim about an all-orders
constant.  All constants may depend on the fixed paths, group
and chosen small-field radius.
\end{proof}

\begin{assessment}[What this removes from the computation]
The Jacobian term does not require a determinant of a growing
coupled matrix: its exact factors have size \(\dim\mathfrak g\).
The general cubic contraction and quadratic Jacobian term can
therefore be computed per coarse edge with its internal stencil.
After substituting the local fibre solve and (67.4), the local
action and density remain finite-range functions of the independent
coordinates, up to field-independent volume constants.
This locality does not imply that the resulting measure is a
product measure or that its correlations decay.  Constraints of
the small-field domain, gauge reduction, and interaction remain.
Nor does a bounded density Hessian guarantee a positive physical
action Hessian or a convergent polymer expansion.
\end{assessment}

\begin{decision}[Next exact local calculation]
Construct each fixed-size \(M_e\) jet with its internal transport
variables and compute \(-D^2\log\det M_e\) directly.  Verify the
section identity and the boundary-fibre value from V61--V62.
Then compare local gauge reduction with the dressed transformation
law; locality alone does not supply a conventional coarse gauge
action or reflection positivity.
\end{decision}

\begin{firewall}[Position after V63]
Exact Jacobian factorization, local fibre coordinates, and a
volume-uniform finite-range density Hessian bound are proved.
The explicit full density coefficients, quotient compatibility,
global compact measure, interacting continuum and Yang--Mills
mass gap remain open.  Only V63 remains active; previous bodies
and archives are retained.  The next companion audit is V65.
\end{firewall}

% END APPENDED POST-TENTH-ARCHIVE V63 MATERIAL

% =============================================================================
\section{V64: fibre-preserving dressed tree reduction and its Haar measure}
\label{sec:v64-dressed-tree}
% =============================================================================

The predecessor V63 has SHA--256
\texttt{ADC1B404B29F44C3BC5809FDBE378007B5DF5EAA5A14AD46283142F60BD41AA3}.
Its body is retained apart from the active-version label.
We go upstream of the mixed density expansion: first remove the
internal tree variables in a way compatible with the nonlinear
block, rather than assuming the linear representative is a
nonlinear gauge slice.  The new result concerns the V57 block
and its measure, not a repetition of the earlier linear tree gauge.

\begin{definition}[Tree-compatible choice of the V57 paths]
In each cell fix a rooted spanning tree on its sixteen sites.
For this section the paths defining \(T_x(U)\) in V57 are the
unique paths of this tree from its root \(r\) to \(x\).
This specializes the previously allowed fixed internal paths;
no assertion that arbitrary path families form a tree is made.
Use the same translated tree in every cell.  There are fifteen
tree links per cell.  All identities below hold on a neighbourhood
of the identity where the V57 logarithms and their transformed
arguments are defined.  Gauge action is
\((g\cdot U)_{xy}=g_xU_{xy}g_y^{-1}\).
Write \(r(x)\) for the cell root of \(x\), and set
\[
 h_r(U)=\exp\left(\frac1{16}\sum_{x\in C_r}\log T_x(U)\right),
 \qquad
 \gamma_x(U)=h_{r(x)}(U)^{-1}T_x(U),\qquad
 Q(U)=\gamma(U)\cdot U.                         \tag{68.1}
\]
Here \(T_r=I\), so \(\gamma_r=h_r^{-1}\).
\end{definition}

\begin{theorem}[Exact coarse-field-preserving tree retraction]
\label{thm:v64-retraction}
The field \(Q(U)\) has every tree link equal to \(I\), satisfies
\[
 T_x(Q(U))=I,\qquad h_r(Q(U))=I,\qquad
 \mathcal B(Q(U))=\mathcal B(U),\qquad S(Q(U))=S(U), \tag{68.2}
\]
where \(S\) is a gauge-invariant Wilson action.  On the tree slice
\(Q\) is the identity map; in particular \(Q(Q(U))=Q(U)\).
For every admissible fine gauge transformation \(g\),
\[
 Q(g\cdot U)=\sigma_U(g)\cdot Q(U),\qquad
 \sigma_{U,r}(g)=h_r(g\cdot U)^{-1}g_rh_r(U),     \tag{68.3}
\]
where on the right \(\sigma_U(g)\) is constant on each fine cell.
\end{theorem}
\begin{proof}
Telescoping on a tree path gives
\(T_x(g\cdot U)=g_rT_x(U)g_x^{-1}\).
Apply this to \(g=\gamma(U)\).  Then
\[
 T_x(Q(U))=h_r^{-1}T_x(U)
                  \big(h_r^{-1}T_x(U)\big)^{-1}=I.
\]
The tree path to each child is its parent's path followed by
the tree edge (or its inverse).  The path identities therefore
force every tree link to be \(I\).  The definition of \(h\)
gives \(h(Q(U))=I\).  The V57 transformation law for
\(\mathcal B\) now has coarse gauge element
\[
 \sigma_{U,r}(\gamma(U))
       =h_r(Q(U))^{-1}\gamma_r(U)h_r(U)=I.
\]
This proves preservation of \(\mathcal B\); ordinary fine gauge
invariance proves preservation of \(S\).
If \(U\) already lies on the tree slice, \(T=h=\gamma=I\),
proving the retraction statement.
Finally direct substitution gives
\[
 \gamma_x(g\cdot U)=
       \sigma_{U,r(x)}(g)\gamma_x(U)g_x^{-1}.
\]
Applying this identity at the two endpoints of each link proves
(68.3).  No equivariance of an average of logarithms under
independent endpoint transformations has been assumed.
\end{proof}

\begin{theorem}[Local product Haar coordinates without a new determinant]
\label{thm:v64-haar}
Let \(E_{\rm nt}\) be the non-tree fine links.  On a sufficiently
small product neighbourhood, the variables
\[
 (T_x)_{x\ne r(x)},\qquad (Q_e)_{e\in E_{\rm nt}}
\]
are smooth coordinates for the fine group links, with inverse
\[
 U_{xy}=T_x^{-1}h_{r(x)}(T)\,Q_{xy}\,
                         h_{r(y)}(T)^{-1}T_y,     \tag{68.4}
\]
where tree \(Q\)'s and root \(T\)'s are \(I\).
For normalized Haar measures the change of variables is exactly
\[
 \prod_{e\ {\rm fine}}dU_e
       =\prod_{x\ne r(x)}dT_x\prod_{e\in E_{\rm nt}}dQ_e .
                                                        \tag{68.5}
\]
Thus a product patch in these coordinates has no additional
field-dependent tree-reduction Jacobian.
\end{theorem}
\begin{proof}
First use ordinary tree coordinates
\(W_{xy}=T_xU_{xy}T_y^{-1}\), with \(W=I\) on tree links.
Order the vertices by distance from each root.  Given a parent's
transport, its child's transport is that transport times the
tree link or its inverse.  Conditional Haar measure is preserved
by translation and inversion.  Induction therefore changes the
product measure on tree edges into the independent product of
the non-root \(dT_x\)'s.  Conditional on these variables,
each non-tree \(U_e\mapsto W_e\) is a fixed left and right
translation.  This proves \(dU=dT\,dW\), for example by Fubini
on arbitrary continuous test functions.

For fixed \(T\), replace each non-tree variable by
\(Q_{xy}=h_{r(x)}(T)^{-1}W_{xy}h_{r(y)}(T)\).
Another application of left and right Haar invariance gives
\(dT\,dW=dT\,dQ\).  Although these translations depend on \(T\),
the \(T\) variables themselves are unchanged; the conditional
measure argument therefore includes all cross derivatives.
Solving for \(U\) gives (68.4).  On tree links its right side
is \(T_x^{-1}T_y\), whose path products recover exactly the
specified \(T\).  This verifies the inverse and smoothness.
Shrinking the two factors of the product patch ensures that
all original and transformed V57 chart arguments stay in their
logarithm neighbourhoods.  Haar factorization holds on that
patch as equality of restricted measures.
\end{proof}

\begin{corollary}[What can be integrated out locally]
For a product patch \(A_T\times A_Q\) as above and an integrable
function \(f\) of the dressed coarse field,
\[
 \int_{U(A_T\times A_Q)}
      e^{-S(U)}f(\mathcal B(U))\,dU
 =\operatorname{vol}_{\rm Haar}(A_T)
       \int_{A_Q}e^{-S(Q)}f(\mathcal B(Q))\,dQ .   \tag{68.6}
\]
This does not give the same formula for an arbitrary fine-link
small-field cutoff.
\end{corollary}
\begin{proof}
Use (68.2), (68.5), and Fubini.  In a general cutoff domain its
indicator becomes a function of both \(T,Q\); integrating it
may produce a \(Q\)-dependent weight instead of the constant
in (68.6).  That dependence has not been bounded here.
\end{proof}

\begin{proposition}[Reduced fibre dimension and second jet]
\label{prop:v64-reduced}
On the tree slice there are 49 independent group links per
cell: 17 non-tree internal and 32 boundary links.
The logarithmic coarse block restricted to this slice is a
local submersion with the same boundary right inverse \(R\).
Its fibre at any sufficiently small coarse field has dimension
\(45\dim\mathfrak g\) per cell.  For a Lie-algebra direction \(a\)
zero on the tree, its second jet at zero is
\[
 D^2F(0)[a,a]_{r\mu}
   =\frac1{16}\sum_{x\in C_r}
                       [a_{x,\mu},a_{x+\hat\mu,\mu}]. \tag{68.7}
\]
\end{proposition}
\begin{proof}
Removing fifteen tree edges from 64 leaves 49.
The injection \(R\) uses boundary edges only, so its image
lies in the slice and \(DF(0)R=LR=I\).
The local implicit-function theorem, or the V58 contraction
restricted to this subspace, eliminates four Lie-algebra
variables per cell, leaving 45.  In exponential coordinates
on the slice the tree links remain identically \(I\).
All tree transports and dressings are then identically \(I\),
so \(P_x=U_{x,\mu}U_{x+\hat\mu,\mu}\).
The V57 second-jet formula has \(\tau_x=\eta_x=0\);
only the two-link commutator remains, giving (68.7).
This is a local finite-dimensional fibre count, not a count
of physical continuum particle degrees of freedom.
\end{proof}

\begin{assessment}[Residual symmetry and the density still to compute]
Cell-constant gauge transformations preserve the tree slice
and act on its coarse field by the ordinary endpoint rule.
At fixed coarse field the corresponding stabilizer remains.
The retraction therefore does not assert a complete global
gauge quotient.  Nor does (68.5) cancel the V58 coarse-constraint
Jacobian: after choosing exponential coordinates for the
remaining links and solving their nonlinear coarse constraint,
their Haar factors and the corresponding \(\det(DFR)\) are
still present.  V59--V63 concern that other change of variables.
The reduced version now needs only the non-tree link Haar
factors and the slice restriction of the block derivative.
The raw and reduced logarithmic densities must not be equated
while retaining different coordinate volume elements.
\end{assessment}

\begin{audit}[Updated companion snapshot requested by the user]
The active YM predecessor was V63, not one of the unrelated
candidate files in the repository root.  SM has advanced to
V66, SHA--256
\texttt{E6A46F6DC67942974FACC24AAEBADDBA9F396ABDC2B3D5FE98147CFEEB5E918F}.
Its final section proves cutoff removal for a specified
classical collective saturated evolution using uniform
physical invariants and an occupation \(L^1\) recovery identity.
It expressly does not establish a spatially local quantum
gauge theory.  Those hypotheses have no proved counterpart
in the present interacting YM measure; no YM gap estimate
is imported.  NS remains V26 with the unchanged hash
\texttt{82C887F8C2C69E4F28FAD7C3DC8DE5B9099CB5A4BF431EBA1FFF3E91935978AD}.
The regular five-version companion checkpoint remains V65.
\end{audit}

\begin{decision}[Next upstream test]
Compute the reduced coarse Jacobian through quadratic order
using \(P_x=Q_{x,\mu}Q_{x+\hat\mu,\mu}\), and distinguish
the boundary-only benchmark from mixed non-tree internal
directions.  The product-patch reduction does not control
the complement of the small-field domain; that issue must
be retained when defining any global blocked measure.
\end{decision}

\begin{firewall}[Position after V64]
The dressed block admits an exact local fibre-preserving
tree retraction, with product Haar coordinates and a
45-coordinate local fibre.  These are finite-lattice results.
Global domain control, interacting scale-uniform estimates,
continuum reconstruction, and a positive physical mass gap
have not been proved.  The target remains the full problem
in the \href{https://www.claymath.org/library/monographs/MPPc.pdf}
{Jaffe--Witten formulation}, not merely a local Hessian bound.
Only V64 is active in this lineage; previous bodies and all
archives are retained.
\end{firewall}

% END APPENDED POST-TENTH-ARCHIVE V64 MATERIAL

% =============================================================================
\section{V65: explicit mixed quadratic density on the tree-reduced fibre}
\label{sec:v65-mixed-density}
% =============================================================================

The predecessor V64 has SHA--256
\texttt{3C3F4BE3CC594EE82669BA6532A7103A41CA0C745370EB2762953D39FBEF2C49}.
Its body is retained apart from the active-version label.
We now compute the third-jet trace needed for the full
quadratic density in the reduced chart; no sign of this
density is identified with a physical mass.

For a fixed coarse edge \(e=(r,\mu)\), average over the sixteen
two-link paths starting at \(x\in C_r\), denoting this average
by brackets \(\langle\cdot\rangle_e\).
For a direction \(a\) on the tree slice write \(i_x\) for its
internal-link value and \(b_x\) for its boundary-link value
along that path.  Put \(\epsilon_x=1\) if the internal link is
first, and \(-1\) if it is second, and define
\[
 p_x=i_x+b_x,\quad \bar p_e=\langle p_x\rangle_e,\quad
 d_x=p_x-\bar p_e,\quad
 k_e=\tfrac12\langle\epsilon_x i_x\rangle_e.       \tag{69.1}
\]
Tree internal-link values are zero.  A boundary link occurs in
two paths; averages here are path averages, not sums over
independent link variables.  Let \(F_{\rm tr}\) be the restriction
of \(F\) to the tree slice, \(L_{\rm tr}=DF_{\rm tr}(0)\), and
\[
 \ell_{\rm tr}(a)
  =\sum_{l\in E_{\rm nt}}\log j_G(a_l)
                -\log\det(DF_{\rm tr}(a)R).       \tag{69.2}
\]
This is the measure-only logarithmic density relative to the
remaining exponential coordinates and the same local splitting,
up to a field-independent constant.

\begin{theorem}[Mixed coarse-Jacobian trace]
\label{thm:v65-trace}
With \(B_2=D^2F_{\rm tr}(0)\), \(B_3=D^3F_{\rm tr}(0)\), the
diagonal coarse block of \(B_2[a,R\,\cdot]\) is
\(\operatorname{ad}_{k_e}\).  Moreover
\[
 \Tr_{\mathfrak g}\big(B_{3,e}[a,a,R_e\,\cdot]\big)
 =\frac16\left\langle
   \Tr_{\mathfrak g}\left(
      \operatorname{ad}_{i_x}^{\,2}
     -\operatorname{ad}_{b_x}\operatorname{ad}_{i_x}
     -\operatorname{ad}_{d_x}^{\,2}\right)
                         \right\rangle_e.        \tag{69.3}
\]
The result is independent of the chosen reference path.
\end{theorem}
\begin{proof}
Polarizing (68.7), and observing that \(R_ev\) changes only
the boundary member of each path by \(v\), gives
\[
 B_{2,e}[a,R_ev]
       =\tfrac12\langle\epsilon_x[i_x,v]\rangle_e=[k_e,v].
\]
For the cubic term write the ordered path values as
\((A_x,B_x)=(i_x,b_x)\) or \((b_x,i_x)\).  The BCH formula and
the V60 group-mean cubic identity give
\[
 \begin{split}
 [t^3]F_{{\rm tr},e}(ta)=\frac1{12}\big\langle&
 [A_x,[A_x,B_x]]+[B_x,[B_x,A_x]]\\
 &+[d_x,[p_*,d_x]]\big\rangle_e .
 \end{split}                                    \tag{69.4}
\]
Here \(p_*\) is the first-order value of the reference path.
The second-order terms in each path logarithm contribute to
the mean of logarithms, already included in BCH; the group-mean
cubic correction depends only on their first-order terms.

Replace \(a\) by \(a+sR_ev\).  Both \(p_x\) and \(\bar p_e\)
increase by \(sv\), so \(d_x\) is unchanged and \(p_*\)
increases by \(sv\).  Twice the derivative in \(s\) of
(69.4) is \(B_{3,e}[a,a,R_ev]\): this follows by polarizing
the symmetric cubic Taylor coefficient \(B_3[a,a,a]/6\).
For an internal-first path the differentiated BCH contribution
is therefore
\[
 \frac16\big([i_x,[i_x,v]]+[v,[b_x,i_x]]
                                      +[b_x,[v,i_x]]\big).
\]
For an internal-second path the same three terms occur with
their order in the sum interchanged.  As a linear map of \(v\),
the middle term is an adjoint map with zero trace, and the
other two have trace
\(\Tr(\operatorname{ad}_{i_x}^2-
       \operatorname{ad}_{b_x}\operatorname{ad}_{i_x})/6\).
The mean correction contributes
\([d_x,[v,d_x]]/6=-\operatorname{ad}_{d_x}^2v/6\).
Adding and averaging proves (69.3).  The derivative removed
\(p_*\) altogether, establishing reference independence.
\end{proof}

\begin{theorem}[Full reduced quadratic measure density]
\label{thm:v65-density}
One has \(D\ell_{\rm tr}(0)=0\), and for every slice direction
\[
 \begin{split}
 D^2\ell_{\rm tr}(0)[a,a]
 ={}&\frac1{12}\sum_{l\in E_{\rm nt}}
                           \Tr_{\mathfrak g}\operatorname{ad}_{a_l}^2\\
 &-\frac16\sum_e\left\langle\Tr_{\mathfrak g}
       (\operatorname{ad}_{i_x}^2
       -\operatorname{ad}_{b_x}\operatorname{ad}_{i_x}
       -\operatorname{ad}_{d_x}^2)\right\rangle_e
       +\sum_e\Tr_{\mathfrak g}\operatorname{ad}_{k_e}^2 .
 \end{split}                                    \tag{69.5}
\]
Its bilinear form is determined by polarization.  On the
linear fibre \(\ker L_{\rm tr}\), the middle line simplifies to
\[
 \frac16\sum_e\left\langle\Tr_{\mathfrak g}
       (3\operatorname{ad}_{b_x}\operatorname{ad}_{i_x}
                     +\operatorname{ad}_{b_x}^2)\right\rangle_e.
                                                        \tag{69.6}
\]
\end{theorem}
\begin{proof}
The V59 determinant differentiation argument applies on this
slice.  First derivatives of the Haar terms vanish, and
\(\Tr\operatorname{ad}_{k_e}=0\).  For the second derivative,
each Haar link contributes
\(\Tr\operatorname{ad}_{a_l}^2/12\).
For \(M_e(a)=D F_{{\rm tr},e}(a)R_e\), \(M_e(0)=I\), the
second derivative of \(-\log\det M_e\) is
\(-\Tr B_{3,e}[a,a,R_e\,\cdot]
 +\Tr(B_{2,e}[a,R_e\,\cdot]^2)\).
Apply (69.3) and the first assertion of the preceding theorem.
On \(\ker L_{\rm tr}\), \(\bar p_e=0\), so \(d_x=i_x+b_x\).
Expand its squared adjoint map and use cyclicity of trace to
obtain (69.6).
\end{proof}

\begin{corollary}[Actual constrained-chart Hessian and regressions]
Let \(\Psi_{\rm tr}(u,c)\) be the local implicit slice chart.
At \(u=c=0\), for \(u\in\ker L_{\rm tr}\), the logarithm of
the Haar/action density has second derivative
\[
 -D^2S(I)[u,u]+D^2\ell_{\rm tr}(0)[u,u].          \tag{69.7}
\]
For boundary-only fibre directions, (69.5) is exactly
\((5/48)\sum_{\rm boundary}\Tr\operatorname{ad}_{u_l}^2\).
On the exact section \(a=Rc\) it equals
\((2/3)\sum_e\Tr\operatorname{ad}_{c_e}^2\).
\end{corollary}
\begin{proof}
The derivative of the implicit chart in a fibre direction is
the inclusion.  Both \(DS(I)\) and \(D\ell_{\rm tr}(0)\) vanish,
so the chart's second derivative makes no contribution.
This proves (69.7).
For the first regression \(i=k=0\), \(\bar p=0\).
The middle term contributes \(1/(6\cdot8)\) per boundary link,
in addition to its Haar coefficient \(1/12\).
For the section, \(i=k=d=0\) and each coarse \(c_e\) occurs
on eight links, giving \(8/12=2/3\).
\end{proof}

\begin{record}[Exact mixed algebra regression]
The script \texttt{scripts/verify\_ym\_v65\_reduced\_density.py}
uses exact rational complex matrices and truncated exponential
and logarithm products to compute the mixed third derivative
independently of (69.3).  A sixteen-path \(SU(2)\) tuple with
both internal-first and internal-second orders, and eight
boundary values each repeated twice, gives
\[
 \Tr B_{3,e}[a,a,R_e\,\cdot]=113/384,\qquad
 \Tr\operatorname{ad}_{k_e}^2=-1/256.
\]
Reference choices 0, 7 and 15 give the same third-jet trace
and agree with (69.3).  Results are saved in
\texttt{artifacts/ym\_v65\_reduced\_density.json}.
This is an algebraic tuple regression, not an assignment
certificate on a full periodic lattice or a mass-gap test.
The general lattice result is the analytic proof above.
\end{record}

\begin{audit}[Required five-version companion check]
During this iteration SM advanced from V66 to V67, with hash
\texttt{1ED7BD7C2FF7C25DB6D1E8CD3688259DD70D4776AC9A69D5010AD0BCD011D8BD}.
Its endpoint constructs a conserved renormalized energy on
a finite-energy class of its classical collective model,
while distinguishing that global energy from a local stress
tensor or microscopic quantum construction.  It does not
provide a spectral estimate for our constrained YM measure.
NS remains V26: its endpoint explicitly retains global
regularity as open after the pointwise rank-one norm result.
Parallel YM remains V5: its extensive-charge tightness
obstruction cautions against global volume-uniform tails
for total charge.  Neither the present local Jacobian nor
the tree reduction bypasses that obstruction.
No companion theorem changes the present calculation's
hypotheses.  Next compulsory audit and ten-version review: V70.
\end{audit}

\begin{decision}[Next meaningful test, not another formal jet]
The complete quadratic reduced measure term is now explicit.
The next issue is whether the constrained Wilson quadratic
form dominates it in a specified small-field regime, with
constants uniform in volume, and whether this extends beyond
the Hessian at the identity.  Such a conditional local
coercivity bound would still require large-field and
scale-iteration control before any continuum conclusion.
Do not infer a sign for the mixed expression (69.6) merely
from compactness of the Lie algebra.
\end{decision}

\begin{firewall}[Position after V65]
The full mixed quadratic reduced density, its reference
independence and two benchmark restrictions are proved.
The exact regression passes.  These results do not establish
an interacting continuum theory or the Yang--Mills mass gap.
Only V65 is active; the earlier bodies and archives remain.
\end{firewall}

% END APPENDED POST-TENTH-ARCHIVE V65 MATERIAL

% =============================================================================
\section{V66: averaged-block coercivity and nonlinear local convexity}
\label{sec:v66-convexity}
% =============================================================================

The V65 predecessor has SHA--256
\texttt{8243E6581BB605975C2205796D80D67CD3BC8357839719C17FE03014FBE343DB}.
Its body is retained except for the version label.
The archives already contain local reduced-filling coercivity.
The extension here verifies it for the averaged block of V57
and the dressed tree chart of V64, then includes the nonlinear
constraint and measure density in the Hessian.  We do not
relabel the old Gaussian result as a new mass-gap theorem.

Use periodic even lattices with at least two cells in each
direction.  Count each positively oriented edge and each
unoriented elementary plaquette once.  Fix the cell tree whose
parent of a nonzero vertex \(v\in\{0,1\}^4\) is obtained by
clearing its highest-index nonzero bit.  This is an explicit
instance of the tree choice in V64.  Norms on the Lie algebra
come from a faithful unitary matrix representation,
\(\|X\|^2=-\operatorname{Re}\Tr X^2\).
Let \(d\) be the ordinary linear plaquette curl.

\begin{lemma}[Internal cell constant]
For a cell field zero on the fifteen tree edges,
\[
 \sum_{\text{internal links of }C_r}\|a_l\|^2
 \leq C_{\rm int}
       \sum_{\text{plaquettes inside }C_r}\|(da)_p\|^2,
 \qquad C_{\rm int}=\frac{179}{12}.               \tag{70.1}
\]
\end{lemma}
\begin{proof}
Let \(D\) be the 24-by-17 scalar curl matrix on the remaining
internal edges.  It is injective: a curl-free edge field on
the cube is a vertex gradient (commuting adjacent steps across
square faces proves path independence); zero tree edges make
that vertex function constant.  Hence \(G=D^TD\) is positive
definite.  Its smallest eigenvalue is at least
\(1/\Tr G^{-1}\), since the largest eigenvalue of \(G^{-1}\)
is at most its trace.  The dependency-free rational script
\texttt{scripts/verify\_ym\_v66\_cell\_coercivity.py}
enumerates the edges and faces, inverts \(G\), and verifies
\(GG^{-1}=I\).  It gives
\(\det G=82944\), \(\Tr G^{-1}=179/12\).
Applying the scalar inequality to an orthonormal Lie-algebra
basis proves (70.1).  The constant is an upper bound, not an
optimal eigenvalue calculation.
\end{proof}

\begin{theorem}[Coercivity for the actual averaged linear fibre]
\label{thm:v66-linear}
On the tree slice with \(La=0\),
\[
 \|a\|_2^2\leq C_*\|da\|_2^2,\qquad
 C_*=\frac{2327}{24},\qquad \kappa=C_*^{-1}.       \tag{70.2}
\]
The constants do not depend on the number of cells.
\end{theorem}
\begin{proof}
Denote the total squared internal and boundary edge norms by
\(I_a,B_a\).  For a coarse edge \(e=(r,\mu)\), its eight
boundary values \(b_s\) form a three-dimensional cube indexed
by the transverse bits.  The graph Laplacian of this cube
has nonzero eigenvalues \(2,4,6\): its sign characters have
eigenvalue twice the number of nonconstant coordinates.
Consequently
\[
 \sum_s\|b_s-\bar b_e\|^2
       \leq\frac12\sum_{s\sim t}\|b_s-b_t\|^2.    \tag{70.3}
\]
Each difference on the right equals, up to orientation,
one crossing plaquette curl and two internal edge values.
The squared norm is at most three times the sum of their
squared norms.  In the sum over all coarse edges, each
internal fine edge occurs in exactly three such crossing
faces, one for each transverse direction.  Each crossing
face used here crosses precisely one cell interface and is
counted once.  Thus, writing \(P_{\rm cross}\) for their
total squared curls,
\[
 \sum_{e,s}\|b_s-\bar b_e\|^2
       \leq \tfrac32 P_{\rm cross}+\tfrac92 I_a. \tag{70.4}
\]
The averaged constraint reads
\(\bar b_e=-\frac1{16}\sum_{x\in C_r}i_x\), with the internal
path members \(i_x\) of V65.  Cauchy--Schwarz gives
\(8\|\bar b_e\|^2\leq\frac12\sum_x\|i_x\|^2\).
Each internal edge appears in two such path sums: for the
outgoing coarse edge of its cell and the incoming one.
Hence \(\sum_e8\|\bar b_e\|^2\leq I_a\).
Orthogonal mean/deviation decomposition in each boundary
family now gives
\[
 I_a+B_a\leq\tfrac32P_{\rm cross}+\tfrac{13}{2}I_a
 \leq\tfrac32P_{\rm cross}
               +\tfrac{13}{2}C_{\rm int}P_{\rm int}.
\]
The internal and crossing plaquette sets are disjoint;
plaquettes crossing two interfaces are not needed.
This proves (70.2) with
\(\max(3/2,13C_{\rm int}/2)=2327/24\).
\end{proof}

\begin{theorem}[Local interacting convexity, uniformly in volume]
\label{thm:v66-nonlinear}
Normalize the Wilson action by
\[
 S_0(Q)=\sum_p(\dim V-\operatorname{Re}\Tr Q_p),
 \qquad S_\beta=\beta S_0 .
\]
Let \(x=\Psi_{\rm tr}(u,c)\) be the V64--V65 implicit
exponential-coordinate chart, \(u\in\ker L_{\rm tr}\).
There exist constants \(\varepsilon>0\), \(M<\infty\),
independent of volume and of \(\beta\), such that for fixed
\(\|c\|_\infty<\varepsilon\) and
\(\|u\|_\infty<\varepsilon\),
\[
 D_u^2\!\left(S_\beta(\exp\Psi_{\rm tr}(u,c))
                 -\ell_{\rm tr}(\Psi_{\rm tr}(u,c))\right)[v,v]
 \geq\frac{\beta\kappa}{4}\|v\|_2^2              \tag{70.5}
\]
for every \(v\in\ker L_{\rm tr}\), whenever
\(\beta\geq\max(1,4M/\kappa)\).
The constants depend on the group representation and chosen
tree and chart.  No numerical value of \(M\) is asserted.
\end{theorem}
\begin{proof}
At the origin \(DS_0=0\), and expansion of a plaquette yields
\(D^2S_0(I)[v,v]=\|dv\|_2^2\).
The chart derivative in a fibre direction is the inclusion,
so (70.2) gives a lower bound \(\kappa\|v\|_2^2\) for the
composed action Hessian at \(u=c=0\).

For clarity, the extension away from zero includes the term
\(DS_0[D_u^2\Psi_{\rm tr}]\); it is not just restriction
of the ambient Hessian.  V63's local implicit solve remains
local after restricting the tree entries to zero.  Each
plaquette action composed with this solve is a smooth function
of a fixed finite stencil of \(u,c\), with uniformly bounded
third derivatives on a smaller closed chart.
The mean-value theorem, followed by the bounded-overlap
Cauchy--Schwarz estimate used in V63, therefore gives a
volume-independent \(C\) such that
\[
 |(D_u^2(S_0\circ\exp\Psi_{\rm tr})(u,c)
    -D_u^2(S_0\circ\exp\Psi_{\rm tr})(0,0))[v,v]|
 \leq C\max(\|u\|_\infty,\|c\|_\infty)\|v\|_2^2 .
\]
These estimates are on local functions before summation;
an extensive bound on the full gradient is not used.
Choose \(\varepsilon\) inside the chart and so that
\(C\varepsilon\leq\kappa/2\).  The composed action Hessian
is then at least \(\kappa/2\).

The non-tree Haar terms and the local invertible coarse
Jacobian blocks likewise give a finite-range smooth sum.
After composition with the local solve its Hessian has
absolute bound \(M\|v\|_2^2\), independent of volume.
This includes its own chart-curvature term.  Subtracting
that bound gives \(\beta\kappa/2-M\), which is at least
\(\beta\kappa/4\) for the stated \(\beta\).
The convex set \(\{u\in\ker L_{\rm tr}:\|u\|_\infty<
\varepsilon\}\) supplies a common fibre-coordinate domain.
\end{proof}

\begin{assessment}[Precisely what has and has not been stabilized]
Equation (70.5) is a strong-convexity result for the exact
local conditional action and density on a product chart.
It includes the mixed density from V65 and the nonlinear
constraint curvature, rather than only a Gaussian truncation.
It does not assert that a stationary point lies inside every
chosen chart, nor that arbitrary transformed cutoff indicators
are convex.  The original fine-link small-field domain may
give the additional \(Q\)-dependent weight identified in V64.
Integrating the coarse fields reintroduces retained modes;
the fibre curvature cannot be identified with a physical gap.
\end{assessment}

\begin{record}[Verification and next obstruction]
The exact cell check passed; its output is
\texttt{artifacts/ym\_v66\_cell\_coercivity.json}.
The global estimate is the boundary-face proof, not a finite
volume numerical extrapolation.  The natural next task is to
control conditional normalization and coarse derivatives for
this actual locally convex measure, keeping the chart boundary
and retained-field dependence explicit.  Scale stability and
the large-field complement remain required inputs.
The preceding turn was progress: it installed V64--V65 and
their exact mixed-density regression.  This turn extends that
state rather than restarting an old shell calculation.
\end{record}

\begin{firewall}[Position after V66]
Volume-uniform averaged-fibre coercivity and local nonlinear
strong convexity at sufficiently large \(\beta\) are proved.
No global interacting continuum construction or physical
Yang--Mills mass gap follows yet.  Only V66 remains active;
all earlier bodies and archives are retained.  The next
companion audit and ten-version review remain V70.
\end{firewall}

% END APPENDED POST-TENTH-ARCHIVE V66 MATERIAL

% =============================================================================
\section{V67: compact root block and the gauge-invariant observable bridge}
\label{sec:v67-compact-bridge}
% =============================================================================

The V66 predecessor has SHA--256
\texttt{CB3B874503BB983D055C0A6DFE3E5F9E4965CE83738A64CBFC7FBDF1572E9A9F}.
Its body is retained apart from the active-version label.
V66 stabilizes a local conditional density, but its product
patch is not a gauge-invariant domain.  Before differentiating
its normalization as if it were the full blocked action, we
construct an exact compact observable-level bridge.  This
implements one of the alternatives left open in V56.  It is
not a claim that the full compact density is strongly convex.

Let \(G\) be compact and connected with an invariant inner
product.  Fix a conjugation-invariant exponential neighbourhood
\(\mathcal O\) of \(I\) on which \(\log\) is single-valued
and smooth.  Such a neighbourhood is the exponential image
of a sufficiently small invariant ball in \(\mathfrak g\).
Keep the trees of V64.  Their transports \(T_x(U)\) are
products of links, so they exist on the entire compact
configuration space, without taking a logarithm.  Define
\[
 W_{xy}(U)=T_x(U)U_{xy}T_y(U)^{-1},\qquad W=I
       \text{ on tree links}.                    \tag{71.1}
\]
For a tree-slice field \(W\) and coarse edge \(e=(r,\mu)\),
put \(P_x(W)=W_{x,\mu}W_{x+\hat\mu,\mu}\), choose the same
reference path \(P_*\) as before, and let
\[
 \mathcal A_e=\{W:P_*^{-1}P_x\in\mathcal O
                         \text{ for all }x\in C_r\}.
\]
Define the total, measurable coarse block edgewise by
\[
 K_e(W)=
 \begin{cases}
 P_*\exp\big(\frac1{16}\sum_x\log(P_*^{-1}P_x)\big),
                                      &W\in\mathcal A_e,\\
 P_*,                                 &W\notin\mathcal A_e.
 \end{cases}                                      \tag{71.2}
\]
The second branch is a declared extension, not an estimate
that the first branch applies with high probability.

\begin{theorem}[Global equivariance and local comparison]
\label{thm:v67-block}
The map \(U\mapsto K(W(U))\) is measurable on the full finite
compact lattice.  It is equivariant for the fine gauge group
with the field-independent coarse transformation \(g\mapsto
(g_r)_r\).  In the common small-field chart of V57 and V64,
\[
 K(W(U))_{r\mu}
       =h_r(U)\,\mathcal B(U)_{r\mu}\,
                                      h_{r+2\hat\mu}(U)^{-1}.
                                                        \tag{71.3}
\]
Consequently every coarse gauge-invariant observable \(f\)
satisfies \(f(K(W(U)))=f(\mathcal B(U))\) on that chart.
\end{theorem}
\begin{proof}
Tree-path telescoping gives
\[
 W_{xy}(g\cdot U)=g_{r(x)}W_{xy}(U)g_{r(y)}^{-1}.
\]
Under a cell-constant transformation \(q\), both \(P_x\)
and \(P_*\) transform as \(q_rPq_{r'}^{-1}\).  Their relative
products transform by conjugation by \(q_{r'}\).
The set \(\mathcal O\) and its logarithm respect conjugation.
Thus \(\mathcal A_e\) is invariant, and both branches of
(71.2) transform as a coarse link.  Products, the smooth
logarithm on its domain, and the piecewise definition are
measurable, proving the global assertions.

Locally, V64 gives \(Q(U)=h(U)^{-1}\cdot W(U)\).
On the tree slice the V57 block is exactly the first branch
of (71.2), after reducing the common chart if necessary.
Its equivariance and V64's identity
\(\mathcal B(Q(U))=\mathcal B(U)\) prove (71.3).
Gauge invariance of \(f\) proves the final assertion.
\end{proof}

\begin{corollary}[The changed linearization is a coarse gradient]
Let \(\tau_r(a)=\frac1{16}\sum_{x\in C_r}\sum_{l\in[r,x]}a_l\),
using signed tree-path sums.  The logarithmic derivative of
the root block at \(I\) is
\[
 L_{\rm root}a=La+\tau_r(a)-\tau_{r'}(a)
                                    \quad\text{on }e=(r,r'). \tag{71.4}
\]
On the tree slice \(L_{\rm root}=L_{\rm tr}\).
\end{corollary}
\begin{proof}
Differentiate (71.3), using \(h_r(I)=I\) and
\(Dh_r(I)[a]=\tau_r(a)\).  Tree-zero directions have
\(\tau_r(a)=0\).  This exhibits the required comparison;
the root block does not secretly retain the averaged
field-independent gauge differential forbidden by V56.
\end{proof}

\begin{theorem}[Exact compact positive blocked probability]
\label{thm:v67-measure}
On any fixed finite periodic lattice and for finite
\(\beta\geq0\), define
\[
 d\mu_\beta(U)=Z_\beta^{-1}e^{-\beta S_0(U)}
                                   \prod_l dU_l .
 \qquad \nu_\beta=(K\circ W)_*\mu_\beta .          \tag{71.5}
\]
This is a normalized positive coarse probability measure,
invariant under the whole coarse gauge group, and
\[
 \nu_\beta(f)=
 \frac{\displaystyle\int_{G^{E_{\rm nt}}}
                 e^{-\beta S_0(W)}f(K(W))\,dW}
      {\displaystyle\int_{G^{E_{\rm nt}}}
                 e^{-\beta S_0(W)}\,dW}.         \tag{71.6}
\]
The tree integrations in this formula are global and exact.
\end{theorem}
\begin{proof}
The continuous Wilson weight on a compact space is positive
and integrable, so its partition function is finite and
strictly positive.  A measurable pushforward preserves
normalization and positivity.  Every coarse gauge
transformation lifts, for example, to a cell-constant fine
one.  Gauge invariance of \(\mu_\beta\) and equivariance
prove invariance of \(\nu_\beta\).
Ordinary tree coordinates \((T,W)\), unlike the logarithmic
dressing, are globally defined.  The Haar disintegration
argument of V64 applies globally to these ordinary
coordinates: \(dU=dT\,dW\).  Also \(S_0(U)=S_0(W)\).
Each normalized tree Haar integral is one, proving (71.6).
\end{proof}

\begin{proposition}[Local extension sensitivity without a global good event]
Let \(\widetilde K\) be any other measurable extension of the
same first branch, edgewise agreeing on each \(\mathcal A_e\).
For any bounded observable \(f\) depending only on a finite
set \(E_f\) of coarse links,
\[
 \begin{split}
 |\mathbb E_{\mu_\beta}f(K(W))
       -\mathbb E_{\mu_\beta}f(\widetilde K(W))|
 &\leq2\|f\|_\infty
       \mu_\beta\big(W\in\bigcup_{e\in E_f}\mathcal A_e^c\big)\\
 &\leq2\|f\|_\infty
               \sum_{e\in E_f}\mu_\beta(W\notin\mathcal A_e).
 \end{split}                                    \tag{71.7}
\]
\end{proposition}
\begin{proof}
Outside the indicated union the two block fields agree
on every argument used by \(f\).  On the union their
observable values differ by at most \(2\|f\|_\infty\).
Integrate and apply the union bound.
\end{proof}

\begin{assessment}[Why this does not close the complement]
The construction populates an exact compact blocked
probability and its gauge-invariant observable comparison,
not a regulator-uniform construction of a continuum state.
The fallback may be discontinuous at the chart boundary;
no smooth global density or reflection positivity is proved.
Its choice can matter until the right side of (71.7) is
controlled in the relevant limiting regime.
In particular we have not estimated those probabilities,
their conditional versions, or their connected cumulants.
An all-volume event that every edge lies in \(\mathcal A_e\)
is not needed for (71.7), and cannot be presumed likely.
This distinction avoids an extensive-good-event assumption.

Equation (71.6) is also not permission to replace the compact
integral by V66's convex patch.  The latter is a subdomain,
and a gauge transformation generally moves its boundary.
A coarse Ward identity for the full measure cannot simply
be imposed on that one patch's normalization.
The local comparison (71.3) applies to gauge-invariant
observables; equality of the two coarse link-valued
pushforward measures has not been asserted.
\end{assessment}

\begin{decision}[Next test on the actual complement]
Relate failure of the relative-path logarithm condition
\(\mathcal A_e\) to a finite collection of plaquette
holonomies, with explicit thresholds, before attempting
probability bounds.  Then separate unconditional local
energy suppression from the conditional estimates needed
to iterate the blocked measure.  Existing archived
large-field estimates must be checked against this precise
event, not imported by name alone.
\end{decision}

\begin{firewall}[Position after V67]
The compact root-block extension, its exact positive
pushforward, and its local agreement with the dressed
block on gauge-invariant observables are proved.
These remove a domain-definition ambiguity but do not
prove continuum reconstruction, reflection positivity,
scale stability, or a physical Yang--Mills mass gap.
Only V67 remains active; preceding bodies and archives
are retained.  Next companion audit: V70.
\end{firewall}

% END APPENDED POST-TENTH-ARCHIVE V67 MATERIAL

% =============================================================================
\section{V68: the relative-path complement has a local action threshold}
\label{sec:v68-complement}
% =============================================================================

The V67 predecessor has SHA--256
\texttt{FD0E0326E858A2E2C731FFC4F519743A251ED8862A1D1E8022B9EA6E123A5468}.
Its body is retained except for the active-version label.
This section checks the actual event of V67, not an unnamed
large-field substitute.  It first proves a deterministic action
cost, then a separate unconditional probability bound.
Neither statement is a conditional influence estimate.

Use the tree specified in V66: its root paths traverse active
coordinates in increasing order.  In the faithful unitary
representation let \(\delta(V)=\|I-V\|_{\rm HS}\), and choose
\(\rho>0\) such that
\[
 \{V\in G:\delta(V)<\rho\}\subset\mathcal O.       \tag{72.1}
\]
Such a \(\rho\) exists by the neighbourhood property and
faithfulness.  Write \(s_p=\dim V-\operatorname{Re}\Tr U_p\);
then \(\delta(U_p)^2=2s_p\).
For \(e=(r,\mu)\), let \(\mathcal R_e\) be the rectangular
cell complex of side lengths \(3\) in direction \(\mu\) and
\(1\) in the other three directions, based at \(r\).
It has sixty elementary plaquettes.  Let
\(S_e=\sum_{p\subset\mathcal R_e}s_p\).
Use periodic lattices with at least two parity cells in
each direction, as in V66.

\begin{lemma}[Six square moves per transported path]
\label{lem:v68-filling}
Let \(s\in\{0,1\}^4\), \(x=r+s\).  The path underlying
\(P_x(W(U))\), from \(r\) to \(r+2\hat\mu\), differs from
the straight two-link path by at most
\(2\sum_{\nu\ne\mu}s_\nu\leq6\) elementary square moves
inside \(\mathcal R_e\).  Consequently
\[
 \delta(P_*^{-1}P_x)
       \leq12\max_{p\subset\mathcal R_e}\delta(U_p). \tag{72.2}
\]
\end{lemma}
\begin{proof}
In original link variables the path is
\(T_x\,(\mu\mu)\,T_{x+2\hat\mu}^{-1}\).
The positive word for \(T_x\) lists the nonzero coordinates
of \(s\) in increasing order; the target tree has the same
word translated by \(2\hat\mu\).  Move the two added
\(\mu\)-steps to the front of this word.  Each crosses
each transverse step at most once; identical \(\mu\)-steps
require no move.  The resulting path is the straight
\(\mu\mu\) followed by the target tree path, which cancels.
Each interchange of distinct adjacent positive steps is
an elementary square and stays inside the specified box.

Replacing one path segment around a square changes its
holonomy by a conjugate of a plaquette holonomy or inverse.
The Hilbert--Schmidt distance to \(I\) is invariant under
unitary conjugation and inversion, and satisfies
\(\delta(AB)\leq\delta(A)+\delta(B)\).
Compare both the reference and the chosen path with the
same straight path; there are at most twelve factors.
This proves (72.2), including paths with fewer moves.
\end{proof}

\begin{theorem}[Action cost of the precise logarithm failure event]
\label{thm:v68-cost}
For the event \(\mathcal A_e\) defined in V67,
\[
 W(U)\notin\mathcal A_e
       \ \Longrightarrow\ S_e\geq c_\rho,
       \qquad c_\rho=\rho^2/288.                 \tag{72.3}
\]
If \(N_{\rm bad}\) counts failed coarse edges, then
\[
 S_0(U)\geq\frac{\rho^2}{2304}N_{\rm bad}.        \tag{72.4}
\]
\end{theorem}
\begin{proof}
Failure means that some relative holonomy is not in
\(\mathcal O\); by (72.1) its distance is at least \(\rho\).
Equation (72.2) forces a plaquette distance at least
\(\rho/12\), and hence action at least \(\rho^2/288\).

For each orientation of a coarse edge, a fixed plaquette
can belong to at most two of its boxes: in the elongated
coordinate there are at most two even root positions,
and in each other coordinate the root is determined if
containment is possible.  Summing over the four orientations
gives multiplicity at most eight.  Thus
\(c_\rho N_{\rm bad}\leq\sum_{\rm bad\ e}S_e
\leq8S_0\), proving (72.4).
No independence of different failures has been used.
\end{proof}

\begin{theorem}[Uniform unconditional weak-coupling bound]
\label{thm:v68-probability}
For the normalized Wilson measure of V67, let
\(d_G=\dim G\).  There is a finite \(C_G\), independent
of lattice volume, such that for all \(\beta\geq1\)
\[
 \mathbb E_{\mu_\beta} S_0
       \leq N\,\frac{2d_G\log\beta+C_G}{\beta},   \tag{72.5}
\]
where \(N\) is the number of fine vertices.  In particular
\[
 \mu_\beta(W\notin\mathcal A_e)
 \leq\min\left(1,\frac{17280}{\rho^2}
                   \frac{2d_G\log\beta+C_G}{\beta}\right).
                                                        \tag{72.6}
\]
The bound is uniform in the chosen coarse edge and volume;
it is unconditional.
\end{theorem}
\begin{proof}
Normalized Haar volume of a sufficiently small
Hilbert--Schmidt ball about \(I\) is at least \(c r^{d_G}\)
for \(0<r\leq r_0\), for some \(c,r_0>0\).
This follows from the smooth positive Haar density in a
coordinate chart and equivalence of its coordinate norm
with the displayed distance near \(I\).  Decrease \(r_0\)
if necessary so \(c r_0^{d_G}\leq1\).
Restrict all \(4N\) independent links to the ball of radius
\(r=r_0/\sqrt{\beta}\).  Every plaquette has distance at
most \(4r\), so \(s_p\leq8r^2\).  There are \(6N\)
plaquettes.  Therefore
\[
 Z_\beta\geq
     e^{-48Nr_0^2}(c r_0^{d_G}\beta^{-d_G/2})^{4N},
 \qquad
 -\log Z_\beta\leq N(2d_G\log\beta+C_G),
\]
with \(C_G=48r_0^2-4\log(c r_0^{d_G})\).
Relative entropy of \(\mu_\beta\) with respect to normalized
product Haar is nonnegative and equals
\(-\beta\mathbb E S_0-\log Z_\beta\).
This proves (72.5).  No thermodynamic limit is needed.

Translation invariance makes the expectation of a plaquette
constant at fixed orientation.  For every such orientation
there are \(N\) plaquettes, with nonnegative action.  Hence
each individual expectation is at most
\(\mathbb E S_0/N\); this argument does not require equality
of the side lengths or symmetry between orientations.
The sixty plaquettes of \(\mathcal R_e\) give
\(\mathbb E S_e\leq60(2d_G\log\beta+C_G)/\beta\).
Apply (72.3) and Markov's inequality.  Since \(60\cdot288
=17280\), the result is (72.6).
\end{proof}

\begin{corollary}[Local observable insensitivity to the extension]
For the two extensions and bounded observable \(f\) of V67,
\[
 |\mathbb E f(K(W))-\mathbb E f(\widetilde K(W))|
 \leq \frac{34560\,|E_f|\,\|f\|_\infty}{\rho^2}
                   \frac{2d_G\log\beta+C_G}{\beta}. \tag{72.7}
\]
In particular the difference tends to zero as
\(\beta\to\infty\), uniformly in volume for fixed \(E_f\).
\end{corollary}
\begin{proof}
Insert (72.6) into (71.7), dropping the optional cap by one.
\end{proof}

\begin{assessment}[The remaining distinction is essential]
The action threshold and probability estimate are separate
arguments.  The former alone would not justify dividing by
the global partition function and declaring an exponential
probability bound.  The latter supplies only
\(O(\log\beta/\beta)\) for a fixed local event and no
uniform conditional estimate.

Equation (72.7) also does not establish continuum
extension-independence for observables whose lattice support
grows while the spacing goes to zero.  A loop of fixed
physical size generally uses an increasing number of
lattice edges; the factor \(|E_f|\) must then be retained.
The relation between bare coupling, spacing and physical
observables has not been populated by this estimate.
Neither local logarithm validity nor local strong convexity
controls the full interacting renormalization trajectory.
\end{assessment}

\begin{decision}[Next upstream bottleneck]
Test whether local modifications can yield a conditional
version of (72.6) under the retained coarse-field constraints.
The elementary Haar-ball lower bound changes all links
and is not such a conditional comparison.  Before using
an archived bad-polymer bound, identify a modification
that preserves the actual root-block value and controls
the neighbouring action.  If that cannot be done, retain
the unconditional result and do not feed it into a
conditional influence theorem.
\end{decision}

\begin{firewall}[Position after V68]
The actual relative-path failure event has a deterministic
local action cost and a volume-uniform unconditional
weak-coupling bound.  The two compact extensions agree
asymptotically on fixed-support bounded observables.
Conditional RG estimates, continuum reconstruction and
the full Yang--Mills mass gap remain open in these notes.
Only V68 remains active; prior bodies and archives are
retained.  Next companion audit and review: V70.
\end{firewall}

% END APPENDED POST-TENTH-ARCHIVE V68 MATERIAL

% =============================================================================
\section{V69: a forced conditional failure and the correct fibre comparison}
\label{sec:v69-conditional}
% =============================================================================

The V68 predecessor has SHA--256
\texttt{8660F7FF3FFA68D4F8DA0880D5D09EB1A4F1D4B9372E2BB244D9CB755CD1A4EF}.
Its body is retained apart from the active-version label.
The preceding turn made progress by proving a bound for
the actual unconditional complement.  We now test a proposed
boundary-only conditional repair, rather than promoting that
bound to an unspecified conditional estimate.

\begin{countertheorem}[Internal data can force logarithm failure]
\label{thm:v69-forced}
Take \(G=SU(2)\) in its defining representation and choose
the V67 logarithm neighbourhood sufficiently small that
\[
 -I\notin\overline{\mathcal O^{-1}\mathcal O}.      \tag{73.1}
\]
For a specified coarse edge \(e\), there is a nonempty open
set of non-tree internal link values in its source and
target cells such that
\[
 W\notin\mathcal A_e
 \quad\hbox{for every choice of all boundary links}. \tag{73.2}
\]
Consequently, at every finite \(\beta\), conditional on all
internal links, the probability of failure is one on a set
of conditioning values of positive marginal probability.
In particular the essential supremum of that conditional
failure probability is one.
\end{countertheorem}
\begin{proof}
Use coordinates relative to the source root, and set
\(\mu=0\), \(s=(0,1,0,0)\).
The internal edge in direction 0 starting at \(s\) is not
in the V66 tree: the parent of \((1,1,0,0)\) clears bit 1,
not bit 0.  Assign this edge the value \(-I\) and assign
all other internal links in the source and target cells
the value \(I\).  Tree links remain \(I\).

The two two-link paths starting at \(s\) and \(s+\hat0\)
share the boundary link \(B\) starting at \(s+\hat0\).
Their products are respectively
\[
 P_s=(-I)B=-B,\qquad P_{s+\hat0}=BI=B.
                                                        \tag{73.3}
\]
If both relative products \(P_*^{-1}P_s\) and
\(P_*^{-1}P_{s+\hat0}\) belonged to \(\mathcal O\), their
relative product would put
\(P_s^{-1}P_{s+\hat0}=-I\) in \(\mathcal O^{-1}\mathcal O\),
contradicting (73.1).  This argument works for every
reference path and every \(B\).

It also survives an open perturbation of the internal data.
Write the two pertinent internal values as \(J_s,J_t\).
The relative product is
\[
 (J_sB)^{-1}(BJ_t)=B^{-1}J_s^{-1}BJ_t.
 \]
As \(J_s\to-I\), \(J_t\to I\), this converges uniformly
in \(B\in SU(2)\) to \(-I\), by unitary invariance of
distance and the triangle inequality.  The closed set
in (73.1) has positive distance from \(-I\); sufficiently
small open neighbourhoods of the two internal values
therefore force (73.2).  Other internal values may vary
in open neighbourhoods of the specified ones.

In ordinary tree coordinates the finite-\(\beta\) Wilson
measure has a strictly positive continuous density with
respect to product Haar on all non-tree links.  Its
internal marginal assigns positive probability to this
open cylinder.  Conditional on these internal variables,
the boundary density is the normalized positive Wilson
weight; its entire support lies in the failure event.
Thus the conditional probability is one, not just at
a measure-zero exceptional configuration.
\end{proof}

\begin{assessment}[Scope of the counterexample]
This disproves a uniform-smallness claim when conditioning
on all internal links, and rules out curing every failure
by changing boundary links alone.  It does not disprove
smallness conditioned only on the retained coarse block:
that conditioning permits internal variables to fluctuate.
Nor does it contradict V68, because the forcing internal
data themselves have a cost and may be rare as \(\beta\)
increases.  A repair capable of treating this example
must change internal variables and account for its effect
on the retained block and nearby plaquettes.
\end{assessment}

\begin{theorem}[Exact conditioning on the measurable compact block]
\label{thm:v69-disintegration}
Let \(\lambda\) be normalized product Haar on the non-tree
variables, \(C=K(W)\), and let \(\lambda_c\) be a regular
conditional probability for \(\lambda\) given \(C=c\).
Write \(\nu_0=K_*\lambda\).
For \(\nu_0\)-almost every \(c\), set
\[
 Z_\beta(c)=\int e^{-\beta S_0(W)}\,d\lambda_c(W).
 \]
Then \(Z_\beta(c)>0\), and the conditional Wilson law is
\[
 d\mu_{\beta,c}(W)=
        Z_\beta(c)^{-1}e^{-\beta S_0(W)}d\lambda_c(W),
 \quad
 \frac{d\nu_\beta}{d\nu_0}(c)=\frac{Z_\beta(c)}{Z_\beta}.
                                                        \tag{73.4}
\]
These assertions require no global smoothness or
submersion property of \(K\).
\end{theorem}
\begin{proof}
The finite products of compact Lie groups are standard
Borel spaces, and \(K\) is measurable, so regular conditional
probabilities exist.  The Wilson action is bounded and
nonnegative at finite volume, giving
\(e^{-\beta\sup S_0}\leq Z_\beta(c)\leq1\).
Apply the defining conditional-integral identity first
to \(e^{-\beta S_0}\), and then to its product with bounded
test functions of \(W\) and \(c\).  The resulting integral
factorization gives both identities (73.4), including
the normalization of each conditional law.
\end{proof}

\begin{proposition}[Relative energy and fibre volume are the missing inputs]
\label{prop:v69-comparison}
In the preceding conditional law let
\[
 m(c)=\mathop{\rm ess\,inf}_{\lambda_c} S_0,\quad
 v_c(\delta)=\lambda_c\{S_0\leq m(c)+\delta\},
 \qquad \delta>0 .
 \]
Then \(v_c(\delta)>0\).  If an event \(B\) satisfies
\(S_0\geq m(c)+\Delta\) \(\lambda_c\)-almost everywhere
on \(B\), with \(\Delta>\delta\), then
\[
 \mu_{\beta,c}(B)
       \leq v_c(\delta)^{-1}
                        e^{-\beta(\Delta-\delta)}. \tag{73.5}
\]
One also has the conditional entropy estimate
\[
 \mathbb E_{\mu_{\beta,c}}S_0-m(c)
 \leq-\frac1\beta
          \log\int e^{-\beta(S_0-m(c))}\,d\lambda_c,
 \qquad \beta>0.                                 \tag{73.6}
\]
\end{proposition}
\begin{proof}
The essential-infimum definition gives positive measure
to every indicated sublevel set.  The denominator of
(73.4) is at least
\(v_c(\delta)e^{-\beta(m(c)+\delta)}\).
The numerator restricted to \(B\) is at most
\(e^{-\beta(m(c)+\Delta)}\lambda_c(B)\), and
\(\lambda_c(B)\leq1\).  Division proves (73.5).
For (73.6), compute the relative entropy of
\(\mu_{\beta,c}\) with respect to the probability
\(\lambda_c\), after subtracting the constant \(m(c)\)
from the action, and use its nonnegativity.
\end{proof}

\begin{assessment}[Why the absolute local cost is insufficient]
The cost \(S_e\geq c_\rho\) in V68 is an absolute local
statement, not a lower bound on
\(S_0-m(c)\) on the conditional bad event.
It supplies neither \(\Delta\) nor a uniform lower bound
on \(v_c(\delta)\) in (73.5).
The Haar-ball denominator proof of V68 changes the coarse
block and therefore cannot provide that fibre volume.
The same comparison applies if one conditions jointly
on the coarse block and chosen exterior variables;
those conditional fibres still need their own relative
energy and volume estimates.
The forced example above shows why a conditioning
class must be specified before seeking uniform constants.
\end{assessment}

\begin{decision}[Repair problem to solve next]
Construct a finite-region internal-and-boundary modification
that removes a relative-path failure while preserving
the root-block values affected by that region.  Its domain
must exclude, or explicitly price, incompatible exterior
constraints.  Quantify the change of Wilson action and
the conditional fibre measure, rather than assuming
ambient Haar preservation is enough.
The compulsory V70 companion audit and ten-version review
should assess whether the accumulated local results now
support such a modification or reveal a different route.
\end{decision}

\begin{firewall}[Position after V69]
An actual positive-measure obstruction to boundary-only
conditional repair is proved, and the exact compact
conditional law identifies the missing relative-energy
and fibre-volume inputs.  No uniform conditional
large-field estimate, continuum theory, or Yang--Mills
mass gap has been proved.  Only V69 remains active;
earlier bodies and archives are retained.
\end{firewall}

% END APPENDED POST-TENTH-ARCHIVE V69 MATERIAL

% =============================================================================
\section{V70: global good section, exact action, and ten-version audit}
\label{sec:v70-review}
% =============================================================================

The V69 predecessor has SHA--256
\texttt{AD8914B424E11BB6849889B0FE9CB250766EB63E5CF8C69F10270A3B4D580557}.
Its body is retained except for the version label.
This is the compulsory five-version companion audit and
ten-version review.  The new calculation extends the local
section of V61 to the compact block of V67 and computes
its exact action; it does not repeat the local section
identity as though that identity were new.

\begin{theorem}[Global all-good section and its Wilson action]
\label{thm:v70-section}
For any compact coarse link field \(C\), define \(J(C)\)
on the tree slice by setting every internal link to \(I\)
and all eight boundary links of each coarse edge \(e\)
to \(C_e\).  Then
\[
 K(J(C))=C,\qquad J(C)\in\bigcap_e\mathcal A_e,
 \qquad S_0(J(C))=4S_{\rm coarse}(C),             \tag{74.1}
\]
where \(S_{\rm coarse}\) is the Wilson action on the
coarse lattice in the same representation and normalization.
The block \(K\) is smooth and a submersion in an open
neighbourhood of this compact section.
\end{theorem}
\begin{proof}
Each two-link path has one identity internal link and
one boundary link \(C_e\), so all its products equal
\(C_e\).  All relative products are \(I\); the first
branch of (71.2) applies and gives \(K_e=C_e\).
This is valid for arbitrarily large compact \(C_e\),
not just its logarithmic neighbourhood.

Classify fine plaquettes by the parity of their base
point in their two directions.  If neither is odd,
all links are internal.  If exactly one is odd, its
two boundary links are equal and cancel, with identity
internal links between them.  If both are odd, its
product is the coarse plaquette product in the same
orientation.  There are four such fine plaquettes
for each coarse plaquette, from the two independent
transverse parity bits.  This proves the action identity.

The relative products stay in \(\mathcal O\) on an
open neighbourhood of the section, so \(K\) is smooth
there.  Differentiating \(K\circ J=\mathrm{id}\) proves
that \(DK\) is onto at every point of the section.
Surjectivity is open.  Compactness of the coarse
configuration space supplies an open neighbourhood
of the whole section on which the first branch is
smooth and \(DK\) remains onto.
\end{proof}

\begin{corollary}[A repair as a map, not as a measure comparison]
The measurable map \(W\mapsto J(K(W))\) preserves the
coarse block and eliminates every relative-path failure.
It is not a fibre diffeomorphism or a Haar-preserving
change of variables: on each fixed coarse fibre it is
constant.  Its action need not be smaller than the
original action, since (74.1) alone compares neither
action with the other.
\end{corollary}
\begin{proof}
Apply (74.1) with \(C=K(W)\).  Constancy on the fibre
shows the loss of its nontrivial coordinates, precluding
an invertible local change of variables there.
\end{proof}

\begin{theorem}[A positive smooth component of the coarse measure]
\label{thm:v70-component}
At fixed finite volume and finite \(\beta\geq0\), there
is a smooth nonnegative cutoff \(\chi\), supported in
the regular all-good neighbourhood just constructed
and equal to one near \(J(G^{E_{\rm coarse}})\), such
that
\[
 K_*(\chi e^{-\beta S_0}\lambda)=h_\beta(C)\,dC,
 \qquad h_\beta(C)>0,                            \tag{74.2}
\]
where \(dC\) is normalized coarse product Haar and
\(h_\beta\) is smooth.  Consequently the full coarse
Wilson probability dominates
\(Z_\beta^{-1}h_\beta(C)dC\).
Its lower density bound at fixed volume is positive,
but no volume-uniform bound is claimed.
\end{theorem}
\begin{proof}
Choose nested open neighbourhoods of the compact section,
with closures inside the regular all-good set, and a
smooth cutoff equal to one on the smaller one.
On the compact support of this cutoff, submersion
coordinates take the form \((C,z)\) with \(K(C,z)=C\).
Smooth Haar volume has a strictly positive coordinate
density.  A finite partition of unity subordinate to
such charts expresses the pushforward density as a
finite sum of integrals in the \(z\) variables of
smooth compactly supported functions.  Differentiation
under these finite integrals proves smoothness.
For every \(C\), a chart through \(J(C)\) has a
nonempty open \(z\)-region with \(\chi=1\), positive
Haar density, and positive Wilson weight.  Its integral
is strictly positive.  Compactness then gives a
positive minimum for this fixed-volume function.
Finally \(0\leq\chi\leq1\) implies domination by the
full pushforward.
\end{proof}

\begin{assessment}[The fibre-volume problem is still quantitative]
The section supplies a genuine regular neighbourhood
and positive measure, not merely a single isolated
competitor.  But (74.2) supplies no useful scaling of
its density with the number of cells or \(\beta\).
Nor does it bound the remaining part of the pushforward
or exclude a singular component away from this tube.
Accordingly it is not a uniform lower bound on the
conditional volumes \(v_c(\delta)\) of V69.
Replacing all configurations by \(J(C)\) would throw
away precisely the fluctuations whose measure must
be controlled.  A local comparison must retain them
or price their multiplicity.
\end{assessment}

\subsection{Compulsory companion audit}

\begin{audit}[Observed versions and transferable scope]
SM is now V72, SHA--256
\texttt{F44F9745D43379E1672C5550411B58E97195A4C9278592D221DF5D3F644DC1F5}.
Its endpoint proves a fixed-time, fixed-mode many-copy
limit to a momentum-resolved covariance flow and
vanishing empirical mixed-parent residual for
concentrated initial data.  Its declared next problem
is cutoff-uniform control of the pairing kernel and
source.  It does not supply a YM conditional fibre
volume or a gauge-field large-deviation estimate.
The auxiliary many-copy scaling cannot be substituted
for the present regulator limit.

NS remains V26, SHA--256
\texttt{82C887F8C2C69E4F28FAD7C3DC8DE5B9099CB5A4BF431EBA1FFF3E91935978AD}.
Its endpoint retains global regularity as unproved
after the pointwise rank-one Oseen norm calculation.
Parallel YM remains V5 and retains the extensive-charge
tightness obstruction.  Its fixed-window warning is
consistent with V68's explicit support-size factor,
but supplies no missing conditional estimate.
No companion theorem changes the hypotheses of
(74.1)--(74.2).
\end{audit}

\subsection{Ten-version review and direction}

\begin{record}[What V61--V70 actually changed]
V61--V63 established section and reference identities,
a boundary density test, and finite-range Jacobian
factorization.  V64--V65 supplied the compatible tree
reduction and full mixed quadratic density.
V66 proved local nonlinear convexity for that specific
averaged fibre, not just an inherited Gaussian model.
V67 defined a compact measurable block and its positive
measure with an observable-level local comparison.
V68 controlled its precise logarithm failure event
unconditionally.  V69 disproved uniform repair with
all internal links frozen and identified the necessary
relative-energy and fibre-volume estimates.
V70 adds the global section's exact action and a
regular positive-measure neighbourhood.

None of these results proves the coupled coarse
effective action closes under repeated blocking.
The current obstacle is not another uncomputed
quadratic density term.  It is a quantitative
comparison of compact conditional fibres, together
with control of the interactions generated by their
integration.  Repeating local Taylor calculations
without addressing this would not close the gap.
\end{record}

\begin{decision}[Next acquisition]
Quantify a neighbourhood of the global section in
group coordinates while retaining a prescribed coarse
field.  First test whether its finite-volume fibre
volume admits an explicit per-cell bound and whether
the Wilson action on that neighbourhood is controlled
by \(4S_{\rm coarse}(C)\) plus a fluctuation cost.
Keep any dependence on coarse curvature and boundary
conditions visible.  Even a successful lower bound
must subsequently be paired with the bad-fibre
numerator and connected interaction estimates.
Next companion audit: V75; next ten-version review: V80.
\end{decision}

\begin{firewall}[Position after V70]
The exact compact section has action \(4S_{\rm coarse}\)
and a regular neighbourhood contributing positive
smooth coarse density at each finite volume.
Quantitative conditional comparison, scale stability,
continuum reconstruction and the physical mass gap
remain unproved.  Only V70 is active; prior bodies
and all archives are retained.  The V100 rollover
policy remains unchanged.
\end{firewall}

% END APPENDED POST-TENTH-ARCHIVE V70 MATERIAL

% =============================================================================
\section{V71: quantitative compact section tube and coarse-density lower bound}
\label{sec:v71-tube}
% =============================================================================

The V70 predecessor has SHA--256
\texttt{00ABD1D0E5271A9AA3EB39C3B53A794F817FACF417CCD47F2ADB95713AE69BB4}.
Its body is retained apart from the active-version label.
V70 proved finite-volume positivity without estimating its
volume dependence.  Here that dependence is made explicit.
Let \(n\) be the number of coarse cells and \(d=\dim G\).
There are \(49n\) non-tree fine links and \(4n\) coarse links.
All constants below may depend on the fixed group, logarithm
neighbourhood and tree, but not on \(n\) or the compact
coarse field \(C\).

\begin{theorem}[Uniform pivot chart about the global section]
\label{thm:v71-chart}
There are \(r_0>0\), \(L<\infty\), and \(a>0\) with the
following properties.  For every coarse field \(C\) and
every collection \(z\) of \(45n\) Lie-algebra vectors with
\(\max_j\|z_j\|<r_0\), a smooth injective parametrization
\(\Phi(C,z)\) satisfies
\[
 K(\Phi(C,z))=C,\quad \Phi(C,0)=J(C),\quad
 \Phi(C,z)\in\bigcap_e\mathcal A_e,\quad
 \|\Phi(C,z)_l-J(C)_l\|_{\rm HS}\leq L\max_j\|z_j\|.
                                                        \tag{75.1}
\]
If \(\lambda\) is product Haar on the non-tree links,
then on this tube its coordinate density obeys
\[
 \Phi^*\lambda=q(C,z)\,dC\,dz,\qquad q(C,z)\geq a^n,
                                                        \tag{75.2}
\]
where \(dz\) is product Euclidean volume for the declared
orthonormal Lie-algebra coordinates.
\end{theorem}
\begin{proof}
Choose one pivot among the eight boundary links for each
coarse edge.  The 17 non-tree internal links per cell are
\(\exp z_l\); the other seven boundary links per coarse
edge are \(C_e\exp z_l\).  Write its pivot as \(C_e\exp v_e\).
There are \(17n+28n=45n\) free variables.

For fixed internal and free boundary variables, \(K_e\)
depends on its own pivot but on no other coarse edge's
pivot.  At \(z=v=0\) all two-link products are \(C_e\).
A pivot increment \(v_e\) appears in two of the sixteen
paths.  Differentiating the group mean in right
coordinates therefore gives
\[
 D_{v_e}\log\big(C_e^{-1}K_e\big)=\tfrac18 I.
                                                        \tag{75.3}
\]
This also holds when the reference path contains the pivot:
the reference derivative and the relative-log derivatives
sum to the same average.

All functions in this equation involve a fixed finite
stencil and compact coarse parameters.  Their derivatives
are uniformly continuous on a sufficiently small common
neighbourhood of zero.  Multiplying the equation by eight
gives a uniformly invertible derivative at zero.
For completeness one may solve it by
\(v_e\mapsto v_e-8\log(C_e^{-1}K_e)\):
shrink the common neighbourhood until its derivative in
\(v_e\) has norm at most \(1/2\); its value at \(v_e=0\)
is bounded by a constant times the local free-variable
radius.  A ball with that enlarged radius is invariant
under the map.  Contraction gives a unique \(v_e\), with
\(\|v_e\|\leq L_1\max_j\|z_j\|\).
The constants are local and independent of volume.
Shrinking \(r_0\) keeps all relative products in
\(\mathcal O\) and all exponential coordinates injective.
The solution is smooth by the implicit-function theorem.
The link bound follows from bounded exponential derivatives.

Injectivity follows by recovering \(C=K(W)\), then the
internal \(z_l=\log W_l\) and free boundary
\(z_l=\log(C_e^{-1}W_l)\).  To check the measure bound,
first hold the internal and free boundary group links fixed
and replace the pivot group links by their coarse outputs.
The corresponding derivative is block diagonal in \(e\).
Each group-volume Jacobian equals \(8^{-d}\) at the
section and stays between two positive constants after
the same uniform shrinking.  Its inverse thus has a
positive lower bound per coarse edge.
Next write internal links as \(\exp z_l\) and free boundary
links as \(C_e\exp z_l\).  Conditional on \(C\), their
Haar densities are bounded below by a fixed positive
constant per link.  Dependence on \(C\) does not change
this Jacobian: \(C\) is unchanged in this second coordinate
change, making its derivative triangular.
There are \(4n\) pivot Jacobian factors and \(45n\)
free-link factors.  Their product is at least \(a^n\).
This proves (75.2) without a growing-matrix determinant.
\end{proof}

\begin{theorem}[Action and density with explicit volume dependence]
\label{thm:v71-density}
For every \(\eta>0\), put
\(A_\eta=768L^2(1+\eta^{-1})\).
For \(0<r\leq r_0\) and \(\max_j\|z_j\|<r\),
\[
 S_0(\Phi(C,z))
 \leq4(1+\eta)S_{\rm coarse}(C)+A_\eta n r^2.     \tag{75.4}
\]
There is \(b>0\), independent of \(C,n,r,\beta\), such that
as measures on coarse configurations,
\[
 K_*(e^{-\beta S_0}\lambda)(dC)
 \ \geq\
 \left(b r^{45d}e^{-\beta A_\eta r^2}\right)^n
 e^{-4(1+\eta)\beta S_{\rm coarse}(C)}\,dC .
                                                        \tag{75.5}
\]
In particular, for \(\beta\geq1\) one may choose
\(r=r_0/\sqrt\beta\), obtaining
\[
 K_*(e^{-\beta S_0}\lambda)(dC)
 \ \geq\
 b_\eta^n\beta^{-45dn/2}
 e^{-4(1+\eta)\beta S_{\rm coarse}(C)}\,dC,       \tag{75.6}
\]
where \(b_\eta=b r_0^{45d}e^{-A_\eta r_0^2}>0\).
\end{theorem}
\begin{proof}
Telescoping a product of four unitary links bounds
the change in a plaquette holonomy by \(4Lr\).
For its Hilbert--Schmidt distance to \(I\), use
\((x+y)^2\leq(1+\eta)x^2+(1+\eta^{-1})y^2\)
and \(s_p=\delta(U_p)^2/2\).
There are \(96n\) fine plaquettes.
The reference action is \(4S_{\rm coarse}(C)\)
by V70.  Summation gives (75.4), with
\(96\cdot(4L)^2/2=768L^2\).

Restrict the nonnegative measure to the injective tube
with each of its \(45n\) variables in a Euclidean
radius-\(r\) ball.  Its coordinate volume is
\((v_d r^d)^{45n}\), where \(v_d\) is unit-ball volume.
Use (75.2), (75.4), and integrate in \(z\).  This proves
(75.5) with \(b=a v_d^{45}\), by testing against arbitrary
nonnegative coarse functions.  It is a measure domination
statement and does not assert the rest of the pushforward
has a density.  Substitution gives (75.6).
\end{proof}

\begin{corollary}[A finite-volume partition-function comparison]
With \(Z_{\rm coarse}(\gamma)=
\int e^{-\gamma S_{\rm coarse}(C)}dC\),
\[
 Z_\beta\geq
 b_\eta^n\beta^{-45dn/2}
          Z_{\rm coarse}(4(1+\eta)\beta),
 \qquad\beta\geq1.                              \tag{75.7}
\]
\end{corollary}
\begin{proof}
Integrate (75.6) over \(C\) and use V67's exact tree
Haar integration.  All normalizations are product Haar.
\end{proof}

\begin{assessment}[What the bound does not cancel]
This proves an explicit per-cell lower contribution,
uniform over the compact coarse field.  It is stronger
than fixed-volume positivity, but is not a local
conditional probability estimate.  The factor
\(b_\eta^n\beta^{-45dn/2}\) is extensive.
Dividing an unmatched global bad-event numerator by
this lower bound would generally lose volume uniformity.
Nor does a lower comparison with coarse coupling
\(4(1+\eta)\beta\) identify the effective coupling or
prove an RG recursion.  In particular it cannot
determine the calibrated marginal flow.

The comparison uses a full-volume tube, not a
block-preserving local repair with exterior links fixed.
Such a repair would have to cancel unchanged exterior
integrations and deal with the forced internal
configurations of V69.  The present proof avoids
declaring its section to be an action minimizer.
\end{assessment}

\begin{decision}[Next numerator problem]
Seek a matched conditional comparison in a finite region,
using the local pivot solve only where its hypotheses
hold and retaining the same exterior integrations on
both sides.  Determine whether a relative-path failure
has excess action above the relevant constrained local
minimum; V68's absolute threshold alone does not answer
this.  If the exterior forces failure, isolate that
exterior event instead of asserting a uniform repair.
\end{decision}

\begin{firewall}[Position after V71]
A coarse-field-uniform pivot tube and explicit per-cell
lower bound for its weighted coarse contribution are
proved.  The unmatched numerator, generated interactions,
continuum limit and full Yang--Mills mass gap remain
unproved.  Only V71 is active; prior bodies and archives
are retained.  Next companion audit: V75.
\end{firewall}

% END APPENDED POST-TENTH-ARCHIVE V71 MATERIAL

% =============================================================================
\section{V72: a local block-preserving repair tube with fixed exterior}
\label{sec:v72-local-repair}
% =============================================================================

The V71 predecessor has SHA--256
\texttt{3217B1F0C88124B99C81F1924DC3749EEF6AB12EF0A550CD29B07DB527AA0A05}.
Its body is retained except for the active-version label.
We localize the previous full-volume tube.  The resulting
comparison has a fixed-dimensional denominator, but does
not assume a bound on the remaining conditional numerator.

Fix a cell \(r\).  Its star \(E_r\) consists of its four
outgoing and four incoming coarse edges.  Let \(\Lambda_r\)
contain its seventeen non-tree internal links and all
eight boundary links of each edge in \(E_r\).
Thus \(\Lambda_r\) has \(17+8\cdot8=81\) group links.
Hold every other non-tree link fixed at \(\xi\).
Assume the internal links in the neighbouring cells
incident to \(E_r\) have logarithms of norm at most
\(\varepsilon\).  No smallness of exterior boundary
links is assumed.  Let \(c=(c_e)_{e\in E_r}\) denote
arbitrary compact coarse values on the star.

\begin{theorem}[Local repair with 73 free variables]
\label{thm:v72-repair}
For sufficiently small fixed \(\varepsilon>0\), there is
a radius \(r_1>0\), uniform in \(c,\xi\) and total volume,
and a smooth injective tube
\[
 V=\Phi_{r,\xi}(c,z),\qquad
       z\in(B_{r_1}^{\mathfrak g})^{73},
                                                        \tag{76.1}
\]
in the variables of \(\Lambda_r\), such that all eight
star blocks equal \(c\), all their logarithm conditions
\(\mathcal A_e\) hold, and every exterior link is unchanged.
No block outside \(E_r\) is changed.
With product Haar \(dV\) on \(\Lambda_r\),
\[
 \Phi_{r,\xi}^*dV=q_{r,\xi}(c,z)\,dc\,dz,
       \qquad q_{r,\xi}(c,z)\geq a_*>0,           \tag{76.2}
\]
where \(a_*\) is independent of total volume and allowed
exterior data.  The central representative is
\(V_0(c,\xi)=\Phi_{r,\xi}(c,0)\).
\end{theorem}
\begin{proof}
Set the central internal links to \(\exp z_l\).
For each of the eight star edges, set seven boundary
links to \(c_e\exp z_l\), and solve its pivot link as
\(c_e\exp v_e\).  The free count is \(17+8\cdot7=73\).
At zero central and neighbouring internal fields, and
zero free boundary fields, the pivot derivative is
\(I/8\), uniformly in \(c_e\), by (75.3).
The exterior internal fields are small parameters.
The same local contraction gives
\(\|v_e\|\leq C(\varepsilon+\max_l\|z_l\|)\).
Choose \(\varepsilon,r_1\) so the contraction ball stays
in the all-good relative-logarithm neighbourhood.
The finite stencils and compact parameters make all
constants uniform.  Exterior boundary fields do not
enter the eight block equations.

Each central internal link affects only a block whose
source or target is the central cell.  Each changed
boundary link belongs to exactly one star edge.
Therefore no non-star block changes.
The recovery of \(c\) from the block and \(z\) from
the free links proves injectivity as in V71.
Eight uniformly bounded inverse pivot Jacobians and
73 positive free-link Haar densities give (76.2).
The smooth dependence on the small exterior internal
parameters also follows from the local implicit solve.
\end{proof}

\begin{theorem}[Matched local denominator with no volume factor]
\label{thm:v72-denominator}
Let \(\mathcal P_r\) be the plaquettes touching
\(\Lambda_r\), and put
\[
 H_\xi(V)=\sum_{p\in\mathcal P_r}s_p(V,\xi),
       \qquad H_0(c,\xi)=H_\xi(V_0(c,\xi)).
 \]
There are constants \(b_*>0\), \(A_\eta<\infty\),
independent of total volume and allowed exterior data,
such that, for \(0<t\leq r_1\) and \(\eta>0\),
\[
 (K_{E_r})_*(e^{-\beta H_\xi}dV)(dc)
 \geq b_*t^{73d}
       e^{-\beta((1+\eta)H_0(c,\xi)+A_\eta t^2)}\,dc.
                                                        \tag{76.3}
\]
All exterior-only action terms cancel from the
conditional normalization for this region.
\end{theorem}
\begin{proof}
Uniform differentiation of the pivot solve gives
\(\|V_l(c,z)-V_l(c,0)\|_{\rm HS}\leq L_*\max_j\|z_j\|\).
An edge belongs to six elementary plaquettes in four
dimensions, so \(|\mathcal P_r|\leq486\).
The plaquette telescoping and squared-distance inequality
of V71 give
\[
 H_\xi(\Phi_{r,\xi}(c,z))
 \leq(1+\eta)H_0(c,\xi)
             +3888L_*^2(1+\eta^{-1})t^2
 \]
when all \(\|z_j\|<t\).  This includes plaquettes containing
arbitrary fixed exterior links, since unitarity suffices
for the difference bound.  Integrate over the 73 balls
using (76.2).  This proves (76.3), with
\(b_*=a_*v_d^{73}\) and the displayed choice of \(A_\eta\).
Finally \(S_0(V,\xi)=H_\xi(V)+S_{\rm ext}(\xi)\);
the last factor is constant and cancels on conditioning.
\end{proof}

\begin{proposition}[Precisely the remaining numerator criterion]
\label{prop:v72-numerator}
Let \(\nu_{0,\xi}=(K_{E_r})_*dV\), and write \(g_\xi(c)\)
for the density of its absolutely continuous part with
respect to \(dc\).  Suppose an event \(B\) has
\(H_\xi(V)\geq E(c,\xi)\) almost everywhere on its
conditional fibre at \(c\).  Then for Haar-almost every
\(c\), on the absolutely continuous part of the
weighted coarse law,
\[
 \mathbb P_{\beta,\xi}(B\mid c)
 \leq
 \frac{g_\xi(c)}{b_*t^{73d}}\,
 \exp\{-\beta[
    E(c,\xi)-(1+\eta)H_0(c,\xi)-A_\eta t^2]\}.    \tag{76.4}
\]
No claim is made on a possible singular coarse component.
\end{proposition}
\begin{proof}
The weighted bad pushforward is dominated by
\(e^{-\beta E(c,\xi)}\nu_{0,\xi}(dc)\).
Its absolutely continuous density is therefore at most
\(e^{-\beta E(c,\xi)}g_\xi(c)\).
The total weighted density is bounded below by (76.3).
Their ratio is the conditional probability on this
component.  A singular part cannot be bounded by
replacing it with a Haar density, which is why the
component qualification is retained.
\end{proof}

\begin{corollary}[The actual failure event gives an explicit candidate energy]
For a selected star edge \(e\), define
\[
 S_{e,\rm ext}(\xi)=
       \sum_{\substack{p\subset\mathcal R_e\\p\notin\mathcal P_r}}
                       s_p(\xi).
 \]
The event \(B=\{W\notin\mathcal A_e\}\) satisfies the
preceding proposition with
\[
 E(c,\xi)=\max(0,c_\rho-S_{e,\rm ext}(\xi)).       \tag{76.5}
\]
\end{corollary}
\begin{proof}
All plaquettes in the displayed exterior sum are fixed.
The other nonnegative terms of \(S_e\) are among the
terms of \(H_\xi\).  Hence \(S_e\leq
H_\xi+S_{e,\rm ext}\).  Apply (72.3) and \(H_\xi\geq0\).
\end{proof}

\begin{assessment}[What was removed, and what was not]
The denominator no longer contains a factor growing
with total volume: the repair changes 81 group links
and integrates 73 free variables.  It changes internal
links, so it is not contradicted by V69's boundary-only
obstruction.  It preserves the exterior and the entire
coarse block of any initial configuration when \(c\)
is chosen to be that configuration's star block.

However (76.4) is not yet a uniform suppression theorem.
It requires control of \(g_\xi(c)\), possible singular
coarse mass, and a positive energy difference relative
to \(H_0(c,\xi)\).  With arbitrary exterior boundary
fields the latter need not follow from (76.5).
The small-internal collar is a genuine hypothesis;
its probability or conditional cost has not been bounded
here.  These restrictions must survive any use of the
repair in a polymer estimate.
\end{assessment}

\begin{decision}[Next geometric obstruction to inspect]
Study the full local coarse pushforward, not just its
regular repair tube: determine whether the fallback
and relative-mean branches admit a uniform density
bound under the small-internal collar, including
configurations with large central internal links.
If critical fibres obstruct this, stratify them or
change the comparison variables with an explicit
observable-level justification.  Do not assume the
local implicit-function theorem covers the full fibre.
\end{decision}

\begin{firewall}[Position after V72]
A finite-region, exterior-fixed, block-preserving
all-good tube and its matched local denominator
are proved.  The conditional numerator and generated
interaction estimates remain incomplete; no continuum
Yang--Mills mass gap has been established.
Only V72 is active; earlier bodies and archives are
retained.  Next companion audit: V75.
\end{firewall}

% END APPENDED POST-TENTH-ARCHIVE V72 MATERIAL

% =============================================================================
\section{V73: bounded full local coarse density under a small internal collar}
\label{sec:v73-density}
% =============================================================================

The V72 predecessor has SHA--256
\texttt{C44E9DBA1DD07CE3DAC3ED992EE520D68580718208AC495C063475664F6D17C3}.
Its body is retained except for the version label.
We now control the full local Haar pushforward, not just
the regular replacement tube.  For this theorem choose
the previously free V67 logarithm neighbourhood to be a
sufficiently small invariant exponential ball.  All
earlier statements apply to this choice.  The result
is not asserted for an arbitrary larger logarithm domain.
Keep the V72 star and small exterior internal collar.
For each star edge choose as pivot the boundary link
contained in its fixed reference path.

\begin{lemma}[A good edge forces its participating internal links small]
\label{lem:v73-internal}
Write a shared-boundary pair of path products as
\[
 P_0=J_sB,\qquad P_1=BJ_t.
 \]
If both relative products \(P_*^{-1}P_0\) and
\(P_*^{-1}P_1\) belong to \(\mathcal O\), and
\(\sup_{V\in\mathcal O}\delta(V)\leq q\), then
\[
 \delta(J_s)\leq\delta(J_t)+2q,\qquad
 \delta(J_t)\leq\delta(J_s)+2q.                  \tag{77.1}
\]
For a star edge one of these internal links is exterior.
Thus on its good branch every participating central
internal link is uniformly close to \(I\) when the
collar and \(\mathcal O\) are small.
\end{lemma}
\begin{proof}
The distance between \(P_0,P_1\) is at most \(2q\).
Multiplying by \(B^{-1}\) on the appropriate sides
gives distance at most \(2q\) between \(J_s\) and
\(BJ_tB^{-1}\).  The two inequalities follow from
the triangle inequality and conjugation invariance.
There are eight such pairs for each edge, covering
all its participating internal links.  One member
of each pair lies in the central cell and the other
in the exterior neighbouring cell.
\end{proof}

\begin{lemma}[Uniform one-pivot density bound]
\label{lem:v73-pivot}
Fix all internal links and the seven nonpivot boundary
links for one star edge.  If its exterior internal
links satisfy the small-collar hypothesis, the image
of normalized Haar on the remaining pivot \(B\)
under the full piecewise map \(B\mapsto K_e\) is
absolutely continuous with density at most a constant
\(M<\infty\).  This bound is uniform over the fixed
links, including unrestricted central internal links.
\end{lemma}
\begin{proof}
On the fallback set the output is \(P_*\), either
\(J_sB\) or \(BJ_t\) with fixed internal factors.
On the entire pivot group this is a Haar-preserving
translation.  Its restriction to any measurable set
therefore has pushforward bounded above by Haar.

It remains to treat the good subset.  If it is empty
there is nothing to prove.  Otherwise (77.1) makes
all participating internal links small.  Choose an
anchor \(A\) among the nonpivot path products; it is
independent of \(B\).  Existence of a good pivot
also makes all the fixed nonpivot path products
close to \(A\).  On the good subset the reference
has the unique coordinate
\[
 P_*(B)=A\exp u,\qquad \|u\|\leq Cq.
 \]
Translations from \(B\) to \(P_*\) preserve Haar.
The partner of the reference path contains the
same pivot.  Since both its internal factors are
small, for all \(u\) in a fixed small convex ball
its product is close to \(A\exp u\).
All the relative logarithms thus extend smoothly
to that full ball, provided \(q\) and the collar
are chosen small enough.  In making this extension
one may use a larger injective logarithm chart;
on the actual good set its logarithms agree with
those defining the block.

Consider in this ball the output coordinate
\[
 F(u)=\log(A^{-1}K_e(B(u))).
 \]
When all participating internal links are \(I\)
and the fourteen nonpivot path products are \(A\),
the two pivot products are \(A\exp u\).
Direct substitution in the reference mean gives
\(F(u)=u/8\), throughout a sufficiently small ball.
The full functions and their first derivatives
depend continuously on the finitely many parameters.
Compactness in \(A\), and smallness of the other
parameters established above, allow a common ball
and choices such that
\[
 \|DF(u)-I/8\|\leq1/16
 \]
on that ball.  This justification uses (77.1);
it would be false to invoke smallness of the
unrestricted central fields without it.

For two points in the convex ball, integrate the
derivative of \(F(u)-u/8\) along their segment.
It follows that
\(\|F(u)-F(v)\|\geq\|u-v\|/16\).
Thus the extended map is injective and has a
Jacobian determinant bounded away from zero.
The input and output Haar chart densities have
uniform positive lower and finite upper bounds.
Change of variables gives a uniform bound \(M_g\)
for the pushforward of the good subset, even if
that subset itself is not convex.
Adding the fallback bound gives \(M=1+M_g\).
No differentiability across the branch boundary
is needed.
\end{proof}

\begin{theorem}[The full star pushforward has bounded density]
\label{thm:v73-full}
For every allowed fixed exterior \(\xi\),
\[
 \nu_{0,\xi}=(K_{E_r})_*dV=g_\xi(c)\,dc,
       \qquad 0\leq g_\xi(c)\leq M^8
       \quad\hbox{almost everywhere}.           \tag{77.2}
\]
In particular there is no singular component.
The weighted pushforward with density
\(e^{-\beta H_\xi}\) is also absolutely continuous.
\end{theorem}
\begin{proof}
Condition product Haar first on all seventeen central
internal links and all 56 nonpivot boundary links.
The eight pivot variables are independent Haar.
The block for each edge depends on its own pivot,
not on the other seven.  The preceding lemma
therefore bounds their joint output density by
\(M^8\), uniformly in the conditioned variables.
Integrating those variables preserves the bound.
Since \(0<e^{-\beta H_\xi}\leq1\), its pushforward
is dominated by \(\nu_{0,\xi}\), proving the last
assertion.  The V72 tube also supplies a strictly
positive lower density for that weighted law.
\end{proof}

\begin{corollary}[A volume-uniform conditional bound with an explicit gap test]
For the event and energy threshold of V72,
\[
 \mathbb P_{\beta,\xi}(B\mid c)
 \leq \frac{M^8}{b_*t^{73d}}
 \exp\{-\beta[
 E(c,\xi)-(1+\eta)H_0(c,\xi)-A_\eta t^2]\},
                                                        \tag{77.3}
\]
for almost every \(c\), with no singular-component
exception.  All constants are independent of total
volume.  If, on a declared set of \(c,\xi\),
\[
 E(c,\xi)-(1+\eta)H_0(c,\xi)\geq\Delta>0,
 \]
then taking \(t=r_1/\sqrt\beta\), \(\beta\geq1\),
gives a bound of the form
\[
 C\,\beta^{73d/2}e^{-\beta\Delta},                \tag{77.4}
\]
with \(C\) independent of total volume.
\end{corollary}
\begin{proof}
Use (77.2) in (76.4) and the positive lower density
from (76.3).  With the indicated \(t\), the term
\(\beta A_\eta t^2=A_\eta r_1^2\) is constant
and is absorbed into \(C\).
\end{proof}

\begin{assessment}[The remaining hypothesis has not disappeared]
The full local density obstacle is removed under the
specified small-internal collar and sufficiently small
choice of logarithm neighbourhood.  Arbitrary exterior
boundary fields are still permitted in the density
theorem.  They can, however, make \(H_0(c,\xi)\) large;
the positive excess \(\Delta\) in (77.4) is not automatic.
Neither the collar nor its stability under conditioning
has been proved typical.  Thus (77.4) is a proved
conditional implication, not an unconditional polymer
or continuum estimate.
\end{assessment}

\begin{decision}[Populate the excess-action regime]
Specify a finite collar of exterior links and coarse
values for which \(H_0\) and \(S_{e,\rm ext}\) are
small enough to make (76.5) exceed \((1+\eta)H_0\).
Then determine how a failure of that collar can be
charged to neighbouring bad regions without assuming
independence.  The next task is this energy and collar
accounting, not further differentiation of the density.
\end{decision}

\begin{firewall}[Position after V73]
The full local coarse Haar pushforward has a uniform
bounded density under the small-internal collar.
The matched conditional estimate now has no unknown
density or singular-component factor.  The excess-action
and collar hypotheses still require control before RG
iteration.  No full Yang--Mills mass gap is proved.
Only V73 is active; prior bodies and archives remain.
Next companion audit: V75.
\end{firewall}

% END APPENDED POST-TENTH-ARCHIVE V73 MATERIAL

% =============================================================================
\section{V74: an explicit flat-collar regime for conditional suppression}
\label{sec:v74-flat-collar}
% =============================================================================

The V73 predecessor has SHA--256
\texttt{01BFC0F8BE2BED28AF437A58D38F0E73A2EAA47DE30EA264C5D46CEF1413DC65}.
Its body is retained except for the active-version label.
We populate the excess-action hypothesis of V73 on
a specified nonempty set of exterior and coarse data.
Controlling the complement of that set is a separate
remaining problem, not an assumption of this theorem.

Fix a star \(r\) and a selected edge \(e\in E_r\).
Let \(\mathcal F_{r,e}\) be the following finite set of
fixed exterior non-tree links: all such links in a
plaquette of \(\mathcal P_r\) or \(\mathcal R_e\), together
with all exterior internal links used by the eight
star block equations.  Tree links are already \(I\).
This is a fixed-size stencil independent of total volume.
For \(a>0\), define the small data condition
\[
 \delta(\xi_l)<a\ (l\in\mathcal F_{r,e}),\qquad
 \delta(c_f)<a\ (f\in E_r).                      \tag{78.1}
\]
Take \(a\) small enough to imply the internal logarithm
bound required in V72--V73, by local equivalence of the
group distance and logarithmic norm.

\begin{lemma}[Uniform action bounds for the local competitor]
\label{lem:v74-energy}
There is a constant \(L_0\geq1\), independent of volume,
such that under (78.1) the central representative
\(V_0(c,\xi)\) of V72 satisfies
\[
 H_0(c,\xi)\leq3888L_0^2a^2,\qquad
 S_{e,\rm ext}(\xi)\leq480a^2.                  \tag{78.2}
\]
\end{lemma}
\begin{proof}
Its central internal links are \(I\), its free boundary
links are \(c_f\), and its solved pivots have distance
from \(I\) at most \(L_0a\).  The last assertion follows
from the uniform pivot estimate at \(z=0\), with the
small exterior internal variables as parameters, and
the triangle inequality for \(c_f\exp v_f\).
Increase \(L_0\) to at least one to cover every link.

Each plaquette in \(\mathcal P_r\) therefore has distance
at most \(4L_0a\) from \(I\); its action is at most
\(8L_0^2a^2\).  Use \(|\mathcal P_r|\leq486\).
For the exterior-only sum, all its links satisfy
(78.1), and there are at most sixty plaquettes.
This gives \(60\cdot8a^2=480a^2\).
\end{proof}

\begin{theorem}[Exponential conditional suppression on the small collar]
\label{thm:v74-suppression}
Choose \(a>0\) within the preceding chart bounds and with
\[
 (480+7776L_0^2)a^2\leq c_\rho/2,\qquad
 c_\rho=\rho^2/288.                             \tag{78.3}
\]
Then for every allowed fixed exterior \(\xi\), and
almost every star coarse value \(c\) satisfying
(78.1), the full compact conditional law obeys
\[
 \mathbb P_{\beta,\xi}
       (W\notin\mathcal A_e\mid K_{E_r}=c)
 \leq Q_\beta:=
       \min(1,C_0\beta^{73d/2}e^{-\beta c_\rho/2}),
       \qquad \beta\geq1,                       \tag{78.4}
\]
where \(C_0\) is independent of total volume and the
particular small data.
\end{theorem}
\begin{proof}
Use \(\eta=1\) in V73.  Equations (76.5), (78.2)
and (78.3) give
\[
 E(c,\xi)-2H_0(c,\xi)
 \geq c_\rho-(480+7776L_0^2)a^2\geq c_\rho/2.
 \]
The maximum by zero in (76.5) does not weaken this
inequality, and (78.3) ensures the untruncated
threshold is positive.  Apply (77.3) with
\(t=r_1/\sqrt\beta\).  The factor
\(e^{A_1r_1^2}M^8/(b_*r_1^{73d})\) may be used
as \(C_0\).  This is the full conditional measure,
not a conditional measure restricted to the repair tube.
\end{proof}

\begin{theorem}[Gauge-saturated flat-collar formulation]
\label{thm:v74-gauge}
The same estimate holds when (78.1) is replaced by
the existence of a cell-constant gauge transformation
on the finite stencil that sends the exterior data
and star coarse links into (78.1).
Thus the criterion is proximity to a locally pure-gauge
field, not proximity to \(I\) in one fixed gauge.
\end{theorem}
\begin{proof}
A cell-constant gauge transformation preserves the tree
slice and the set of changed versus fixed links.
It acts by conjugation on internal links, by endpoint
translations on boundary links, and by the ordinary
endpoint rule on coarse links.  Product Haar and
Wilson action are invariant; \(\mathcal A_e\) is
invariant by V67.  Consequently the conditional law
is carried to the conditional law for the transformed
exterior and coarse data, and (78.4) has unchanged
constants.

To make the almost-everywhere statement precise, use
a countable dense set of gauge transformations on
the finitely many cells in the stencil, extended
arbitrarily by the identity elsewhere.  Strictness
of (78.1) means any witnessing transformation can be
approximated by one from this set.  For fixed exterior
data the exceptional sets of coarse values for these
countably many transformations have Haar measure zero.
Their union remains null.  This yields one
almost-everywhere bound on the saturated data set,
without taking an uncountable union of null sets.
\end{proof}

\begin{corollary}[A precise separation of collar failure]
Let \(\mathcal G_{r,e}\) be the event that the exterior
and star coarse data satisfy the gauge-saturated
condition.  Under the original finite-volume Wilson law,
\[
 \mu_\beta(\{W\notin\mathcal A_e\}\cap\mathcal G_{r,e})
      \leq Q_\beta\,\mu_\beta(\mathcal G_{r,e}),
 \qquad
 \mu_\beta(W\notin\mathcal A_e)
      \leq Q_\beta+\mu_\beta(\mathcal G_{r,e}^c).
                                                        \tag{78.5}
\]
\end{corollary}
\begin{proof}
The collar event is measurable with respect to the
fixed exterior links and the star coarse values.
Condition on these data and use (78.4) and its
gauge-saturated version, then integrate.  The
absolute continuity established in V73 handles
the coarse exceptional sets on this event.
\end{proof}

\begin{assessment}[No independence or global small-field assumption]
Equation (78.5) is not an exponential unconditional
bound: the collar complement is still present.
It uses neither independence of neighbouring collars
nor a claim that all links of the whole lattice
can be gauged close to \(I\).  Global toron and
topological issues cannot be dismissed by a local
gauge saturation.  Nor does one conditional local
tail by itself control connected coarse interactions.

The saturated criterion removes a coordinate artefact,
but converting its failure into a controlled collection
of gauge-invariant plaquette events requires a
finite-stencil filling argument with the actual
exterior/coarse geometry.  That conversion has
not been proved in this section.
\end{assessment}

\begin{decision}[Next audit and collar geometry]
At V75 perform the required companion audit, then
test whether small plaquette curvature on a declared
enlarged stencil forces \(\mathcal G_{r,e}\).
Account for the root-block values, not just fine
boundary links, and isolate any noncontractible loops
if the stencil wraps the periodic lattice.
Only after this implication is proved should collar
failure be charged to neighbouring curvature defects.
\end{decision}

\begin{firewall}[Position after V74]
A nonempty gauge-saturated flat-collar regime now
has exponential conditional suppression of the actual
relative-path failure event, uniformly in total volume.
The collar-complement and connected-interaction bounds
needed for iteration remain incomplete.  The full
Yang--Mills mass gap is not proved.
Only V74 remains active; prior bodies and archives
are retained.  Next companion audit: V75.
\end{firewall}

% END APPENDED POST-TENTH-ARCHIVE V74 MATERIAL

% =============================================================================
\section{V75: curvature forces the flat collar on a non-wrapping box}
\label{sec:v75-collar}
% =============================================================================

The V74 predecessor has SHA--256
\texttt{9C6FE088A445C6166279EECB0E537140BCC96F0DDE47EF1639EB303ACAB93BE0}.
Its body is retained except for the version label.
This section proves the missing finite-stencil geometric
implication and performs the compulsory companion audit.
No probabilistic independence is inferred from that implication.

For this section require each fine lattice period to be
at least twelve.  Around the central root \(r\), choose
the translated box
\[
 \mathcal B_r=r+[-4,5]^4
 \]
in an injective lift to the lattice, including all its
internal edges and faces.  It has side length nine,
is a union of complete parity cells, and contains the
V74 stencil, all its cell-tree paths, and the two-link
paths defining the star blocks.  The deliberately enlarged
box avoids a wrap-around loop.  Let
\[
 D_r=\max_{p\subset\mathcal B_r}\delta(W_p).
 \]
Plaquette distances are unchanged by the tree gauge
used to obtain \(W\) from the original fine field.

\begin{lemma}[Cell-constant approximate flattening]
\label{lem:v75-flatten}
For every tree-slice field \(W\), there are cell-constant
gauge values \(q\) on \(\mathcal B_r\) such that every
edge contained in that box satisfies
\[
 \delta((q\cdot W)_l)\leq243D_r.                 \tag{79.1}
\]
No logarithm of an individual link is needed.
\end{lemma}
\begin{proof}
First choose monotone coordinate-ordered paths from
the lower corner of the box to every vertex, and let
\(g_x\) be their holonomies.  For a positive edge
\(x\to x+\hat\mu\), compare its path prefixed by \(g_x\)
with the ordered path to its endpoint.  Moving the
added step across later coordinate steps requires at
most \(3\cdot9=27\) square moves.  Telescoping these
plaquette factors, using unitary invariance and the
triangle inequality, gives
\(\delta(g_xW_{xy}g_y^{-1})\leq27D_r\).

If \(s\) is the root of the parity cell containing
\(x\), its cell-tree path has length at most four
and has original holonomy \(I\), since \(W\) is
on the tree slice.  Applying the last bound to
the edges of that path gives
\(\delta(g_sg_x^{-1})\leq108D_r\).
Set \(q_s=g_s\) and use this value on the whole cell.
For an edge \(x\to y\), with cell roots \(s,t\),
\[
 q_sW_{xy}q_t^{-1}
 =(g_sg_x^{-1})(g_xW_{xy}g_y^{-1})(g_yg_t^{-1}).
 \]
Its distance is at most
\((108+27+108)D_r=243D_r\).
All paths stay inside the box.  Cell-constant \(q\)
preserves the tree slice.
\end{proof}

\begin{theorem}[Small curvature implies the actual saturated collar]
\label{thm:v75-implication}
There is a fixed \(\zeta>0\), independent of volume,
such that
\[
 D_r<\zeta\quad\Longrightarrow\quad
                       \mathcal G_{r,e}
       \quad\text{for every }e\in E_r,           \tag{79.2}
\]
where \(\mathcal G_{r,e}\) is the V74 condition on the
exterior and the actual star coarse values \(K(W)\).
\end{theorem}
\begin{proof}
Use \(q\) from the lemma.  The required exterior links
are within distance \(243D_r\) of \(I\).
Each two-link star product is within \(486D_r\);
each relative product is within \(972D_r\).
Take \(972\zeta<\rho\), with \(\rho\) as in (72.1).
Then all star blocks use the good branch.

On a smaller fixed neighbourhood choose \(L_{\log}\)
such that \(\|\log V\|\leq L_{\log}\delta(V)\).
For anti-Hermitian \(X\), integration of the exponential
gives \(\delta(e^X)\leq\|X\|\).
The defining mean therefore obeys
\[
 \delta(K_f(q\cdot W))
       \leq (486+972L_{\log})D_r
       \quad(f\in E_r).
 \]
Choose \(\zeta\) still smaller so both this bound and
\(243\zeta\) are strictly below the V74 threshold \(a\).
Equivariance identifies these outputs with \(q\cdot K(W)\).
Thus this single cell-constant gauge witnesses the
required small exterior and coarse data.
It can be extended to the other cells arbitrarily.
This proves (79.2), including the coarse-value part
of the condition rather than only its fine-link part.
\end{proof}

\begin{corollary}[Local and collective action cost of collar failure]
For each \(r\), select a fixed translated incident edge
\(e(r)\), and count the failed conditions
\(\mathcal G_{r,e(r)}\) by \(N_{\rm col}\).  Then
\[
 \mathcal G_{r,e}^c
   \Longrightarrow
   \sum_{p\subset\mathcal B_r}s_p\geq\zeta^2/2,
 \qquad
 S_0\geq\frac{\zeta^2}{1250}N_{\rm col}.         \tag{79.3}
\]
\end{corollary}
\begin{proof}
The contrapositive of (79.2) gives a plaquette distance
at least \(\zeta\), hence the first inequality.
A fixed plaquette is contained in at most \(5^4=625\)
boxes with even roots: in each coordinate there are
at most five allowed root positions.  Sum the first
inequality over failed collars and use this overlap
bound.  No disjointness or independence is required.
\end{proof}

\begin{corollary}[What unconditional control presently follows]
The box contains
\(6\cdot9^2\cdot10^2=48600\) plaquettes.  The V68
energy estimate gives
\[
 \mu_\beta(\mathcal G_{r,e}^c)
 \leq\min\left(1,\frac{97200}{\zeta^2}
          \frac{2d\log\beta+C_G}{\beta}\right).  \tag{79.4}
\]
Together with (78.5) this bounds relative-path failure
by \(Q_\beta\) plus the right side of (79.4).
\end{corollary}
\begin{proof}
Apply Markov's inequality to the first sum in (79.3).
At each fixed plaquette orientation, translation
invariance and nonnegativity give
\(\mathbb E s_p\leq(2d\log\beta+C_G)/\beta\),
as proved in V68.  Sum over the displayed face count.
\end{proof}

\begin{audit}[Required V75 companion check]
SM remains V72, with unchanged SHA--256
\texttt{F44F9745D43379E1672C5550411B58E97195A4C9278592D221DF5D3F644DC1F5}.
Its endpoint is the fixed-mode many-copy limit, with
growing-cutoff covariance and source estimates still
required.  NS remains V26 with SHA--256
\texttt{82C887F8C2C69E4F28FAD7C3DC8DE5B9099CB5A4BF431EBA1FFF3E91935978AD};
its rank-one norm result does not close global regularity.
The parallel YM active version remains V5.
No updated companion estimate supplies a conditional
defect-cluster bound.  Their scopes therefore do not
alter the finite-box proof or its probability limitations.
Next companion audit and ten-version review: V80.
\end{audit}

\begin{assessment}[No cluster estimate has been smuggled in]
The geometry closes the implication left open in V74.
It does not improve the currently derived unconditional
rate beyond order \(\log\beta/\beta\).
The collective deterministic cost (79.3) alone does
not supply the partition-function cancellation or
conditional entropy bound for a defect cluster.
Iterating the local estimate across failed collars
without such a comparison would be circular.

The period restriction is essential to this particular
path proof.  Smaller tori require a different treatment
of wrapping loops; no assertion here removes global
holonomies or topological sectors.  The constants are
deliberately conservative finite-stencil bounds, not
optimized physical scales.
\end{assessment}

\begin{decision}[Next noncircular step]
Analyze a connected region of collar failures as one
conditional integration problem, rather than repeatedly
conditioning a cell on an already failed collar.
Identify which coarse constraints and exterior holonomies
must be retained in that region and whether a common
repair has a matched determinant and relative-energy
gain.  Compare with archived polymer hypotheses before
claiming those hypotheses have been populated.
\end{decision}

\begin{firewall}[Position after V75]
Small curvature on a declared non-wrapping box forces
the actual gauge-saturated collar; collar failures
have a bounded-overlap action cost.  Conditional
cluster suppression and iterative continuum control
remain unproved, as does the Yang--Mills mass gap.
Only V75 remains active; prior bodies and archives
are retained.
\end{firewall}

% END APPENDED POST-TENTH-ARCHIVE V75 MATERIAL

% =============================================================================
\section{V76: root-reference pivots remove the upper-density collar hypothesis}
\label{sec:v76-root-pivot}
% =============================================================================

The V75 predecessor has SHA--256
\texttt{B6825F71444BAA8B1BB82D0DAC1A6FE6CA826128BCEC89EA6FE94EE1F424B7B3}.
Its body is retained except for the active-version label.
In a connected region both endpoint cells of an interior
coarse edge are integrated, so V73's exterior-partner
argument is not automatically available.  The fixed tree
nevertheless supplies a distinguished pivot that avoids
that argument altogether.

\begin{definition}[Reference convention for this continuation]
For every \(e=(r,\mu)\), choose the reference path starting
at \(x=r\), and its boundary link as pivot.  Denote the
compact block (71.2) with this convention by \(K^0\).
This specializes the previously free reference choice.
Earlier analytic results stated for arbitrary references
remain applicable with this choice.  Different reference
means need not agree as nonlinear maps; no equality
with a previously fixed different map is asserted.
For a dressed local comparison use the V57 block with
this same root reference.  Keep the sufficiently small
logarithm neighbourhood required in V73.
\end{definition}

\begin{lemma}[Exact duplicate reference product]
\label{lem:v76-duplicate}
On the tree slice, for arbitrary internal links, the
two paths starting at \(r\) and \(r+\hat\mu\) have
identical product equal to the pivot boundary link \(B_e\):
\[
 P_r=P_{r+\hat\mu}=B_e.                          \tag{80.1}
\]
None of the other fourteen path products uses \(B_e\).
\end{lemma}
\begin{proof}
The edge from a cell root to that root plus \(\hat\mu\)
is in the coordinate-ordered tree for every \(\mu\).
It is therefore \(I\), both in the source and target
cells.  The first path consists of this source tree
edge followed by \(B_e\); the second consists of
\(B_e\) followed by the target root tree edge.
Each boundary link occurs in exactly this pair of
two-link paths, proving the last assertion.
\end{proof}

\begin{theorem}[Uniform one-pivot bound without an internal collar]
\label{thm:v76-density}
Fix all internal links and the seven nonpivot boundary
links of a coarse edge, arbitrarily in \(G\).
The pushforward of normalized pivot Haar under
\(B_e\mapsto K^0_e\) has a density bounded above by
a constant \(M_0<\infty\), independent of these data.
\end{theorem}
\begin{proof}
Let the other fourteen products be fixed \(P_j\).
On the good subset the map is exactly
\[
 B\longmapsto
 B\exp\left(\frac1{16}\sum_{j=1}^{14}
                             \log(B^{-1}P_j)\right). \tag{80.2}
\]
The two omitted logarithms are identically zero,
by (80.1).  On the fallback subset the map is \(B\);
its restricted Haar pushforward is at most Haar.

If the good subset is empty, the assertion follows.
Otherwise choose one fixed product \(A=P_1\).
Existence of a good \(B\) puts every fixed \(P_j\)
close to \(A\), uniformly when \(\mathcal O\) is small.
Also every good pivot is \(B=A\exp u\) in a common
small coordinate ball.  Extend (80.2) smoothly to
a slightly larger convex \(u\)-ball using an injective
logarithm chart containing all its arguments.
In output coordinates put \(F(u)=\log(A^{-1}K^0_e)\).
When every \(P_j=A\), formula (80.2) gives
\[
 F(u)=\log(\exp u\,\exp(-14u/16))=u/8
 \]
throughout this ball.  Uniform continuous dependence
of first derivatives on the fourteen nearby fixed
products, and compactness in \(A\), allow the
logarithm neighbourhood to be chosen such that
\(\|DF-I/8\|\leq1/16\) on the whole ball.
The segment argument of V73 makes this extension
injective with inverse derivative norm at most sixteen.
Uniform upper and lower Haar chart densities then
bound the image density of the good subset by a
constant \(M_g\).
Adding the fallback contribution gives \(M_0=1+M_g\).
The proof does not bound, or require smallness of,
any internal link.  It only uses fixed path products
being close when a good pivot exists.
\end{proof}

\begin{corollary}[Full and regional Haar pushforwards are bounded]
\label{cor:v76-regional}
Let \(E\) be any finite set of coarse edges, and
integrate all their pivot links independently with
normalized Haar.  Condition on any other links.
Then
\[
 (K^0_E)_*(dB_E)\leq M_0^{|E|}\,dc_E.           \tag{80.3}
\]
Integrating any additional internal or nonpivot
variables with normalized product Haar preserves
the same bound.  In particular on the full finite
lattice with \(n\) coarse cells,
\[
 (K^0)_*\lambda=g_0(C)\,dC,\qquad
                 g_0(C)\leq M_0^{4n}.           \tag{80.4}
\]
There is no singular component, without any collar
hypothesis.  The same absolute continuity holds
after multiplication by a Wilson weight.
\end{corollary}
\begin{proof}
With the other variables fixed, the output for each
edge depends on its own pivot only.  Their conditional
joint law is a product of the one-pivot laws, giving
(80.3).  Integration preserves domination.
The Wilson weight lies between zero and one before
normalization, so its pushforward is dominated by
the unweighted pushforward.
\end{proof}

\begin{proposition}[Connected-region dimension and the remaining asymmetry]
Let \(R\) contain \(s\) coarse cells, and let \(E(R)\)
be all coarse edges incident to them.  If \(i(R)\)
is the number of edges with both endpoints in \(R\),
then
\[
 m:=|E(R)|=8s-i(R),\qquad
 \dim(\hbox{local repair parameters})=
                     d(17s+7m)\leq73ds.         \tag{80.5}
\]
The upper Haar pushforward bound for these edges is
\(M_0^m\), even when the internal fields in \(R\)
are unrestricted.
\end{proposition}
\begin{proof}
Counting eight incident edges per cell counts an
internal edge twice and an exterior incident edge
once, so \(8s=2i(R)+(m-i(R))\).
For a repair tube, there are seventeen free internal
links per cell and seven free boundary links per
affected coarse edge, with one pivot solved per edge.
This gives the dimension count.  The upper bound
follows from (80.3) after integrating the internal
and nonpivot links.
This count alone is not a proof that the repair
tube has favourable energy for arbitrary exterior data.
\end{proof}

\begin{assessment}[A useful obstruction removed, not a cluster theorem]
The upper-density obstruction for edges internal to
a connected region is removed by the specified
reference convention.  V73's collar is no longer
needed for that upper bound.  It remains relevant
to the existing lower repair construction and its
energy comparison; those hypotheses cannot be
dropped as a consequence of (80.3).

Both the upper factor \(M_0^m\) and the prospective
lower tube factors now scale with the region rather
than total volume.  To obtain cluster suppression
one still needs a relative action gain proportional
to the defects after preserving the coarse data,
and control of exterior compatibility.  A dimension
count and an ambient Haar density bound do not
establish either property.
\end{assessment}

\begin{decision}[Next matched region estimate]
Use the root-reference convention to extend the lower
repair tube to a connected region with an explicitly
small exterior internal collar.  Track the action
of its central representative, the number of affected
plaquettes and the free-coordinate volume.
Compare this with a declared defect count in that
same region; do not use the absolute global Wilson
action as though it were conditional excess energy.
\end{decision}

\begin{firewall}[Position after V76]
For the root-reference block, the full local and global
Haar pushforwards have bounded densities with explicit
edge-count dependence and no internal-collar requirement.
Connected-region energy comparisons, scale iteration
and the full Yang--Mills mass gap remain incomplete.
Only V76 is active; prior bodies and archives remain.
Next companion audit and ten-version review: V80.
\end{firewall}

% END APPENDED POST-TENTH-ARCHIVE V76 MATERIAL

% =============================================================================
\section{V77: matched suppression for defect-dense regions with flat data}
\label{sec:v77-region}
% =============================================================================

The V76 predecessor has SHA--256
\texttt{EE0404F0BC57D46BCB03A1803E8EFBB6AE8E2369B18A8D14A3B68AF372AA238F}.
Its body is retained except for the version label.
Use the root-reference block \(K^0\).
Let \(R\) be a finite connected set of \(s\) coarse cells,
\(E(R)\) its \(m\) incident coarse edges, and \(\Lambda_R\)
the \(17s\) non-tree internal links together with all
\(8m\) boundary links on those edges.  Fix the remaining
links at \(\xi\), and condition on the affected coarse
values \(c=K^0_{E(R)}\).  All unaffected coarse values
depend only on \(\xi\).
Let \(\mathcal P_R\) be the plaquettes touching
\(\Lambda_R\), and \(H_\xi\) their Wilson action.
Write \(D=17s+7m\leq73s\) for the number of free
Lie-algebra coordinates in the repair.

\begin{theorem}[Regional comparison with explicit size dependence]
\label{thm:v77-comparison}
Assume the exterior internal links used by affected
block equations are sufficiently small.  There are
constants \(b>0,r_*>0,A_\eta<\infty\), independent
of \(R,c,\xi\) and total volume, such that a local
repair tube has central representative \(V_0(c,\xi)\)
and
\[
 (K^0_{E(R)})_*(e^{-\beta H_\xi}dV)(dc)
 \geq b^s t^{dD}
 e^{-\beta((1+\eta)H_0(c,\xi)+A_\eta s t^2)}dc,
 \quad 0<t\leq r_*,                             \tag{81.1}
\]
where \(H_0=H_\xi(V_0)\).
If an event \(B\) has \(H_\xi\geq E\) on its conditional
fibre, then for almost every \(c\),
\[
 \mathbb P_{\beta,\xi}(B\mid c)
 \leq
 M_0^m b^{-s}t^{-dD}
 e^{-\beta(E-(1+\eta)H_0-A_\eta s t^2)}.         \tag{81.2}
\]
\end{theorem}
\begin{proof}
Give all central internal links free exponential
coordinates; give seven boundary links per affected
edge free coordinates about \(c_e\); solve its pivot.
Every block equation still involves only its own pivot.
At zero free variables and zero exterior internal
parameters its derivative is \(I/8\).
The uniform local contraction of V71--V72 therefore
works simultaneously on the region, with
\(D=17s+7m\) free coordinates.  It preserves the
exterior and every coarse value.  The product
Jacobian has a lower bound \(a^s\), since the
number of its bounded local factors is at most
a constant times \(s\).  Constants can be decreased
to at most one when taking this bound.

There are \(17s+8m\leq81s\) changed group links,
so \(|\mathcal P_R|\leq486s\).
The same link-difference and plaquette estimate as
V72 gives
\(H_\xi(V(c,z))\leq(1+\eta)H_0+A_\eta s t^2\).
Integrating the \(D\) radius-\(t\) balls proves
(81.1), absorbing their unit volumes into \(b^s\).
Such an absorption is uniform because \(D/s\)
lies between fixed positive constants.
The unweighted full pushforward is bounded by
\(M_0^m dc\) by V76.  Its weighted bad part is
therefore at most \(e^{-\beta E}M_0^m dc\).
Division proves (81.2).  Exterior-only action
terms cancel exactly.
\end{proof}

\begin{lemma}[Defect cost inside the integrated region]
Let \(F\) be a specified set of \(k\) coarse edges
whose source and target cells both belong to \(R\).
On the event \(B_F\) that all these edges fail
their relative-logarithm conditions,
\[
 H_\xi\geq\gamma k,\qquad
                   \gamma=c_\rho/8=\rho^2/2304. \tag{81.3}
\]
\end{lemma}
\begin{proof}
The box \(\mathcal R_e\) for such an edge lies in
its two endpoint cells and the interface between them.
Every plaquette in it contains a non-tree link
changed in \(\Lambda_R\): an elementary square
cannot consist solely of edges of the cell trees.
Thus its action is part of \(H_\xi\).
Each failed edge has box action at least \(c_\rho\)
by V68; each plaquette occurs in at most eight
of these boxes.  Summing proves (81.3).
\end{proof}

\begin{theorem}[A populated defect-dense conditional regime]
\label{thm:v77-dense}
Fix \(\theta>0\), and assume \(k\geq\theta s\).
Suppose all exterior links in \(\mathcal P_R\),
all exterior internal links used by the block
equations, and all affected coarse links \(c_e\)
are within distance \(a\) of \(I\).
There is \(a_\theta>0\), independent of region size,
such that for \(a\leq a_\theta\) and \(\beta\geq1\),
\[
 \mathbb P_{\beta,\xi}(B_F\mid c)
 \leq
 \left(C_\theta\,
       \beta^{73d/(2\theta)}e^{-\beta\gamma/2}\right)^k
                                                        \tag{81.4}
\]
for almost every such coarse value, with constants
independent of total volume and \(R\).
\end{theorem}
\begin{proof}
The central repair has each changed link within
\(L_0a\) of \(I\), with uniform \(L_0\).
As in V74,
\[
 H_0\leq A_0a^2s,\qquad A_0=3888L_0^2.
 \]
Choose \(a_\theta\) within the chart radii and so
\(2A_0a_\theta^2\leq\gamma\theta/2\).
With \(\eta=1\), (81.3) gives
\(\gamma k-2H_0\geq\gamma k/2\).
Use (81.2) with \(t=r_*/\sqrt\beta\).
Because \(m\leq8s\) and \(D\leq73s\), all constant
factors, including \(e^{A_1r_*^2s}\), are bounded
by \(C^s\), with \(C\geq1\).
The power of \(\beta\) is at most \(73ds/2\).
Finally \(s\leq k/\theta\) yields (81.4).
The radius can be decreased so \(r_*\leq1\),
making the constant-factor bounds immediate.
\end{proof}

\begin{assessment}[Gauge saturation and the unpopulated hypothesis]
The flat-data hypothesis can be replaced by existence
of a cell-constant gauge transformation making these
data small.  Haar measure, block equivariance and
the countable-witness argument of V74 give the same
conditional estimate.  This is a sufficient regime,
not a statement that every failed-collar cluster
has such data.

In particular the affected coarse values include
edges throughout the region.  A good exterior
collar does not ensure those retained values are
close to a pure gauge.  Large retained curvature
can make \(H_0\) extensive and destroy the excess
in (81.2).  The lower repair comparison must then
be made relative to an appropriate coarse-dependent
background, not forced into the flat competitor.
Neither conditioning nor the defect count alone
removes this retained curvature.
\end{assessment}

\begin{decision}[Next upstream comparison]
Determine whether the action cost associated with
the retained coarse curvature can be shared between
the bad numerator and a better repair denominator.
First seek a deterministic fine-to-coarse curvature
comparison for the actual root-reference map, with
its good and fallback branches both included.
Its constants and direction must be checked before
using it to claim a relative defect energy.
\end{decision}

\begin{firewall}[Position after V77]
The upper and lower region factors are matched,
and defect-dense regions with specified flat data
have exponential conditional suppression.
Arbitrary retained coarse curvature and generated
interactions remain uncontrolled.  The full
Yang--Mills mass gap is not proved.
Only V77 is active; prior bodies and archives
remain.  Next companion audit and review: V80.
\end{firewall}

% END APPENDED POST-TENTH-ARCHIVE V77 MATERIAL

% =============================================================================
\section{V78: full-branch curvature comparison and its uncancelled cost}
\label{sec:v78-curvature}
% =============================================================================

The V77 predecessor has SHA--256
\texttt{0DB1AEC17C27114A76EBE16412B193EAFA637CA709CABC1A53A19FE9A7BFAEDC}.
Its body is retained except for the version label.
This section tests the deterministic comparison requested
there.  Use the root-reference block \(K^0\), the same
Wilson normalization, and a constant \(L_{\log}\) such that
\(\|\log V\|\leq L_{\log}\delta(V)\) on \(\mathcal O\).
The latter constant is finite for the chosen sufficiently
small logarithm neighbourhood.

\begin{lemma}[Distance from the straight reference block]
\label{lem:v78-edge}
Let \(B_e=P_r(W)\), the straight root two-link product,
and put \(d_e=\|K^0_e(W)-B_e\|_{\rm HS}\).
For every compact tree-slice field, on both branches,
\[
 d_e^2\leq72L_{\log}^2 S_e.                     \tag{82.1}
\]
\end{lemma}
\begin{proof}
On the fallback branch \(K^0_e=B_e\), so the left side
is zero.  On the good branch,
\[
 d_e=\delta\left(\exp\left[
       \tfrac1{16}\sum_x\log(B_e^{-1}P_x)\right]\right)
 \leq L_{\log}\max_x\delta(B_e^{-1}P_x).
 \]
The straight reference needs no filling moves.
The proof of V68 therefore bounds each relative
distance by six, rather than twelve, times the
largest plaquette distance in \(\mathcal R_e\).
That largest squared distance is at most \(2S_e\).
Squaring gives (82.1).  No continuity across the
fallback boundary is required.
\end{proof}

\begin{theorem}[Uniform fine-to-coarse Wilson action comparison]
\label{thm:v78-action}
For every field on the finite lattices considered above,
\[
 S_{\rm coarse}(K^0(W))
       \leq C_{\rm curv}S_0(W),\qquad
 C_{\rm curv}=20+8640L_{\log}^2.                 \tag{82.2}
\]
This constant is deliberately conservative and
independent of volume.
\end{theorem}
\begin{proof}
For a coarse plaquette \(P\), let \(B_P\) be the
plaquette product of the straight reference links.
Its fine path is the boundary of a two-by-two square.
Four elementary plaquette moves give
\[
 \delta(B_P)^2
       \leq4\sum_{p\ {\rm in\ that\ square}}\delta(W_p)^2
       =8\sum_{p\ {\rm in\ that\ square}}s_p.
 \]
Telescoping the four link replacements yields
\(\delta(K^0_P)\leq\delta(B_P)+\sum_{e\in\partial P}d_e\).
The square of a sum of five nonnegative terms is at
most five times their squared sum.  Hence
\[
 s_P(K^0)\leq
 20\sum_{p\ {\rm in\ the\ square}}s_p
       +180L_{\log}^2\sum_{e\in\partial P}S_e .
                                                        \tag{82.3}
\]
In the sum over coarse plaquettes, each fine plaquette
belongs to at most one of the straight two-by-two
squares of its orientation.  Each coarse edge is in
six coarse plaquettes.  Finally \(\sum_eS_e\leq8S_0\)
by the box-overlap count of V68.  Thus the second
term sums to at most
\(180\cdot6\cdot8\,L_{\log}^2S_0\).
This proves (82.2).
\end{proof}

\begin{corollary}[Bounds on the conditional ground energy]
For Haar-almost every coarse field \(C\), the
essential infimum \(m(C)\) of the fine action under
the conditional Haar law for \(K^0=C\) satisfies
\[
 C_{\rm curv}^{-1}S_{\rm coarse}(C)
        \leq m(C)\leq4S_{\rm coarse}(C).         \tag{82.4}
\]
\end{corollary}
\begin{proof}
The lower bound follows pointwise from (82.2).
For the upper bound, the section \(J(C)\) has action
\(4S_{\rm coarse}(C)\).  A single point would not
suffice for an essential-infimum assertion.
Instead use V71's regular tube: for each positive
integer \(j\), its sufficiently small free-coordinate
neighbourhood has action less than
\(4S_{\rm coarse}(C)+1/j\), positive fibre volume,
and a positive coarse density.
Such neighbourhoods may be chosen uniformly over
the compact coarse space at this fixed volume.
The full unweighted coarse density is finite and
bounded above by V76, and positive by V71.
Disintegration of the tube therefore gives positive
conditional mass to these sublevel sets for almost
every \(C\).  Intersect the countably many full-measure
sets and let \(j\) tend to infinity.
\end{proof}

\begin{proposition}[What combining the curvature and defect costs gives]
If a declared family of \(k\) failed coarse edges has
the V68 global action bound \(S_0\geq\gamma k\),
where \(\gamma=\rho^2/2304\), then for every
\(0\leq\lambda\leq1\),
\[
 S_0\geq
  \frac{\lambda}{C_{\rm curv}}S_{\rm coarse}(K^0)
                      +(1-\lambda)\gamma k.    \tag{82.5}
\]
This inequality does not give a positive defect
excess above the section action for arbitrary \(C\).
\end{proposition}
\begin{proof}
Multiply the two lower bounds by \(\lambda\) and
\(1-\lambda\) and add.  Subtracting the section
action from the resulting right side gives
\[
 (1-\lambda)\gamma k
       -\left(4-\lambda/C_{\rm curv}\right)
                                      S_{\rm coarse}(C).
 \]
The coarse term has the wrong sign for a uniform
positive excess when coarse curvature is unrestricted.
This is a limitation of the proved comparison, not
a counterexample to every possible sharper estimate.
Adding the two original lower bounds without the
convex weights would be unjustified.
\end{proof}

\begin{assessment}[A deterministic comparison is not the effective action]
The bound applies to both branches of the declared
compact map and controls its retained curvature.
It does not identify the conditional minimizer,
the effective coupling, or a renormalization law.
The interval (82.4) is far too wide to cancel
retained curvature in V77's section-based denominator.
Improving a constant is useful only if it provides
the needed relative-energy statement on the same fibre.
The current estimates do not do that.
\end{assessment}

\begin{decision}[Next constrained-minimum question]
Examine whether the section is stationary on a
nonflat coarse fibre.  If not, quantify the action
decrease obtainable by a block-preserving variation
and use that to identify an improved local background.
Do not continue treating the section as a vacuum:
it was constructed as an exact right inverse, not
as a constrained minimizer.  A matched repair for
curved retained fields must account for this distinction.
\end{decision}

\begin{firewall}[Position after V78]
A volume-uniform full-branch curvature comparison
and conditional ground-energy bounds are proved.
Their constants do not cancel arbitrary retained
curvature in the existing defect comparison.
The constrained-background and iterative continuum
problems remain open; no full mass gap is proved.
Only V78 is active; previous bodies and archives
remain.  Next companion audit and review: V80.
\end{firewall}

% END APPENDED POST-TENTH-ARCHIVE V78 MATERIAL

% =============================================================================
\section{V79: an exact nonstationary section on a curved coarse fibre}
\label{sec:v79-descent}
% =============================================================================

The V78 predecessor has SHA--256
\texttt{82B73E9E9724CC131816EF99B90485E6D25925D59DE359A8006196734FF3F8BC}.
Its body is retained except for the version label.
We test the section as requested there.  The example
uses a commuting subgroup of \(SU(2)\), but is an
actual nonabelian-group configuration and an exact
nonlinear block fibre, not just a linearized constraint.

Let \(T=\operatorname{diag}(i,-i)\), and let the fine
periods be even and at least six.  Put all coarse
links equal to \(I\), except the two direction-0
links based at coarse roots \(0\) and \(2\hat1\),
which are both \(e^{\theta T}\), with \(0<\theta<\pi\).
Let \(W(0)=J(C)\).  In the boundary family for the
direction-0 coarse edge based at \(0\), change the
four links with transverse bit \(s_1=0\) to
\(e^{(\theta-t)T}\), and the four with \(s_1=1\)
to \(e^{(\theta+t)T}\).  Leave every other fine
link unchanged.  Denote this field by \(W(t)\).

\begin{theorem}[Exact fibre preservation and action descent]
\label{thm:v79-descent}
For all sufficiently small \(|t|\),
\[
 K^0(W(t))=C,\qquad
 \left.\frac{d}{dt}S_0(W(t))\right|_{t=0}
                                      =-8\sin\theta. \tag{83.1}
\]
Thus the section is not a constrained stationary
point for these nonflat retained data.
\end{theorem}
\begin{proof}
All internal links remain \(I\), so the sixteen
path products are the eight boundary values, each
repeated twice.  On the changed edge their angle
average is exactly \(\theta\).
The root reference has angle \(\theta-t\);
the relative angles are zero or \(2t\).
For sufficiently small \(t\) these lie in the
declared logarithm neighbourhood, so the reference
mean is exactly \(e^{\theta T}\).  Other coarse
edges are unaffected.  This proves the nonlinear
fibre identity without an implicit correction.

For \(e^{\alpha T}\) the plaquette action is
\(2(1-\cos\alpha)\).  The changed plaquettes can
be counted by their angle, using evenness of cosine.
The net multiplicity changes from \(t=0\) are
\[
 \begin{array}{c|rrrrrr}
 \text{angle}&0&t&2t&\theta-t&\theta&\theta+t\\ \hline
 \text{change}&-8&4&4&12&-20&8 .
 \end{array}                                    \tag{83.2}
\]
To see the count directly, the four direction-\((0,1)\)
faces between the two changed transverse layers
have angle \(2t\).  Their four outgoing faces meet
the neighbouring coarse link with the same angle
\(\theta\), giving angle \(t\).  The four incoming
faces have angle \(\theta-t\) instead of \(\theta\).
For each of the other two transverse directions,
the crossing faces have four angles \(\theta-t\)
and four angles \(\theta+t\) instead of eight
angles \(\theta\).  The remaining faces either
are unchanged or have equal perturbations that
cancel.  This gives (83.2).

Consequently the exact action difference is
\[
 \begin{split}
 \Delta S(t)={}&16+40\cos\theta
       -8\cos t-8\cos(2t)\\
       &-24\cos(\theta-t)-16\cos(\theta+t).
 \end{split}                                    \tag{83.3}
\]
It vanishes at zero and its derivative there is
\(-8\sin\theta\), proving the theorem.
\end{proof}

\begin{corollary}[Quantified strict decrease]
For \(\theta=\pi/6\), and every sufficiently small
\(0<t\leq1/20\) for which the logarithm branch above
is valid,
\[
 S_0(W(t))\leq4S_{\rm coarse}(C)-2t.             \tag{83.4}
\]
More generally the leading small-angle improvement
with \(t=\theta/10\) is
\(\Delta S=-2\theta^2/5+O(\theta^4)\).
\end{corollary}
\begin{proof}
The absolute value of the second derivative in
(83.3) is at most \(8+32+24+16=80\).
For \(\theta=\pi/6\), Taylor's theorem thus gives
\(\Delta S(t)\leq-4t+40t^2\leq-2t\).
Use \(S_0(W(0))=4S_{\rm coarse}(C)\).
For the last assertion, expand (83.3) at
\((\theta,t)=(0,0)\):
\(\Delta S=-8\theta t+40t^2+
O((|\theta|+|t|)^4)\), and substitute \(t=\theta/10\).
For small enough \(\theta\), this also stays on
the required logarithm branch.
\end{proof}

\begin{record}[Independent exact count check]
The dependency-free script
\texttt{scripts/verify\_ym\_v79\_section\_descent.py}
enumerates a fine period-six lattice and represents
every link angle by its integer coefficients of
\(\theta,t\).  It computes every oriented plaquette
angle, identifies opposite signs only because
cosine is even, and checks exactly the six
multiplicity differences in (83.2).
All checks passed; the output is
\texttt{artifacts/ym\_v79\_section\_descent.json}.
The check uses no numerical trigonometric evaluations.
It verifies the finite incidence count, not a
continuum or mass-gap assertion.
This example does not use V75's separate period-twelve
non-wrapping-box theorem.
\end{record}

\begin{assessment}[Consequence for the repair programme]
The gap between the section action and the constrained
minimum is genuine, already at arbitrarily small
nonzero coarse curvature.  It is not merely looseness
in the bound of V78.  Section neighbourhoods remain
valid lower contributions to the conditional integral,
but may pay an unnecessary bulk action cost.
Calling the section a vacuum would therefore be
incorrect.

The present variation gives a better competitor for
these data, not a minimizer for general coarse fields.
It neither proves uniqueness nor controls the measure
near a minimizing configuration.  Constructing and
using a curved constrained background remains the
next substantive task.
\end{assessment}

\begin{decision}[V80 review and constrained background]
Perform the required companion audit and ten-version
review.  Then use the actual constrained stationarity
equation to construct a moving local minimum near
small coarse fields, keeping track of the boundary
of the chart and volume-uniform norms.
The existing local Hessian bound can help, but does
not alone put a critical point inside the chart
uniformly over arbitrary coarse data.
\end{decision}

\begin{firewall}[Position after V79]
An exact block-preserving descent disproves constrained
stationarity of the section for explicit curved data.
The analytic action formula and independent integer
count agree.  A suitable general moving background,
cluster control and the full Yang--Mills mass gap
remain unproved.  Only V79 is active; earlier bodies
and archives remain.  Next audit and review: V80.
\end{firewall}

% END APPENDED POST-TENTH-ARCHIVE V79 MATERIAL

% =============================================================================
\section{V80: a volume-uniform moving minimum for the actual local fibre}
\label{sec:v80-minimum}
% =============================================================================

The V79 predecessor has SHA--256
\texttt{E508B98F82A0A7056128E5913495F35A22F9F18462ED97C3F37AD33EB3E1DBD7}.
Its body is retained except for the version label.
The archives contain moving-background constructions for
other declared blockings.  Here the input is specifically
the root-reference nonlinear block and its V64--V66 chart.
The new step is a volume-uniform max-norm inverse and
existence radius, not merely a finite-dimensional
implicit-function assertion.

Use independent linear fibre coordinates \(z\), consisting
of seventeen internal and twenty-eight nonpivot boundary
Lie-algebra values per cell.  Let \(E z\in\ker L_{\rm tr}\)
be their completion by the V63 pivot formula, with tree
entries zero, and write
\[
 X(z,c)=\Psi_{\rm tr}(Ez,c),\qquad
 f(z,c)=S_0(\exp X(z,c)).
 \]
The map \(E\) is finite range, has uniformly bounded
operator norm, and satisfies \(\|Ez\|_2\geq\|z\|_2\),
because its free entries are exactly \(z\).  Coordinate
max norms below use an orthonormal real Lie-algebra basis.

\begin{lemma}[A finite-range gap gives a max-norm inverse]
\label{lem:v80-inverse}
The matrix \(H=D_z^2 f(0,0)\) has a bound
\[
 \|H^{-1}\|_{\infty\to\infty}\leq C_\infty
                                                        \tag{84.1}
\]
independent of volume.
\end{lemma}
\begin{proof}
V66 and the property of \(E\) give
\(H=E^*d^*dE\geq\kappa I\).
Finite-range entries and bounded stencil overlap give
a volume-independent upper spectral bound \(\Lambda\).
Increase \(\Lambda\) if necessary so \(\Lambda>\kappa\).
The self-adjoint matrix \(T=I-H/\Lambda\) has
\(\|T\|_{2\to2}\leq q=1-\kappa/\Lambda<1\)
and a fixed finite propagation range.
A row of \(T^j\) is supported in a graph ball with
at most \(C(j+1)^4\) scalar coordinates, including
the fixed number of coordinates per cell.  The same
upper bound holds on periodic lattices.
Its row Euclidean norm is at most \(q^j\), hence
its row sum is at most \(C^{1/2}(j+1)^2q^j\).
The convergent Neumann series gives
\[
 H^{-1}=\Lambda^{-1}\sum_{j=0}^\infty T^j,\qquad
 C_\infty=\frac{C^{1/2}}{\Lambda}
                       \sum_{j=0}^\infty(j+1)^2q^j<\infty .
 \]
This proves (84.1) without replacing a max norm by
a volume-dependent Euclidean norm.
\end{proof}

\begin{theorem}[Uniform small-coarse-field constrained minimum]
\label{thm:v80-minimum}
There exist \(r,\delta,A>0\), independent of volume,
such that for every \(\|c\|_\infty<\delta\) the equation
\[
 \nabla_z f(z,c)=0                              \tag{84.2}
\]
has a unique solution \(z_*(c)\) in
\(\|z\|_\infty<r\), with
\(\|z_*(c)\|_\infty\leq A\|c\|_\infty\).
The solution is smooth in \(c\), lies inside the
common nonlinear chart, and is a strict local
constrained minimum.  In that ball,
\[
 f(z,c)-f(z_*(c),c)
                  \geq\frac{\kappa}{4}\|z-z_*(c)\|_2^2 .
                                                        \tag{84.3}
\]
\end{theorem}
\begin{proof}
The composed action is a sum of smooth finite-stencil
functions on a common chart.  Their derivatives are
uniformly bounded.  Therefore, in independent coordinates,
\[
 \|D_z^2f(z,c)-H\|_{\infty\to\infty}
       \leq C_1(\|z\|_\infty+\|c\|_\infty),\qquad
 \|\nabla_z f(0,c)\|_\infty\leq C_2\|c\|_\infty .
 \]
These are row bounds from local functions, not bounds
on an extensive total derivative.  The second inequality
uses \(\nabla_z f(0,0)=0\).
Consider \(T_c(z)=z-H^{-1}\nabla_z f(z,c)\).
Choose \(r,\delta\) inside the chart so that
\(C_\infty C_1(r+\delta)\leq1/2\) and
\(C_\infty C_2\delta<r/2\).
Then \(T_c\) contracts the closed radius-\(r\) ball
into itself.  Its unique fixed point is strictly
inside and satisfies
\(\|z_*\|_\infty\leq2C_\infty C_2\|c\|_\infty\).

Uniform local Hessian continuity in Euclidean operator
norm, again by bounded stencil overlap, allows the
radii to be reduced so \(D_z^2f\geq\kappa I/2\).
Integrating this bound along the segment from \(z_*\)
to \(z\) proves (84.3) and uniqueness on the ball.
Invertibility of the Hessian and the ordinary smooth
implicit-function theorem give smooth dependence.
The uniform contraction, rather than that theorem
alone, supplies the volume-independent domain.
\end{proof}

\begin{corollary}[Measured local minimum at large coupling]
Let \(\ell(z,c)=\ell_{\rm tr}(X(z,c))\).
For sufficiently large \(\beta\), uniformly in volume,
the measured potential
\[
 f_\beta(z,c)=f(z,c)-\beta^{-1}\ell(z,c)
 \]
has a unique small critical point \(z_{\beta,*}(c)\)
on a possibly smaller common chart, and
\[
 \|z_{\beta,*}(c)-z_*(c)\|_\infty
                  \leq A'\beta^{-1}\|c\|_\infty. \tag{84.4}
\]
\end{corollary}
\begin{proof}
The density is a finite-range smooth sum by V63--V65,
and \(D_z\ell(0,0)=0\).  Its Hessian has uniform
row and Euclidean operator bounds.  The contraction
proof therefore tolerates its \(\beta^{-1}\) perturbation
for all sufficiently large \(\beta\).  At \(z_*(c)\),
the perturbed gradient is
\(-\beta^{-1}\nabla_z\ell(z_*(c),c)\), with max norm
at most \(C\beta^{-1}\|c\|_\infty\).
The inverse Hessian has a uniform max-norm bound
by a Neumann perturbation of (84.1).  Applying the
same contraction estimate about \(z_*(c)\) proves
(84.4).  Positive Euclidean Hessian makes the
critical point a strict local minimum.
\end{proof}

\begin{assessment}[Claim boundary for the new background]
These minima solve the actual stationarity equations
on a uniform neighbourhood of small coarse fields.
They are not asserted to minimize the entire compact
fibre, which may have configurations outside the chart.
Nor is a uniform neighbourhood of every curved compact
coarse field obtained.  The mass-gap route still needs
comparison of remote defects with the local minimum,
and control of the effective interactions after integration.
V79's explicit descent is consistent with \(z_*(c)\)
generally differing from the section value \(z=0\).
\end{assessment}

\begin{audit}[Required V80 companion check]
SM remains V72, SHA--256
\texttt{F44F9745D43379E1672C5550411B58E97195A4C9278592D221DF5D3F644DC1F5};
NS remains V26, SHA--256
\texttt{82C887F8C2C69E4F28FAD7C3DC8DE5B9099CB5A4BF431EBA1FFF3E91935978AD}.
Their fixed-mode many-copy and rank-one kernel results
do not provide the missing YM compact-fibre comparison.
The parallel YM active note remains V5.  No change
of direction is justified by a new companion theorem.
Next audit: V85; next ten-version review: V90.
\end{audit}

\begin{record}[Ten-version review]
V71--V73 made the compact conditional denominator
and local upper density quantitative.
V74--V75 populated a flat-collar suppression regime
and related its complement to curvature.
V76's explicit root-reference choice removed the
upper-density collar restriction.
V77 matched factors on connected regions, but only
under suitable retained data.
V78 proved a curvature comparison too weak to cancel
the section cost.  V79 demonstrated an actual
constrained descent from that section.
V80 supplies a local moving minimum with a uniform
existence radius for the declared nonlinear block.
The advance is beyond a repeated section estimate,
but remains local in retained field space.
\end{record}

\begin{decision}[Next estimate around the minimum]
Quantify spatial dependence of \(z_*(c)\) and the
Hessian of the minimized action, using finite-range
operators and their gapped fibre inverses.
Keep the retained coarse gauge directions visible;
the resulting Schur complement need not be gapped.
This should identify the correct curved quadratic
background before attempting a compact defect repair.
\end{decision}

\begin{firewall}[Position after V80]
A volume-uniform local moving constrained minimum
and its measured perturbation are proved.
Global fibre comparison, scale iteration and the
full Yang--Mills mass gap remain incomplete.
Only V80 is active; prior bodies and archives remain.
The V100 rollover policy is unchanged.
\end{firewall}

% END APPENDED POST-TENTH-ARCHIVE V80 MATERIAL

% =============================================================================
\section{V81: spatial response of the moving minimum and retained gauge modes}
\label{sec:v81-response}
% =============================================================================

The V80 predecessor has SHA--256
\texttt{0DE5CF121968B6F3BED6BE71C0AEEAA8BB86A90F212797DA95C0CA9928EF18D0}.
Its body is retained except for the version label.
Keep its independent fibre coordinates and shrink
the coarse chart once more if necessary.  All estimates
below are on that uniform chart, with periodic cell
graph distance.  They concern the moving nonlinear
minimum, not a claim that the entire coarse theory
has the eliminated fibre's gap.

At \((z_*(c),c)\), write the second derivative blocks
of \(f(z,c)\) as
\[
 H_c=f_{zz},\qquad A_c=f_{zc},\qquad B_c=f_{cc}.
 \]
They have finite range in their coordinate indices,
even though their values depend on the moving background.
There are constants \(0<h\leq L<\infty\), independent
of volume and \(c\), with \(hI\leq H_c\leq LI\).
These follow from V80 and bounded local derivatives.

\begin{theorem}[Exponential inverse and response bounds]
\label{thm:v81-response}
For constants \(C,\alpha>0\) independent of volume
and \(c\),
\[
 |(H_c^{-1})_{ij}|\leq C e^{-\alpha\,{\rm dist}(i,j)},
 \qquad
 \left|\frac{\partial z_{*,i}}{\partial c_e}\right|
                  \leq C e^{-\alpha\,{\rm dist}(i,e)}.
                                                        \tag{85.1}
\]
Scalar indices include fixed Lie-algebra components;
equivalent block-norm bounds follow.
\end{theorem}
\begin{proof}
Choose \(L'>L\) uniformly and put
\(T_c=I-H_c/L'\).  It has spectral norm at most
\(q=1-h/L'<1\) and a fixed propagation range \(R\geq1\).
In the Neumann series
\(H_c^{-1}=(L')^{-1}\sum_{j\geq0}T_c^j\),
the \(i,j\) entry is zero until the number of factors
is at least their distance divided by \(R\).
Each remaining entry is bounded by the spectral norm
of that power, at most \(q\) to that power.
Summing this geometric tail proves the first bound.

Differentiate the stationarity equation of V80:
\[
 D_c z_*=-H_c^{-1}A_c.                          \tag{85.2}
\]
Each column of \(A_c\) has a uniformly bounded
finite stencil and bounded entries.  The first
bound applied to this finite sum gives the second.
\end{proof}

\begin{corollary}[Nonlinear remote-data stability]
If two small coarse fields \(c,\widetilde c\) agree
within distance \(R_0\) of a fibre coordinate \(i\),
then
\[
 |z_{*,i}(c)-z_{*,i}(\widetilde c)|
 \leq C'e^{-\alpha' R_0}
                   \|c-\widetilde c\|_\infty,   \tag{85.3}
\]
for some \(\alpha'>0\), uniformly in volume.
\end{corollary}
\begin{proof}
The straight segment between the two fields stays
inside the convex coarse chart.  Integrate (85.2)
along it.  In each distance shell the number of
coarse coordinates grows at most polynomially
in the distance, with degree bounded by four.
The tail sum of the exponential in (85.1) is
bounded by a slightly slower exponential.
\end{proof}

\begin{theorem}[Minimized action and its spatial Hessian]
\label{thm:v81-schur}
For \(G(c)=f(z_*(c),c)\),
\[
 D_cG=f_c,\qquad
 D_c^2G=B_c-A_c^*H_c^{-1}A_c.                   \tag{85.4}
\]
There are volume-independent constants such that
\[
 |(D_c^2G)_{ef}|\leq C e^{-\alpha\,{\rm dist}(e,f)}.
                                                        \tag{85.5}
\]
\end{theorem}
\begin{proof}
The fibre-gradient term in the first chain rule
vanishes by stationarity.  Differentiating again
and substituting (85.2) gives (85.4).
The matrices \(A_c,B_c\) have finite range and
uniformly bounded entries.  Each entry of the
second term is a sum over two fixed-size stencils;
apply (85.1) between them.  The finite-range term
also satisfies an exponential bound after increasing
the constant.
\end{proof}

\begin{proposition}[Local gauge compatibility and the zero modes at the origin]
\label{prop:v81-gauge}
For sufficiently small cell-constant gauge transformations
whose transformed fields remain in the uniqueness chart,
the moving constrained minima transform into one another
and \(G\) is invariant under the induced coarse action.
In particular, with
\((d_{\rm c}\omega)_{rs}=\omega_r-\omega_s\),
\[
 D_c^2G(0)\,d_{\rm c}\omega=0.                  \tag{85.6}
\]
\end{proposition}
\begin{proof}
The tree slice is preserved by cell-constant gauge
transformations, and \(K^0\) is equivariant.
Such a transformation carries a fine fibre to the
corresponding transformed coarse fibre and preserves
the Wilson action.  A constrained critical point
is therefore carried to a constrained critical point.
For transformations small enough to remain in the
uniform uniqueness neighbourhood of V80, that point
is the corresponding moving minimum.

At \(c=0\), the identity is a zero-action minimum.
The small pure-gauge curve
\(C_{rs}(t)=e^{t\omega_r}e^{-t\omega_s}\)
has a zero-action fine representative obtained by
the cell-constant gauge transformation of the identity.
It stays in the chart and hence has minimized action
zero.  Its tangent in logarithmic coordinates is
\(d_{\rm c}\omega\).  Since \(DG(0)=0\), the Hessian
quadratic form on this tangent vanishes.
Finally \(G\geq0=G(0)\), so the Hessian at zero is
positive semidefinite.  A null quadratic direction
of a positive semidefinite symmetric matrix lies
in its kernel, proving (85.6).
\end{proof}

\begin{assessment}[What agrees with, and what exceeds, the Gaussian calculation]
At the origin (85.4) is the familiar harmonic
Schur complement already encountered in the
Gaussian work.  Its existence there is not a new
Gaussian mass-gap result.  The new control is the
uniform spatial response of the actual nonlinear
moving minimum throughout its small coarse chart.
The fibre inverse is gapped; the retained Hessian
still has gauge directions.  At nonzero fields
one must also retain the derivative-of-gauge-vector
term in an off-shell Ward identity, as in V51--V55.
Equation (85.6) is not being extended off shell by
discarding that term.

Most importantly, \(G\) is a minimized classical
action, not the logarithm of the full conditional
partition function.  Its exponential Hessian decay
does not establish reflection positivity, clustering
of the quantum measure, or a physical Hamiltonian gap.
\end{assessment}

\begin{decision}[Next fluctuation contribution]
Study the fluctuation determinant around this moving
minimum, including its dependence through \(z_*(c)\).
Separate the finite-dimensional quadratic determinant
from the exact interacting fibre integral and keep
the non-Gaussian remainder explicit.
Only a controlled remainder can turn the moving
background into a blocked-action estimate.
\end{decision}

\begin{firewall}[Position after V81]
The moving minimum and its minimized Hessian have
volume-uniform exponential spatial response bounds.
Local gauge compatibility and retained zero modes
are explicit.  The full interacting fibre integral,
scale iteration and Yang--Mills mass gap remain
unproved.  Only V81 is active; prior bodies and
archives remain.  Next companion audit: V85.
\end{firewall}

% END APPENDED POST-TENTH-ARCHIVE V81 MATERIAL

% =============================================================================
\section{V82: determinant expansion and the exact interacting remainder}
\label{sec:v82-fluctuations}
% =============================================================================

The V81 predecessor has SHA--256
\texttt{3E06CC81F5E729029FD5C67AB1247F226D99B2F3264B976AEC197F574AA73BC8}.
Its body is retained except for the version label.
Let \(N=45dn\) be the scalar dimension of the local
fibre coordinates, and keep \(hI\leq H_c\leq LI\)
from V81.  All results below concern the declared
small-field chart, not the whole compact fibre.

\begin{theorem}[Convergent finite-range determinant expansion]
\label{thm:v82-det}
Choose \(\Lambda>L\), \(T_c=I-H_c/\Lambda\),
and \(q=1-h/\Lambda<1\).  Then
\[
 \log\det H_c
 =N\log\Lambda-\sum_{j=1}^\infty
                              \frac{\Tr T_c^j}{j}. \tag{86.1}
\]
Truncation after order \(J\) has absolute error at most
\[
 \frac{Nq^{J+1}}{(J+1)(1-q)}.                    \tag{86.2}
\]
The diagonal contribution \((T_c^j)_{ii}\) uses only
matrix entries within distance \(jR\) of \(i\), where
\(R\) is the fixed range of \(H_c\).
\end{theorem}
\begin{proof}
Every eigenvalue of \(T_c\) lies in \([0,q]\).
Apply the absolutely convergent scalar series
\(\log(1-t)=-\sum_{j\geq1}t^j/j\) to these eigenvalues
and sum.  The tail is at most
\(N\sum_{j>J}q^j/j\), bounded by (86.2).
Matrix multiplication expresses a diagonal power
as a sum over closed sequences of \(j\) finite-range
steps, proving the locality statement.
The entries themselves are evaluated at \(z_*(c)\);
their dependence on distant coarse fields is not
strictly finite range, and is instead governed by V81.
\end{proof}

Choose a convex coordinate box
\(\mathcal D=\{z:\|z\|_\infty<r\}\) inside the common
chart, with \(\|z_*(c)\|_\infty\leq r/2\).
Write \(\ell(z,c)\) as in V80, and define
\[
 Z_{\rm loc}(c)=\int_{\mathcal D}
                            e^{-\beta f(z,c)+\ell(z,c)}dz .
 \]
Any fixed linear coordinate-volume constant is understood
to be included consistently in \(\ell\) or factored out.
Set \(G(c)=f(z_*(c),c)\), \(\ell_*=\ell(z_*(c),c)\),
\[
 V_c(u)=f(z_*+u,c)-G(c)-\tfrac12u^TH_cu,\qquad
 L_c(u)=\ell(z_*+u,c)-\ell_*.
 \]

\begin{theorem}[Exact Gaussian representation, with cutoff retained]
\label{thm:v82-exact}
For the centered Gaussian law \(\gamma_{\beta,c}\)
with covariance \(\beta^{-1}H_c^{-1}\),
\[
 \begin{split}
 Z_{\rm loc}(c)={}&
 e^{-\beta G(c)+\ell_*}
 (2\pi/\beta)^{N/2}(\det H_c)^{-1/2}\\
 &\times
 \mathbb E_{\gamma_{\beta,c}}\!
 \left[\mathbf1_{\{z_*+u\in\mathcal D\}}
                       e^{-\beta V_c(u)+L_c(u)}\right].
 \end{split}                                    \tag{86.3}
\]
The integrand in the expectation is defined to be
zero outside the chart; no extension of its remainder
there is assumed.
\end{theorem}
\begin{proof}
Translate the integration variable by \(z_*(c)\),
insert the definitions of \(V_c,L_c\), and divide
and multiply by the normalized Gaussian density.
All factors are finite and positive on the finite
chart.  No asymptotic approximation is involved.
\end{proof}

\begin{proposition}[Uniform local remainder bounds and their volume scaling]
\label{prop:v82-remainder}
For constants independent of \(c,n,\beta\),
on \(z_*+u\in\mathcal D\),
\[
 |V_c(u)|\leq C\sum_{i=1}^N|u_i|^3
                \leq C\|u\|_\infty\|u\|_2^2,    \tag{86.4}
\]
\[
 |L_c(u)|\leq C\sum_{i=1}^N|u_i|,\qquad
 |L_c(u)-\nabla\ell_*\cdot u|\leq C\|u\|_2^2.
                                                        \tag{86.5}
\]
Consequently
\[
 \mathbb E_{\gamma_{\beta,c}}
 [\mathbf1_{\{z_*+u\in\mathcal D\}}
                (\beta|V_c(u)|+|L_c(u)|)]
                       \leq C N\beta^{-1/2},    \tag{86.6}
\]
and
\[
 \gamma_{\beta,c}(z_*+u\notin\mathcal D)
                  \leq2N e^{-\beta h r^2/8}.    \tag{86.7}
\]
\end{proposition}
\begin{proof}
Expand each finite-stencil action term to second
order about \(z_*\).  Uniform third derivatives
bound its remainder by a constant times the cube
of the Euclidean norm of \(u\) on that stencil.
Each stencil has bounded size; each coordinate
occurs in boundedly many stencils.  Summing and
using equivalence of norms in that fixed size
gives the first inequality of (86.4).
Stationarity eliminates the total linear term.
The finite-stencil first and second derivative
bounds for the density similarly give (86.5).

Every scalar Gaussian variance is at most
\((\beta h)^{-1}\).  Its first absolute moment
is bounded by \(C\beta^{-1/2}\), and its third
by \(C\beta^{-3/2}\).  Apply these bounds to the
nonnegative upper sums in (86.4)--(86.5);
the chart indicator may be dropped from those sums.
This proves (86.6), without evaluating \(V_c,L_c\)
outside their domain.
Leaving the box requires some \(|u_i|\geq r/2\).
The scalar Gaussian tail bound and a union bound
give (86.7).
\end{proof}

\begin{assessment}[What cannot be concluded from these estimates]
The last factor of (86.3) is not proved close to one
uniformly in volume.  An \(L^1\) bound on its exponent
does not bound its exponential moment, and the factor
\(N\) in (86.6)--(86.7) cannot be silently discarded.
The determinant is therefore a well-controlled
quadratic factor, not yet a controlled replacement
for the interacting fibre integral.

The exact blocked action contains the logarithm of
that expectation as well as \(\beta G\), the density
at the minimum and the determinant.  Controlling
its connected contributions per region, including
the cutoff, is the missing step.  Even a uniform
per-volume free-energy error would not by itself
control the coarse derivatives or prove a physical gap.
\end{assessment}

\begin{decision}[Next connected estimate]
Test whether a local interpolation between the
quadratic and exact potentials retains enough convexity
and spatial covariance control to bound the logarithm
of the expectation per cell.  Keep the finite-domain
boundary term explicit.  Do not infer an exponential
moment bound from (86.6), or treat a small fixed-volume
error as a uniform interaction estimate.
\end{decision}

\begin{firewall}[Position after V82]
The determinant has an explicit convergent expansion,
and the local interacting integral has an exact
Gaussian representation with quantified elementary
remainder bounds.  A volume-uniform connected remainder,
compact-fibre completion and the full mass gap remain
unproved.  Only V82 is active; earlier bodies and
archives remain.  Next companion audit: V85.
\end{firewall}

% END APPENDED POST-TENTH-ARCHIVE V82 MATERIAL

% =============================================================================
\section{V83: a one-sided fluctuation free-energy bound with the cutoff retained}
\label{sec:v83-lower}
% =============================================================================

The V82 predecessor has SHA--256
\texttt{3A95E04FFC115400A40E972CBBDB277D614439CD6027705A39CBFBE7035F4B5D}.
Its body is retained except for the version label.
We improve the denominator estimate by controlling
the logarithm of the exact fluctuation factor from
below.  This is explicitly one-sided.

\begin{record}[External Gaussian inequality used]
We use the Gaussian rectangle inequality of
Z.~\v{S}id\'ak, \emph{Rectangular Confidence Regions
for the Means of Multivariate Normal Distributions},
JASA 62 (1967), 626--633, Corollary 1:
for a centered Gaussian vector \(Y\) and positive \(a_i\),
\[
 \mathbb P(\,|Y_i|\leq a_i\ \hbox{for all }i\,)
           \geq\prod_i\mathbb P(|Y_i|\leq a_i).  \tag{87.1}
\]
The original paper was checked at
\href{https://www.math.utoronto.ca/mccann/assignments/477/Sidak67.pdf}
{the primary-paper PDF};
DOI \href{https://doi.org/10.1080/01621459.1967.10482935}
{10.1080/01621459.1967.10482935}.
This theorem allows arbitrary correlations.  Its
application below is to a positive-definite centered
Gaussian and a symmetric coordinate rectangle;
no independence hypothesis is introduced.
\end{record}

Let \(\mathcal R_{\beta,c}\) denote the expectation
on the second line of (86.3), and keep the chart
margin \(\|z_*(c)\|_\infty\leq r/2\).
Take constants \(C_V,C_\ell\) from V82 such that
\[
 V_c(u)\leq C_V\|u\|_\infty\|u\|_2^2,\qquad
 L_c(u)\geq g_c\cdot u-C_\ell\|u\|_2^2,
 \quad g_c=\nabla\ell_* .
 \]
Both bounds hold throughout the shifted chart.

\begin{theorem}[Explicit lower bound for the full local fluctuation factor]
\label{thm:v83-lower}
For \(0<a\leq r/2\), define
\[
 \epsilon_a=2C_Va+2C_\ell/\beta,\qquad
             \widetilde H=H_c+\epsilon_a I .
 \]
If \(p_a=2e^{-\beta h a^2/2}<1\), then
\[
 \frac1N\log\mathcal R_{\beta,c}
 \geq-\frac{\epsilon_a}{2h}+\log(1-p_a).        \tag{87.2}
\]
This is uniform in volume and in the small coarse chart.
\end{theorem}
\begin{proof}
Restrict the positive expectation to the symmetric
cube \(\|u\|_\infty\leq a\), which is contained in
the shifted chart.  Combining the Gaussian density
with the two remainder bounds gives
\[
 \mathcal R_{\beta,c}\geq
 \left(\frac{\det H_c}{\det\widetilde H}\right)^{1/2}
 \mathbb E_{\widetilde\gamma}
       [\mathbf1_{\{\|u\|_\infty\leq a\}}e^{g_c\cdot u}],
                                                        \tag{87.3}
\]
where \(\widetilde\gamma\) has covariance
\(\beta^{-1}\widetilde H^{-1}\).
Its density and the cube are invariant under
\(u\mapsto-u\).  Averaging the integrand at \(u\)
and \(-u\) replaces the exponential by
\(\cosh(g_c\cdot u)\geq1\).
This step does not require a small norm for the
extensive vector \(g_c\).

Each scalar variance is at most \(1/(\beta h)\).
The scalar Gaussian tail bound and (87.1) give
\[
 \widetilde\gamma(\|u\|_\infty\leq a)
                         \geq(1-p_a)^N.
 \]
Finally
\(\log\det\widetilde H-\log\det H_c
=\sum_i\log(1+\epsilon_a/\lambda_i(H_c))
\leq N\epsilon_a/h\).
Take logarithms in (87.3) and divide by \(N\).
\end{proof}

\begin{corollary}[A vanishing one-sided error per fibre coordinate]
\label{cor:v83-rate}
For all sufficiently large \(\beta\), independently
of volume, choose
\[
 a_\beta=\sqrt{\frac{8\log\beta}{\beta h}}\leq r/2.
 \]
Then for a volume-independent constant \(C\),
\[
 \log\mathcal R_{\beta,c}
 \geq-CN\left(\sqrt{\frac{\log\beta}{\beta}}
                       +\beta^{-1}+\beta^{-4}\right).
                                                        \tag{87.4}
\]
Thus, with the Gaussian prefactor
\[
 Z_{\rm quad}(c)=e^{-\beta G(c)+\ell_*}
                   (2\pi/\beta)^{N/2}(\det H_c)^{-1/2},
 \]
the exact local integral satisfies
\[
 Z_{\rm loc}(c)\geq Z_{\rm quad}(c)
 \exp\left[-CN\left(\sqrt{\frac{\log\beta}{\beta}}
                       +\beta^{-1}+\beta^{-4}\right)\right].
                                                        \tag{87.5}
\]
\end{corollary}
\begin{proof}
The chosen radius gives \(p_a=2\beta^{-4}\).
For sufficiently large \(\beta\), use
\(\log(1-2\beta^{-4})\geq-4\beta^{-4}\)
in (87.2).  Substitute the displayed
\(\epsilon_{a_\beta}\), and then use (86.3).
\end{proof}

\begin{assessment}[Why this avoids the previous volume-loss argument]
We did not subtract a union-bound error from one.
The retained cube probability is bounded below by
a product whose logarithm has an explicit per-coordinate
bound even when \(N\) is arbitrarily large.
Nor did we exponentiate the first-moment estimate
of V82.  The proof uses a pointwise quadratic
majorant on a smaller cube and an exact Gaussian
change of precision.

The result is only a lower bound for
\(\mathcal R_{\beta,c}\), or equivalently an upper
bound for its negative logarithm.  It does not
bound \(\mathcal R_{\beta,c}\) from above, its
coarse derivatives, or its connected spatial terms.
It also concerns the local chart contribution,
not all configurations in the compact fibre.
It therefore cannot yet be used as a two-sided
effective-action approximation.
\end{assessment}

\begin{decision}[Next missing side]
Seek an upper bound for the local integral relative
to the same quadratic prefactor.  The part outside
the shrinking cube must be controlled within the
interacting convex measure, not discarded using
the unweighted Gaussian union bound.
Track whether the resulting estimate is per volume
only or also controls coarse derivatives and local
connected contributions; these are different targets.
\end{decision}

\begin{firewall}[Position after V83]
A vanishing one-sided fluctuation free-energy error
per fibre coordinate is proved, uniformly in volume,
with the cutoff retained.  Two-sided remainder and
connected-interaction control, compact completion,
and the full Yang--Mills mass gap remain unproved.
Only V83 is active; prior bodies and archives remain.
Next companion audit: V85.
\end{firewall}

% END APPENDED POST-TENTH-ARCHIVE V83 MATERIAL

% =============================================================================
\section{V84: two-sided local free-energy control by Gaussian product bounds}
\label{sec:v84-upper}
% =============================================================================

The V83 predecessor has SHA--256
\texttt{6D6921B83E17068C7ACCC937ABABDC32A90F8734A24F66562D2ECC32D1FF5A2A}.
Its body is retained except for the version label.
Here we choose the fixed local box radius \(r\)
sufficiently small, and shrink the coarse chart so
\(\|z_*(c)\|_\infty\leq r/2\) as before.
All references to \(Z_{\rm loc}\) and
\(\mathcal R_{\beta,c}\) in this section use this
chosen box.  No bound for a larger chart or the
full compact fibre is inferred by renaming it.

\begin{record}[External Gaussian product inequality]
We use Theorem 1(i) of W.-K.~Chen, N.~Dafnis and
G.~Paouris, \emph{Improved H\"older and reverse
H\"older inequalities for Gaussian random vectors},
\href{https://arxiv.org/pdf/1306.2410}{arXiv:1306.2410}.
For a centered Gaussian scalar vector with covariance
\(\Sigma\), if \(\Sigma\leq p\,\operatorname{diag}(\Sigma_{ii})\),
then for nonnegative measurable functions \(F_i\),
\[
 \mathbb E\prod_iF_i(X_i)
       \leq\prod_i(\mathbb E F_i(X_i)^p)^{1/p}.  \tag{88.1}
\]
The primary-paper theorem and its covariance condition
were checked.  This is not ordinary independence
or H\"older with an exponent growing with dimension.
\end{record}

\begin{lemma}[The required exponent is volume independent]
For \(\Sigma=\beta^{-1}H_c^{-1}\), with
\(hI\leq H_c\leq LI\), one may take \(p=L/h\geq1\)
in (88.1).
\end{lemma}
\begin{proof}
\(\Sigma\leq(\beta h)^{-1}I\), whereas every
\(\Sigma_{ii}\geq(\beta L)^{-1}\).
Thus \(p\,\operatorname{diag}(\Sigma_{ii})
\geq(\beta h)^{-1}I\geq\Sigma\).
\end{proof}

\begin{lemma}[Uniform one-dimensional truncated exponential estimate]
\label{lem:v84-scalar}
Fix \(C,h,L,p>0\).  For a sufficiently small
\(a>0\), a centered scalar Gaussian \(X\) with
\((\beta L)^{-1}\leq\operatorname{Var}X
\leq(\beta h)^{-1}\) satisfies
\[
 \mathbb E\left[
  \mathbf1_{\{|X|\leq a\}}
  e^{pC(\beta|X|^3+|X|)}\right]
                    \leq1+C'\beta^{-1/2},
                    \qquad\beta\geq1,           \tag{88.2}
\]
where \(C'\) is independent of \(\beta\) and the
particular variance.
\end{lemma}
\begin{proof}
Choose \(pCa\leq h/8\).  On the interval in question
the cubic term is at most \(\beta h x^2/8\).
Also \(pC|x|\leq\beta h x^2/8+C_1/\beta\)
by the elementary quadratic inequality.
The Gaussian density is at most
\(\sqrt{\beta L/(2\pi)}e^{-\beta h x^2/2}\).
For \(y=pC(\beta|x|^3+|x|)\geq0\), use
\(e^y-1\leq y e^y\).  Multiplying by the density
and using the preceding bounds, the excess over
the interval probability is at most a constant
times
\[
 \sqrt\beta\int_{\mathbb R}
    (\beta|x|^3+|x|)e^{-\beta h x^2/4}\,dx
                         \leq C'\beta^{-1/2}.
 \]
The interval probability itself is at most one.
\end{proof}

\begin{theorem}[Upper fluctuation bound on the fixed small chart]
\label{thm:v84-upper}
Choose \(r\) so \(2r\) is an allowed \(a\) in the
preceding lemma, with the uniform remainder constant
of V82 and \(p=L/h\).  Then
\[
 \log\mathcal R_{\beta,c}\leq C N\beta^{-1/2},
                         \qquad\beta\geq1,      \tag{88.3}
\]
uniformly in volume and the resulting small coarse chart.
\end{theorem}
\begin{proof}
On the shifted chart, V82 gives the pointwise bound
\[
 -\beta V_c(u)+L_c(u)
                 \leq C\sum_i(\beta|u_i|^3+|u_i|).
 \]
Every point there has \(|u_i|\leq2r\).
Extend the original integrand by zero off its
domain.  It is bounded above everywhere by
\(\prod_iF_i(u_i)\), where
\[
 F_i(x)=\mathbf1_{\{|x|\leq2r\}}
                             e^{C(\beta|x|^3+|x|)}.
 \]
Use (88.1), the covariance lemma, and (88.2).
This gives
\(\mathcal R_{\beta,c}\leq
(1+C'\beta^{-1/2})^{N/p}\).
Taking logarithms proves (88.3).
The cubic exponential is never extended without
its coordinate cutoff.
\end{proof}

\begin{corollary}[Two-sided error per fibre coordinate]
For sufficiently large \(\beta\), on this same chart,
\[
 -C_1N\sqrt{\frac{\log\beta}{\beta}}
 \leq\log\frac{Z_{\rm loc}(c)}{Z_{\rm quad}(c)}
 \leq C_2N\beta^{-1/2}.                         \tag{88.4}
\]
In particular the absolute logarithmic error divided
by \(N\) tends to zero uniformly in volume and \(c\).
\end{corollary}
\begin{proof}
Combine (88.3) with V83's lower bound, absorbing
its smaller terms for large \(\beta\), and use
the exact identity \(Z_{\rm loc}/Z_{\rm quad}
=\mathcal R_{\beta,c}\).
The radius \(a_\beta\) in that lower proof
eventually lies inside the fixed smaller chart.
\end{proof}

\begin{assessment}[Precisely the approximation now justified]
The quadratic prefactor approximates the local
conditional free energy per coordinate, with a
two-sided, volume-uniform error tending to zero.
This does not say the ratio of partition functions
tends to one uniformly in volume: the exponent
still scales with \(N\).
The theorem supplies no bounds on coarse derivatives
of that error, no connected spatial decomposition,
and no estimate for configurations outside the
chosen compact-fibre chart.

Uniform bounds on function values cannot simply
be differentiated.  Accordingly (88.4) is not
yet a controlled marginal coupling calculation,
an RG contraction, or a mass-gap theorem.
\end{assessment}

\begin{decision}[Next audit and response control]
At V85 perform the compulsory companion audit.
Then examine derivatives of the local partition
function using its exact score and covariance
identities, retaining the fixed-domain convention.
Determine which volume-uniform local response
bounds can be proved from convexity and finite
range, rather than differentiating (88.4).
\end{decision}

\begin{firewall}[Position after V84]
Two-sided local fluctuation free-energy control
per coordinate is proved on a fixed small chart.
Connected response, compact-fibre completion,
scale iteration and the full Yang--Mills mass gap
remain unproved.  Only V84 is active; earlier
bodies and archives remain.  Next audit: V85.
\end{firewall}

% END APPENDED POST-TENTH-ARCHIVE V84 MATERIAL

% =============================================================================
\section{V85: exact interacting score bounds on the fixed convex chart}
\label{sec:v85-scores}
% =============================================================================

The V84 predecessor has SHA--256
\texttt{634E7B4590D4E617A69E0EACD284A56B3298E2F8D7432F0E4F7E561A92344766}.
Its body is retained except for the version label.
This is the required companion-audit version.
We derive response bounds from the local measure,
not by differentiating the free-energy error estimate.
Use the fixed box \(\mathcal D\) and define
\[
 V_{\beta,c}(z)=\beta f(z,c)-\ell(z,c),\qquad
 d\pi_{\beta,c}=Z_{\rm loc}(c)^{-1}
                       e^{-V_{\beta,c}(z)}dz,\qquad
 W_\beta(c)=-\log Z_{\rm loc}(c).
 \]
For sufficiently large \(\beta\), the preceding
local convexity estimates give
\(D_z^2V_{\beta,c}\geq m\beta I\) on the box,
with \(m>0\) independent of volume and \(c\).

\begin{lemma}[Poincar\'e inequality with the box boundary retained]
\label{lem:v85-poincare}
For every smooth real \(F\) on the box,
\[
 \operatorname{Var}_{\pi_{\beta,c}}F
       \leq(m\beta)^{-1}
                   \mathbb E_{\pi_{\beta,c}}|\nabla_zF|^2.
                                                        \tag{89.1}
\]
\end{lemma}
\begin{proof}
Here is the boundary justification.  On a smooth
bounded convex domain with density \(e^{-V}\),
use the nonnegative weighted Neumann operator
\(\mathcal L=-\Delta+\nabla V\cdot\nabla\).
Integration by parts and the differentiated
Neumann condition give
\[
 \int(\mathcal Lu)^2\,d\pi
 =\int\big(\|\nabla^2u\|_{\rm HS}^2+
       \nabla u^T(D^2V)\nabla u\big)d\pi
       +\int_{\partial\mathcal D}
          {\rm II}(\nabla_{\rm tan}u,\nabla_{\rm tan}u)
                                    \frac{e^{-V}}{Z}\,dS .
 \]
The second fundamental form with the convex
outward-normal convention is nonnegative.
For a nonconstant Neumann eigenfunction of
eigenvalue \(\lambda\), this identity gives
\(\lambda^2\|u\|_2^2\geq m\beta\lambda\|u\|_2^2\),
hence \(\lambda\geq m\beta\).
The spectral variational characterization proves
the variance bound on the smooth domain.
Approximate the box by smooth convex domains inside
it with volumes converging to that of the box.
The potential is smooth on a neighbourhood of the
closed chart, the Hessian bound is unchanged, and
integrals of each fixed smooth \(F\), \(F^2\) and
\(|\nabla F|^2\) converge by bounded convergence
at each finite dimension.  Passing the inequality
to the limit proves (89.1).  Its constant was
independent of the approximating domains and volume.
\end{proof}

\begin{theorem}[Fixed-domain score identities]
\label{thm:v85-identities}
For constant coarse-coordinate directions \(a,b\), let
\(S_a=D_cV_{\beta,c}[a]\) and
\(T_{ab}=D_c^2V_{\beta,c}[a,b]\).  Then
\[
 D_cW_\beta[a]=\mathbb E S_a,\qquad
 D_c^2W_\beta[a,b]
       =\mathbb E T_{ab}
                         -\operatorname{Cov}(S_a,S_b).
                                                        \tag{89.2}
\]
\end{theorem}
\begin{proof}
The coordinate box is independent of \(c\).
At fixed finite volume all required derivatives
are bounded on a smaller closed parameter chart,
so differentiation under the integral is valid.
Differentiating the normalized expectation of
\(S_a\) gives its derivative expectation minus
its covariance with \(S_b\), proving both formulas.
There is no moving-domain boundary term in these
coordinates.  If one instead translates by \(z_*(c)\),
the translated domain must also be differentiated;
(89.2) has not discarded that dependence.
\end{proof}

\begin{theorem}[Interacting response bounds independent of volume]
\label{thm:v85-response}
For constants independent of volume, \(c\), and
sufficiently large \(\beta\),
\[
 \operatorname{Var}(S_a)\leq C\beta\|a\|_2^2,
 \qquad
 |D_c^2W_\beta[a,b]|
                    \leq C\beta\|a\|_2\|b\|_2. \tag{89.3}
\]
For a single normalized coarse coordinate \(e\),
\(|\partial_eW_\beta|\leq C\beta\).
\end{theorem}
\begin{proof}
The functions \(f,\ell\) are finite-stencil sums
in the independent \(z,c\) coordinates, with
uniform local derivatives.  The mixed derivative
matrix \(D_zD_cV_{\beta,c}\) has bounded range
and operator norm at most \(C\beta\).
Likewise \(D_c^2V_{\beta,c}\) has operator norm
at most \(C\beta\).  These statements follow
from bounded row and column sums; they do not
bound an extensive derivative by a constant.
Thus pointwise
\[
 |\nabla_z S_a|\leq C\beta\|a\|_2,\qquad
 |T_{ab}|\leq C\beta\|a\|_2\|b\|_2 .
 \]
Use (89.1) for the variance.  Cauchy--Schwarz
bounds the covariance in (89.2) by
\(C\beta\|a\|_2\|b\|_2\), proving (89.3).
A derivative in one coarse coordinate involves
only a bounded number of local terms, each
bounded by \(C\beta\), proving the final assertion.
\end{proof}

\begin{assessment}[The covariance has not been approximated away]
The Hessian formula is for the interacting local
measure.  Its covariance is not replaced by the
Gaussian covariance or the inverse Hessian at the
minimum.  The \(O(\beta)\) estimate is on the
natural scale of the classical coarse Hessian;
it is not a vanishing derivative error relative
to the quadratic approximation.
It proves neither positivity of the coarse Hessian
nor exponential decay of its entries.

The box itself is not a gauge-invariant compact
domain.  These fixed-chart identities do not
establish a global Ward identity for its partition
function or for the full compact blocked measure.
\end{assessment}

\begin{audit}[Required V85 companion check]
SM remains V72 with SHA--256
\texttt{F44F9745D43379E1672C5550411B58E97195A4C9278592D221DF5D3F644DC1F5};
NS remains V26 with SHA--256
\texttt{82C887F8C2C69E4F28FAD7C3DC8DE5B9099CB5A4BF431EBA1FFF3E91935978AD}.
Their active conclusions and hashes are unchanged.
Neither supplies the interacting YM spatial covariance
or compact-chart completion needed here.
The parallel YM active note remains V5.
Next companion audit and ten-version review: V90.
\end{audit}

\begin{decision}[Next spatial covariance question]
Derive a spatially weighted covariance estimate
for this strongly convex interacting measure,
including the product-box boundary conditions.
The scalar Poincar\'e bound alone does not encode
distance.  A valid estimate must use finite-range
Hessian coupling and must not replace the
interacting response operator by \(H_c^{-1}\).
\end{decision}

\begin{firewall}[Position after V85]
The fixed-chart interacting score and covariance
give a volume-uniform coarse-response operator bound.
Spatial response, compact-fibre completion, RG
iteration and the full Yang--Mills mass gap remain
unproved.  Only V85 is active; prior bodies and
archives remain.
\end{firewall}

% END APPENDED POST-TENTH-ARCHIVE V85 MATERIAL

% =============================================================================
\section{V86: interacting spatial covariance on the product box}
\label{sec:v86-covariance}
% =============================================================================

The V85 predecessor has SHA--256
\texttt{5E8A5119747BF15DAC14A5B7DEE9D48140342394B5BFF454F303A9033DA10195}.
Its body is retained except for the active-version label.
It remains the authoritative active predecessor.
We now add distance to its interacting covariance estimate.
Keep the measure \(\pi=\pi_{\beta,c}\) on
\(\mathcal D=(-r,r)^N\), with potential \(V=V_{\beta,c}\).
There are uniform constants \(m,K,R>0\) such that
\[
 D^2V\geq m\beta I,\qquad
 \sup_z\max_i\sum_j|V_{ij}(z)|\leq K\beta,\qquad
 V_{ij}=0\quad\text{if }{\rm dist}(i,j)>R.        \tag{90.1}
\]
The same column bound holds by symmetry.
These are the previously proved local-derivative and
convexity bounds in independent fibre coordinates.
Distance refers to their associated cell positions.

\begin{lemma}[Interacting covariance representation with flat-face conditions]
\label{lem:v86-representation}
Let \(\mathcal L=-\Delta+\nabla V\cdot\nabla\)
be the weighted scalar Neumann operator.
For vector fields define the closed quadratic form
\[
 \mathfrak a(v,w)=
 \mathbb E_\pi\left[\sum_i\nabla v_i\cdot\nabla w_i
                         +\sum_{ij}v_iV_{ij}w_j\right],
                                                        \tag{90.2}
\]
with \(v_i=0\) on the two faces \(z_i=\pm r\);
there are natural Neumann conditions on its other faces.
Let \(\mathcal A\) be its associated positive operator.
Then \(\mathcal A\geq m\beta I\), and for smooth \(F,G\),
\[
 \operatorname{Cov}_\pi(F,G)=
       \langle\nabla F,\mathcal A^{-1}\nabla G\rangle_{L^2(\pi)}.
                                                        \tag{90.3}
\]
\end{lemma}
\begin{proof}
The lower bound follows immediately from (90.1)
and the nonnegative derivative term in (90.2).
Solve \(\mathcal Lu=G-\mathbb E G\) with zero mean
and Neumann conditions; V85's Poincar\'e inequality
gives existence and uniqueness in the energy space.
Differentiating its equation gives, in the weak sense,
\[
 (\mathcal L\otimes I+D^2V)\nabla u=\nabla G.
 \]
On a face \(z_i=\pm r\), the scalar Neumann condition
is \(u_i=0\); on a different face \(z_j=\pm r\),
its tangential derivative gives \(\partial_j u_i=0\).
These are precisely the mixed conditions encoded
by (90.2).  Thus \(\nabla u=\mathcal A^{-1}\nabla G\).

The differentiation can equivalently be performed
by tangential difference quotients in the scalar
weak equation and integration by parts against
vector tests vanishing on their own coordinate
faces.  All face terms then vanish as above.
Rectangular Neumann regularity, or reflection across
the flat faces followed by difference quotients,
gives \(\nabla u\) in the form domain for smooth data.
Edges and corners have zero face measure and impose
no extra term in the resulting weak identity.
Approximation in the energy norm extends it to
the solution just constructed.
Finally scalar Neumann integration by parts gives
\(\operatorname{Cov}(F,G)=
\langle\nabla F,\nabla u\rangle\), proving (90.3).
The operator contains the random field-dependent
matrix \(D^2V(z)\), not just its value at a minimum.
\end{proof}

\begin{theorem}[Volume-uniform exponential covariance bound]
\label{thm:v86-decay}
Suppose the coordinate supports of \(\nabla F\)
and \(\nabla G\) are contained in sets \(A,B\).
There is \(\alpha>0\), independent of volume,
\(\beta\) and the small coarse field, such that
\[
 |\operatorname{Cov}_\pi(F,G)|
 \leq\frac{2}{m\beta}
       e^{-\alpha\,{\rm dist}(A,B)}
       \|\nabla F\|_{L^2(\pi)}
       \|\nabla G\|_{L^2(\pi)}.                 \tag{90.4}
\]
\end{theorem}
\begin{proof}
Let \(w_i=e^{\alpha\,{\rm dist}(i,B)}\), and let
\(W\) multiply vector component \(i\) by \(w_i\).
These weights are constant in the integration
variables \(z\).  They commute with the differential
part of \(\mathcal A\) and preserve every mixed
boundary condition.  Only the matrix term changes:
\[
 W\mathcal AW^{-1}=\mathcal A+D_\alpha,\qquad
 (D_\alpha)_{ij}=(w_i/w_j-1)V_{ij}.
 \]
For nonzero off-diagonal entries, the distance
difference is at most \(R\).  Both row and column
absolute sums of \(D_\alpha\) are bounded by
\(K\beta(e^{\alpha R}-1)\); the diagonal entries
are zero.  The row-column norm estimate therefore
gives the same bound on its \(L^2(\pi)\) operator norm.
Choose \(\alpha>0\) so \(K(e^{\alpha R}-1)\leq m/2\).
A bounded-perturbation Neumann series about
\(\mathcal A\geq m\beta I\) yields
\[
 \|(W\mathcal AW^{-1})^{-1}\|\leq2/(m\beta).
 \]
At each finite volume the weights and inverses
are bounded; the resulting estimate does not
depend on their largest value or on volume.

For \(v=\mathcal A^{-1}\nabla G\), the source
is supported on \(B\), where the weights are one.
Hence \(\|Wv\|_2\leq2\|\nabla G\|_2/(m\beta)\).
On \(A\), \(w_i\geq e^{\alpha\,{\rm dist}(A,B)}\).
Use (90.3) and Cauchy--Schwarz to conclude (90.4).
This proof has kept the box boundary conditions
through the conjugation.
\end{proof}

\begin{corollary}[Exponential spatial Hessian bound for the interacting local action]
\label{cor:v86-coarse}
For individual coarse scalar coordinates \(e,f\),
\[
 |\partial_e\partial_f W_\beta(c)|
              \leq C\beta e^{-\alpha'\,{\rm dist}(e,f)}
                                                        \tag{90.5}
\]
uniformly in volume and the small coarse chart,
for some \(\alpha'>0\).
\end{corollary}
\begin{proof}
The score \(S_e=\partial_e V\) has a gradient
supported on a fixed-size fibre stencil near \(e\),
with \(L^2\) norm at most \(C\beta\), by V85.
Apply (90.4) to the two scores.  Their support
distance differs from the coarse-edge distance
by at most a fixed constant.
The direct term \(\mathbb E\partial_e\partial_f V\)
in (89.2) vanishes outside a fixed range and is
bounded by \(C\beta\) inside that range.
Combining the two terms proves (90.5).
\end{proof}

\begin{assessment}[This is an interacting, but still local-chart, result]
Equation (90.5) strengthens V85 by controlling
distance, using the actual interacting response
operator and its mixed boundary conditions.
It is not obtained from the Gaussian covariance
\(\beta^{-1}H_c^{-1}\).
It does not show that derivatives of the fluctuation
error vanish as \(\beta\) grows: the estimate
remains \(O(\beta)\) on the classical response scale.
Nor does it establish a global gauge identity for
the fixed box or control the rest of the compact fibre.

Exponential dependence on cell distance at one
blocking step is not a positive continuum mass
in physical units.  The relevant constants must
survive scale iteration and reconstruction, neither
of which has been proved here.
\end{assessment}

\begin{decision}[Next non-Gaussian derivative comparison]
Compare the interacting response in (89.2) with
the moving-minimum Schur response using fluctuation
moments and the exact covariance representation.
Seek a small error on the natural \(\beta\) scale,
not merely two separate \(O(\beta)\) bounds.
Retain the fixed-chart boundary and the separate
compact-fibre complement in that comparison.
\end{decision}

\begin{firewall}[Position after V86]
The interacting local measure has a volume-uniform
exponential spatial covariance bound and hence
an exponentially decaying coarse Hessian.
Derivative approximation, compact completion,
scale iteration and the full Yang--Mills mass gap
remain unproved.  Only V86 is active; prior
bodies and archives remain.  Next audit and
ten-version review: V90.
\end{firewall}

% END APPENDED POST-TENTH-ARCHIVE V86 MATERIAL

% =============================================================================
\section{V87: mean-square fluctuations and an intensive first-derivative comparison}
\label{sec:v87-moments}
% =============================================================================

The V86 predecessor has SHA--256
\texttt{73F37AA16A41499D9D4602BF11AF52B34B40592B8F5C855D8C82C902D6F50D22}.
Its body is retained except for the active-version label.
We compare actual derivatives rather than differentiate
the free-energy bounds.  The conclusion is in a
volume-normalized Euclidean norm, not a pointwise
bound at each coarse link.

Keep the fixed convex box and write
\(b=z_{\beta,*}(c)\) for the measured local minimum
of V80, with the coarse chart and lower bound on
\(\beta\) chosen so \(b\) lies strictly inside.
For \(V=\beta f-\ell\), assume the established
bound \(D_z^2V\geq m\beta I\).
The scalar fibre dimension is \(N=45dn\).

\begin{theorem}[Boundary-signed second-moment estimate]
\label{thm:v87-moment}
Under the exact local interacting law,
\[
 \mathbb E_\pi\|z-b\|_2^2\leq\frac{N}{m\beta},
 \qquad
 \mathbb E_\pi\|z-z_*(c)\|_2^2\leq C N/\beta.   \tag{91.1}
\]
The constants are uniform in volume and the small
coarse chart.
\end{theorem}
\begin{proof}
Apply the divergence theorem to
\((z-b)e^{-V(z)}\) on the box.  Dividing by its
partition function gives
\[
 N-\mathbb E_\pi[(z-b)\cdot\nabla V(z)]
 =Z_{\rm loc}^{-1}\int_{\partial\mathcal D}
              (z-b)\cdot n\,e^{-V(z)}\,dS\geq0.
 \]
The sign holds because \(b\) is inside the convex
box.  Corners have zero face measure.
Since \(\nabla V(b)=0\), integration of the Hessian
along the segment from \(b\) to \(z\) yields
\((z-b)\cdot\nabla V(z)\geq m\beta\|z-b\|_2^2\).
This proves the first bound; the boundary flux
was not set to zero.
For the second, V80 gives
\(\|b-z_*\|_\infty\leq C\beta^{-1}\|c\|_\infty\).
Use \(\|x+y\|_2^2\leq2\|x\|_2^2+2\|y\|_2^2\),
the fixed coarse-chart radius, and \(\beta\geq1\).
\end{proof}

\begin{theorem}[Vanishing first-derivative error in normalized Euclidean norm]
\label{thm:v87-score}
For \(W_\beta=-\log Z_{\rm loc}\) and
\(G(c)=f(z_*(c),c)\),
\[
 \frac1{\sqrt N}
 \left\|\beta^{-1}\nabla_cW_\beta(c)-\nabla_cG(c)\right\|_2
                 \leq C(\beta^{-1/2}+\beta^{-1}).
                                                        \tag{91.2}
\]
\end{theorem}
\begin{proof}
The exact score identity and stationarity give
\[
 \beta^{-1}\nabla_cW_\beta-\nabla_cG
 =\mathbb E_\pi[f_c(z,c)-f_c(z_*,c)]
                          -\beta^{-1}\mathbb E_\pi\ell_c(z,c).
 \]
Bounded finite-stencil mixed derivatives imply
\(\|f_c(z,c)-f_c(z_*,c)\|_2\leq C\|z-z_*\|_2\).
All segments used here stay in the convex box.
Each component of \(\ell_c\) is bounded by a
volume-independent constant, and there are
\(4dn\) coarse scalar coordinates, so
\(\|\ell_c\|_2\leq C\sqrt N\).
Apply Jensen and Cauchy--Schwarz to the first
term and use (91.1).  Divide by \(\sqrt N\).
No Gaussian replacement of \(\pi\) is used.
\end{proof}

\begin{proposition}[Direct Hessian term, but not yet the full response]
The explicit second-derivative matrix has the bound
\[
 \frac1N\mathbb E_\pi
       \|f_{cc}(z,c)-f_{cc}(z_*,c)\|_{\rm HS}^2
                                      \leq C/\beta. \tag{91.3}
\]
\end{proposition}
\begin{proof}
Each nonzero matrix entry depends on a fixed
stencil and has bounded first derivatives in
those fibre variables.  Each fibre variable
appears in only a bounded number of such entries.
The mean-value theorem and bounded overlap give
\(\|f_{cc}(z,c)-f_{cc}(z_*,c)\|_{\rm HS}^2
\leq C\|z-z_*\|_2^2\).  Apply (91.1).
\end{proof}

\begin{assessment}[The norm and covariance qualifications matter]
Equation (91.2) gives a small intensive first-derivative
error.  It does not by itself bound the maximum
error at one coordinate uniformly in volume.
For example, a bounded normalized Euclidean norm
can hide an error concentrated at a few coordinates.

Likewise (91.3) controls only the direct term of
the exact Hessian identity.  The scaled covariance
\(\beta^{-1}\operatorname{Cov}(S_a,S_b)\) is
generally of order one and is responsible for
the Schur correction in a Gaussian limit.
It cannot be discarded on the strength of
the second-moment estimate.
The spatial bounds of V86 help with that comparison
but do not perform it automatically.
\end{assessment}

\begin{corollary}[An averaged fluctuation-density consequence]
For every \(t>0\),
\[
 \mathbb E_\pi\left[
 \frac1N\#\{i:|z_i-z_{*,i}|>t\}\right]
                                \leq C/(\beta t^2).
                                                        \tag{91.4}
\]
\end{corollary}
\begin{proof}
Bound each indicator by
\(|z_i-z_{*,i}|^2/t^2\), sum, and apply (91.1).
This is a bound on an expected fraction, not
a claim that no coordinate leaves that range.
\end{proof}

\begin{decision}[Next response comparison]
Compare the interacting vector response operator
of V86 with its frozen quadratic counterpart
using the fluctuation moments.  Account for
the mixed box boundary conditions: a constant
vector is not in every component's Dirichlet
form domain.  Any Gaussian-limit comparison
must control that boundary mismatch rather
than formally applying the operator to an
inadmissible constant test field.
\end{decision}

\begin{firewall}[Position after V87]
Exact local mean-square fluctuation bounds yield
a vanishing first-derivative error in a declared
intensive norm.  Full Hessian approximation,
compact completion and continuum scale control
remain incomplete; no Yang--Mills mass gap is proved.
Only V87 is active; previous bodies and archives
remain.  Next audit and review: V90.
\end{firewall}

% END APPENDED POST-TENTH-ARCHIVE V87 MATERIAL

% =============================================================================
\section{V88: boundary-compatible covariance and Hessian approximation}
\label{sec:v88-comparison}
% =============================================================================

The V87 predecessor has SHA--256
\texttt{78DF4CCFCA85637F259B7D18DCDD9ACF96EEF76E6118431509597651047ECD86}.
Its body is retained except for the active-version label.
All norms and expectations below refer to the fixed
local chart.  Let \(u=z-z_*(c)\), \(H=f_{zz}(z_*,c)\),
and \(M(z)=\beta^{-1}D_z^2V(z)\).
The normalized interacting vector operator is
\(\mathcal B=\mathcal A/\beta\), with form
\[
 b(v,w)=\beta^{-1}\mathbb E\sum_i\nabla v_i\cdot\nabla w_i
                          +\mathbb E v^TMw .
 \]
It has the mixed conditions of V86 and
\(b(v,v)\geq m\|v\|_{L^2(\pi)}^2\).

\begin{lemma}[Coefficient fluctuation and admissible trial fields]
\label{lem:v88-trials}
Uniformly in volume and the small coarse chart,
\[
 \mathbb E\|M-H\|_{\rm HS}^2\leq CN/\beta.       \tag{92.1}
\]
There is a diagonal matrix \(Q(z)\) such that every
column of \(QH^{-1}\) belongs to the mixed form domain,
and
\[
 \mathbb E\|I-Q\|_{\rm HS}^2\leq CN/\beta,\qquad
 \sum_j\|\nabla(QH^{-1}e_j)\|_{L^2(\pi)}^2\leq CN.
                                                        \tag{92.2}
\]
\end{lemma}
\begin{proof}
The finite-range Hessian coefficients of \(f\) are
locally Lipschitz with bounded overlap.  Consequently
\(\|f_{zz}(z,c)-H\|_{\rm HS}^2\leq C\|u\|_2^2\).
The density Hessian has squared Hilbert--Schmidt
norm at most \(CN\).  Apply V87's moment bound
to \(M=f_{zz}-\beta^{-1}\ell_{zz}\) to get (92.1).

Choose smooth scalar cutoffs \(\chi_i(z_i)\) equal
to one for \(|u_i|\leq r/8\), zero for
\(|u_i|\geq r/4\), and with bounded derivative
depending only on the fixed radius.  Since
\(|z_{*,i}|\leq r/2\), each cutoff vanishes near
its own coordinate faces.  Put \(Q_{ii}=\chi_i\).
Its \(i\)-th component is independent of the other
coordinates and satisfies their natural Neumann
conditions.  Moreover \(|1-\chi_i|^2\leq C|u_i|^2\),
giving the first part of (92.2).
The derivative bound and the uniform norm of
\(H^{-1}\) give the second, since
\(\|H^{-1}\|_{\rm HS}^2\leq CN\).
\end{proof}

\begin{theorem}[Covariance approximation in an intensive matrix norm]
\label{thm:v88-covariance}
For the actual interacting coordinate covariance,
\[
 \frac1{\sqrt N}
       \|\beta\operatorname{Cov}_\pi(z)-H^{-1}\|_{\rm HS}
                                      \leq C\beta^{-1/2}.
                                                        \tag{92.3}
 \]
\end{theorem}
\begin{proof}
Let \(w_j=\mathcal B^{-1}e_j\) and use the admissible
trial field \(v_j=QH^{-1}e_j\).  Its weak residual is
\[
 \langle e_j,w\rangle-b(v_j,w)
 =\langle F_j,w\rangle
        -\beta^{-1}\langle\nabla v_j,\nabla w\rangle,
 \quad F_j=(I-MQH^{-1})e_j .
 \]
Here
\(I-MQH^{-1}=(H-M)H^{-1}+M(I-Q)H^{-1}\).
The matrices \(M,H^{-1}\) have uniformly bounded
operator norm.  Thus (92.1)--(92.2) imply
\(\sum_j\|F_j\|_{L^2}^2\leq CN/\beta\).
The residual dual norm for the \(b\)-energy is at most
\[
 m^{-1/2}\|F_j\|_{L^2}
                  +\beta^{-1/2}\|\nabla v_j\|_{L^2}.
 \]
Coercivity and summation give
\(\sum_j\|w_j-v_j\|_{L^2}^2\leq CN/\beta\).

V86's representation for coordinate observables
identifies the \(j\)-th column of
\(\beta\operatorname{Cov}_\pi(z)\) with \(\mathbb E w_j\).
Jensen bounds the squared matrix error from
\(\mathbb E QH^{-1}\) by \(CN/\beta\).
The error between \(\mathbb E QH^{-1}\) and
\(H^{-1}\) has the same bound by (92.2).
Take square roots.  No inadmissible constant field
was used as a trial vector in the form domain.
\end{proof}

\begin{theorem}[Interacting coarse Hessian approaches the Schur Hessian]
\label{thm:v88-hessian}
For \(W_\beta=-\log Z_{\rm loc}\) and
\(G(c)=f(z_*(c),c)\),
\[
 \frac1{\sqrt N}
 \|\beta^{-1}D_c^2W_\beta(c)-D_c^2G(c)\|_{\rm HS}
                                      \leq C\beta^{-1/2}.
                                                        \tag{92.4}
 \]
\end{theorem}
\begin{proof}
Put \(A=f_{zc}(z_*,c)\).  The vector of scores has
the exact decomposition
\[
 S=\beta f_c(z_*,c)+Y+R,\qquad Y=\beta A^Tu,
 \quad R=\beta[f_c(z,c)-f_c(z_*,c)-A^Tu]-\ell_c(z,c).
 \]
Finite-range derivative bounds and (91.1) imply
\[
 \mathbb E\|D_zR\|_{\rm HS}^2\leq C\beta N .
 \]
Indeed the Taylor derivative difference is bounded
in squared Hilbert--Schmidt norm by \(C\|u\|_2^2\)
before multiplication by \(\beta^2\).
Poincar\'e therefore gives
\(\Tr\operatorname{Cov}(R)\leq CN\).
The pointwise operator norms of \(D_zR\) and
\(D_zY\) are at most \(C\beta\), so the same
inequality gives
\(\|\operatorname{Cov}(R)\|_{\rm op}
+\|\operatorname{Cov}(Y)\|_{\rm op}\leq C\beta\).

For any two random vectors, positivity of their
joint covariance gives
\(\|\operatorname{Cov}(Y,R)\|_{\rm HS}^2
\leq\|\operatorname{Cov}(Y)\|_{\rm op}
                  \Tr\operatorname{Cov}(R)\).
This follows, for example, by the covariance
Cauchy--Schwarz contraction factorization.
Also
\(\|\operatorname{Cov}(R)\|_{\rm HS}^2
\leq\|\operatorname{Cov}(R)\|_{\rm op}
                  \Tr\operatorname{Cov}(R)\).
It follows that
\[
 \|\beta^{-1}\operatorname{Cov}(S)
                  -\beta A^T\operatorname{Cov}(z)A\|_{\rm HS}
                                  \leq C\sqrt{N/\beta}.
 \]
Use (92.3) and the bounded operator norm of \(A\)
to replace the second term by \(A^TH^{-1}A\).
For the direct term in (89.2), V87 and bounded
density derivatives give
\[
 \|\mathbb E f_{cc}(z,c)-f_{cc}(z_*,c)
                 -\beta^{-1}\mathbb E\ell_{cc}(z,c)\|_{\rm HS}
                                  \leq C\sqrt{N/\beta}.
 \]
Finally apply the exact score identity and
\(D_c^2G=f_{cc}(z_*,c)-A^TH^{-1}A\).
\end{proof}

\begin{assessment}[What this comparison actually controls]
The full local Hessian, including its interacting
score covariance, now has a vanishing error on
the natural \(\beta\) scale in the stated intensive
Hilbert--Schmidt norm.  The boundary cutoff error
is included.  This is stronger than comparing only
the direct Hessian term or bounding both responses
separately by \(O(\beta)\).

It is still an averaged matrix estimate.  It does
not automatically give an operator-norm or uniform
entrywise error tending to zero over arbitrary
inhomogeneous coarse fields.  Nor does it bound
all higher connected derivatives, justify a change
of local chart in the compact measure, or establish
continuum reflection positivity.
\end{assessment}

\begin{decision}[Next norm and locality test]
Determine whether the intensive comparison can be
upgraded to a uniform local response estimate by
combining the spatial bounds with stronger local
moments, or identify precisely where concentration
on a sparse set of coordinates remains possible.
Do not replace the Hilbert--Schmidt conclusion
by operator-norm convergence without that argument.
\end{decision}

\begin{firewall}[Position after V88]
The interacting coordinate covariance and coarse
Hessian approach their frozen/Schur counterparts
in volume-normalized Hilbert--Schmidt norms.
Uniform local derivative control, compact completion
and scale iteration remain incomplete; the full
Yang--Mills mass gap is not proved.
Only V88 is active; prior bodies and archives
remain.  Next audit and review: V90.
\end{firewall}

% END APPENDED POST-TENTH-ARCHIVE V88 MATERIAL

% =============================================================================
\section{V89: local moments and uniform response-operator convergence}
\label{sec:v89-local}
% =============================================================================

The V88 predecessor has SHA--256
\texttt{4DEE95E0F9DD0592AAFF5008FFDAEB4FF2161797BA4B124605AB2FD5E8AF2D1B}.
Its body is retained except for the version label.
We now supply the local moment estimate missing there,
rather than infer it from an intensive matrix norm.
All results remain on the fixed small-field box.

\begin{theorem}[A second moment at every coordinate]
\label{thm:v89-moment}
Let \(b=z_{\beta,*}(c)\) be the measured minimum.
There is a constant independent of coordinate,
volume and small coarse field such that
\[
 \mathbb E_\pi|z_i-b_i|^2\leq C/\beta,\qquad
 \mathbb E_\pi|z_i-z_{*,i}(c)|^2\leq C/\beta .
                                                        \tag{93.1}
\]
\end{theorem}
\begin{proof}
Fix a coordinate \(i_0\), and put
\(w_i=e^{-\alpha\,{\rm dist}(i,i_0)}\), \(D=\operatorname{diag}w_i\).
The finite-range Hessian bounds (90.1) imply
\[
 \|D^{1/2}(D^2V)D^{-1/2}-D^2V\|_{\rm op}
       \leq K\beta(e^{\alpha R/2}-1).
 \]
Choose \(\alpha>0\) so this is at most \(m\beta/2\).
Although the conjugated matrix need not be symmetric,
its quadratic form on real vectors is bounded below
by \(m\beta/2\).  With \(v=z-b\), integrate the
Hessian on the segment from \(b\) to \(z\) to get
\[
 v^TD\nabla V(z)\geq(m\beta/2)\sum_iw_iv_i^2 .
 \]
The weighted divergence identity gives
\[
 \sum_iw_i-\mathbb E_\pi v^TD\nabla V
 =Z_{\rm loc}^{-1}\int_{\partial\mathcal D}
                  (Dv)\cdot n\,e^{-V}\,dS\geq0.
 \]
Each face contribution is nonnegative since \(b\)
is inside and the weights are positive.
The sum \(\sum_iw_i\) is uniformly bounded:
cell balls have at most polynomial growth of
degree four and a fixed number of coordinates
per cell.  Since \(w_{i_0}=1\), the first
inequality follows.  The uniform max-norm
bound \(b-z_*=O(\beta^{-1})\) from V80 proves
the second.  The argument does not assume
translation-invariant coarse data.
\end{proof}

\begin{theorem}[Uniform entrywise covariance and Hessian comparison]
\label{thm:v89-entries}
With \(H=f_{zz}(z_*,c)\),
\[
 \sup_{i,j}|\beta\operatorname{Cov}_\pi(z_i,z_j)
                                     -(H^{-1})_{ij}|
                      \leq C\beta^{-1/2},       \tag{93.2}
\]
\[
 \sup_{e,f}|\beta^{-1}\partial_e\partial_fW_\beta
                                      -\partial_e\partial_fG|
                      \leq C\beta^{-1/2}.       \tag{93.3}
\]
Also
\(\sup_e|\beta^{-1}\partial_eW_\beta-\partial_eG|
\leq C\beta^{-1/2}\).
\end{theorem}
\begin{proof}
Apply the cutoff trial construction of V88 one
column at a time.  For each fixed \(j\),
\(H^{-1}e_j\) decays exponentially in distance
from \(j\), by V81.  Each nonzero coefficient
of \(M-H\) has second moment at most \(C/\beta\):
use its finite-stencil Lipschitz bound and (93.1).
The same is true of \(1-\chi_i\).  Summing against
the exponentially decaying column gives
\[
 \|(I-MQH^{-1})e_j\|_{L^2(\pi)}^2\leq C/\beta .
 \]
For example each row is a fixed-size sum of
coefficient errors times nearby column entries;
Cauchy--Schwarz bounds it by
\(C\beta^{-1}e^{-\alpha\,{\rm dist}(i,j)}\),
whose sum is uniform.  The trial gradient norm
is at most \(C\), also uniformly in \(j\).
The energy residual argument of V88 gives
\(\|\mathcal B^{-1}e_j-QH^{-1}e_j\|_{L^2}
\leq C/\sqrt\beta\).
The expected cutoff error has the same bound
by (93.1) and exponential summability.
This proves (93.2), indeed with a uniform
Euclidean bound for each column.

For the score decomposition \(S=\mathrm{constant}+Y+R\)
of V88, the gradient of a single remainder \(R_e\)
has finite support and
\(\mathbb E|\nabla R_e|^2\leq C\beta\), by (93.1).
Poincar\'e gives \(\operatorname{Var}(R_e)\leq C\),
while \(\operatorname{Var}(Y_e)\leq C\beta\).
Each covariance cross term is therefore at most
\(C\sqrt\beta\), and the remainder-remainder
term at most \(C\).  Divide by \(\beta\) and
use (93.2) between the finite supports of
the frozen mixed-derivative columns.
The direct Hessian term differs from its value
at \(z_*\) by at most \(C/\sqrt\beta\) in
expectation, again by locality and (93.1).
The exact score and Schur identities prove (93.3).
The first-derivative assertion follows in the
same way from the identity in V87, now using
the finitely many local moments rather than
a total Euclidean moment.
\end{proof}

\begin{corollary}[Operator-norm convergence]
\label{cor:v89-operator}
For \(\beta\geq2\), after increasing a constant,
\[
 \|\beta\operatorname{Cov}_\pi(z)-H^{-1}\|_{\rm op}
 +\|\beta^{-1}D_c^2W_\beta-D_c^2G\|_{\rm op}
       \leq C\beta^{-1/2}(1+\log\beta)^4.        \tag{93.4}
\]
\end{corollary}
\begin{proof}
Each error matrix has uniformly small entries
by the theorem.  It also has an exponential
entry bound independent of \(\beta\):
for covariance use V86 on coordinate functions
and V81 for \(H^{-1}\); for the coarse Hessian
use V86 and V81's minimized Hessian bound.
In a row, sum the small-entry bound inside
radius \(L_\beta=(2\alpha)^{-1}\log\beta\).
There are at most \(C(1+L_\beta)^4\) entries.
Outside that radius, sum the exponential bound;
its tail is at most
\(C(1+L_\beta)^4e^{-\alpha L_\beta}\).
Both contributions have the asserted order.
The matrices are symmetric, so the row-sum
bound also controls their Euclidean operator norm.
\end{proof}

\begin{assessment}[The upgraded conclusion and its limit]
The sparse-coordinate loophole left in V88 is
closed on this local chart.  The natural-scale
coarse Hessian now converges in operator norm
to the minimized Schur Hessian, uniformly over
volume and the declared small coarse fields.
This does not make that Schur Hessian uniformly
positive on retained gauge directions, nor
provide a sign or beta-independent value for
the next-order marginal correction.

The proof also does not estimate higher connected
coarse derivatives or configurations outside the
chart.  In particular it does not identify the
full compact blocked action with \(W_\beta\).
These distinctions remain essential for a
continuum mass-gap claim.
\end{assessment}

\begin{decision}[V90 review and the next scale-sensitive target]
Perform the required companion audit and review.
Determine which next-order response, beyond the
leading Schur Hessian, is required by the actual
scale iteration.  Separate the classical nonlinear
coarse action, the density/determinant correction,
and the controlled but not yet asymptotically
identified interacting remainder.  Avoid declaring
a running coupling from leading-order convergence alone.
\end{decision}

\begin{firewall}[Position after V89]
Uniform local moments upgrade the covariance and
coarse Hessian approximation to operator-norm
convergence on the fixed small-field chart.
Next-order scale response, compact completion
and the full Yang--Mills mass gap remain unproved.
Only V89 is active; prior bodies and archives
remain.  Next audit and review: V90.
\end{firewall}

% END APPENDED POST-TENTH-ARCHIVE V89 MATERIAL

% =============================================================================
\section{V90: the next-order local action and the uniformity audit}
\label{sec:v90-next-order}
% =============================================================================

The V89 predecessor has SHA--256
\texttt{50D900CE1BD775922EEAE2CAB32A80DCD9DC0428C83BC9B9A770953D14B47A76}.
Its body is retained except for the version label.
This is the required audit and ten-version review.
We identify the next asymptotic term for the actual
local density.  Fixed-volume asymptotics below are
not promoted to a volume-uniform expansion.

All fibre derivatives in this section are evaluated
at \((z_*(c),c)\).  Let \(Y\) be centered Gaussian
with covariance \(H_c^{-1}\), and define the polynomials
\[
 a_c(Y)=D_z\ell[Y]-\tfrac16D_z^3f[Y,Y,Y],
 \qquad
 b_c(Y)=\tfrac12D_z^2\ell[Y,Y]
                         -\tfrac1{24}D_z^4f[Y,Y,Y,Y],
 \]
\[
 A_1(c)=\mathbb E\left[b_c(Y)+\tfrac12a_c(Y)^2\right],
 \qquad
 Q(c)=\tfrac12\log\det H_c-\ell(z_*(c),c).       \tag{94.1}
\]
The linear density term is retained in \(a_c\);
it need not vanish at the classical minimum.

\begin{theorem}[Fixed-volume expansion with two coarse derivatives]
\label{thm:v90-laplace}
Fix a finite lattice and a closed smaller coarse
chart on which the minimum stays away from the
box boundary.  Then
\[
 W_\beta(c)=\beta G(c)+\frac N2\log(\beta/2\pi)
              +Q(c)-\beta^{-1}A_1(c)+O(\beta^{-2})
                                                        \tag{94.2}
\]
in \(C^2\) on that parameter chart.  The constant
in the remainder may depend on the lattice volume.
\end{theorem}
\begin{proof}
Translate to \(z_*(c)\) and rescale \(u=\beta^{-1/2}Y\)
in the exact formula (86.3).  The remainder exponent
has expansion
\[
 -\beta V_c(\beta^{-1/2}Y)+L_c(\beta^{-1/2}Y)
       =\beta^{-1/2}a_c(Y)+\beta^{-1}b_c(Y)
                          +O(\beta^{-3/2})
 \]
with the indicated polynomial Taylor coefficients.
The polynomial \(a_c\) is odd, so its Gaussian
expectation is zero.  At order \(\beta^{-1}\),
expansion of the exponential gives precisely
\(b_c+a_c^2/2\).

For control of the remainder, insert an even smooth
cutoff in \(u\) about zero, equal to one on a
smaller ball and supported inside the shifted box
for every parameter in the chosen closed chart.
Uniform positive Hessian bounds make the integral
outside that ball exponentially small in \(\beta\),
also after two coarse differentiations, at this
fixed dimension.  This can be checked in the
original fixed coordinates, where the domain has
no parameter derivative and the differentiated
integrand has only polynomial powers of \(\beta\).

Inside the cutoff, Taylor-expand through the
order needed for a \(\beta^{-2}\) bound, including
two parameter derivatives.  The finite-stencil
functions are smooth to every required order.
On a smaller fixed ball their remainders are
bounded by polynomials in \(Y\) times a Gaussian
with a strictly positive quadratic decay rate.
One may first restrict to a slowly growing
\(Y\)-ball to expand the exponential, and bound
its complement with this Gaussian majorant.
All half-integer terms through order
\(\beta^{-3/2}\) are odd polynomials; their
integrals vanish with the even cutoff.
The difference from full Gaussian polynomial
moments is exponentially small.  This proves
\(\mathcal R_{\beta,c}=1+\beta^{-1}A_1(c)
+O_{C^2}(\beta^{-2})\) at fixed volume.
Taking its logarithm and using (86.3) yields
(94.2).  None of these remainder estimates
has been asserted independent of dimension.
\end{proof}

\begin{proposition}[The first coefficient is extensive]
\label{prop:v90-extensive}
There is a volume-independent constant such that
\[
 |A_1(c)|\leq C N                               \tag{94.3}
\]
throughout the uniform small coarse chart.
\end{proposition}
\begin{proof}
Write \(b_c\) as a sum of bounded-coefficient local
polynomials of degree at most four, with a fixed
number per cell.  Gaussian variances and all
fixed-order local moments are uniformly bounded,
so \(|\mathbb E b_c|\leq CN\).
Likewise \(a_c=\sum_x a_x\), where each \(a_x\)
is an odd local polynomial of degree at most three
with uniformly bounded coefficients.  Each has
zero Gaussian mean.

Wick expansion of \(\mathbb E a_xa_y\) has a
bounded number of terms.  Every nonzero pairing
contains a covariance joining the two local
supports, since each polynomial has odd degree.
V81's exponential inverse bound gives
\(|\mathbb E a_xa_y|\leq Ce^{-\alpha\,{\rm dist}(x,y)}\),
after a fixed stencil enlargement.  Summing over
\(y\) and then \(x\) gives
\(\mathbb E a_c^2\leq CN\).
This proves (94.3), without treating the local
polynomials as independent.
\end{proof}

\begin{corollary}[The actual fixed-volume order-one Hessian target]
At each fixed volume,
\[
 D_c^2W_\beta-\beta D_c^2G
                       \longrightarrow D_c^2Q
                                                        \tag{94.4}
\]
on the smaller coarse chart.  All derivatives
of \(Q\) include the dependence through \(z_*(c)\).
\end{corollary}
\begin{proof}
Differentiate the proved \(C^2\) expansion.
The normalization \(N\log(\beta/2\pi)/2\) is
independent of \(c\).  The chain rule in (94.1)
includes the moving minimum and cannot be
replaced by derivatives at the section.
\end{proof}

\begin{assessment}[What remains before a scale-response claim]
V89 established a volume-uniform leading-order
response limit.  Equation (94.4) identifies the
next term only at fixed volume.  The extensive
coefficient bound (94.3) does not imply a
volume-uniform \(C^2\) remainder in (94.2).
Interchanging the large-\(\beta\), volume and
regulator limits remains unjustified.

The order-one response is also not merely a
Wilson quartic trace: it includes the measured
Jacobian through \(-\ell_*\), the full constrained
Hessian determinant, and the moving background.
Its sign in any calibrated marginal direction
has not been evaluated here.
\end{assessment}

\begin{audit}[Required V90 companion check]
SM remains V72, SHA--256
\texttt{F44F9745D43379E1672C5550411B58E97195A4C9278592D221DF5D3F644DC1F5};
NS remains V26, SHA--256
\texttt{82C887F8C2C69E4F28FAD7C3DC8DE5B9099CB5A4BF431EBA1FFF3E91935978AD}.
These hashes match the previously read companion
conclusions.  The parallel YM active note remains V5.
No new companion estimate supplies the uniform
next-order or compact-fibre remainder.
Next companion audit: V95; next ten-version review
and mandatory rollover: V100.
\end{audit}

\begin{record}[Ten-version review]
V81 localized the moving minimum and its Schur
response.  V82 separated the exact fluctuation
factor from the quadratic determinant.
V83--V84 proved two-sided intensive local free-energy
control.  V85--V86 bounded the interacting spatial
response with the box boundary retained.
V87--V89 progressed from intensive moments to
uniform local moments and operator-norm leading
response convergence.  V90 identifies the next
fixed-volume coefficient without claiming the
uniformity still missing.
The local chart now has substantially more control;
it has not been identified with the entire compact
blocked measure or iterated to a continuum theory.
\end{record}

\begin{decision}[Next concrete coefficient calculation]
Evaluate the combined order-one quadratic response
\(D_c^2Q(0)\) for a specified finite regulator,
including the moving-minimum derivatives and
the reduced density.  Compare it with any earlier
frozen-background card only after verifying that
its block map, fibre coordinates and measure
are the same.  In parallel with that logical
comparison, retain the uniform-remainder problem
as an explicit prerequisite for iteration.
\end{decision}

\begin{firewall}[Position after V90]
The next-order local action and its first correction
are identified with fixed-volume \(C^2\) control,
and the correction coefficient is bounded extensively.
Uniform next-order response, compact completion
and the full Yang--Mills mass gap remain unproved.
Only V90 is active; earlier bodies and archives remain.
\end{firewall}

% END APPENDED POST-TENTH-ARCHIVE V90 MATERIAL

% =============================================================================
\section{V91: coarse-Haar normalization of the order-one response}
\label{sec:v91-haar}
% =============================================================================

The V90 predecessor has SHA--256
\texttt{C2218CC35569AD0267CB1FA689FCE0747A5B80FAEB955FB7CCAE2458EF2F61DA}.
Its body is retained except for the version label.
Before evaluating a quadratic response, its reference
measure must be fixed.  The local density used so
far is relative to Lebesgue coarse coordinates.
Let
\[
 j_{\rm c}(c)=\prod_e j_G(c_e),\qquad
 \widehat Q(c)=Q(c)+\log j_{\rm c}(c).           \tag{95.1}
\]
Field-independent normalization constants have no
effect on the derivatives below.  Work with a compact
connected semisimple gauge group and the root-reference
block on the common small chart.

\begin{theorem}[Local transformation law of the determinant-density coefficient]
\label{thm:v91-transform}
Let a fixed sufficiently small cell-constant gauge
transformation induce \(c'=\psi(c)\), and assume
the relevant minima and their neighbourhoods stay
inside the chart.  Then
\[
 Q(\psi(c))=Q(c)+\log|\det D\psi(c)|,\qquad
                    \widehat Q(\psi(c))=\widehat Q(c).
                                                        \tag{95.2}
\]
\end{theorem}
\begin{proof}
The gauge transformation preserves the tree slice,
the Wilson action and its Haar measure.  Block
equivariance makes its coordinate map triangular:
\((z,c)\mapsto(z',c')=(\phi(z,c),\psi(c))\).
At the moving minimum put \(J=D_z\phi\).
Uniqueness of that local minimum carries it to
the minimum over \(\psi(c)\), as in V81.

The joint local fine measure is
\(e^{-\beta f(z,c)+\ell(z,c)}dz\,dc\).
Its invariance gives, at the corresponding minima,
\[
 \ell_*'=\ell_*-\log|\det J|-\log|\det D\psi|.
 \]
Constrained stationarity eliminates the coordinate
curvature term in transformation of the fibre
Hessian, giving
\(H'=J^{-T}HJ^{-1}\).  Consequently
\(\frac12\log\det H'=\frac12\log\det H-\log|\det J|\).
Subtracting the density identity proves the
first assertion of (95.2).

Coarse product Haar is invariant under its
endpoint gauge action, so
\(j_{\rm c}(\psi(c))|\det D\psi(c)|=j_{\rm c}(c)\).
Adding its logarithm proves the second assertion.
This is a local transformation of the leading
Laplace coefficient and its regular measure.
It does not assert invariance of the full
fixed-box partition function, whose boundary
generally moves under the transformation.
\end{proof}

\begin{theorem}[Exact gauge-sector Hessian benchmark]
\label{thm:v91-benchmark}
Let \(d_{\rm c}\omega\) be a coarse gauge gradient.
At the identity, for every coarse direction \(v\),
\[
 D^2\widehat Q(0)[d_{\rm c}\omega,v]=0,\qquad
 D^2Q(0)[d_{\rm c}\omega,v]
 =-\frac1{12}\sum_e
       \Tr_{\mathfrak g}
          (\operatorname{ad}_{(d_{\rm c}\omega)_e}
                         \operatorname{ad}_{v_e}). \tag{95.3}
\]
\end{theorem}
\begin{proof}
Global constant gauge transformations act by
simultaneous adjoint conjugation and fix \(c=0\).
Invariance of \(\widehat Q\) implies that its
first derivative at zero is an invariant covector
in each coarse Lie-algebra component.  Such
covectors vanish for a semisimple Lie algebra.
Thus \(D\widehat Q(0)=0\).

Differentiate the infinitesimal local gauge
identity \(D\widehat Q(c)[R_c\omega]=0\) with
respect to \(c\) in direction \(v\).
At zero the term containing \(D\widehat Q(0)\)
vanishes and \(R_0\omega=d_{\rm c}\omega\),
proving the first equality.  No positivity of
the Hessian has been assumed.

The one-link Haar expansion already used in V59 is
\[
 \log j_G(x)=\log j_G(0)
               +\tfrac1{24}\Tr_{\mathfrak g}\operatorname{ad}_x^2
               +O(\|x\|^4).
 \]
Polarizing its quadratic term and using (95.1)
gives the second equality.
\end{proof}

\begin{corollary}[A concrete \(SU(2)\) regression value]
For \(T=\operatorname{diag}(i,-i)\), let \(\omega\)
equal \(T\) at one coarse vertex and zero elsewhere,
on a torus with distinct incoming and outgoing edges.
Then
\[
 D^2Q(0)[d_{\rm c}\omega,d_{\rm c}\omega]=16/3,
 \qquad
 D^2\widehat Q(0)[d_{\rm c}\omega,d_{\rm c}\omega]=0.
                                                        \tag{95.4}
\]
\end{corollary}
\begin{proof}
There are eight nonzero incident coarse gradients,
each \(T\) or \(-T\).
On \(\mathfrak{su}(2)\), \(\operatorname{ad}_T\)
has eigenvalues \(0,2i,-2i\), so
\(\Tr_{\mathfrak g}\operatorname{ad}_T^2=-8\).
Equation (95.3) gives \(-8(-8)/12=16/3\);
the Haar correction cancels it.
\end{proof}

\begin{assessment}[How to use, and not use, the benchmark]
The nonzero value in (95.4) is a coordinate-density
effect in the Lebesgue-normalized coefficient.
It is not a physical gauge-boson mass or a violation
of gauge invariance.  Conversely, a computation
of the combined coefficient that fails (95.3)
has a normalization, moving-background, or
Jacobian inconsistency to resolve.

This theorem fixes the gauge-sector benchmark.
It does not evaluate the transverse order-one
response or establish its sign in a marginal
running-coupling direction.  It also does not
make the fixed-volume next-order expansion
uniform in volume.
\end{assessment}

\begin{decision}[Next transverse coefficient target]
Compute a declared transverse component of
\(D^2\widehat Q(0)\), with the coarse Haar
term, moving-minimum derivatives, constrained
determinant and reduced fine density all included.
Use (95.3)--(95.4) as exact normalization checks.
Do not compare with a frozen-background card
until that card's measure and fibre convention
have been matched.
\end{decision}

\begin{firewall}[Position after V91]
The order-one coefficient now has an explicit
coarse-Haar normalization and an exact gauge-sector
quadratic benchmark.  The transverse next-order
response, uniform remainder, compact completion
and full Yang--Mills mass gap remain unproved.
Only V91 is active; preceding bodies and archives
remain.  Next companion audit: V95.
\end{firewall}

% END APPENDED POST-TENTH-ARCHIVE V91 MATERIAL

% =============================================================================
\section{V92: the identity fourth-jet reduction and order-one spatial locality}
\label{sec:v92-fourth-jet}
% =============================================================================

The V91 predecessor has SHA--256
\texttt{B9FE6287A8FD885AC2EEDDC937C5D0EADD2F9F09F4AD12EAF6ECB01D810555B4}.
Its body is retained except for the version label.
The pending transverse calculation involves the moving
minimum, not the section.  Before forming its full jets,
we remove a term that vanishes by symmetry and establish
an order-one, rather than order-\(\beta\), locality bound
for the identified coefficient itself.  This does not
differentiate the error estimate in V89 or make V90's
asymptotic remainder uniform.

Use precisely the root-reference block and independent
tree-fibre coordinates of V80.  Write \(x=(z,c)\),
\[
 H=f_{zz}(0,0),\qquad P=H^{-1},\qquad A=f_{zc}(0,0),
 \qquad J a=(-PAa,a).
                                                        \tag{96.1}
\]
Here \(J\) is a map from coarse directions to the joint
coordinate space, not the fibre-change Jacobian in V91.
The Lie algebra is compact semisimple with the invariant
inner product and orthonormal coordinates used above.
Every derivative in the following displayed jet formula
is evaluated at \(x=0\).

\begin{lemma}[Stationarity of the off-minimum coefficient at the identity]
\label{lem:v92-stationary}
On a small joint chart define
\[
 q(z,c)=\tfrac12\log\det f_{zz}(z,c)-\ell(z,c).
                                                        \tag{96.2}
\]
Then \(Dq(0,0)=0\).  In particular, if
\(\iota(c)=(z_*(c),c)\), then
\[
 D^2(q\circ\iota)(0)[a,b]=D^2q(0)[Ja,Jb].
                                                        \tag{96.3}
\]
\end{lemma}
\begin{proof}
The fibre Hessian is positive throughout a sufficiently
small joint chart by V80, so (96.2) is defined there.
A global constant gauge transformation acts by the same
adjoint map on each free fibre and coarse coordinate.
The linear maps \(E,R\) commute with this action;
uniqueness of the local implicit block solution implies
the same equivariance for \(X(z,c)\).
Consequently \(f\) is invariant.  Under this linear
orthogonal coordinate action its \(zz\) Hessian is
conjugated by the orthogonal action on fibre coordinates,
so its determinant is invariant.

The reduced Haar density is invariant as well: the joint
coordinate transformation has absolute determinant one,
and the underlying fine Haar measure is invariant.
Thus \(q\) is invariant as a germ at zero.  This argument
does not require the entire fixed coordinate box to be
preserved by a finite adjoint transformation.
An invariant linear functional on each copy of a
semisimple Lie algebra vanishes: differentiating its
adjoint invariance annihilates every bracket, and
\([\mathfrak g,\mathfrak g]=\mathfrak g\).
Applying this to the derivative of \(q\) gives \(Dq(0)=0\).

Differentiating \(f_z(z_*(c),c)=0\) once gives
\(D\iota(0)=J\).  The second chain rule contains the
additional term \(Dq(0)[D^2\iota(0)[a,b]]\), which is
zero.  This proves (96.3).  It does not say that
\(D^2z_*(0)\) is zero or that the minimum can be frozen.
\end{proof}

For joint vectors \(\xi,\eta\), define fibre matrices
\[
 B(\xi)_{ij}=D^3f(0)[e_i^z,e_j^z,\xi],\qquad
 C(\xi,\eta)_{ij}=D^4f(0)[e_i^z,e_j^z,\xi,\eta].
                                                        \tag{96.4}
\]
The vectors \(e_i^z\) are the scalar fibre-coordinate basis.
Traces below are finite-dimensional fibre traces; the
last term explicitly uses the adjoint Lie-algebra trace.

\begin{theorem}[Complete fourth-jet formula at the identity]
\label{thm:v92-jet}
For every pair of coarse directions,
\[
\begin{split}
 D^2\widehat Q(0)[a,b]
 ={}&\tfrac12\Tr\!\left(P C(Ja,Jb)
                         -P B(Ja)P B(Jb)\right)\\
 &-D^2\ell(0)[Ja,Jb]
   +\tfrac1{12}\sum_e\Tr_{\mathfrak g}
                (\operatorname{ad}_{a_e}\operatorname{ad}_{b_e}).
\end{split}                                             \tag{96.5}
\]
No second derivative of the minimizing map is required
in this formula.  The linear response \(-PA\) is required.
\end{theorem}
\begin{proof}
The determinant identity
\(D\log\det M[\dot M]=\Tr(M^{-1}\dot M)\), followed
by \(D(M^{-1})[\dot M]=-M^{-1}\dot M M^{-1}\),
gives
\[
 D^2q(0)[\xi,\eta]
 =\tfrac12\Tr\!\left(PC(\xi,\eta)-PB(\xi)PB(\eta)\right)
                              -D^2\ell(0)[\xi,\eta].
\]
Apply the lemma, use \(Q=q\circ\iota\), and add the
coarse-Haar Hessian from V91.  Cyclicity of the trace
makes the expression symmetric in \(a,b\).
This also shows exactly where the terms involving
\(D^2z_*\) in a direct differentiation cancel:
their combined coefficient is \(D_zq(0)\).
\end{proof}

\begin{proposition}[Finite input jet order]
\label{prop:v92-input-order}
Computing (96.5) requires the action through fourth
order, the logarithmic block through third order,
and the logarithmic reduced density through second
order.  Fourth derivatives of the implicit fibre map
and nonlinear minimizer solves are unnecessary.
\end{proposition}
\begin{proof}
Let \(F\) denote the logarithmic root block on the
tree slice.  In these coordinates
\(X(z,c)=Ez+Rv(z,c)\) and \(F(X(z,c))=c\).
Put \(U=D X(0)\); then \(U(z,c)=Ez+Rc\), since
\(DF(0)=L\), \(LE=0\), and \(LR=I\).
For arbitrary joint directions \(u,v,w\), implicit
differentiation gives
\[
 X_2[u,v]=-R\,F_2[Uu,Uv],
\]
\[
\begin{split}
 X_3[u,v,w]=-R\big(&F_3[Uu,Uv,Uw]
                   +F_2[X_2[u,v],Uw]\\
                  &+F_2[X_2[u,w],Uv]
                   +F_2[X_2[v,w],Uu]\big).
\end{split}                                             \tag{96.6}
\]
There is no coarse second or third derivative on the
right side of \(F(X)=c\).  All \(X_k\), \(k\geq2\),
take values in the range of \(R\), which justifies
solving these equations using \(LR=I\).

Write the log-coordinate Wilson action as \(s(X)\).
It has \(Ds(0)=0\).  The only term involving \(X_4\)
in the fourth chain rule for \(f=s\circ X\) is
\(Ds(0)[X_4]\), hence it vanishes.  The remaining
terms use only \(s_2,s_3,s_4,U,X_2,X_3\).
The second derivative of the density uses at most
its fine-coordinate second derivative and \(X_2\).
Its determinant representation contains \(DFR\);
two differentiations of that matrix need at most
\(F_3\).  Constant linear coordinate-volume factors
have no derivatives.  This proves the stated input
order, with the same root-reference convention in
every jet.
\end{proof}

\begin{theorem}[Volume-uniform locality of the identity order-one coefficient]
\label{thm:v92-locality}
Let \(K=D^2\widehat Q(0)\), with scalar coarse indices
including their fixed Lie components.  There are
constants \(C_0,\alpha_0>0\), independent of periodic
volume, such that
\[
 |K_{ef}|\leq C_0 e^{-\alpha_0\,{\rm dist}(e,f)},
 \qquad \|K\|_{\rm op}\leq C_0'
                                                        \tag{96.7}
\]
for another volume-independent constant \(C_0'\).
These are bounds on the identified coefficient,
not on the next-order error of the exact integral.
\end{theorem}
\begin{proof}
V63 and V80 give finite-stencil representations of
\(f,\ell\), with bounded derivatives through the orders
needed here and a bounded number of terms per cell.
V81 gives exponential bounds on \(P\) and on the
columns of \(J\); the coarse part of a column of
\(J\) is a coordinate vector and satisfies the same
bound after increasing its constant.

Choose one common finite interaction radius \(r_0\)
and reduce a positive exponent \(\alpha\) as needed.
Writing \(J_e=Je_e\), the finite-stencil derivative
tensors imply
\[
 |B(J_e)_{ij}|
 \leq C\,{\bf1}_{\{{\rm dist}(i,j)\leq r_0\}}
                   e^{-\alpha\,{\rm dist}(i,e)},
\]
\[
 |C(J_e,J_f)_{ij}|
 \leq C\,{\bf1}_{\{{\rm dist}(i,j)\leq r_0\}}
 e^{-\alpha({\rm dist}(i,e)+{\rm dist}(i,f))}.
                                                        \tag{96.8}
\]
Indeed, a nonzero tensor entry has every index in one
bounded stencil; contracting each extra index against
a column of \(J\) gives the indicated bound.  The
number of tensor entries at a fixed anchor is bounded.
The same argument bounds the contracted density Hessian
by
\[
 |D^2\ell(0)[J_e,J_f]|
 \leq C\sum_i
       e^{-\alpha({\rm dist}(i,e)+{\rm dist}(i,f))}.
\]
The first trace in (96.5) obeys this last bound as
well, using (96.8), the bounded entries of \(P\),
and the bounded number of allowed \(j\) for each \(i\).

For completeness the trace with two inverses is a
sum of \(P_{ij}B(J_e)_{jk}P_{kl}B(J_f)_{li}\).
Here \(k\) is within \(r_0\) of \(j\), and \(i\)
within \(r_0\) of \(l\).  Its absolute sum is bounded by
\[
 C\sum_{j,l}
 e^{-\alpha({\rm dist}(j,e)+{\rm dist}(l,f)
                                      +2{\rm dist}(j,l))}.
                                                        \tag{96.9}
\]
The triangle inequality bounds \({\rm dist}(e,f)\)
by \({\rm dist}(e,j)+{\rm dist}(j,l)+{\rm dist}(l,f)\).
Use half of the exponent to extract
\(e^{-\alpha{\rm dist}(e,f)/2}\).
The remaining sum is bounded by the product of two
uniform exponential row sums centered at \(e\) and
\(f\).  Such sums are uniform because balls in the
periodic cell graph contain at most \(C(1+r)^4\)
scalar coordinates.  The single-anchor sums above
are bounded by the same argument.
Finally, the coarse-Haar term is diagonal in its
edge index and has bounded finite-dimensional entries.
This proves the entry bound in (96.7).

Summing that bound over \(f\) gives a uniform absolute
row sum.  Symmetry gives the same column sum, so the
Schur bound proves the Euclidean operator estimate.
No sign of \(K\), and no spectral gap for \(K\), was
used in this proof.
\end{proof}

\begin{assessment}[What this changes for the transverse calculation]
The order-one identity Hessian is now a bounded
exponentially local operator uniformly in volume,
and its finite-regulator value is reduced to the
explicit contractions in (96.5).
The second-minimizer-derivative term is eliminated
by a proved symmetry, not dropped as an approximation.
In particular, one must still use \(Ja=(-PAa,a)\);
substituting the section tangent \((0,a)\) is not justified.

The V91 gauge test remains mandatory: for
\(a=d_{\rm c}\omega\), the full right side of (96.5)
must vanish against every \(b\).
This test does not determine the value on transverse
directions.  Neither (96.7) nor that Ward identity
fixes the sign of a running-coupling correction.
Also, a uniform bound for the limiting coefficient
does not imply uniform convergence to it in V90.
\end{assessment}

\begin{decision}[Next coefficient work]
Assemble the root-reference jets in (96.6) and (96.5)
for a declared non-pure-gauge coarse direction, with
the V91 gauge regression in the same calculation.
The remaining input is explicit finite algebra,
not an uncomputed nonlinear minimum.
Keep the volume-uniform asymptotic remainder and the
comparison with remote compact-fibre configurations
as separate unresolved estimates.
\end{decision}

\begin{firewall}[Position after V92]
The order-one identity response has a complete
fourth-jet formula and a volume-uniform exponential
locality bound.  Its transverse value and sign,
the uniform next-order remainder, compact completion
and the full Yang--Mills mass gap remain unproved.
Only V92 is active; preceding bodies and archives
remain.  Next companion audit: V95.
\end{firewall}

% END APPENDED POST-TENTH-ARCHIVE V92 MATERIAL

% =============================================================================
\section{V93: the coefficient limit, the zero-momentum test, and a transverse target}
\label{sec:v93-coefficient-limit}
% =============================================================================

The V92 predecessor has SHA--256
\texttt{FF8D86874D978F8CA74A0289B8E483803CE2271B4834D776F5187121AA9DED27}.
Its body is retained except for the version label.
We now fix cubic periodic coarse lattices of side \(L\),
in cell units, with the same translated rooted tree,
boundary-pivot choice and root-reference block at every
cell.  Lie coordinates and the bare action normalization
are fixed as \(L\) varies.  All statements concern the
identity coefficient \(K_L=D^2\widehat Q_L(0)\).
They do not exchange the volume limit with the
large-\(\beta\) limit of the exact blocked integral.

\begin{lemma}[Finite-range inverse approximants]
\label{lem:v93-polynomial}
There is a positive bounded finite-range operator
\(H_\infty\) on the infinite cell lattice with the
same local coefficients as \(H_L=f_{zz,L}(0)\), and
\(hI\leq H_\infty\leq bI\) for fixed \(h,b>0\).
Choose \(b'>b\), set \(T=I-H/b'\), and put
\[
 P^{(m)}=(b')^{-1}\sum_{j=0}^m T^j .
                                                        \tag{97.1}
\]
On both the infinite lattice and every periodic lattice,
\[
 \|P-P^{(m)}\|_{\rm row}
 +\|P-P^{(m)}\|_{\rm col}\leq C e^{-a m},
 \qquad P=H^{-1}.                              \tag{97.2}
\]
Here \(\|M\|_{\rm row}=\sup_i\sum_j|M_{ij}|\), and
the column norm is defined analogously.
All constants are independent of \(L,m\).
\end{lemma}
\begin{proof}
Define \(H_\infty\) by the common finite stencil.
A finitely supported vector embeds in a sufficiently
large torus without crossing its boundary; its
quadratic form and norm there are unchanged.
V80's lower and upper bounds therefore hold on the
dense subspace of finitely supported vectors, hence
on its completion.  In particular \(H_\infty\) is
invertible.  The same argument identifies its stencil
with that of every sufficiently large torus.

The spectral norm of \(T\) is at most
\(q=1-h/b'<1\).  A row of \(T^j\) is supported on
at most \(C(j+1)^4\) scalar indices.  Its Euclidean
norm is at most \(q^j\), so its absolute sum is at
most \(C(j+1)^2q^j\).  Sum over \(j>m\), and absorb
the polynomial in a slightly slower exponential.
Symmetry proves the column estimate as well.
This argument applies verbatim to the infinite lattice.
\end{proof}

For a translation-invariant infinite kernel \(k\), write
\[
 (\Pi_L k)(r)=\sum_{n\in\mathbb Z^4}k(r+Ln)
                                                        \tag{97.3}
\]
for its periodization.  Kernel entries retain their
finite edge-orientation and Lie-component indices.

\begin{theorem}[Thermodynamic limit of the identified coefficient]
\label{thm:v93-limit}
There is a translation-invariant real symmetric
operator \(K_\infty\), with kernel \(k_\infty(r)\),
such that
\[
 |k_\infty(r)|\leq C e^{-a|r|_1},\qquad
 \|K_L-\Pi_L K_\infty\|_{\rm row}\leq C e^{-a' L}.
                                                        \tag{97.4}
\]
Matrix norms on the fixed internal components can
be used in the first bound.  The constants may depend
on the fixed blocking but not on the volume.
\end{theorem}
\begin{proof}
Use V92's formula (96.5), with all local tensors on
the infinite lattice, \(P=H_\infty^{-1}\), and
\(J=(-PA,I)\).  The sums defining each entry converge
absolutely by precisely the anchored exponential sums
in (96.8)--(96.9).  That proof gives the first bound
in (97.4); translation symmetry and real symmetry
are inherited from the tensors and the contractions.
No infinite-dimensional determinant is being defined.

Let \(K^{(m)}\) be the same finite tensor expression
with every occurrence of \(P\) replaced by \(P^{(m)}\),
including those in \(J\).  Then, on either lattice,
\[
 \|K-K^{(m)}\|_{\rm row}\leq C e^{-a_1m}.
                                                        \tag{97.5}
\]
To justify this estimate for the trace terms, expand
their difference by replacing one propagator at a
time.  The local tensors have bounded entries,
bounded stencil size, and a bounded number of
terms per anchor.  Each resulting contraction is a
connected finite graph: for the quartic trace it has
one local vertex with a propagator loop and two
external lift factors; for the pair of cubic tensors
it has two local vertices joined by two propagators
and the external lift factors.  The density term
has one local vertex and two external lifts.

Fix the first external index and sum absolute values
over the second external index and all internal ones.
Along a spanning tree each summation costs a bounded
row or column norm.  Any remaining loop edge costs
at most a bounded supremum entry.  A loop at a local
vertex also has a bounded supremum entry.  For the
single replaced factor these bounds cost at most
\(Ce^{-a_1m}\) by (97.2); all other factors have
uniformly bounded row and column norms.  If the
replaced factor is not in the chosen spanning tree,
use its supremum bound instead.  The graph has a
fixed number of edges, independent of \(m,L\).
This proves (97.5), including the absolute row sum,
without using an extensive bound on a trace.

Each \(P^{(m)}\) has propagation radius at most
\(R(m+1)\).  Expand it into the finite path sums
in (97.1).  Every graph in \(K^{(m)}\) then has
total path length at most \(R_1(m+1)\), for a
fixed \(R_1\).  If \(L>2R_1(m+1)\), it cannot
have a cycle with nonzero winding around the torus.
Its paths lift consistently to the infinite lattice;
the possible endpoints give exactly periodization.
Thus
\[
 K_L^{(m)}=\Pi_L K_\infty^{(m)}
       \quad\hbox{when }L>2R_1(m+1).            \tag{97.6}
\]
Periodization cannot increase an absolute row sum.
Apply (97.5) on each lattice and choose
\(m=\lfloor L/(4R_1)\rfloor-1\) for sufficiently
large \(L\).  This proves the second assertion of
(97.4).  Increase the constant for the finitely
many smaller admissible sides.
\end{proof}

Use the convention
\[
 (K_\infty a)(x)=\sum_r k_\infty(r)a(x+r),
 \qquad
 \mathcal K(p)=\sum_r k_\infty(r)e^{ip\cdot r}.
                                                        \tag{97.7}
\]
The symbol is Hermitian for real \(p\).  The linear
coarse gauge gradient has symbol
\(g_\mu(p)=1-e^{ip_\mu}\), tensored with the identity
on the Lie algebra.

\begin{theorem}[Analytic Ward identity and absence of a constant quadratic term]
\label{thm:v93-zero}
The symbol in (97.7) is analytic in a complex strip
around the real momentum torus and satisfies
\[
 \mathcal K(p)g(p)=0,\qquad \mathcal K(0)=0.
                                                        \tag{97.8}
\]
Moreover the finite-torus constant-mode matrix obeys
\[
 \|\mathcal K_L(0)\|\leq C e^{-a'L}.             \tag{97.9}
\]
\end{theorem}
\begin{proof}
Exponential kernel summability permits the Fourier
series and all its derivatives in a sufficiently
small complex strip.
V91 gives \(K_Ld_{\rm c}=0\).  Apply it to any
finitely supported gauge parameter embedded in
successively larger tori.  Equation (97.4) and
exponential decay give \(K_\infty d_{\rm c}=0\)
entrywise.  Fourier transformation proves the first
identity in (97.8).

Differentiate that identity in \(p_\nu\) at zero.
The term with \(D\mathcal K\) vanishes because
\(g(0)=0\), while
\(\partial_\nu g_\mu(0)=-i\delta_{\mu\nu}I_{\mathfrak g}\).
Thus \(\mathcal K(0)\) annihilates each orientation
and every Lie component, proving \(\mathcal K(0)=0\).
Finally the constant-mode symbol of a periodized
absolutely summable kernel equals the infinite
kernel's symbol at zero.  Apply the row estimate
in (97.4) to the constant mode to obtain (97.9).
Finite-torus Ward identities alone would not suffice:
the constant torus mode is not the gradient of a
periodic gauge parameter.
\end{proof}

\begin{proposition}[A specified transverse small-momentum coefficient]
\label{prop:v93-transverse}
Choose distinct orientations \(1,2\), a real unit
Lie vector \(T\), and the polarization \(v=e_2\otimes T\).
For \(p=t e_1\), this polarization is exactly transverse
to \(g(p)\), and
\[
 v^*\mathcal K(t e_1)v
      =t^2\theta_{12,T}+O(t^4),\qquad
 \theta_{12,T}=-\tfrac12\sum_{r\in\mathbb Z^4}
                          r_1^2 v^T k_\infty(r)v .
                                                        \tag{97.10}
\]
The sum converges absolutely.  Its sign is not fixed
by the Ward identity or the locality bound.
\end{proposition}
\begin{proof}
The only nonzero component of \(g(t e_1)\) is its
orientation \(1\) component, proving transversality.
Real symmetry of the kernel says
\(k_\infty(-r)=k_\infty(r)^T\).
Therefore the scalar Fourier series for this real
polarization is even in \(t\).  Its constant term
is zero by (97.8).  Taylor expansion gives (97.10),
with the remainder bounded by a constant times
\(t^4\sum_r|r_1|^4|k_\infty(r)|\).
Exponential decay makes this moment finite.
No reflection symmetry of the rooted blocking,
nor a scalar tensor decomposition of the full
symbol, is assumed in this argument.
\end{proof}

\begin{assessment}[Interpretation and remaining calculation]
The coefficient itself now has a controlled volume
limit, and constant-mode finite-size effects are
exponentially small.  This is not a volume limit
for the full compact quantum theory.
Equation (97.10) specifies a non-pure-gauge momentum
response whose value can be computed from V92's
root-reference jets.  Its numerical value and sign
have not been computed here, and it has not been
identified with a universal beta-function coefficient.

The zero in (97.8) concerns an analytic quadratic
coefficient in a weak-field local expansion.
It neither disproves a nonperturbative physical mass
gap nor supplies one.  Interacting large-field control,
uniform asymptotic remainders, continuum construction
and reconstruction remain separate requirements.
\end{assessment}

\begin{firewall}[Position after V93]
A thermodynamic limit and a zero-momentum Ward test
are proved for the identified order-one coefficient,
with exponentially small finite-volume errors.
A transverse second-moment target is fixed but not
evaluated.  The full Yang--Mills mass gap remains
unproved.  Only V93 is active; earlier bodies and
archives remain.  Next companion audit: V95.
\end{firewall}

% END APPENDED POST-TENTH-ARCHIVE V93 MATERIAL

% =============================================================================
\section{V94: explicit root-chart inversion and an exact Wilson fourth-order check}
\label{sec:v94-root-jets}
% =============================================================================

The V93 predecessor has SHA--256
\texttt{CBC760D0A005F5E0874AA4ACE59493839D9B36D19E09E1E0EB7CB4BBFA134425}.
Its body is retained except for the version label.
V60 already proved the cubic reference correction,
and V92 gave the abstract implicit-jet reduction.
This section does not claim those results again:
it supplies their explicit chronological two-link
implementation on the current tree slice, a fourth-order
Wilson action formula, and a full finite-lattice
exact-arithmetic regression of those inputs.

All coefficients below are Taylor coefficients,
not derivatives multiplied by factorials.
The tree links are identity.  For a coarse edge \(e\),
write the sixteen chronological two-link paths as
\((l_j,k_j)\), with \(j=0\) the root path, and let
the bar denote their average.  For a fine first-order
field \(u\), set
\[
 a_j=u_{l_j},\quad b_j=u_{k_j},\quad
 p_j=a_j+b_j,\quad \delta p_j=p_j-\overline p,
\]
\[
 \mathcal Q_e(u)=\tfrac12\overline{[a_j,b_j]},
 \qquad
 \mathcal C_e(u)=\tfrac1{12}\overline{
 [a_j,[a_j,b_j]]+[b_j,[b_j,a_j]]
                       +[\delta p_j,[p_0,\delta p_j]]}.
                                                        \tag{98.1}
\]
Thus the ordering is chronological even for paths
whose boundary link comes before their internal link.
For two fine fields define the polarized cross coefficient
\[
 \mathcal B_e(u,w)=\tfrac12\overline{
       [u_{l_j},w_{k_j}]+[w_{l_j},u_{k_j}]}.
                                                        \tag{98.2}
\]

\begin{proposition}[Explicit cubic inverse in a joint direction]
\label{prop:v94-inverse}
Fix a joint direction \((z,c)\) and put
\[
 u=Ez+Rc,\qquad
 v_2=-\mathcal Q(u),\quad x_2=Rv_2,\qquad
 v_3=-\mathcal C(u)-\mathcal B(u,x_2),\quad x_3=Rv_3.
                                                        \tag{98.3}
\]
Then the actual implicit fibre map has expansion
\[
 X(tz,tc)=t u+t^2x_2+t^3x_3+O(t^4).
                                                        \tag{98.4}
\]
\end{proposition}
\begin{proof}
For a first-order input, the two-link logarithm is
\[
 \log(e^{ta_j}e^{tb_j})
 =tp_j+\tfrac{t^2}{2}[a_j,b_j]
   +\tfrac{t^3}{12}
       ([a_j,[a_j,b_j]]+[b_j,[b_j,a_j]])+O(t^4).
 \]
Adding V60's reference correction gives
\(F(tu)=tLu+t^2\mathcal Q(u)+t^3\mathcal C(u)+O(t^4)\)
for the root-reference tree-slice block.
For an input \(tu+t^2w+t^3y\), the same polynomial
expansion gives
\[
 F(tu+t^2w+t^3y)
 =tLu+t^2(Lw+\mathcal Q(u))
    +t^3(Ly+\mathcal B(u,w)+\mathcal C(u))+O(t^4).
 \]
Use \(LE=0\), \(LR=I\) and (98.3).  The output is
\(tc+O(t^4)\), and the free fibre component is exactly
\(tEz\), since the corrections lie in the range of \(R\).
Uniqueness of the smooth implicit solution identifies
its first three coefficients with those in (98.4).
This is also the diagonal specialization of (96.6),
with the Taylor-versus-derivative factors accounted for.
\end{proof}

\begin{theorem}[Wilson fourth-order coefficient from cubic link logs]
\label{thm:v94-action}
Let the logarithm of the oriented plaquette holonomy
for the links in (98.4) be
\[
 \log U_p(t)=tP_p+t^2Q_p+t^3R_p+O(t^4).
 \]
Use the real invariant inner product
\(\langle A,B\rangle=-\operatorname{Re}\Tr(AB)\).
The composed Wilson action has Taylor coefficients
\[
\begin{split}
 f(tz,tc)={}&t^2 s_2+t^3s_3+t^4s_4+O(t^5),\\
 s_2={}&\tfrac12\sum_p\|P_p\|^2,\qquad
 s_3=\sum_p\langle P_p,Q_p\rangle,\\
 s_4={}&\sum_p\left(\langle P_p,R_p\rangle
             +\tfrac12\|Q_p\|^2
             -\tfrac1{24}\operatorname{Re}\Tr(P_p^4)\right).
\end{split}                                             \tag{98.5}
\]
The coefficient \(s_4\) does not require the fourth-order
link or plaquette logarithm.
\end{theorem}
\begin{proof}
The logarithm of each plaquette is anti-Hermitian.
Its trace has zero real part; traces of its odd powers
also have zero real part.  In the quadratic term of
the exponential, the degree-four contribution is
\(\Tr(P_pR_p)+\Tr(Q_p^2)/2\), using trace cyclicity.
In the fourth power it is \(\Tr(P_p^4)/24\).
The degree-four cubic contribution is purely imaginary:
for example \(P_p^2\) is Hermitian and \(Q_p\) is
anti-Hermitian, so \(\operatorname{Re}\Tr(P_p^2Q_p)=0\).
Subtract the real trace of the exponential from the
representation dimension.  The resulting coefficients
are (98.5).  A fourth-order logarithm could only enter
the linear trace at this order, whose real part is zero.
Smoothness on the local chart gives the stated remainder
at each finite volume.
\end{proof}

The plaquette coefficients in this formula can be
computed by sequentially composing triples with
the following truncated BCH rule:
\[
\begin{split}
 (a,b,c)\star(d,e,f)=\big(&a+d,\ b+e+\tfrac12[a,d],\\
 &c+f+\tfrac12([a,e]+[b,d])
        +\tfrac1{12}([a,[a,d]]+[d,[d,a]])\big).
\end{split}                                             \tag{98.6}
\]
An inverse link negates its entire logarithmic triple.
Equation (98.6) follows by collecting degrees through
three in the BCH formula; it contains no truncated
group multiplication in the verification below.

\begin{record}[Independent exact finite-lattice regression]
\label{rec:v94-regression}
The script
\texttt{scripts/verify\_ym\_v94\_root\_chart\_action\_jets.py}
uses \(SU(2)\) unit quaternions and exact rational
arithmetic.  Its imaginary-coordinate bracket is
\(2\) times the vector cross product; the invariant
norm squared is \(2|v|^2\).
It uses a fine period-four torus, hence sixteen
coarse cells and 64 coarse edges, with the fixed
highest-active-bit parent tree.  Each boundary pivot
has its transverse parity bits zero.

The coarse test field is
\(c_{r,2}=(-1)^{r_1}T/5\), with other orientations
zero, for a quaternion unit \(T\).
The implementation numbers these two orientations
by \(0,1\).  Fourteen explicitly listed nonzero free
fibre coordinates, using all three Lie directions,
are completed by the exact \(E\) pivot formula.
Thus this is a noncommuting test of the joint chart,
not a purely abelian or zero-fibre test.

The script first constructs (98.3).
Independently, it multiplies truncated quaternion
exponential series, evaluates
\(P_0\exp(\overline{\log(P_0^{-1}P_j)})\), and
checks every block logarithm through degree three.
All 64 give precisely \((c_e,0,0)\).
Two coarse edges have nonzero quadratic inverse
corrections, and four have nonzero cubic corrections.

For each of the 1,536 fine plaquettes it then compares
(98.5), computed by the Lie-algebra rule (98.6), with
the real trace of the directly multiplied group series
through degree four.  Every coefficient agrees exactly.
The total coefficients are
\[
 s_2=\frac{119631173}{6350400},\qquad
 s_3=\frac1{360},\qquad
 s_4=-\frac{522279305527811}{1935723847680000}.
                                                        \tag{98.7}
\]
The machine-readable output is
\texttt{artifacts/ym\_v94\_root\_chart\_action\_jets.json}.
The executed script has SHA--256
\texttt{50cbd945ab688b03dde0b2fffae492ecd29325f0d41d3a6a003a460acab69345}.
\end{record}

\begin{assessment}[Scope of the calculation]
The analytic identities (98.3)--(98.6) hold for all
joint directions on the declared chart; the exact
script tests their implementation on one specified
finite-lattice direction and every plaquette of
that lattice.  It is not an exhaustive tensor check.
In particular its fourteen free coordinates are not
the minimizing lift \(-PAc\).  The number \(s_4\) in
(98.7) is a bare-action Taylor coefficient along
that test curve, not the order-one transverse
coefficient in (97.10).  Its negative sign does not
contradict the positive quadratic action.

The next calculation still requires the mixed
polarizations and inverse-Hessian contractions in
(96.5), together with the reduced density and
coarse-Haar terms.  These action and block routines
now provide tested inputs for that calculation.
The mandatory V91 gauge-sector regression must
be applied to the assembled coefficient, not
inferred from the block-only tests above.
\end{assessment}

\begin{firewall}[Position after V94]
The current root chart has an explicit cubic
directional inverse and a Wilson quartic action
implementation tested with exact rational arithmetic.
The minimizing-lift loop contraction, transverse
coefficient value, uniform next-order remainder,
compact completion and full Yang--Mills mass gap
remain unproved.  Only V94 is active; earlier bodies
and archives remain.  Next companion audit: V95.
\end{firewall}

% END APPENDED POST-TENTH-ARCHIVE V94 MATERIAL

% =============================================================================
\section{V95: an exact transverse minimizing lift and its measure contribution}
\label{sec:v95-lift}
% =============================================================================

The V94 predecessor has SHA--256
\texttt{0A420F0791AC35E6AFBC32F9CA1B95E867FAA95002D204C6C6166CF2461CC1D1}.
Its body is retained except for the version label.
We perform the compulsory companion audit below.
The new calculation replaces V94's arbitrary free
fibre direction by the actual linear minimizing lift
for one nonzero-momentum transverse mode.  It also
evaluates the measure part of (96.5), but not its
log-determinant part.

In this section, orientations are numbered \(0,1,2,3\),
as in the scripts.  The chosen pair is the same pair
called \(1,2\) in (97.10), with shifted numbering.
Work on an even-sided periodic coarse lattice with
\(n\) cells.  Write a fine vertex as \(x=2r+b\),
\(b\in\{0,1\}^4\), and put \(\chi_r=(-1)^{r_0}\).
Take \(G=SU(2)\), with imaginary quaternion generator
\(T\) of norm squared \(2\), and coarse direction
\[
 c_{r,1}=\chi_r T,\qquad c_{r,\mu}=0\quad(\mu\ne1).
                                                        \tag{99.1}
\]
This is the transverse mode \(p=(\pi,0,0,0)\);
it is not the small-momentum limit in (97.10).

\begin{theorem}[Closed-form harmonic minimizing lift]
\label{thm:v95-lift}
The first derivative of the fine constrained minimum
in the direction (99.1) is the following tree-slice
fine field:
\[
\begin{split}
 a_{2r+b,0}&=\chi_r\,{\bf1}_{\{b_0=1\}}
                                      (b_1-\tfrac12)T,\\
 a_{2r+b,1}&=\chi_r\,{\bf1}_{\{b_1=1\}}T,\\
 a_{2r+b,2}&=a_{2r+b,3}=0 .
\end{split}                                             \tag{99.2}
\]
It obeys \(La=c\), and is precisely
\(a=E(-H^{-1}Ac)+Rc\) in V92's notation.
The minimized and section quadratic actions are
\[
 [t^2]G(tc)=8n,\qquad
 [t^2]S_0(\exp(tRc))=16n,\qquad
 D^2G(0)[c,c]=16n.                              \tag{99.3}
\]
\end{theorem}
\begin{proof}
All internal links in (99.2) vanish, so the tree
condition holds.  Every boundary link appears
twice in its sixteen-path block average.
In orientation \(0\), the eight boundary values
have four values \(-\chi_rT/2\) and four
\(\chi_rT/2\), so their average is zero.
In orientation \(1\), all eight are \(\chi_rT\);
the remaining orientations vanish.  Thus \(La=c\).

It remains to check minimization, not just the block
constraint.  Remove the common generator \(T\) and
write \(D\) for the scalar oriented plaquette curl.
Direct differencing of (99.2) gives
\[
 (Da)_{2r+b,01}=-\chi_r{\bf1}_{\{b_0=1\}},
 \qquad (Da)_{2r+b,\mu\nu}=0
                         \quad((\mu,\nu)\ne(0,1)).
                                                        \tag{99.4}
\]
For example, when \(b_0=1\) the two orientation-\(0\)
terms contribute \(2\chi_r(b_1-\tfrac12)\);
the two orientation-\(1\) terms contribute
\(-2\chi_rb_1\).  When \(b_0=0\) they cancel.
All other curls vanish because the field is
independent of the other coordinates.

The scalar adjoint curl therefore gives zero
in orientation \(0\) (the two adjacent \(01\)
faces agree), and gives \(\chi_r\) on every
orientation-\(1\) fine edge, whether internal
or boundary.  The other components vanish.
This is exactly \(L^*\lambda\), with
\(\lambda_{r,1}=8\chi_r\) and other components zero:
each boundary coefficient in \(L\) is \(1/8\);
an internal orientation-\(1\) edge receives
\(1/16\) from each of its two adjacent coarse
orientation-\(1\) blocks, whose signs agree.
Restricting to the tree-slice variables preserves
this identity.  Hence
\[
 E^*D^*Da=E^*L^*\lambda=0 .
 \]
The scalar quadratic action is \(\|Da\|^2\)
because \(\|T\|^2=2\).
The fibre quadratic form is strictly positive
by V66, so this constrained stationary field
is its unique minimum.  V81 identifies it with
the linear derivative of the nonlinear minimum.

There are eight nonzero unit-magnitude curls
per cell in (99.4), proving the first equality
in (99.3).  For the section alone, the nonzero
curls have magnitude \(2\) when \(b_0=b_1=1\),
four per cell, giving \(16n\).
Twice the minimized Taylor coefficient is the
Hessian in (99.3).
\end{proof}

\begin{theorem}[Exact measure contribution on the minimizing lift]
\label{thm:v95-density}
Let \(Jc=(-H^{-1}Ac,c)\) be the joint lift and let
\(\ell\) be the reduced density used in (96.5).
Then
\[
 D^2\ell(0)[Jc,Jc]=-7n,\qquad
 D^2\log j_{\rm c}(0)[c,c]=-\tfrac23n .
                                                        \tag{99.5}
\]
Consequently the measure contribution to the
coarse-Haar normalized order-one Hessian is
\[
 -D^2\ell(0)[Jc,Jc]
             +D^2\log j_{\rm c}(0)[c,c]=\tfrac{19}{3}n .
                                                        \tag{99.6}
\]
\end{theorem}
\begin{proof}
The fine lift has no internal components.
Its eight orientation-\(0\) boundary values have
squared scalar magnitudes \(1/4\), and its
eight orientation-\(1\) values have magnitude one.
Thus the sum of squared scalar link magnitudes
is \(10\) per cell.
For our generator,
\(\Tr_{\mathfrak{su}(2)}\operatorname{ad}_T^2=-8\).
The fine-Haar term in (69.5) is therefore
\(-8(10)/12=-20/3\) per cell.

All \(k_e\) in (69.1) vanish, so the squared
second-jet trace is zero.  The third-jet trace
(69.3) reduces, since \(i_x=0\), to
\(-\langle\Tr\operatorname{ad}_{d_x}^2\rangle/6\).
In orientation \(0\), the path mean is zero
and every squared magnitude of \(d_x\) is \(1/4\).
This trace is \(1/3\).
In orientation \(1\), every path value equals
its mean \(\chi_rT\), so \(d_x=0\).
Other orientations also contribute zero.
Subtracting the third-jet trace gives
\(-20/3-1/3=-7\) per cell.

The pullback to joint coordinates has no extra
second-chart term because \(D\ell_{\rm tr}(0)=0\),
as proved in V65.  This establishes the first
identity of (99.5) on \(Jc\), not on a frozen
section.  There is one nonzero coarse link per
cell with generator magnitude one.  Its Haar
Hessian is \(-8/12=-2/3\), proving the second
identity and (99.6).
\end{proof}

\begin{record}[Exact rational independent-coordinate verification]
The script
\texttt{scripts/verify\_ym\_v95\_transverse\_lift\_density.py}
builds the scalar Fourier matrices at
\(p=(\pi,0,0,0)\): a \(96\)-by-\(49\) curl,
a \(49\)-by-\(45\) independent fibre inclusion,
and a \(45\)-by-\(45\) positive Gram matrix.
It verifies \(LE=0\), \(LR=I\), and solves the
stationarity equation with rational elimination.
All 45 unpivoted Schur pivots are positive.
The retained exact vector satisfies both the block
constraint and all 45 stationarity equations.
The resulting field agrees with (99.2).

The density trace is also evaluated by polarizing
V94's cubic root-block polynomial in each of the
three Lie directions; it agrees with (69.3)
edge by edge.  The squared second-jet trace is
independently formed as a \(3\)-by-\(3\) matrix
trace and is zero.  The output values per cell are
\[
 \begin{array}{c|r}
 \text{minimized scalar curl square}&8\\
 \text{section scalar curl square}&16\\
 \text{fine-Haar Hessian}&-20/3\\
 \text{third block-jet trace}&1/3\\
 \text{squared second block-jet trace}&0\\
 \text{reduced-density Hessian}&-7\\
 \text{measure contribution to }D^2\widehat Q&19/3
 \end{array}
 \]
The rational vector, pivot list and trace details
are saved in
\texttt{artifacts/ym\_v95\_transverse\_lift\_density.json}.
The executed script has SHA--256
\texttt{118f302c5c8b313d083393c11d9906ec90ba1aecaccdcfb7c793a1caaed5deeb}.
The general proofs above do not rely solely on
a finite numerical sample.
\end{record}

\begin{audit}[Required V95 companion check]
The active SM note remains V72, SHA--256
\texttt{F44F9745D43379E1672C5550411B58E97195A4C9278592D221DF5D3F644DC1F5}.
Its fixed-time, fixed-mode many-copy covariance
limit does not supply the growing-cutoff or
physical finite-species bounds still required.
The active NS note remains V26, SHA--256
\texttt{82C887F8C2C69E4F28FAD7C3DC8DE5B9099CB5A4BF431EBA1FFF3E91935978AD};
its endpoint provides pointwise rank-one kernel
norms, not full stability or global regularity.
The parallel YM note remains V5, SHA--256
\texttt{306F1BB84CFA8A8D47EA569EC17E9545F9867C8D32B548EC9D9C444B4F3C0512}.
Its extensive-charge obstruction continues to
favor fixed-window or density-based control
instead of a global total-charge tail.
All three hashes agree with the preceding audit;
their endpoint conclusions were reread.
No new companion estimate changes this calculation.
Next companion audit and ten-version review: V100,
also the compulsory archive rollover.
\end{audit}

\begin{assessment}[The remaining term is still essential]
For the direction (99.1), V92 now reads
\[
 \frac1nD^2\widehat Q(0)[c,c]
 =\frac1{2n}\Tr\!\left(PC(Jc,Jc)
                          -PB(Jc)PB(Jc)\right)+\frac{19}{3}.
                                                        \tag{99.7}
\]
The first term is not evaluated in this version.
It can cancel or outweigh the positive measure
term; no sign for the total follows from (99.6).
Although all background link values in (99.2)
commute, the determinant traces include
noncommuting fibre fluctuation directions.
Discarding them would be an abelian substitution,
not an evaluation of the \(SU(2)\) coefficient.
Also, a response at \(p=\pi\) is not the
small-momentum coefficient \(\theta_{12,T}\).
\end{assessment}

\begin{firewall}[Position after V95]
An actual transverse minimizing lift and its
measure contribution are now known exactly,
and the required companion audit is complete.
The log-determinant contraction, full transverse
response, uniform next-order remainder, compact
completion and full Yang--Mills mass gap remain
unproved.  Only V95 is active; earlier bodies
and archives remain.
\end{firewall}

% END APPENDED POST-TENTH-ARCHIVE V95 MATERIAL

% =============================================================================
\section{V96: the neutral-colour determinant and the charged remainder}
\label{sec:v96-neutral}
% =============================================================================

The V95 predecessor has SHA--256
\texttt{D660516336CC617C81D82CD434C1714E29DB8EBB4404A9C6130385F1D859B88D}.
Its body is retained except for the version label.
Retain the \(SU(2)\) transverse direction (99.1).
We compute one colour sector of the determinant
that remains in (99.7), and identify precisely
the sector not included.  The background is
commuting; the full fluctuation space is not.

Let \(D\) be the real scalar fine plaquette curl
on the tree slice, \(E\) the independent scalar
fibre inclusion, and define
\[
 B=DE,\qquad G=B^*B .
                                                        \tag{100.1}
\]
The letter \(B\) in this section is this curl matrix,
not the third-action-derivative matrix in (96.4).
Let \(M\) be the diagonal plaquette projection onto
orientation \(01\) faces with fine parity \(b_0=1\).
It has rank \(8n\) on an \(n\)-cell torus.

\begin{lemma}[Neutral and charged colour splitting at the moving minimum]
\label{lem:v96-colour}
For sufficiently small \(t\), the minimum over \(tc\)
has every Lie coordinate in \(\mathbb RT\).
The real fibre Hessian at that minimum preserves
the neutral subspace \((\mathbb RT)^{45n}\) and
its two-colour orthogonal complement.
Thus its determinant factors as
\[
 \det H_t=\det H_{0,t}\det H_{{\rm ch},t}.
                                                        \tag{100.2}
\]
Here the neutral block has \(45n\) real dimensions
and the charged block has \(90n\).
\end{lemma}
\begin{proof}
Global conjugation by every element of the one-parameter
subgroup generated by \(T\) fixes the coarse data \(tc\).
It preserves the action and maps a small constrained
critical point to one in the same fibre.
For sufficiently small transformations, uniqueness in
V80 therefore fixes the minimum.  Differentiating this
invariance makes every coordinate commute with \(T\).
In \(\mathfrak{su}(2)\) its centralizer is \(\mathbb RT\).

The derivative of the same linear adjoint action
commutes with the fibre Hessian at this invariant
point.  The neutral line is a trivial representation;
its orthogonal plane is a nontrivial rotation
representation and has no invariant linear functional.
Consequently the mixed neutral/charged bilinear block
vanishes.  Orthogonality gives (100.2).  No assertion
that the charged block equals two copies of the
neutral block is made.
\end{proof}

\begin{theorem}[Projection-trace formula for the neutral response]
\label{thm:v96-neutral}
Define the neutral determinant contribution by
\[
 \mathcal D_0=\left.\frac{d^2}{dt^2}
                       \frac12\log\det H_{0,t}\right|_{t=0}.
 \]
On every even-sided torus in V95,
\[
 \mathcal D_0=-\tfrac12\Tr\!\left(G^{-1}B^*MB\right)
             =-\tfrac12\Tr(M\Pi),\qquad
 \Pi=B(B^*B)^{-1}B^*,
 \qquad -4n\leq\mathcal D_0\leq0 .
                                                        \tag{100.3}
\]
\end{theorem}
\begin{proof}
When all fine and coarse coordinates lie in the
same Cartan line, the local block is exactly linear:
the two-link products commute and their local
logarithmic mean is the ordinary mean.  Thus on
this subspace \(X(z,c)=Ez+Rc\).
Let \(\theta_p(t)\) be the scalar plaquette angle
at the actual moving minimum, measured along \(T\).
The neutral Hessian, in scalar fibre coordinates
whose unit is \(T\), is exactly
\[
 H_{0,t}=2B^*\operatorname{diag}(\cos\theta_p(t))B .
                                                        \tag{100.4}
\]
The factor \(2\) comes from
\(2-\operatorname{Re}\Tr e^{sT}=2-2\cos s\).
Changing to orthonormal neutral coordinates only
changes the determinant by a field-independent factor.

V95 gives \(\theta_p(t)=t(Da)_p+O(t^2)\), and
(99.4) gives \((Da)_p^2=M_{pp}\).
Therefore
\[
 H_{0,0}=2G,\qquad H_{0,0}'=0,\qquad
 H_{0,0}''=-2B^*MB .
 \]
The coefficient of \(t^2\) in the angles themselves
does not contribute because the first derivative
of cosine at zero vanishes.
Twice differentiating the log determinant gives
the first expression in (100.3).
Since \(G>0\) by V66, \(\Pi\) is the orthogonal
projection onto the range of \(B\).
Cyclicity of the finite trace gives the second
expression.  Finally \(0\leq\Pi\leq I\) implies
\(0\leq\Tr(M\Pi)\leq\Tr M=8n\), proving the bounds.
\end{proof}

\begin{corollary}[A volume-uniform bound for a partial, not total, response]
The measure-plus-neutral contribution satisfies
\[
 \frac{19}{3}+\frac{\mathcal D_0}{n}
                           \in[\,7/3,\,19/3\,].
                                                        \tag{100.5}
\]
\end{corollary}
\begin{proof}
Combine (99.6) and (100.3).
This leaves the full charged determinant contribution
out of the expression.
\end{proof}

\begin{proposition}[Exact side-two neutral determinant value]
\label{prop:v96-side-two}
On the coarse side-two torus, \(n=16\),
\[
 \Tr(M\Pi)=\frac{689}{12},\qquad
 \frac{\mathcal D_0}{n}=-\frac{689}{384},\qquad
 \frac{19}{3}+\frac{\mathcal D_0}{n}=\frac{581}{128}.
                                                        \tag{100.6}
\]
\end{proposition}
\begin{proof}
Both \(B\) and \(M\) are invariant under cell
translations.  The orthonormal real character
transform of the side-two torus decomposes them
into sixteen sectors with phases
\(\varepsilon_\mu=e^{ip_\mu}\in\{1,-1\}\).
In each sector the curl has 96 rows and 49
non-tree fine columns; the fibre inclusion has
49 rows and 45 columns.  Put \(B_\varepsilon
=D_\varepsilon E_\varepsilon\) and
\(G_\varepsilon=B_\varepsilon^*B_\varepsilon\).
Let \(R_\varepsilon\) be the eight selected rows
of \(B_\varepsilon\).
Then
\[
 \Tr(M\Pi)=\sum_\varepsilon
     \Tr\!\left(R_\varepsilon G_\varepsilon^{-1}
                                      R_\varepsilon^*\right).
                                                        \tag{100.7}
\]
The required exact sector traces are
\[
 \begin{array}{c|r|r}
 (\varepsilon_1,\varepsilon_2,\varepsilon_3)
       &\varepsilon_0=1&\varepsilon_0=-1\\ \hline
 (+,+,+)&15/4&367/105\\
 (+,+,-)&368/105&7/2\\
 (+,-,+)&368/105&7/2\\
 (+,-,-)&10/3&17/5\\
 (-,+,+)&839/210&11/3\\
 (-,+,-)&15/4&18/5\\
 (-,-,+)&15/4&18/5\\
 (-,-,-)&107/30&7/2
 \end{array}
 \]
Here is an exact construction and verification of
these finite rational operations.  For each phase,
shift a parity vertex across a face of the cell
with the appropriate factor \(\varepsilon_\mu\);
use these signs to form every plaquette curl and
the sixteen-path linear block.  The free-coordinate
columns of \(E_\varepsilon\) have their chosen free
entry one and boundary pivot entries
\(-8(L_\varepsilon)_{\mu,\mathrm{free}}\).
This gives \(L_\varepsilon E_\varepsilon=0\)
exactly.  Solve the eight right-hand sides
\[
 G_\varepsilon Y_\varepsilon=R_\varepsilon^*
 \]
by rational Gaussian elimination.  All 45
unpivoted Schur pivots are positive.
Multiply back to check this equation entrywise,
then sum the diagonal of \(R_\varepsilon Y_\varepsilon\).
These operations give the displayed table;
its sum is \(689/12\).
Substitution in (100.3), with \(n=16\), yields
the remaining identities in (100.6).

The implementation of these exact operations is
\texttt{scripts/verify\_ym\_v96\_neutral\_determinant.py}.
It retains every sector's diagonal and pivot list in
\texttt{artifacts/ym\_v96\_neutral\_determinant.json},
checks the solved matrix equations, and uses no
floating-point inversion.  Its executed SHA--256 is
\texttt{d9d265833afa5ecbef5b89dad2328c81fa4238169bed14c9b4b7b84b22c3e7cd}.
\end{proof}

\begin{assessment}[What the colour split leaves unresolved]
Define
\[
 \mathcal D_{\rm ch}
 =\left.\frac{d^2}{dt^2}
                 \frac12\log\det H_{{\rm ch},t}\right|_{0}.
 \]
For the side-two torus the full response is now
\[
 \frac1nD^2\widehat Q(0)[c,c]
                       =\frac{581}{128}+\frac{\mathcal D_{\rm ch}}n .
                                                        \tag{100.8}
\]
The charged term is not computed or sign-controlled
here.  The positive subtotal in (100.6) must not
be reported as a positive full transverse response.
The charged block can have a nonzero first
background derivative, so its log determinant
requires both terms in V92's trace formula.

The next algebraic target is this charged block,
including the nonlinear fibre-coordinate correction.
Even its successful evaluation would yield a
fixed-regulator coefficient at \(p=\pi\), not yet
the small-momentum running coefficient, a uniform
interacting expansion, or a mass-gap theorem.
\end{assessment}

\begin{firewall}[Position after V96]
The neutral determinant has an exact projection
formula, a volume-uniform bound, and a verified
finite-regulator value.  Two charged real colour
directions remain in the determinant.
The full transverse response, uniform next-order
remainder, compact completion and Yang--Mills
mass gap remain unproved.  Only V96 is active;
earlier bodies and archives remain.
Next companion audit and rollover review: V100.
\end{firewall}

% END APPENDED POST-TENTH-ARCHIVE V96 MATERIAL

% =============================================================================
\section{V97: an exact nonlinear background and the charged Hessian input}
\label{sec:v97-background}
% =============================================================================

The V96 predecessor has SHA--256
\texttt{B47BF0D793381D21DDF197186D327728FD887C8EBB5D8FB3B540EF13E729EFE1}.
Its body is retained except for the version label.
The checkerboard transverse field of V95 has a stronger
property than its harmonic minimization: it is an exact
small-amplitude nonlinear constrained minimum.
We prove that fact before forming the charged Hessian.
This does not replace the charged fluctuations by
commuting ones.

Keep \(c,a\) from (99.1)--(99.2).  Let
\(s(x)=S_0(\exp x)\) be the Wilson action in fine
tree-slice logarithmic coordinates, and let \(F(x)\)
be the logarithmic root-reference block.
Adjoints and gradients in the next proof use the
invariant Lie inner product, including \(\|T\|^2=2\).

\begin{theorem}[Exact local constrained minimum along the checkerboard ray]
\label{thm:v97-exact-minimum}
There is a volume-independent \(t_0>0\) such that
for \(|t|<t_0\) the V80 classical local minimum obeys
\[
 X(z_*(tc),tc)=ta,\qquad z_*(tc)=t z_1,\qquad
 z_1=-H^{-1}Ac .
                                                        \tag{101.1}
\]
Its Lagrange multiplier and minimized action are
\[
 \lambda(t)=8\sin(t)c,\qquad
 G(tc)=16n(1-\cos t).
                                                        \tag{101.2}
\]
These are local minimum statements on the uniform
chart, not global minimization over the compact fibre.
\end{theorem}
\begin{proof}
All components of \(a\) lie in \(\mathbb RT\).
On that subspace the local block is exactly \(L\);
hence \(F(ta)=tc\), and the joint coordinates of
\(ta\) are \((tz_1,tc)\).
The fine plaquette angles have coefficients zero
or \(\pm1\), by (99.4).  Thus, in scalar notation,
\(\sin(tDa)=\sin(t)Da\) entrywise.
Differentiation of \(2-2\cos\theta\), with the
factor \(2\) absorbed by the Lie inner product,
gives the neutral part of the fine gradient as
\[
 \nabla s(ta)=\sin(t)D^*Da
             =8\sin(t)L^*c .                    \tag{101.3}
\]
There is no charged part: the action and the
background are invariant under global conjugation
by \(\exp(uT)\), so the charged gradient would have
to be a fixed vector in a nontrivial rotation plane.
Such a vector is zero.

The same equivariance implies that \(DF(ta)\)
maps neutral directions to neutral directions
and charged directions to charged directions.
On the neutral directions it is exactly \(L\).
It follows that
\[
 DF(ta)^*\lambda(t)=L^*(8\sin(t)c)=\nabla s(ta).
 \]
This is the full constrained stationarity equation,
including charged variations.  The field \(ta\)
and its joint coordinates remain in V80's uniqueness
neighbourhood for a volume-independent interval,
since their coordinate max norms are uniformly
bounded by a constant times \(|t|\).
The strictly positive fibre Hessian on that chart
then identifies this critical point as its unique
minimum, proving (101.1).
There are eight nonzero plaquette angles per cell,
each of magnitude \(|t|\), so the action is
\(8n(2-2\cos t)\), proving (101.2).
\end{proof}

\begin{theorem}[Constrained charged Hessian without a nonlinear minimum solve]
\label{thm:v97-charged}
Put \(Z(t)=D_zX(z_*(tc),tc)\).  At the exact
background in (101.1), the fibre Hessian is
\[
 H(t)=Z(t)^*\mathcal L(t)Z(t),\qquad
 \mathcal L(t)[u,v]
 =s''(ta)[u,v]-\langle\lambda(t),F''(ta)[u,v]\rangle .
                                                        \tag{101.4}
\]
This formula can be restricted to the fixed charged
real subspace of V96.
Write \(s_k=D^ks(0)\) and \(F_k=D^kF(0)\);
subscripts denote derivatives, not Taylor coefficients.
Then
\[
 Z(t)=E+tZ_1+t^2Z_2+O(t^3),\quad
 Z_1=-R F_2[a,E\,\cdot],\quad
 Z_2=-\tfrac12R F_3[a,a,E\,\cdot].
                                                        \tag{101.5}
\]
Define symmetric fine bilinear forms
\[
\begin{split}
 L_0&=s_2,\\
 L_1[u,v]&=s_3[a,u,v]-8\langle c,F_2[u,v]\rangle,\\
 L_2[u,v]&=\tfrac12s_4[a,a,u,v]
                            -8\langle c,F_3[a,u,v]\rangle .
\end{split}                                             \tag{101.6}
\]
The Taylor coefficients \(H(t)=H_0+tH_1+t^2H_2+O(t^3)\)
are
\[
\begin{split}
 H_0={}&E^*L_0E,\\
 H_1={}&E^*L_1E+Z_1^*L_0E+E^*L_0Z_1,\\
 H_2={}&E^*L_2E+Z_1^*L_1E+E^*L_1Z_1
           +Z_2^*L_0E+E^*L_0Z_2+Z_1^*L_0Z_1 .
\end{split}                                             \tag{101.7}
\]
\end{theorem}
\begin{proof}
Differentiate \(s(X(z,c))\) twice in \(z\).
The result is \(Z^*s''Z+\langle\nabla s,X_{zz}\rangle\).
Differentiate \(F(X(z,c))=c\) twice, giving
\(DF\,X_{zz}+F''[Z,Z]=0\).
Use the exact stationarity
\(\nabla s=DF^*\lambda\) to replace the second
chain-rule term by \(-\langle\lambda,F''[Z,Z]\rangle\).
This proves (101.4) without discarding the curvature
of the fibre coordinates.

For the tangent, direct differentiation of the
implicit chart gives
\[
 Z(t)=E-R(DF(ta)R)^{-1}DF(ta)E .
 \]
Now \(LE=0\) and \(LR=I\).  In addition,
\(F_2[a,R\,\cdot]=0\): the diagonal block of this
map is \(\operatorname{ad}_{k_e}\) by V65,
and every internal component of \(a\) is zero,
so every \(k_e=0\).
Consequently \(DF(ta)R=I+O(t^2)\), whereas
\[
 DF(ta)E=tF_2[a,E\,\cdot]
                   +\tfrac12t^2F_3[a,a,E\,\cdot]+O(t^3).
 \]
The inverse factor first changes this product at
order \(t^3\), proving (101.5).
Expand (101.4), using
\(\lambda(t)=8tc+O(t^3)\), to obtain (101.6).
Multiplying the three series gives every term
in (101.7).  In particular the terms containing
\(Z_1,Z_2\) and the multiplier are not optional.
\end{proof}

\begin{corollary}[The charged trace to be evaluated]
After restricting (101.7) to the charged subspace
and putting \(P_{\rm ch}=H_0^{-1}\) there,
\[
 \mathcal D_{\rm ch}
 =\Tr(P_{\rm ch}H_2)
        -\tfrac12\Tr(P_{\rm ch}H_1P_{\rm ch}H_1).
                                                        \tag{101.8}
\]
The second term is nonpositive, and is strictly
negative whenever \(H_1\ne0\).
\end{corollary}
\begin{proof}
Differentiate \(\frac12\log\det H(t)\) twice,
remembering that \(H''(0)=2H_2\).
Since \(H_1\) is real symmetric and \(H_0>0\),
\[
 \Tr(P_{\rm ch}H_1P_{\rm ch}H_1)
  =\Tr\big((P_{\rm ch}^{1/2}H_1P_{\rm ch}^{1/2})^2\big)
  \geq0 ,
 \]
with equality only if \(H_1=0\).
\end{proof}

\begin{record}[Exact background and nonzero charged first-jet regression]
The script
\texttt{scripts/verify\_ym\_v97\_exact\_background.py}
checks \(D^*Da=L^*(8c)\) on all 1,024 fine edges
of the fine period-four torus.  It counts 1,408
zero-angle plaquettes and 64 of each sign \(\pm1\),
in agreement with (101.2).

It then chooses two charged free fibre coordinates:
fine edge \(((1,0,0,1),0)\) has colour \(X\),
and \(((2,0,0,1),1)\) has colour \(Y\), where
\(X,Y,T\) are imaginary quaternion units.
All other free entries are zero and the pivot
entries are completed by the exact \(E\) formula.
Denote the resulting free direction by \(w\).
Exact cubic polarization of the composed action
from V94 gives
\[
 \langle w,H_1w\rangle=-3.                      \tag{101.9}
\]
The independent decomposition into the three
terms of \(H_1\) in (101.7) gives raw action
third-jet contribution \(-3\), multiplier contribution
zero, and tangent contribution zero for this
particular test.  The equality is checked rationally.
Thus the charged \(H_1\) is genuinely nonzero,
and its second trace term in (101.8) is strictly
negative on this regulator.

This single test does not check all matrix entries;
in particular its vanishing correction terms do not
justify deleting them for other directions.
The output is
\texttt{artifacts/ym\_v97\_exact\_background.json}.
The executed script has SHA--256
\texttt{2ae4a43e854b6bcc6eaa7675eae0070154eaf511610fb4afa67b2cd48759b12d}.
\end{record}

\begin{assessment}[What was removed from the bottleneck]
The checkerboard background no longer requires any
nonlinear minimizing-map approximation: it is exact
on the uniform small-field interval.
Equations (101.5)--(101.8) now specify the charged
matrix inputs through the already available fine
action fourth jet and block third jet.
The nonzero test (101.9) rules out omitting the
quadratic first-derivative trace.

The full charged traces have not yet been assembled
or evaluated, and their total sign is not determined.
The first term in (101.8) has no sign established here.
Nor does the exact background imply global compact
fibre dominance, a uniform quantum remainder,
continuum construction or a physical mass gap.
\end{assessment}

\begin{firewall}[Position after V97]
The transverse checkerboard ray has an exact local
nonlinear minimum.  Its charged Hessian now has
an explicit constrained expansion and a verified
nonzero first-jet test.  The full charged determinant,
uniform next-order remainder, compact completion
and Yang--Mills mass gap remain unproved.
Only V97 is active; earlier bodies and archives remain.
Next companion audit and rollover review: V100.
\end{firewall}

% END APPENDED POST-TENTH-ARCHIVE V97 MATERIAL

% =============================================================================
\section{V98: closed charged local coefficients and a corrected second-order test}
\label{sec:v98-local-charged}
% =============================================================================

The V97 predecessor has SHA--256
\texttt{EA6BC97B1B65C8F9F1F7EF865BAEFD725C282AA51BA1065CC453FC99AAF3317A}.
Its body is retained except for the version label.
The next matrix assembly needs explicit charged
coefficients, not just names for third and fourth
tensors.  We give those local quadratic forms and
test (101.7) in a direction where the constraint
and tangent corrections are all nonzero.

Use imaginary quaternion coordinates with \(T\)
the third unit and charged vectors in its first
two coordinate directions.  Dots and crosses in
this section are ordinary Euclidean vector
operations; the Lie inner product is twice the dot,
and the Lie bracket is twice the cross.

\begin{theorem}[Charged coefficients of the cubic root block]
\label{thm:v98-block}
Consider one coarse edge with its sixteen ordered
two-link paths.  In a Cartan background \(a\) and
charged direction \(w\), write their values on
path \(j\) as
\[
 (a_{l_j},a_{k_j})=(\alpha_jT,\beta_jT),
 \qquad (w_{l_j},w_{k_j})=(u_j,v_j).
 \]
Put \(p_j=\alpha_j+\beta_j\), \(q_j=u_j+v_j\),
\(\delta p_j=p_j-\overline p\),
\(\delta q_j=q_j-\overline q\), with index zero
the actual root reference.
The degree-three terms of \(F(ta+sw)\) with
degrees \((2,1)\) and \((1,2)\) are
\[
\begin{split}
 C_{21}=\tfrac13\Big\langle
   &(\alpha_j\beta_j-\beta_j^2)u_j
    +(\alpha_j\beta_j-\alpha_j^2)v_j\\
   &+(\delta p_j)^2q_0
                      -\delta p_jp_0\delta q_j\Big\rangle,
\end{split}                                             \tag{102.1}
\]
\[
\begin{split}
 C_{12}=\tfrac13\Big\langle
   &(\alpha_j+\beta_j)u_j\cdot v_j
          -\beta_j|u_j|^2-\alpha_j|v_j|^2\\
   &+p_0|\delta q_j|^2
                    -\delta p_j\,\delta q_j\cdot q_0\Big\rangle T .
\end{split}                                             \tag{102.2}
\]
Here angle brackets denote the sixteen-path average.
In derivative notation,
\(F_3[a,a,w]=2C_{21}\) and \(F_3[a,w,w]=2C_{12}\).
For the V97 background, in particular,
\(Z_2w=-RC_{21}\).
\end{theorem}
\begin{proof}
Use the cubic root-block expression (98.1) and
the quaternion identity
\[
 [x,[y,z]]=4\big(y(x\cdot z)-z(x\cdot y)\big).
 \]
In a nested bracket with Cartan first-order
coefficients \(a_x,a_y,a_z\) and charged
coefficients \(w_x,w_y,w_z\), the \(t^2s\)
coefficient is
\(4a_x(a_z w_y-a_y w_z)\).
Its \(ts^2\) coefficient along \(T\) is
\(4(a_y w_x\cdot w_z-a_z w_x\cdot w_y)T\).
Apply these two identities to the two ordered
BCH brackets and the reference bracket in
(98.1).  The factor \(4/12\) gives the displayed
averages.  The factors of two relating the
coefficients to \(F_3\) follow from the
multinomial expansion of \(F_3[x,x,x]/6\).
Equation (101.5) then gives the tangent assertion.
\end{proof}

\begin{theorem}[Raw Wilson charged Hessian coefficients in local form]
\label{thm:v98-raw}
On a plaquette, write its four signed links in
chronological order as \((\sigma_j,e_j)\),
\(\sigma_j\in\{1,-1\}\), \(j=0,1,2,3\).
For a Cartan background \(a_e=a_e^{\rm sc}T\)
and a charged field \(w\), set
\[
 A_j=\sigma_j a_{e_j}^{\rm sc},\quad v_j=w_{e_j},
 \quad \theta=\sum_j A_j,\quad
 \gamma_{jk}=\theta-A_j-A_k-2\sum_{j<h<k}A_h .
 \]
Let \(\mathcal R_1(w)=s_3[a,w,w]\) and
\(\mathcal R_2(w)=s_4[a,a,w,w]/2\).
Their contributions from this plaquette are
\[
 \mathcal R_{1,p}(w)=
 4\sum_{j<k}\sigma_j\sigma_k\gamma_{jk}
                               (v_j\times v_k)\cdot T ,
                                                        \tag{102.3}
\]
\[
\begin{split}
 \mathcal R_{2,p}(w)={}&
 -\sum_j\left((\theta-A_j)^2+\frac{A_j^2}{3}
                     +\frac{2A_j(\theta-A_j)}3\right)|v_j|^2\\
 &-4\sum_{j<k}\sigma_j\sigma_k
       \left(\frac{\gamma_{jk}^2}{2}
                       +\frac{A_j^2+A_k^2}{6}\right)v_j\cdot v_k .
\end{split}                                             \tag{102.4}
\]
Sum these expressions over plaquettes to obtain
the full raw coefficients.  Bilinear coefficients
are recovered by polarization.
\end{theorem}
\begin{proof}
For one signed link, its logarithm is
\(tA_jT+s\sigma_jv_j\).
The term linear in \(s\) in its exponential is
\[
 s\sigma_j(1-t^2A_j^2/6)v_j+O(t^3s).
 \]
The terms quadratic in \(s\) have scalar part
\(s^2(-1/2+t^2A_j^2/12)|v_j|^2\)
and \(T\) part
\(-ts^2A_j|v_j|^2T/6\), through the needed orders.
These follow by expanding the quaternion cosine
and sine-over-radius functions.

For two distinct linear insertions, cycle the
neutral factors through the scalar trace and use
that a charged imaginary vector anticommutes with
\(T\).  The remaining neutral angle is
\(t\gamma_{jk}\).  The scalar part of the
product of charged vectors is \(-v_j\cdot v_k\)
and its \(T\) part is \((v_j\times v_k)\cdot T\).
After taking minus twice the scalar part, their
action contribution is
\[
 2s^2\sigma_j\sigma_k
  \left[v_j\cdot v_k+t\gamma_{jk}(v_j\times v_k)\cdot T
   -t^2\left(\frac{\gamma_{jk}^2}{2}
               +\frac{A_j^2+A_k^2}{6}\right)v_j\cdot v_k\right].
 \]
For a quadratic insertion on one link, multiplying
by the remaining neutral angle \(t(\theta-A_j)\)
gives the diagonal term in (102.4).
Finally, the coefficients of \(ts^2\) and \(t^2s^2\)
in the action are respectively \(\mathcal R_1/2\)
and \(\mathcal R_2/2\).  This proves all factors
and signs in (102.3)--(102.4).
\end{proof}

\begin{corollary}[Scalar test of the full charged second Hessian coefficient]
For a free charged direction \(w\), put
\(b=Ew\), \(u_1=Z_1w\), \(u_2=Z_2w\), and use
\(L_0,L_1\) from (101.6).  Then
\[
\begin{split}
 \langle w,H_2w\rangle={}&
 \mathcal R_2(b)-16\sum_e\langle c_e,C_{12,e}(a,b)\rangle\\
 &+2L_1[b,u_1]+2L_0[b,u_2]+L_0[u_1,u_1].
\end{split}                                             \tag{102.5}
\]
\end{corollary}
\begin{proof}
Insert (102.1)--(102.2) into (101.6)--(101.7).
The first line is \(L_2[b,b]\) because
\(F_3[a,b,b]=2C_{12}(a,b)\).
All pairings in this formula are Lie pairings,
not the unscaled Euclidean dots in (102.1)--(102.4).
\end{proof}

\begin{record}[Exact corrected charged regression]
The script
\texttt{scripts/verify\_ym\_v98\_charged\_second\_jet.py}
uses the V97 background on the fine period-four
torus.  Its four charged free entries are
\[
 \begin{array}{c|c}
 \text{fine edge}&\text{value}\\ \hline
 ((1,0,0,1),0)&X\\
 ((2,0,0,1),1)&Y\\
 ((0,0,0,1),0)&X/3\\
 ((1,1,0,1),1)&2Y/5
 \end{array}
 \]
and all remaining free entries are zero.
The pivot entries are completed using \(E\).
It checks (102.1)--(102.2) on all 64 coarse
edges by exact polarization of the unreduced
cubic polynomial from V94.
The raw action formulas are checked independently
by cubic and quartic action polarization.

The five terms of (102.5), in order, are
\[
 -\frac{5203}{450},\qquad
 -\frac{82}{75},\qquad
 -\frac14,\qquad
 -\frac{23}{48},\qquad
 \frac{37}{96}.
 \]
They sum to
\[
 \langle w,H_2w\rangle=-\frac{18719}{1440}.
                                                        \tag{102.6}
\]
For a separate check of the composition, let
\(\Phi_4(\xi)=[s^4]f(s\xi)\), computed by
V94's cubic chart inverse and plaquette-log formula.
With the joint background direction \(\xi\) whose
fine first lift is \(a\), exact arithmetic verifies
\[
 \Phi_4(\xi+w)+\Phi_4(\xi-w)
                  -2\Phi_4(\xi)-2\Phi_4(w)
                    =-\frac{18719}{1440}.
 \]
This is the required \(H_2\) quadratic value:
the polarization equals \(f_4[\xi,\xi,w,w]/2\).
Both first and second tangent corrections are
nonzero on 32 fine links.  Unlike the V97 first
test, every correction term displayed in (102.5)
is nonzero in this regression.

The output is
\texttt{artifacts/ym\_v98\_charged\_second\_jet.json}.
The executed script has SHA--256
\texttt{b912f9ded28a66f9291308766e29215ee16c4537468c5e1e5320300b1fdd98f0}.
This verifies one corrected quadratic test and
the associated local formulas, not all entries
of a completed determinant matrix.
\end{record}

\begin{assessment}[Ready inputs, but not a completed determinant]
The charged block now has local dot/cross formulas
for the needed block jets and raw action Hessian
coefficients.  Polarization yields the matrix
entries for (101.7), and the nonzero correction
test (102.6) checks that assembly's factor conventions.
The next step is to assemble and contract those
matrices with the gapped fibre inverse, retaining
both terms of (101.8).
No determinant value, total response sign,
small-momentum limit, or mass gap follows from
the negative quadratic test alone.
\end{assessment}

\begin{firewall}[Position after V98]
Explicit charged local coefficients and a
nontrivially corrected second-order regression
are complete.  The full charged determinant,
uniform next-order remainder, compact completion
and full Yang--Mills mass gap remain unproved.
Only V98 is active; earlier bodies and archives remain.
Next companion audit, ten-version review and
compulsory rollover: V100.
\end{firewall}

% END APPENDED POST-TENTH-ARCHIVE V98 MATERIAL

% =============================================================================
\section{V99: the complete finite-torus charged trace and the gauge cancellation}
\label{sec:v99-charged-trace}
% =============================================================================

The V98 predecessor has SHA--256
\texttt{638CA93DD249EB96099AEECCB2FCE85F3CC982DC5DE5E4357902B08EC1E3C012}.
Its body is retained except for the version label.
We now complete the charged determinant on the
side-two coarse torus for the transverse direction
of V95.  The same assembly is tested on a separate
pure-gauge background.  This is an exact
finite-regulator result at external momentum
\((\pi,0,0,0)\), not a small-momentum or continuum
running-coupling calculation.

There are \(n=16\) coarse cells.  Let
\(\varepsilon\in\{1,-1\}^4\) label their real
Fourier characters and put
\(\varepsilon^\sharp=(-\varepsilon_0,
\varepsilon_1,\varepsilon_2,\varepsilon_3)\).
Identify the two charged real colours with one
complex coordinate \(X+iY\).

\begin{lemma}[Momentum-transfer form and the real-colour factor]
\label{lem:v99-transfer}
The complex charged Hessian coefficients in V97
take the form
\[
 H^{\mathbb C}_0(\varepsilon,\varepsilon)=M_\varepsilon,
 \quad H^{\mathbb C}_1(\varepsilon^\sharp,\varepsilon)
       =iK_\varepsilon,\quad
 H^{\mathbb C}_2(\varepsilon,\varepsilon)=A_\varepsilon,
                                                        \tag{103.1}
\]
with other momentum blocks zero.  Here
\(M_\varepsilon=2(DE)_\varepsilon^*(DE)_\varepsilon>0\),
\(A_\varepsilon=A_\varepsilon^T\), and
\[
 K_{\varepsilon^\sharp}=-K_\varepsilon^T.
 \]
All displayed matrices are real rational
\(45\)-by-\(45\) matrices.  With
\(P_\varepsilon=M_\varepsilon^{-1}\), the charged
real determinant contribution is
\[
 \frac{\mathcal D_{\rm ch}}n
 =\frac1n\sum_\varepsilon
 \left(2\Tr(P_\varepsilon A_\varepsilon)
 -\Tr(P_{\varepsilon^\sharp}K_\varepsilon
                            P_\varepsilon K_\varepsilon^T)\right).
                                                        \tag{103.2}
\]
\end{lemma}
\begin{proof}
Each first-order coefficient in the background
contains one factor \(\chi_r=(-1)^{r_0}\), so it
transfers momentum by \((\pi,0,0,0)\).
Each second-order coefficient contains two
such factors and is invariant under cell translations.
The charged dot/cross formulas of V98 give a
purely imaginary first-order complex matrix
and a real second-order one.  Hermiticity gives
the stated transpose identities.
All phases here are signs, and every local
Taylor and fibre-coordinate coefficient is rational.

For a positive Hermitian complex matrix, its
realification has determinant equal to the square
of its complex determinant.  Therefore
\(\frac12\log\det H_{\rm ch}^{\mathbb R}
=\log\det H_{\rm ch}^{\mathbb C}\).
Twice differentiating the latter gives
\(2\Tr(PH_2)-\Tr(PH_1PH_1)\).
Insert (103.1), using
\(iK_{\varepsilon^\sharp}=-iK_\varepsilon^T\),
to obtain (103.2).  The sum already includes
both members of every momentum-transfer pair;
there is no additional pair or colour factor.
\end{proof}

\begin{record}[Explicit matrix assembly convention]
The real first-order fine complex coefficient
is denoted \(U_\varepsilon\), so the fine
Lagrangian Hessian has first block
\(iU_\varepsilon\) from \(\varepsilon\) to
\(\varepsilon^\sharp\).  Its second-order
diagonal coefficient is \(V_\varepsilon\);
its zeroth-order one is \(L_{0,\varepsilon}\).
These include the multiplier terms in (101.6).
Write the tangent coefficients as
\(Z_1(\varepsilon^\sharp,\varepsilon)=iW_{1,\varepsilon}\)
and \(Z_2(\varepsilon,\varepsilon)=W_{2,\varepsilon}\).
Then the implementation uses
\[
\begin{split}
 K_\varepsilon={}&
 E_{\varepsilon^\sharp}^TU_\varepsilon E_\varepsilon
 +E_{\varepsilon^\sharp}^TL_{0,\varepsilon^\sharp}W_{1,\varepsilon}\\
 &-W_{1,\varepsilon^\sharp}^TL_{0,\varepsilon}E_\varepsilon ,
\end{split}                                             \tag{103.3}
\]
\[
\begin{split}
 A_\varepsilon={}&E_\varepsilon^TV_\varepsilon E_\varepsilon
 +W_{1,\varepsilon}^TU_\varepsilon E_\varepsilon
 +E_\varepsilon^TU_\varepsilon^TW_{1,\varepsilon}\\
 &+W_{2,\varepsilon}^TL_{0,\varepsilon}E_\varepsilon
 +E_\varepsilon^TL_{0,\varepsilon}W_{2,\varepsilon}
 +W_{1,\varepsilon}^TL_{0,\varepsilon^\sharp}W_{1,\varepsilon}.
\end{split}                                             \tag{103.4}
\]
They follow by multiplying the three series in
(101.4); in particular the minus sign in
(103.3) comes from the adjoint of \(iW_1\).

The local fine coefficients are assembled from
(102.1)--(102.4), with the constraint terms of
(101.6).  A fine edge at \(2r+b\) in an input
mode \(\varepsilon\) contributes its cell
character \(\prod_\mu\varepsilon_\mu^{r_\mu}\)
to the coordinate labelled by \((b,\mu)\).
Thus links crossing the parity cell are not
treated as if they were in the source cell.
For first-order bilinear entries the output
character is \(\varepsilon^\sharp\), not
\(\varepsilon\).  The scalar fibre matrices
are the same \(E_\varepsilon\) as in V96.
\end{record}

\begin{theorem}[Exact full transverse coefficient on the side-two torus]
\label{thm:v99-value}
For the transverse direction \(c\) in (99.1),
\[
\begin{split}
 \frac1n\sum_\varepsilon2\Tr(P_\varepsilon A_\varepsilon)
                 &=\frac{79227}{4480},\\
 \frac1n\sum_\varepsilon
   \Tr(P_{\varepsilon^\sharp}K_\varepsilon
                            P_\varepsilon K_\varepsilon^T)
                 &=\frac{166901}{5760},\\
 \frac{\mathcal D_{\rm ch}}n&=-\frac{14227}{1260}.
\end{split}                                             \tag{103.5}
\]
Including the neutral determinant and coarse-Haar
normalized measure terms gives
\[
 \boxed{\frac1nD^2\widehat Q(0)[c,c]
             =-\frac{272249}{40320}.}           \tag{103.6}
\]
\end{theorem}
\begin{proof}[Exact finite-dimensional computation]
For each of the sixteen phases, form \(M_\varepsilon\)
and solve all 45 identity right-hand sides using
rational elimination.  All elimination pivots
are positive, and multiplication verifies
\(M_\varepsilon P_\varepsilon=I\) entrywise.
Assemble (103.3)--(103.4).
The computation checks
\(U_{\varepsilon^\sharp}=-U_\varepsilon^T\),
\(K_{\varepsilon^\sharp}=-K_\varepsilon^T\)
and \(A_\varepsilon=A_\varepsilon^T\) exactly.

As independent assembly regressions, transform
the V97 two-entry free field into the orthonormal
character basis, whose factor is \(1/4\).
Its first-order quadratic value is \(-3\), as
in (101.9).  The V98 four-entry field gives
second-order value \(-18719/1440\), as in (102.6).
Thus the assembled matrices reproduce the
earlier physical-space polarizations, including
the nonzero second-order chart corrections.

Evaluate the traces in (103.2) with rational
matrix arithmetic.  Their sums divided by
sixteen are the first two numbers in (103.5);
their difference is the third.
Finally add the neutral contribution
\(-689/384\) from V96 and the measure contribution
\(19/3\) from V95.  Their sum is (103.6).
The script and retained sector traces are
specified below; this calculation uses no
floating-point inverse, finite differences,
or fitted coefficients.
\end{proof}

\begin{proposition}[Full normalization regression on a pure-gauge mode]
\label{prop:v99-gauge}
Use instead
\[
 c^{\rm g}_{r,0}=2\chi_r T,\qquad
 c^{\rm g}_{r,\mu}=0\quad(\mu\ne0),\qquad
 a^{\rm g}=Rc^{\rm g}.
 \]
The same assembly gives, per cell,
\[
 \begin{array}{c|r}
 \text{charged second-order trace}&38597/420\\
 \text{charged first-order square trace}&15479/140\\
 \text{charged determinant difference}&-56/3\\
 \text{neutral determinant}&0\\
 \text{coarse-Haar normalized measure term}&56/3\\ \hline
 \text{full response}&0
 \end{array}                                            \tag{103.7}
\]
\end{proposition}
\begin{proof}
This is the coarse gradient of
\(\omega_r=\chi_rT\).  Its fine field is the
corresponding cell-constant gauge gradient,
and its exponentiated links are a flat gauge
transform of the identity for small \(t\).
Hence the classical multiplier is zero.
The same local coefficient assembly is run
with this background and zero multiplier,
without replacing any charged matrix by zero.
It gives the first three entries of the table.

The neutral determinant is constant because
every plaquette angle is zero.
There are eight nonzero fine boundary values
of scalar magnitude two per cell, so their
Haar Hessian is \(-64/3\).
All path values within each block agree;
the third block-jet trace and squared
second-jet trace are zero by V65.
The coarse-Haar Hessian is \(-8/3\).
Thus the measure contribution is
\(64/3-8/3=56/3\).
The total vanishes exactly, agreeing with
V91's Ward identity.  This is a regression
of the full normalized coefficient, not
only of its neutral or bare-action part.
\end{proof}

\begin{record}[Reproducible exact trace record]
The executed verifier is
\texttt{scripts/verify\_ym\_v99\_charged\_determinant.py},
SHA--256
\texttt{225bbeb70c349993c33881133b804a4735fef50e19222b4624093cdc61750a22}.
Its output is
\texttt{artifacts/ym\_v99\_charged\_determinant.json},
SHA--256
\texttt{7F406068B15CF6EF9378B074D3437D0AF28939CF621AACD297359B163F403848}.
The output retains all sixteen sector traces
for both backgrounds and the hashes of the
three imported local verifier sources.
All checks passed on the recorded run.
The charged real dimension is 1,440; its
complex representation has sixteen blocks
of dimension 45.
\end{record}

\begin{corollary}[Fixed-regulator local integral interpretation]
Let \(\widehat W_\beta=W_\beta+\log j_{\rm c}\)
be the local-box effective action relative to
coarse Haar.  On this fixed side-two torus,
\[
 \frac1nD^2\widehat W_\beta(0)[c,c]
   =16\beta-\frac{272249}{40320}+O(\beta^{-1})
                 \quad(\beta\longrightarrow\infty).
                                                        \tag{103.8}
\]
\end{corollary}
\begin{proof}
Use V90's fixed-volume \(C^2\) Laplace expansion,
V95's \(D^2G(0)[c,c]=16n\), and (103.6).
This does not assert a volume-uniform constant
for the remainder.
\end{proof}

\begin{assessment}[What the negative coefficient does and does not establish]
The previously unevaluated finite-regulator
transverse coefficient is now complete,
including the moving tangent, constraint
curvature, both charged trace terms,
the neutral determinant and the measure.
The pure-gauge cancellation is a substantial
normalization check.

The sign in (103.6) is not by itself a
universal beta-function sign or a physical
mass statement.  It is at momentum \(\pi\),
on a specified finite lattice and blocking.
At large \(\beta\), its negative order-one
value coexists with the positive leading
term \(16\beta\) in (103.8).
The small-momentum coefficient (97.10),
uniform interacting remainder, compact-fibre
comparison and continuum mass gap are not
resolved by this calculation.
\end{assessment}

\begin{decision}[V100 review and rollover]
At V100, perform the compulsory companion
audit and ten-version review, then freeze
this lineage with the next archive suffix
and restart a concise V1 with the proved
interfaces and remaining bottlenecks.
The next scale-sensitive target is the
small-momentum response, not another
interpretation of the fixed \(\pi\)-mode value.
Retain the uniform remainder and compact
completion as explicit independent tasks.
\end{decision}

\begin{firewall}[Position after V99]
The full order-one transverse response is
evaluated exactly on the side-two regulator,
and the full pure-gauge regression vanishes.
The small-momentum response, uniform
next-order remainder, compact completion
and full Yang--Mills mass gap remain unproved.
Only V99 is active; earlier bodies and
archives remain.  V100 audit and rollover
are due next.
\end{firewall}

% END APPENDED POST-TENTH-ARCHIVE V99 MATERIAL

% =============================================================================
\section{V100: companion audit, ten-version review, and eleventh archive}
\label{sec:v100-rollover}
% =============================================================================

The V99 predecessor has SHA--256
\texttt{C86B6C0D1F37BD88B6E5BD8657C5BAADA1C38A5331285046B12957E03B10E77C}.
Its body is retained except for the version label.
This is the compulsory V100 audit and rollover.
The complete current lineage is frozen as
\texttt{Renewal\_Yang\_Mills\_Mass\_Gap\_Research\_Notes\_V100\_f11.tex}.
The next active file restarts at V1, with context,
a proof-indexed summary, and a new small-momentum
acquisition result.  The archive is retained.

\begin{audit}[Companion state at V100]
The active SM note is still V72, SHA--256
\texttt{F44F9745D43379E1672C5550411B58E97195A4C9278592D221DF5D3F644DC1F5}.
Its many-copy fixed-mode covariance limit does
not provide a physical ultraviolet completion.
NS remains V26, SHA--256
\texttt{82C887F8C2C69E4F28FAD7C3DC8DE5B9099CB5A4BF431EBA1FFF3E91935978AD},
with pointwise rank-one kernel norms, not global
regularity.  Parallel YM remains V5, SHA--256
\texttt{306F1BB84CFA8A8D47EA569EC17E9545F9867C8D32B548EC9D9C444B4F3C0512}.
Its extensive-total-charge obstruction remains
relevant to the choice of infinite-volume topology
estimates, but supplies no missing charged or
continuum estimate.  The hashes agree with V95.
No companion update changes the direction.
\end{audit}

\begin{record}[Ten-version review: V91--V100]
V91 corrected the order-one coefficient to coarse
Haar and proved its exact gauge-sector benchmark.
V92 removed the second-minimizer-derivative term
at the identity and proved volume-uniform locality
for the complete fourth-jet coefficient.
V93 constructed its thermodynamic limit and
zero-momentum Ward identity, without constructing
the full thermodynamic quantum theory.

V94 provided exact local block/action jet tests.
V95 found the actual transverse harmonic lift
and evaluated its measure term.
V96 evaluated the neutral determinant.
V97 proved that this special ray has an exact
nonlinear local minimum and exposed the nonzero
charged first-derivative term.
V98 supplied closed local charged coefficients
with a test retaining nonzero chart corrections.
V99 completed both charged traces, obtained the
full fixed-regulator response
\(-272249/40320\) per cell, and passed a full
pure-gauge cancellation check.

The major new endpoint is therefore a complete
specified finite-regulator coefficient, not just
an unassembled tensor formula.  Its momentum is
\(\pi\), and its lattice has only two coarse
sites per direction.  It is not the infrared
running coefficient.
\end{record}

\begin{assessment}[Remaining bottlenecks after the successful finite calculation]
Three gaps must not be conflated:
\begin{enumerate}
\item The small-momentum coefficient in the
thermodynamic limit still needs evaluation and
calibration to the chosen scale evolution.
\item The order-one coefficient has a uniform
volume limit, but the exact interacting local
integral has not been given a uniform next-order
response remainder.
\item The local chart has not been shown to
dominate the full compact fibre in the curved
retained data required by iteration.
\end{enumerate}
Beyond them, the continuum observable state,
required positivity/reconstruction, and a
physical spectral gap remain unproved.
An additional decimal place in the V99 corner
coefficient would not resolve any of these gaps.
\end{assessment}

\begin{decision}[Restart direction]
Carry the precise local chart, coefficient
normalization, exact test values, and proof
boundaries into a compact V1.
Give a finite-volume estimator for the
small-momentum second-order coefficient with
both volume error and arithmetic error explicit.
This turns the next target into a specified
calculation and prevents the corner sign from
being substituted for an infrared conclusion.
The next companion audit is new V5; the next
ten-version review is new V10.
\end{decision}

\begin{firewall}[Archive f11 closure status]
This archive records progress toward, not a
solution of, the full Yang--Mills mass-gap
problem.  It is frozen after V100.  The
programme continues in the next active V1;
all older archives remain preserved.
\end{firewall}

% END APPENDED POST-TENTH-ARCHIVE V100 MATERIAL

\end{document}
